\documentclass{article}
\usepackage{readmasters}
\begin{document}

\origpage{i}

\section*{EXERCICES DE CALCUL INTÉGRAL SUR DIVERS ORDRES DE TRANSCENDANTES ET SUR LES QUADRATURES;}

Par A. M. LEGENDRE, Chevalier de l'Empire, Membre de l'Institut et de la Légion d'honneur, Conseiller-Titulaire de l'Université impériale, etc.

PARIS,

M\textsuperscript{me} V\textsuperscript{e} COURCIER, IMPRIMEUR-LIBRAIRE POUR LES MATHÉMATIQUES.

1811.

\origpage{1}

\section*{EXERCICES DE CALCUL INTÉGRAL.}

\section*{PREMIÈRE PARTIE.}

\section*{DES FONCTIONS ELLIPTIQUES.}

Après avoir épuisé les formules différentielles qui s'intègrent tant algébriquement que par arcs de cercle ou par logarithmes, les Géomètres s'occupèrent de rechercher toutes celles qui sont intégrables par les arcs d'ellipse ou par les arcs d'hyperbole${}^{(1)}$. On avait lieu de croire que ces transcendantes tenaient le premier rang après les fonctions circulaires et logarithmiques, et il importait au progrès des nouveaux calculs, de ramener à un point de difficulté bien connu, toutes les intégrales qui étaient susceptibles de cette réduction.

Les formules qu'on peut intégrer par cette voie se trouvèrent très-nombreuses; mais il n'y avait point de liaison entre les résultats, et ils étaient loin de former une théorie.

Un Géomètre italien, d'une grande sagacité, ouvrit la route à des spéculations plus profondes${}^{(2)}$. Il prouva que sur toute ellipse ou

\textbf{(1)} \emph{Maclaurin.} Traité des Fluxions.~--- \emph{D'Alembert.} Mém.\ de Berlin. 1746.

\textbf{(2)} \emph{Fagnani.} Produzioni matematiche, Tom.\ II.

\origpage{2}
sur toute hyperbole donnée, on peut assigner, d'une infinité de manières, deux arcs dont la différence soit égale à une quantité algébrique. Il démontra en même temps que la \emph{lemniscate} jouit de cette singulière propriété, que ses arcs peuvent être multipliés ou divisés algébriquement, comme les arcs de cercle, quoique chacun d'eux soit une transcendante d'un ordre supérieur.

Euler, par une combinaison qu'on peut regarder comme fort heureuse, quoique ces hasards n'arrivent jamais qu'à ceux qui savent les faire naître, trouva l'intégrale algébrique complète d'une équation différentielle composée de deux termes séparés, mais semblables, dont chacun n'est intégrable que par des arcs de sections coniques${}^{(1)}$.

Cette découverte importante donna lieu à son auteur de comparer d'une manière plus générale qu'on ne l'avait fait avant lui, non-seulement les arcs d'une même ellipse ou d'une même hyperbole, mais en général toutes les transcendantes comprises dans la formule $\displaystyle\int\frac{P\,dx}{R}$, où $P$ est une fonction rationnelle de $x$, et $R$ un radical de la forme $\sqrt{(\alpha+\beta x+\gamma x^{2}+\delta x^{3}+\varepsilon x^{4})}$, $\alpha$, $\beta$, $\gamma$, $\delta$, $\varepsilon$, étant constans.

L'intégrale trouvée par Euler était trop remarquable pour ne pas fixer particulièrement l'attention des Géomètres. Lagrange voulut faire rentrer cette intégration dans les procédés ordinaires de l'analyse; il y réussit par une méthode fort ingénieuse${}^{(2)}$, dont l'application s'élève graduellement des transcendantes inférieures aux transcendantes Eulériennes; mais il essaia inutilement de parvenir à un résultat plus général que celui d'Euler.

Peu de temps après, Landen, géomètre anglais, démontra que tout arc d'hyperbole peut être mesuré par deux arcs d'ellipse${}^{(3)}$; découverte mémorable qui réduit aux seuls arcs d'ellipse toutes les intégrales qu'on n'avait pu exprimer jusques-là que par la rectification des deux courbes.

Enfin Lagrange se signala de nouveau dans la même carrière,

\textbf{(1)} \emph{Euler.} Novi Com.\ Petrop.\ Tom.\ VI et VII.

\textbf{(2)} Mém.\ de Turin, Tom.\ IV.

\textbf{(3)} Mathematical Memoirs, by John Landen, 1780.

\origpage{3}
en donnant une méthode générale${}^{(1)}$ pour ramener, par des transformations successives, l'intégrale $\displaystyle\int\frac{P\,dx}{R}$ à l'intégrale d'une formule semblable qui, par la disposition de ses coefficiens, est facile à évaluer par approximation. Ces transformations ont le double but de servir à la comparaison d'une suite de transcendantes formées d'après la même loi, et de conduire aux approximations les plus rapides dont ces fonctions sont susceptibles.

Telles étaient les principales découvertes des Géomètres dans la théorie des transcendantes désignées par $\displaystyle\int\frac{P\,dx}{R}$, lorsque je publiai mes Recherches sur l'intégration par arcs d'ellipse${}^{(2)}$. La première partie avait été composée avant que j'eusse connaissance du théorème de Landen; elle contenait des vues nouvelles sur l'usage des arcs d'ellipse, et particulièrement un moyen d'éviter l'emploi des arcs d'hyperbole dans le calcul intégral, en y suppléant par une table d'arcs d'ellipse dressée convenablement. Je donnai ensuite une nouvelle démonstration du théorème de Landen, et je prouvai par la même méthode, que toute ellipse donnée fait partie d'une suite infinie d'ellipses tellement liées entre elles, que par la rectification de deux de ces ellipses, prises à volonté, on obtient la rectification de toutes les autres. Ces ellipses ayant un demi-grand axe commun égal à l'unité, et leurs excentricités variant suivant une loi connue, depuis zéro jusqu'à l'unité, on peut par ce théorème réduire la rectification d'une ellipse donnée à celle de deux autres ellipses aussi peu différentes du cercle qu'on voudra. C'était un pas de plus dans une carrière difficile.

Mais cette matière, et en général la théorie des transcendantes désignée par $\displaystyle\int\frac{P\,dx}{R}$, demandait à être traitée d'une manière plus méthodique et plus approfondie. C'est ce que j'essayai de faire dans mon \emph{Mémoire sur les Transcendantes elliptiques}, publié en 1793. Je me proposai dans cet ouvrage, de comparer entre elles toutes les fonctions comprises sous cette dénomination, de les classer en différentes espèces, de réduire chacune à la forme la plus simple

\textbf{(1)} Nouveaux Mémoires de Turin, ann.\ 1784 et 1785, Tom.\ II.

\textbf{(2)} Mém.\ de l'Acad.\ des Sciences de Paris, ann.\ 1786.

\origpage{4}
dont elle est susceptible, de les évaluer par les approximations les plus promptes et les plus faciles; enfin, de former de l'ensemble de cette théorie une sorte d'algorithme qui pût contribuer à étendre le domaine de l'analyse.

Ayant repris la suite de ces recherches, après une longue interruption, j'ai réussi à perfectionner cette théorie dans quelques parties, principalement dans celle qui concerne les fonctions elliptiques de la troisième espèce. Ces améliorations étaient de nature à apporter quelque changement à mon premier travail; j'ai donc cru devoir traiter cette matière dans un nouvel ordre, en lui donnant de plus grands développemens et l'éclaircissant par des exemples choisis. C'est le résultat de ce dernier travail que je présente dans ce moment aux Géomètres; j'espère qu'ils voudront bien l'accueillir comme une nouvelle branche d'analyse qui peut offrir de belles et d'utiles applications.

\section*{Idée générale des différentes sortes de transcendantes contenues dans la formule intégrale $\displaystyle\int\frac{P\,dx}{R}$.}

\subsection*{(1).}

Nous représentons par $P$ une fonction rationnelle quelconque de $x$, et par $R$ le radical $\sqrt{(\alpha+\beta x+\gamma x^{2}+\delta x^{3}+\varepsilon x^{4})}$: ce radical restant le même, on peut donner une infinité de valeurs à $P$; mais il n'en résulte pas pour cela une infinité de transcendantes de nature différente. On peut toujours par des intégrations partielles, réduire l'intégrale $\displaystyle\int\frac{P\,dx}{R}$ à une partie algébrique, plus un certain nombre de transcendantes, qui sont toujours de la même forme et de la même nature. C'est ce qu'il s'agit de développer.

Supposons d'abord que $P$ soit une fonction entière de $x$, ensorte qu'on ait $P=A+Bx+Cx^{2}+\ldots+Kx^{k}$. Si on représente, pour abréger, l'intégrale $\displaystyle\int\frac{x^{m}\,dx}{R}$ par $\Pi^{m}$, il est clair qu'on aura
\[ \int\frac{P\,dx}{R}=A\Pi^{0}+B\Pi^{1}+C\Pi^{2}+\ldots+K\Pi^{k}: \]
or, en différentiant la quantité $x^{m-3}R$, et revenant de la différentielle à l'intégrale, on trouve cette formule
\[ \begin{aligned} x^{m-3}R={}&(m-3)\alpha\Pi^{m-4}+(m-\tfrac{5}{2})\beta\Pi^{m-3}+(m-2)\gamma\Pi^{m-2}\\ &+(m-\tfrac{3}{2})\delta\Pi^{m-1}+(m-1)\varepsilon\Pi^{m}; \end{aligned} \]

\origpage{5}
d'où il suit que $\Pi^{m}$, $\Pi^{m-1}$, etc. peuvent s'abaisser successivement à des degrés inférieurs, et qu'on peut continuer la réduction jusqu'à ce qu'on n'ait que des quantités au-dessous de $\Pi^{3}$; car la réduction de $\Pi^{3}$ est encore possible, parce que dans le cas de $m=3$, le terme qui semblerait devoir contenir $\Pi^{-1}$ s'évanouit. Donc $P$ étant une fonction entière de $x$, l'intégrale $\displaystyle\int\frac{P\,dx}{R}$ pourra toujours se réduire à une quantité algébrique, plus une transcendante qui sera constamment de la forme
\[ \int(A+Bx+Cx^{2})\frac{dx}{R}. \]

\subsection*{(2).}

Considérons maintenant la formule dans toute sa généralité, et soit $P$ une fonction rationnelle quelconque de $x$: on fera comme s'il était question d'intégrer la fraction rationnelle $P\,dx$; on extraira d'abord la partie entière qui peut y être contenue, et cette partie sera traitée comme il vient d'être dit. On décomposera ensuite le reste en plusieurs fractions partielles, selon le nombre des facteurs du dénominateur; de là résulteront autant de termes dans l'intégrale totale, et chacun de ces termes pourra être représenté par l'expression générale $\displaystyle N\int\frac{dx}{(1+nx)^{k}R}$. Or, il est facile de réduire tous les termes de cette espèce à des termes semblables, où $k=1$; c'est ce que nous allons prouver dans un instant. Observons auparavant que si le dénominateur, au lieu d'être complexe, était simplement $x^{k}$, l'intégrale $\displaystyle\int\frac{dx}{x^{k}R}$, et toutes les semblables où $k>1$, pourraient se déterminer par le moyen de l'intégrale $\displaystyle\int\left(\frac{A}{x}+B+Cx+Dx^{2}\right)\frac{dx}{R}$, il suffirait pour cela de donner à $m-4$ des valeurs négatives dans la formule de l'article (1).

Revenons au terme général $\displaystyle\int\frac{dx}{(1+nx)^{k}R}$ que nous appellerons $\Gamma^{k}$. Si on fait $1+nx=\omega$, et qu'on prenne les coefficiens $\alpha'$, $\beta'$, etc., de manière qu'on ait l'équation identique
\[ \begin{aligned} &\alpha+\beta\left(\frac{\omega-1}{n}\right)+\gamma\left(\frac{\omega-1}{n}\right)^{2}+\delta\left(\frac{\omega-1}{n}\right)^{3}+\varepsilon\left(\frac{\omega-1}{n}\right)^{4}\\ &=\alpha'+\beta'\omega+\gamma'\omega^{2}+\delta'\omega^{3}+\varepsilon'\omega^{4}, \end{aligned} \]

\origpage{6}
on trouvera par la différentiation de la quantité $\omega^{-k+1}R$ cette formule générale:
\[ \begin{aligned} \frac{-R}{n(1+nx)^{k-1}}={}&(k-1)\alpha'\Gamma^{k}+(k-\tfrac{3}{2})\beta'\Gamma^{k-1}+(k-2)\gamma'\Gamma^{k-2}\\ &+(k-\tfrac{5}{2})\delta'\Gamma^{k-3}+(k-3)\varepsilon'\Gamma^{k-4}. \end{aligned} \]
D'où il suit que les quantités $\Gamma^{2}$, $\Gamma^{3}$, etc., peuvent se déterminer au moyen de $\Gamma^{1}$, $\Gamma^{0}$, $\Gamma^{-1}$, $\Gamma^{-2}$; or, $\displaystyle\Gamma^{0}=\int\frac{dx}{R}$, $\displaystyle\Gamma^{-1}=\int(1+nx)\frac{dx}{R}$, $\displaystyle\Gamma^{-2}=\int(1+nx)^{2}\frac{dx}{R}$. Donc en général la formule $\displaystyle\int\frac{dx}{(1+nx)^{k}R}$, et toutes les formules semblables dans lesquelles $k$ est plus grand que l'unité, se réduiront toujours à une partie algébrique, plus une intégrale de la forme
\[ \int\left(\frac{A}{1+nx}+B+Cx+Dx^{2}\right)\frac{dx}{R}. \]

\subsection*{(3).}

Nous conclurons de là que, quelle que soit la fonction rationnelle de $x$ représentée par $P$, l'intégrale $\displaystyle\int\frac{P\,dx}{R}$ sera toujours décomposable en trois parties principales, la première algébrique, la seconde de la forme $\displaystyle\int(A+Bx+Cx^{2})\frac{dx}{R}$, et la troisième renfermant un ou plusieurs termes de la forme $\displaystyle N\int\frac{dx}{(1+nx)R}$, où le coefficient $n$ peut être réel ou imaginaire.

On voit par cette analyse que le nombre des transcendantes comprises dans la formule $\displaystyle\int\frac{P\,dx}{R}$ est très-limité. Il n'en existe que de deux espèces principales, l'une de la forme $\displaystyle\int(A+Bx+Cx^{2})\frac{dx}{R}$, l'autre de la forme $\displaystyle\int\frac{dx}{(1+nx)R}$. J'observe même que tant que $n$ est réel, cette seconde espèce est comprise dans la première, et s'y ramène immédiatement, en faisant $1+nx=\dfrac{1}{z}$.

Ces réductions générales une fois aperçues, nous allons suivre une autre route pour parvenir à une connaissance plus précise des mêmes transcendantes.

\origpage{7}

\section*{Manière de faire disparaître les puissances impaires de la variable sous le radical.}

\subsection*{(4).}

Le radical étant toujours $\sqrt{(\alpha+\beta x+\gamma x^{2}+\delta x^{3}+\varepsilon x^{4})}$, supposons que la quantité sous le signe soit décomposée en deux facteurs réels $\zeta+2\eta x+\theta x^{2}$, $\lambda+2\mu x+\nu x^{2}$. Ces facteurs doivent être de même signe pour que le radical soit réel; ainsi on peut supposer
\[ \zeta+2\eta x+\theta x^{2}=(\lambda+2\mu x+\nu x^{2})y^{2}. \]
De là on tire,
\[ x=\frac{\mu y^{2}-\eta+\sqrt{[(\mu y^{2}-\eta)^{2}+(\lambda y^{2}-\zeta)(\theta-\nu y^{2})]}}{\theta-\nu y^{2}}; \]
et si on appelle, pour abréger, $Y$ le radical de cette expression, on aura $\dfrac{dx}{R}=\dfrac{dy}{Y}$. La valeur de $x$ étant substituée dans $P$, on voit que la différentielle $\dfrac{P\,dx}{R}$ ou $\dfrac{P\,dy}{Y}$ se décomposera toujours en deux parties, l'une rationnelle et intégrable par les règles ordinaires, l'autre affectée du radical $Y$, mais telle, qu'il n'y aura que des puissances paires de $y$ sous le radical et hors du radical. Ainsi l'intégration de la formule $\dfrac{P\,dx}{R}$ se réduit toujours à celle d'une formule de la même nature, dans laquelle $P$ et $R$ ne contiennent que des puissances paires de $x$.

\subsection*{(5).}

Cette méthode est générale; cependant, comme la valeur de $x$, qui doit être substituée dans $P$, est un peu compliquée, il ne sera pas inutile de faire voir qu'on peut parvenir au même but par une substitution beaucoup plus simple, qui consiste à faire $x=\dfrac{p+qy}{1+y}$, $p$ et $q$ étant deux constantes indéterminées.

Reprenons les facteurs $\zeta+2\eta x+\theta x^{2}$, $\lambda+2\mu x+\nu x^{2}$, et substituons dans chacun d'eux la valeur de $x$, on pourra faire abstraction du dénominateur commun qui sort du radical, et pour que les

\origpage{8}
puissances impaires de $y$ disparaissent dans les numérateurs, il faudra qu'on ait
\[ \begin{aligned} &\zeta+\eta(p+q)+\theta pq=0,\\ &\lambda+\mu(p+q)+\nu pq=0. \end{aligned} \]
Ces deux équations donneront des valeurs rationnelles de $p+q$ et $pq$; mais pour que $p$ et $q$ soient réels, il faut encore que $(p+q)^{2}-4pq$, ou $(p-q)^{2}$ soit positif. Or, on peut distinguer deux cas:

1°. Si la quantité $\alpha+\beta x+\gamma x^{2}+\text{etc.}$ n'a pas tous ses facteurs réels, on pourra supposer que $\lambda+2\mu x+\nu x^{2}$ en contient deux imaginaires, et qu'on a en conséquence $\lambda\nu>\mu^{2}$. Mais la seconde des équations en $p$ et $q$ donne
\[ (p-q)^{2}=\left(\frac{\lambda+\nu pq}{\mu}-\frac{2\mu}{\nu}\right)^{2}+\frac{4\lambda\nu-4\mu^{2}}{\nu^{2}}; \]
donc dans ce premier cas le second membre est positif, et les valeurs de $p$ et $q$ seront réelles.

2°. Si la quantité $\alpha+\beta x+\text{etc.}$ a tous ses facteurs réels, et qu'en la décomposant en deux facteurs du second degré, comme on vient de faire, il n'en résulte pas des valeurs réelles pour $p$ et $q$; alors on observera qu'il y a deux autres manières de décomposer la quantité du quatrième degré en deux facteurs du second, et on peut être assuré que ces deux nouvelles combinaisons donneront des valeurs réelles pour $p$ et $q$. En voici la démonstration.

Soient $a$, $b$, $c$, $d$ les quatre valeurs de $x$ qu'on a en faisant $R=0$, les équations qui déterminent $p$ et $q$ pourront s'écrire ainsi:
\[ \begin{aligned} &ab-\tfrac{1}{2}(a+b)(p+q)+pq=0.\\ &cd-\tfrac{1}{2}(c+d)(p+q)+pq=0, \end{aligned} \]
\ednote{The print ends the first of these two equations with a period, where a comma is expected.}
et il en résulte
\[ \begin{aligned} &\frac{p+q}{2}=\frac{ab-cd}{a+b-c-d},\\ &\left(\frac{p-q}{2}\right)^{2}=\frac{a-c.a-d.b-c.b-d}{(a+b-c-d)^{2}}. \end{aligned} \]

\origpage{9}
Or, on est maître de rendre positif le numérateur de cette dernière quantité; car supposons que $a$, $b$, $c$, $d$, soient écrites dans l'ordre de leur grandeur, les racines négatives venant après les positives, alors les différences $a-b$, $a-c$, etc. seront toutes positives, et $p-q$ sera réel. On aura encore $p-q$ réel, si on prend pour $a$ et $b$ les extrêmes en grandeur des racines $a$, $b$, $c$, $d$. Enfin, la troisième combinaison donnerait $p-q$ imaginaire.

Donc par la substitution de $x=\dfrac{p+qy}{1+y}$, il est toujours possible de faire disparaître les puissances impaires de la variable sous le radical, et en même temps on obtient cet avantage, que les deux facteurs, sous le nouveau radical, sont réels et de la forme $f+gy^{2}$, ce qui est un point essentiel, et qu'on n'obtiendrait pas toujours par la première méthode.

Remarquons qu'il y a deux cas particuliers à examiner.

1°. Lorsque l'une des racines $a$ et $b$ est égale à l'une des racines $c$ et $d$, le radical $R$ se simplifie, et l'intégrale ne dépend plus que des arcs de cercle et des logarithmes.

2°. Si l'on a $a+b=c+d$, il suffit d'une simple permutation pour cesser d'avoir $a+b=c+d$; mais on peut aussi profiter de cette circonstance pour faciliter la transformation. En effet, les deux facteurs de la quantité sous le radical étant alors de la forme $\lambda+\nu(x^{2}+2mx)$, $\zeta+\theta(x^{2}+2mx)$, il est clair que si on fait $x+m=y$, les puissances impaires de la variable disparaîtront.

\section*{Réduction de la différentielle $\dfrac{P\,dx}{R}$ à la forme $\dfrac{Q\,d\varphi}{\sqrt{(1-c^{2}\sin.^{2}\varphi)}}$.}

\subsection*{(6).}

Puisque par une première préparation on peut faire ensorte qu'il n'y ait pas de puissances impaires de la variable sous le radical, nous ferons désormais $R=\sqrt{(\alpha+\beta x^{2}+\gamma x^{4})}$. Nous supposerons aussi que $P$ est une fonction paire de $x$; car quelle que soit la fonction rationnelle $P$, on peut toujours faire $P=M+Nx$, $M$ et $N$ étant des fonctions paires de $x$: or la partie $\dfrac{Nx\,dx}{R}$ se ramène aux règles ordinaires, en faisant $x^{2}=y$; ainsi toute la difficulté se réduit à intégrer la formule $\dfrac{M\,dx}{R}$, dans laquelle $M$ est une fonction paire de $x$.

\origpage{10}

Cela posé, nous allons prouver par l'énumération des différens cas, que la différentielle $\dfrac{dx}{R}$ peut toujours se ramener à la forme $\dfrac{m\,d\varphi}{\sqrt{(1-c^{2}\sin^{2}.\varphi)}}$, où l'on a $c$ plus petit que l'unité.

\emph{Premier cas.} Supposons les facteurs de $\alpha+\beta x^{2}+\gamma x^{4}$ imaginaires, et représentons cette quantité par $\lambda^{2}+2\lambda\mu x^{2}\cos\theta+\mu^{2}x^{4}$, $\lambda$ et $\mu$ étant positifs, et $\cos\theta$ pouvant être positif ou négatif. On fera $x=\sqrt{\left(\dfrac{\lambda}{\mu}\right)}\cdot\operatorname{tang}\tfrac{1}{2}\varphi$, et prenant $c=\sin\tfrac{1}{2}\theta$, on aura la transformée
\[ \frac{dx}{R}=\frac{1}{2\sqrt{\lambda\mu}}\cdot\frac{d\varphi}{\sqrt{(1-c^{2}\sin^{2}\varphi)}}. \]

\emph{Second cas.} Soit $\alpha+\beta x^{2}+\gamma x^{4}=m^{2}(1+p^{2}x^{2})(1-q^{2}x^{2})$; la limite de $x$ étant $\dfrac{1}{q}$, on fera $x=\dfrac{1}{q}\cos\varphi$ et $\dfrac{p^{2}}{p^{2}+q^{2}}=c^{2}$, ce qui donnera
\[ \frac{dx}{R}=\frac{c}{mp}\cdot\frac{-d\varphi}{\sqrt{(1-c^{2}\sin^{2}\varphi)}}. \]
Si on voulait que la transformée fût positive, il faudrait faire $\cot\varphi=\sqrt{(1-c^{2})}\operatorname{tang}\psi$, ce qui donne directement
\[ x^{2}=\frac{\sin^{2}\psi}{q^{2}+p^{2}\cos^{2}\psi}, \]
et la transformée serait
\[ \frac{dx}{R}=\frac{c}{mp}\cdot\frac{d\psi}{\sqrt{(1-c^{2}\sin^{2}\psi)}}. \]

\emph{Troisième cas.} Soit $\alpha+\beta x^{2}+\gamma x^{4}=m^{2}(1+p^{2}x^{2})(x^{2}-q^{2})$, on fera $x=\dfrac{q}{\cos\varphi}$, $\dfrac{1}{1+p^{2}q^{2}}=c^{2}$, et on aura
\[ \frac{dx}{R}=\frac{c}{m}\cdot\frac{d\varphi}{\sqrt{(1-c^{2}\sin^{2}\varphi)}}. \]

\emph{Quatrième cas.} Soit $\alpha+\beta x^{2}+\gamma x^{4}=m^{2}(1+p^{2}x^{2})(1+q^{2}x^{2})$, on supposera $p>q$, et faisant $x=\dfrac{\operatorname{tang}\varphi}{p}$, $\dfrac{p^{2}-q^{2}}{p^{2}}=c^{2}$, on aura
\[ \frac{dx}{R}=\frac{1}{mp}\cdot\frac{d\varphi}{\sqrt{(1-c^{2}\sin^{2}\varphi)}}. \]

\emph{Cinquième cas.} Soit $\alpha+\beta x^{2}+\gamma x^{4}=m^{2}(1-p^{2}x^{2})(1-q^{2}x^{2})$,

\origpage{11}
on supposera $p>q$, et faisant $px=\sin\varphi$, $\dfrac{q}{p}=c$, la transformée sera
\[ \frac{dx}{R}=\frac{1}{mp}\cdot\frac{d\varphi}{\sqrt{(1-c^{2}\sin^{2}\varphi)}}. \]

Cette formule servira depuis $x=0$ jusqu'à $x=\dfrac{1}{p}$; le radical $R$ serait imaginaire depuis $x=\dfrac{1}{p}$ jusqu'à $x=\dfrac{1}{q}$; mais il redevient réel depuis $x=\dfrac{1}{q}$ jusqu'à $x=\infty$. Dans ce dernier cas, il faut écrire $R^{2}=m^{2}(p^{2}x^{2}-1)(q^{2}x^{2}-1)$, et faisant $qx=\dfrac{1}{\sin\varphi}$, $\dfrac{q}{p}=c$, la transformée sera $\dfrac{-d\varphi}{mp\sqrt{(1-c^{2}\sin^{2}\varphi)}}$. Si on veut qu'elle soit positive, il faudra, comme ci-dessus, faire $\cot\varphi=\sqrt{(1-c^{2})}\cdot\operatorname{tang}\psi$, ce qui donnera directement $x^{2}=\dfrac{1-c^{2}\sin^{2}\psi}{q^{2}\cos^{2}\psi}$, et la transformée sera
\[ \frac{dx}{R}=\frac{1}{mp}\cdot\frac{d\psi}{\sqrt{(1-c^{2}\sin^{2}\psi)}}, \]
formule absolument semblable à la première; d'où il suit que l'intégrale de $\dfrac{dx}{R}$ pour le cas de $x>\dfrac{1}{q}$, se déduira toujours de la même intégrale, prise en supposant $x<\dfrac{1}{p}$.

\emph{Sixième cas.} Enfin, soit $\alpha+\beta x^{2}+\gamma x^{4}=m^{2}(x^{2}-q^{2})(p^{2}-x^{2})$; alors $x$ doit être compris entre $p$ et $q$. Soit $p>q$, et soit fait $x=\dfrac{q}{\cos\psi}$, on aura d'abord $\dfrac{dx}{R}=\dfrac{d\psi}{m\sqrt{(p^{2}-q^{2}-p^{2}\sin^{2}\psi)}}$. Dans cette formule, l'angle $\psi$ a une limite; pour en introduire un indéfini, soit $\sin\psi=c\sin\varphi$ et $c^{2}=\dfrac{p^{2}-q^{2}}{p^{2}}$, on aura la transformée
\[ \frac{dx}{R}=\frac{1}{mp}\cdot\frac{d\varphi}{\sqrt{(1-c^{2}\sin^{2}\varphi)}}, \]
à laquelle on serait parvenu directement en faisant
\[ x^{2}=\frac{q^{2}}{1-c^{2}\sin^{2}\varphi}. \]

\subsection*{(7).}

On peut remarquer maintenant, qu'abstraction faite du premier cas, les valeurs de $x^{2}$ qui opèrent la réduction cherchée, sont toujours

\origpage{12}
de la forme $\dfrac{A+B\sin^{2}\varphi}{C+D\sin^{2}\varphi}$, où les coefficiens sont constans. Il ne s'agit plus que de substituer cette valeur de $x^{2}$ dans $P$, et la formule $\displaystyle\int\frac{P\,dx}{R}$ sera transformée en une autre $\displaystyle\int\frac{Q\,d\varphi}{\sqrt{(1-c^{2}\sin^{2}\varphi)}}$, dans laquelle $c$ sera plus petit que l'unité, et $Q$ sera une fonction rationnelle paire de $\sin\varphi$, laquelle contiendra $\sin\varphi$ au même degré que $P$ contient $x$.

Nous faisons abstraction du premier cas, parce que la forme de la valeur de $x$ est un peu différente, et que d'ailleurs en suivant la méthode de l'article 5, on évite ce cas, puisqu'on tombe toujours sur des facteurs réels. Mais nous reviendrons dans un autre article sur le cas des facteurs imaginaires.

\section*{Développement de la formule $\displaystyle\int\frac{Q\,d\varphi}{\sqrt{(1-c^{2}\sin^{2}\varphi)}}$.}

\subsection*{(8).}

Nous représenterons dorénavant le radical $\sqrt{(1-c^{2}\sin^{2}\varphi)}$ par $\Delta$, et aussi par $\Delta(\varphi)$ et $\Delta(c,\varphi)$, lorsqu'on le regardera comme fonction de $\varphi$, ou comme fonction de $c$ et $\varphi$.

Considérons d'abord le cas où $Q$ serait une fonction entière, et soit $Q=A+B\sin^{2}\varphi+C\sin^{4}\varphi+\text{etc.}$; si on représente pour abréger, l'intégrale $\displaystyle\int\frac{\sin^{2k}\varphi\cdot d\varphi}{\Delta}$ par $Z^{2k}$, on aura
\[ \int\frac{Q\,dx}{\Delta}=AZ^{0}+BZ^{2}+CZ^{4}+\text{etc.} \]
\ednote{Printed $Q\,dx$ in the numerator; the integrand is in $\varphi$, so $Q\,d\varphi$ is expected. An apparent misprint.}
Mais en différentiant la quantité $\Delta\cos\varphi\sin^{2k-3}\varphi$, on trouve aisément la formule
\[ \Delta\cos\varphi\sin^{2k-3}\varphi=(2k-3)Z^{2k-4}-(1+c^{2})(2k-2)Z^{2k-2}+c^{2}(2k-1)Z^{2k}; \]
d'où il suit que $Z^{4}$, $Z^{6}$, etc.\ peuvent s'exprimer au moyen de $Z^{0}$ et $Z^{2}$, et qu'ainsi dans le cas où $Q$ est une fonction entière, l'intégrale $\displaystyle\int\frac{Q\,d\varphi}{\Delta}$ est égale à une quantité algébrique, plus une intégrale de la forme $\displaystyle\int(A+B\sin^{2}\varphi)\frac{d\varphi}{\Delta}$.

\subsection*{(9).}

Soit maintenant $Q$ une fonction quelconque rationnelle de $\sin^{2}\varphi$; après avoir extrait la partie entière, et l'avoir traitée comme

\origpage{13}
il vient d'être dit, le développement de la partie fractionnaire pourra toujours se faire de manière que chaque fraction partielle soit de la forme $\dfrac{N}{(1+n\sin^{2}\varphi)^{k}}$, $n$ et $N$ étant des coefficiens constans, réels ou imaginaires. Il s'agira donc de réduire la formule $\displaystyle\int\frac{d\varphi}{(1+n\sin^{2}\varphi)^{k}\Delta}$, et toutes les semblables aux formules les plus simples de la même espèce. Or, si on fait pour un moment $\sin\varphi=x$, cette formule deviendra
\[ \int\frac{dx}{(1+nx^{2})^{k}\sqrt{[1-(1+c^{2})x^{2}+c^{2}x^{4}]}}. \]
Considérons, pour plus de généralité, la formule
\[ \Pi^{k}=\int\frac{dx}{(1+nx^{2})^{k}R}, \]
dans laquelle $R$ représente le radical $\sqrt{(\alpha+\beta x^{2}+\gamma x^{4})}$, on trouvera par les moyens déjà indiqués,
\[ \begin{aligned} \frac{xR}{(1+nx^{2})^{k-1}}={}&(2k-2)\left(\alpha-\frac{\beta}{n}+\frac{\gamma}{n^{2}}\right)\Pi^{k}-(2k-3)\left(\alpha-\frac{2\beta}{n}+\frac{3\gamma}{n^{2}}\right)\Pi^{k-1}\\ &+(2k-4)\left(-\frac{\beta}{n}+\frac{3\gamma}{n^{2}}\right)\Pi^{k-2}-(2k-5)\frac{\gamma}{n^{2}}\Pi^{k-3}. \end{aligned} \]
De là il résulte que le terme $\Pi^{k}$, et tous les semblables où $k$ est plus grand que l'unité, peuvent s'exprimer en partie algébriquement, en partie par les quantités $\Pi^{1}$, $\Pi^{0}$, $\Pi^{-1}$, qui sont les plus simples de leur espèce. On peut même observer que si $1+nx^{2}$ était diviseur de la quantité sous le radical, auquel cas on aurait $\alpha-\dfrac{\beta}{n}+\dfrac{\gamma}{n^{2}}=0$, alors la formule précédente aurait un terme de moins; ainsi la quantité $\Pi^{k}$ ne dépendrait plus que des deux $\Pi^{0}$ et $\Pi^{-1}$.

Donc la détermination de $\Pi^{k}$ dépendra en général de la formule
\[ \int\left(\frac{A}{1+nx^{2}}+B+Cx^{2}\right)\frac{dx}{R}, \]
et en particulier lorsque $1+nx^{2}$ sera diviseur de $R^{2}$, on pourra faire disparaître ce dénominateur, et la détermination de $\Pi^{k}$ ne dépendra plus que de la formule $\displaystyle\int(B+Cx^{2})\frac{dx}{R}$.

\origpage{14}

\subsection*{(10).}

L'application de ce résultat à l'intégrale $\displaystyle\int\frac{d\varphi}{(1+n\sin^{2}\varphi)^{k}\Delta}$ fait voir qu'en général cette intégrale dépendra de la suivante,
\[ \int\left(\frac{A}{1+n\sin^{2}\varphi}+B+C\sin^{2}\varphi\right)\frac{d\varphi}{\Delta}. \]
Et dans le cas particulier où $1+n\sin^{2}\varphi$ sera facteur de $\cos^{2}\varphi\cdot\Delta^{2}$, l'intégrale pourra être débarrassée du dénominateur, et ne dépendra plus que de la formule entière
\[ \int(B+C\sin^{2}\varphi)\frac{d\varphi}{\Delta}. \]
Ce cas particulier aura lieu si $n=-1$, et si $n=-c^{2}$, alors les formules sont $\displaystyle\int\frac{d\varphi}{\cos^{2k}\varphi\cdot\Delta}$, et $\displaystyle\int\frac{d\varphi}{\Delta^{2k+1}}$; il est facile en effet de réduire l'une et l'autre à une partie algébrique, plus une intégrale de la forme $\displaystyle\int(A+B\sin^{2}\varphi)\frac{d\varphi}{\Delta}$. On peut ajouter à ces deux cas celui de $\displaystyle\int\frac{d\varphi}{\sin^{2k}\varphi\cdot\Delta}$, qui est susceptible d'une semblable réduction, et qu'on effectuera par la formule de l'article 8.

Il résulte de cette analyse, que la formule générale $\displaystyle\int\frac{Q\,d\varphi}{\Delta}$, réduite convenablement, contiendra, 1°.\ une partie algébrique; 2°.\ une intégrale de la forme $\displaystyle\int(A+B\sin^{2}\varphi)\frac{d\varphi}{\Delta}$; 3° une ou plusieurs parties de la forme $\displaystyle\int\frac{N}{1+n\sin^{2}\varphi}\cdot\frac{d\varphi}{\Delta}$, dans chacune desquelles les coefficiens $N$ et $n$ auront des valeurs quelconques, réelles ou imaginaires.

\section*{Définition des fonctions elliptiques, et leur division en trois espèces.}

\subsection*{(11).}

Puisque les transcendantes que nous avons considérées, se réduisent toujours à ces deux formes $\displaystyle\int(A+B\sin^{2}\varphi)\frac{d\varphi}{\Delta}$, $\displaystyle\int\frac{d\varphi}{(1+n\sin^{2}\varphi)\Delta}$, il est clair qu'elles sont comprises dans la formule générale
\[ H=\int\frac{A+B\sin^{2}\varphi}{1+n\sin^{2}\varphi}\cdot\frac{d\varphi}{\Delta}. \]

\origpage{15}
Nous appellerons désormais \emph{fonctions} ou \emph{transcendantes elliptiques}, les intégrales comprises dans cette formule. La transcendante $H$ sera supposée s'évanouir ou commencer lorsque $\varphi=0$; son étendue sera déterminée par la variable $\varphi$, qu'on appellera \emph{l'amplitude}; la constante $c$, toujours plus petite que l'unité, s'appellera \emph{le module}; enfin $\sqrt{(1-c^{2})}$ que nous désignerons constamment par $b$, pourra s'appeler \emph{le complément du module}.

La quantité $H$, lorsqu'on ne considère que la variabilité de $\varphi$, peut se désigner par $H(\varphi)$; si on considère à la fois la variation du module et de l'amplitude, on peut la représenter par $H(c,\varphi)$, et ainsi de suite, si on voulait avoir égard à l'inégalité des coefficiens.

\subsection*{(12).}

Il suffit de connaître toutes les valeurs de $H$, depuis $\varphi=0$ jusqu'à $\varphi=90^{\circ}=\dfrac{\pi}{2}$, et on connaîtra facilement la valeur de cette transcendante pour une amplitude quelconque. En effet, soit $\varphi=i\pi\pm\alpha$, $i$ étant un entier, et $\alpha$ un arc plus petit que $\dfrac{\pi}{2}$, alors on aura
\[ H(\varphi)=2iH\left(\frac{\pi}{2}\right)\pm H(\alpha). \]
Il en est à cet égard de la fonction $H$ comme des arcs d'ellipse, et en général des arcs de toutes les courbes ovales composées de quatre parties égales et semblables: un arc, quelque grand qu'il soit, et renfermant, si l'on veut, plusieurs circonférences, s'exprime toujours sans difficulté par le quart de la courbe et une portion de ce quart. La fonction déterminée $H\left(\dfrac{\pi}{2}\right)$ est en quelque sorte l'unité des fonctions $H$, ou la fonction devenue complète: nous la désignerons par $H^{1}$.

Il ne paraît pas que la fonction $H$, prise dans toute sa généralité, puisse se réduire à des arcs d'ellipse; cela n'a lieu que lorsque $n=0$, ou lorsque quelque substitution peut faire disparaître le dénominateur $1+n\sin^{2}\varphi$, ce que nous ne croyons pas possible en général. Ainsi la dénomination de fonction elliptique est impropre à quelques égards; nous l'adoptons néanmoins à cause de la grande analogie qu'on trouvera entre les propriétés de cette fonction et celles des arcs d'ellipse.

\origpage{16}

\subsection*{(13).}

Puisque toute intégrale désignée par $\displaystyle\int\frac{P\,dx}{R}$, se réduit à une partie algébrique, plus un certain nombre de termes qui peuvent chacun être assimilés à la fonction $H$, il s'ensuit qu'on pourrait n'admettre, pour les intégrales dont il s'agit, qu'une seule espèce de transcendantes, représentée par la fonction $H$, et dans laquelle les coefficiens $A$, $B$, $n$, seraient à volonté réels ou imaginaires. Mais pour bien pénétrer la nature de ces intégrales et pouvoir établir entre elles les comparaisons et les réductions dont elles sont susceptibles, il est nécessaire de diviser la fonction $H$ en plusieurs espèces distinctes, dont les propriétés deviendront plus sensibles, lorsqu'on les considérera chacune isolément.

J'observe d'abord que les arcs d'ellipse sont contenus dans la formule $H$. En effet, l'équation de l'ellipse étant $y^{2}=\dfrac{b^{2}}{a^{2}}(a^{2}-x^{2})$, si on fait $x=a\sin\varphi$, on aura
\[ y=b\cos\varphi,\quad\text{et}\quad\sqrt{(dx^{2}+dy^{2})}=d\varphi\sqrt{(a^{2}\cos^{2}\varphi+b^{2}\sin^{2}\varphi)}, \]
formule qui se réduit à $d\varphi\sqrt{(1-c^{2}\sin^{2}\varphi)}$, en faisant $a=1$, et $c^{2}=1-b^{2}$.

\rmfigure{figures/fig-1.png}{Fig.~1.}{A quadrant of the ellipse BMKNA with centre C and semi-axes CA and CB, inside the dotted circular quadrant DM'K'N'A, with the ordinates PMM', IKK' and QNN' standing on CA, the lines CO and CZ from the centre, the segment OM and the segment ZN.}

Soit donc le demi-grand axe $CA=1$, le demi-petit axe $CB=b$; si sur le cercle circonscrit $DM'A$, on prend l'arc $DM'=\varphi$, et qu'on abaisse du point $M'$ la perpendiculaire $M'P$ sur le grand axe, on déterminera sur l'ellipse un arc $BM=E$, dont la valeur sera
\[ E=\int\Delta\,d\varphi. \]
Cette fonction ou transcendante $E$ constitue l'une des espèces dans lesquelles nous diviserons la fonction $H$; elle se déduit de la formule générale en faisant $n=0$, $A=1$, $B=-c^{2}$.

\rmfigure{figures/fig-2.png}{Fig.~2.}{A hyperbola branch rising from A through M, N and K towards O, with the axis CAPF, the perpendicular CB, the straight line CV from C, the ordinate PM, the tangent at M extended down to Z, and the segment CZ.}

L'équation de l'hyperbole étant $y^{2}=\dfrac{\beta^{2}}{\alpha^{2}}(x^{2}-\alpha^{2})$, si l'on fait semblablement $x=\dfrac{\alpha}{\cos\varphi}$, on aura
\[ y=\frac{\beta\sin\varphi}{\cos\varphi},\quad\text{et}\quad\sqrt{(dx^{2}+dy^{2})}=\frac{d\varphi}{\cos^{2}\varphi}\sqrt{(\beta^{2}+\alpha^{2}\sin^{2}\varphi)}; \]
mais pour avoir un radical entièrement semblable à celui de l'arc d'ellipse, il faut prendre d'autres dénominations. Soit donc le demi-axe transverse $CA=c$, son conjugué $CB=b$, la distance du

\origpage{17}
centre au foyer $CF=1$, l'ordonnée $y=b^{2}\operatorname{tang}\varphi$, on aura l'abscisse $x=\dfrac{c}{\cos\varphi}\sqrt{(1-c^{2}\sin^{2}\varphi)}$, d'où résulte $\sqrt{(dx^{2}+dy^{2})}=\displaystyle\int\frac{b^{2}\,d\varphi}{\Delta\cos^{2}\varphi}$. Mais en différentiant la quantité $\Delta\operatorname{tang}\varphi$, on a
\[ d(\Delta\operatorname{tang}\varphi)=\frac{b^{2}\,d\varphi}{\Delta\cos^{2}\varphi}-\frac{b^{2}\,d\varphi}{\Delta}+\Delta\,d\varphi. \]
Donc si on appelle $\Upsilon$ l'arc d'hyperbole $AM$ qui répond à l'amplitude $\varphi$, ou dont l'ordonnée extrême $PM=b^{2}\operatorname{tang}\varphi$, on aura
\[ \Upsilon=\Delta\operatorname{tang}\varphi-\int\Delta\,d\varphi+b^{2}\int\frac{d\varphi}{\Delta}: \]
et l'on peut remarquer que la partie algébrique $\Delta\operatorname{tang}\varphi$ n'est autre chose que la tangente $MZ$ terminée par la perpendiculaire $CZ$ abaissée du centre sur cette tangente.

La fonction $\Upsilon$ est comprise dans la formule $H$, puisqu'elle s'en déduit en faisant $A=b^{2}$, $B=-b^{2}c^{2}$, $n=-1$; si on prend cette fonction pour la seconde espèce des transcendantes contenues dans la formule $H$, l'intégrale $\displaystyle\int\frac{d\varphi}{\Delta}$ pourra s'exprimer par les arcs $E$ et $\Upsilon$ au moyen de l'équation
\[ b^{2}\int\frac{d\varphi}{\Delta}=E+\Upsilon-\Delta\operatorname{tang}\varphi. \]
Enfin, comme on peut mettre $H$ sous la forme
\[ H=A'\int\frac{d\varphi}{\Delta}+B'\int\Delta\,d\varphi+C'\int\frac{d\varphi}{(1+n\sin^{2}\varphi)\Delta}; \]
il ne restera plus à considérer qu'une troisième espèce de transcendantes, représentée par la formule
\[ \Pi=\int\frac{d\varphi}{(1+n\sin^{2}\varphi)\Delta}. \]

Telle serait donc la division des fonctions elliptiques en trois espèces, si l'on admettait les arcs $E$ et $\Upsilon$ comme constituant les deux premières espèces, et c'est en effet la première idée qui se présente dans ce genre de recherches.

\subsection*{(14).}

Mais, par un examen plus approfondi des propriétés de ces

\origpage{18}
fonctions, on trouve que la fonction $F=\displaystyle\int\frac{d\varphi}{\Delta}$ doit être préférée à l'arc d'hyperbole $\Upsilon$, pour en faire l'une des trois espèces de fonctions elliptiques, et que même cette fonction $\displaystyle\int\frac{d\varphi}{\Delta}$ doit être considérée comme plus simple que l'arc d'ellipse $\displaystyle\int\Delta\,d\varphi$.

En effet, on prouvera ci-après que la fonction $F$ peut s'exprimer par deux arcs d'ellipse; mais l'inverse n'a pas lieu, et un arc d'ellipse ne peut pas s'exprimer par deux fonctions telles que $F$, ce qui indique déjà que la fonction $E$ est d'une nature plus composée que la fonction $F$. Mais cette conséquence se manifeste plus clairement encore par l'examen des propriétés respectives de ces fonctions.

La propriété la plus remarquable des fonctions $F$, est qu'on peut déterminer par des opérations purement algébriques, une fonction égale à la somme ou à la différence de deux autres fonctions; d'où il suit qu'on peut déterminer algébriquement une fonction multiple, sous-multiple ou en général qui soit dans un rapport rationnel avec une fonction donnée; propriété que les fonctions $F$ partagent avec les arcs de cercle et les logarithmes, et qui a lieu quand même ces fonctions, considérées comme des intégrales ou des arcs de courbe, n'auraient pas l'origine commune $\varphi=0$, et commenceraient à des points quelconques.

Les arcs d'ellipse et les arcs d'hyperbole sont de toutes les autres transcendantes, celles qui approchent le plus de jouir de la même propriété, mais elles n'en jouissent pas d'une manière absolue. Ainsi deux arcs étant donnés sur l'une de ces courbes, à compter du même point où $\varphi=0$, on peut trouver algébriquement, non pas un arc égal à leur somme, mais un arc égal à cette somme, plus ou moins une quantité algébrique, ce qui prouve que les arcs $\Upsilon$ et $E$ sont d'une nature plus composée que les fonctions $F$. On doit donc regarder ces dernières comme tenant le premier rang après les arcs de cercle et les logarithmes~${}^{(1)}$.

\textbf{(1)} Ces fonctions réunissent un si grand nombre de propriétés, que quand elles seront plus généralement connues, on jugera sans doute nécessaire de leur imposer un nom particulier, et de désigner la fonction de $c$ et $\varphi$ égale à $\displaystyle\int\frac{d\varphi}{\Delta}$, comme on désigne l'arc dont le sinus est $x$, ou le nombre dont le logarithme est $y$. Il semble qu'on caractériserait assez bien la fonction $F$ en lui donnant le nom de \emph{Nome}, parce que cette fonction a la propriété de régler tout ce qui concerne la comparaison des fonctions elliptiques. Peut-être conviendrait-il en même temps de donner les noms d'\emph{Epinome} et de \emph{Paranome} aux fonctions $E$ et $\Pi$ qui constituent les deux autres espèces.

\origpage{19}

Enfin un motif qui suffirait seul pour faire préférer les fonctions $F$ aux fonctions $\Upsilon$ dans le classement des fonctions elliptiques, c'est qu'on ne peut supposer $\varphi>\tfrac{1}{2}\pi$ dans la fonction $\Upsilon$, tandis qu'on peut supposer $\varphi$ d'une grandeur quelconque dans les fonctions $F$ et $E$ qui, en général, croissent presque proportionnellement à l'angle $\varphi$. Or les applications du calcul intégral, principalement celles qui concernent la mécanique, donnent souvent lieu d'attribuer à la variable principale $\varphi$ des valeurs quelconques, composées, si l'on veut, de plusieurs circonférences. L'emploi des fonctions $\Upsilon$ ferait donc naître des embarras ou même des erreurs, ce qui n'est jamais à craindre dans celui des fonctions $F$.

\subsection*{(15).}

Cela posé, les fonctions ou transcendantes elliptiques comprises dans la formule $H$, seront divisées en trois espèces:

La première et la plus simple est représentée par la formule $F=\displaystyle\int\frac{d\varphi}{\Delta}$;

La seconde est l'arc d'ellipse, compté depuis le petit axe et dont l'expression est $E=\displaystyle\int\Delta\,d\varphi$;

Enfin la troisième espèce est représentée par la formule $\Pi=\displaystyle\int\frac{d\varphi}{(1+n\sin^{2}\varphi)\Delta}$; elle contient, outre le module $c$ commun aux deux autres espèces, un paramètre $n$ qui peut être à volonté positif ou négatif, réel ou imaginaire.

On pourrait croire que le cas où $n$ est imaginaire, diffère essentiellement du cas où $n$ est réel, et qu'il exige la formation d'une quatrième espèce de fonctions elliptiques; mais par des réductions et des transformations que nous exposerons ci-après, on verra que cette quatrième espèce est inutile à considérer, et qu'on peut ainsi restreindre la troisième espèce aux seuls cas où $n$ est réel. Nous regarderons seulement comme conditions nécessaires et communes

\origpage{20}
aux trois espèces de fonctions, que l'amplitude $\varphi$ et le module $c$ soient réels, et qu'en même temps $c$ soit plus petit que l'unité.

\section*{Comparaison des fonctions elliptiques de la première espèce.}

\subsection*{(16).}

Tous les Géomètres connaissent l'intégrale algébrique complète qu'Euler a donnée de l'équation différentielle
\[ \frac{dx}{\sqrt{(\alpha+\beta x+\gamma x^{2}+\delta x^{3}+\varepsilon x^{4})}}+\frac{dy}{\sqrt{(\alpha+\beta y+\gamma y^{2}+\delta y^{3}+\varepsilon y^{4})}}=0. \]
D'après les réductions indiquées dans l'article 6, cette équation peut, sans perdre de sa généralité, être mise sous la forme
\[ \frac{d\varphi}{\sqrt{(1-c^{2}\sin^{2}\varphi)}}+\frac{d\psi}{\sqrt{(1-c^{2}\sin^{2}\psi)}}=0; \]
et alors son intégrale est $F(\varphi)+F(\psi)=F(\mu)$, $\mu$ étant une constante arbitraire. Mais la même intégrale trouvée par la méthode d'Euler, s'exprime ainsi:
\[ \cos\varphi\cos\psi-\sin\varphi\sin\psi\sqrt{(1-c^{2}\sin^{2}\mu)}=\cos\mu\ldots(a'), \]
et voici comment on peut vérifier ce résultat \emph{à posteriori}.

J'observe d'abord que l'équation $(a')$ peut être mise sous l'une ou l'autre des deux formes suivantes:
\[ \left.\begin{aligned} \cos\mu\cos\varphi+\sin\mu\sin\varphi\,\Delta(\psi)&=\cos\psi\\ \cos\mu\cos\psi+\sin\mu\sin\psi\,\Delta(\varphi)&=\cos\varphi \end{aligned}\right\}\ldots(b'); \]
car ces équations dégagées chacune du radical qu'elles contiennent, conduisent au même résultat que l'équation $(a')$ dégagée de son radical.

Cela posé, si on différentie l'équation $(a')$ après avoir divisé chaque membre par $\sin\varphi\sin\psi$, afin de faire disparaître le radical, on aura
\[ \frac{d\varphi}{\sin\varphi}(\cos\psi-\cos\mu\cos\varphi)+\frac{d\psi}{\sin\psi}(\cos\varphi-\cos\mu\cos\psi)=0. \]
Substituant dans celle-ci les valeurs de $\cos\psi-\cos\mu\cos\varphi$ et $\cos\varphi-\cos\mu\cos\psi$, données par les équations $(b')$, on aura
\[ \frac{d\varphi}{\Delta(\varphi)}+\frac{d\psi}{\Delta(\psi)}=0. \]

\origpage{21}
De là on voit que l'équation transcendante $F(\varphi)+F(\psi)=F(\mu)$ est satisfaite, soit par l'équation algébrique $(a')$, soit par l'une des équations $(b')$; car ces trois équations peuvent être employées indistinctement l'une pour l'autre.

Si on avait $c=0$, la fonction $F(\varphi)$ se réduirait à l'arc $\varphi$, et alors ayant $\varphi+\psi=\mu$, on en conclurait
\[ \begin{gathered} \cos\varphi\cos\psi-\sin\varphi\sin\psi=\cos\mu\\ \cos\mu\cos\varphi+\sin\mu\sin\varphi=\cos\psi\\ \cos\mu\cos\psi+\sin\mu\sin\psi=\cos\varphi. \end{gathered} \]
C'est en effet ce que donnent les équations $(a')$ et $(b')$ dans le cas de $c=0$, qui réduit à l'unité les radicaux $\Delta(\mu)$, $\Delta(\varphi)$, $\Delta(\psi)$.

\subsection*{(17).}

\rmfigure{figures/fig-3.png}{Fig.~3.}{A spherical triangle ABC with the side AC prolonged beyond C and a second arc drawn from the point D on BC up to meet that prolongation, together with a dashed arc from a to beta beside AB and short dotted marks lettered alpha, a at A and beta, b at B.}

C'est ici le lieu de faire observer, d'après Lagrange~${}^{(1)}$, que si l'on construit un triangle sphérique dont les côtés $AB$, $AC$, $BC$ soient respectivement égaux aux amplitudes $\mu$, $\varphi$, $\psi$, alors les angles opposés $C$, $B$, $A$, seront tels qu'on aura
\[ \begin{aligned} \cos C&=\frac{\cos\mu-\cos\varphi\cos\psi}{\sin\varphi\sin\psi}=-\sqrt{(1-c^{2}\sin^{2}\mu)}\\ \cos B&=\frac{\cos\varphi-\cos\mu\cos\psi}{\sin\mu\sin\psi}=\sqrt{(1-c^{2}\sin^{2}\varphi)}\\ \cos A&=\frac{\cos\psi-\cos\mu\cos\varphi}{\sin\mu\sin\varphi}=\sqrt{(1-c^{2}\sin^{2}\psi)}; \end{aligned} \]
d'où l'on déduit
\[ \begin{gathered} \sin^{2}C=c^{2}\sin^{2}\mu,\quad\sin^{2}B=c^{2}\sin^{2}\varphi,\quad\sin^{2}A=c^{2}\sin^{2}\psi,\\ \text{ou}\quad c=\frac{\sin C}{\sin\mu}=\frac{\sin B}{\sin\varphi}=\frac{\sin A}{\sin\psi}. \end{gathered} \]
Ces équations s'accordent avec les propriétés connues des triangles sphériques; elles prouvent en même temps que dans tout triangle sphérique formé par trois côtés $\varphi$, $\psi$, $\mu$, qui satisfont à l'équation transcendante $F(\varphi)+F(\psi)=F(\mu)$, le rapport du sinus de chaque

\textbf{(1)} \emph{Théorie des Fonct.\ analyt.}, pag.\ 85.

\origpage{22}
angle au sinus du côté opposé, sera égal au module~${}^{(1)}$. On voit de plus que dans le triangle $ABC$, les deux angles $A$ et $B$ sont toujours aigus, et l'angle $C$ toujours obtus.

Si on fait ensuite $F(\mu)+F(\theta)=F(\nu)$, il faudra observer que dans le nouveau triangle sphérique qui aura pour côtés $\mu$, $\theta$, $\nu$, l'angle opposé au côté $\mu$ devra être aigu et avoir le même sinus que l'angle opposé au côté $\mu$ dans le premier triangle, pour que son rapport avec $\sin\mu$ soit égal à $c$ dans les deux triangles; ainsi il faudra que l'angle opposé au côté $\mu$ dans le second triangle, soit supplément de l'angle opposé au côté $\mu$ dans le premier; d'où naît cette construction.

Prolongez le côté $AC$ vers $E$, faites $CD=\theta$; du point $D$ et d'un intervalle $DE=\mu$, décrivez un arc qui rencontrera $CE$ en $E$, et déterminera le second triangle $CDE$, dans lequel on aura $CE=\nu$, et le côté $\nu$ satisfera à l'équation $F(\nu)=F(\varphi)+F(\psi)+F(\theta)$.

Cette construction peut être continuée aussi loin que les limites des côtés des triangles sphériques peuvent le permettre, c'est-à-dire tant que le grand côté n'est pas plus grand qu'une demi-circonférence. Et il est évident qu'on pourrait se servir de la même construction pour trouver successivement les angles $\varphi_{2}$, $\varphi_{3}$, $\varphi_{4}$, etc., tels qu'on eût $F(\varphi_{2})=2F(\varphi)$, $F(\varphi_{3})=3F(\varphi)$, etc., ce qui servirait à la multiplication de la fonction $F$. Mais, nous le répétons, ces constructions ne peuvent s'étendre que jusqu'aux limites des triangles sphériques, tandis que l'analyse ne connaît aucune borne.

Il est bon de remarquer que la considération du triangle sphérique $ABC$ offre un moyen fort simple de vérifier l'intégrale $(a')$. En effet si en regardant $\mu$ et $C$ comme constans, on fait varier infiniment peu les côtés $\varphi$ et $\psi$, on aura $A\alpha=\beta b$, ou $-d\varphi\cos A=d\psi\cos B$; mais on a trouvé $\cos A=\Delta(\psi)$ et $\cos B=\Delta(\varphi)$; donc $\dfrac{d\varphi}{\Delta(\varphi)}+\dfrac{d\psi}{\Delta(\psi)}=0$.

Ainsi une application fort simple de la trigonométrie sphérique

\textbf{(1)} Un autre triangle sphérique formé par les trois angles $\varphi$, $\psi$, $180^{\circ}-\mu$, et par les trois côtés $B$, $A$, $180^{\circ}-C$, satisferait également à l'équation $F(\varphi)+F(\psi)=F(\mu)$. Mais le rapport qui était $c$ dans le premier triangle, deviendrait $\dfrac{1}{c}$ dans celui-ci.

\origpage{23}
aurait suffi pour trouver l'intégrale algébrique complète de l'équation transcendante $\dfrac{d\varphi}{\Delta(\varphi)}+\dfrac{d\psi}{\Delta(\psi)}=0$.

\subsection*{(18).}

Etant données deux fonctions elliptiques de première espèce $F(\varphi)$, $F(\psi)$, si on veut trouver une troisième fonction $F(\mu)$ égale à leur somme, il faut déterminer $\mu$ par l'équation $(a')$, ce qui donnera les formules
\[ \begin{aligned} \sin\mu &= \frac{\sin\varphi\cos\psi\,\Delta(\psi)+\sin\psi\cos\varphi\,\Delta(\varphi)}{1-c^{2}\sin^{2}\varphi\sin^{2}\psi} \\ \cos\mu &= \frac{\cos\varphi\cos\psi-\sin\varphi\sin\psi\,\Delta(\varphi)\,\Delta(\psi)}{1-c^{2}\sin^{2}\varphi\sin^{2}\psi} \\ \Delta(\mu) &= \frac{\Delta(\varphi)\,\Delta(\psi)-c^{2}\sin\varphi\sin\psi\cos\varphi\cos\psi}{1-c^{2}\sin^{2}\varphi\sin^{2}\psi} \\ \operatorname{tang}\mu &= \frac{\operatorname{tang}\varphi\,\Delta(\psi)+\operatorname{tang}\psi\,\Delta(\varphi)}{1-\operatorname{tang}\varphi\operatorname{tang}\psi\,\Delta(\varphi)\,\Delta(\psi)}. \end{aligned} \]
D'après cette dernière formule, si on prenait deux angles auxiliaires $\varphi'$, $\psi'$, tels que
\[ \operatorname{tang}\varphi'=\operatorname{tang}\varphi\,\Delta(\psi) \quad\text{et}\quad \operatorname{tang}\psi'=\operatorname{tang}\psi\,\Delta(\varphi), \]
il en résulterait
\[ \mu=\varphi'+\psi', \]
ce qui est un moyen de calculer aisément l'angle $\mu$ par les tables des sinus.

Si l'on fait $\mu=\frac{1}{2}\pi$, c'est-à-dire, si l'on a $F(\varphi)+F(\psi)=F^{1}$, les deux fonctions $F(\varphi)$, $F(\psi)$ seront en quelque sorte complémens l'une de l'autre, puisque leur somme est égale à la fonction complète $F^{1}$; alors on aura immédiatement par l'équation $(a')$,
\[ b\operatorname{tang}\varphi\operatorname{tang}\psi=1: \]
c'est la relation nécessaire entre les angles $\varphi$ et $\psi$, pour qu'on ait $F(\varphi)+F(\psi)=F(\frac{1}{2}\pi)=F^{1}$. On en déduit pour l'expression de $\psi$ en $\varphi$,
\[ \sin\psi=\frac{\cos\varphi}{\Delta(\varphi)}, \quad \cos\psi=\frac{b\sin\varphi}{\Delta(\varphi)}, \quad \Delta(\psi)=\frac{b}{\Delta(\varphi)}. \]

\subsection*{(19).}

Etant données deux fonctions $F(\varphi)$, $F(\psi)$, si on veut connaître

\origpage{24}
une troisième fonction $F(\mu)$ qui soit égale à leur différence, ensorte qu'on ait
\[ F(\varphi)-F(\psi)=F(\mu), \]
la solution de ce problème se déduira aisément de celle du problème précédent, et on aura pour résultat les formules
\[ \begin{aligned} \sin\mu &= \frac{\sin\varphi\cos\psi\,\Delta(\psi)-\sin\psi\cos\varphi\,\Delta(\varphi)}{1-c^{2}\sin^{2}\varphi\sin^{2}\psi} \\ \cos\mu &= \frac{\cos\varphi\cos\psi+\sin\varphi\sin\psi\,\Delta(\varphi)\,\Delta(\psi)}{1-c^{2}\sin^{2}\varphi\sin^{2}\psi} \\ \Delta(\mu) &= \frac{\Delta(\varphi)\,\Delta(\psi)+c^{2}\sin\varphi\sin\psi\cos\varphi\cos\psi}{1-c^{2}\sin^{2}\varphi\sin^{2}\psi} \\ \operatorname{tang}\mu &= \frac{\operatorname{tang}\varphi\,\Delta(\psi)-\operatorname{tang}\psi\,\Delta(\varphi)}{1+\operatorname{tang}\varphi\operatorname{tang}\psi\,\Delta(\varphi)\,\Delta(\psi)}. \end{aligned} \]
Si dans cette dernière formule on fait $\operatorname{tang}\varphi'=\operatorname{tang}\varphi\,\Delta(\psi)$, et $\operatorname{tang}\psi'=\operatorname{tang}\psi\,\Delta(\varphi)$, on aura $\mu=\varphi'-\psi'$.

Ces formules pour la différence se déduiraient des formules pour la somme, en changeant simplement dans celles-ci le signe de $\psi$, et conservant $\Delta(\psi)$ positive. Il est inutile d'ailleurs de faire observer l'analogie qui règne entre les valeurs de $\sin\mu$, $\cos\mu$, et celles de $\sin(\varphi\pm\psi)$, $\cos(\varphi\pm\psi)$; elles coïncideraient entièrement si l'on avait $c=0$.

\subsection*{(20).}

Puisqu'on connaît algébriquement l'amplitude de la fonction égale à la somme ou à la différence de deux fonctions données, il est clair qu'on peut trouver algébriquement une fonction multiple d'une fonction donnée, et qu'en général on peut résoudre sur la multiplication et la division des fonctions elliptiques de la première espèce, les mêmes problèmes qu'on résout sur la multiplication et la division des arcs de cercle. Désignons par $\varphi_{n}$ l'amplitude de la fonction qui contient $n$ fois la fonction dont l'amplitude est $\varphi$, ensorte qu'on ait $F(\varphi_{n})=nF(\varphi)$; il s'agirait d'avoir l'expression générale de $\sin$ et $\cos\varphi_{n}$, par le moyen de $\sin$ et $\cos\varphi$.

Les cas extrêmes n'ont aucune difficulté. Si on a $c=0$, alors $\varphi_{n}=n\varphi$, et l'expression de $\sin\varphi_{n}$, ainsi que celle de $\cos\varphi_{n}$, se déduisent de la formule connue
\[ \cos\varphi_{n}+\sqrt{-1}\cdot\sin\varphi_{n}=(\cos\varphi+\sqrt{-1}\sin\varphi)^{n}. \]

\origpage{25}
Si on a $c=1$, alors la fonction $F(\varphi)$ devient $\displaystyle\int\frac{d\varphi}{\cos\varphi}$, de sorte qu'on a $F(\varphi)=\frac{1}{2}\log\left(\dfrac{1+\sin\varphi}{1-\sin\varphi}\right)$, et la formule pour la multiplication des fonctions est
\[ \frac{1+\sin\varphi_{n}}{1-\sin\varphi_{n}}=\left(\frac{1+\sin\varphi}{1-\sin\varphi}\right)^{n}. \]
La difficulté est de trouver l'expression générale de $\sin\varphi_{n}$, lorsque $c$ a une valeur quelconque entre $0$ et $1$.

\subsection*{(21).}

Considérons d'abord le cas le plus simple, qui est celui de la duplication; il suffira de faire $\psi=\varphi$ dans les formules de l'article 18, ce qui donnera
\[ \begin{aligned} \sin\varphi_{2} &= \frac{2\sin\varphi\cos\varphi\,\Delta(\varphi)}{1-c^{2}\sin^{4}\varphi} \\ \cos\varphi_{2} &= \frac{1-2\sin^{2}\varphi+c^{2}\sin^{4}\varphi}{1-c^{2}\sin^{4}\varphi} \\ \Delta(\varphi_{2}) &= \frac{1-2c^{2}\sin^{2}\varphi+c^{2}\sin^{4}\varphi}{1-c^{2}\sin^{4}\varphi} \\ \operatorname{tang}\tfrac{1}{2}\varphi_{2} &= \operatorname{tang}\varphi\,\Delta(\varphi). \end{aligned} \]
La dernière de ces formules se déduirait immédiatement de la considération du triangle sphérique $ABC$ qui, étant isoscèle dans le cas dont il s'agit, donne $\operatorname{tang}\frac{1}{2}BA=\cos B\operatorname{tang}BC=\operatorname{tang}\varphi\,\Delta(\varphi)$.

Ainsi en prenant l'auxiliaire $B$, telle que $\sin B=c\sin\varphi$, on aura $\varphi_{2}$ par l'équation $\operatorname{tang}\frac{1}{2}\varphi_{2}=\cos B\operatorname{tang}\varphi$; on trouvera semblablement $\varphi_{4}$ au moyen de $\varphi_{2}$, $\varphi_{8}$ au moyen de $\varphi_{4}$, etc., de sorte que la fonction $F$ sera multipliée par 2, 4, 8, 16, etc.

Réciproquement, étant donné $\varphi_{2}$, on trouvera $\varphi$ par l'équation
\[ \sin\varphi=\frac{\sin\frac{1}{2}\varphi_{2}}{\sqrt{[\frac{1}{2}+\frac{1}{2}\Delta(\varphi_{2})]}}, \]
ou bien, appelant par analogie $\varphi_{\frac{1}{2}}$ l'amplitude de la fonction qui est égale à la moitié de $F(\varphi)$, on aura
\[ \sin\varphi_{\frac{1}{2}}=\frac{\sin\frac{1}{2}\varphi}{\sqrt{(\frac{1}{2}+\frac{1}{2}\Delta)}}. \]
Si l'on fait $c\sin\varphi=\sin\beta$, ce qui donne $\Delta=\cos\beta$, on aura plus

\origpage{26}
simplement
\[ \sin\varphi_{\frac{1}{2}}=\frac{\sin\frac{1}{2}\varphi}{\cot\frac{1}{2}\beta}, \]
\ednote{Printed $\cot\frac{1}{2}\beta$ in the denominator; since $\sqrt{(\frac{1}{2}+\frac{1}{2}\cos\beta)}=\cos\frac{1}{2}\beta$, the expected reading is $\cos\frac{1}{2}\beta$. An apparent misprint.} et on satisfera ainsi à l'équation $F(\varphi_{\frac{1}{2}})=\frac{1}{2}F(\varphi)$.

On trouvera semblablement les amplitudes $\varphi_{\frac{1}{4}}$, $\varphi_{\frac{1}{8}}$, etc., par lesquelles la bisection de la fonction $F(\varphi)$ peut être continuée indéfiniment.

\subsection*{(22).}

Venons à la multiplication par un nombre quelconque. Pour cela, considérons trois amplitudes consécutives $\varphi_{n-1}$, $\varphi_{n}$, $\varphi_{n+1}$, qui, suivant l'indication, répondent aux fonctions multiples $(n-1)F$, $nF$, $(n+1)F$. Les formules pour la somme et la différence de deux fonctions, s'appliqueront aux équations $F(\varphi_{n+1})=F(\varphi_{n})+F(\varphi)$, $F(\varphi_{n-1})=F(\varphi_{n})-F(\varphi)$, et il en résultera
\[ \begin{aligned} \sin\varphi_{n+1}+\sin\varphi_{n-1} &= \frac{2\Delta\cos\varphi\sin\varphi_{n}}{1-c^{2}\sin^{2}\varphi\sin^{2}\varphi_{n}} \\ \cos\varphi_{n+1}+\cos\varphi_{n-1} &= \frac{2\cos\varphi\sin\varphi_{n}}{1-c^{2}\sin^{2}\varphi\sin^{2}\varphi_{n}}. \end{aligned} \]
\ednote{In the second formula the numerator is printed $2\cos\varphi\sin\varphi_{n}$; the sum of the cosines requires $2\cos\varphi\cos\varphi_{n}$, as the division on p. 27 confirms. An apparent misprint.} Ces formules où $\Delta$ et $\varphi$ restent constamment les mêmes, tandis que $n$ varie, paraissent aussi commodes qu'il est possible pour en tirer les valeurs successives de $\sin\varphi_{2}$, $\sin\varphi_{3}$, $\cos\varphi_{2}$, $\cos\varphi_{3}$, etc.

Soit, pour abréger, $2\Delta\cos\varphi=p$, $1-c^{2}\sin^{4}\varphi=q$, $c^{2}\sin^{4}\varphi=r$, on trouvera pour la suite des sinus,
\[ \begin{aligned} \sin\varphi_{2} &= \frac{p}{q}\sin\varphi \\ \sin\varphi_{3} &= \frac{p^{2}-q^{2}}{q^{2}-rp^{2}}\sin\varphi \\ \sin\varphi_{4} &= \frac{(1+r)p^{2}-2q^{2}}{q^{4}-rp^{4}}\,pq\sin\varphi \\ \sin\varphi_{5} &= \frac{q^{6}-3p^{2}q^{4}+(1+2r)p^{4}q^{2}-r^{2}p^{6}}{q^{6}-3rp^{2}q^{4}+r(r+2)p^{4}q^{2}-rp^{6}}\sin\varphi \\ &\text{etc.}, \end{aligned} \]
et pour celle des cosinus, on a, en faisant de plus $s=1-2\sin^{2}\varphi+r=2\cos^{2}\varphi-q$,

\origpage{27}
\[ \begin{aligned} \cos\varphi_{2} &= \frac{1-2\sin^{2}\varphi+r}{q}=\frac{s}{q} \\ \cos\varphi_{3} &= \frac{2sq-q^{2}+p^{2}r}{q^{2}-rp^{2}}\cos\varphi \\ \cos\varphi_{4} &= \frac{q^{4}-2q^{2}p^{2}\sin^{2}\varphi+rp^{4}}{q^{4}-rp^{4}} \\ \cos\varphi_{5} &= \left(\frac{2sq(q^{2}-rp^{2})(2sq-q^{2}+rp^{2})}{q^{6}-3rp^{2}q^{4}+(r^{2}+2r)p^{4}q^{2}-rp^{6}}-1\right)\cos\varphi \\ &\text{etc.} \end{aligned} \]
Mais ces expressions étant entièrement développées, deviennent fort prolixes, et leur loi est très-difficile à apercevoir.

Lorsqu'il s'agira de calculer trigonométriquement $\varphi_{n}$ par le moyen de $\varphi$, on y parviendra aisément de cette manière. Les deux formules générales étant divisées l'une par l'autre donneront
\[ \frac{\sin\varphi_{n+1}+\sin\varphi_{n-1}}{\cos\varphi_{n+1}+\cos\varphi_{n-1}}=\Delta\operatorname{tang}\varphi_{n}, \]
ce qui se réduit à cette formule très-simple,
\[ \operatorname{tang}(\tfrac{1}{2}\varphi_{n+1}+\tfrac{1}{2}\varphi_{n-1})=\Delta\operatorname{tang}\varphi_{n}; \]
d'où l'on tire successivement
\[ \begin{aligned} &\operatorname{tang}\tfrac{1}{2}\varphi_{2}=\Delta\operatorname{tang}\varphi \\ &\operatorname{tang}(\tfrac{1}{2}\varphi_{3}+\tfrac{1}{2}\varphi)=\Delta\operatorname{tang}\varphi_{2} \\ &\operatorname{tang}(\tfrac{1}{2}\varphi_{4}+\tfrac{1}{2}\varphi_{2})=\Delta\operatorname{tang}\varphi_{3} \\ &\qquad\qquad\text{etc.} \end{aligned} \]
On connaîtra ainsi successivement par un calcul fort simple, les valeurs de $\varphi_{2}$, $\varphi_{3}$, etc. On pourrait aussi procéder par de plus grands intervalles pour arriver plus tôt à un terme éloigné; car on a semblablement
\[ \operatorname{tang}(\tfrac{1}{2}\varphi_{n+i}+\tfrac{1}{2}\varphi_{n-i})=\Delta i\operatorname{tang}\varphi_{n}, \]
$\Delta i$ étant le $\Delta$ qui répond à l'amplitude $\varphi i$.

\subsection*{(23).}

La division d'une fonction elliptique donnée de première espèce en un certain nombre de parties égales, est un problème

\origpage{28}
algébrique qu'on résoudra par le développement des formules qui servent à la multiplication. On a déjà vu les formules pour diviser par 2 ou par une puissance de 2. Supposons qu'on veuille diviser en trois parties égales la fonction $F$; on appellera $\varphi_{3}$ l'amplitude de la fonction donnée, et $\varphi$ celle de la fonction qui en est le tiers. Faisant donc $\sin\varphi_{3}=a$ et $\sin\varphi=x$, on aura pour déterminer $x$, l'équation
\[ a=\frac{3x-4(1+c^{2})x^{3}+6c^{2}x^{5}-c^{4}x^{9}}{1-6c^{2}x^{4}+4c^{2}(1+c^{2})x^{6}-3c^{4}x^{8}}, \]
équation qui est du neuvième degré. Elle serait du degré vingt-cinquième pour la quintisection, et ainsi de suite.

Les équations sont moins élevées de moitié lorsqu'il s'agit de diviser la fonction complète $F^{1}$ en un nombre impair de parties; il en est à cet égard de la division des fonctions $F$, comme de celle des arcs de cercle; tandis que la division d'un arc quelconque en $n$ parties égales, exige la résolution d'une équation du $n^{\text{me}}$ degré, celle du quart de circonférence n'exige que la résolution d'une équation du degré $\dfrac{n-1}{2}$.

Supposons en général $\varphi_{n}=\frac{1}{2}\pi$, afin qu'on ait $F(\varphi)=\dfrac{1}{n}F^{1}$, on aura donc aussi $F(\varphi_{n-i})+F(\varphi_{i})=nF(\varphi)=F^{1}$; ce qui donne la formule $\operatorname{tang}(\varphi_{n-i})=\dfrac{1}{b}\cot\varphi_{i}$, d'où l'on déduit
\[ \begin{gathered} \operatorname{tang}\varphi_{n-1}=\frac{1}{b}\cot\varphi \\ \operatorname{tang}\varphi_{n-2}=\frac{1}{b}\cot\varphi_{2} \\ \operatorname{tang}\varphi_{n-3}=\frac{1}{b}\cot\varphi_{3} \\ \text{etc.} \end{gathered} \]
Mais à cause de $\varphi_{n}=\frac{1}{2}\pi$, on a
\[ \begin{gathered} \operatorname{tang}(\tfrac{1}{4}\pi+\tfrac{1}{2}\varphi_{n-2})=\Delta\operatorname{tang}\varphi_{n-1}=\frac{\Delta}{b}\cot\varphi \\ \operatorname{tang}(\tfrac{1}{2}\varphi_{n-1}+\tfrac{1}{2}\varphi_{n-3})=\Delta\operatorname{tang}\varphi_{n-2} \\ \text{etc.} \end{gathered} \]
De là résultent des formules assez simples pour déterminer $\varphi_{n-1}$, $\varphi_{n-2}$, etc.; développant celles qui donnent $\varphi_{2}$, $\varphi_{3}$, etc., on aura

\origpage{29}
par la rencontre de ces deux suites, l'équation qui doit déterminer $\varphi$, et le calcul sera moins compliqué que par le développement de $\sin\varphi_{n}$ ou de $\cos\varphi_{n}$.

Lorsque $n=3$, on aura immédiatement $\operatorname{tang}(\frac{1}{4}\pi+\frac{1}{2}\varphi)=\dfrac{\Delta}{b}\cot\varphi$, ou $b\sin\varphi=\Delta(1-\sin\varphi)$. Soit $\sin\varphi=x$, et l'équation pour déterminer $x$ sera
\[ 0=1-2x+2c^{2}x^{3}-c^{2}x^{4}: \]
c'est l'équation qui donne la trisection de la fonction $F^{1}$, et on voit que son degré $=\dfrac{9-1}{2}$.

Lorsque $n=5$, il faudra éliminer $\varphi_{3}$ des deux équations
\[ \begin{gathered} \operatorname{tang}(\tfrac{1}{4}\pi+\tfrac{1}{2}\varphi_{3})=\frac{\Delta}{b}\cot\varphi \\ \operatorname{tang}(\tfrac{1}{2}\varphi_{3}+\tfrac{1}{2}\varphi)=\Delta\operatorname{tang}\varphi_{2}, \end{gathered} \]
et ensuite mettre au lieu de $\operatorname{tang}\varphi_{2}$ sa valeur $\dfrac{2\Delta\sin\varphi\cos\varphi}{1-2\sin^{2}\varphi+c^{2}\sin^{4}\varphi}$; on obtiendra ainsi, en faisant $\sin\varphi=x$,
\[ \frac{1+x}{bx}\sqrt{(1-c^{2}x^{2})}=\frac{1+2x-2c^{2}x^{3}-c^{2}x^{4}}{1-2x+2c^{2}x^{3}-c^{2}x^{4}}, \]
équation pour la quintisection de la fonction $F^{1}$, laquelle étant entièrement développée, montera au degré $12=\dfrac{25-1}{2}$.

\subsection*{(24).}

Revenons à l'équation de la trisection; elle offre dans sa résolution quelques particularités qui méritent d'être remarquées.

Soit d'abord $x=y+\frac{1}{2}$, cette équation deviendra
\[ y^{4}-\tfrac{3}{2}y^{2}+\left(\frac{2}{c^{2}}-1\right)y-\tfrac{3}{16}=0. \]
Supposons que le premier membre soit le produit des deux facteurs $y^{2}-py+q$, $y^{2}+py-r$, on aura pour déterminer $p$, $q$, $r$, les équations $q-r=p^{2}-\frac{3}{2}$, $qr=\frac{3}{16}$, $p(q+r)=\dfrac{2}{c^{2}}-1$. Des deux premières on tire $(q+r)^{2}=p^{4}-3p^{2}+3$, et cette valeur étant

\origpage{30}
substituée dans la troisième, il en résulte $p^{6}-3p^{4}+3p^{2}=\left(\dfrac{2}{c^{2}}-1\right)^{2}$, ou $(p^{2}-1)^{3}=\dfrac{4b^{2}}{c^{4}}$. Soit donc
\[ k=\sqrt[3]{\left(\frac{4b^{2}}{c^{4}}\right)}, \]
et on aura $p^{2}-1=k$, ou $p=\sqrt{(1+k)}$; d'où résulte
\[ \begin{aligned} q &= \tfrac{1}{2}k-\tfrac{1}{4}+\tfrac{1}{2}\sqrt{(k^{2}-k+1)} \\ r &= \tfrac{1}{2}k-\tfrac{1}{4}-\tfrac{1}{2}\sqrt{(k^{2}-k+1)}. \end{aligned} \]
Ainsi les quatre valeurs de $y$ seront
\[ \begin{aligned} y &= \tfrac{1}{2}\sqrt{(1+k)}\pm\tfrac{1}{2}\sqrt{[2-k-2\sqrt{(1-k+k^{2})}]} \\ y &= -\tfrac{1}{2}\sqrt{(1+k)}\pm\tfrac{1}{2}\sqrt{[2-k+2\sqrt{(1-k+k^{2})}]}. \end{aligned} \]
Les deux premières sont imaginaires, et des deux autres, il n'y a que la racine positive qui convienne à la question, ainsi on aura $y+\frac{1}{2}$, ou
\[ \sin\varphi=\tfrac{1}{2}-\tfrac{1}{2}\sqrt{(1+k)}+\tfrac{1}{2}\sqrt{[2-k+2\sqrt{(1-k+k^{2})}]}; \]
de sorte que cette valeur pourra se construire géométriquement, lorsque $k$ sera donné. Or, quel que soit $k$, on a
\[ c^{2}=\frac{2}{1+\sqrt{(1+k^{3})}}, \]
valeur qui est toujours comprise entre les limites 1 et 0.

Soit $k=1$, on aura
\[ \begin{gathered} c^{2}=\frac{2}{1+\sqrt{2}}=2(\sqrt{2}-1) \quad\text{et} \\ \sin\varphi=\tfrac{1}{2}-\tfrac{1}{2}\sqrt{2}+\tfrac{1}{2}\sqrt{3}. \end{gathered} \]

Soit $k=2$, on aura
\[ \begin{gathered} c^{2}=\tfrac{1}{2}=b^{2} \quad\text{et} \\ \sin\varphi=\tfrac{1}{2}-\tfrac{1}{2}\sqrt{3}+\sqrt{(\tfrac{1}{2}\sqrt{3})} \\ \cos^{2}\varphi=(-1+\sqrt{3})\sqrt{(\tfrac{1}{2}\sqrt{3})}=\sqrt{(2\sqrt{3}-3)}. \end{gathered} \]
\ednote{The sign before the 1 in $(-1+\sqrt{3})$ is faint and partly obscured by a reader's pencil stroke through the following $+$; it is read as $-1$, which the value of $\cos^{2}\varphi$ requires.} Voilà deux cas dans lesquels la trisection de la fonction $F^{1}$ se fait par de simples extractions de racines carrées.

Les applications que nous aurons occasion de faire par la suite,

\origpage{31}
exigent que nous considérions encore le cas de $c=\frac{1}{2}\sqrt{(2+\sqrt{3})}=\cos\dfrac{\pi}{12}$, et celui de $c=\frac{1}{2}\sqrt{(2-\sqrt{3})}=\sin\dfrac{\pi}{12}$.

Soit d'abord $c=\frac{1}{2}\sqrt{(2+\sqrt{3})}$, on aura $k^{3}=\dfrac{4b^{2}}{c^{4}}=\dfrac{4b^{2}c^{2}}{c^{6}}=\dfrac{1}{4c^{6}}$; donc $k=\dfrac{1}{2c^{2}}\sqrt[3]{2}$; d'après cette valeur on trouvera
\[ \sin\varphi=\sqrt{3}-1. \]
Mais comme les réductions pour parvenir à ce résultat seraient fort pénibles, il est plus simple de revenir à l'équation primitive
\[ 0=1-2x+\frac{2+\sqrt{3}}{4}(2x^{3}-x^{4}), \]
et on vérifiera facilement que cette équation est satisfaite en faisant $x=\sqrt{3}-1$; alors on a $\sin^{2}\varphi=4-2\sqrt{3}$, $\cos^{2}\varphi=2\sqrt{3}-3$, et $\operatorname{tang}^{2}\varphi=\dfrac{2}{\sqrt{3}}$, ou $\operatorname{tang}\varphi=\sqrt{\left(\dfrac{2}{\sqrt{3}}\right)}$. C'est la valeur de $\varphi$ qui donne $F(\varphi)=\frac{1}{3}F^{1}$, lorsque $c=\cos\dfrac{\pi}{12}$.

Soit maintenant $c=\frac{1}{2}\sqrt{(2-\sqrt{3})}$, on aura encore $k^{3}=\dfrac{1}{4c^{6}}$, et $k=\dfrac{1}{2c^{2}}\sqrt[3]{2}=(4+2\sqrt{3})\sqrt[3]{2}$.
La substitution de cette valeur dans la formule générale, fera connaître $\sin\varphi$; mais pour parvenir au résultat le plus simple, voici la route qu'il faut tenir.

Je reprends l'équation générale $0=1-2x+c^{2}(2x^{3}-x^{4})$; je fais $x^{2}=1-y$, ce qui revient à supposer $y=\cos^{2}\varphi$, la transformée sera
\[ 0=c^{4}y^{4}+6b^{2}c^{2}y^{2}+4b^{2}(b^{2}-c^{2})y-3b^{4}. \]
Soit $y=\dfrac{b}{c}z$, on aura de nouveau
\[ 0=z^{4}+6z^{2}+\left(\frac{4b}{c}-\frac{4c}{b}\right)z-3. \]
Dans le cas présent, on a $\dfrac{b}{c}-\dfrac{c}{b}=\dfrac{b^{2}-c^{2}}{bc}=4(b^{2}-c^{2})=2\sqrt{3}$; de sorte que l'équation devient
\[ 0=z^{4}+6z^{2}+8\sqrt{3}\cdot z-3. \]

\origpage{32}
Cette équation étant divisée par $z+\sqrt{3}$, facteur inutile pour notre objet, le quotient est
\[ 0=z^{3}-z^{2}\sqrt{3}+9z-\sqrt{3}. \]
Soit enfin $z=\dfrac{1+2v}{\sqrt{3}}$, et la transformée sera
\[ 0=v^{3}+6v+2, \]
d'où l'on tire $v=\sqrt[3]{2}-\sqrt[3]{4}$. Soit $\sqrt[3]{4}=m$, et on aura $v=\frac{1}{2}m^{2}-m$; donc $z=\dfrac{m^{2}-2m+1}{\sqrt{3}}$, $y=\dfrac{b}{c}z=(2+\sqrt{3})z=\dfrac{m^{2}-2m+1}{2\sqrt{3}-3}=\cos^{2}\varphi$; donc enfin $\cos\varphi=\dfrac{m-1}{\sqrt{(2\sqrt{3}-3)}}=(m-1)\sqrt{\left(\dfrac{2+\sqrt{3}}{\sqrt{3}}\right)}$. C'est l'expression la plus simple par laquelle on peut déterminer $\varphi$ pour satisfaire à l'équation $F(\varphi)=\frac{1}{3}F^{1}$, dans le cas de $c=\frac{1}{2}\sqrt{(2-\sqrt{3})}=\sin\dfrac{\pi}{12}$.

Pour effectuer entièrement la trisection de la fonction $F^{1}$, il faut encore avoir la valeur de $\varphi_{2}$ qui donne $F(\varphi_{2})=\frac{2}{3}F^{1}$: or $\varphi_{2}$ est facile à déduire de $\varphi$, soit par les formules de la duplication, soit par les formules particulières
\[ \begin{gathered} \operatorname{tang}\varphi_{2}=\frac{\cot\varphi}{b} \\ \cos\varphi_{2}=1-\sin\varphi. \end{gathered} \]
La première résulte de l'équation $b\operatorname{tang}\varphi\operatorname{tang}\psi=1$ qui a lieu (art.~18) lorsque $F(\varphi)+F(\psi)=F^{1}$; la seconde résulte de la combinaison des équations $\operatorname{tang}\varphi_{2}=\dfrac{\cot\varphi}{b}$, $\operatorname{tang}\frac{1}{2}\varphi_{2}=\Delta(\varphi)\operatorname{tang}\varphi=\dfrac{1-\cos\varphi_{2}}{\sin\varphi_{2}}$, et $b\sin\varphi=\Delta(1-\sin\varphi)$. La formule $\cos\varphi_{2}=1-\sin\varphi$ est d'autant plus remarquable qu'elle ne contient point le module $c$, et qu'ainsi elle a lieu quel que soit le module. On aura donc aussi rationnellement $\sin\varphi_{2}=\dfrac{\cot\varphi}{b}(1-\sin\varphi)$.

La valeur de $\sin\varphi_{2}$ peut être déterminée directement par la trisection de la fonction $F(\pi)$, car il est évident qu'on a $F(\varphi_{2})=\frac{2}{3}F^{1}=\frac{1}{3}F(\pi)$. Soit donc $\sin\varphi_{2}=y$, et en faisant dans la formule de la triplication (art.~23) $a=\sin\pi=0$, on aura pour déterminer $y$, l'équation
\[ 0=3-4(1+c^{2})y^{2}+6c^{2}y^{4}-c^{4}y^{8}, \]
qui n'a que la difficulté du quatrième degré.

\origpage{33}
\subsection*{(25).}

Occupons-nous maintenant de l'équation pour la quintisection, qui étant entièrement développée, devient
\[ \begin{aligned} 0={}&1-4x^{2}+5c^{2}x^{4}+4b^{2}c^{2}x^{6}-5c^{4}x^{8}+4c^{6}x^{10}-c^{6}x^{12} \\ &-2x[1-5c^{2}x^{2}+(6c^{2}+4c^{4})x^{4}-(4c^{2}+6c^{4})x^{6}+5c^{4}x^{8}-c^{6}x^{10}]. \end{aligned} \]

Pour donner au moins quelques exemples particuliers de la résolution de cette équation, on pourrait attribuer une valeur à $x$ (en ayant soin qu'elle fût plus grande que la racine de l'équation $0=1-2x-4x^{2}$); et d'après la valeur supposée de $x$, on voit que $c^{2}$ se déterminerait par une équation du troisième degré.

Mais la recherche des cas particuliers de solution, est une question d'analyse indéterminée qu'on peut résoudre plus facilement de la manière suivante.

Soit $p=1-c^{2}x^{4}$, $q=2x(1-c^{2}x^{2})$, et l'équation de l'article 23 sera
\[ \frac{p+q}{p-q}=\frac{(1+x)\sqrt{(1-c^{2}x^{2})}}{bx}. \]
Élevant chaque membre au quarré, et retranchant de part et d'autre l'unité, on aura
\[ \frac{4pq}{(p-q)^{2}}=\frac{p+q}{b^{2}x^{2}}. \]
Soit $p=mq$, cette équation donnera $b^{2}x^{2}=\dfrac{(m+1)(m-1)^{2}}{4m}q$; soit donc encore
\[ n=\frac{(m+1)(m-1)^{2}}{2m}, \]
et on aura, en remettant la valeur de $q$, $b^{2}x=n(1-c^{2}x^{2})$, ce qui donne
\[ c^{2}=\frac{x-n}{x-nx^{2}}. \]
Mais de l'équation $p=mq$, ou $1-c^{2}x^{4}=2mx(1-c^{2}x^{2})$, on tire
\[ c^{2}=\frac{2mx-1}{2mx^{3}-x^{4}}. \]
Égalant ces deux valeurs de $c^{2}$, il viendra $0=1-(2m+n)x+x^{2}$.

De là se déduit une solution générale fort simple du problème que nous nous sommes proposé. Ayant pris pour $m$ une valeur quelconque

\origpage{34}
comprise entre 1 et $\dfrac{\sqrt{5}+1}{2}=1.618$ (car ces valeurs répondent aux limites $c=1$ et $c=0$), on en déduit $n=\dfrac{(m+1)(m-1)^{2}}{2m}$; ensuite
\[ \begin{aligned} x &= m+\tfrac{1}{2}n-\sqrt{[(m+\tfrac{1}{2}n)^{2}-1]}, \\ c^{2} &= \frac{1-n(m+\frac{1}{2}n)-n\sqrt{[(m+\frac{1}{2}n)^{2}-1]}}{1-n(m+\frac{1}{2}n)+n\sqrt{[(m+\frac{1}{2}n)^{2}-1]}}. \end{aligned} \]

Soit, par exemple, $m=\dfrac{3}{2}$, on aura $n=\dfrac{5}{24}$, de là
\[ \begin{aligned} x &= \frac{77-5\sqrt{145}}{48} \\ c^{2} &= \frac{767-25\sqrt{145}}{767+25\sqrt{145}}. \end{aligned} \]
D'après cette valeur du module $c$ et celle de $x=\sin\varphi$, on satisfera à l'équation $F(\varphi)=\frac{1}{5}F^{1}$; ensuite il sera facile d'avoir les valeurs de $\varphi_{2}$, $\varphi_{3}$, $\varphi_{4}$, par lesquelles on achevera de diviser en cinq parties égales la fonction complète $F^{1}$.

\subsection*{(26).}

Nous avons fait voir comment on trouve la relation entre $\varphi_{n}$ et $\varphi$, pour que $F(\varphi_{n})=nF(\varphi)$, $n$ étant un nombre entier. S'il fallait trouver la relation entre $\varphi$ et $\psi$ pour que $F(\psi)=\dfrac{m}{n}F(\varphi)$, $m$ et $n$ étant des entiers, on prendrait un angle auxiliaire $\omega$, tel qu'on eût à la fois $nF(\psi)=F(\omega)$ et $mF(\varphi)=F(\omega)$. La première condition donnerait une équation algébrique entre les sinus des angles $\psi$ et $\omega$, la seconde en donnerait une entre ceux des angles $\varphi$ et $\omega$; d'où éliminant $\omega$, on aura la relation cherchée entre $\varphi$ et $\psi$.

En général on voit qu'il sera toujours possible de satisfaire par une équation algébrique à l'équation transcendante
\[ 0=mF(\varphi)+nF(\psi)+pF(\omega)+\text{etc.}, \]
les coefficiens $m$, $n$, $p$, etc. étant des nombres entiers à volonté, qui n'aient pas tous le même signe. C'est une conséquence toute simple de ce que l'équation transcendante $F(\varphi)+F(\psi)=F(\mu)$ est représentée par une équation algébrique.

\subsection*{(27).}

De là il résulte qu'on peut toujours trouver l'intégrale algébrique

\origpage{35}
complète de l'équation différentielle
\[ \frac{m\,d\varphi}{\sqrt{(1-c^{2}\sin^{2}\varphi)}}\pm\frac{n\,d\psi}{\sqrt{(1-c^{2}\sin^{2}\psi)}}=0, \]
$m$ et $n$ étant des nombres entiers; car l'intégrale est d'abord $mF(\varphi)\pm nF(\psi)=\text{const.}$ Et si on suppose que lorsque $\psi=0$, on ait $\varphi=\mu$, la constante sera $mF(\mu)$, et on aura l'intégrale
\[ mF(\varphi)\pm nF(\psi)=mF(\mu). \]
Soit $\omega$ une auxiliaire telle que $F(\mu)=F(\varphi)\pm F(\omega)$, on aura en même temps $mF(\omega)=nF(\psi)$; si on exprime ensuite ces deux équations en termes algébriques, et qu'on élimine $\omega$, on aura l'intégrale algébrique cherchée dans laquelle $\mu$ sera la constante arbitraire.

En général;\ednote{Printed with a semicolon after ``général''; a comma is expected, an apparent misprint.} si on avait l'équation suivante, dans laquelle $m$, $n$, $p$, etc. sont des entiers positifs ou négatifs,
\[ 0=\frac{m\,d\varphi}{\sqrt{(1-c^{2}\sin^{2}\varphi)}}+\frac{n\,d\psi}{\sqrt{(1-c^{2}\sin^{2}\psi)}}+\frac{p\,d\omega}{\sqrt{(1-c^{2}\sin^{2}\omega)}}+\text{etc.}, \]
l'intégrale complète sera $F(\mu)=mF(\varphi)+nF(\psi)+pF(\omega)+\text{etc.}$, $\mu$ étant la constante arbitraire, et cette intégrale pourra toujours être représentée par une équation algébrique, quel que soit le nombre des termes, pourvu qu'il ne soit pas infini.

Si on désigne par $R(x)$ le radical $\sqrt{(\alpha+\beta x+\gamma x^{2}+\delta x^{3}+\varepsilon x^{4})}$, et par $R(y)$, $R(z)$, etc. des radicaux semblables en $y$, $z$, etc.; si de plus $m$, $n$, $p$, etc. désignent des nombres entiers positifs ou négatifs, il est clair que l'équation
\[ 0=\frac{m\,dx}{R(x)}+\frac{n\,dy}{R(y)}+\frac{p\,dz}{R(z)}+\text{etc.}, \]
pourra toujours être réduite à la forme précédente, et qu'ainsi elle aura toujours une intégrale algébrique complète.

Rien n'empêcherait de supposer que $z$ et les variables suivantes fussent des fonctions algébriques données de $x$ et $y$, alors l'équation précédente ne renfermerait que deux variables, et malgré l'infinité de formes dont elle est susceptible, elle admettrait toujours une intégrale algébrique complète. C'est peut-être la seule manière de

\origpage{36}
généraliser le résultat d'Euler, concernant l'équation $\dfrac{dx}{R(x)}+\dfrac{dy}{R(y)}=0$. Lagrange s'est proposé de trouver des cas d'intégrabilité de l'équation $\dfrac{dx}{\sqrt{X}}+\dfrac{dy}{\sqrt{Y}}=0$, sans supposer que les deux polynomes $X$ et $Y$ sont entièrement semblables${}^{(1)}$; mais il ne paraît pas que les recherches de ce grand Géomètre l'aient conduit au-delà de l'équation d'Euler; car l'équation qu'il donne page 119, comme étant plus générale, s'y ramène immédiatement en faisant $v=ky$, et donnant au coefficient $k$ une valeur convenable. Ainsi il est très-douteux qu'avec deux termes seulement, l'équation d'Euler puisse être généralisée; mais avec un plus grand nombre, on voit qu'elle admet une grande extension.

\subsection*{(28).}

Puisque les fonctions $F$ peuvent être multipliées ou divisées à volonté, cette propriété fournit un moyen assez simple de les évaluer par approximation. D'abord nous supposerons que $\varphi$ ne surpasse pas $\frac{1}{2}\pi$, car suivant l'article 12, tous les cas se réduisent à celui-là.

Cela posé, on déterminera $\varphi_{\frac{1}{2}}$ par l'article 21, de manière que $F(\varphi_{\frac{1}{2}})=\frac{1}{2}F(\varphi)$, et l'intégrale $\displaystyle\int\frac{d\varphi_{\frac{1}{2}}}{\Delta(\varphi_{\frac{1}{2}})}$ aura une étendue moindre que $\displaystyle\int\frac{d\varphi}{\Delta(\varphi)}$, de près de moitié, si $c$ n'est pas trop près de l'unité. Par une seconde bisection, on peut faire ensorte que $F(\varphi_{\frac{1}{4}})=\frac{1}{4}F(\varphi)$, et l'intégrale que représente $F(\varphi_{\frac{1}{4}})$ aura encore une étendue près de deux fois moindre, ainsi de suite. Mais lorsque l'amplitude $\psi$ de la formule qu'on considère est devenue très-petite, la fonction $F(\psi)$ se réduit sensiblement à l'arc $\psi$. Donc quelle que soit la première valeur de $\varphi$, la fonction correspondante $F(\varphi)$ sera égale au dernier terme de la suite $\varphi$, $2\varphi_{\frac{1}{2}}$, $4\varphi_{\frac{1}{4}}$, $8\varphi_{\frac{1}{8}}$, etc., et dans la plupart des cas, on obtiendra cette limite par un calcul assez court.

\textbf{(1)} Mémoires de Turin, tome IV.

\origpage{37}
\section*{EXEMPLE.}

Soit proposé de trouver la valeur de $F$ lorsque $c=\sqrt{\left(\dfrac{2+\sqrt{3}}{4}\right)}=\sin 75^{\circ}$ et $\operatorname{tang}\varphi=\sqrt{\left(\dfrac{2}{\sqrt{3}}\right)}$, on trouvera par les formules de l'art.\ 21, ce qui suit
\[ \begin{aligned} \varphi &= 47^{\circ}\;3'\;30''91 & \qquad \beta &= 45^{\circ}\;0'\;0''00 \\ \varphi_{\frac{1}{2}} &= 25\;\;36\;\;\;5.64 & \beta_{\frac{1}{2}} &= 24\;\;40\;\;10.94 \\ \varphi_{\frac{1}{4}} &= 13\;\;\;6\;\;30.98 & \beta_{\frac{1}{4}} &= 12\;\;39\;\;15.83 \\ \varphi_{\frac{1}{8}} &= \;\;6\;\;35\;\;40.74 & \beta_{\frac{1}{8}} &= \;\;6\;\;22\;\;\;8.40 \\ \varphi_{\frac{1}{16}} &= \;\;3\;\;18\;\;\;8.75 & &\text{etc.} \\ &\text{etc.} & & \end{aligned} \]
\ednote{The first row is printed $30''91$; the same angle is given as $30''94$ on p.~91, and the exact value is $30.95''$.}
Les deux dernières valeurs donnent
\[ \begin{aligned} 8\varphi_{\frac{1}{8}} &= 52^{\circ}\;45'\;25''.92 \\ 16\varphi_{\frac{1}{16}} &= 52\;\;\;50\;\;\;20\;.00 \end{aligned} \]
Leur différence est $4'\;54''.08$; et comme par la nature de ces approximations, un résultat doit approcher de la limite environ quatre fois plus que le précédent, nous ajouterons au dernier résultat le tiers de la différence $4'\;54''.08$, qui est $1'\;38''.03$, et nous aurons ainsi pour la valeur de $F$, l'arc très-approché
\[ 52^{\circ}\;51'\;58''.03, \]
qui en parties du rayon $=\frac{\pi}{2}\times 0.5874013=0.9226878$.

Cette méthode pourrait devenir très-longue lorsque $c$ est fort près de l'unité, et $\varphi$ peu différent de $90^{\circ}$: nous en donnerons ci-après une plus expéditive.

\section*{Application à la Lemniscate.}

\subsection*{(29).}

\rmfigure{figures/fig-4.png}{Fig.~4.}{A lemniscate with its double point C on the axis AB, the points L, M, H, K, O, I on the upper left loop, M' on the lower left loop and N', N on the right loop, with the chord CM and the ordinate MP to the axis.}

La Lemniscate est, comme on sait, une courbe du quatrième ordre qui a pour équation $(x^{2}+y^{2})^{2}=a^{2}(x^{2}-y^{2})$. On suppose le demi-axe $CA=a$, l'abscisse $CP=x$, et l'ordonnée $PM=y$; si de plus on fait la corde $CM=z$, on aura en fonctions de $z$, $x=\dfrac{z}{a}\sqrt{\left(\dfrac{a^{2}+z^{2}}{2}\right)}$, $y=\dfrac{z}{a}\sqrt{\left(\dfrac{a^{2}-z^{2}}{2}\right)}$, d'où résulte l'arc $AM$,

\origpage{38}
\[ s=\int\frac{-a^{2}\,dz}{\sqrt{(a^{4}-z^{4})}}. \]

Soit encore $z=a\cos\varphi$, et $c^{2}=\frac{1}{2}$, on aura
\[ s=\frac{a}{\sqrt{2}}\int\frac{d\varphi}{\sqrt{(1-c^{2}\sin^{2}\varphi)}}=\frac{a}{\sqrt{2}}F(\varphi), \]
résultat auquel on serait parvenu directement en faisant
\[ \begin{aligned} x &= a\Delta\cos\varphi \\ y &= \frac{a}{\sqrt{2}}\sin\varphi\cos\varphi. \end{aligned} \]
Avec ces valeurs où $\Delta$ est toujours pris positivement, suivant notre usage, on suit la courbe dans tout son contour, de cette manière:

Dans le premier quart $AMC$, les valeurs de $\varphi$ s'étendent depuis $\varphi=0$ jusqu'à $\varphi=\frac{1}{2}\pi$; dans le second quart $CNB$, elles s'étendent depuis $\varphi=\frac{1}{2}\pi$ jusqu'à $\varphi=\pi$; dans le troisième quart $BN'C$, depuis $\varphi=\pi$ jusqu'à $\varphi=\frac{3}{2}\pi$, et enfin dans le quatrième $CM'A$, depuis $\varphi=\frac{3}{2}\pi$ jusqu'à $\varphi=2\pi$.

Cela posé, si l'on fait $\dfrac{a}{\sqrt{2}}=1$, on aura $s=F(\varphi)$; mais quelle que soit la ligne qu'on prend pour unité, on voit que les arcs de la Lemniscate jouissent de toutes les propriétés des fonctions elliptiques de première espèce, c'est-à-dire, qu'ils peuvent être ajoutés, retranchés ou divisés algébriquement comme les arcs de cercle. Ainsi étant donné un arc quelconque $MO$, on peut, à compter d'un point donné $H$, déterminer algébriquement un autre arc $HI$ ou $HL$ qui soit dans un rapport rationnel donné avec l'arc $MO$. Il suffit pour cela de déterminer $\varphi$ d'après une équation de la forme $F(\beta)-F(\alpha)=\pm\dfrac{m}{n}[F(\varphi)-F(\gamma)]$, à laquelle correspond toujours une équation algébrique.

A plus forte raison peut-on diviser un arc donné $MO$ en un certain nombre de parties égales. Par exemple, s'il s'agit du quart de la courbe $AMC$, on déterminera son milieu $K$ au moyen de l'équation $\sin^{2}\varphi=\dfrac{1}{1+b}=\dfrac{1}{1+c}$, ou par la corde $CK=\dfrac{1}{c}\cos\varphi=\sqrt{(2\sqrt{2}-2)}$; si l'on veut déterminer l'arc $AM$ égal au tiers de $AOC$, il faudra,

\origpage{39}
d'après l'art.\ 24, faire $\cos\varphi=\sqrt[4]{(2\sqrt{3}-3)}$, ou prendre la corde $CM=\sqrt{[2\sqrt{(2\sqrt{3}-3)}]}$.

Les arcs de la Lemniscate représentent les fonctions elliptiques de la première espèce dans le cas où le module $c=\sqrt{\frac{1}{2}}=b$. Il serait assez curieux de rechercher s'il y a quelque autre courbe algébrique dont les arcs représentent la fonction $F$ pour une autre valeur de $c$; mais cette recherche ne laisse pas d'être difficile.

Elle ne présenterait aucune difficulté si on admettait les arcs de cercle et les logarithmes dans l'expression des coordonnées. En effet il s'agit de satisfaire à l'équation $dx^{2}+dy^{2}=\dfrac{d\varphi^{2}}{1-c^{2}\sin^{2}\varphi}$, laquelle peut être mise sous cette forme
\[ dx^{2}+dy^{2}=\frac{d\varphi^{2}(\cos^{2}\varphi+b^{2}\sin^{2}\varphi)(\cos^{2}\psi+\sin^{2}\psi)}{(1-c^{2}\sin^{2}\varphi^{2})^{2}}; \]
\ednote{The denominator is printed with $\sin^{2}\varphi^{2}$; $(1-c^{2}\sin^{2}\varphi)^{2}$ is meant, an apparent misprint.}
c'est ce qu'on obtient généralement par les valeurs
\[ \begin{aligned} dx &= \frac{d\varphi(\cos\varphi\cos\psi-b\sin\varphi\sin\psi)}{1-c^{2}\sin^{2}\varphi} \\ dy &= -\frac{d\varphi(\cos\varphi\sin\psi+b\sin\varphi\cos\psi)}{1-c^{2}\sin^{2}\varphi}, \end{aligned} \]
dans lesquelles on peut prendre à volonté $\sin\psi$ en fonction de $\sin$ et $\cos\varphi$. Pour nous borner à un cas particulier, soit $\psi=0$, on aura
\[ \begin{aligned} dx &= \frac{d\varphi\cos\varphi}{1-c^{2}\sin^{2}\varphi} \\ dy &= -\frac{b\,d\varphi\sin\varphi}{1-c^{2}\sin^{2}\varphi}, \end{aligned} \]
d'où l'on déduit en intégrant
\[ \begin{aligned} x &= \frac{1}{2c}\log\left(\frac{1+c\sin\varphi}{1-c\sin\varphi}\right) \\ y &= \frac{1}{c}\operatorname{arc}\operatorname{tang}\left(\frac{c\cos\varphi}{b}\right). \end{aligned} \]
La courbe décrite d'après ces deux équations, sera donc telle qu'un arc quelconque $s$, compté depuis $\varphi=0$, aura pour valeur $F(\varphi)$, et représentera généralement une fonction elliptique de première espèce, quel que soit son module.

\origpage{40}
\section*{Application au mouvement du pendule simple.}

\subsection*{(30).}

Les fonctions elliptiques de la première espèce reçoivent une application immédiate dans la détermination du mouvement du pendule simple. Soit $1$ la gravité, $L$ longueur du pendule, $H$ la hauteur due à la vitesse dans le point le plus bas, $\psi$ l'angle dont le pendule est écarté de la verticale au bout du temps $t$, on aura
\[ t=\int\frac{\tfrac{1}{2}d\psi\sqrt{L}}{\sqrt{\left(\dfrac{H}{2L}-\sin^{2}\tfrac{1}{2}\psi\right)}}. \]
Cette formule générale offre deux cas à considérer, selon que $\dfrac{H}{2L}$ sera plus grand ou plus petit que l'unité.

Dans le premier cas, il est clair que le corps tournera sans cesse dans le même sens, et aura dans ses révolutions successives les mêmes vitesses aux mêmes points de la circonférence. Soit $\dfrac{2L}{H}=c^{2}$, et on aura
\[ t=c\sqrt{L}\int\frac{\tfrac{1}{2}d\psi}{\sqrt{(1-c^{2}\sin^{2}\tfrac{1}{2}\psi)}}=c\sqrt{L}\cdot F(\tfrac{1}{2}\psi); \]
d'où résulte le temps d'une révolution entière
\[ T=2c\sqrt{L}\cdot F^{1}. \]

Dans le second cas qui est proprement celui des oscillations, on fera $\dfrac{H}{2L}=c^{2}$, $\sin\frac{1}{2}\psi=c\sin\varphi$, et on aura
\[ t=\sqrt{L}\int\frac{d\varphi}{\sqrt{(1-c^{2}\sin^{2}\varphi)}}=\sqrt{L}\cdot F(\varphi); \]
donc le temps d'une demi-oscillation $=\sqrt{L}\cdot F^{1}$, et le temps de l'oscillation entière $T=2\sqrt{L}\cdot F^{1}$.

Lorsque les oscillations sont infiniment petites, on a $F^{1}=\frac{1}{2}\pi$, et $T=\pi\sqrt{L}$, ce qui s'accorde avec les formules connues.

Puisque le temps employé à parcourir un arc quelconque, à compter de la verticale, est représenté par une fonction elliptique de la première espèce, il s'ensuit qu'étant donné un arc parcouru

\origpage{41}
dans le temps $t$, on pourra trouver algébriquement un autre arc parcouru dans un temps multiple de $t$, ou en général commensurable avec $t$. On peut aussi trouver un arc tel, que le temps par cet arc, soit égal à la somme ou à la différence des temps par deux ou plusieurs arcs donnés, et cela, soit que ces arcs aboutissent à la verticale, soit qu'ils n'y aboutissent pas.

\rmfigure{figures/fig-5.png}{Fig.~5.}{A circular arc through its lowest point A rising to B, with the vertical line AC, the points I and O on the arc between A and B, and the chords AI and AO.}

Si on veut diviser en deux parties égales le temps de la demi-oscillation, il faudra, suivant les formules connues, faire $\sin^{2}\varphi=\dfrac{1}{1+b}$, ce qui donnera $\sin\frac{1}{2}\psi=c\sin\varphi=\sqrt{(1-b)}$. Donc si $AB$ est l'arc de la demi-oscillation, et qu'après avoir divisé l'arc $AB$ en deux également au point $I$, on prenne l'arc $AO$ tel que la corde $AO$ soit à la corde $AI$ comme $\sqrt{2}$ est à $1$, le temps de la demi-oscillation sera partagé en deux également au point $O$.

\section*{Comparaison des fonctions elliptiques de la seconde espèce.}

\subsection*{(31).}

Supposons que les amplitudes $\varphi$, $\psi$, $\mu$ soient telles qu'on ait $F(\varphi)+F(\psi)-F(\mu)=0$, je dis qu'on aura en même temps $E(\varphi)+E(\psi)-E(\mu)=P$, $P$ étant une quantité algébrique. En effet, si l'on différentie cette équation, en regardant $\mu$ comme constante, on aura
\[ dP=d\varphi\,\Delta(\varphi)+d\psi\,\Delta(\psi). \]
Mettant au lieu de $\Delta(\varphi)$ et $\Delta(\psi)$, leurs valeurs tirées des équations $(b')$, il viendra
\[ dP=d\varphi\left(\frac{\cos\varphi-\cos\psi\cos\mu}{\sin\psi\sin\mu}\right)+d\psi\left(\frac{\cos\psi-\cos\varphi\cos\mu}{\sin\varphi\sin\mu}\right), \]
ou
\[ dP=\frac{\tfrac{1}{2}d(\sin^{2}\varphi+\sin^{2}\psi+2\cos\mu\cos\varphi\cos\psi)}{\sin\mu\sin\varphi\sin\psi}; \]
Mais de l'équation $(a')$, on déduit
\[ \sin^{2}\varphi+\sin^{2}\psi+2\cos\mu\cos\varphi\cos\psi=1+\cos^{2}\mu+c^{2}\sin^{2}\mu\sin^{2}\varphi\sin^{2}\psi; \]
donc
\[ dP=\frac{\tfrac{1}{2}d(c^{2}\sin^{2}\mu\sin^{2}\varphi\sin^{2}\psi)}{\sin\mu\sin\varphi\sin\psi}=c^{2}d(\sin\mu\sin\varphi\sin\psi); \]

\origpage{42}
donc $P=c^{2}\sin\mu\sin\varphi\sin\psi$ sans constante, parce que $P$ doit s'évanouir lorsque $\varphi=0$.

Ainsi pendant que les fonctions elliptiques de la première espèce ont entre elles la relation $F(\varphi)+F(\psi)-F(\mu)=0$, les fonctions correspondantes de la seconde espèce satisfont à l'équation
\[ E(\varphi)+E(\psi)-E(\mu)=c^{2}\sin\mu\sin\varphi\sin\psi\ldots(c'): \]
c'est la formule générale qui servira à comparer entre elles les fonctions de la seconde espèce, comme nous avons comparé celles de la première.

Si l'on considérait plus généralement la fonction $G(\varphi)$ composée de la première et de la seconde espèce, savoir,
\[ G(\varphi)=E(\varphi)+kF(\varphi), \]
$k$ étant un coefficient constant quelconque, il est visible qu'on aurait semblablement
\[ G(\varphi)+G(\psi)-G(\mu)=c^{2}\sin\mu\sin\varphi\sin\psi; \]
de sorte que toutes les conséquences qu'on tirera de l'équation $(c')$ pour les fonctions $E$, s'appliqueront généralement aux fonctions $G$.

\subsection*{(32).}

Désignons comme ci-dessus par $\varphi_{2}$, $\varphi_{3}$, $\varphi_{4}$, etc. les amplitudes qui donnent $F(\varphi_{2})=2F(\varphi)$, $F(\varphi_{3})=3F(\varphi)$, etc., on aura, en vertu de l'équation $(c')$,
\[ \begin{aligned} 2E(\varphi)-E(\varphi_{2}) &= c^{2}\sin\varphi\sin\varphi\sin\varphi_{2} \\ 3E(\varphi)-E(\varphi_{3}) &= c^{2}\sin\varphi(\sin\varphi\sin\varphi_{2}+\sin\varphi_{2}\sin\varphi_{3}) \\ 4E(\varphi)-E(\varphi_{4}) &= c^{2}\sin\varphi(\sin\varphi\sin\varphi_{2}+\sin\varphi_{2}\sin\varphi_{3}+\sin\varphi_{3}\sin\varphi_{4}) \\ &\text{etc.} \end{aligned} \]
Donc la même relation entre $\varphi_{n}$ et $\varphi$, qui donne $F(\varphi_{n})=nF(\varphi)$, donnera $nE(\varphi)-E(\varphi_{n})$ égale à une quantité algébrique.

En général, si $m$, $n$, $p$, etc. sont des nombres entiers positifs ou négatifs, on peut faire ensorte que
\[ mE(\varphi)+nE(\psi)+pE(\omega)+\text{etc.} \]
soit égale à une quantité algébrique: il faut pour cela établir entre

\origpage{43}
les angles $\varphi$, $\psi$, $\omega$, etc., la relation qui donne
\[ 0=mF(\varphi)+nF(\psi)+pF(\omega)+\text{etc.} \]
Nous observerons que les fonctions $F(\varphi)$, $E(\varphi)$ sont en général de même signe que $\varphi$; lorsque $\varphi$ change de signe, elles en changent aussi, et conservent la même valeur. Cela posé, il semblerait qu'on peut satisfaire à l'équation
\[ 0=mF(\varphi)+nF(\psi)+pF(\omega)+\text{etc.} \]
de bien des manières différentes; car on est maître de changer le signe de chaque coefficient, pourvu qu'on change en même temps le signe de l'amplitude correspondante. Mais on peut se borner à considérer les fonctions $F(\varphi)$, $F(\psi)$, $F(\omega)$, etc. comme toujours positives; et dans ce cas, il n'y aura jamais qu'une relation entre les angles $\varphi$, $\psi$, $\omega$, etc., qui donnera $0=mF(\varphi)+nF(\psi)+pF(\omega)+\text{etc.}$; alors on voit que les coefficiens $m$, $n$, $p$, etc. ne sauraient être tous de même signe.

\subsection*{(33).}

Telles sont en général les relations qu'ont entre elles les fonctions elliptiques de la première et de la deuxième espèce: il faut maintenant entrer dans quelques détails sur les nombreux corollaires qu'on peut tirer de l'équation des arcs
\[ E(\varphi)+E(\psi)-E(\mu)=c^{2}\sin\varphi\sin\psi\sin\mu, \]
combinée avec l'équation algébrique qu'elle suppose et qui peut se mettre sous l'une de ces trois formes,
\[ \begin{aligned} \cos\mu &= \cos\varphi\cos\psi-\sin\varphi\sin\psi\Delta(\mu) \\ \cos\psi &= \cos\mu\cos\varphi+\sin\mu\sin\varphi\Delta(\psi) \\ \cos\varphi &= \cos\mu\cos\psi+\sin\mu\sin\psi\Delta(\varphi). \end{aligned} \]
On déduit de ces équations les valeurs suivantes qui serviront à exprimer le second membre $c^{2}\sin\varphi\sin\psi\sin\mu$ en fonction de deux seulement des amplitudes $\varphi$, $\psi$, $\mu$.
\[ \begin{aligned} \sin\mu &= \frac{\sin\varphi\cos\psi\Delta(\psi)+\sin\psi\cos\varphi\Delta(\varphi)}{1-c^{2}\sin^{2}\varphi\sin^{2}\psi} \\ \sin\psi &= \frac{\sin\mu\cos\varphi\Delta(\varphi)-\sin\varphi\cos\mu\Delta(\mu)}{1-c^{2}\sin^{2}\mu\sin^{2}\varphi} \\ \sin\varphi &= \frac{\sin\mu\cos\psi\Delta(\psi)-\sin\psi\cos\mu\Delta(\mu)}{1-c^{2}\sin^{2}\mu\sin^{2}\psi}, \end{aligned} \]

\origpage{44}

Cela posé, voici les principaux corollaires qui résultent des équations mentionnées pour la comparaison des arcs d'ellipse.

I. Si l'on fait $\mu=\frac{1}{2}\pi$, l'équation des arcs deviendra
\[ E(\varphi)+E(\psi)-E^{1}=c^{2}\sin\varphi\sin\psi, \]
et l'équation algébrique correspondante sera
\[ b\operatorname{tang}\varphi\operatorname{tang}\psi=1. \]
On en tire $\operatorname{tang}\psi=\dfrac{1}{b}\cot\varphi$, $\sin\psi=\dfrac{\cos\varphi}{\Delta(\varphi)}$, $\cos\psi=\dfrac{b\sin\varphi}{\Delta(\varphi)}$, et $c^{2}\sin\varphi\sin\psi=\dfrac{c^{2}\sin\varphi\cos\varphi}{\Delta(\varphi)}$. On aurait semblablement $\operatorname{tang}\varphi=\dfrac{1}{b}\cot\psi$, $\sin\varphi=\dfrac{\cos\psi}{\Delta(\psi)}$, $\cos\varphi=\dfrac{b\sin\psi}{\Delta(\psi)}$, et $c^{2}\sin\varphi\sin\psi=\dfrac{c^{2}\sin\psi\cos\psi}{\Delta(\psi)}$.

Si d'après cette relation entre les amplitudes $\varphi$ et $\psi$, on détermine sur l'ellipse les arcs $BM=E(\varphi)$, $BN=E(\psi)$, on aura
\[ BM-AN=c^{2}\sin\varphi\sin\psi. \]
C'est dans cette équation que consiste le théorème de Fagnani. Il en résulte qu'étant donné l'arc $BM$ terminé au petit axe, on peut trouver un second arc $AN$ terminé au grand axe, tel que la différence de ces arcs $BM-AN$ soit égale à une quantité algébrique $c^{2}\sin\varphi\sin\psi$; et cette quantité est représentée par la partie de la tangente $OM$ ou $ZN$ terminée à la perpendiculaire $CO$ ou $CZ$, abaissée du centre de l'ellipse.

Lorsque $\varphi$ et $\psi$ sont égaux à un même angle $\theta$, les deux points $M$ et $N$ coïncident en un même point $K$. Alors on a $\operatorname{tang}^{2}\theta=\dfrac{1}{b}$, $\sin^{2}\theta=\dfrac{1}{1+b}$, $\cos^{2}\theta=\dfrac{b}{1+b}$, ce qui donne
\[ BK-AK=1-b. \]
Donc alors la différence des arcs $BK$, $AK$ est égale à la différence des demi-axes $CA$, $CB$; en même temps on a
\[ \begin{aligned} BK &= \tfrac{1}{2}E^{1}+\tfrac{1}{2}(1-b) \\ AK &= \tfrac{1}{2}E^{1}-\tfrac{1}{2}(1-b). \end{aligned} \]
Il y a donc au point $K$ une sorte de bisection du quart d'ellipse,

\origpage{45}
puisque chacune des parties $BK$, $AK$ peut se mesurer par la moitié de ce quart d'ellipse.

II. Soit dans l'équation générale des arcs $\varphi=\psi=\theta$, on aura
\[ 2E(\theta)-E(\mu)=c^{2}\sin^{2}\theta\sin\mu, \]
et l'équation algébrique correspondante sera $\cos^{2}\theta-\sin^{2}\theta\Delta(\mu)=\cos\mu$. C'est l'équation qui sert à la duplication des arcs elliptiques ou à leur bisection.

On en tire pour la duplication
\[ \sin\mu=\frac{2\sin\theta\cos\theta\Delta(\theta)}{1-c^{2}\sin^{4}\theta},\qquad \operatorname{tang}\tfrac{1}{2}\mu=\Delta(\theta)\operatorname{tang}\theta, \]
et pour la bisection,
\[ \sin^{2}\theta=\frac{1-\cos\mu}{1+\Delta(\mu)}. \]
De là on voit qu'étant donné l'arc $BM=E(\theta)$, on peut trouver un autre arc $BN=E(\mu)$ qui soit mesuré par le double de $BM$, de sorte qu'on ait $2BM-BN=c^{2}\sin^{2}\theta\sin\mu$.

\rmfigure{figures/fig-6.png}{Fig.~6.}{A quadrant of an ellipse with centre C and semi-axes CA and CB, with the points M, I, N, K, P, Q marked in turn along the arc from B to A.}

Et réciproquement, étant donné l'arc $BN=E(\mu)$, on peut trouver un second arc $BM=E(\theta)$ qui soit mesuré par la moitié de l'arc $BN$; on aura en effet $BM=\frac{1}{2}BN+\frac{1}{2}c^{2}\sin^{2}\theta\sin\mu=\frac{1}{2}BN+\dfrac{1-\Delta(\mu)}{2}\operatorname{tang}\frac{1}{2}\mu$.

Lorsqu'on fait $\mu=\frac{1}{2}\pi$, le point $M$ tombe en $K$, et on a comme ci-dessus $BK=\frac{1}{2}E^{1}+\frac{1}{2}(1-b)$.

Pour avoir de nouveau le point de bisection de l'arc $BK$, il faudra faire $\operatorname{tang}^{2}\mu=\dfrac{1}{b}$, $\sin^{2}\mu=\dfrac{1}{1+b}$, $\Delta(\mu)=\sqrt{b}$, et on aura
\[ \sin^{2}\theta=\frac{1-\sqrt{\left(\dfrac{1}{1+b}\right)}}{1+\sqrt{b}}. \]
\ednote{Printed $\sqrt{\left(\frac{1}{1+b}\right)}$; since $\sin^{2}\mu=\frac{1}{1+b}$, the numerator should be $1-\cos\mu=1-\sqrt{\left(\frac{b}{1+b}\right)}$; an apparent misprint.}
Cette valeur détermine l'arc $BI=E(\theta)$ qui se mesure par la moitié de $BK$ ou par le quart de $E^{1}$; on a en effet
\[ BI=\tfrac{1}{2}BK+\tfrac{1}{2}\cdot\frac{1-\sqrt{b}}{\sqrt{(1+b)}+\sqrt{b}}, \]

\origpage{46}
ou
\[ BI=\tfrac{1}{4}E^{1}+\tfrac{1}{4}(1-b)+\tfrac{1}{2}\cdot\frac{1-\sqrt{b}}{\sqrt{(1+b)}+\sqrt{b}}, \]
et on peut continuer à l'infini cette sorte de bisection.

\rmfigure{figures/fig-7.png}{Fig.~7.}{An elliptic arc starting at B, with the points M, O, N, D marked in turn along it.}

III. Ayant pris à volonté l'arc $BD$ compté depuis le petit axe, avec un point quelconque $M$ sur cet arc, il y aura toujours un autre point $N$ sur le même arc, tel que la différence des arcs $BM$, $DN$ sera égale à une quantité algébrique.

Il suffit pour cela de faire $BD=E(\varphi)$, $BM=E(\psi)$, $BN=E(\mu)$, et de déterminer $\mu$ en fonction de $\varphi$ et de $\psi$, comme si on voulait satisfaire à l'équation $F(\mu)=F(\varphi)-F(\psi)$, ce qui se fera par la formule de l'article 19. On aura par ce moyen $E(\mu)+E(\psi)-E(\varphi)=c^{2}\sin\mu\sin\varphi\sin\psi$, ou $BM-DN=c^{2}\sin\mu\sin\varphi\sin\psi$.

Lorsque $\mu=\psi$, les deux arcs $BM$, $BN$ coïncident en un seul $BO$, et chacun des arcs $BO$, $OD$ se mesure par la moitié de l'arc $BD$. C'est un cas qui a été examiné dans le Corollaire II.

\rmfigure{figures/fig-8.png}{Fig.~8.}{An elliptic arc starting at B, with the points N, D, M, P marked in turn along it.}

IV. Etant donné un arc quelconque $BD$, terminé au petit axe, avec un point $M$ pour servir d'origine à un second arc, on peut déterminer ce second arc $MP$ ou $NM$, dans le sens qu'on voudra, de manière que sa différence avec l'arc $BD$ soit égale à une quantité algébrique.

Dans le premier cas, si l'on fait $BM=E(\varphi)$, $BD=E(\psi)$, $BP=E(\mu)$, et qu'on détermine $\mu$ d'après l'équation $F(\varphi)+F(\psi)=F(\mu)$, ou d'après les formules de l'article 18, on aura $BD-MP=c^{2}\sin\varphi\sin\psi\sin\mu$.

Dans le second cas, si l'on fait $BM=E(\varphi)$, $BD=E(\psi)$, $BN=E(\mu)$, et qu'on détermine $\mu$ d'après l'équation $F(\varphi)-F(\psi)=F(\mu)$, ou d'après les formules de l'article 19, on aura $BD-MN=c^{2}\sin\mu\sin\varphi\sin\psi$.

\rmfigure{figures/fig-9.png}{Fig.~9.}{An elliptic arc starting at B, with the points N, O, H, M, I, D, P marked in turn along it.}

V. Etant donné un arc quelconque $OD$ dont l'origine ne soit plus à l'extrémité du petit axe, avec un point $M$ pour servir d'origine à un second arc, on pourra déterminer ce second arc $MN$ ou $MP$, de manière que sa différence avec l'arc $OD$ soit égale à une quantité algébrique.

Car, 1°. par le Corollaire III, on peut trouver $BH-OD=$ à une quantité algébrique, et ensuite $BH-MN=$ à une quantité algébrique; donc $MN-OD$ sera encore une quantité algébrique; 2°. ayant trouvé le point $H$ par le Coroll.\ III, on peut trouver le

\origpage{47}
point $P$ par le Coroll.\ IV, de sorte que $BH-MP$ soit égal à une quantité algébrique, et alors $OD-MP$ sera égal aussi à une quantité algébrique.

Ainsi on peut trouver sur l'ellipse une infinité d'arcs égaux à un arc donné, plus ou moins une différence assignable géométriquement, de sorte que l'arc prendra son origine à volonté sur tous les points de l'ellipse, et sera dirigé dans le sens qu'on voudra.

VI. Quel que soit l'arc $OP$ et le point $M$ pris sur cet arc, il y aura toujours sur le même arc un autre point $D$ tel que la différence des arcs $OM$, $DP$ sera égale à une quantité algébrique.

Cela suit immédiatement du Coroll.\ V.

Lorsque les points $M$ et $D$ coïncident en un seul point $I$, chacun des arcs $OI$, $IP$ est mesuré par la moitié de $OP$, et on a une première bisection de cet arc. On pourra de même en trouver une seconde, une troisième, etc.\ à l'infini.

VII. Etant donné l'arc $BM$ dont l'origine est au petit axe, on peut trouver un arc $BN$ qui soit égal à un multiple quelconque de l'arc $BM$, plus ou moins une quantité algébrique.

Car en faisant $BM=E(\varphi)$, $BN=E(\psi)$, si on satisfait à l'équation $F(\psi)=nF(\varphi)$, on aura en même temps $nE(\varphi)-E(\psi)=$ à une quantité algébrique. Dans ce cas, $\psi$ serait ce que nous avons désigné ci-dessus par $\varphi_{n}$, et on a vu la manière de déterminer $\varphi_{n}$ par le moyen de $\varphi$.

VIII. Réciproquement, étant donné un arc quelconque $BN=E(\psi)$, on pourra, par la résolution d'une équation algébrique, déterminer l'arc $BM=E(\varphi)$, qui soit égal à un sous-multiple de l'arc $BN$, plus une quantité algébrique.

Par exemple, pour la trisection de l'arc $BN$, il faudra déterminer $\sin\varphi$ par l'équation
\[ \sin\psi=\frac{3\sin\varphi-4(1+c^{2})\sin^{3}\varphi+6c^{2}\sin^{5}\varphi-c^{4}\sin^{9}\varphi}{1-6c^{2}\sin^{4}\varphi+4c^{2}(1+c^{2})\sin^{6}\varphi-3c^{4}\sin^{8}\varphi}, \]
équation qui est en général du neuvième degré, mais qui se réduit au quatrième lorsque $\psi=\frac{1}{2}\pi$.

IX. En général on pourra trouver par la résolution d'une équation algébrique, un arc $E(\varphi)$ qui soit égal à une partie rationnelle

\origpage{48}
$\frac{m}{n}$ de l'arc donné $E(\psi)$, plus ou moins une quantité algébrique. L'équation sera la même que celle qui donnerait $F(\varphi)=\frac{m}{n}F(\psi)$.

Il en sera de même si l'arc donné n'est pas terminé au petit axe. Car si $OP$ est cet arc, on cherchera par le Coroll.\ III, l'arc $BM$ égal à $OP+$ une quantité algébrique, et il ne s'agira plus que de trouver un arc $BN=\frac{m}{n}BM\pm$ une quantité algébrique. On pourra ensuite donner à l'arc $BM$ une autre origine à volonté.

X. Deux arcs étant donnés partout où l'on voudra sur l'ellipse, on pourra trouver un arc égal à leur somme ou à leur différence, plus ou moins une ligne droite assignable géométriquement. C'est une suite des Coroll.\ III et V.

Ainsi toutes les comparaisons qu'on fait ordinairement des arcs de cercle par voie d'analyse, ont lieu également pour les arcs d'ellipse, à la différence près qui affecte tous les résultats, mais qu'on peut faire disparaître dans beaucoup de cas, lorsque l'origine de l'arc cherché est arbitraire. La disparition de cette différence, lorsqu'elle peut avoir lieu, ajoute aux problèmes un degré d'intérêt de plus; elle rapproche alors les propriétés des arcs d'ellipse de celles des fonctions elliptiques de la première espèce. C'est pourquoi nous croyons devoir en apporter ici quelques exemples.

\subsection*{(34).}

Problème I. \emph{Déterminer sur le quart d'ellipse $BKA$ un arc $MP$ qui soit précisément égal à la moitié de l'arc $BKA$.}

Soit $\varphi$ l'amplitude du point $M$, $\psi$ celle du point $P$, $\theta$ l'amplitude du point $K$, premier point de bisection de l'arc $BKA$, pour lequel on a $\sin^{2}\theta=\frac{1}{1+b}$ et $E(\theta)=\frac{1}{2}E^{1}+\frac{1}{2}(1-b)$. Supposons qu'on ait $F(\varphi)+F(\theta)-F(\psi)=0$, il s'ensuivra $E(\varphi)+E(\theta)-E(\psi)=c^{2}\sin\varphi\sin\psi\sin\theta$; donc
\[ E(\psi)-E(\varphi)=\tfrac{1}{2}E^{1}+\tfrac{1}{2}(1-b)-c^{2}\sin\varphi\sin\psi\sin\theta. \]
Donc si on veut qu'on ait $MP=\frac{1}{2}E^{1}$, il faudra faire
\[ c^{2}\sin\varphi\sin\psi\sin\theta=\tfrac{1}{2}(1-b), \]
ce qui donnera $\sin\varphi\sin\psi=\frac{1}{2}\sin\theta$. Mais d'un autre côté l'équation

\origpage{49}
$F(\varphi)+F(\theta)-F(\psi)=0$ donne $\cos\varphi\cos\psi+\sin\varphi\sin\psi\Delta(\theta)=\cos\theta$, et puisque $\Delta(\theta)=\sqrt{b}$, on aura à la fois $\sin\varphi\sin\psi=\frac{1}{2}\sin\theta$, et $\cos\varphi\cos\psi=\frac{1}{2}\cos\theta$. De ces équations, on tire
\[ \begin{aligned} \cos(\psi-\varphi)&=\tfrac{1}{2}(\cos\theta+\sin\theta)=\cos\frac{\pi}{4}\cos\left(\theta-\frac{\pi}{4}\right)\\ \cos(\psi+\varphi)&=\tfrac{1}{2}(\cos\theta-\sin\theta)=\cos\frac{\pi}{4}\cos\left(\theta+\frac{\pi}{4}\right). \end{aligned} \]
Ainsi les angles $\psi-\varphi$ et $\psi+\varphi$ seront connus. On pourrait aussi déterminer directement les valeurs de $\sin\varphi$ et $\sin\psi$ au moyen des formules
\[ \begin{aligned} \sin\varphi&=\tfrac{1}{4}\sqrt{(3+4\sin\theta+2\sin^{2}\theta)}-\tfrac{1}{4}\sqrt{(3-4\sin\theta+2\sin^{2}\theta)}\\ \sin\psi&=\tfrac{1}{4}\sqrt{(3+4\sin\theta+2\sin^{2}\theta)}+\tfrac{1}{4}\sqrt{(3-4\sin\theta+2\sin^{2}\theta)}, \end{aligned} \]
et ces sinus seront les abscisses des points cherchés $M$ et $P$.

\subsection*{(35).}

Problème II. \emph{Déterminer sur le quart d'ellipse $BMA$ un arc $NQ$ qui soit égal au tiers de $BMA$.}

Si l'on suppose $F(\varphi)=\frac{1}{3}F^{1}$ et $F(\varphi_{2})=2F(\varphi)$, on aura $\sin\varphi$ par la résolution de l'équation
\[ 0=1-2\sin\varphi+2c^{2}\sin^{3}\varphi-c^{2}\sin^{4}\varphi, \]
et $\varphi_{2}$ se déduira de $\varphi$, soit par l'équation $\cos\varphi_{2}=1-\sin\varphi$, soit par l'équation $\sin\varphi_{2}=\frac{\cos\varphi}{\Delta}$.

Cela posé, on aura $3E(\varphi)-E^{1}=c^{2}\sin\varphi\sin\varphi_{2}(1+\sin\varphi)$, ou
\[ E(\varphi)=\tfrac{1}{3}E^{1}+\tfrac{1}{3}c^{2}\sin\varphi\sin\varphi_{2}(1+\sin\varphi). \]
Supposons de nouveau $F(\varphi)+F(\psi)-F(\omega)=0$, on aura
\[ E(\varphi)+E(\psi)-E(\omega)=c^{2}\sin\varphi\sin\psi\sin\omega. \]
Donc si l'on veut que $E(\omega)-E(\psi)=\frac{1}{3}E^{1}$, il faudra faire $\sin\psi\sin\omega=\frac{1}{3}\sin\varphi_{2}(1+\sin\varphi)$. Mais d'ailleurs en vertu de la supposition faite, on a l'équation $\cos\psi\cos\omega+\sin\psi\sin\omega\Delta(\varphi)=\cos\varphi$; donc les inconnues $\psi$ et $\omega$ devront être déterminées par les équations
\[ \begin{aligned} \sin\psi\sin\omega&=\tfrac{1}{3}\sin\varphi_{2}(1+\sin\varphi)=\frac{\cos^{3}\varphi}{3b\sin\varphi}\\ \cos\psi\cos\omega&=\tfrac{1}{3}\cos\varphi(2-\sin\varphi) \end{aligned} \]

\origpage{50}
Elles donneront immédiatement les valeurs de $\cos(\omega-\psi)$ et $\cos(\omega+\psi)$, d'où l'on conclura celles de $\psi$ et $\omega$. Connaissant ces valeurs, on fera $BN=E(\psi)$, $BQ=E(\omega)$, et l'arc $NQ$ sera égal au tiers du quart d'ellipse $E^{1}$.

On peut démontrer que le problème est toujours possible par cette considération. Soit pris $BI=E(\varphi)$, on aura $BI>\frac{1}{3}E^{1}$. Soit pris ensuite l'arc $AP$ correspondant à $BI$, de manière que la différence $BI-AP$ soit égale à une quantité algébrique, c'est-à-dire, soit prise l'amplitude $\sigma$ du point $P$, de manière qu'on ait $b\operatorname{tang}\varphi\operatorname{tang}\sigma=1$, on aura $AP<\frac{1}{3}E^{1}$. Mais si on imagine que l'arc $\frac{1}{3}E^{1}$ soit transporté le long du quart d'ellipse, de manière que son origine parcoure successivement tout l'intervalle de $B$ en $P$, cet arc étant représenté en $B$ par $BI>\frac{1}{3}E^{1}$, et en $P$ par $AP<\frac{1}{3}E^{1}$, il faudra, en vertu de la loi de continuité, qu'il soit représenté exactement en un point intermédiaire $N$ par l'arc $NQ=\frac{1}{3}E^{1}$.

Pour appliquer la solution générale à un cas particulier, soit $c^{2}=\frac{1}{2}=b^{2}$, on a trouvé ci-dessus (art.~24) $\cos\varphi=\sqrt[4]{(2\sqrt{3}-3)}=\sqrt[4]{(8\cos 30^{\circ}\sin^{2}15^{\circ})}$; on aura donc par le calcul trigonométrique,
\[ \begin{aligned} L\cos\varphi&=9.9166532 & &\\ L\sin\varphi&=9.7517253\ldots\ldots & \sin\varphi&=0.5645797\\ & & 2-\sin\varphi&=1.4354203\\ L\sin\psi\sin\omega&=9.6716281 & \sin\psi\sin\omega&=0.4694920\\ L\cos\psi\cos\omega&=9.5965110 & \cos\psi\cos\omega&=0.3949217\\ & & \cos(\omega+\psi)&=-0.0745703\\ & & \cos(\omega-\psi)&=0.8644137\\ \omega+\psi&=94^{\circ}16'35''5 & \omega&=62^{\circ}13'49''3\\ \omega-\psi&=30.11.3.1 & \psi&=32^{\circ}2'46''2 \end{aligned} \]
Ces valeurs ne sont qu'approchées, mais les formules trouvées donnent la solution rigoureuse qui peut même être construite géométriquement.

\rmfigure{figures/fig-10.png}{Fig.~10.}{An elliptic arc starting at B, with the points M, N, D, P, Q, R marked in turn along it.}

\subsection*{(36).}

Problème III. \emph{Etant donnés les deux arcs $MN$ et $PQ$, situés comme on voudra sur la circonférence de l'ellipse, trouver un troisième arc $DR$ égal à leur somme $MN+PQ$.}

Soient $\alpha$, $\beta$, $\delta$, $\varepsilon$ les amplitudes des points donnés $M$, $N$, $P$, $Q$,

\origpage{51}
ensorte qu'on ait $MN=E(\beta)-E(\alpha)$, $PQ=E(\varepsilon)-E(\delta)$; soient $\psi$ et $\omega$ les amplitudes des points cherchés $D$, $R$, ensorte qu'on ait $DR=E(\omega)-E(\psi)$.

Soient encore $\lambda$ et $\mu$ deux amplitudes telles qu'on ait $F(\lambda)=F(\beta)-F(\alpha)$, $F(\mu)=F(\varepsilon)-F(\delta)$; on aura en même temps
\[ \begin{aligned} E(\lambda)+E(\alpha)-E(\beta)&=c^{2}\sin\alpha\sin\beta\sin\lambda\\ E(\mu)+E(\delta)-E(\varepsilon)&=c^{2}\sin\delta\sin\varepsilon\sin\mu; \end{aligned} \]
donc
\[ MN+PQ=E(\lambda)+E(\mu)-c^{2}\sin\alpha\sin\beta\sin\lambda-c^{2}\sin\delta\sin\varepsilon\sin\mu. \]
Soit enfin $\varphi$ une nouvelle amplitude telle que $F(\varphi)=F(\lambda)+F(\mu)$, on aura
\[ E(\lambda)+E(\mu)-E(\varphi)=c^{2}\sin\varphi\sin\lambda\sin\mu; \]
donc
\[ MN+PQ-E(\varphi)=c^{2}\sin\varphi\sin\lambda\sin\mu-c^{2}\sin\alpha\sin\beta\sin\lambda-c^{2}\sin\delta\sin\varepsilon\sin\mu. \]
Faisons maintenant $F(\varphi)=F(\omega)-F(\psi)$, il en résultera
\[ E(\varphi)+E(\psi)-E(\omega)=c^{2}\sin\varphi\sin\psi\sin\omega; \]
mais par hypothèse, on a $E(\omega)-E(\psi)=MN+PQ$; donc
\[ \sin\varphi\sin\psi\sin\omega=\sin\alpha\sin\beta\sin\lambda+\sin\delta\sin\varepsilon\sin\mu-\sin\lambda\sin\mu\sin\varphi. \]
Soit pour abréger $\sin\alpha\sin\beta\sin\lambda+\sin\delta\sin\varepsilon\sin\mu=M$, et on aura
\[ \sin\omega\sin\psi=\frac{M}{\sin\varphi}-\sin\lambda\sin\mu. \]
D'ailleurs l'équation $F(\varphi)=F(\omega)-F(\psi)$, donne $\cos\omega\cos\psi+\sin\omega\sin\psi\Delta(\varphi)=\cos\varphi$; donc on aura pour déterminer $\psi$ et $\omega$, les deux équations
\[ \begin{aligned} \sin\omega\sin\psi&=\frac{M}{\sin\varphi}-\sin\lambda\sin\mu\\ \cos\omega\cos\psi&=\cos\varphi-\frac{\Delta(\varphi)}{\sin\varphi}M+\sin\lambda\sin\mu\Delta(\varphi). \end{aligned} \]
Si de plus on observe que l'équation $F(\varphi)=F(\lambda)+F(\mu)$ donne $\cos\varphi=\cos\lambda\cos\mu-\sin\lambda\sin\mu\Delta(\varphi)$, on déduira des équations

\origpage{52}
précédentes,
\[ \begin{aligned} \cos(\omega+\psi)&=\cos(\lambda-\mu)-\frac{M}{\sin\varphi}[1+\Delta(\varphi)]\\ \cos(\omega-\psi)&=\cos(\lambda+\mu)+\frac{Mc^{2}\sin\varphi}{1+\Delta(\varphi)}. \end{aligned} \]
Les données immédiates étant $\alpha$, $\beta$, $\delta$, $\varepsilon$, on en déduit $\lambda$ et $\mu$ par les équations $F(\lambda)=F(\beta)-F(\alpha)$, $F(\mu)=F(\varepsilon)-F(\delta)$, qui donnent
\[ \begin{aligned} \sin\lambda&=\frac{\sin\beta\cos\alpha\Delta(\alpha)-\sin\alpha\cos\beta\Delta(\beta)}{1-c^{2}\sin\alpha\sin^{2}\beta}\\ \sin\mu&=\frac{\sin\varepsilon\cos\delta\Delta(\delta)-\sin\delta\cos\varepsilon\Delta(\varepsilon)}{1-c^{2}\sin\delta\sin^{2}\varepsilon}. \end{aligned} \]
\ednote{In both denominators the print sets the first sine with no exponent ($\sin\alpha$, $\sin\delta$); by the addition formula $\sin^{2}\alpha$ and $\sin^{2}\delta$ are expected (cf.\ $\sin^{2}\mu\sin^{2}\lambda$ below). Apparent misprints or dropped superscripts.}
La valeur de $M$ est donc connue, ensuite on déduira $\sin\varphi$ et $\Delta(\varphi)$ de l'équation $F(\varphi)=F(\lambda)+F(\mu)$, qui donne
\[ \begin{aligned} \sin\varphi&=\frac{\sin\lambda\cos\mu\Delta(\mu)+\sin\mu\cos\lambda\Delta(\lambda)}{1-c^{2}\sin^{2}\mu\sin^{2}\lambda}\\ \Delta(\varphi)&=\frac{\Delta(\mu)\Delta(\lambda)-c^{2}\sin\lambda\sin\mu\cos\lambda\cos\mu}{1-c^{2}\sin^{2}\mu\sin^{2}\lambda}. \end{aligned} \]
Or les valeurs de $\sin\lambda$, $\cos\lambda$, $\Delta(\lambda)$ sont données en fonctions de $\beta$ et $\alpha$ par les formules de l'art.~19; il en est de même des valeurs de $\sin\mu$, $\cos\mu$, $\Delta(\mu)$ exprimées en fonctions de $\varepsilon$ et $\delta$. On connaîtra donc toutes les quantités qui composent les valeurs de $\cos(\omega-\psi)$ et $\cos(\omega+\psi)$.

Ce problème peut servir à en résoudre beaucoup d'autres, et notamment à trouver un arc qui soit exactement dans un rapport rationnel avec un arc donné; mais il faut pour cela que dans chaque application les valeurs trouvées pour $\cos(\omega+\psi)$ et $\cos(\omega-\psi)$, soient renfermées chacune entre les limites $+1$ et $-1$, sans quoi le problème deviendrait impossible.

\section*{Comparaison des arcs d'hyperbole.}

\subsection*{(37).}

Nous avons déjà trouvé (art.~13) que l'arc $AM$ désigné par $\Upsilon$, a pour expression
\[ \Upsilon=\Delta\operatorname{tang}\varphi-E(\varphi)+b^{2}F(\varphi). \]

\origpage{53}
Le point $M$, extrémité de l'arc $AM$, est censé avoir pour amplitude $\varphi$. Considérons deux autres points $N$ et $O$ dont les amplitudes soient $\psi$ et $\omega$, et supposons qu'on ait l'équation $F(\varphi)+F(\psi)-F(\omega)=0$; alors les arcs $AM$, $AN$, $AO$, désignés respectivement par $\Upsilon(\varphi)$, $\Upsilon(\psi)$, $\Upsilon(\omega)$, auront entre eux cette relation
\[ \begin{aligned} \Upsilon(\varphi)+\Upsilon(\psi)-\Upsilon(\omega)={}&\Delta(\varphi)\operatorname{tang}\varphi+\Delta(\psi)\operatorname{tang}\psi-\Delta(\omega)\operatorname{tang}\omega\\ &-E(\varphi)-E(\psi)+E(\omega), \end{aligned} \]
ou, en mettant la valeur connue de $E(\varphi)+E(\psi)-E(\omega)$,
\[ \begin{aligned} \Upsilon(\varphi)+\Upsilon(\psi)-\Upsilon(\omega)={}&\Delta(\varphi)\operatorname{tang}\varphi+\Delta(\psi)\operatorname{tang}\psi-\Delta(\omega)\operatorname{tang}\omega\\ &-c^{2}\sin\varphi\sin\psi\sin\omega. \end{aligned} \]
C'est l'équation fondamentale d'après laquelle on peut faire sur les arcs d'hyperbole les mêmes comparaisons que nous avons faites sur les arcs d'ellipse, mais en observant que dans l'hyperbole on ne peut donner aux amplitudes une valeur plus grande que $\frac{1}{2}\pi$, et que lorsque $\varphi=\frac{1}{2}\pi$, l'arc $AM$ devient infini.

La quantité $\Delta\operatorname{tang}\varphi$ n'est autre chose que la tangente $MZ$, terminée par la perpendiculaire $CZ$ abaissée du centre sur cette tangente. Ainsi $\Delta\operatorname{tang}\varphi-\Upsilon(\varphi)$ est l'excès de la tangente $MZ$ sur l'arc $AZ$\ednote{Printed $AZ$; the arc meant is $AM$, an apparent misprint.}. Si on appelle $G(\varphi)$ cette fonction, on aura pour chaque point $M$,
\[ G(\varphi)=E(\varphi)-b^{2}F(\varphi), \]
et lorsque les trois points $M$, $N$, $O$, déterminés par les amplitudes $\varphi$, $\psi$, $\omega$ sont liés entre eux par la relation $F(\varphi)+F(\psi)-F(\omega)=0$, les fonctions correspondantes $G(\varphi)$, $G(\psi)$, $G(\omega)$ relatives à l'hyperbole, satisferont à l'équation
\[ G(\varphi)+G(\psi)-G(\omega)=c^{2}\sin\varphi\sin\psi\sin\omega, \]
équation entièrement semblable à celle qui a lieu dans l'ellipse, et d'où l'on déduira de semblables corollaires, sauf les restrictions particulières à l'hyperbole et dont nous avons déjà parlé.

Nous avons trouvé que lorsqu'on fait $\sin^{2}\theta=\frac{1}{1+b}$, on a $F(\theta)=\frac{1}{2}F^{1}$, et $E(\theta)=\frac{1}{2}E^{1}+\frac{1}{2}(1-b)$, on aura donc semblablement pour l'hyperbole,
\[ G(\theta)=\tfrac{1}{2}G^{1}+\tfrac{1}{2}(1-b), \]

\origpage{54}
$G^{1}$ est la différence entre la branche infinie d'hyperbole $AMO$ et son asymptote $CV$, qui est censée la rencontrer dans un point infiniment éloigné. Cette différence estimée au moyen des fonctions $E^{1}$, $F^{1}$, a pour valeur $G^{1}=E^{1}-b^{2}F^{1}$; mais sans recourir à ces fonctions, on peut la déduire de l'équation précédente qui donne
\[ G^{1}=2G(\theta)-(1-b). \]
Ainsi la quantité $G^{1}$, différence de deux infinis, se déterminera par la quantité $G(\theta)$, relative au point $K$ dont l'amplitude est $\theta$, et qui a pour coordonnées $y=b^{2}\operatorname{tang}\theta=b\sqrt{b}$, $x=c\sqrt{(1+b)}$.

L'amplitude $\theta$ est celle qui donne $F(\theta)=\frac{1}{2}F^{1}$; si on cherche successivement par les formules de la bisection les amplitudes $\theta'$, $\theta''$, etc.\ telles qu'on ait $F(\theta')=\frac{1}{2}F(\theta)=\frac{1}{4}F^{1}$, $F(\theta'')=\frac{1}{2}F(\theta')=\frac{1}{8}F^{1}$, etc., on aura en même temps
\[ \begin{gathered} G(\theta)=2G(\theta')-c^{2}\sin^{2}\theta'\sin\theta\\ G(\theta')=2G(\theta'')-c^{2}\sin^{2}\theta''\sin\theta'\\ \text{etc.} \end{gathered} \]
D'où il suit que la quantité $G^{1}$ se déterminera par le dernier terme de la suite $G(\theta)$, $G(\theta')$, $G(\theta'')$, etc., prolongée aussi loin qu'on voudra. Or lorsque $\varphi$ est devenu très-petit, la quantité $G(\varphi)$ a pour valeur très-approchée $c^{2}\sin\varphi$; ainsi la bisection répétée de la fonction $G(\theta)$ fournit un moyen de déterminer par approximation, la valeur de la transcendante $G^{1}$, et la même méthode s'applique à toute fonction $G(\varphi)$ dont on voudrait avoir une valeur approchée. Mais nous donnerons ci-après pour cet objet des méthodes plus expéditives.

\section*{Développement particulier de la formule $Z=\displaystyle\int\frac{(f+gx^{2})\,dx}{\sqrt{(\alpha^{2}+2\alpha\beta x^{2}\cos\theta+\beta^{2}x^{4})}}$,}

\subsection*{(38).}

Cette formule se rencontre assez souvent dans les applications, et d'ailleurs il est nécessaire d'examiner particulièrement le cas des facteurs imaginaires dont nous avons parlé (art.~7).

\origpage{55}

La variable $x$ est susceptible de toutes les valeurs, depuis $x=0$ jusqu'à $x=\infty$; mais comme en faisant $x=\infty$, on a $Z=\displaystyle\int\frac{g}{\beta}\,dx$, pour nous débarrasser du terme qui peut devenir infini, nous considérerons simplement la formule
\[ X=\int\left(\frac{(f+gx^{2})\,dx}{\sqrt{(\alpha^{2}+2\alpha\beta x^{2}\cos\theta+\beta^{2}x^{4})}}-\frac{g}{\beta}\,dx\right), \]
qui par ce moyen aura toujours une valeur finie.

Il s'agit maintenant de transformer cette expression de manière que les facteurs binomes de la quantité sous le radical deviennent réels. Pour cela on peut faire indifféremment l'une des quatre suppositions suivantes:
\[ \begin{gathered} \sqrt{(\alpha^{2}+2\alpha\beta x^{2}\cos\theta+\beta^{2}x^{4})}=2xy\sqrt{\alpha\beta},\\ \alpha+\beta x^{2}=2xy\sqrt{\alpha\beta},\\ \beta x^{2}+\alpha\cos\theta+\sqrt{(\alpha^{2}+2\alpha\beta x^{2}\cos\theta+\beta^{2}x^{4})}=2\alpha y^{2},\\ x=\frac{1-y}{1+y}\sqrt{\frac{\alpha}{\beta}}, \end{gathered} \]
et la transformée en $y$ aura les conditions requises. Bornons-nous à présenter le résultat de la troisième supposition: elle donne $x^{2}=\frac{\alpha}{\beta}\left(y^{2}-\cos\theta-\frac{\sin^{2}\theta}{4y^{2}}\right)$, et
\[ X=\int\frac{\frac{dy}{\sqrt{\alpha\beta}}\left(f-\frac{g\alpha}{\beta}\cos\theta-\frac{g\alpha\sin^{2}\theta}{2\beta y^{2}}\right)}{\sqrt{[(y^{2}-\cos^{2}\tfrac{1}{2}\theta)(y^{2}+\sin^{2}\tfrac{1}{2}\theta)]}}, \]
où l'on voit qu'en effet les deux facteurs de la quantité sous le radical sont réels. On voit aussi que la moindre valeur de $y$ étant $\cos\frac{1}{2}\theta$, on peut faire $y=\frac{\cos\frac{1}{2}\theta}{\cos\varphi}$, ce qui donnera la transformée suivante, où l'on a fait $c=\sin\frac{1}{2}\theta$,
\[ X=\int\frac{f\beta-g\alpha+2g\alpha c^{2}\sin^{2}\varphi}{\beta\sqrt{\alpha\beta}}\cdot\frac{d\varphi}{\sqrt{(1-c^{2}\sin^{2}\varphi)}}, \]
et il faut remarquer que dans cette formule nous supposons $\sqrt{\alpha\beta}$ réel; car si le terme $2\alpha\beta x^{2}\cos\theta$ était négatif, on ferait tomber le signe $-$ sur $\cos\theta$.

\origpage{56}

Cela posé, on aura, suivant les expressions accoutumées,
\[ X=\frac{f\beta+g\alpha}{\beta\sqrt{\alpha\beta}}F-\frac{2g\alpha}{\beta\sqrt{\alpha\beta}}E; \]
c'est l'intégrale demandée prise depuis $x=0$, car on a
\[ \cos^{2}\varphi=\frac{-\beta x^{2}-\alpha\cos\theta+\sqrt{(\alpha^{2}+2\alpha\beta x^{2}\cos\theta+\beta^{2}x^{4})}}{2\alpha\sin^{2}\tfrac{1}{2}\theta}, \]
et réciproquement $x\sqrt{\frac{\beta}{\alpha}}=\frac{\sin\varphi}{\cos\varphi}\sqrt{(1-c^{2}\sin^{2}\varphi)}$; d'où l'on voit qu'en faisant $x=0$, on a $\varphi=0$, et qu'en faisant $x=\infty$, on a $\varphi=\frac{1}{2}\pi$.

La valeur totale de l'intégrale, prise depuis $x=0$ jusqu'à $x=\infty$, sera donc
\[ X^{1}=\frac{f\beta+g\alpha}{\beta\sqrt{\alpha\beta}}F^{1}-\frac{2g\alpha}{\beta\sqrt{\alpha\beta}}E^{1}. \]
Voici quelques applications de ces formules qui conduiront à des résultats assez remarquables.

\section*{EXEMPLE I.}

\subsection*{(39).}

Soit proposé d'évaluer les deux intégrales
\[ M=\int\frac{z^{-\frac{2}{3}}\,dz}{\sqrt{(1-z^{2})}},\quad N=\int\frac{z^{\frac{1}{3}}\,dz}{\sqrt{(1-z^{2})}} \]
entre les limites $z=0$, $z=1$. On sait d'ailleurs que le produit de ces intégrales $=\frac{3}{2}\pi$ (\emph{Calc.\ intég.\ d'Euler}, tom.~1, pag.~244).

Si on fait dans la première $z=(1-x^{2})^{-\frac{3}{2}}$\ednote{Printed $1-x^{2}$; the transformed integral and the limits $x=0$ to $x=\infty$ require $z=(1+x^{2})^{-\frac{3}{2}}$, an apparent misprint.}, on aura la transformée $M=\displaystyle\int\frac{3\,dx}{\sqrt{(3+3x^{2}+x^{4})}}$, qu'il faudra intégrer depuis $x=0$ jusqu'à $x=\infty$. Cette formule étant comparée avec la formule $X$, on aura $f=3$, $g=0$, $\alpha=\sqrt{3}$, $\beta=1$, $\cos\theta=\frac{1}{2}\sqrt{3}$, $c=\sin\frac{1}{2}\theta=\frac{1}{2}\sqrt{(2-\sqrt{3})}=\sin 15^{\circ}$. Ainsi l'intégrale indéfinie est $M=\frac{3}{\sqrt[4]{3}}F$, et l'intégrale complète, en faisant $x=\infty$ ou $\varphi=\frac{1}{2}\pi$, est
\[ M^{1}=\frac{3}{\sqrt[4]{3}}F^{1}. \]

\origpage{57}

Si dans la seconde formule on fait $z=(1-x^{2})^{\frac{3}{2}}$, la transformée sera $N=\displaystyle\int\frac{3\,dx\,(1-x^{2})}{\sqrt{(3-3x^{2}+x^{4})}}$, et l'intégrale devra être prise depuis $x=0$ jusqu'à $x=1$. Comparant cette formule à la formule $X$, on aura $f=3$, $g=-3$, $\alpha=\sqrt{3}$, $\beta=1$, $\cos\theta=-\frac{1}{2}\sqrt{3}$, $c'=\cos 15^{\circ}=\frac{1}{2}\sqrt{(2+\sqrt{3})}$; ce module étant différent de celui de l'autre formule, je le distingue par un accent; j'ai donc par la substitution
\[ N=-3x+\frac{3-3\sqrt{3}}{\sqrt[4]{3}}F(c',\varphi)+\frac{6\sqrt{3}}{\sqrt[4]{3}}E(c',\varphi). \]
Or la relation entre $\varphi$ et $x$ étant
\[ \frac{3+2\sqrt{3}}{2}\cos^{2}\varphi=\tfrac{3}{2}-x^{2}+\sqrt{(3-3x^{2}+x^{4})}, \]
si on fait $x=1$, on aura $\cos^{2}\varphi=2\sqrt{3}-3$, ou bien $\operatorname{tang}\varphi=\sqrt{\left(\frac{2}{\sqrt{3}}\right)}$. Mais nous avons déjà vu (art 24) que dans le cas de $c=\frac{1}{2}\sqrt{(2+\sqrt{3})}$ et $\operatorname{tang}\varphi=\sqrt{\frac{2}{\sqrt{3}}}$, on a exactement $F(\varphi)=\frac{1}{3}F^{1}$; il s'ensuit par les formules de l'art.~32, qu'on a en même temps $E(\varphi)=\frac{1}{3}E^{1}+\frac{1}{2\sqrt[4]{3}}$. Donc l'intégrale $N$ prise depuis $x=0$ jusqu'à $x=1$, aura cette valeur
\[ N^{1}=\frac{1-\sqrt{3}}{\sqrt[4]{3}}F^{1}(c')+\frac{2\sqrt{3}}{\sqrt[4]{3}}E^{1}(c'). \]
Puis donc qu'on a déjà trouvé $M^{1}=\frac{3}{\sqrt[4]{3}}F^{1}(c)$, et que le produit $M^{1}N^{1}=\frac{3}{2}\pi$, on aura entre les trois fonctions $F^{1}(c)$, $F^{1}(c')$, $E^{1}(c')$, cette relation fort remarquable:
\[ \tfrac{1}{2}\pi=F^{1}(c)\cdot\left[\frac{1-\sqrt{3}}{\sqrt{3}}F^{1}(c')+2E^{1}(c')\right], \]
d'où il suit que dans ce cas particulier l'arc d'ellipse $E^{1}(c')$ qui est une fonction de la seconde espèce, peut s'exprimer par les deux fonctions de la première espèce $F^{1}(c)$, $F^{1}(c')$. D'ailleurs on peut observer que $c'=\sqrt{(1-c^{2})}=b$, c'est-à-dire que les modules $c$ et $c'$ sont complémens l'un de l'autre.

\origpage{58}

\section*{EXEMPLE II.}

\subsection*{(40).}

On propose d'évaluer par les fonctions elliptiques, les deux intégrales $P=\displaystyle\int\frac{z^{-\frac{1}{3}}\,dz}{\sqrt{(1-z^{2})}}$, $Q=\displaystyle\int\frac{z^{\frac{2}{3}}\,dz}{\sqrt{(1-z^{2})}}$, depuis $z=0$ jusqu'à $z=1$: on sait d'ailleurs que leur produit $=\frac{3}{4}\pi$.

Si on fait dans la première, $z=(1-x^{2})^{\frac{3}{2}}$, on aura la transformée $P=\displaystyle\int\frac{3\,dx}{\sqrt{(3-3x^{2}+x^{4})}}$, qu'il faut intégrer depuis $x=0$ jusqu'à $x=1$. Celle-ci étant traitée comme l'intégrale $N$ de l'exemple précédent, on aura pour résultat
\[ P^{1}=\frac{1}{\sqrt[4]{3}}F^{1}(c'). \]

Venons à la formule $Q$, et faisons d'abord $z^{2}=y^{-3}$, nous aurons la transformée $Q=\frac{3}{2}\displaystyle\int\frac{dy}{y^{2}\sqrt{(y^{3}-1)}}$ qu'il faudra intégrer depuis $y=1$ jusqu'à $y=\infty$. On a d'abord, en intégrant par partie,
\[ Q=\frac{3}{2}\cdot\frac{\sqrt{(y^{3}-1)}}{y}-\frac{3}{4}\int\frac{y\,dy}{\sqrt{(y^{3}-1)}}. \]
Soit ensuite $y=1+x^{2}$, et on aura $\displaystyle\int\frac{y\,dy}{\sqrt{(y^{3}-1)}}=\displaystyle\int\frac{2(1+x^{2})\,dx}{\sqrt{(3+3x^{2}+x^{4})}}$. Cela posé, la valeur de $Q$ pourra se mettre sous la forme
\[ Q=\frac{3}{2}\cdot\frac{x\sqrt{(3+3x^{2}+x^{4})}}{1+x^{2}}-\frac{3}{2}x+\frac{3}{2}\int\left[dx-\frac{(1+x^{2})\,dx}{\sqrt{(3+3x^{2}+x^{4})}}\right]. \]
Cette intégrale doit être prise depuis $x=0$ jusqu'à $x=\infty$, et comme la partie hors du signe s'évanouit dans ces deux limites, on aura simplement
\[ Q=\frac{3}{2}\int\left[dx-\frac{(1+x^{2})\,dx}{\sqrt{(3+3x^{2}+x^{4})}}\right]. \]
Comparant avec la formule $X$, on aura $f=-\frac{3}{2}$, $g=-\frac{3}{2}$, $\alpha=\sqrt{3}$, $\beta=1$, $\cos\theta=\frac{1}{2}\sqrt{3}$, $c=\frac{1}{2}\sqrt{(2-\sqrt{3})}$; donc l'intégrale cherchée est
\[ Q^{1}=-\left(\frac{3+2\sqrt{3}}{2\sqrt[4]{3}}\right)F^{1}(c)+\frac{3\sqrt{3}}{\sqrt[4]{3}}E^{1}(c). \]
\ednote{The numerator is printed $3+2\sqrt{3}$; with $f=g=-\frac{3}{2}$, $\alpha=\sqrt{3}$, $\beta=1$ the formula for $X^{1}$ gives $3+3\sqrt{3}$. An apparent misprint.}

\origpage{59}
Maintenant puisqu'on a $P^{1}Q^{1}=\frac{3}{4}\pi$, on aura une relation entre trois fonctions elliptiques, qui, avec celle qu'on a trouvée, offre ces deux résultats:
\[ \begin{gathered} \frac{\pi}{4}=F^{1}(b)\cdot\left[E^{1}(c)-\left(\frac{1+\sqrt{3}}{2\sqrt{3}}\right)F^{1}(c)\right]\\ \frac{\pi}{4}=F^{1}(c)\cdot\left[E^{1}(b)-\left(\frac{\sqrt{3}-1}{2\sqrt{3}}\right)F^{1}(b)\right]; \end{gathered} \]
d'où l'on voit que les fonctions de seconde espèce $E^{1}(c)$, $E^{1}(b)$ peuvent, dans ce cas particulier, s'exprimer par les deux fonctions de première espèce $F^{1}(b)$, $F^{1}(c)$. A ces deux relations qui sont déjà fort remarquables, il faut en joindre une troisième $F^{1}(b)=\sqrt{3}\,F^{1}(c)$, qui sera démontrée dans l'exemple suivant.

\section*{EXEMPLE III.}

\subsection*{(41).}

Soit proposé d'évaluer l'intégrale $R=\displaystyle\int dz\,(1-z^{3})^{-\frac{2}{3}}$, prise depuis $z=0$ jusqu'à $z=1$.

On fera d'abord $1-z^{3}=\left(\frac{z}{y}\right)^{\frac{3}{2}}$, ce qui donnera la transformée $R=\displaystyle\int\frac{dy}{\sqrt{(4y^{3}+1)}}$ à intégrer depuis $y=0$ jusqu'à $y=\infty$. Soit ensuite $m=\sqrt[3]{4}$ et $my=x^{2}-1$, on aura la nouvelle transformée $R=\frac{2}{m}\displaystyle\int\frac{dx}{\sqrt{(x^{4}-3x^{2}+3)}}$, qu'il faut intégrer depuis $x=1$ jusqu'à $x=\infty$. Cette intégrale est, au coefficient près, la même que l'intégrale $P$ de l'exemple précédent; ainsi ayant égard aux limites de $R$, et faisant $\sqrt[4]{3}=n$, on aura la valeur cherchée
\[ R^{1}=\frac{4}{3mn}F^{1}(b); \]
mais il y a une autre manière de trouver la valeur de $R$.

Soit $1-z^{3}=\left(1-\frac{z}{y}\right)^{3}$, on trouvera d'abord la transformée $R=\displaystyle\int\frac{dy\sqrt{3}}{\sqrt{(4y^{3}-1)}}$, qu'il faut intégrer depuis $y=1$ jusqu'à $y=\infty$. Soit ensuite $my=1+x^{2}$, on aura $R=\frac{2\sqrt{3}}{m}\displaystyle\int\frac{dx}{\sqrt{(x^{4}+3x^{2}+3)}}$, nouvelle formule qu'il faut intégrer depuis $x=\sqrt{(m-1)}$ jusqu'à

\origpage{60}
$x=\infty$. Or cette intégrale est semblable à la formule $M$ de l'exemple I, et on obtiendra de même $R=\frac{2\sqrt{3}}{mn}F(c,\varphi)+\text{const.}$, en observant que cette intégrale doit être prise depuis la valeur de $\varphi$ qui donne $x=\sqrt{(m-1)}$ jusqu'à la valeur de $\varphi$ qui donne $x=\infty$; celle-ci est $\varphi=\frac{1}{2}\pi$, l'autre étant nommée $\theta$, on aura $R^{1}=\frac{2\sqrt{3}}{mn}[F^{1}(c)-F(c,\theta)]$. Mais en général,
\[ \cos^{2}\varphi=\frac{-x^{2}-\frac{3}{2}+\sqrt{(x^{4}+3x^{2}+3)}}{2c^{2}\sqrt{3}}: \]
donc en faisant $x^{2}=m-1$, il viendra
\[ \cos^{2}\theta=\frac{-m-\frac{1}{2}+\sqrt{(m^{2}+m+1)}}{2c^{2}\sqrt{3}}=\frac{(m-1)^{2}}{2\sqrt{3}-3}. \]
Or, d'après l'article 24, cette valeur de $\theta$ est celle qui pour le module $c=\frac{1}{2}\sqrt{(2-\sqrt{3})}$, donne $F(\theta)=\frac{1}{3}F^{1}$; donc
\[ R^{1}=\frac{4\sqrt{3}}{3mn}F^{1}(c). \]
Comparant cette valeur à celle qu'on a trouvée par l'autre méthode, il en résulte cette nouvelle relation
\[ F^{1}(b)=\sqrt{3}\cdot F^{1}(c), \]
laquelle étant jointe aux deux déjà trouvées, fait voir qu'une seule des quatre transcendantes $F^{1}(c)$, $F^{1}(b)$, $E^{1}(c)$, $E^{1}(b)$ suffit pour déterminer les trois autres. On a, par exemple, les équations
\[ \begin{gathered} \frac{\pi}{4\sqrt{3}}=F^{1}(c)\left[E^{1}(c)-\left(\frac{\sqrt{3}+1}{2\sqrt{3}}\right)F^{1}(c)\right]\\ \frac{\pi\sqrt{3}}{4}=F^{1}(b)\left[E^{1}(b)-\left(\frac{\sqrt{3}-1}{2\sqrt{3}}\right)F^{1}(b)\right], \end{gathered} \]
qui servent à déterminer la fonction $E^{1}(c)$ par le moyen de $F^{1}(c)$, et la fonction $E^{1}(b)$ par le moyen de $F^{1}(b)$.

\origpage{61}

\section*{Théorème sur les fonctions de première et de seconde espèce, dont les modules sont complémens l'un de l'autre.}

\subsection*{(42).}

Dans les comparaisons qu'on vient d'établir entre les fonctions $F^{1}(c)$, $E^{1}(c)$, qui se rapportent au module $c=\frac{1}{2}\sqrt{(2-\sqrt{3})}=\sin 15^{\circ}$, et les fonctions $F^{1}(b)$, $E^{1}(b)$, qui se rapportent au module complémentaire $b=\frac{1}{2}\sqrt{(2+\sqrt{3})}=\cos 15^{\circ}$, les trois équations trouvées conduisent à ce résultat remarquable
\[ \frac{\pi}{2}=F^{1}(c)E^{1}(b)+F^{1}(b)E^{1}(c)-F^{1}(b)F^{1}(c)\ldots(d') \]
où l'on voit que les deux quantités $b$, $c$ peuvent être échangées entre elles, et qu'ainsi cette équation est vérifiée dans deux cas, celui de $c=\sin 15^{\circ}$, et celui de $c=\sin 75^{\circ}$. Il serait facile de démontrer directement qu'elle est encore vraie dans deux autres cas, lorsque $c$ est infiniment petit, et lorsque $c=\sqrt{\frac{1}{2}}=b$; mais nous allons prouver généralement qu'elle a lieu quel que soit $c$.

Pour abréger la notation, désignons simplement par $F$, $E$, les quantités $F^{1}(c)$, $E^{1}(c)$, et par $F'$, $E'$ les quantités $F^{1}(b)$, $E^{1}(b)$, et supposons
\[ P=FE'+F'E-FF', \]
$P$ étant une fonction de $c$ encore inconnue.

Je différentie les deux membres par rapport à $c$ qui est la seule variable qu'ils contiennent. Or ayant $E(\varphi)=\displaystyle\int\Delta\,d\varphi$, $F(\varphi)=\displaystyle\int\frac{d\varphi}{\Delta}$, $\Delta^{2}=1-c^{2}\sin^{2}\varphi$, la différentiation donne
\[ \begin{aligned} \frac{dE}{dc}&=-\int\frac{c\,d\varphi\sin^{2}\varphi}{\Delta}=\frac{1}{c}(E-F)\\ \frac{dF}{dc}&=\int\frac{c\,d\varphi\sin^{2}\varphi}{\Delta^{3}}=\frac{1}{c}\int\frac{d\varphi}{\Delta^{3}}-\frac{1}{c}\int\frac{d\varphi}{\Delta}. \end{aligned} \]
Mais par les formules de l'art.\ 9, on a $\displaystyle\int\frac{d\varphi}{\Delta^{3}}=\frac{1}{b^{2}}\displaystyle\int\Delta\,d\varphi-\frac{c^{2}\sin\varphi\cos\varphi}{b^{2}\Delta}$, et dans le cas de $\varphi=\frac{1}{2}\pi$ dont il s'agit, le second terme s'évanouit: ainsi on aura
\[ \frac{dF}{dc}=\frac{1}{b^{2}c}(E-b^{2}F). \]

\origpage{62}
On aura semblablement $\frac{dE'}{db}=\frac{1}{b}(E'-F')$, $\frac{dF'}{db}=\frac{1}{c^{2}b}(E'-c^{2}F')$, et parce que $b\,db+c\,dc=0$, on en déduira
\[ \begin{aligned} \frac{dE'}{dc}&=-\frac{c}{b^{2}}(E'-F')\\ \frac{dF'}{dc}&=-\frac{1}{b^{2}c}(E'-c^{2}F'). \end{aligned} \]
Substituant ces valeurs dans celle de $dP$, on aura $dP=0$; donc $P=\text{const.}$ Mais on a trouvé dans un cas particulier $P=\frac{1}{2}\pi$; donc l'équation $(d')$ a lieu généralement, quel que soit $c$.

Lorsque $c=\sqrt{\frac{1}{2}}=b$, l'équation $(d')$ donne
\[ \frac{\pi}{2}=F^{1}(2E^{1}-F^{1}). \]
Ainsi dans ce cas particulier, $E^{1}$ se détermine encore par $F^{1}$.

Il peut y avoir quelques autres cas particuliers où la fonction de seconde espèce $E^{1}(c)$ s'exprime par la fonction de seconde\ednote{Printed ``seconde''; $F^{1}(c)$ is the function of the first kind, so ``première'' is meant. An apparent misprint.} espèce $F^{1}(c)$; nous ferons voir en effet que chaque cas particulier connu en fait connaître une infinité d'autres; mais la détermination générale paraît impossible d'après les recherches suivantes.

\section*{Equations différentielles qui expriment la liaison mutuelle des fonctions $E$ et $F$.}

\subsection*{(43).}

Si l'on différentie par rapport à $c$ les deux formules $E=\displaystyle\int\Delta\,d\varphi$, $F=\displaystyle\int\frac{d\varphi}{\Delta}$, on aura comme ci-dessus,
\[ \begin{aligned} \frac{dE}{dc}&=\int\frac{-c\,d\varphi\sin^{2}\varphi}{\Delta}=\frac{1}{c}(E-F)\\ \frac{dF}{dc}&=\frac{1}{b^{2}c}(E-b^{2}F)-\frac{c}{b^{2}}\cdot\frac{\sin\varphi\cos\varphi}{\Delta}. \end{aligned} \]
Il en résulte les deux formules
\[ \begin{aligned} F&=E-c\,\frac{dE}{dc}\\ E&=b^{2}\left(E+c\,\frac{dF}{dc}\right)+\frac{c^{2}\sin\varphi\cos\varphi}{\Delta}, \end{aligned}\ldots(e') \]
\ednote{Inside the bracket the print reads $E$, though the glyph also carries a reader's pen mark and the reading is not certain; the derivation and the restatement on p. 64 require $F$. An apparent misprint.}
qui contiennent les relations mutuelles des fonctions $E$ et $F$.

\origpage{63}

On peut simplifier encore ces équations en faisant $F=\frac{M}{c}$ et $E=Nc$, ce qui donne
\[ \begin{aligned} M&=-c^{3}\frac{dN}{dc}\\ N&=\frac{b^{2}}{c}\cdot\frac{dM}{dc}+\frac{c\sin\varphi\cos\varphi}{\Delta}. \end{aligned} \]
Mais quelque simples que soient celles-ci, on ne saurait les intégrer par les moyens ordinaires, même quand on supposerait $\varphi=\frac{1}{2}\pi$, ce qui changerait les fonctions $F$ et $E$ en $F^{1}$ et $E^{1}$. Ainsi il n'y a guère d'espérance de déterminer généralement $E^{1}$ par $F^{1}$, et encore moins la fonction indéfinie $E$ par la fonction $F$; ce qui maintient et justifie la distinction que nous avons faite entre les fonctions elliptiques de la première et de la seconde espèce.

Il se présente à ce sujet une observation assez importante. On verra par les méthodes qui seront ci-après exposées, que la fonction $F$ peut être exprimée en termes finis par la fonction $E$ et une autre fonction de même espèce, dont le module et l'amplitude se déduisent suivant une loi connue du module et de l'amplitude qui conviennent à la fonction $E$. Cette expression de $F$ par deux fonctions de l'espèce de $E$, est une véritable intégrale qui doit satisfaire aux équations $(e')$, et qui ne rentre pas dans les procédés ordinaires de l'intégration. Le succès qu'on obtient dans ce cas particulier, peut donc devenir un exemple utile dans d'autres recherches.

\subsection*{(44).}

Revenons aux équations $(e')$, et voyons quel pourrait être leur usage.

Si on avait une table d'arcs d'ellipse dressée pour toutes les valeurs de $c$ et $\varphi$, à des intervalles égaux et suffisamment rapprochés, cette même table pourrait offrir pour chaque arc $E$, le coefficient différentiel $\frac{dE}{dc}$, puisque si $\alpha$ est la différence entre deux valeurs de $c$ consécutives, le coefficient différentiel se déduit des différences finies par la formule
\[ \alpha\frac{dE}{dc}=\Delta E-\tfrac{1}{2}\Delta^{2}E+\tfrac{1}{3}\Delta^{3}E-\text{etc.}, \]
$\Delta E$, $\Delta^{2}E$, etc. étant les différences premières, secondes, etc. relatives

\origpage{64}
à $c$, données immédiatement par les nombres de la Table. Une Table ainsi disposée, ferait connaître la fonction $F$ pour toute valeur de $c$ et $\varphi$, puisqu'on a $F=E-c\frac{dE}{dc}$. C'est le moyen que nous avions proposé dans les Mémoires cités de 1786, pour éviter l'emploi des arcs d'hyperbole, ou celui des fonctions $F$.

Une Table des valeurs de la fonction $F$, dressée de la même manière avec le coefficient différentiel $\frac{dF}{dc}$ correspondant à chaque terme, serait propre également à remplacer les arcs d'ellipse, puisque la seconde des équations $(e')$ donne
\[ E=b^{2}\left(F+c\frac{dF}{dc}\right)+\frac{c^{2}\sin\varphi\cos\varphi}{\Delta}. \]

Le travail pour former une Table un peu étendue, telle que nous venons de la proposer, est trop considérable pour croire qu'il soit jamais entrepris; si cependant les intervalles par lesquels on fait croître $c$ et $\varphi$ n'étaient pas trop rapprochés, la Table pourrait encore être fort utile, sans que le travail de sa construction excédât la patience d'un calculateur zélé.

Je proposerais donc comme chose fort praticable, que l'on construisît une Table d'arcs d'ellipse ou de fonctions $E(c,\varphi)$ pour toutes les valeurs de $\varphi$, de degrés en degrés, et pour toutes les valeurs du module $c$ mis sous la forme $\sin\theta$, aussi de degrés en degrés. Cette Table ne contiendrait que 8100 nombres, calculés avec sept chiffres significatifs; et en y joignant les différences premières et secondes, on pourrait étendre son usage à toutes les valeurs de $c$ et $\varphi$. Il serait facile dans chaque cas particulier, de tirer de cette Table la valeur du coefficient $\frac{dE}{dc}$, au moins pour les valeurs de $c$ et de $\varphi$ comprises dans la Table. Enfin on pourrait à la rigueur éviter entièrement la formation des coefficiens $\frac{dE}{dc}$, ou celle d'une Table particulière pour les fonctions $F$, puisqu'il sera démontré ci-après que toute fonction $F$ ou tout coefficient $\frac{dE}{dc}$, peut s'exprimer par deux arcs d'ellipse.

\origpage{65}

\section*{Développement des fonctions $F^{1}$ et $E^{1}$ en séries.}

\subsection*{(45).}

On peut déduire des équations $(e')$ deux équations différentielles du second ordre propres à déterminer séparément les fonctions $F$ et $E$; ces équations sont
\[ \begin{aligned} (1-c^{2})\frac{ddF}{dc^{2}}+\frac{1-3c^{2}}{c}\cdot\frac{dF}{dc}-F+\frac{\sin\varphi\cos\varphi}{\Delta^{3}}&=0\\ (1-c^{2})\frac{ddE}{dc^{2}}+\frac{1-c^{2}}{c}\cdot\frac{dE}{dc}+E-\frac{\sin\varphi\cos\varphi}{\Delta}&=0; \end{aligned} \]
et lorsqu'on considère les fonctions complètes, ces équations deviennent
\[ \begin{aligned} (1-c^{2})\frac{ddF^{1}}{dc^{2}}+\frac{1-3c^{2}}{c}\cdot\frac{dF^{1}}{dc}-F^{1}&=0\\ (1-c^{2})\frac{ddE^{1}}{dc^{2}}+\frac{1-c^{2}}{c}\cdot\frac{dE^{1}}{dc}+E^{1}&=0. \end{aligned} \]
Celles-ci se simplifieront encore en faisant $c=\sin\theta$, ce qui donnerait
\[ \begin{aligned} \frac{ddF^{1}}{d\theta^{2}}+\cot 2\theta\cdot\frac{dF^{1}}{d\theta}-F^{1}&=0\\ \frac{ddE^{1}}{d\theta^{2}}+\frac{1}{\sin 2\theta}\cdot\frac{dE^{1}}{d\theta}+E^{1}&=0. \end{aligned} \]
\ednote{Misprints, kept as printed: the substitution $c=\sin\theta$ gives $2\cot 2\theta$ in the first equation and $\frac{2}{\sin 2\theta}$ in the second. The print has no 2 before $\cot 2\theta$ and a 1 over $\sin 2\theta$; a reader has added a handwritten 2 in both places.}
On peut enfin regarder $F^{1}$ et $E^{1}$ comme des fonctions du module complémentaire $b$, ce qui donnera les équations différentielles
\[ \begin{aligned} (1-b^{2})\frac{ddF^{1}}{db^{2}}+\frac{1-3b^{2}}{b}\cdot\frac{dF^{1}}{db}-F^{1}&=0\\ (1-b^{2})\frac{ddE^{1}}{db^{2}}-\left(\frac{1+b^{2}}{b}\right)\cdot\frac{dE^{1}}{db}+E^{1}&=0; \end{aligned} \]
d'où l'on voit que l'équation en $F^{1}$ est de la même forme, soit que l'on considère $F^{1}$ comme fonction de $b$ ou comme fonction de $c$.

Ces équations sont utiles pour faire connaître la loi du développement des fonctions $F^{1}$, $E^{1}$ en séries. Il n'y a aucune difficulté à développer ces fonctions suivant les puissances de $c$; car les expressions $F=\displaystyle\int d\varphi\,(1-c^{2}\sin^{2}\varphi)^{-\frac{1}{2}}$, $E=\displaystyle\int d\varphi\,(1-c^{2}\sin^{2}\varphi)^{\frac{1}{2}}$ étant intégrées

\origpage{66}
depuis $\varphi=0$ jusqu'à $\varphi=\frac{1}{2}\pi$, on en tire immédiatement
\[ \begin{aligned} F^{1}&=\frac{\pi}{2}\left(1+\frac{1^{2}}{2^{2}}c^{2}+\frac{1^{2}\cdot 3^{2}}{2^{2}\cdot 4^{2}}\cdot c^{4}+\frac{1^{2}\cdot 3^{2}\cdot 5^{2}}{2^{2}\cdot 4^{2}\cdot 6^{2}}c^{6}+\text{etc.}\right)\\ E^{1}&=\frac{\pi}{2}\left(1-\frac{1}{2^{2}}c^{2}-\frac{1^{2}}{2^{2}\cdot 4^{2}}\cdot 3c^{4}-\frac{1^{2}\cdot 3^{2}}{2^{2}\cdot 4^{2}\cdot 6^{2}}\cdot 5c^{6}-\text{etc.}\right) \end{aligned} \]
Mais ces séries ne sont plus suffisamment convergentes, lorsque $c$ est fort près de l'unité, et alors il convient de les ordonner suivant les puissances de $b$, en considérant $b$ comme très-petit.

Or d'après l'équation $E^{1}=b^{2}\left(F^{1}+c\frac{dF^{1}}{dc}\right)=b^{2}\cdot\frac{d(cF^{1})}{dc}$, on peut pour première approximation, faire $E^{1}=1$, ce qui donnera $d(cF^{1})=\frac{dc}{b^{2}}=\frac{dc}{1-c^{2}}$, et par conséquent $cF^{1}=\frac{1}{2}\log\left(\frac{1+c}{1-c}\right)$; mais dans le même cas, on a $c=1-\frac{1}{2}b^{2}$, donc la première valeur approchée de $F^{1}$ est $F^{1}=\log\left(\frac{4}{b}\right)$.

Soit maintenant $F^{1}=P\log\left(\frac{4}{b}\right)+Q$, $P$ et $Q$ étant des suites ordonnées suivant les puissances de $b$; si on substitue cette valeur dans l'équation différentielle
\[ (1-b^{2})\frac{ddF^{1}}{db^{2}}+\frac{1-3b^{2}}{b}\cdot\frac{dF^{1}}{db}-F^{1}=0, \]
on trouvera que l'équation pour déterminer $P$ est absolument semblable à celle qui détermine $F^{1}$; celle-ci est de la même forme, soit que $F$ soit considéré comme fonction de $c$ ou comme fonction de $b$; ainsi on aura d'abord
\[ P=1+\frac{1^{2}}{2^{2}}b^{2}+\frac{1^{2}\cdot 3^{2}}{2^{2}\cdot 4^{2}}b^{4}+\frac{1^{2}\cdot 3^{2}\cdot 5^{2}}{2^{2}\cdot 4^{2}\cdot 6^{2}}b^{6}+\text{etc.} \]
Désignons les coefficiens successifs par $1$, $m'$, $m''$, etc., en sorte qu'on ait
\[ P=1+m'b^{2}+m''b^{4}+m'''b^{6}+\text{etc.}; \]
nous supposerons ensuite
\[ \begin{aligned} F^{1}={}&(1+m'b^{2}+m''b^{4}+m'''b^{6}+\text{etc.})\log\frac{4}{b}\\ &-m'A'b^{2}-m''A''b^{4}-m'''A'''b^{6}-\text{etc.} \end{aligned} \]
Substituant cette valeur dans l'équation différentielle ci-dessus, on

\origpage{67}
aura pour déterminer les nouveaux coefficiens $A'$, $A''$, etc., l'équation suivante, dont les termes suivent une loi très-simple:
\[ \begin{aligned} 0={}&2&&+6m'b^{2}&&+10m''b^{4}&&+14m'''b^{6}&&+18m^{\text{iv}}b^{8}&&+\text{etc.}\\ &-4m'&&-6m''b^{2}&&-12m'''b^{4}&&-16m^{\text{iv}}b^{6}&&-20m^{\text{v}}b^{8}&&-\text{etc.}\\ &-2^{2}m'A'&&-4^{2}m''A''b^{2}&&-6^{2}m'''A'''b^{4}&&-8^{2}m^{\text{iv}}A^{\text{iv}}b^{6}&&-10^{2}m^{\text{v}}A^{\text{v}}b^{8}&&-\text{etc.}\\ &&&+3^{2}m'A'b^{2}&&+5^{2}m''A''b^{4}&&+7^{2}m'''A'''b^{6}&&+9^{2}m^{\text{iv}}A^{\text{iv}}b^{8}&&+\text{etc.} \end{aligned} \]
\ednote{Second row: $-6m''b^{2}$ as printed; the row's law ($-4k\,m^{(k)}b^{2k-2}$: $-4m'$, $-12m'''b^{4}$, $-16m^{\text{iv}}b^{6}$) requires $-8m''b^{2}$.}
Observant ensuite qu'on a $2^{2}m'=1$, $4^{2}m''=3^{2}m'$, $6^{2}m'''=5^{2}m''$, etc., on trouve successivement
\[ \begin{aligned} A'&=1\\ A''&=1+\frac{2}{3\cdot 4}\\ A'''&=1+\frac{2}{3\cdot 4}+\frac{2}{5\cdot 6}\\ A^{\text{iv}}&=1+\frac{2}{3\cdot 4}+\frac{2}{5\cdot 6}+\frac{2}{7\cdot 8}\\ &\text{etc.} \end{aligned} \]
de sorte qu'on a en général $A^{(n)}=A^{(n-1)}+\frac{2}{(2n-1)2n}$. D'après cette loi, il est facile de voir ce que devient $A^{(n)}$ lorsque $n$ est très-grand; considérons pour cet effet la suite
\[ y=\frac{2x^{2}}{1\cdot 2}+\frac{2x^{4}}{3\cdot 4}+\frac{2x^{6}}{5\cdot 6}+\frac{2x^{8}}{7\cdot 8}+\text{etc.}, \]
laquelle peut être mise sous cette forme
\[ \begin{aligned} y={}&2x\left(x+\frac{x^{3}}{3}+\frac{x^{5}}{5}+\frac{x^{7}}{7}+\text{etc.}\right)\\ &-\left(x^{2}+\frac{x^{4}}{2}+\frac{x^{6}}{3}+\frac{x^{8}}{4}+\text{etc.}\right); \end{aligned} \]
on aura en sommant ces suites, $y=x\log\left(\frac{1+x}{1-x}\right)+\log(1-x^{2})=(1+x)\log(1+x)+(1-x)\log(1-x)$. Soit $x=1$, on aura $y=2\log 2$; donc la limite des quantités $A'$, $A''$, $A'''$, etc. est $2\log 2$ ou 1,386, etc. Cela posé, la valeur complète de $F^{1}$ se développe ainsi:

\origpage{68}
\[ \begin{aligned} F^{1}={}&\log\frac{4}{b}+\frac{1^{2}}{2^{2}}b^{2}\left(\log\frac{4}{b}-1\right)\\ &+\frac{1^{2}\cdot 3^{2}}{2^{2}\cdot 4^{2}}b^{4}\left(\log\frac{4}{b}-1-\frac{2}{3\cdot 4}\right)\\ &+\frac{1^{2}\cdot 3^{2}\cdot 5^{2}}{2^{2}\cdot 4^{2}\cdot 6^{2}}b^{6}\left(\log\frac{4}{b}-1-\frac{2}{3\cdot 4}-\frac{2}{5\cdot 6}\right)\\ &+\frac{1^{2}\cdot 3^{2}\cdot 5^{2}\cdot 7^{2}}{2^{2}\cdot 4^{2}\cdot 6^{2}\cdot 8^{2}}b^{8}\left(\log\frac{4}{b}-1-\frac{2}{3\cdot 4}-\frac{2}{5\cdot 6}-\frac{2}{7\cdot 8}\right)\\ &+\text{etc.} \end{aligned} \]
Connaissant $F^{1}$, on aura immédiatement $E^{1}$ d'après l'équation $E^{1}=b^{2}F^{1}+b^{2}c\frac{dF^{1}}{dc}=b^{2}F^{1}-b(1-b^{2})\frac{dF^{1}}{db}$, et il en résulte
\[ \begin{aligned} E^{1}={}&1+\frac{1}{2}b^{2}\left(\log\frac{4}{b}-\frac{1}{1\cdot 2}\right)\\ &+\frac{1^{2}}{2^{2}}\cdot\frac{3}{4}b^{4}\left(\log\frac{4}{b}-A'-\frac{1}{3\cdot 4}\right)\\ &+\frac{1^{2}\cdot 3^{2}}{2^{2}\cdot 4^{2}}\cdot\frac{5}{6}b^{6}\left(\log\frac{4}{b}-A''-\frac{1}{5\cdot 6}\right)\\ &+\frac{1^{2}\cdot 3^{2}\cdot 5^{2}}{2^{2}\cdot 4^{2}\cdot 6^{2}}\cdot\frac{7}{8}b^{8}\left(\log\frac{4}{b}-A'''-\frac{1}{7\cdot 8}\right)\\ &+\text{etc.}, \end{aligned} \]
expression dont la loi est manifeste et qui s'accorde avec celle que nous avons donnée sous une autre forme (\emph{Mém. de l'Acad.}, 1786, pag.~630).

\section*{Des changemens qu'on peut faire subir au paramètre dans les fonctions elliptiques de la troisième espèce.}

\subsection*{(46).}

Soit $p=\frac{\operatorname{tang}\varphi}{\Delta}$, et soit $\alpha$ une constante indéterminée, on trouve par la différentiation,
\[ \frac{dp}{1+\alpha p^{2}}=\frac{d\varphi}{\Delta}\cdot\frac{1-c^{2}\sin^{4}\varphi}{(1-c^{2}\sin^{2}\varphi)\cos^{2}\varphi+\alpha\sin^{2}\varphi}. \]
Supposons que le dénominateur $(1-c^{2}\sin^{2}\varphi)\cos^{2}\varphi+\alpha\sin^{2}\varphi$ soit égal au produit des deux facteurs $(1+n\sin^{2}\varphi)\left(1+\frac{c^{2}}{n}\sin^{2}\varphi\right)$, la

\origpage{69}
valeur de $\alpha$ devra être
\[ \alpha=(1+n)\left(1+\frac{c^{2}}{n}\right) \]
et alors l'équation différentielle se décomposant ainsi,
\[ \frac{dp}{1+\alpha p^{2}}=\frac{d\varphi}{\Delta}\left[\frac{1}{1+n\sin^{2}\varphi}+\frac{1}{1+\frac{c^{2}}{n}\sin^{2}\varphi}-1\right], \]
son intégrale est
\[ \Pi(n)+\Pi\left(\frac{c^{2}}{n}\right)=F+\frac{1}{\sqrt{\alpha}}\operatorname{arc}\operatorname{tang}\frac{\sqrt{\alpha}\operatorname{tang}\varphi}{\Delta}\ldots(f'), \]
formule très-remarquable au moyen de laquelle toute fonction $\Pi$ dont le paramètre est plus grand que le module $c$, peut toujours être transformée en une autre dont le paramètre sera plus petit que $c$; car des deux paramètres $n$ et $\frac{c^{2}}{n}$, il est évident que l'un est toujours plus grand que $c$ et l'autre plus petit. Les fonctions $\Pi$ qui se transforment ainsi l'une dans l'autre, ont d'ailleurs le même module $c$ et la même amplitude $\varphi$; c'est pourquoi, au lieu de les désigner par $\Pi(n,c,\varphi)$, $\Pi\left(\frac{c^{2}}{n},c,\varphi\right)$, nous les désignons simplement par $\Pi(n)$, $\Pi\left(\frac{c^{2}}{n}\right)$, en ne signalant que l'élément par lequel les deux fonctions diffèrent l'une de l'autre.

La formule générale $(f')$ appliquée aux cas de $n=c$ et de $n=-c$, fournit ces deux corollaires,
\[ \begin{aligned} \Pi(c)&=\tfrac{1}{2}F+\frac{\frac{1}{2}}{1+c}\operatorname{arc}\operatorname{tang}\frac{(1+c)\operatorname{tang}\varphi}{\Delta}\\ \Pi(-c)&=\tfrac{1}{2}F+\frac{\frac{1}{2}}{1-c}\operatorname{arc}\operatorname{tang}\frac{(1-c)\operatorname{tang}\varphi}{\Delta}. \end{aligned} \]
Ainsi dans ces deux cas, la fonction de troisième espèce se réduit immédiatement à une fonction de première espèce.

\subsection*{(47).}

Il y a trois cas à considérer dans la formule $(f')$, selon que $\alpha$ est positif, zéro ou négatif.

Le coefficient $\alpha$ sera positif toutes les fois que le paramètre $n$ sera positif, ou toutes les fois que $n$ étant négatif, il sera compris entre

\origpage{70}
$-1$ et $-c^{2}$; en d'autres termes $\alpha$ sera positif si le paramètre $n$ est de l'une des deux formes $n=\cot^{2}\theta$, $n=-1+b^{2}\sin^{2}\theta$. Alors l'équation $(f')$ contiendra un arc de cercle réel, et par cette équation on ramènera l'une à l'autre les deux fonctions $\Pi(n)$, $\Pi\left(\frac{c^{2}}{n}\right)$.

Dans ce premier cas, si l'on fait $\varphi=\frac{1}{2}\pi$, on aura entre les fonctions complètes cette relation
\[ \Pi^{1}(n)+\Pi^{1}\left(\frac{c^{2}}{n}\right)=F^{1}+\frac{\pi}{2\sqrt{\alpha}}. \]

\subsection*{(48).}

Si l'on fait $\alpha=0$, on aura $n=-1$ ou $n=-c^{2}$, alors l'intégrale $\displaystyle\int\frac{dp}{1+\alpha p^{2}}$ se réduisant à $p$, l'équation $(f')$ devient
\[ \Pi(-1)+\Pi(-c^{2})=F+\frac{\operatorname{tang}\varphi}{\Delta}. \]
Ce résultat est facile à vérifier; on a en effet par les réductions déjà connues,
\[ \begin{aligned} \Pi(-1)&=\int\frac{d\varphi}{\Delta\cos^{2}\varphi}=F-\frac{1}{b^{2}}E+\frac{1}{b^{2}}\Delta\operatorname{tang}\varphi\\ \Pi(-c^{2})&=\int\frac{d\varphi}{\Delta^{3}}=\frac{1}{b^{2}}E-\frac{c^{2}\sin\varphi\cos\varphi}{b^{2}\Delta}; \end{aligned} \]
et la somme de ces deux quantités se réduit à $F+\frac{\operatorname{tang}\varphi}{\Delta}$.

La fonction $\Pi(-1)$ qui se ramène immédiatement aux fonctions de la première et de la seconde espèce, est celle qui donne la rectification de l'hyperbole (art.\ 13). Elle a une valeur infinie lorsque $\varphi=\frac{1}{2}\pi$, parce qu'alors elle représente la longueur totale de la courbe jusqu'à son extrémité infinie; et c'est aussi ce que donne la formule précédente; mais cette formule serait en défaut si on voulait faire $\varphi>\frac{1}{2}\pi$; elle semblerait donner une valeur finie pour $\Pi(-1)$, tandis que cette valeur, composée de la partie où $\varphi=\frac{1}{2}\pi$ et d'une autre partie, est nécessairement infinie. C'est du moins ce qui paraît résulter de la formule intégrale $\Pi(-1)$ considérée en elle-même, et sans rapport à aucune courbe; car nous supposons toujours $\Delta$ positif.

\subsection*{(49).}

Il reste à examiner le cas où $\alpha$ est négatif; ce cas aura lieu

\origpage{71}
toutes les fois que le paramètre sera négatif, mais non compris entre $-1$ et $-c^{2}$; de sorte que $n$ devra être représenté par l'une ou l'autre des formules $n=-\frac{1}{\sin^{2}\theta}$, $n=-c^{2}\sin^{2}\theta$.

Soit donc $n=-c^{2}\sin^{2}\theta$, et $\alpha=-\beta$, on aura $\beta=\frac{\cos^{2}\theta}{\sin^{2}\theta}(1-c^{2}\sin^{2}\theta)$, et l'intégrale $\displaystyle\int\frac{dp}{1+\alpha p^{2}}$ devient $\displaystyle\int\frac{dp}{1-\beta p^{2}}$ ou $\frac{1}{2\sqrt{\beta}}\log\left(\frac{1+p\sqrt{\beta}}{1-p\sqrt{\beta}}\right)$; donc alors l'équation $(f')$ devra être remplacée par la suivante:
\[ \Pi(n)+\Pi\left(\frac{c^{2}}{n}\right)=F+\frac{1}{2\sqrt{\beta}}\log\left(\frac{\Delta+\sqrt{\beta}\operatorname{tang}\varphi}{\Delta-\sqrt{\beta}\operatorname{tang}\varphi}\right). \]
Mettant dans cette équation les valeurs de $n$ et de $\beta$, et observant que $\sqrt{(1-c^{2}\sin^{2}\theta)}$ peut être désigné par $\Delta(\theta)$, tandis que $\sqrt{(1-c^{2}\sin^{2}\varphi)}$ l'est par $\Delta(\varphi)$, on aura cette formule générale pour le cas de $\alpha$ négatif:
\[ \Pi(-c^{2}\sin^{2}\theta)+\Pi\left(-\frac{1}{\sin^{2}\theta}\right)=F+\frac{\operatorname{tang}\theta}{2\Delta(\theta)}\log\left(\frac{\Delta(\varphi)\operatorname{tang}\theta+\Delta(\theta)\operatorname{tang}\varphi}{\Delta(\varphi)\operatorname{tang}\theta-\Delta(\theta)\operatorname{tang}\varphi}\right). \]
Il est à remarquer que le second membre de cette équation devient infini lorsque $\varphi=\theta$: en effet, comme on a
\[ \Pi\left(-\frac{1}{\sin^{2}\theta}\right)=\int\frac{d\varphi}{\left(1-\frac{\sin^{2}\varphi}{\sin^{2}\theta}\right)\Delta}, \]
on voit que le dénominateur de la différentielle est zéro lorsque $\varphi=\theta$, ce qui rend, pour ce cas, l'intégrale infinie.

Lorsqu'ensuite on suppose $\varphi>\theta$, le dénominateur devient négatif, et la valeur de $\Pi\left(-\frac{1}{\sin^{2}\theta}\right)$ redevient finie par la destruction mutuelle des parties infinies et de signes contraires. Mais alors la formule a besoin d'être rectifiée pour ne pas offrir le logarithme d'une quantité négative, et il faudra l'écrire ainsi:
\[ \Pi(-c^{2}\sin^{2}\theta)+\Pi\left(-\frac{1}{\sin^{2}\theta}\right)=F+\frac{\operatorname{tang}\theta}{2\Delta(\theta)}\log\frac{\Delta(\theta)\operatorname{tang}\varphi+\Delta(\varphi)\operatorname{tang}\theta}{\Delta(\theta)\operatorname{tang}\varphi-\Delta(\varphi)\operatorname{tang}\theta}: \]
elle aura lieu depuis $\varphi=\theta$ jusqu'à $\varphi=\frac{1}{2}\pi$.

Lorsque $\varphi=\frac{1}{2}\pi$, cette formule donne entre les fonctions complètes

\origpage{72}
la relation
\[ \Pi^{1}(-c^{2}\sin^{2}\theta)+\Pi^{1}\left(-\frac{1}{\sin^{2}\theta}\right)=F^{1}, \]
où l'on voit que les logarithmes ont entièrement disparu.

J'observe sur cette dernière équation que la fonction $\Pi(-c^{2}\sin^{2}\theta)$, qui représente l'intégrale $\displaystyle\int\frac{d\varphi}{(1-c^{2}\sin^{2}\theta\sin^{2}\varphi)\Delta}$, est nécessairement plus grande que $\displaystyle\int\frac{d\varphi}{\Delta}$ ou $F$, et qu'ainsi $\Pi^{1}(-c^{2}\sin^{2}\theta)$ est plus grande que $F^{1}$; l'équation précédente ne peut donc subsister à moins que $\Pi^{1}\left(-\frac{1}{\sin^{2}\theta}\right)$ ne soit négatif. Or, ce résultat s'explique facilement par la nature de la fonction $\Pi\left(-\frac{1}{\sin^{2}\theta}\right)$ qui est positive depuis $\varphi=0$, jusqu'à $\varphi=\theta$, et négative depuis $\varphi=\theta$ jusqu'à $\varphi=\frac{1}{2}\pi$; il faut par conséquent que la partie négative soit plus grande que la partie positive.

\subsection*{(50).}

\emph{Au reste pour éviter toute anomalie étrangère à l'objet dont nous nous occupons, et pour ne considérer des fonctions $\Pi$ que celles qui sont positives et finies, nous ferons abstraction, dans tout ce qui suit, du cas où le paramètre $n$ est à la fois négatif et plus grand que l'unité.}

Si ce cas se rencontrait, on pourrait, au moyen de la formule de l'art.\ 49, réduire la fonction proposée $\Pi\left(-\frac{1}{\sin^{2}\theta}\right)$ à la fonction $\Pi(-c^{2}\sin^{2}\theta)$ qui n'est sujette à aucune difficulté; l'anomalie tomberait alors sur le terme logarithmique joint à cette fonction.

Cela posé, les fonctions $\Pi$, eu égard aux diverses valeurs du paramètre, se rapportent à trois cas principaux qui exigent des développemens particuliers.

\emph{Premier cas.} Si le paramètre $n$ est positif, on pourra toujours lui donner la forme $n=\cot^{2}\theta$, et on aura alors $\sqrt{\alpha}=\frac{\sqrt{(1-b^{2}\sin^{2}\theta)}}{\sin\theta\cos\theta}=\frac{\Delta(b,\theta)}{\sin\theta\cos\theta}$.

\emph{Second cas.} Si le paramètre $n$ est négatif, mais qu'il soit compris entre $-1$ et $-c^{2}$, on pourra le désigner par la formule $n=-1+b^{2}\sin^{2}\theta$, et alors on aura $\sqrt{\alpha}=\frac{b^{2}\sin\theta\cos\theta}{\sqrt{(1-b^{2}\sin^{2}\theta)}}=\frac{b^{2}\sin\theta\cos\theta}{\Delta(b,\theta)}$.

\origpage{73}

\emph{Troisième cas.} Si le paramètre $n$ est négatif et plus petit que $c^{2}$, on pourra faire $n=-c^{2}\sin^{2}\theta$, et on aura $\sqrt{(\alpha)}=\frac{\cos\theta}{\sin\theta}\Delta(\theta)\cdot\sqrt{-1}$.

On verra ci-après que les deux premiers cas pourraient être censés n'en faire qu'un, attendu qu'une fonction qui appartient à l'un de ces cas, peut être transformée en une fonction qui appartienne à l'autre cas. Il n'en est pas de même du troisième cas qui diffère essentiellement des deux autres; car les fonctions qui se rapportent aux deux premiers cas, entraînent toujours dans leur comparaison des arcs de cercle; tandis que les fonctions qui se rapportent au troisième, n'admettent dans leur comparaison que des logarithmes.

Il semblerait naturel d'ajouter à ces trois cas celui où l'on supposerait le paramètre $n$ imaginaire; mais nous prouverons ci-après que ce quatrième cas est inutile à considérer, et qu'on peut toujours y suppléer par des transformations convenables.

\subsection*{(51).}

Soit maintenant $p=\frac{\sin\varphi\cos\varphi}{\Delta}$, et $k$ un coefficient indéterminé, on aura par la différentiation,
\[ \frac{dp}{1+kp^{2}}=\frac{1-2\sin^{2}\varphi+c^{2}\sin^{4}\varphi}{1+(k-c^{2})\sin^{2}\varphi-k\sin^{4}\varphi}\cdot\frac{d\varphi}{\Delta}; \]
si on fait ensuite le dénominateur
\[ 1+(k-c^{2})\sin^{2}\varphi-k\sin^{4}\varphi=(1+n\sin^{2}\varphi)(1-m\sin^{2}\varphi), \]
il en résultera $k=mn$, $(1+n)(1-m)=b^{2}$, et l'équation différentielle prendra la forme
\[ \frac{dp}{1+kp^{2}}=\frac{d\varphi}{\Delta}\left[\frac{1+n}{n}\cdot\frac{1}{1+n\sin^{2}\varphi}-\frac{1-m}{m}\cdot\frac{1}{1-m\sin^{2}\varphi}-\frac{c^{2}}{mn}\right], \]
d'où l'on tire en intégrant,
\[ \frac{1+n}{n}\Pi(n)-\left(\frac{1-m}{m}\right)\Pi(-m)=\frac{c^{2}}{mn}F+\frac{1}{\sqrt{mn}}\operatorname{arc}\operatorname{tang}\frac{\sqrt{mn}\sin\varphi\cos\varphi}{\Delta}\ldots(g'). \]

Par cette formule les deux fonctions $\Pi(n)$, $\Pi(-m)$, peuvent être réduites l'une à l'autre, pourvu qu'entre les paramètres $n$ et $-m$, on ait la relation
\[ (1+n)(1-m)=b^{2}. \]

\origpage{74}
Cette relation est telle, que si l'on fait $n=\cot^{2}\theta$, on aura $-m=-1+b^{2}\sin^{2}\theta$; d'où l'on voit que les deux premiers cas du n° 50, n'en font à proprement parler qu'un seul, puisque la réduction d'une fonction à l'autre peut se faire immédiatement au moyen de la formule
\[ \begin{aligned} \Pi(\cot^{2}\theta)&=\frac{b^{2}\sin^{2}\theta\cos^{2}\theta}{1-b^{2}\sin^{2}\theta}\Pi(-1+b^{2}\sin^{2}\theta)\\ &\quad+\frac{c^{2}\sin^{2}\theta}{1-b^{2}\sin^{2}\theta}F+\frac{\sin\theta\cos\theta}{\Delta(b,\theta)}\operatorname{arc}\operatorname{tang}\frac{\sin\varphi\cos\varphi\Delta(b,\theta)}{\operatorname{tang}\theta\Delta(c,\varphi)}. \end{aligned} \]

\subsection*{(52).}

Si dans la formule $(g')$, on fait $m=c^{2}\sin^{2}\theta$, l'équation de condition $(1+n)(1-m)=b^{2}$, donnera $n(\cos^{2}\theta+b^{2}\sin^{2}\theta)=-c^{2}\cos^{2}\theta$, d'où il suit que $n$ pourra aussi être représenté par la valeur $n=-c^{2}\sin^{2}\lambda$, et alors on aura entre $\theta$ et $\lambda$ la relation $\sin^{2}\lambda(\cos^{2}\theta+b^{2}\sin^{2}\theta)=\cos^{2}\theta$, d'où résulte
\[ 1=b\operatorname{tang}\theta\operatorname{tang}\lambda; \]
c'est la même relation qui donne $F(\theta)+F(\lambda)=F^{1}$.

On aura donc dans ce cas la formule générale
\[ \begin{aligned} &\frac{1-c^{2}\sin^{2}\theta}{c^{2}\sin^{2}\theta}\Pi(-c^{2}\sin^{2}\theta)+\frac{1-c^{2}\sin^{2}\lambda}{c^{2}\sin^{2}\lambda}\Pi(-c^{2}\sin^{2}\lambda)\\ &=\frac{F}{\sin^{2}\theta\sin^{2}\lambda}-\frac{1}{2\sin\theta\sin\lambda}\log\left(\frac{\Delta+c^{2}\sin\theta\sin\lambda\sin\varphi\cos\varphi}{\Delta-c^{2}\sin\theta\sin\lambda\sin\varphi\cos\varphi}\right), \end{aligned} \]
\ednote{The right-hand side is printed without $c^{2}$ in its denominators; formula $(g')$ with $m=c^{2}\sin^{2}\theta$, $n=-c^{2}\sin^{2}\lambda$ gives $\frac{F}{c^{2}\sin^{2}\theta\sin^{2}\lambda}$ and $\frac{1}{2c^{2}\sin\theta\sin\lambda}$, and only these yield the case $\theta=\lambda$ on p.~75. An apparent misprint.}
au moyen de laquelle on pourra réduire l'une à l'autre les fonctions $\Pi(-c^{2}\sin^{2}\theta)$, $\Pi(-c^{2}\sin^{2}\lambda)$, pourvu qu'entre les paramètres on ait la relation $1=b\operatorname{tang}\theta\operatorname{tang}\lambda$.

Ainsi, non seulement les fonctions $\Pi$ dont le paramètre est négatif et plus grand que l'unité, peuvent se réduire aux fonctions dont le paramètre est plus petit que $c^{2}$ (art.\ 49), mais celles-ci peuvent encore être réduites au cas où le paramètre, toujours de forme $-c^{2}\sin^{2}\theta$, n'excède pas $1-b$; car dans le cas où l'on aurait $\theta=\lambda$, l'équation $1=b\operatorname{tang}\theta\operatorname{tang}\lambda$ donne $\sin^{2}\theta=\frac{1}{1+b}$, et par conséquent $c^{2}\sin^{2}\theta=1-b$; dans tout autre cas, l'un des paramètres $-c^{2}\sin^{2}\theta$, $-c^{2}\sin^{2}\lambda$, abstraction faite de son signe, sera nécessairement plus petit que $1-b$.

Le cas de $\theta=\lambda$ mérite d'être remarqué; alors la formule générale

\origpage{75}
donne
\[ \Pi(-1+b)=\frac{1+b}{2b}F-\frac{1}{4b}\log\left(\frac{\Delta+(1-b)\sin\varphi\cos\varphi}{\Delta-(1-b)\sin\varphi\cos\varphi}\right). \]
Ainsi la fonction de troisième espèce $\Pi(-1+b)$ se réduit indéfiniment à la première espèce, et on a en particulier pour l'expression de la fonction complète,
\[ \Pi^{1}(-1+b)=\frac{1+b}{2b}F^{1}. \]

On pourrait, en choisissant d'autres valeurs de $p$ ou d'autres différentielles que $\frac{dp}{1+kp^{2}}$, parvenir à d'autres formules de réduction pour les fonctions $\Pi$; mais la plupart de ces formules rentreraient dans les deux que nous avons trouvées, ou ne seraient qu'une combinaison de ces deux formules et de celles que nous exposerons dans le chapitre suivant; c'est pourquoi nous ne nous y arrêterons pas davantage.

\section*{Comparaison des fonctions elliptiques de la troisième espèce.}

\subsection*{(53).}

Considérons les deux fonctions semblables
\[ \Pi(\varphi)=\int\frac{d\varphi}{(1+n\sin^{2}\varphi)\Delta(\varphi)},\quad \Pi(\psi)=\int\frac{d\psi}{(1+n\sin^{2}\psi)\Delta\psi}, \]
et supposons, comme nous l'avons fait pour les fonctions de la première et de la seconde espèce, qu'on a l'équation $F(\varphi)+F(\psi)-F(\mu)=0$, $\mu$ étant une constante; alors si on fait $\Pi(\varphi)+\Pi(\psi)-\Pi(\mu)=Q$, on aura
\[ Q=\int\left[\frac{d\varphi}{(1+n\sin^{2}\varphi)\Delta(\varphi)}+\frac{d\psi}{(1+n\sin^{2}\psi)\Delta(\psi)}\right], \]
cette intégrale étant prise de manière qu'elle s'évanouisse lorsque $\varphi=0$ ou lorsque $\psi=\mu$.

Mais puisque $\mu$ est constant, on a $\frac{d\varphi}{\Delta(\varphi)}+\frac{d\psi}{\Delta(\psi)}=0$, d'où résulte
\[ Q=\int\frac{d\varphi}{\Delta(\varphi)}\cdot\frac{n(\sin^{2}\psi-\sin^{2}\varphi)}{1+n(\sin^{2}\psi+\sin^{2}\varphi)+n^{2}\sin^{2}\psi\sin^{2}\varphi}. \]

\origpage{76}
Or on a trouvé (art.\ 31) $d\varphi\Delta(\varphi)+d\psi\Delta(\psi)=c^{2}d(\sin\mu\sin\varphi\sin\psi)$; le premier membre se réduit à $\frac{d\varphi}{\Delta(\varphi)}[\Delta^{2}(\varphi)-\Delta^{2}(\psi)]=\frac{c^{2}d\varphi}{\Delta(\varphi)}(\sin^{2}\psi-\sin^{2}\varphi)$; donc en faisant $\sin^{2}\varphi+\sin^{2}\psi=p$, $\sin\varphi\sin\psi=q$, on a~$\ldots\ \frac{d\varphi}{\Delta(\varphi)}(\sin^{2}\psi-\sin^{2}\varphi)=\sin\mu\cdot dq$; donc enfin
\[ Q=\int\frac{n\sin\mu\cdot dq}{1+np+n^{2}q^{2}}. \]
Mais de l'équation $\cos\varphi\cos\psi-\sin\varphi\sin\psi\Delta(\mu)=\cos\mu$, on déduit $1-p+q^{2}=[\cos\mu+q\Delta(\mu)]^{2}$, et par conséquent
\[ p=\sin^{2}\mu-2q\cos\mu\Delta(\mu)+c^{2}q^{2}\sin^{2}\mu; \]
mettant cette valeur de $p$ dans la formule précédente, on a
\[ Q=\int\frac{n\sin\mu\cdot dq}{1+n\sin^{2}\mu-2nq\cos\mu\Delta(\mu)+q^{2}(n^{2}+nc^{2}\sin^{2}\mu)}. \]
Effectuant donc l'intégration, et faisant comme ci-dessus $\alpha=(1+n)\left(1+\frac{c^{2}}{n}\right)$, on aura $Q$ ou
\[ \Pi(\varphi)+\Pi(\psi)-\Pi(\mu)=\frac{1}{\sqrt{\alpha}}\operatorname{arc}\operatorname{tang}\left(\frac{n\sqrt{\alpha}\sin\mu\sin\varphi\sin\psi}{1+n-n\cos\mu\cos\varphi\cos\psi}\right)\ldots(h'). \]
C'est la formule générale qui, pour les fonctions de troisième espèce, correspond à la formule $F(\varphi)+F(\psi)-F(\mu)=0$ pour les fonctions de la première espèce, et à la formule $E(\varphi)+E(\psi)-E(\mu)=c^{2}\sin\mu\sin\varphi\sin\psi$ pour les fonctions de la seconde espèce.

Ainsi la différence qui est zéro dans les fonctions de première espèce, et algébrique dans celles de la seconde espèce, est exprimée par un arc de cercle ou par un logarithme dans les fonctions de troisième espèce. Je dis \emph{arc de cercle} ou \emph{logarithme}, car on voit que le second membre de l'équation $(h')$ sera un arc de cercle ou un logarithme, selon que $\alpha$ sera positif ou négatif, c'est-à-dire, selon que la fonction $\Pi$ se rapportera aux deux premiers cas ou au troisième de l'article 50.

\subsection*{(54).}

De l'équation $(h')$ et de toutes celles qu'on peut former semblablement entre trois fonctions $\Pi$, nous conclurons que si $i$, $k$, $l$, etc. sont des entiers positifs, on pourra toujours faire ensorte

\origpage{77}
que l'on ait
\[ i\Pi(\varphi)+k\Pi(\psi)+l\Pi(\omega)+\text{etc.}=W, \]
$W$ étant une quantité déterminable par arcs de cercle ou par logarithmes. Il faut pour cela établir entre les amplitudes $\varphi$, $\psi$, $\omega$, etc., la relation qui donne
\[ iF(\varphi)+kF(\psi)+lF(\omega)+\text{etc.}=0, \]
et cette relation peut toujours être exprimée par une équation algébrique entre les sinus des angles $\varphi$, $\psi$, $\omega$, etc. Ce résultat ne souffre aucune exception, et aurait lieu même quand le paramètre $n$ serait imaginaire. Mais on peut considérer ces propriétés sous un point de vue encore plus général.

\subsection*{(55).}

Soit $P$ ou $P(\varphi)$ une fonction rationnelle paire de $\sin\varphi$, et soit
\[ Z(\varphi)=\int\frac{P(\varphi)\cdot d\varphi}{\sqrt{(1-c^{2}\sin^{2}\varphi)}}. \]
Soit $Z(\psi)$ une fonction semblable de $\psi$, ces fonctions ou intégrales étant prises de manière qu'elles s'évanouissent lorsque les amplitudes $\varphi$ et $\psi$ sont nulles. Supposons qu'on ait toujours l'équation $F(\varphi)+F(\psi)-F(\mu)=0$, laquelle en prenant $\mu$ constante donne $\frac{d\varphi}{\Delta(\varphi)}+\frac{d\psi}{\Delta(\psi)}=0$, on aura donc
\[ Z(\varphi)+Z(\psi)-Z(\mu)=\int\frac{d\varphi}{\Delta(\varphi)}[P(\varphi)-P(\psi)]. \]
Faisons comme ci-dessus $\sin^{2}\varphi+\sin^{2}\psi=p$, $\sin\varphi\sin\psi=q$, il en résultera
\[ \begin{aligned} \sin^{2}\varphi&=\tfrac{1}{2}p+\tfrac{1}{2}\sqrt{(p^{2}-4q^{2})}\\ \sin^{2}\psi&=\tfrac{1}{2}p-\tfrac{1}{2}\sqrt{(p^{2}-4q^{2})}. \end{aligned} \]
Substituons maintenant la valeur de $\sin^{2}\varphi$ dans $P(\varphi)$, le résultat sera de la forme
\[ P(\varphi)=M+N\sqrt{(p^{2}-4q^{2})}, \]
$M$ et $N$ étant des fonctions rationnelles de $p$ et $q$. On aura de même
\[ P(\psi)=M-N\sqrt{(p^{2}-4q^{2})}; \]
donc $P(\varphi)-P(\psi)=2N\sqrt{(p^{2}-4q^{2})}=2N(\sin^{2}\varphi-\sin^{2}\psi)$, et

\origpage{78}
\[ \int\frac{d\varphi}{\Delta(\varphi)}[P(\varphi)-P(\psi)]=\int\frac{2N\,d\varphi(\sin^{2}\varphi-\sin^{2}\psi)}{\Delta(\varphi)}=-2\int N\,dq\sin\mu; \]
donc enfin
\[ Z(\varphi)+Z(\psi)-Z(\mu)=-2\sin\mu\int N\,dq. \]
Dans cette formule, $N$ est une fonction rationnelle de $p$ et de $q$; si on y substitue la valeur de $p$ en $q$, savoir $p=\sin^{2}\mu-2q\cos\mu\Delta(\mu)+c^{2}q^{2}\sin^{2}\mu$, $N$ sera une fonction rationnelle de $q$ seule, et ainsi la valeur de $Z(\varphi)+Z(\psi)-Z(\mu)$ pourra toujours se déterminer par arcs de cercle et par logarithmes.

En général si $i$, $k$, $l$, etc. désignent des nombres entiers positifs ou négatifs, et qu'on établisse entre les angles $\varphi$, $\psi$, $\omega$, etc. la relation qui donne $iF(\varphi)+kF(\psi)+lF(\omega)+\text{etc.}=0$, on aura en même temps
\[ iZ(\varphi)+kZ(\psi)+lZ(\omega)+\text{etc.}=W, \]
$W$ étant une quantité déterminable par arcs de cercle et par logarithmes. La même propriété aura lieu quand même $P$ contiendrait des puissances impaires de $\sin\varphi$; car la partie de $dZ$ affectée des puissances impaires, s'intégrerait par arcs de cercle et par logarithmes.

Il en est absolument de même de la fonction
\[ Z(x)=\int\frac{P\,dx}{\sqrt{(\alpha+\beta x+\gamma x^{2}+\delta x^{3}+\varepsilon x^{4})}}, \]
$P$ étant une fonction rationnelle de $x$, et on pourra toujours trouver une équation algébrique entre $x$, $y$, $z$, etc., telle que la quantité
\[ iZ(x)+kZ(y)+lZ(z)+\text{etc.} \]
soit déterminable par les arcs de cercle et les logarithmes.

\subsection*{(56).}

Revenons à la formule $(h')$ d'où l'on doit déduire tout ce qui concerne la comparaison des fonctions de troisième espèce. Nous avons déjà observé que dans cette formule l'arc de cercle doit être remplacé par un logarithme, lorsque le paramètre est de la forme $n=-c^{2}\sin^{2}\theta$; alors on a $\sqrt{\alpha}=\frac{\cos\theta}{\sin\theta}\Delta(\theta)\cdot\sqrt{-1}$, et parce qu'en général $\frac{1}{\sqrt{-1}}\operatorname{arc}\operatorname{tang}z\sqrt{-1}=\frac{1}{2}\log\frac{1+z}{1-z}$, si l'on fait

\origpage{79}
pour abréger,
\[ \begin{aligned} \frac{\sin\theta\cos\mu\Delta(\mu)-\cos\theta\sin\mu\Delta(\theta)}{1-c^{2}\sin^{2}\theta\sin^{2}\mu}&=\sin\mu',\\ \frac{\sin\theta\cos\mu\Delta(\mu)+\cos\theta\sin\mu\Delta(\theta)}{1-c^{2}\sin^{2}\theta\sin^{2}\mu}&=\sin\mu'', \end{aligned} \]
la formule $(h')$ se changera en celle-ci:
\[ \Pi(\varphi)+\Pi(\psi)-\Pi(\mu)=\frac{\operatorname{tang}\theta}{2\Delta(\theta)}\log\left(\frac{1+c^{2}\sin\theta\sin\mu'\sin\varphi\sin\psi}{1+c^{2}\sin\theta\sin\mu''\sin\varphi\sin\psi}\right), \]
et l'on peut remarquer que les angles auxiliaires $\mu'$, $\mu''$ sont ceux qui donneraient
\[ F(\theta)-F(\mu)=F(\mu'),\quad F(\theta)+F(\mu)=F(\mu''). \]
Au reste, comme on peut immédiatement convertir en logarithmes les arcs dont les tangentes sont de la forme $t\sqrt{-1}$, on pourra se borner à n'employer que l'équation $(h')$ dans toutes les comparaisons qu'on aura à faire des fonctions $\Pi$, selon les diverses valeurs de $n$; mais pour faciliter les applications, nous croyons devoir rapporter ici les formules qui ont lieu dans quelques-uns des cas les plus simples, en y joignant celles qui concernent les fonctions de la première et de la deuxième espèce.

\subsection*{(57).}

Soit, 1°. $\mu=\frac{1}{2}\pi$, on aura successivement pour les trois espèces de fonctions elliptiques:
\[ \begin{aligned} F(\varphi)+F(\psi)-F^{1}&=0\\ E(\varphi)+E(\psi)-E^{1}&=c^{2}\sin\varphi\sin\psi\\ \Pi(\varphi)+\Pi(\psi)-\Pi^{1}&=\frac{1}{\sqrt{\alpha}}\operatorname{arc}\operatorname{tang}\frac{n\sqrt{\alpha}\sin\varphi\sin\psi}{1+n}. \end{aligned} \]
La relation entre $\varphi$ et $\psi$ est donnée par l'équation $b\operatorname{tang}\varphi\operatorname{tang}\psi=1$, d'où l'on tire $\sin\psi=\frac{\cos\varphi}{\Delta(\varphi)}$ et $\sin\varphi=\frac{\cos\psi}{\Delta(\psi)}$. Quant à la valeur de $\alpha$, elle s'exprime suivant les différentes formes de $n$, comme on l'a vu (art.\ 50).

Si en particulier on a $\varphi=\psi$, ce qui donne $\operatorname{tang}\varphi=\frac{1}{\sqrt{b}}$, $\sin\varphi=\frac{1}{\sqrt{(1+b)}}$, les formules précédentes deviennent

\origpage{80}
\[ \begin{aligned} F(\varphi)&=\tfrac{1}{2}F^{1}\\ E(\varphi)&=\tfrac{1}{2}E^{1}+\tfrac{1}{2}(1-b)\\ \Pi(\varphi)&=\tfrac{1}{2}\Pi^{1}+\frac{1}{2\sqrt{\alpha}}\operatorname{arc}\operatorname{tang}\frac{n\sqrt{\alpha}}{(1+n)(1+b)}; \end{aligned} \]
elles servent ainsi à la bisection de la fonction complète.

Soit, 2°. $\psi=\varphi$ et $\mu=\varphi_{2}$, l'amplitude $\varphi_{2}$ se déduira de $\varphi$ par les formules de l'article 21, et on aura pour la duplication des fonctions, les formules
\[ \begin{aligned} 2F(\varphi)-F(\varphi_{2})&=0\\ 2E(\varphi)-E(\varphi_{2})&=c^{2}\sin^{2}\varphi\sin\varphi_{2}\\ 2\Pi(\varphi)-\Pi(\varphi_{2})&=\frac{1}{\sqrt{\alpha}}\operatorname{arc}\operatorname{tang}\left(\frac{n\sqrt{\alpha}\sin^{2}\varphi\sin\varphi_{2}}{1+n-n\cos^{2}\varphi\cos\varphi_{2}}\right). \end{aligned} \]
Les mêmes formules serviront à la bisection, en déterminant $\varphi$ par le moyen de $\varphi_{2}$.

Soit, 3°. $\psi=\varphi_{2}$ et $\mu=\varphi_{3}$, $\varphi_{3}$ étant l'amplitude qui donne $F(\varphi_{3})=3F(\varphi)$, et qu'on détermine par les formules de l'article 22, on aura pour la triplication des fonctions, les formules
\[ \begin{aligned} 3F(\varphi)-F(\varphi_{3})&=0\\ 3E(\varphi)-E(\varphi_{3})&=c^{2}\sin\varphi\sin\varphi_{2}(\sin\varphi+\sin\varphi_{3})\\ 3\Pi(\varphi)-\Pi(\varphi_{3})&=\frac{1}{\sqrt{\alpha}}\operatorname{arc}\operatorname{tang}\left(\frac{n\sqrt{\alpha}\sin\varphi\sin\varphi_{2}\sin\varphi_{3}}{1+n-n\cos\varphi\cos\varphi_{2}\cos\varphi_{3}}\right)\\ &\quad+\frac{1}{\sqrt{\alpha}}\operatorname{arc}\operatorname{tang}\left(\frac{n\sqrt{\alpha}\sin^{2}\varphi\sin\varphi_{2}}{1+n-n\cos^{2}\varphi\cos\varphi_{2}}\right). \end{aligned} \]
Dans le cas où $\varphi_{3}=\frac{1}{2}\pi$, on a pour la trisection des fonctions complètes, les formules
\[ \begin{aligned} 3F(\varphi)&=F^{1}\\ 3E(\varphi)&=E^{1}+c^{2}\sin\varphi\sin\varphi_{2}(1+\sin\varphi)\\ 3\Pi(\varphi)&=\Pi^{1}+\frac{1}{\sqrt{\alpha}}\operatorname{arc}\operatorname{tang}\left(\frac{n\sqrt{\alpha}\sin\varphi\sin\varphi_{2}}{1+n}\right)\\ &\quad+\frac{1}{\sqrt{\alpha}}\operatorname{arc}\operatorname{tang}\left(\frac{n\sqrt{\alpha}\sin\varphi\sin\varphi_{2}}{1+n-n\cos^{2}\varphi\cos\varphi_{2}}\right). \end{aligned} \]
\ednote{In the last line the numerator is printed $n\sqrt{\alpha}\sin\varphi\sin\varphi_{2}$; the second term of the triplication formula above, of which this is the case $\varphi_{3}=\frac{1}{2}\pi$, gives $n\sqrt{\alpha}\sin^{2}\varphi\sin\varphi_{2}$. An apparent misprint.}
Quant aux valeurs de $\varphi$ et $\varphi_{2}$, elles se trouveront par les articles 23 et 24.

\origpage{81}
\section*{Formation d'une suite infinie de fonctions elliptiques de la première espèce, liées entre elles par des rapports constans.}

\subsection*{(58).}

Jusqu'ici nous n'avons comparé entre elles les fonctions elliptiques de la première espèce, qu'autant qu'elles avaient le même module, ou qu'elles pouvaient être considérées comme représentant différens arcs d'une même courbe; ces comparaisons ont ensuite été étendues, d'après le même principe, aux fonctions de la seconde et de la troisième espèce; et les théorèmes contenus dans les formules $(f')$ et $(g')$ supposent encore que le module est le même dans les deux fonctions comparées.

Nous allons faire voir qu'on peut, par une loi très-simple, former une infinité de fonctions elliptiques de première espèce, qui diffèrent les unes des autres tant par le module que par l'amplitude, mais qui ont la propriété fort remarquable d'être entre elles dans des rapports constans.

Considérons les deux fonctions $F(c,\varphi)=\displaystyle\int\frac{d\varphi}{\sqrt{(1-c^{2}\sin^{2}\varphi)}}$, $F(c',\varphi')=\displaystyle\int\frac{d\varphi'}{\sqrt{(1-c'^{2}\sin^{2}\varphi')}}$; je dis que si l'on fait $c'=\frac{2\sqrt{c}}{1+c}$, et que l'on détermine $\varphi'$ d'après l'équation $\sin(2\varphi'-\varphi)=c\sin\varphi$, on aura généralement
\[ F(c',\varphi')=\frac{1+c}{2}F(c,\varphi). \]
En effet l'équation supposée donne d'abord $\cos(2\varphi'-\varphi)=\Delta$; ainsi on aura successivement:
\[ \begin{aligned} \cos2\varphi'&=\Delta\cos\varphi-c\sin^{2}\varphi\\ 2\cos^{2}\varphi'&=1-c\sin^{2}\varphi+\Delta\cos\varphi\\ 2\sin^{2}\varphi'&=1+c\sin^{2}\varphi-\Delta\cos\varphi\\ 2\sin\varphi'\cos\varphi'&=\sin\varphi(c\cos\varphi+\Delta)\\ 2d\varphi'&=\frac{d\varphi}{\Delta}(c\cos\varphi+\Delta)\\ \sqrt{(1-c'^{2}\sin^{2}\varphi')}&=\frac{c\cos\varphi+\Delta}{1+c}. \end{aligned} \]

\origpage{82}
Des deux dernières on tire $\frac{d\varphi'}{\sqrt{(1-c'^{2}\sin^{2}\varphi')}}=\frac{1+c}{2}\cdot\frac{d\varphi}{\Delta}$, ce qui donne en intégrant, $F(c',\varphi')=\frac{1+c}{2}F(c,\varphi)$. Nous n'ajoutons point de constante, parce que les amplitudes $\varphi$ et $\varphi'$ s'évanouissent en même temps. On voit donc par ce résultat que les fonctions $F(c',\varphi')$, $F(c,\varphi)$ seront entre elles dans un rapport constant, quelles que soient les amplitudes $\varphi'$ et $\varphi$, pourvu qu'elles soient liées entre elles par l'équation $\sin(2\varphi'-\varphi)=c\sin\varphi$.

Observons que comme l'arc $\varphi$ croît indéfiniment, et peut être de tant de circonférences qu'on voudra, l'arc $\varphi'$ croît aussi indéfiniment, mais de manière que $2\varphi'-\varphi$ est toujours renfermé entre les limites $+\theta$ et $-\theta$, $\theta$ étant l'arc dont $c$ est le sinus. En effet, puisqu'on a $\sin(2\varphi'-\varphi)=c\sin\varphi$ et $\cos(2\varphi'-\varphi)=\Delta$, $\Delta$ étant toujours positif, on voit que $2\varphi'-\varphi$ est toujours égal au plus petit arc $\delta$, positif ou négatif, déterminé par l'équation $\sin\delta=c\sin\varphi=\sin\theta\sin\varphi$; ainsi on a toujours $\varphi'=\frac{1}{2}\varphi+\frac{1}{2}\delta$ ou $\varphi=2\varphi'-\delta$. D'après cette observation, on n'aura jamais aucune ambiguité à craindre dans la détermination des valeurs respectives de $\varphi'$ et $\varphi$.

Si l'on fait $\varphi'=\frac{1}{2}\pi$, on aura $\varphi=\pi$ et $F(c,\varphi)=2F^{1}(c)$; ainsi les fonctions complètes $F^{1}(c')$, $F^{1}(c)$ ont entre elles cette relation
\[ F^{1}(c')=(1+c)F^{1}(c). \]

\subsection*{(59).}

Concevons maintenant qu'à partir du terme donné $c$, on forme une suite infinie de modules $c$, $c'$, $c''$, $c'''$, etc., d'après la loi,
\[ c'=\frac{2\sqrt{c}}{1+c},\quad c''=\frac{2\sqrt{c'}}{1+c'},\quad c'''=\frac{2\sqrt{c''}}{1+c''},\quad\text{etc.} \]
Cette suite de modules qui est continuellement croissante, aura pour limite l'unité, et atteindra sensiblement cette limite au bout d'un assez petit nombre de termes.

Si on appelle par analogie $b'$, $b''$, $b'''$, etc. les complémens des modules $c'$, $c''$, $c'''$, etc., la suite $b'$, $b''$, $b'''$, etc. sera continuellement décroissante, et chaque terme se déduira du module précédent suivant cette loi:
\[ b'=\frac{1-c}{1+c},\quad b''=\frac{1-c'}{1+c'},\quad b'''=\frac{1-c''}{1+c''},\quad\text{etc.} \]

\origpage{83}

Soit ensuite $\varphi$, $\varphi'$, $\varphi''$, $\varphi'''$, etc. la série des amplitudes qui se déduisent chacune de la précédente par les formules
\[ \begin{gathered} \sin(2\varphi'-\varphi)=c\sin\varphi\\ \sin(2\varphi''-\varphi')=c'\sin\varphi'\\ \sin(2\varphi'''-\varphi'')=c''\sin\varphi''\\ \text{etc.} \end{gathered} \]
on formera de cette manière une suite infinie de fonctions de première espèce $F(c,\varphi)$, $F(c',\varphi')$, $F(c'',\varphi'')$, etc., entre lesquelles on aura les équations
\[ \begin{aligned} &F(c',\varphi')=\frac{1+c}{2}F(c,\varphi)\\ &F(c'',\varphi'')=\frac{1+c'}{2}F(c',\varphi')=\frac{1+c}{2}\cdot\frac{1+c'}{2}F(c,\varphi)\\ &F(c''',\varphi''')=\frac{1+c''}{2}F(c'',\varphi'')=\frac{1+c}{2}\cdot\frac{1+c'}{2}\cdot\frac{1+c''}{2}F(c,\varphi),\\ &\qquad\text{etc.} \end{aligned} \]
d'où il suit que deux quelconques de ces fonctions sont toujours entre elles dans un rapport constant pour toutes les valeurs des amplitudes correspondantes.

Quant aux fonctions complètes, leurs rapports seront également constans, et l'équation $F^{1}(c')=(1+c)F^{1}(c)$, déjà trouvée, donnera successivement
\[ \begin{aligned} &F^{1}(c')=(1+c)F^{1}(c)\\ &F^{1}(c'')=(1+c')F^{1}(c')=(1+c)(1+c')F^{1}(c)\\ &F^{1}(c''')=(1+c'')F^{1}(c'')=(1+c)(1+c')(1+c'')F^{1}(c)\\ &\qquad\text{etc.} \end{aligned} \]

\subsection*{(60).}

On peut encore donner à ces résultats une plus grande extension. En effet, la suite infinie de modules $c$, $c'$, $c''$, etc., qui est croissante dans un sens, et qui a pour limite l'unité, peut être prolongée à l'infini dans le sens contraire où elle sera décroissante et aura pour limite zéro. Désignons par $c$, $c^{0}$, $c^{00}$, $c^{000}$, etc. cette suite décroissante; la loi qui lie deux termes consécutifs sera semblablement
\[ c=\frac{2\sqrt{c^{0}}}{1+c^{0}},\quad c^{0}=\frac{2\sqrt{c^{00}}}{1+c^{00}},\quad c^{00}=\frac{2\sqrt{c^{000}}}{1+c^{000}},\quad\text{etc.} \]

\origpage{84}
Mais pour former chaque terme au moyen du précédent, il sera plus simple de se servir des valeurs suivantes, où $b^{0}$, $b^{00}$, $b^{000}$, etc. désignent les complémens des modules $c^{0}$, $c^{00}$, $c^{000}$, etc.,
\[ c^{0}=\frac{1-b}{1+b},\quad c^{00}=\frac{1-b^{0}}{1+b^{0}},\quad c^{000}=\frac{1-b^{00}}{1+b^{00}},\quad\text{etc.} \]
Cela posé, si on désigne par $\varphi$, $\varphi^{0}$, $\varphi^{00}$, $\varphi^{000}$, etc. la série des amplitudes correspondantes aux modules $c$, $c^{0}$, $c^{00}$, $c^{000}$, etc., on aura d'abord $\sin(2\varphi-\varphi^{0})=c^{0}\sin\varphi^{0}$, équation qu'il faut résoudre pour déduire $\varphi^{0}$ de $\varphi$; on en tire successivement
\[ \begin{aligned} \sin\varphi^{0}&=(1+b)\frac{\sin\varphi\cos\varphi}{\Delta}\\ \cos\varphi^{0}&=\frac{1-(1+b)\sin^{2}\varphi}{\Delta}\\ \Delta(c^{0},\varphi^{0})&=\frac{1-(1-b)\sin^{2}\varphi}{\Delta}\\ \operatorname{tang}\varphi^{0}&=\frac{(1+b)\operatorname{tang}\varphi}{1-b\operatorname{tang}^{2}\varphi}. \end{aligned} \]
Cette dernière peut se mettre sous la forme très-simple $\operatorname{tang}(\varphi^{0}-\varphi)=b\operatorname{tang}\varphi$, et on en déduirait sans ambiguité,
\[ \varphi^{0}=2\varphi-c^{0}\sin2\varphi+\tfrac{1}{2}c^{02}\sin4\varphi-\tfrac{1}{3}c^{03}\sin6\varphi+\text{etc.}, \]
ce qui s'accorde avec la remarque que nous avons faite sur la valeur toujours limitée de $2\varphi'-\varphi$, qui s'applique à $2\varphi-\varphi^{0}$, $2\varphi^{0}-\varphi^{00}$, etc.

Supposons donc qu'après avoir déterminé les modules décroissans $c^{0}$, $c^{00}$, $c^{000}$, etc., ainsi que leurs complémens $b^{0}$, $b^{00}$, $b^{000}$, etc. comme il vient d'être dit, on calcule successivement les amplitudes $\varphi^{0}$, $\varphi^{00}$, $\varphi^{000}$, etc. par les formules
\[ \begin{gathered} \operatorname{tang}(\varphi^{0}-\varphi)=b\operatorname{tang}\varphi\\ \operatorname{tang}(\varphi^{00}-\varphi^{0})=b^{0}\operatorname{tang}\varphi^{0}\\ \operatorname{tang}(\varphi^{000}-\varphi^{00})=b^{00}\operatorname{tang}\varphi^{00}\\ \text{etc.} \end{gathered} \]
Alors on aura une suite infinie de fonctions $F(c,\varphi)$, $F(c^{0},\varphi^{0})$, $F(c^{00},\varphi^{000})$\ednote{Printed $\varphi^{000}$; the series and the next page give $F(c^{00},\varphi^{00})$, so probably a misprint for $\varphi^{00}$.}, etc. liées entre elles par des rapports constans, en cette sorte:

\origpage{85}
\[ \begin{gathered} F(c,\varphi)=\frac{1+c^{0}}{2}F(c^{0},\varphi^{0})\\ F(c^{0},\varphi^{0})=\frac{1+c^{00}}{2}F(c^{00},\varphi^{00})\\ F(c^{00},\varphi^{00})=\frac{1+c^{000}}{2}F(c^{000},\varphi^{000})\\ \text{etc.} \end{gathered} \]
Lorsqu'on fait $\varphi=\frac{1}{2}\pi$, on a $\varphi^{0}=\pi$, $\varphi^{00}=2\pi$, $\varphi^{000}=4\pi$, etc.; d'où il suit que les fonctions complètes ont entre elles ces relations
\[ \begin{gathered} F^{1}(c)=(1+c^{0})F^{1}(c^{0})\\ F^{1}(c^{0})=(1+c^{00})F^{1}(c^{00})\\ F^{1}(c^{00})=(1+c^{000})F^{1}(c^{000})\\ \text{etc.} \end{gathered} \]

\section*{Application de la même loi aux fonctions elliptiques de la seconde espèce.}

\subsection*{(61).}

Considérons présentement les deux fonctions $E(c,\varphi)=\displaystyle\int d\varphi\,\Delta(c,\varphi)$, $E(c',\varphi')=\displaystyle\int d\varphi'\,\Delta(c',\varphi')$; si dans la seconde formule on substitue pour $d\varphi'$ et $\Delta(c',\varphi')$ les valeurs trouvées art.~58, on aura
\[ \begin{aligned} (2+c)\,E(c',\varphi')&=\int\frac{d\varphi}{\Delta}(c\cos\varphi+\Delta)^{2}\\ &=2E(c,\varphi)-b^{2}F(c,\varphi)+2c\sin\varphi; \end{aligned} \]
\ednote{The factor reads $(2+c)$, its right part covered by a reader's ink stroke; the display below and the computation require $(2+2c)$. Whether a second $2$ is hidden under the mark or the print omits it cannot be settled from the scan.}
d'où l'on voit que la fonction de première espèce $F(c,\varphi)$ peut s'exprimer par les deux arcs d'ellipse $E(c,\varphi)$, $E(c',\varphi')$, puisqu'on a
\[ b^{2}F(c,\varphi)=2E(c,\varphi)-(2+2c)E(c',\varphi')+2c\sin\varphi. \]

\subsection*{(62).}

On a trouvé (art.~13) que l'expression de l'arc d'hyperbole est
\[ \Upsilon=\Delta\operatorname{tang}\varphi-E(c,\varphi)+b^{2}F(c,\varphi); \]
si on y substitue la valeur de $b^{2}F(c,\varphi)$, cette expression deviendra
\[ \Upsilon=\Delta\operatorname{tang}\varphi+E(c,\varphi)-2(1+c)E(c',\varphi')+2c\sin\varphi; \]
d'où il suit qu'un arc d'hyperbole peut toujours s'exprimer par deux

\origpage{86}
arcs d'ellipse, ce qui est le beau théorème dont Landen a enrichi la Géométrie${}^{(1)}$.

Nous avons déjà appelé $G(\varphi)$ la différence entre l'arc d'hyperbole $\Upsilon$ et sa tangente $\Delta\operatorname{tang}\varphi$, terminée à la perpendiculaire abaissée du centre; cette différence s'exprime donc en général par deux arcs d'ellipse, de sorte qu'on a
\[ G(\varphi)=2(1+c)E(c',\varphi')-E(c,\varphi)-2c\sin\varphi; \]
et en particulier lorsqu'on fait $\varphi=\frac{1}{2}\pi$, on a pour l'expression de la différence entre l'asymptote et la courbe,
\[ G^{1}=2(1+c)E(c',\varphi')-E^{1}(c)-2c. \]
Mais puisque $\varphi=\frac{1}{2}\pi$, on a $\sin(2\varphi'-\frac{1}{2}\pi)=c\sin\varphi=c$, ou $\cos2\varphi'=-c$, ce qui donne $\sin^{2}\varphi'=\dfrac{1+c}{2}=\dfrac{1}{1+b'}$; donc l'arc $E(c',\varphi')$ est celui qui se mesure par la moitié de $E^{1}(c')$, et on a $E'(c',\varphi')$\ednote{Printed $E'$ with a prime, where the same arc is $E(c',\varphi')$ in the line above; probably a misprint or a damaged glyph for $E$.}${}=\frac{1}{2}E^{1}(c')+\frac{1}{2}(1-b')$. Substituant cette valeur dans celle de $G^{1}$, il viendra
\[ G^{1}=(1+c)E^{1}(c')-E^{1}(c); \]
donc la transcendante $G^{1}$ est égale à la différence des deux quarts d'ellipse qui ont pour demi-axes, l'un $1+c$ et $1-c$, l'autre $1$ et $\sqrt{(1-c^{2})}$.

\subsection*{(63).}

La valeur de $F$ trouvée n°~61, est l'intégrale de forme particulière qui satisfait aux équations différentielles $(e')$ de l'article~43; nous allons en déduire de nouvelles propriétés des fonctions $E$.

Considérons une suite infinie de fonctions $E(c,\varphi)$, $E(c',\varphi')$, $E(c'',\varphi'')$, etc. formées d'après la même loi que les fonctions de première espèce $F(c,\varphi)$, $F(c',\varphi')$, $F(c'',\varphi'')$, etc., on aura d'abord les deux équations
\[ \begin{aligned} &\tfrac{1}{2}b^{2}F(c,\varphi)=E(c,\varphi)-(1+c)E(c',\varphi')+c\sin\varphi\\ &\tfrac{1}{2}b'^{2}F(c',\varphi')=E(c',\varphi')-(1+c')E(c'',\varphi'')+c'\sin\varphi'; \end{aligned} \]

\textbf{(1)} On déduirait aisément des mêmes formules que tout arc d'ellipse peut s'exprimer par deux arcs d'hyperbole.

\origpage{87}
mais on a de plus $F(c',\varphi')=\dfrac{1+c}{2}F(c,\varphi)$; éliminant donc de ces trois équations les fonctions de première espèce, on aura entre les trois fonctions consécutives de seconde espèce $E(c,\varphi)$, $E(c',\varphi')$, $E(c'',\varphi'')$, cette équation
\[ \begin{aligned} 0={}&\tfrac{1}{2}b'(1+b')E(c,\varphi)-(2+b')E(c',\varphi')+2(1+c')E(c'',\varphi'')\\ &+\tfrac{1}{2}b'(1-b')\sin\varphi-2c'\sin\varphi', \end{aligned} \]
équation qui peut être considérée comme une sorte d'intégrale particulière de l'équation différentielle du second ordre $0=(1+c^{2})\dfrac{ddE}{dc^{2}}+\text{etc.}$ donnée art.~45.

La même équation peut être appliquée à trois ellipses consécutives, prises non-seulement dans la suite infinie $E(c,\varphi)$, $E(c',\varphi')$, $E(c'',\varphi'')$, $E(c''',\varphi''')$, etc., mais en général dans la suite doublement infinie
\[ \ldots E(c'',\varphi''),\ E(c',\varphi'),\ E(c,\varphi),\ E(c^{0},\varphi^{0}),\ E(c^{00},\varphi^{00}),\ldots \]
dont les extrêmes sont, d'une part, l'ellipse qui a pour excentricité $1$ et qui se réduit à son grand axe; d'autre part, l'ellipse qui a pour excentricité $0$ et qui se confond avec le cercle: il en résulte donc que par la rectification indéfinie de deux ellipses de cette suite, on obtient la rectification indéfinie de toutes les autres.

La formule générale se simplifie lorsqu'il s'agit de la rectification définie de ces ellipses. En effet si on fait $\varphi''=\frac{1}{2}\pi$, on aura $\varphi'=\pi$ et $\varphi=2\pi$, ce qui donnera $E(c'',\varphi'')=E^{1}(c'')$, $E(c',\varphi')=2E^{1}(c')$, $E(c,\varphi)=4E^{1}(c)$; donc on aura entre les trois quarts d'ellipse $E^{1}(c)$, $E^{1}(c')$, $E^{1}(c'')$, cette équation
\[ 0=b'(1+b')E^{1}(c)-(2+b')E^{1}(c')+(1+c')E^{1}(c''). \]
On aura une équation semblable entre trois termes consécutifs quelconques pris dans la série générale des ellipses. Ainsi la circonférence d'une ellipse proposée pourra toujours se déterminer exactement par les circonférences de deux ellipses, aussi peu différentes de la ligne droite qu'on voudra, en prolongeant les séries dans un sens, ou aussi peu différentes du cercle qu'on voudra, en prolongeant les séries dans l'autre sens.

\origpage{88}

Dans ce dernier cas, il faudrait faire usage des équations successives
\[ \begin{aligned} &0=(1+c^{0})E^{1}(c)-(2+b^{0})E^{1}(c^{0})+b^{0}(1+b^{0})E^{1}(c^{00})\\ &0=(1+c^{00})E^{1}(c^{0})-(2+b^{00})E^{1}(c^{00})+b^{00}(1+b^{00})E^{1}(c^{000})\\ &\text{etc.}, \end{aligned} \]
qu'on prolongerait aussi loin qu'on aurait cru devoir prolonger la suite des modules $c$, $c^{0}$, $c^{00}$, etc.

La détermination exacte et absolue dont on vient de parler peut être regardée comme un théorème fort remarquable dans la théorie des transcendantes; mais si on a seulement pour but d'obtenir des approximations, on y parviendra plus facilement par la méthode que nous donnerons ci-après.

\subsection*{(64).}

Nous avons trouvé (art.~41) que les deux fonctions $E^{1}(c)$, $F^{1}(c)$, tant pour le module $c=\sin15^{\circ}$, que pour son complément $c=\sin75^{\circ}$, peuvent se déterminer par l'une des quatre fonctions supposée connue. Donc les deux séries d'ellipses formées, l'une d'après le module $c=\sin15^{\circ}=\sqrt{\left(\dfrac{2-\sqrt{3}}{4}\right)}$, l'autre d'après le module $c=\sin75^{\circ}=\sqrt{\left(\dfrac{2+\sqrt{3}}{4}\right)}$, sont telles, que connaissant la circonférence d'une seule de ces ellipses, on pourra trouver la circonférence de toutes les autres. Il en est de même des deux suites de fonctions de première espèce $F^{1}(c)$ formées d'après les mêmes modules, et un seul terme connu dans ces quatre séries, suffira pour faire connaître tous les autres.

Nous avons également trouvé (art.~42) que lorsque $c=\sin45^{\circ}=\sqrt{\tfrac{1}{2}}$, les fonctions $F^{1}(c)$, $E^{1}(c)$ peuvent se déterminer l'une par l'autre; mais $F^{1}(c)$ se détermine généralement par les deux quantités $E^{1}(c)$, $E^{1}(c')$, ou par les deux $E^{1}(c)$, $E^{1}(c^{0})$, car d'après nos formules on trouve aisément
\[ bF^{1}(c)=-E^{1}(c)+(1+b)E^{1}(c^{0}). \]
Donc en général toutes les ellipses qui composent la série formée d'après le module $c=\sin45^{\circ}$, sont telles, que la circonférence de l'une d'elles étant connue, on pourra déterminer celle de toutes les autres.

\origpage{89}

A ces trois cas généraux on peut en joindre un quatrième. Supposons qu'on ait $b=c^{0}=\dfrac{1-b}{1+b}$, ce qui donne $b=-1+\sqrt{2}$; et $c^{2}=-2+2\sqrt{2}$, on aura $F^{1}(b)=F^{1}(c^{0})=\dfrac{1}{1+c^{0}}F^{1}(c)=\dfrac{1}{\sqrt{2}}F^{1}(c)$, $E^{1}(b)=E^{1}(c^{0})$. Substituant ces valeurs dans la formule $(d')$, on aura
\[ \frac{\pi}{\sqrt{2}}=F^{1}(c)\,[E^{1}(c)+\sqrt{2}E^{1}(c^{0})-F^{1}(c)]. \]
Mais on a d'ailleurs $bF^{1}(c)=(1+b)E^{1}(c^{0})-E^{1}(c)$; ainsi on pourra déterminer $E^{1}(c)$ par $E^{1}(c^{0})$, et en général dans la suite d'ellipses, formée d'après le module $c=\sqrt{(2\sqrt{2}-2)}$, la circonférence d'une seule ellipse étant connue, on pourra déterminer celle de toutes les autres.

\section*{Méthode d'approximation appliquée aux fonctions elliptiques de la première espèce.}

\subsection*{(65).}

Etant proposée la fonction $F(c,\varphi)$ dont on veut avoir une valeur approchée, on calculera les modules décroissans $c^{0}$, $c^{00}$, $c^{000}$, etc. et les amplitudes croissantes $\varphi^{0}$, $\varphi^{00}$, $\varphi^{000}$, etc. par les formules de l'art.~60; on aura ainsi successivement
\[ \begin{aligned} F(c,\varphi)&=\frac{1+c^{0}}{2}F(c^{0},\varphi^{0})\\ &=\frac{1+c^{0}}{2}\cdot\frac{1+c^{00}}{2}F(c^{00},\varphi^{00})\\ &=\frac{1+c^{0}}{2}\cdot\frac{1+c^{00}}{2}\cdot\frac{1+c^{000}}{2}F(c^{000},\varphi^{000})\\ &\quad\text{etc.} \end{aligned} \]
Mais lorsque $c$ est devenu très-petit, on a $\Delta=1$ et $\displaystyle\int\frac{d\varphi}{\Delta}=\varphi$. Soit donc $\Phi$ la limite des angles $\frac{1}{2}\varphi^{0}$, $\frac{1}{4}\varphi^{00}$, $\frac{1}{8}\varphi^{000}$, etc., limite qu'ils atteindront toujours sensiblement au bout d'un certain nombre de termes, et on aura la fonction demandée
\[ F(c,\varphi)=\Phi(1+c^{0})(1+c^{00})(1+c^{000}),\ \text{etc.} \]
Lorsque $\varphi=\frac{1}{2}\pi$, la limite $\Phi$ sera pareillement $\frac{1}{2}\pi$, de sorte

\origpage{90}
qu'on aura
\[ F^{1}(c)=\frac{\pi}{2}(1+c^{0})(1+c^{00})(1+c^{000}),\ \text{etc.} \]

Le produit constant $(1+c^{0})(1+c^{00})(1+c^{000})$, etc. que nous représenterons par $K$, peut aussi s'exprimer de cette manière:
\[ K=\frac{2\sqrt{c^{0}}}{c}\cdot\frac{2\sqrt{c^{00}}}{c^{0}}\cdot\frac{2\sqrt{c^{000}}}{c^{00}},\ \text{etc.}; \]
et sous cette forme il est aisé à calculer par logarithmes. $K$ étant connu, on aura $F(c,\varphi)=K\Phi$ et $F^{1}(c)=K\cdot\dfrac{\pi}{2}$.

Pour faciliter le calcul des modules décroissans, on déterminera un angle auxiliaire $\mu$ par l'équation $\sin\mu=c$, ce qui donnera $b=\cos\mu$ et $c^{0}=\operatorname{tang}^{2}\frac{1}{2}\mu$. Faisant de même $c^{0}=\operatorname{tang}^{2}\frac{1}{2}\mu=\sin\mu^{0}$, ce qui déterminera un nouvel angle $\mu^{0}$, on aura $c^{00}=\operatorname{tang}^{2}\frac{1}{2}\mu^{00}$\ednote{Printed $\mu^{00}$; by the rule just stated the angle is $\mu^{0}$, so probably a misprint.}, et ainsi de suite.

Lorsqu'on sera parvenu à un $c$ fort petit, on pourra, pour éviter les angles trop petits, calculer le terme suivant $c^{0}$ par la formule
\[ c^{0}=\frac{1}{4}c^{2}+\frac{1\cdot3}{4\cdot6}c^{4}+\frac{1\cdot3\cdot5}{4\cdot6\cdot8}c^{6}+\text{etc.}, \]
dont le premier ou tout au plus les deux premiers termes suffiront.

Quant au calcul des angles $\varphi^{0}$, $\varphi^{00}$, etc., nous n'avons rien à ajouter à la simplicité de la formule $\operatorname{tang}(\varphi^{0}-\varphi)=b\operatorname{tang}\varphi$, qui est très-propre au calcul trigonométrique. Nous rappellerons seulement, ce qui a été dit art.~60, que l'angle $\varphi^{0}-\varphi$ est presque égal à $\varphi$ lorsque $c$ est très-petit, et qu'en général sa différence avec $\varphi$ est moindre que l'angle qui a pour sinus $c$. Il faut donc prendre pour l'angle $\varphi^{0}-\varphi$, non pas toujours le plus petit angle que donnent les tables des sinus, mais celui qui approche beaucoup de $\varphi$, et qui peut être de plusieurs circonférences.

\section*{EXEMPLE.}

\subsection*{(66).}

On demande la valeur de la fonction $F(c,\varphi)$, lorsque $c=\frac{1}{2}\sqrt{(2+\sqrt{3})}=\sin75^{\circ}$ et $\operatorname{tang}\varphi=\sqrt{\left(\dfrac{2}{\sqrt{3}}\right)}$.

\origpage{91}

Voici d'abord un tableau qui offre le calcul des modules et de leurs complémens:
\[ \begin{aligned} &\textit{Valeurs des c et b.} && \textit{Leurs logarithmes.}\\ c&=\sin\ 75^{\circ}\ 0'\ 0''00\ldots\ldots && 9.9849438\\ b&=\cos\ 75.\ 0.\ 0.00\ldots\ldots && 9.4129962\\ c^{0}&=\left\{\begin{array}{l}\operatorname{tang}^{2}37.30.\ 0.00\\ \sin\ 36.\ 4.16.47\end{array}\right\}\ldots\ldots && 9.7699610\\ b^{0}&=\cos\ 36.\ 4.16.47\ldots\ldots && 9.9075648\\ c^{00}&=\left\{\begin{array}{l}\operatorname{tang}^{2}18.\ 2.\ 8.235\\ \sin\ 6.\ 5.\ 9.38\end{array}\right\}\ldots\ldots && 9.0253880\\ b^{00}&=\cos\ 6.\ 5.\ 9.38\ldots\ldots && 9.9975452\\ c^{000}&=\left\{\begin{array}{l}\operatorname{tang}^{2}\ 3.\ 2.34.69\\ \sin\ 0.\ 9.42.90\end{array}\right\}\ldots\ldots && 7.4511672\\ b^{000}&=\cos\ 0.\ 9.42.90\ldots\ldots && 9.9999982\\ c^{0000}&=\tfrac{1}{4}(c^{000})^{2}+\text{etc.}\ldots\ldots && 4.3002761\\ b^{0000}&\ldots\ldots && 0.0000000 \end{aligned} \]
Ces logarithmes donneront aisément la valeur de $K$, pour laquelle il suffira d'employer quatre facteurs. On aura ainsi $\log K=0.2460561$ et $K=1.7622037$. De là résulte d'abord $F^{1}=\dfrac{\pi}{2}K=2.7680631$; on trouvera ensuite par le calcul des amplitudes,
\[ \begin{aligned} \varphi&=47^{\circ}\ 3'\ 30''94\\ \varphi^{0}&=62.36.\ 3.10\\ \varphi^{00}&=119.55.47.67\\ \varphi^{000}&=240.\ 0.\ 0.19\\ \varphi^{0000}&=480.\ 0.\ 0.00. \end{aligned} \]
Les autres valeurs de $\varphi$ augmenteraient en raison double; d'où il suit que la limite des angles $\varphi$, $\dfrac{\varphi^{0}}{2}$, $\dfrac{\varphi^{00}}{4}$, etc. est $30^{\circ}$ ou $\dfrac{\pi}{6}$; donc on a $F(c,\varphi)=K\dfrac{\pi}{6}=0.9226877$, ce qui s'accorde avec la valeur trouvée art.~28.

Remarquons que puisque la limite $\Phi=30^{\circ}=\dfrac{1}{3}\cdot\dfrac{\pi}{2}$, on a $F=\frac{1}{3}F^{1}$; et en effet cette égalité a rigoureusement lieu d'après la valeur que nous avons prise pour $\operatorname{tang}\varphi$. Voyez l'article~24.

\origpage{92}
\subsection*{(67).}

Etant donnée la valeur de l'amplitude $\varphi$, on voit qu'il est facile de trouver la fonction $F$ avec toute l'exactitude nécessaire; réciproquement il peut être utile de déterminer l'amplitude en supposant connue la valeur de la fonction $F$. Pour cet effet, il faudra calculer les termes de la série $c^{0}$, $c^{00}$, etc. jusqu'à un terme assez petit pour être négligé. Soit, par exemple, ce terme $c^{00000}$, ou pour abréger $c^{05}$; puisqu'on a $\varphi^{05}=2\varphi^{04}-c^{05}\sin^{2}\varphi^{04}+\text{etc.}$\ednote{Printed with a raised $2$ after sin; the series of art.\ 60 gives $\sin 2\varphi^{04}$, so the raised figure is probably a misplaced $2$.}, on aura d'une manière suffisamment exacte $\varphi^{05}=2\varphi^{04}$, et à plus forte raison $\varphi^{06}=2\varphi^{05}$, etc. Donc la limite des angles $\varphi$, $\frac{1}{2}\varphi^{0}$, $\frac{1}{4}\varphi^{00}$, etc. sera $\frac{1}{16}\varphi^{04}$: cette limite a été désignée par $\Phi$, ainsi on aura $\varphi^{04}=16\Phi$. D'ailleurs la valeur de $F$ étant donnée, on connaît $\Phi$ par l'équation $\Phi=\dfrac{F}{K}$; donc on connaîtra aussi $\varphi^{04}=16\Phi$. Cela posé, on calculera successivement les valeurs de $\varphi^{000}$, $\varphi^{00}$, $\varphi^{0}$, $\varphi$, au moyen des équations
\[ \begin{aligned} \sin(2\varphi^{03}-\varphi^{04})&=c^{04}\sin\varphi^{04}\\ \sin(2\varphi^{02}-\varphi^{03})&=c^{03}\sin\varphi^{03}\\ \sin(2\varphi^{0}-\varphi^{02})&=c^{02}\sin\varphi^{02}\\ \sin(2\varphi-\varphi^{0})&=c^{0}\sin\varphi^{0}, \end{aligned} \]
et on aura l'amplitude cherchée $\varphi$.

Cette méthode s'applique particulièrement à la résolution de l'équation $F(\psi)=nF(\varphi)$, quel que soit $n$. Etant donné $\varphi$, on connaîtra $F(\varphi)$ et $nF(\varphi)$, ou $F(\psi)$; ensuite de $F(\psi)$, on déduira l'amplitude $\psi$, comme on vient de l'expliquer. Ce moyen n'exigera jamais qu'un petit nombre d'opérations, à moins que $1-c$ ne soit d'une petitesse excessive; au lieu que si $n$ était un peu grand ou seulement fractionnaire, les méthodes algébriques que nous avons données pour déterminer $\varphi_{n}$ d'après l'équation $F(\varphi_{n})=nF(\varphi)$ deviendraient très-longues ou même impraticables.

\subsection*{(68).}

La méthode que nous venons d'exposer est en général très-expéditive; elle exigera seulement qu'on calcule quelques termes de plus, tant dans la suite des modules que dans celle des amplitudes, lorsque $c$ sera extrêmement près de l'unité; mais on peut s'assurer que ce nombre de termes ne sera jamais bien considérable, car dans l'hypothèse où on devrait continuer la suite $c$, $c^{0}$, $c^{00}$, etc. jusqu'au dixième terme pour que ce terme fût d'une unité

\origpage{93}
décimale du dixième ordre, il faudrait que la valeur primitive de $b$ fût plus petite que $10^{-21}$, et par conséquent que celle de $1-c$ fût plus petite que $10^{-42}$.

L'universalité de la méthode est suffisamment établie par cette observation; cependant lorsque $1-c$ est extrêmement petit, on peut profiter de cette circonstance pour simplifier les calculs, et procéder d'une autre manière aux approximations.

Dans le cas dont il s'agit, la quantité $b$ est très-petite, et comme on a $F=\displaystyle\int\frac{d\varphi}{\sqrt{(\cos^{2}\varphi+b^{2}\sin^{2}\varphi)}}$, si l'amplitude $\varphi$ est telle que $\operatorname{tang}\varphi$ soit beaucoup plus petite que $\dfrac{1}{b}$, on aura d'une manière suffisamment approchée
\[ F=\int\frac{d\varphi}{\cos\varphi}=\log\operatorname{tang}(45^{\circ}+\tfrac{1}{2}\varphi). \]

Si $\operatorname{tang}\varphi$ est plus grande que $\dfrac{1}{b}$, il faudra transformer la formule proposée $F(c,\varphi)$ en une autre où $b$ soit beaucoup plus petit, afin que dans la transformée $b\operatorname{tang}\varphi$ devienne une quantité négligeable. C'est ce qu'il est facile d'obtenir en calculant jusqu'au terme convenable les modules croissans $c'$, $c''$, etc. et les amplitudes $\varphi'$, $\varphi''$, etc., par les formules de l'article~59.

Soit pour cet effet $b=\sin\lambda$, on aura $b'=\operatorname{tang}^{2}\frac{1}{2}\lambda$; soit de nouveau $b'=\operatorname{tang}^{2}\frac{1}{2}\lambda=\sin\lambda'$, on aura $b''=\operatorname{tang}^{2}\frac{1}{2}\lambda'$, et ainsi de suite jusqu'à un terme $b^{\mu}$ qui ait le degré de petitesse exigé; on calculera ensuite les amplitudes $\varphi$, $\varphi'$, $\varphi''$, $\varphi'''$, etc., jusqu'à un terme $\varphi^{\mu}$ qui corresponde à la dernière valeur de $b$; et ce terme étant nommé $\Phi'$, si l'on fait
\[ K'=\frac{2}{1+c}\cdot\frac{2}{1+c'}\cdot\frac{2}{1+c''}\,\text{etc.}=\sqrt{\left(\frac{c'c''c'''\ldots\text{etc.}}{c}\right)}, \]
on aura
\[ F(c,\varphi)=K'\log\operatorname{tang}(45^{\circ}+\tfrac{1}{2}\Phi'). \]
Cette valeur sera exacte si on prolonge à l'infini les facteurs dont est composé $K'$, et la suite dont $\Phi''$\ednote{Printed $\Phi''$; the term was named $\Phi'$ just above, so probably a misprint for $\Phi'$.} est le dernier terme; mais au bout d'un assez petit nombre de termes, on obtiendra en général tout le degré d'approximation qu'on peut desirer. On aurait en même temps

\origpage{94}
pour la fonction complète $F^{1}(c)$, soit la valeur donnée par la formule précédente en calculant l'angle $\Phi'$ par le moyen de $\varphi=90^{\circ}$, soit la valeur $F^{1}(c)=K'\cdot\dfrac{1}{2^{\mu}}\log\dfrac{4}{b^{\mu}}$.

\section*{EXEMPLE.}

\subsection*{(69).}

Soit comme ci-dessus $c=\sin75^{\circ}$, $\operatorname{tang}\varphi=\sqrt{\left(\dfrac{2}{\sqrt{3}}\right)}$; ce cas est peu favorable à l'application de la méthode précédente, parce que $b$ n'est pas une quantité très-petite.

On calculera d'abord les modules croissans $c'$, $c''$, $c'''$, etc., et leurs complémens $b'$, $b''$, $b'''$, etc., comme il suit:
\[ \begin{aligned} &\textit{Valeurs des b et c.} && \textit{Leurs logarithmes.}\\ b&=\sin\ 15^{\circ}\ 0'\ 0''00\ldots\ldots && 9.4129962\\ c&=\cos\ 15.\ 0.\ 0.00\ldots\ldots && 9.9849438\\ b'&=\left\{\begin{array}{l}\operatorname{tang}^{2}\ 7.30.\ 0.00\\ \sin\ 0.59.35.24\end{array}\right\}\ldots\ldots && 8.2388582\\ c'&\ldots\ldots && 9.9999348\\ b''&=\left\{\begin{array}{l}\operatorname{tang}^{2}\ 0.29.47.62\\ \sin\ 0.\ 0.\text{etc.}\end{array}\right\}\ldots\ldots && 3.8757219\\ c''&\ldots\ldots && 0.0000000\\ b'''&=(\tfrac{1}{2}b'')^{2}\ldots\ldots && 1.1493838\\ c'''&\ldots\ldots && 0.0000000. \end{aligned} \]
\ednote{The logarithm of $b''$ is printed $3.8757219$; from $\log b'=8.2388582$ it should be $5.8757219$ (as the $b'''$ row, $1.1493838$, presupposes), so the first figure is a misprint.}
Ensuite le calcul des amplitudes $\varphi'$, $\varphi''$, etc. donne les résultats suivans:
\[ \begin{aligned} \varphi&=47^{\circ}\ 3'\ 30''95,\quad 2\varphi'-\varphi=45^{\circ}\\ \varphi'&=46.\ 1.45.475\\ \varphi''&=46.\ 1.29.41\\ \varphi'''&=46.\ 1.29.41. \end{aligned} \]
La valeur du facteur $K'$ se réduit dans cet exemple à $\sqrt{\dfrac{c'}{c}}$, ainsi on a $\log K'=0.0074955$; et comme on a trouvé $\Phi'=46^{\circ}1'29''41$, on aura
\[ F(\varphi)=K'\log\operatorname{tang}68^{\circ}\ 0'\ 44''705. \]

\origpage{95}
Ce log-tangente pris dans les Tables, et multiplié ensuite par le module pour en faire un logarithme hyperbolique, donnera
\[ \log F=9{,}9650547, \]
ou $F=0.9226877$, comme on l'a déjà trouvé n° 66.

Si dans le même exemple on veut avoir la valeur de la fonction complète, il faudra faire $\varphi=90^{\circ}$, et calculer successivement les valeurs de $\varphi'$, $\varphi''$, etc., ce qui donnera
\[ \begin{aligned} \varphi'&=82^{\circ}\ 30'\ 0''\ 00\\ \varphi''&=82.28.\ 2.84\\ \varphi'''&=\varphi''. \end{aligned} \]
Donc la limite $\Phi'=82^{\circ}28'2''84$, et ainsi la fonction complète
\[ F^{1}=K'\log\operatorname{tang}86^{\circ}\ 14'\ 1''42, \]
ou $\log F^{1}=0.4421761$, ce qui s'accorde avec la valeur trouvée n° 66.

Nous remarquerons que dans cet exemple il a été nécessaire de recourir à un moyen particulier pour déterminer avec la précision convenable, l'amplitude $\varphi''$ qui se confond sensiblement avec la limite $\Phi'$. L'équation $\sin(2\varphi''-\varphi')=c'\sin\varphi'$, la plus directe pour déterminer $\varphi''$, n'est pas propre à donner bien exactement les fractions de seconde contenues dans $\varphi''$, parce que ces fractions influent très-peu sur le sinus d'un angle de $82^{\circ}$, trop rapproché de l'angle droit. Dans ce cas, et dans tous les semblables qui peuvent se rencontrer, on déterminera $\varphi''$ beaucoup plus exactement au moyen de l'équation $\operatorname{tang}(\varphi'-\varphi'')=b''\operatorname{tang}\varphi''$, ou $\varphi'-\varphi''=Rb''\operatorname{tang}\varphi''$, $R$ étant le nombre de secondes contenues dans le rayon: pour cela, il faudra mettre dans le second membre $\varphi'$ au lieu de $\varphi''$, ce qui donnera une première valeur approchée de $\varphi'-\varphi''$, et par conséquent une de $\varphi''$. Substituant de nouveau cette valeur à la place de $\varphi''$ dans le second membre, on aura une seconde valeur de $\varphi'-\varphi''$ qui devra être approchée au moins jusqu'aux centièmes de seconde.

Dans l'exemple dont il s'agit, on fera donc d'abord $\varphi'-\varphi''=Rb''\operatorname{tang}82^{\circ}30'$, ce qui donnera $\varphi'-\varphi''=1'57''68$ et $\varphi''=82^{\circ}28'2''32$. Substituant de nouveau cette valeur au lieu de $\varphi''$, on aura plus exactement $\varphi'-\varphi''=Rb''\operatorname{tang}82^{\circ}\,28'\,2''.32=1'\,57''\,167$, d'où $\varphi''=82^{\circ}\,28'\,2''\,833$.

\origpage{96}
\subsection*{(70).}

On voit maintenant qu'il y a deux méthodes pour déterminer la valeur approchée d'une fonction de première espèce $F(c,\varphi)$ dont le module et l'amplitude sont connus. Si $c$ est plus petit que $\sqrt{\frac{1}{2}}$ ou $\sin45^{\circ}$, il conviendra de se servir de la méthode du n° 65, suivant laquelle calculant les modules décroissans $c^{0}$, $c^{00}$, $c^{000}$, etc. et les amplitudes croissantes $\varphi^{0}$, $\varphi^{00}$, $\varphi^{000}$, etc., on obtient la formule $F(c,\varphi)=K\Phi$.

Si le module $c$ est plus grand que $\sin45^{\circ}$, il conviendra de se servir de la méthode du n° 68, qui consiste à calculer les modules croissans $c'$, $c''$, $c'''$, etc., et les amplitudes décroissantes $\varphi'$, $\varphi''$, $\varphi'''$, etc., d'où l'on conclut $F(c,\varphi)=K'\log\operatorname{tang}(45^{\circ}+\frac{1}{2}\Phi')$.

Ces deux méthodes s'étendent à volonté l'une et l'autre au-delà de la limite que nous avons fixée. Mais pour diminuer autant qu'il est possible le nombre des transformées, il est bon de s'en tenir à cette limite, et de cette manière on n'aura jamais besoin de calculer plus de trois termes tant de la série des modules que de celle des amplitudes, pour obtenir un résultat approché jusqu'au septième rang de décimales.

Il est fort remarquable que la fonction $F$ puisse toujours s'exprimer à volonté par un arc de cercle ou par un logarithme; mais on voit qu'elle se rapproche davantage des arcs de cercle si on a $c^{2}<\frac{1}{2}$, et qu'elle a plus d'affinité avec les logarithmes, si on a $c^{2}>\frac{1}{2}$.

\subsection*{(71).}

Nous avons fait voir dans l'article 67 comment on détermine l'amplitude $\varphi$ qui répond à une valeur donnée de la fonction $F(c,\varphi)$. La méthode exposée dans cet article, a toute la perfection qu'on peut desirer lorsque $c^{2}$ est $<\frac{1}{2}$, et elle peut s'appliquer avec succès lorsque $c^{2}$ s'approche beaucoup plus de l'unité. Cependant si l'on veut que le nombre des transformées soit le plus petit possible, il faudra, lorsque $c^{2}$ sera $>\frac{1}{2}$, appliquer la méthode d'approximation de l'article 68.

Pour cela on commencera par calculer la suite $c'$, $c''$, $c'''$, etc. jusqu'à un terme qui ne diffère pas sensiblement de l'unité. Cette limite est $c'''$ ou même $c''$, lorsqu'on ne veut pas pousser l'exactitude au-delà de six ou sept décimales. De là on tirera $K'=\sqrt{\left(\dfrac{c'c''c'''\,\text{etc.}}{c}\right)}$; $K'$ étant connu, on connaîtra $\log\operatorname{tang}(45^{\circ}+\frac{1}{2}\Phi')=\dfrac{F}{K'}$. De ce

\origpage{97}
logarithme hyperbolique ou du logarithme vulgaire qui lui correspond, on tirera la valeur de l'angle $\Phi'$, limite des angles $\varphi'$, $\varphi''$, $\varphi'''$, etc. Si on s'est arrêté à $c'''$ en négligeant la différence $1-c'''$, alors $\varphi'''$ pourra être pris pour la limite $\Phi'$. Connaissant $\varphi'''$, on remontera successivement aux valeurs de $\varphi''$, $\varphi'$, $\varphi$ par les équations
\[ \begin{aligned} \operatorname{tang}(\varphi''-\varphi''')&=b'''\operatorname{tang}\varphi'''\\ \operatorname{tang}(\varphi'-\varphi'')&=b''\operatorname{tang}\varphi''\\ \operatorname{tang}(\varphi-\varphi')&=b'\operatorname{tang}\varphi'; \end{aligned} \]
on connaîtra donc l'amplitude cherchée $\varphi$.

On peut par ces méthodes déterminer la fonction $F$ en connaissant son amplitude, ou réciproquement, déterminer l'amplitude en connaissant la fonction; de manière qu'il suffira de calculer quatre à cinq termes au plus de la série des modules, et un pareil nombre de termes de la série des amplitudes, pour avoir dix décimales exactes, si les Tables dont on fait usage sont à dix décimales. Si elles en avaient vingt, il suffirait de calculer un terme de plus dans les séries mentionnées pour avoir des résultats exacts jusqu'à la vingtième décimale, et ainsi de suite.

\section*{Propriétés particulières des fonctions $F(c)$, $F(b)$, dont les modules sont complémens l'un de l'autre.}

\subsection*{(72).}

Nous avons vu qu'en prolongeant la suite indéfinie des modules dérivés d'un module donné $c$, tant dans le sens où ils convergent vers la limite 0, que dans le sens où ils convergent vers la limite 1, on a pour les fonctions indéfinies et pour les fonctions définies, ces deux suites d'équations
\[ \begin{aligned} F(c,\varphi)&=\frac{1+c^{0}}{2}F(c^{0},\varphi^{0}) &\qquad& F^{1}(c)=(1+c^{0})F^{1}(c^{0})\\ F(c,\varphi)&=\frac{1+c^{0}}{2}\cdot\frac{1+c^{00}}{2}F(c^{00},\varphi^{00}) && F^{1}(c)=(1+c^{0})(1+c^{00})F^{1}(c^{00})\\ F(c,\varphi)&=\frac{1+c^{0}}{2}\cdot\frac{1+c^{00}}{2}\cdot\frac{1+c^{000}}{2}F(c^{000},\varphi^{000}) && F^{1}(c)=(1+c^{0})(1+c^{00})(1+c^{000})F^{1}(c^{000})\\ &\text{etc.} && \text{etc.} \end{aligned} \]

\origpage{98}
\[ \begin{aligned} F(c,\varphi)&=\frac{2}{1+c}F(c',\varphi') &\qquad& F^{1}(c)=\frac{1}{1+c}F^{1}(c')\\ F(c,\varphi)&=\frac{2}{1+c}\cdot\frac{2}{1+c'}F(c'',\varphi'') && F^{1}(c)=\frac{1}{1+c}\cdot\frac{1}{1+c'}F^{1}(c'')\\ F(c,\varphi)&=\frac{2}{1+c}\cdot\frac{2}{1+c'}\cdot\frac{2}{1+c''}F(c''',\varphi''') && F^{1}(c)=\frac{1}{1+c}\cdot\frac{1}{1+c'}\cdot\frac{1}{1+c''}F^{1}(c''')\\ &\text{etc.} && \text{etc.} \end{aligned} \]
Mettons dans ces dernières formules $C$, $B$, $C'$, $B'$, etc. au lieu de $c$, $b$, $c'$, $b'$, nous aurons d'abord
\[ F(C,\varphi)=\frac{2}{1+C}F(C',\varphi'),\quad F^{1}(C)=\frac{1}{1+C}F^{1}(C'). \]
Supposons ensuite qu'on ait $C=b$, et par suite $B=c$, on aura $C'=\dfrac{2\sqrt{C}}{1+C}=\dfrac{2\sqrt{b}}{1+b}=b^{0}$, par conséquent $B'=c^{0}$; de même on aurait dans les transformées suivantes $C''=b^{00}$, $B''=c^{00}$, $C'''=b^{000}$, $B'''=c^{000}$, etc. Mais on a en même temps l'équation $c^{0}=\dfrac{1-b}{1+b}$ qui donne $\dfrac{1}{1+C}=\dfrac{1}{1+b}=\dfrac{1+c^{0}}{2}$, $\dfrac{1}{1+C'}=\dfrac{1}{1+b^{0}}=\dfrac{1+c^{00}}{2}$, etc; donc la seconde série d'équations transférée au module complémentaire donnera
\[ \begin{aligned} F(b,\varphi)&=(1+c^{0})F(b^{0},\varphi') &\qquad& F^{1}(b)=\frac{1+c^{0}}{2}F^{1}(b^{0})\\ F(b,\varphi)&=(1+c^{0})(1+c^{00})F(b^{00},\varphi'') && F^{1}(b)=\frac{1+c^{0}}{2}\cdot\frac{1+c^{00}}{2}F^{1}(b^{00})\\ F(b,\varphi)&=(1+c^{0})(1+c^{00})(1+c^{000})F(b^{000},\varphi''') && F^{1}(b)=\frac{1+c^{0}}{2}\cdot\frac{1+c^{00}}{2}\cdot\frac{1+c^{000}}{2}F^{1}(b^{000})\\ &\text{etc.} && \text{etc.} \end{aligned} \]
Or à mesure que $b^{i}$ devient plus grand, on a de plus en plus exactement $F(b^{i},\varphi^{i})=\log\operatorname{tang}(45^{\circ}+\frac{1}{2}\varphi^{i})$ et $F^{1}(b^{i})=\log\left(\dfrac{4}{c^{i}}\right)$; donc on aura par une approximation toujours croissante,
\[ \begin{aligned} F(b,\varphi)&=(1+c^{0})\log\operatorname{tang}(45^{\circ}+\tfrac{1}{2}\varphi') &\qquad& F^{1}(b)=\frac{1+c^{0}}{2}\log\frac{4}{c^{0}}\\ F(b,\varphi)&=(1+c^{0})(1+c^{00})\log\operatorname{tang}(45^{\circ}+\tfrac{1}{2}\varphi'') && F^{1}(b)=\frac{1+c^{0}}{2}\cdot\frac{1+c^{00}}{2}\log\frac{4}{c^{00}}\\ F(b,\varphi)&=(1+c^{0})(1+c^{00})(1+c^{000})\log\operatorname{tang}(45^{\circ}+\tfrac{1}{2}\varphi''') && F^{1}(b)=\frac{1+c^{0}}{2}\cdot\frac{1+c^{00}}{2}\cdot\frac{1+c^{000}}{2}\log\frac{4}{c^{000}}\\ &\text{etc.} && \text{etc.} \end{aligned} \]

\origpage{99}
La première série d'équations donne semblablement, pour le module $c$, les formules d'approximation
\[ \begin{aligned} F(c,\varphi)&=\frac{1+c^{0}}{2}\varphi^{0} &\qquad& F^{1}(c)=(1+c^{0})\frac{\pi}{2}\\ F(c,\varphi)&=\frac{1+c^{0}}{2}\cdot\frac{1+c^{00}}{2}\varphi^{00} && F^{1}(c)=(1+c^{0})(1+c^{00})\frac{\pi}{2}\\ F(c,\varphi)&=\frac{1+c^{0}}{2}\cdot\frac{1+c^{00}}{2}\cdot\frac{1+c^{000}}{2}\varphi^{000} && F^{1}(c)=(1+c^{0})(1+c^{00})(1+c^{000})\frac{\pi}{2}\\ &\text{etc.} && \text{etc.} \end{aligned} \]
Il y a diverses conséquences à tirer de ces formules.

\subsection*{(73).}

Dans le cas de $b=c=\sqrt{\frac{1}{2}}$, où l'on a $c'=b^{0}$, $b'=c^{0}$, $c''=b^{00}$, $b''=c^{00}$, etc., ces formules offrent pour la même fonction $F(c,\varphi)$ deux séries de valeurs qui étant comparées terme à terme, donnent
\[ \begin{aligned} F(c^{0},\varphi^{0})&=2F(c',\varphi') &\qquad& F^{1}(c^{0})=\tfrac{1}{2}F^{1}(c')\\ F(c^{00},\varphi^{00})&=4F(c'',\varphi'') && F^{1}(c^{00})=\tfrac{1}{4}F^{1}(c'')\\ F(c^{000},\varphi^{000})&=8F(c''',\varphi''') && F^{1}(c^{000})=\tfrac{1}{8}F^{1}(c''')\\ &\text{etc.} && \text{etc.} \end{aligned} \]
Ces formules ont lieu rigoureusement: les suivantes n'ont lieu que d'une manière approchée, mais l'approximation est toujours croissante:
\[ \begin{aligned} \tfrac{1}{2}\varphi^{0}&=\log\operatorname{tang}(45^{\circ}+\tfrac{1}{2}\varphi') &\qquad& \frac{\pi}{2}=\tfrac{1}{2}\log\frac{4}{c^{0}}\\ \tfrac{1}{4}\varphi^{00}&=\log\operatorname{tang}(45^{\circ}+\tfrac{1}{2}\varphi'') && \frac{\pi}{2}=\tfrac{1}{4}\log\frac{4}{c^{00}}\\ \tfrac{1}{8}\varphi^{000}&=\log\operatorname{tang}(45^{\circ}+\tfrac{1}{2}\varphi''') && \frac{\pi}{2}=\tfrac{1}{8}\log\frac{4}{c^{000}}\\ &\text{etc.} && \text{etc.} \end{aligned} \]
Il résulte de là que si la série $\varphi$, $\dfrac{\varphi^{0}}{2}$, $\dfrac{\varphi^{00}}{4}$, $\dfrac{\varphi^{000}}{8}$, etc. a pour limite $\Phi$, et que dans le sens contraire la série $\varphi$, $\varphi'$, $\varphi''$, $\varphi'''$, etc. ait pour limite $\Phi'$, on aura exactement (dans l'hypothèse du module $c=\sin45^{\circ}$),
\[ \Phi=\log\operatorname{tang}(45^{\circ}+\tfrac{1}{2}\Phi'). \]

\origpage{100}
Si on fait en particulier $\varphi=\frac{1}{2}\pi$, on trouvera successivement
\[ \begin{aligned} \varphi'&=67^{\circ}30'0''00\\ \varphi''&=66.30.54.36\\ \varphi'''&=66.30.47.74\\ \varphi^{\text{iv}}&=\varphi'''. \end{aligned} \]
Ainsi la limite $\Phi'$, autant qu'elle peut être déterminée avec une table à sept décimales, est $66^{\circ}30'47''74$, et il en résulte
\[ \frac{\pi}{2}=\log\operatorname{tang}78^{\circ}15'23''87; \]
ce qui est en effet une valeur fort approchée de $\dfrac{\pi}{2}$. On pourrait augmenter indéfiniment le degré d'approximation en calculant la limite $\Phi'$ par des tables plus étendues.

Mais il résulte de l'autre série d'équations un moyen encore plus facile d'avoir la valeur logarithmique du nombre $\pi$, puisqu'on voit que ce nombre est égal à la limite vers laquelle convergent rapidement les termes successifs
\[ \log\frac{4}{c^{0}},\quad \tfrac{1}{2}\log\frac{4}{c^{00}},\quad \tfrac{1}{4}\log\frac{4}{c^{000}},\quad \text{etc.} \]
Pour savoir jusqu'à quel point chaque terme de cette suite approche de la vraie valeur de $\pi$, je désigne par $c^{n}$ le $n^{\text{ième}}$ terme de la suite $c^{0}$, $c^{00}$, etc., et j'appelle $x$ la valeur approchée de $\pi$ donnée par ce terme, ensorte qu'on aura $x=\dfrac{1}{2^{n-1}}\log\dfrac{4}{c^{n}}$. Je suppose ensuite que pour le terme $c^{n+1}$, la quantité $x$ devienne $x-\omega$, on aura $x-\omega=\dfrac{1}{2^{n}}\log\dfrac{4}{c^{n+1}}$. Mais suivant l'article 65, on a $c^{n+1}=\frac{1}{4}(c^{n})^{2}+\dfrac{1\cdot3}{4\cdot6}(c^{n})^{4}$, donc $\dfrac{4}{c^{n+1}}=\left(\dfrac{4}{c^{n}}\right)^{2}[1-\frac{1}{2}(c^{n})^{2}]$, et par conséquent $\log\dfrac{4}{c^{n+1}}=2\log\left(\dfrac{4}{c^{n}}\right)-2c^{n+1}$; donc
\[ x-\omega=\frac{1}{2^{n-1}}\log\frac{4}{c^{n}}-\frac{1}{2^{n-1}}c^{n+1}; \]
et ainsi $\omega=\dfrac{1}{2^{n-1}}c^{n+1}$. Mais puisqu'on a d'une manière très-approchée

\origpage{101}
$\pi=\dfrac{1}{2^{n}}\log\dfrac{4}{c^{n+1}}$, il s'ensuit $c^{n+1}=4e^{-2^{n}\pi}$; donc $\omega=2^{3-n}e^{-2^{n}\pi}$, et $\log\omega=-(n-3)\log2-2^{n}\pi$, ou en logarithmes vulgaires,
\[ \log\omega=-(n-3)\log2-2^{n-3}(10.915). \]
Lorsque $n=3$, cette formule donne $\log\omega=-10.915$; ainsi le troisième terme $\frac{1}{4}\log\dfrac{4}{c^{000}}$ donne déjà la valeur de $\pi$ exacte jusqu'à dix décimales.

Lorsque $n=4$, on a $\log\omega=-22.131$; ainsi le quatrième terme donne 22 décimales exactes, le cinquième en donne 44, le sixième 88, et le septième 175. D'où l'on voit que $\dfrac{1}{64}\log\dfrac{4}{c^{07}}$ est une valeur beaucoup plus approchée de $\pi$ que celle qu'a donnée Wéga avec 140 décimales. Quant à la valeur de $c^{07}$, elle se déduit successivement de celle de $c$, par de simples extractions de racines quarrées, puisque d'un module $c$ on passe au module suivant $c^{0}$, par le moyen de la formule $c^{0}=\dfrac{1-\sqrt{(1-c^{2})}}{1+\sqrt{(1-c^{2})}}$.

\subsection*{(74).}

Supposons maintenant que les deux fonctions complètes $F^{1}(b)$, $F^{1}(c)$ soient entre elles dans un rapport connu, ensorte qu'on ait $F^{1}(b)=mF^{1}(c)$, il suit des formules du n° 72, qu'on aura, 1°. $\dfrac{m\pi}{2}=\log\operatorname{tang}(45^{\circ}+\frac{1}{2}\Phi')$, $\Phi'$ étant la limite des angles $\varphi'$, $\varphi''$, $\varphi'''$, etc. calculés en supposant $\varphi=\frac{1}{2}\pi$; c'est-à-dire que $\dfrac{m\pi}{2}$ sera la limite de la suite
\[ \log\operatorname{tang}(45^{\circ}+\tfrac{1}{2}\varphi'),\quad \log\operatorname{tang}(45^{\circ}+\tfrac{1}{2}\varphi''),\quad \log\operatorname{tang}(45^{\circ}+\tfrac{1}{2}\varphi'''),\quad \text{etc.}; \]
2°. que $\dfrac{m\pi}{2}$ sera également la limite de la suite
\[ \tfrac{1}{2}\log\frac{4}{c^{0}},\quad \tfrac{1}{4}\log\frac{4}{c^{00}},\quad \tfrac{1}{8}\log\frac{4}{c^{000}},\quad \text{etc.} \]
Le cas de $b=c$ est compris dans ces formules en supposant $m=1$; mais il y a d'autres cas où l'on peut en faire l'application.

Ainsi nous avons trouvé que lorsque $c=\sin15^{\circ}$, on a $F^{1}(b)=\sqrt{3}\,F^{1}(c)$. Donc en calculant la suite $c^{0}$, $c^{00}$, etc. d'après le module $c=\sin15^{\circ}=\frac{1}{2}\sqrt{(2-\sqrt{3})}$, on voit que $\dfrac{\pi\sqrt{3}}{2}$ sera égal à la limite

\origpage{102}
de la suite
\[ \tfrac{1}{2}\log\frac{4}{c^{0}},\quad \tfrac{1}{4}\log\frac{4}{c^{00}},\quad \tfrac{1}{8}\log\frac{4}{c^{000}},\quad \text{etc.}; \]
le premier terme donnerait $\pi=\dfrac{2}{\sqrt{3}}\log\dfrac{2}{\operatorname{tang}7^{\circ}\frac{1}{2}}=3.141636$, valeur qui ne diffère de la véritable que dans la cinquième décimale.

Nous verrons ci-après qu'en faisant $c=\sin\alpha$, $\alpha$ étant déterminé par l'équation $\sin2\alpha=\operatorname{tang}^{2}15^{\circ}$, on a $F^{1}(b)=3F^{1}(c)$; ainsi en calculant la suite des modules $c^{0}$, $c^{00}$, etc. d'après la valeur $c=\sin\alpha$, on aura $\dfrac{3\pi}{2}=\frac{1}{2}\log\dfrac{4}{c^{0}}$, $\frac{1}{4}\log\dfrac{4}{c^{00}}$, etc.; et par le premier terme de cette suite, on obtient $\pi=3.1415926627$, l'erreur n'étant que d'une unité décimale du huitième ordre; d'où l'on peut conclure qu'au cinquième terme, l'approximation équivaudrait à 128 décimales environ.

Enfin nous avons vu qu'en prenant $b=\sqrt{2}-1$, $c=\sqrt{(2\sqrt{2}-2)}$, les fonctions $F^{1}(c)$, $F^{1}(b)$ sont telles qu'on a $F^{1}(c)=\sqrt{2}\,F^{1}(b)$; dans ce cas, on ferait $m=\dfrac{1}{\sqrt{2}}$, et les formules précédentes offriront une nouvelle application en donnant la valeur approchée de $\dfrac{\pi}{2\sqrt{2}}$.

\section*{Méthode d'approximation appliquée aux fonctions elliptiques de la seconde espèce.}

\subsection*{(75).}

Considérons généralement une fonction $G$ composée de la première et de la seconde espèce, ensorte qu'on ait
\[ G=\int(A+B\sin^{2}\varphi)\frac{d\varphi}{\Delta}. \]
Si on y substitue les valeurs $\sin^{2}\varphi=\frac{1}{2}(1+c^{0}\sin^{2}\varphi^{0}-\Delta^{0}\cos\varphi^{0})$, $\dfrac{d\varphi}{\Delta}=\dfrac{1+c^{0}}{2}\cdot\dfrac{d\varphi^{0}}{\Delta^{0}}$, on aura la transformée
\[ G=\frac{1+c^{0}}{2}(G^{0}-\tfrac{1}{2}B\sin\varphi^{0}), \]
où l'on suppose $G^{0}=\displaystyle\int(A^{0}+B^{0}\sin^{2}\varphi^{0})\frac{d\varphi^{0}}{\Delta^{0}}$, $A^{0}=A+\frac{1}{2}B$, $B^{0}=\frac{1}{2}Bc^{0}$.

\origpage{103}
On trouvera ensuite par de semblables transformations,
\[ \begin{aligned} G^{0}&=\frac{1+c^{00}}{2}(G^{00}-\tfrac{1}{2}B^{0}\sin\varphi^{00})\\ G^{00}&=\frac{1+c^{000}}{2}(G^{000}-\tfrac{1}{2}B^{00}\sin\varphi^{000})\\ &\text{etc.} \end{aligned} \]
Ainsi la formule intégrale $G$ se ramène successivement aux intégrales semblables $G^{0}$, $G^{00}$, $G^{000}$, etc., et une seule de ces intégrales suffit pour déterminer toutes les autres.

Supposons que cette suite soit prolongée jusqu'à un terme $G^{\mu}$ assez éloigné pour que le module correspondant $c^{\mu}$ soit de l'ordre des quantités qu'on veut négliger; alors on aura $\Delta^{\mu}=1$, et la valeur de $G^{\mu}$ se réduit d'abord à $\displaystyle\int(A^{\mu}+B^{\mu}\sin^{2}\varphi^{\mu})\,d\varphi^{\mu}$; observons ensuite qu'on a $B^{0}=\frac{1}{2}Bc^{0}$, $B^{00}=\frac{1}{2}B^{0}c^{00}=\frac{1}{4}Bc^{0}c^{00}$, $B^{000}=\frac{1}{8}Bc^{0}c^{00}c^{000}$, etc.; donc la suite $B^{0}$, $B^{00}$, $B^{000}$, etc. décroît plus rapidement que la suite $c^{0}$, $c^{00}$, $c^{000}$, etc. Donc on peut faire, après un certain nombre de termes, $B^{\mu}=0$, ce qui donnera $G^{\mu}=A^{\mu}\varphi^{\mu}=A^{\mu}2^{\mu}\Phi$, $\Phi$ étant la limite des angles $\varphi$, $\frac{1}{2}\varphi^{0}$, $\frac{1}{4}\varphi^{00}$, etc.

A l'égard de $A^{\mu}$, puisqu'on a $A^{0}=A+\frac{1}{2}B$, $A^{00}=A^{0}+\frac{1}{2}B^{0}$, etc., l'expression générale de $A^{\mu}$ sera
\[ A^{\mu}=A+\tfrac{1}{2}B\left(1+\frac{c^{0}}{2}+\frac{c^{0}c^{00}}{4}+\frac{c^{0}c^{00}c^{000}}{8}+\text{etc.}\right). \]
Maintenant il est aisé de voir que la valeur de $G$, déterminée par celle de $G^{\mu}$, est composée de deux parties; l'une $KA^{\mu}\Phi$, $K$ désignant le produit $(1+c^{0})(1+c^{00})(1+c^{000})$, etc., l'autre algébrique ou périodique, savoir:
\[ -\left(\frac{1+c^{0}}{2}\right)\frac{B}{2}\sin\varphi^{0}-\left(\frac{1+c^{0}}{2}\right)\left(\frac{1+c^{00}}{2}\right)\frac{Bc^{0}}{4}\sin\varphi^{00}-\text{etc.} \]
Mais comme on a $1+c^{0}=\dfrac{2\sqrt{c^{0}}}{c}$, $1+c^{00}=\dfrac{2\sqrt{c^{00}}}{c^{0}}$, etc., cette seconde partie peut se mettre sous la forme
\[ -\frac{B}{c}\left(\frac{\sqrt{c^{0}}}{2}\sin\varphi^{0}+\frac{\sqrt{c^{0}c^{00}}}{4}\sin\varphi^{00}+\text{etc.}\right). \]

\origpage{104}
Donc la valeur totale de $G$ sera
\[ G=KA^{\mu}\Phi-\frac{B}{c}\left(\frac{\sqrt{c^{0}}}{2}\sin\varphi^{0}+\frac{\sqrt{c^{0}c^{00}}}{4}\sin\varphi^{00}+\frac{\sqrt{c^{0}c^{00}c^{000}}}{8}\sin\varphi^{000}+\text{etc.}\right). \]
Dans le cas où $\varphi=\frac{1}{2}\pi$, la partie algébrique disparaît, et on a simplement $G^{1}=KA^{\mu}\cdot\frac{1}{2}\pi$.

Les suites qui composent la valeur de $G$ seront fort convergentes si on a $c^{2}<\frac{1}{2}$, et alors trois termes suffisent pour donner la valeur de $G$ approchée, jusqu'à sept décimales environ. Les mêmes formules pourront être employées avec avantage, quand même le module $c$ serait beaucoup plus près de l'unité; car si on avait, par exemple, $c=0.985$, il en résulterait à peu près $c^{0}=\sqrt{\frac{1}{2}}$; d'où l'on voit qu'il n'y aurait qu'un terme de plus à calculer, c'est-à-dire quatre termes en tout, pour avoir la valeur de $G$ avec le même degré d'approximation.

Mais si $c$ est beaucoup plus près de l'unité, ou si ayant $c^{2}>\frac{1}{2}$, on veut obtenir un même degré d'approximation avec le moindre nombre possible de transformées, alors il conviendra de se servir de la seconde méthode dont nous avons fait usage en pareil cas, pour les fonctions de la première espèce.

\subsection*{(76).}

Cette méthode consiste à prolonger dans un ordre inverse la série des transformées. Ainsi l'équation entre $G$ et $G^{0}$ qui donne
\[ G^{0}=\frac{2}{1+c^{0}}G+\tfrac{1}{2}B\sin\varphi^{0}, \]
s'applique semblablement aux deux fonctions $G$ et $G^{1}$\ednote{Printed $G^{1}$ (superscript one); the function meant is $G'$, the first of the inverse transforms, as the next display shows; an apparent misprint.}, et on en déduit
\[ G=\frac{2}{1+c}G'+\tfrac{1}{2}B'\sin\varphi, \]
équation où l'on doit supposer $G'=\displaystyle\int(A'+B'\sin^{2}\varphi')\frac{d\varphi'}{\Delta'}$, $A'=A-\frac{1}{2}B'$, $B'=\dfrac{2B}{c}$. On aura semblablement par des transformées ultérieures,
\[ \begin{aligned} G'&=\frac{2}{1+c'}G''+\tfrac{1}{2}B''\sin\varphi''\\ G''&=\frac{2}{1+c''}G'''+\tfrac{1}{2}B'''\sin\varphi''\\ &\text{etc.} \end{aligned} \]
\ednote{The first line prints $\sin\varphi''$ (a reader's pencil stroke crosses it); by the preceding formula $G=\frac{2}{1+c}G'+\frac{1}{2}B'\sin\varphi$ the term should be $\sin\varphi'$, so an apparent misprint.}

\origpage{105}
Mais lorsque la suite $G'$, $G''$, etc. sera prolongée jusqu'à un terme suffisamment éloigné $G^{\mu}$, on pourra faire $c^{\mu}=1$, $\Delta^{\mu}=\cos\varphi^{\mu}$, et alors la fonction $G^{\mu}$ deviendra $\displaystyle\int(A^{\mu}+B^{\mu}\sin^{2}\varphi^{\mu})\frac{d\varphi^{\mu}}{\cos\varphi^{\mu}}$, ce qui donnera en intégrant,
\[ G^{\mu}=(A^{\mu}+B^{\mu})\log\operatorname{tang}(45^{\circ}+\tfrac{1}{2}\varphi^{\mu})-B^{\mu}\sin\varphi^{\mu}, \]
et $\varphi^{\mu}$ sera égal alors à la limite $\Phi'$.

On connaîtra donc la valeur de $G$ avec toute l'exactitude qu'on peut desirer, en prolongeant la suite des modules $c'$, $c''$, $c'''$, etc., et celle de leurs complémens $b'$, $b''$, $b'''$, etc., jusqu'à ce que le dernier terme de ceux-ci $b^{\mu}$ soit assez petit pour être négligé. Alors on pourra faire $c^{\mu}=1$, et calculant un pareil nombre de termes de la série des amplitudes $\varphi'$, $\varphi''$, $\varphi'''$, etc., on aura le dernier terme $\varphi^{\mu}$ qui pourra être pris pour la limite $\Phi'$.

Faisant donc les substitutions nécessaires pour avoir la valeur de $G$ au moyen de $G^{\mu}$, et observant qu'on a $\dfrac{2}{1+c}=\dfrac{c'}{\sqrt{c}}$, $\dfrac{2}{1+c'}=\dfrac{c''}{\sqrt{c'}}$, etc., on trouvera cette formule générale
\[ \begin{aligned} G=K'(A^{\mu}+B^{\mu})\log\operatorname{tang}(45^{\circ}+\tfrac{1}{2}\varphi^{\mu})&-\frac{B}{c}\cdot\frac{2^{\mu}c^{\mu}\sin\varphi^{\mu}}{\sqrt{(cc'c''\ldots c^{\mu-1})}}\\ &+\frac{B}{c}\left(\sin\varphi+\frac{2}{\sqrt{c}}\sin\varphi'+\frac{4}{\sqrt{cc'}}\sin\varphi''\ldots+\frac{2^{\mu-1}\sin\varphi^{\mu-1}}{\sqrt{(cc'c''\ldots c^{\mu-2})}}\right). \end{aligned} \]
On a fait comme ci-dessus $K'=c^{\mu}\sqrt{\left(\dfrac{c'c''\ldots c^{\mu-1}}{c}\right)}$, ou simplement $K'=\sqrt{\left(\dfrac{c'c''\ldots c^{\mu-1}}{c}\right)}$, parce que $c^{\mu}$ peut être pris pour l'unité. Quant à la valeur de $A^{\mu}+B^{\mu}$, si on l'appelle $L^{\mu}$, on aura
\[ L^{\mu}=A+\frac{2^{\mu}B}{cc'c''\ldots c^{\mu-1}}-\frac{B}{c}-\frac{2B}{cc'}-\frac{4B}{cc'c''}\ldots-\frac{2^{\mu-1}B}{cc'c''\ldots c^{\mu-1}}. \]
Mais cette quantité exprimée ainsi par une suite divergente, serait peu commode pour les calculs d'approximation, et il convient de la mettre sous une autre forme. Or, en exprimant $L^{\mu}$ au moyen de $L$ et des différences successives $L'-L$, $L''-L'$, etc., on trouve

\origpage{106}
aisément
\[ L^{\mu}=A+\frac{B}{c}+\frac{2B}{cc'}(1-c')+\frac{4B}{cc'c''}(1-c'')+\text{etc.}; \]
ou, ce qui revient au même,
\[ L^{\mu}=A+B+\frac{Bb}{c}\left[\sqrt{b'}+\sqrt{\left(\frac{b'b''}{c}\right)}+\sqrt{\left(\frac{b'b''b'''}{cc'}\right)}\ldots+\sqrt{\left(\frac{b'b''\ldots b^{\mu}}{cc'\ldots c^{\mu-2}}\right)}\right], \]
et sous cette forme la convergence de la suite est manifeste.

On peut également rendre convergente la suite des quantités algébriques comprises dans la valeur de $G$. Pour cela, soit
\[ \psi(\mu)=\frac{2^{\mu}c^{\mu}\sin\varphi^{\mu}}{\sqrt{(cc'c''\ldots c^{\mu-1})}}-\sin\varphi-\frac{2}{\sqrt{c}}\sin\varphi'-\frac{4}{\sqrt{cc'}}\sin\varphi''\ldots-\frac{2^{\mu-1}\sin\varphi^{\mu-1}}{\sqrt{(cc'c''\ldots c^{\mu-2})}}, \]
on déduit de là
\[ \begin{aligned} \psi(\mu+1)-\psi(\mu)&=\frac{2^{\mu+1}c^{\mu+1}\sin\varphi^{\mu+1}}{\sqrt{(cc'\ldots c^{\mu})}}-\frac{2^{\mu}c^{\mu}\sin\varphi^{\mu}}{\sqrt{(cc'\ldots c^{\mu-1})}}-\frac{2^{\mu}\sin\varphi^{\mu}}{\sqrt{(cc'\ldots c^{\mu-1})}}\\ &=\frac{2^{\mu+1}}{\sqrt{(cc'c''\ldots c^{\mu})}}\left(c^{\mu+1}\sin\varphi^{\mu+1}-\frac{c^{\mu}}{c^{\mu+1}}\sin\varphi^{\mu}\right). \end{aligned} \]
Or lorsque $b^{\mu}$ est devenu très-petit, l'équation entre $\varphi^{\mu}$ et $\varphi^{\mu+1}$, qui est $\operatorname{tang}(\varphi^{\mu}-\varphi^{\mu+1})=b^{\mu+1}\operatorname{tang}\varphi^{\mu+1}$, donne à très-peu près $\varphi^{\mu}=\varphi^{\mu+1}+b^{\mu+1}\operatorname{tang}\varphi^{\mu+1}$, ou $\sin\varphi^{\mu}=(1+b^{\mu+1})\sin\varphi^{\mu+1}$. Donc $c^{\mu+1}\sin\varphi^{\mu+1}-\dfrac{c^{\mu}}{c^{\mu+1}}\sin\varphi^{\mu}=\dfrac{b^{\mu+1}}{c^{\mu+1}}(1-b^{\mu+1})\sin\varphi^{\mu+1}$. Cette quantité est de l'ordre $b^{\mu+1}$; donc on aura pour exprimer $\psi(\mu)$, cette suite très-convergente:
\[ \begin{aligned} \psi(\mu)=c\sin\varphi&+\frac{2}{\sqrt{c}}\left(c'\sin\varphi'-\frac{c}{c'}\sin\varphi\right)\\ &+\frac{4}{\sqrt{cc'}}\left(c''\sin\varphi''-\frac{c'}{c''}\sin\varphi'\right)\\ &+\frac{8}{\sqrt{cc'c''}}\left(c'''\sin\varphi'''-\frac{c''}{c'''}\sin\varphi''\right)\\ &\ \ \vdots\\ &+\frac{2^{\mu}}{\sqrt{(cc'\ldots c^{\mu-1})}}\left(c^{\mu}\sin\varphi^{\mu}-\frac{c^{\mu-1}}{c^{\mu}}\sin\varphi^{\mu-1}\right). \end{aligned} \]

\origpage{107}
Cela posé, la valeur de $G$ exprimée toute entière en suites convergentes sera
\[ G=K'L^{\mu}\log\operatorname{tang}(45^{\circ}+\tfrac{1}{2}\varphi^{\mu})-\frac{B}{c}\psi(\mu). \]
Il ne faut pas perdre de vue que cette approximation est fondée sur ce que dans la valeur de $G^{\mu}$ on a substitué $\cos\varphi^{\mu}$ au radical $\Delta^{\mu}$ dont la vraie valeur est $\sqrt{[\cos^{2}\varphi^{\mu}+(b^{\mu}\sin\varphi^{\mu})^{2}]}$. Cette substitution est permise tant que $b^{\mu}\operatorname{tang}\varphi^{\mu}$ est une quantité négligeable; mais elle ne pourrait plus avoir lieu si $\varphi^{\mu}$ était assez près de $90^{\circ}$, pour que $b^{\mu}\operatorname{tang}\varphi^{\mu}$ cessât d'être très-petit. Supposons, ce qui est le cas le plus défavorable, qu'on ait $\varphi^{\mu}=90^{\circ}$; alors il faudra prolonger les suites d'un terme de plus, ou jusqu'au $(\mu+1)^{\text{ième}}$ terme, ce qui donnera $\operatorname{tang}\varphi^{\mu+1}=\dfrac{1}{\sqrt{(b^{\mu+1})}}$, ou $b^{\mu+1}\operatorname{tang}\varphi^{\mu+1}=\sqrt{(b^{\mu+1})}$, quantité dont le quarré $b^{\mu+1}$ est très-négligeable par rapport à l'unité, puisque $b^{\mu}$ était déjà supposé négligeable. La formule s'appliquera donc sans aucune difficulté; ainsi le moyen de prévenir toute exception, est de prendre $\mu$ assez grand pour que $\varphi^{\mu-1}$ ne surpasse pas $90^{\circ}$, condition qui sera toujours facile à remplir.

Pour avoir la valeur de la fonction complète $G^{1}$, on pourrait faire $\varphi=90^{\circ}$ dans la formule générale; mais cette formule ne se simplifierait pas sensiblement, et elle contiendrait toujours un nombre de termes indéfini. Il sera plus simple, en se conformant à l'observation précédente, de faire $\varphi^{\mu-1}=90^{\circ}$, ce qui donne $\operatorname{tang}\varphi^{\mu}=\dfrac{1}{\sqrt{b^{\mu}}}$, $\sin\varphi^{\mu}=\dfrac{1}{\sqrt{(1+b^{\mu})}}$, et en rétrogradant, $\varphi^{\mu-2}=\pi$, $\varphi^{\mu-3}=2\pi\ldots\varphi^{\mu-i}=2^{i-2}\pi$, $\varphi'=2^{\mu-3}\pi$, $\varphi=2^{\mu-2}\pi=2^{\mu-1}\dfrac{\pi}{2}$. On aura donc $G(\varphi)=2^{\mu-1}G^{1}$; mais dans l'hypothèse que $b^{\mu}$ est de l'ordre des quantités négligeables, on a $\sin\varphi^{\mu}=1-\frac{1}{2}b^{\mu}$, et $\log\operatorname{tang}(45^{\circ}+\frac{1}{2}\varphi^{\mu})=\frac{1}{2}\log\left(\dfrac{1+\sin\varphi^{\mu}}{1-\sin\varphi^{\mu}}\right)=\frac{1}{2}\log\left(\dfrac{4}{b^{\mu}}\right)$; on aura en même temps la quantité $\psi(\mu)=\dfrac{2^{\mu}c^{\mu}}{cK'}\sin\varphi^{\mu}-\dfrac{2^{\mu-1}}{cK'}=\dfrac{2^{\mu-1}}{cK'}$; donc
\[ 2^{\mu-1}G^{1}=\frac{K'L^{\mu}}{2}\log\left(\frac{4}{b^{\mu}}\right)-\frac{2^{\mu-1}B}{c^{2}K'}, \]

\origpage{108}
et enfin
\[ G^{1}=\frac{K'L^{\mu}}{2^{\mu}}\log\left(\frac{4}{b^{\mu}}\right)-\frac{B}{c^{2}K'}. \]
Ces formules pour déterminer tant la fonction indéfinie $G$ que la fonction complète $G^{1}$, permettent de pousser l'approximation aussi loin qu'on voudra. Il faudra prolonger le calcul des modules croissans $c'$, $c''$, $c'''$, etc., et celui de leurs complémens $b'$, $b''\ldots b^{\mu}$, jusqu'à ce que $b^{\mu}$ appartienne à l'ordre de décimales qu'on veut négliger, ou soit inférieur à cet ordre. Il faudra en même temps, s'il s'agit de la fonction non-complète $G$, ne pas donner à $\varphi$ une valeur plus grande que $2^{\mu-2}\pi$. On calculera alors jusqu'à la même limite les valeurs de $K'$, $L^{\mu}$ et $\psi(\mu)$.

\subsection*{(77).}

Pour appliquer ces formules aux arcs d'ellipse ou aux fonctions proprement dites de la seconde espèce, soit $G=E=\displaystyle\int\Delta\,d\varphi$, on aura $A=1$, $B=-c^{2}$.

Par la première méthode, l'expression de l'arc indéfini sera
\[ E(c,\varphi)=KL\Phi+\frac{c\sqrt{c^{0}}}{2}\sin\varphi^{0}+\frac{c\sqrt{c^{0}c^{00}}}{4}\sin\varphi^{00}+\frac{c\sqrt{c^{0}c^{00}c^{000}}}{8}\sin\varphi^{000}+\text{etc.} \]
et celle du quart d'ellipse
\[ E^{1}(c)=KL\frac{\pi}{2}. \]
Dans ces formules, $K$ représente le produit $(1+c^{0})(1+c^{00})(1+c^{000})$, etc., continué jusqu'à un facteur $1+c^{\mu}$, où $c^{\mu}$ soit de l'ordre des quantités négligeables. Cette même quantité peut se mettre sous la forme $K=\dfrac{2\sqrt{c^{0}}}{c}\cdot\dfrac{2\sqrt{c^{00}}}{c^{0}}\cdot\dfrac{2\sqrt{c^{000}}}{c^{00}}\cdot\text{etc.}$, plus commode pour le calcul logarithmique. Enfin la quantité $L$ est donnée par la suite convergente
\[ L=1-\frac{c^{0}}{2}-\frac{c^{2}c^{0}}{4}-\frac{c^{2}c^{0}c^{00}}{4}-\text{etc.}, \]
\ednote{Printed $\frac{c^{0}}{2}$ in the first term and $4$ in the third denominator; the series that agrees with the second form below (and numerically) is $1-\frac{c^{2}}{2}-\frac{c^{2}c^{0}}{4}-\frac{c^{2}c^{0}c^{00}}{8}-\text{etc.}$; apparent misprints.}
à laquelle on peut donner cette autre forme
\[ L=\frac{c^{2}}{4c^{0}}\left(1-\frac{c^{02}c^{00}}{2}-\frac{c^{02}c^{00}c^{000}}{4}-\frac{c^{02}c^{00}c^{000}c^{0000}}{8}-\text{etc.}\right); \]

\origpage{109}
de sorte qu'en faisant $KL=M$, on aurait
\[ M=\left(\frac{c^{2}}{2\sqrt{c^{0}}}\right)\cdot\frac{2\sqrt{c^{00}}}{c^{0}}\cdot\frac{2\sqrt{c^{000}}}{c^{00}}\,\text{etc.}\left(1-\frac{c^{02}c^{00}}{2}-\frac{c^{02}c^{00}c^{000}}{4}-\text{etc.}\right) \]
Ces formules offrent les suites les plus convergentes qu'on puisse desirer pour la rectification de l'ellipse. Elles s'appliquent, comme nous l'avons déjà dit, non-seulement à des valeurs de $c$ plus petites que $\sqrt{\frac{1}{2}}$, mais à des valeurs beaucoup plus grandes et très-rapprochées de l'unité.

\section*{EXEMPLE.}

\subsection*{(78).}

Soit proposé de trouver l'arc $E(c,\varphi)$ lorsque $c=\sin 75^{\circ}$ et $\operatorname{tang}\varphi=\sqrt{\dfrac{2}{\sqrt{3}}}$.

Par les valeurs déjà calculées (art.~66), on a $\log K=0.2460561$, $\Phi=30^{\circ}$, $\varphi^{0}=62^{\circ}36'3''.10$, $\varphi^{00}=119^{\circ}55'47''67$, $\varphi^{000}=240^{\circ}0'0''19$, ensuite on trouvera
\[ \begin{aligned} L&=0.3888658\\ \frac{c\sqrt{c^{0}}}{2}\sin\varphi^{0}&=0.3290186\\ \frac{c\sqrt{c^{0}c^{00}}}{4}\sin\varphi^{00}&=0.0522872\\ \frac{c\sqrt{c^{0}c^{00}c^{000}}}{8}\sin\varphi^{000}&=-0.0013888\\ \frac{c\sqrt{c^{0}c^{00}c^{000}c^{0000}}}{16}\sin\varphi^{0000}&=0.0000010 \end{aligned} \]
La somme des parties algébriques est 0.3799180, ainsi on a $E(c,\varphi)=KL\cdot\frac{\pi}{6}+0.3799180$; on aura en même temps la fonction complète $E^{1}(c)=KL\cdot\frac{\pi}{2}$, d'où résulterait $E(c,\varphi)=\frac{1}{3}E^{1}(c)+0.3799180$. Mais puisque l'arc $E(c,\varphi)$ se mesure par le tiers de $E^{1}(c)$, on trouvera par les formules de la trisection (art.~57) $E(c,\varphi)=\frac{1}{3}E^{1}(c)+\dfrac{c^{2}\cos^{3}\varphi}{3b}$; le résultat précédent s'accorde parfaitement avec cette formule, car on a en effet $\dfrac{c^{2}\cos^{3}\varphi}{3b}=0.3799179$.

Le même calcul donne de plus $\log E^{1}(c)=0.0319757$ et $E^{1}(c)=1.0764049$; donc dans le cas proposé, l'arc $E(c,\varphi)=0.7387196$; et on voit que la partie déterminée algébriquement, forme plus de la moitié de l'arc. Au reste la valeur trouvée de $E^{1}(c)$ s'accorde avec

\origpage{110}
celle qu'on déduirait de la formule de l'article 41, en y mettant pour $F^{1}(c)$ sa valeur trouvée art.~66.

\subsection*{(79).}

Si le module $c$ est assez peu différent de l'unité pour qu'on juge nécessaire d'appliquer la seconde méthode, il faudra, dans les formules de l'art.~76, faire $A=1$, $B=-c^{2}$, ce qui donnera
\[ L^{\mu}=b^{2}-bc\left[\sqrt{b'}+\sqrt{\left(\frac{b'b''}{c}\right)}+\sqrt{\left(\frac{b'b''b'''}{cc'}\right)}\ldots+\sqrt{\left(\frac{b'b''\ldots b^{\mu}}{cc'\ldots c^{\mu-2}}\right)}\right]. \]
Quant à la valeur de $K'$ et à celle de la fonction $\psi(\mu)$, elles ne dépendent point des valeurs de $A$ et de $B$; ainsi elles restent les mêmes que dans l'article cité. On aura donc, d'après ces valeurs, la fonction indéfinie
\[ E(c,\varphi)=K'L^{\mu}\operatorname{tang}(45^{\circ}+\tfrac{1}{2}\Phi')+c\psi(\mu), \]
\ednote{Printed without $\log$ and with $\Phi'$; by the general formula of art.\ 76 and the case $\mu=2$ below, one expects $\log\operatorname{tang}(45^{\circ}+\frac{1}{2}\varphi^{\mu})$; an apparent misprint.}
et la fonction complète
\[ E^{1}(c)=\frac{K'L^{\mu}}{2^{\mu}}\log\frac{4}{b^{\mu}}+\frac{1}{K'}. \]
Ces formules sont ce que l'analyse peut offrir de plus simple pour le calcul des arcs d'ellipse, lorsque le module est très-près de l'unité; elles supposent qu'on prenne pour $\mu$ un nombre d'unités assez grand pour que $b^{\mu}$ appartienne à l'ordre de décimales qu'on veut négliger.

Si on veut avoir les valeurs de $E(c,\varphi)$ et de $E^{1}(c)$ approchées seulement jusqu'au septième ordre de décimales, il suffira de faire $\mu=2$ dans les formules précédentes; on aura d'abord $L^{\mu}=b^{2}-bc\sqrt{b'}-bc\sqrt{\left(\dfrac{b'b''}{c}\right)}$, et parce qu'on a en général $b\sqrt{b'}=1-c$, $b'\sqrt{b''}=1-c'$, $c'=\dfrac{2\sqrt{c}}{1+c}$, on trouve plus simplement $L^{\mu}=\dfrac{b^{2}}{1+\sqrt{c}}$: mais dans le même cas on a $K'=\sqrt{\dfrac{c'}{c}}$. Donc les deux formules générales se réduisent à celles-ci
\[ \begin{aligned} E(c,\varphi)&=\frac{b^{2}}{1+\sqrt{c}}\sqrt{\frac{c'}{c}}\log\operatorname{tang}(45^{\circ}+\tfrac{1}{2}\varphi'')+c^{2}\sin\varphi\\ &\quad+2\sqrt{c}\left(c'\sin\varphi'-\frac{c}{c'}\sin\varphi\right)+4\sqrt{\frac{c}{c'}}(\sin\varphi''-c'\sin\varphi'),\\ E^{1}(c)&=\frac{b^{2}}{1+\sqrt{c}}\cdot\frac{1}{4}\sqrt{\frac{c'}{c}}\cdot\log\frac{4}{b''}+\sqrt{\frac{c}{c'}}. \end{aligned} \]

\origpage{111}
La dernière est la même que $E^{1}(c)=\dfrac{b^{2}}{1+\sqrt{c}}F^{1}(c)+\sqrt{\dfrac{c}{c'}}$, et sous cette forme, elle offre une relation fort simple entre les fonctions $F^{1}(c)$, $E^{1}(c)$, et d'autant plus approchée que $b$ sera plus petit.

\subsection*{(80).}

Appliquons ces formules à l'exemple de l'article 78, ce qui est un cas peu favorable, puisque la valeur de $c$ diffère sensiblement de l'unité. On connaît déjà par l'article 69 les valeurs de $c'$, $b'$, $c''$, $b''$, ainsi que celles des amplitudes $\varphi$, $\varphi'$, $\varphi''$. Au moyen de ces valeurs, le calcul de la fonction complète $E^{1}(c)$ se fera ainsi:
\[ \begin{aligned} &1+\sqrt{c}= && 1.9828152 &\qquad& \log(1+\sqrt{c})\ldots && 0.2972822\\ &\log 4\ldots\ldots && 0.6020600 && \log b^{2}\ldots\ldots && 8.8259924\\ &\log b''\ldots\ldots && 5.8757219 && \text{différ.}\ldots\ldots && 8.5287102\\ &\log\frac{4}{b''}\ldots && 4.7263381 && \log K'\ldots\ldots && 0.0074955\\ &\tfrac{1}{4}\log\frac{4}{b''}= && 1.1815845 && \log 1.1815\ \text{etc.}\ldots && 0.0724648\\ &&&&& \log 2.3025\ \text{etc.}\ldots && 0.3622157\\ &&&&& \log x\ldots\ldots && 8.9708862 \end{aligned} \]
\[ \begin{aligned} x&=0.0935161\\ \frac{1}{K'}&=0.9828891\\ \text{Somme}\ldots\ E^{1}(c)&=1.0764052 \end{aligned} \]
Cette valeur s'accorde avec celle qu'on a déjà trouvée, autant que les petites parties négligées peuvent le permettre.

Quant au calcul de l'arc $E(c,\varphi)$, nous avons déjà trouvé la valeur de $E(c,\varphi)$\ednote{Printed $E(c,\varphi)$; the logarithm taken next is that of $F(c,\varphi)$, so $F(c,\varphi)$ is expected here; an apparent misprint.} $=K'\log\operatorname{tang}(45^{\circ}+\frac{1}{2}\varphi'')$, il en résulte
\[ \begin{aligned} \log F(c,\varphi)&=9.9650547\\ \log\frac{b^{2}}{1+\sqrt{c}}&=8.5287102\\ \log y&=8.4937649,\quad y=0.0311720. \end{aligned} \]
On trouve ensuite les parties algébriques

\origpage{112}
\[ \begin{aligned} &c^{2}\sin\varphi && =0.6830127\\ &2\sqrt{c}\left(c'\sin\varphi'-\frac{c}{c'}\sin\varphi\right) && =0.0243225\\ &4\sqrt{\frac{c}{c'}}(\sin\varphi''-c'\sin\varphi') && =0.0002123\\ &\text{somme}\ldots\ldots && 0.7075475\\ &\text{ajoutant}\ldots\ y && =0.0311720\\ &\text{on a}\ E(c,\varphi) && =0.7387195 \end{aligned} \]
ce qui s'accorde encore avec la valeur trouvée par la première méthode.

\subsection*{(81).}

Pour appliquer les formules précédentes à la rectification de l'hyperbole, il ne s'agit que d'évaluer l'intégrale $G=\displaystyle\int\frac{c^{2}\,d\varphi\cos^{2}\varphi}{\Delta}$ qui représente la quantité $\Delta\operatorname{tang}\varphi-\Upsilon$, différence entre l'arc d'hyperbole et sa tangente. Or en faisant dans la formule générale de l'art.~75, $A=c^{2}$, $B=-c^{2}$, on a la fonction indéfinie
\[ \begin{aligned} G(c,\varphi)=\tfrac{1}{2}Kc^{2}\Phi&\left(1-\frac{c^{0}}{2}-\frac{c^{0}c^{00}}{4}-\frac{c^{0}c^{00}c^{000}}{8}-\text{etc.}\right)\\ &+\frac{c\sqrt{c^{0}}}{2}\sin\varphi+\frac{c\sqrt{c^{0}c^{00}}}{4}\sin\varphi^{00}+\frac{c\sqrt{c^{0}c^{00}c^{000}}}{8}\sin\varphi^{000}+\text{etc.} \end{aligned} \]
\ednote{The second line prints $\sin\varphi$ in its first term; the corresponding series on p. 108 has $\sin\varphi^{0}$, so either a misprint or a superscript lost in printing.}
et la fonction définie
\[ G^{1}(c)=Kc^{2}\cdot\frac{\pi}{4}\left(1-\frac{c^{0}}{2}-\frac{c^{0}c^{00}}{4}-\frac{c^{0}c^{00}c^{000}}{8}-\text{etc.}\right): \]
celle-ci est la différence entre l'asymptote et la courbe, qui s'exprime ainsi très-simplement par une suite fort convergente.

Lorsque l'hyperbole est équilatère, la quantité $G^{1}(c)$ qui en général a pour valeur $E^{1}(c)-b^{2}F^{1}(c)$, se réduit à $E^{1}(c)-\frac{1}{2}F^{1}(c)$; mais comme dans le cas particulier où $c^{2}=\frac{1}{2}=b^{2}$, on a $\dfrac{\pi}{2}=(2E^{1}-F^{1})F^{1}$, la valeur de $G^{1}$ se réduira ultérieurement à $\dfrac{\frac{1}{4}\pi}{F^{1}}$, et parce que $F^{1}=\frac{1}{2}\pi K$, on aura encore plus simplement $G^{1}=\dfrac{1}{2K}$. On a trouvé ci-dessus $\log K=0.0720074$, ainsi on aura $\log G^{1}=9.6269626$.

On peut regarder comme applicable à tous les cas, la formule

\origpage{113}
$\Upsilon=\Delta\operatorname{tang}\varphi-G(c,\varphi)$; cependant si $c$ était extrêmement près de l'unité, on pourrait, pour déterminer la fonction $G$, faire usage des formules de l'article 76, après avoir fait $A=c^{2}$, $B=-c^{2}$.

\section*{Formules remarquables pour déterminer les fonctions complètes $F^{1}(c)$, $E^{1}(c)$, lorsque $c$ est peu éloigné de l'une de ses limites.}

\subsection*{(82).}

Lorsque le module $c$ est fort peu différent de l'unité, nous avons trouvé (art.~76) les deux formules
\[ \begin{aligned} F^{1}(c)&=\tfrac{1}{4}\sqrt{\left(\frac{c'}{c}\right)}\log\frac{4}{b''}\\ E^{1}(c)&=\frac{b^{2}}{1+\sqrt{c}}F^{1}(c)+\sqrt{\left(\frac{c}{c'}\right)}. \end{aligned} \]
Ces formules donneront environ dix décimales exactes si $c=\sin 80^{\circ}$; elles en donneraient un plus grand nombre si $c$ était plus grand; elles n'en donneront que sept lorsqu'on fera $c=\sin 64^{\circ}$, et ainsi à proportion dans les autres cas. Du reste leur usage est très-commode puisqu'il exige seulement qu'on calcule avec précision les valeurs de $c'$, $b'$ et $b''$. Il faut pour cet effet se rappeler qu'en faisant $c=\cos\mu$, ce qui donne $b=\sin\mu$, on aura $c'=\dfrac{\sqrt{\cos\mu}}{\cos^{2}\frac{1}{2}\mu}$, $b'=\operatorname{tang}^{2}\frac{1}{2}\mu$. Si $\mu$ est fort petit, les tables donneront avec précision $\cos\mu$ et $\cos\frac{1}{2}\mu$; connaissant $\cos\frac{1}{2}\mu$ et $\sin\mu$, on en déduira $\sin\frac{1}{2}\mu=\dfrac{\sin\mu}{2\cos\frac{1}{2}\mu}$, et $\operatorname{tang}\frac{1}{2}\mu=\dfrac{\sin\frac{1}{2}\mu}{\cos\frac{1}{2}\mu}$, ce qui fera connaître $b'$; enfin pour éviter les trop petits angles, on observera que $b''$ se déduit de $b'$ comme $c^{0}$ de $c$, et qu'ainsi on a $b''=\frac{1}{4}b'^{2}+\dfrac{1\cdot 3}{4\cdot 6}b'^{4}+\text{etc.}$; d'où résulte $\dfrac{4}{b''}=\left(\dfrac{4}{b'}\right)^{2}(1-\frac{1}{2}b'^{2})$, et par conséquent $\frac{1}{2}\log\dfrac{4}{b''}=\log\dfrac{4}{b'}-\frac{1}{4}b'^{2}m$, $m$ étant mis pour le nombre $0.43429$ etc.; d'où l'on voit qu'avec une légère correction facile à calculer, on déduira $\log\dfrac{4}{b''}$ de $\log\dfrac{4}{b'}$. Le reste n'a aucune difficulté, en se rappelant que les logarithmes de la formule sont des logarithmes hyperboliques qu'il faudra réduire en logarithmes vulgaires.

\origpage{114}
\subsection*{(83).}

Si on substitue la valeur de $E^{1}(c)$ donnée par la seconde des formules précédentes, dans l'équation de l'article 42, on en déduira
\[ F^{1}(c)=\frac{\tfrac{1}{2}\pi-\sqrt{\left(\frac{1+c}{2}\sqrt{c}\right)}\cdot F^{1}(b)}{E^{1}(b)-\left(\frac{c^{2}+\sqrt{c}}{1+\sqrt{c}}\right)F^{1}(b)}, \]
et si on observe que la valeur de $F^{1}(c)$ contient un logarithme que le second membre ne saurait offrir par le développement des quantités $F^{1}(b)$, $E^{1}(b)$, en séries ordonnées suivant les puissances de $b$ supposé très-petit, on en conclura que l'équation précédente ne peut subsister à moins que le numérateur et le dénominateur du second membre ne soient sensiblement égaux à zéro. Ainsi on devra avoir les deux équations
\[ \begin{aligned} F^{1}(b)&=\frac{\tfrac{1}{2}\pi}{\sqrt{\left(\frac{1+c}{2}\sqrt{c}\right)}}=\tfrac{1}{2}\pi\sqrt{\left(\frac{c'}{c}\right)}\\ E^{1}(b)&=\frac{1+c\sqrt{c}}{1+\sqrt{c}}\cdot\tfrac{1}{2}\pi\sqrt{c'}. \end{aligned} \]
Changeons $b$ en $c$, et réciproquement, afin de rapporter ces nouvelles formules au module $c$, nous aurons
\[ \begin{aligned} F^{1}(c)&=\frac{\tfrac{1}{2}\pi}{\sqrt{\left(\frac{1+b}{2}\sqrt{b}\right)}}=\tfrac{1}{2}\pi\sqrt{\left(\frac{b^{0}}{b}\right)}\\ E^{1}(c)&=\frac{1+b\sqrt{b}}{1+\sqrt{b}}\cdot\tfrac{1}{2}\pi\sqrt{b^{0}}=(2-b^{0})\frac{\pi}{2}\sqrt{\left(\frac{b}{b^{0}}\right)}. \end{aligned} \]
Ces formules sont encore plus simples que celles de l'article précédent; elles ne supposent d'autre calcul préliminaire que celui de $b$ et $b^{0}$; or en faisant $c=\sin\mu$, on a $b=\cos\mu$, et $b^{0}=\dfrac{\sqrt{\cos\mu}}{\cos^{2}\frac{1}{2}\mu}$: d'ailleurs $\log(2-b^{0})$ est très-facile à déduire de $\log b^{0}$; car comme $1-b^{0}$ est très-petit, si l'on fait $\log b^{0}=-\delta$, on aura $\log(2-b^{0})=\delta-\delta^{2}(2.3025,\text{etc.})$.

Ces formules ne sont pas moins exactes pour les petites valeurs de $c$, que celles de l'article précédent pour les valeurs de $c$ peu différentes de l'unité; les unes et les autres s'accordent avec les formules de l'article 45, et on peut s'assurer qu'elles équivalent au développement d'un assez grand nombre de termes de ces dernières.

\origpage{115}
\subsection*{(84).}

Appliquons ces formules au calcul des fonctions $F^{1}(c)$, $F^{1}(b)$, en supposant $c=\dfrac{\sqrt[4]{4}-\sqrt[4]{3}}{1+\sqrt{3}}$ et $b=\dfrac{\sqrt[4]{4}+\sqrt[4]{3}}{1+\sqrt{3}}$. Soit $c=\sin\alpha$, $b=\cos\alpha$, on aura $\sin 2\alpha=\operatorname{tang}^{2}15^{\circ}$ et $\operatorname{tang}(45^{\circ}+\alpha)=\sqrt{\left(\dfrac{2}{\sqrt{3}}\right)}$; de l'une ou de l'autre de ces équations on tirera
\[ \alpha=2^{\circ}\,3'\,30''.94970. \]
Cet angle étant assez petit, il faudra user de quelques précautions pour calculer les valeurs des quantités cherchées jusqu'à dix décimales. On trouve d'abord directement par les tables, $\log\sin 2\alpha=8.8561049048$; ensuite la valeur connue de $\alpha$ donne $\log\cos\alpha=9.9997196211$; de là on déduit
\[ \log\sin\alpha=8.55535\ 52880. \]
Pour le calcul de la formule $F^{1}(c)=\frac{1}{2}\pi\sqrt{\dfrac{b^{0}}{b}}$, je fais $k=\sqrt{\dfrac{b^{0}}{b}}=\dfrac{1}{\cos\frac{1}{2}\alpha\sqrt[4]{\cos\alpha}}$, et je trouve
\[ \begin{aligned} &\cos\tfrac{1}{2}\alpha\ldots\ldots && 9.99992\ 99167\\ &\sqrt[4]{\cos\alpha}\ldots\ldots && 9.99992\ 99053\\ &\text{somme}\ldots\ldots && 9.99985\ 98220\\ &k\ldots\ldots && 0.00014\ 01780\\ &\tfrac{1}{2}\pi\ldots\ldots && 0.19611\ 98770\\ &\text{Donc}\ \log F^{1}(c) && =0.19626\ 00550 \end{aligned} \]
Le calcul de $F^{1}(b)$ se fera par la première formule qui donne $F^{1}(b)=\frac{1}{4}k\log\dfrac{4}{c^{00}}$. Or on a $c^{0}=\operatorname{tang}^{2}\frac{1}{2}\alpha$; ainsi $\log c^{0}$ se déduit des logarithmes déjà trouvés en cette sorte:

\origpage{116}
\[ \begin{aligned} &\cos\tfrac{1}{2}\alpha\ldots\ldots && 9.99992\ 99167\\ &2\ldots\ldots && 0.30102\ 99957\\ &2\cos\tfrac{1}{2}\alpha\ldots\ldots && 0.30095\ 99124\\ &\sin\alpha\ldots\ldots && 8.55535\ 52880\\ &\sin\tfrac{1}{2}\alpha\ldots\ldots && 8.25439\ 53756\\ &\cos\tfrac{1}{2}\alpha\ldots\ldots && 9.99992\ 99167\\ &\operatorname{tang}\tfrac{1}{2}\alpha\ldots\ldots && 8.25446\ 54589\\ &c^{0}\ldots\ldots && 6.50893\ 09178\\ &2\ldots\ldots && 0.30102\ 99957\\ &\tfrac{1}{2}c^{0}\ldots\ldots && 6.20790\ 09221\\ &\tfrac{1}{4}c^{02}\ldots\ldots && 2.41580\ 18442. \end{aligned} \]
Il faut de plus chercher la valeur de $c^{00}$; mais, comme on l'a déjà observé, on peut suppléer au calcul de $c^{00}$ par la formule $\frac{1}{2}\log\dfrac{4}{c^{00}}=\log\dfrac{4}{c^{0}}-\frac{1}{4}c^{02}(0.434\ \text{etc.})$. De cette manière on obtient en logarithmes vulgaires $\frac{1}{4}\log\dfrac{4}{c^{00}}=2.0465645311$. Prenant le tiers de ce nombre, et faisant pour abréger $\beta=0.6821881770$, on aura $\frac{1}{3}F^{1}(b)=\beta Mk$, $M$ étant mis pour le facteur 2.3025 etc.
\[ \begin{aligned} &\beta\ldots\ldots && 9.83390\ 41883\\ &M\ldots\ldots && 0.36221\ 56887\\ &k\ldots\ldots && 0.00014\ 01780\\ &\tfrac{1}{3}F^{1}(b)\ldots\ldots && 0.19626\ 00550. \end{aligned} \]
Cette valeur est exactement la même qu'on a trouvée pour le logarithme de $F^{1}(c)$; ainsi on a $F^{1}(b)=3F^{1}(c)$. Cette égalité est au moins très-approchée, puisque les formules d'où on l'a tirée pourraient, avec des tables plus étendues, donner les valeurs de $F^{1}(b)$ et $F^{1}(c)$, approchées au moins jusqu'à 15 décimales; mais il est très-vraisemblable qu'elle est rigoureuse, et c'est ce que nous aurons occasion de confirmer par une démonstration particulière.

L'égalité $F^{1}(b)=3F^{1}(c)$ étant supposée exacte d'après ce résultat, on trouve par d'autres considérations, que les fonctions $F^{1}(c)$ et $F^{1}(\sin 45^{\circ})$ ont entre elles cette relation $F^{1}(c)=\dfrac{2\sqrt[4]{3}}{3}\cos 15^{\circ}\cdot F^{1}(\sin 45^{\circ})$;

\origpage{117}
c'est ce qu'on vérifiera aisément par la valeur trouvée de $F^{1}(c)$, et par celle de $F^{1}(\sin 45^{\circ})$ tirée de la petite Table qui suit, calculée avec dix décimales.
\[ \begin{aligned} c && \log F^{1}(c) &\qquad& c && \log F^{1}(c)\\ \sin 0^{\circ} && 0.19611\ 98770 && \sin 90^{\circ} && \ldots\text{infini}\ldots\\ \sin 15^{\circ} && 0.20361\ 53664 && \sin 75^{\circ} && 0.44217\ 59938\\ \sin 30^{\circ} && 0.22679\ 32597 && \sin 60^{\circ} && 0.33375\ 26135\\ \sin 45^{\circ} && 0.26812\ 72224 && \sin 45^{\circ} && 0.26812\ 72224 \end{aligned} \]

Nous joignons ici une Table plus étendue des logarithmes des fonctions complètes $F^{1}(c)$, $E^{1}(c)$, calculées à sept décimales, pour tous les angles du module, de degré en degré, le module étant mis sous la forme $c=\sin\theta$. Cette Table sera fort utile dans la théorie des fonctions elliptiques; on l'a calculée par les formules abrégées que nous venons de rapporter, depuis $\theta=0^{\circ}$ jusqu'à $\theta=20^{\circ}$, et depuis $\theta=90^{\circ}$ jusqu'à $\theta=70^{\circ}$; pour les autres valeurs de $\theta$, on a suivi les formules des articles 65 et 77.

\origpage{118}
\[ \begin{aligned} \text{Angle du module.} && \text{Log. }F^{1} && \text{Log. }E^{1} &\qquad& \text{Angle du module.} && \text{Log. }F^{1} && \text{Log. }E^{1}\\ 0^{\circ} && 0.1961199 && 0.1961199 && 90^{\circ} && \text{Infini.} && 0.0000000\\ 1 && 0.1961530 && 0.1960868 && 89 && 0.7351923 && 0.0003264\\ 2 && 0.1962522 && 0.1959876 && 88 && 0.6760272 && 0.0012533\\ 3 && 0.1964176 && 0.1958222 && 87 && 0.6373550 && 0.0022773\\ 4 && 0.1966493 && 0.1955908 && 86 && 0.6077507 && 0.0037396\\ 5 && 0.1969474 && 0.1952932 && 85 && 0.5834041 && 0.0054653\\ 6 && 0.1973119 && 0.1949295 && 84 && 0.5625137 && 0.0074221\\ 7 && 0.1977430 && 0.1944998 && 83 && 0.5441205 && 0.0095837\\ 8 && 0.1982409 && 0.1940041 && 82 && 0.5276129 && 0.0119271\\ 9 && 0.1988057 && 0.1934423 && 81 && 0.5125914 && 0.0144321\\ 10 && 0.1994378 && 0.1928147 && 80 && 0.4987770 && 0.0170811\\ 11 && 0.2001373 && 0.1921212 && 79 && 0.4859667 && 0.0198581\\ 12 && 0.2009044 && 0.1913618 && 78 && 0.4740076 && 0.0227487\\ 13 && 0.2017396 && 0.1905367 && 77 && 0.4627820 && 0.0257397\\ 14 && 0.2026431 && 0.1896460 && 76 && 0.4521964 && 0.0288191\\ 15 && 0.2036154 && 0.1886896 && 75 && 0.4421760 && 0.0319757\\ 16 && 0.2046567 && 0.1876678 && 74 && 0.4326595 && 0.0351996\\ 17 && 0.2057675 && 0.1865805 && 73 && 0.4235961 && 0.0384811\\ 18 && 0.2069484 && 0.1854280 && 72 && 0.4149432 && 0.0418119\\ 19 && 0.2081997 && 0.1842104 && 71 && 0.4066648 && 0.0451835\\ 20 && 0.2095220 && 0.1829278 && 70 && 0.3987298 && 0.0485886\\ 21 && 0.2109158 && 0.1815802 && 69 && 0.3911076 && 0.0520197\\ 22 && 0.2123818 && 0.1801680 && 68 && 0.3837870 && 0.0554720\\ 23 && 0.2139206 && 0.1786913 && 67 && 0.3767358 && 0.0589377\\ 24 && 0.2155328 && 0.1771502 && 66 && 0.3699400 && 0.0624119\\ 25 && 0.2172193 && 0.1755451 && 65 && 0.3633838 && 0.0658895\\ 26 && 0.2189809 && 0.1738761 && 64 && 0.3570533 && 0.0693649\\ 27 && 0.2208180 && 0.1721435 && 63 && 0.3509356 && 0.0728337\\ 28 && 0.2227319 && 0.1703476 && 62 && 0.3450196 && 0.0762928\\ 29 && 0.2247233 && 0.1684886 && 61 && 0.3392950 && 0.0797376\\ 30 && 0.2267933 && 0.1665669 && 60 && 0.3337526 && 0.0831643\\ 31 && 0.2289427 && 0.1645829 && 59 && 0.3283839 && 0.0865695\\ 32 && 0.2311728 && 0.1625368 && 58 && 0.3231814 && 0.0899500\\ 33 && 0.2334846 && 0.1604292 && 57 && 0.3181381 && 0.0933029\\ 34 && 0.2358794 && 0.1582605 && 56 && 0.3132474 && 0.0966256\\ 35 && 0.2383586 && 0.1560313 && 55 && 0.3085036 && 0.0999152\\ 36 && 0.2409233 && 0.1537418 && 54 && 0.3039013 && 0.1031694\\ 37 && 0.2435751 && 0.1513928 && 53 && 0.2994353 && 0.1063861\\ 38 && 0.2463154 && 0.1489848 && 52 && 0.2951012 && 0.1095630\\ 39 && 0.2491459 && 0.1465186 && 51 && 0.2908945 && 0.1126982\\ 40 && 0.2520685 && 0.1439949 && 50 && 0.2868114 && 0.1157899\\ 41 && 0.2550848 && 0.1414143 && 49 && 0.2828480 && 0.1188364\\ 42 && 0.2581965 && 0.1387776 && 48 && 0.2790011 && 0.1218362\\ 43 && 0.2614061 && 0.1360858 && 47 && 0.2752673 && 0.1247878\\ 44 && 0.2647155 && 0.1333399 && 46 && 0.2716436 && 0.1276898\\ 45 && 0.2681272 && 0.1305409 && 45 && 0.2681272 && 0.1305409 \end{aligned} \]
\ednote{Three entries are printed as given but differ from the exact values by more than rounding: Log.\ $F^{1}$ at $85^{\circ}$ is printed 0.5834041 (exact 0.5833963), Log.\ $F^{1}$ at $69^{\circ}$ is printed 0.3911076 (exact 0.3911115), and Log.\ $E^{1}$ at $88^{\circ}$ is printed 0.0012533 (exact 0.0011208). The other entries agree with the exact values to within a few units of the seventh decimal.}

\origpage{119}
\section*{Méthode d'approximation appliquée aux fonctions elliptiques de la troisième espèce.}

\subsection*{(85).}

\ednote{The article number is printed (85), and the next article is also numbered (85); art.\ 88 cites this equation as n° 84, so (84) is meant. An apparent misprint.} Au lieu de considérer la simple formule $\Pi=\displaystyle\int\frac{d\varphi}{(1+n\sin^{2}\varphi)\Delta}$ qui appartient à la troisième espèce, nous considérerons plus généralement la formule
\[ H=\int\left(A+\frac{B\sin^{2}\varphi}{1+n\sin^{2}\varphi}\right)\frac{d\varphi}{\Delta}, \]
qui est composée d'une fonction de la troisième espèce et d'une de la première, affectées de coefficiens constans. Si on appliquait le calcul à la simple fonction $\Pi$, la première transformée contiendrait une partie affectée de la fonction de première espèce $F(c^{0},\varphi^{0})$; c'est pourquoi il convient, pour l'uniformité des résultats, d'admettre les deux sortes de fonctions dans la formule primitive, ainsi que nous l'avons déjà fait pour les fonctions $G$ (art.\ 76).

Cela posé, si on fait $\sin^{2}\varphi=\frac{1}{2}(1+c^{0}\sin^{2}\varphi^{0}-\Delta^{0}\cos\varphi^{0})$, et $\frac{d\varphi}{\Delta}=\frac{1+c^{0}}{2}\cdot\frac{d\varphi^{0}}{\Delta^{0}}$, on aura pour première transformée
\[ H=\frac{1+c^{0}}{2}\left[H^{0}-\frac{\tfrac{1}{2}B}{1+n}\int\frac{d\varphi^{0}\cos\varphi^{0}}{1+n^{0}\sin^{2}\varphi^{0}}\right], \]
les valeurs de $H^{0}$ et de ses coefficiens étant
\[ \begin{aligned} H^{0}&=\int\left(A^{0}+\frac{B^{0}\sin^{2}\varphi^{0}}{1+n^{0}\sin^{2}\varphi^{0}}\right)\frac{d\varphi^{0}}{\Delta^{0}}.\\ n^{0}&=\frac{n(n+c^{2})}{(1+b)^{2}(1+n)}\\ A^{0}&=A+\frac{\tfrac{1}{2}B}{1+n}\\ B^{0}&=B\cdot\frac{c^{2}+2n+n^{2}}{2(1+b)^{2}(1+n)^{2}}. \end{aligned} \]
Ainsi la fonction $H(n,c,\varphi)$ est ramenée à une fonction semblable $H(n^{0},c^{0},\varphi^{0})$, dans laquelle les trois élémens $n$, $c$, $\varphi$ ont subi des changemens propres à faciliter les approximations.

La même loi aura lieu dans les transformées ultérieures; de sorte

\origpage{120}
qu'en faisant
\[ \begin{aligned} H^{00}&=\int\left(A^{00}+\frac{B^{00}\sin^{2}\varphi^{00}}{1+n^{00}\sin^{2}\varphi^{00}}\right)\frac{d\varphi^{00}}{\Delta^{00}}\\ n^{00}&=\frac{n^{0}(n^{0}+c^{02})}{(1+b^{0})^{2}(1+n^{0})}\\ A^{00}&=A^{0}+\frac{\tfrac{1}{2}B^{0}}{1+n^{0}}\\ B^{00}&=B^{0}\cdot\frac{c^{02}+2n^{0}+n^{02}}{2(1+b^{0})^{2}(1+n^{0})^{2}}; \end{aligned} \]
on aura la seconde transformée
\[ H^{0}=\frac{1+c^{00}}{2}\left[H^{00}-\frac{\tfrac{1}{2}B^{0}}{1+n^{0}}\int\frac{d\varphi^{00}\cos\varphi^{00}}{1+n^{00}\sin^{2}\varphi^{00}}\right]. \]
On peut continuer indéfiniment ces transformations, d'où résulteront une suite infinie de fonctions $H(n,c,\varphi)$, $H(n^{0},c^{0},\varphi^{0})$, $H(n^{00},c^{00},\varphi^{00})$, etc., telles, qu'une seule étant connue, on pourra déterminer toutes les autres.

\subsection*{(85).}

Examinons d'abord avec quelque détail la loi de progression des paramètres $n^{0}$, $n^{00}$, etc. Si on fait $n=mc$, et semblablement $n^{0}=m^{0}c^{0}$, l'équation qui détermine $n^{0}$ donnera $m^{0}=\frac{m(m+c)}{1+mc}$. Mais on peut toujours supposer le paramètre $n<c$, ce qui donne $m<1$; on aura donc aussi $m^{0}<1$, et même $m^{0}<m$; car de la valeur précédente on déduit $1+\frac{m^{0}}{m}=\frac{(1+c)(1+m)}{1+cm}$, $1-\frac{m^{0}}{m}=\frac{(1-c)(1-m)}{1+cm}$, ce qui prouve que $\frac{m^{0}}{m}$, positif ou négatif, est toujours plus petit que l'unité. Donc les termes $m$, $m^{0}$, $m^{00}$, etc. sont tous plus petits que l'unité et forment une suite décroissante; et par conséquent la suite des paramètres $n$, $n^{0}$, $n^{00}$, etc. décroît plus rapidement que la suite des modules $c$, $c^{0}$, $c^{00}$, etc., chaque terme de la première étant plus petit que le terme correspondant de la seconde.

A l'égard des signes de ces paramètres, il y a deux cas à considérer. 1°. Si $n$ est positif ou si $n$ est négatif et plus grand que $c^{2}$, ensorte qu'il soit compris dans la forme $-1+b^{2}\sin^{2}\theta$, alors $n^{0}$ sera positif, et il en sera de même de tous les paramètres suivans $n^{00}$, $n^{000}$, etc. Dans ce premier cas se trouve compris celui de $n=\pm c$

\origpage{121}
qui donnera $n^{0}=c^{0}$, $n^{00}=c^{00}$, etc.; mais ce cas est inutile à considérer, parce que d'après l'article 46, la fonction $H$ se réduit alors à la première espèce.

2°. Si $n$ est négatif et plus petit que $c^{2}$, ensorte qu'on ait $n=-c^{2}\sin^{2}\theta$, on pourra faire semblablement $n^{0}=-c^{02}\sin^{2}\theta^{0}$, et l'angle $\theta^{0}$ se déterminera par l'équation
\[ \operatorname{tang}(\theta^{0}-\theta)=b\operatorname{tang}\theta, \]
ce qui est la même loi suivant laquelle l'amplitude $\varphi^{0}$ se déduit de $\varphi$.

Calculant donc successivement $\theta^{0}$, $\theta^{00}$, $\theta^{000}$, etc. d'après cette loi, on formera la suite des paramètres $n^{0}=-c^{02}\sin^{2}\theta^{0}$, $n^{00}=-c^{002}\sin^{2}\theta^{00}$, etc., qui décroît, comme on voit, avec encore plus de rapidité que dans le premier cas.

Connaissant ainsi la loi de progression des paramètres $n$, venons à celle des coefficiens $A$ et $B$ dans les transformées successives. Soit pour abréger,
\[ k=\frac{c^{2}+2n+n^{2}}{(1+b)^{2}(1+n)^{2}}=\frac{(1+n)^{2}-b^{2}}{(1+b)^{2}(1+n)^{2}}; \]
et soit désigné par $k^{0}$ une quantité composée de $c^{0}$, $b^{0}$, $n^{0}$, comme $k$ l'est de $c$, $b$, $n$, et ainsi de suite, on aura
\[ B^{0}=\tfrac{1}{2}kB,\quad B^{00}=\tfrac{1}{4}kk^{0}B,\quad B^{000}=\tfrac{1}{8}kk^{0}k^{00}B,\quad\text{etc.} \]
Mais puisqu'après un certain nombre $\mu$ de termes, les $c^{\mu}$ et $n^{\mu}$ pris dans les suites $c^{0}$, $c^{00}$, etc., $n^{0}$, $n^{00}$, etc., sont assez petits pour être regardés comme nuls, il est clair qu'il en sera de même du terme correspondant $k^{\mu}$ de la suite $k$, $k^{0}$, $k^{00}$, etc., et qu'ainsi on peut faire en toute sureté $B^{\mu}=0$. Quant à la valeur de $A^{\mu}$, elle sera donnée par la suite
\[ A^{\mu}=A+\frac{\tfrac{1}{2}B}{1+n}+\frac{\tfrac{1}{4}kB}{1+n^{0}}+\frac{\tfrac{1}{8}kk^{0}B}{1+n^{00}}+\text{etc.} \]
qu'il faudra prolonger jusqu'au terme qui contiendrait $k^{\mu}$ exclusivement.

Appelons toujours $\Phi$ la limite des angles $\varphi$, $\frac{\varphi^{0}}{2}$, $\frac{\varphi^{00}}{4}$, etc., cette limite étant censée atteinte après le nombre de termes $\mu$, on aura $\varphi^{\mu}=2^{\mu}\Phi$, et parce que $c^{\mu}$ est de l'ordre des quantités négligeables, on aura $H^{\mu}=2^{\mu}A^{\mu}\Phi$.

\origpage{122}
\subsection*{(86).}

Il est facile maintenant d'avoir le résultat des transformées successives jusqu'à la $\mu^{\text{ième}}$ inclusivement. Soit toujours
\[ K^{\mu}=(1+c^{0})(1+c^{00})\ldots(1+c^{\mu}),\quad\text{ou}\quad K^{\mu}=\frac{2\sqrt{c^{0}}}{c}\cdot\frac{2\sqrt{c^{00}}}{c^{0}}\cdot\frac{2\sqrt{c^{000}}}{c^{00}}\cdot\text{etc.}, \]
la valeur de la fonction $H$ sera donnée en général par la formule
\[ \begin{aligned} H=K^{\mu}A^{\mu}\Phi&-\frac{\sqrt{c^{0}}}{c}\cdot\frac{\tfrac{1}{2}B}{1+n}\cdot\frac{\operatorname{arc}\operatorname{tang}(\sqrt{n^{0}}\sin\varphi^{0})}{\sqrt{n^{0}}}\\ &-\frac{\sqrt{c^{0}}}{c}\cdot\frac{\sqrt{c^{00}}}{c^{0}}\cdot\frac{\tfrac{1}{4}kB}{1+n^{0}}\cdot\frac{\operatorname{arc}\operatorname{tang}(\sqrt{n^{00}}\sin\varphi^{00})}{\sqrt{n^{00}}}\\ &-\frac{\sqrt{c^{0}}}{c}\cdot\frac{\sqrt{c^{00}}}{c^{0}}\cdot\frac{\sqrt{c^{000}}}{c^{00}}\cdot\frac{\tfrac{1}{8}kk^{0}B}{1+n^{00}}\cdot\frac{\operatorname{arc}\operatorname{tang}(\sqrt{n^{000}}\sin\varphi^{000})}{\sqrt{n^{000}}}\\ &-\text{etc.} \end{aligned} \]
Cette suite étant prolongée jusqu'au terme dont le rang est $\mu$ inclusivement, l'erreur de la formule, mesurée par les termes suivans, ne pourra être que de l'ordre $c^{\mu+1}$; de sorte que si on ne veut pousser la précision que jusqu'aux quantités de l'ordre $c^{\mu}$, on aura un terme de moins à calculer, tant dans la suite précédente que dans la valeur de $K$ et dans celle de $A^{\mu}$.

Au reste les arcs de cercle qui entrent dans la formule générale se changeraient en logarithmes, si $n$ était de la forme $-c^{2}\sin^{2}\theta$; mais ce changement n'est sujet à aucune difficulté.

Lorsque $\varphi=\frac{1}{2}\pi$ ou $=$ un multiple de $\frac{1}{2}\pi$, tous les arcs de cercle disparaissent, et on a simplement $H=KA^{\mu}\Phi$. Donc la valeur de la fonction complète est
\[ H^{1}=KA^{\mu}\cdot\tfrac{1}{2}\pi. \]
Remarquez que la valeur de $H$ exprimée par les transformées successives, se terminerait d'elle-même, si on avait l'une des égalités $n=-1+b$, $n^{0}=-1+b^{0}$, $n^{00}=-1+b^{00}$, etc.: alors la fonction $H$ se ramènerait indéfiniment à la première espèce.

Pour qu'on ait en général $-1+b^{\mu}=n^{\mu}=-(c^{\mu}\sin\theta^{\mu})^{2}$, il faut que $\cot\theta^{\mu}=\sqrt{b^{\mu}}$, ou qu'on ait $F(c^{\mu},\theta^{\mu})=\frac{1}{2}F^{1}(c^{\mu})$. Mais les paramètres successifs $n=-c^{2}\sin^{2}\theta$, $n^{0}=-c^{02}\sin^{2}\theta^{0}$, $n^{00}=-c^{002}\sin^{2}\theta^{00}$, etc. étant formés d'après la même loi qui lie entre elles les amplitudes $\theta$, $\theta^{0}$, $\theta^{00}$, etc., on a $F(c,\theta)=\frac{1+c^{0}}{2}F(c^{0},\theta^{0})$, $F(c^{0},\theta^{0})=\frac{1+c^{00}}{2}F(c^{00},\theta^{00})$, etc.;

\origpage{123}
et en général $F(c,\theta)=\frac{(1+c^{0})(1+c^{00})\ldots(1+c^{\mu})}{2^{\mu}}F(c^{\mu},\theta^{\mu})=\frac{1}{2^{\mu+1}}F^{1}(c)$. Les cas qui donnent lieu à la réduction de la fonction $H$, sont donc tous ceux où le paramètre $n$, de forme $-c^{2}\sin^{2}\theta$, est tel qu'on a $F(c,\theta)=\frac{1}{2^{\mu+1}}F^{1}(c)$, $\mu$ étant un entier. La même réduction aurait encore lieu si on avait $F(c,\theta)=\frac{2\nu+1}{2^{\mu+1}}F^{1}(c)$, $2\nu+1$ étant un nombre impair quelconque.

Les formules d'approximation que nous venons de donner pour déterminer la valeur des fonctions $H$ et $H^{1}$, peuvent s'appliquer à tous les cas, et le plus souvent il ne faudra calculer que fort peu de termes pour obtenir un grand degré de précision. Cependant si la différence $1-c$ était tellement petite qu'il fallût prolonger assez loin la suite $c^{0}$, $c^{00}$, $c^{000}$, etc. pour parvenir à un terme négligeable $c^{\mu}$, il pourrait être préférable de suivre la seconde méthode, c'est-à-dire, de faire les transformations dans un ordre inverse, ainsi que nous l'avons pratiqué en pareil cas relativement aux fonctions de la première et de la seconde espèce.

\subsection*{(87).}

Et d'abord si la différence $1-c$ est déjà assez petite pour qu'elle puisse être négligée, on pourra intégrer immédiatement la formule $H=\displaystyle\int\left(A+\frac{B\sin^{2}\varphi}{1+n\sin^{2}\varphi}\right)\frac{d\varphi}{\Delta}$, en mettant $\cos\varphi$ au lieu de $\Delta$, ce qui donnera
\[ H=\left(A+\frac{B}{1+n}\right)\log\operatorname{tang}(45^{\circ}+\tfrac{1}{2}\varphi)-\frac{B}{1+n}\cdot\frac{\operatorname{arc}\operatorname{tang}(\sqrt{n}\sin\varphi)}{\sqrt{n}}. \]
Cette intégrale est rigoureuse lorsque $c=1$; mais s'il y a une différence, quelque petite qu'elle soit, entre 1 et $c$, elle ne pourra s'appliquer sans erreur à des valeurs de $\varphi$ plus grandes qu'une certaine limite. En effet, la quantité $\Delta$ qui en général est $\sqrt{(\cos^{2}\varphi+b^{2}\sin^{2}\varphi)}$, ne se réduit à $\cos\varphi$ que lorsque $\cot\varphi$ est censé beaucoup plus grand que $b$. C'est pourquoi il convient de chercher une formule qui donne la valeur de $H$ pour toute valeur de $\varphi$, dans l'hypothèse seulement qu'on puisse négliger les termes affectés du facteur $b$ élevé au quarré ou à une puissance supérieure.

Pour cet effet nous supposerons qu'il a été fait dans la fonction $H$

\origpage{124}
une préparation relative au paramètre. Si $n$ est négatif, nous le supposons toujours ou de la forme $-c^{2}\sin^{2}\theta$, ou de la forme $-1+b^{2}\sin^{2}\theta$; si ce dernier cas a lieu, on changera la fonction $H$ en une autre où le paramètre soit positif, ce qui se fera par la formule du n° 51.

Cela posé, pour avoir la valeur de la fonction $H$, je la mets sous la forme
\[ H=\left(A+\frac{B}{1+n}\right)F(c,\varphi)-\frac{B}{1+n}P, \]
et il restera à déterminer $P$ par la formule
\[ P=\int\frac{d\varphi\cos^{2}\varphi}{(1+n\sin^{2}\varphi)\Delta}. \]
J'observe ensuite qu'on a $\frac{\cos\varphi}{\Delta}=1-\frac{b^{2}\sin^{2}\varphi}{\Delta(\cos\varphi+\Delta)}$; donc si on fait
\[ Q=\int\frac{2d\varphi\cos\varphi\sin^{2}\varphi}{(1+n\sin^{2}\varphi)\Delta(\cos\varphi+\Delta)}; \]
on aura
\[ P=\frac{\operatorname{arc}\operatorname{tang}(\sqrt{n}\sin\varphi)}{\sqrt{n}}-\tfrac{1}{2}b^{2}Q. \]
Tout se réduit donc à trouver une valeur approchée de l'intégrale $Q$. Pour cet effet, j'observe qu'on a $\cos\varphi+\Delta<2\Delta$, et $\cos\varphi+\Delta>2\cos\varphi$; donc si on fait
\[ \begin{aligned} Q'&=\int\frac{d\varphi\sin^{2}\varphi}{(1+n\sin^{2}\varphi)\Delta}\\ Q''&=\int\frac{d\varphi\cos\varphi\sin^{2}\varphi}{(1+n\sin^{2}\varphi)(1-c^{2}\sin^{2}\varphi)}, \end{aligned} \]
on aura $Q<Q'$ et $Q>Q''$. Or en premier lieu, la combinaison des intégrales $P$ et $Q'$ donne $(1+n)Q'+P=\displaystyle\int\frac{d\varphi}{\Delta}$, d'où résulte $Q'=\frac{F(c,\varphi)-P}{1+n}$, donc on a
\[ P>\frac{\operatorname{arc}\operatorname{tang}(\sqrt{n}\sin\varphi)}{\sqrt{n}}-\frac{\tfrac{1}{2}b^{2}}{1+n}F(c,\varphi)+\frac{\tfrac{1}{2}b^{2}}{1+n}P, \]
ou
\[ P>\left(1+\frac{\tfrac{1}{2}b^{2}}{1+n}\right)\frac{\operatorname{arc}\operatorname{tang}(\sqrt{n}\sin\varphi)}{\sqrt{n}}-\frac{\tfrac{1}{2}b^{2}}{1+n}F(c,\varphi) \]
On trouve ensuite par l'intégration directe,
\[ (c^{2}+n)Q''=\frac{1}{2c}\log\left(\frac{1+c\sin\varphi}{1-c\sin\varphi}\right)-\frac{1}{\sqrt{n}}\operatorname{arc}\operatorname{tang}(\sqrt{n}\sin\varphi); \]

\origpage{125}
donc on a
\[ P<\left(1+\frac{\tfrac{1}{2}b^{2}}{1+n}\right)\frac{\operatorname{arc}\operatorname{tang}(\sqrt{n}\sin\varphi)}{\sqrt{n}}-\frac{\tfrac{1}{4}b^{2}}{1+n}\log\left(\frac{1+c\sin\varphi}{1-c\sin\varphi}\right). \]
La plus grande différence entre ces deux limites de $P$ a lieu lorsque $\varphi=90^{\circ}$, et comme on a dans ce cas $F^{1}(c)=\log\frac{4}{b}$, $\log\left(\frac{1+c}{1-c}\right)=2\log\frac{2}{b}$, cette différence $=\frac{b^{2}}{1+n}\cdot\frac{1}{2}\log 2$. Elle est, comme on voit, de l'ordre des quantités qu'on veut négliger. Ainsi on a avec une exactitude suffisante,
\[ P=\frac{1}{\sqrt{n}}\operatorname{arc}\operatorname{tang}(\sqrt{n}\sin\varphi)-\frac{\tfrac{1}{2}b^{2}}{1+n}F(c,\varphi), \]
ce qui donne la fonction cherchée
\[ H=\left(A+\frac{B}{c+n}\right)F(c,\varphi)-\frac{B}{1+n}\cdot\frac{\operatorname{arc}\operatorname{tang}(\sqrt{n}\sin\varphi)}{\sqrt{n}}; \]
et il en résulte pour la fonction complète,
\[ H^{1}=\left(A+\frac{B}{1+n}\right)(1+\tfrac{1}{4}b^{2})\log\frac{4}{b}-\frac{B}{1+n}\cdot\frac{\operatorname{arc}\operatorname{tang}(\sqrt{n})}{\sqrt{n}}. \]
Ces valeurs de $H$ et de $H^{1}$ sont celles qu'il convient d'employer si les quantités de l'ordre $b^{2}$ sont négligeables, et alors il n'y a pas lieu de recourir aux transformations; mais il convient d'examiner particulièrement le cas où $n$ est négatif.

Le cas de $n=-1+b^{2}\sin^{2}\theta$, où $n+1$ serait de l'ordre $b^{2}$, ne permet pas qu'on lui applique avec succès les formules précédentes, c'est pourquoi nous l'avons évité par une transformation qui rend $n$ positif.

Le cas de $n=-c^{2}\sin^{2}\theta$ ne souffre aucune difficulté, si $\theta$ n'est pas trop près de $90^{\circ}$, parce qu'alors $\frac{b^{2}}{1+n}$ reste une quantité très-petite; seulement dans ce cas, il faut changer l'expression $\frac{\operatorname{arc}\operatorname{tang}(\sqrt{n}\sin\varphi)}{\sqrt{n}}$ en celle-ci $\frac{1}{2c\sin\theta}\log\left(\frac{1+c\sin\theta\sin\varphi}{1-c\sin\theta\sin\varphi}\right)$. Mais si en même temps $\theta$ est très-près de $90^{\circ}$, le dénominateur $1+n$ équivalant à $b^{2}+c^{2}\cos^{2}\theta$, devient très-petit; cependant tant que $\cos\theta$ sera beaucoup plus grand que $b$, la fraction $\frac{b^{2}}{1+n}$ ou $\frac{b^{2}}{b^{2}+c^{2}\cos^{2}\theta}$ demeurera très-petite, et la

\origpage{126}
méthode précédente sera applicable; mais si $\cos\theta$ est une quantité très-petite de l'ordre $b$ ou d'un ordre supérieur, la fraction $\frac{b^{2}}{1+n}$ cesse d'être négligeable. Pour obvier à cet inconvénient, il faut recourir à la formule du n° 52, au moyen de laquelle toute fonction de troisième espèce dont le paramètre $n=-c^{2}\sin^{2}\theta$, peut être transformée en une autre fonction dont le paramètre $n'=-c^{2}\sin^{2}\lambda$, pourvu qu'entre les angles $\theta$ et $\lambda$ on ait la relation $b\operatorname{tang}\theta\operatorname{tang}\lambda=1$. Alors on a $\frac{b^{2}}{1+n}\cdot\frac{b^{2}}{1+n'}=b^{2}$, donc des deux facteurs $\frac{b^{2}}{1+n}$, $\frac{b^{2}}{1+n'}$, il y en aura toujours un plus petit que $b$; d'où il suit qu'à l'aide d'une transformation, si elle est nécessaire, l'erreur de la formule n'excédera pas les quantités de l'ordre $b$; et parce que dans le cas le plus défavorable, qui est celui de $1+n=1+n'=b$, la valeur de $H^{1}$ devient très-grande et se rapporte à l'ordre $b^{-1}$, ou plus exactement à l'ordre $b^{-1}\log\left(\frac{4}{b}\right)$ qui est très-peu différent, il s'ensuit que l'erreur absolue étant de l'ordre $b$, l'erreur relative est en effet de l'ordre $b^{2}$, ainsi que dans les cas où $n$ est positif.

Nous sommes entrés dans ces détails pour prouver que lorsque $1-c$ est assez petit pour être négligeable, les formules trouvées pour $H$ et $H^{1}$ peuvent donner une approximation suffisante, sans exiger de transformations autres que celles qui ont été indiquées par les articles 51 et 52. Mais la méthode que nous allons exposer a l'avantage de donner une approximation indéfinie, qui s'applique généralement à tous les cas.

\subsection*{(88).}

L'équation trouvée n° 84 entre $H$ et $H^{0}$, s'applique aux deux fonctions consécutives $H'$ et $H$, et donne la transformée
\[ H=\frac{2}{1+c}H'+\frac{\tfrac{1}{2}B'}{1+n'}\cdot\int\frac{d\varphi\cos\varphi}{1+n\sin^{2}\varphi}; \]
on aura en même temps
\[ H'=\int\left(A'+\frac{B'\sin^{2}\varphi'}{1+n'\sin^{2}\varphi'}\right)\frac{d\varphi'}{\Delta'}, \]
et les coefficiens $n'$, $A'$, $B'$ devront être déduits des équations
\[ n=\frac{n'(n'+c'^{2})}{(1+b')^{2}(1+n')},\quad B=\frac{B'}{2}\cdot\frac{(1+n')^{2}-b'^{2}}{(1+b')^{2}(1+n')^{2}},\quad A=A'+\frac{\tfrac{1}{2}B'}{1+n'}. \]

\origpage{127}
On sait comment se calculent $c'$, $b'$, $\varphi'$, au moyen de $c$ et $\varphi$; il ne s'agit donc que d'examiner la loi de formation des quantités $n'$, $A'$, $B'$. Et d'abord pour déduire $n'$ de $n$, on a la formule
\[ n'=\frac{2}{(1+c)^{2}}[n-c+\sqrt{(1+n)(c^{2}+n)}]; \]
d'où l'on voit qu'il faut faire abstraction du cas où l'on aurait $n=-1+b^{2}\sin^{2}\theta$; car si ce cas avait lieu, le paramètre $n'$ deviendrait imaginaire, et on tomberait ainsi dans une difficulté plus grande que celle qu'on veut résoudre. Heureusement il n'y a aucun inconvénient à faire cette exception, puisque par la formule du n° 51, toute fonction proposée de troisième espèce dont le paramètre est de la forme $-1+b^{2}\sin^{2}\theta$, peut être transformée en une autre dont le paramètre soit positif. Nous pouvons donc nous borner à considérer deux cas principaux, celui où $n$ est positif et celui où $n$ est de la forme $-c^{2}\sin^{2}\theta$. Dans ces deux cas, on peut s'assurer que la valeur de $n'$ sera toujours de même forme que celle de $n$, ainsi rien ne pourra arrêter le calcul des transformées successives.

\subsection*{(89).}

Soit donc en premier lieu $n=-c^{2}\sin^{2}\theta$, et semblablement $n'=-c'^{2}\sin^{2}\theta'$, on aura par la formule précédente,
\[ \sin^{2}\theta'=\tfrac{1}{2}+\tfrac{1}{2}c\sin^{2}\theta-\tfrac{1}{2}\cos\theta\Delta(\theta). \]
C'est la même loi suivant laquelle l'amplitude $\varphi'$ se déduit de $\varphi$; elle peut être représentée plus simplement par l'équation $\sin(2\theta'-\theta)=c\sin\theta$. En vertu de cette loi, les angles $\theta$, $\theta'$, $\theta''$, etc. forment une suite décroissante qui s'approche sans cesse d'une limite $\Theta$, et qui l'atteint sensiblement après un petit nombre de termes. La suite des paramètres $n$, $n'$, $n''$, etc. converge donc de même vers la limite $-\sin^{2}\Theta$, puisque d'ailleurs la suite $c$, $c'$, $c''$, etc. a pour limite l'unité.

En second lieu, si $n$ est positif, soit $1+n=\frac{b^{2}}{\sin^{2}\lambda}$, $\lambda$ étant un angle très-petit déterminé par la valeur $\sin\lambda=\frac{b}{\sqrt{(1+n)}}$, de sorte que si on faisait $n=\cot^{2}\theta$, on aurait $\sin\lambda=b\sin\theta$. Cela posé, la valeur de $n'$ sera donnée par l'équation
\[ 1+n'=b'\cot^{2}\tfrac{1}{2}\lambda; \]

\origpage{128}
et si l'on fait semblablement $1+n'=\frac{b'^{2}}{\sin^{2}\lambda'}$, l'angle $\lambda'$ se déduira de $\lambda$ par la formule très-simple
\[ \sin\lambda'=\sqrt{b'}\cdot\operatorname{tang}\tfrac{1}{2}\lambda. \]
On déduira semblablement $\lambda''$ de $\lambda'$, $\lambda'''$ de $\lambda''$, etc., de sorte qu'on formera la suite des paramètres $n$, $n'$, $n''$, etc. d'après les équations $1+n'=\frac{b'^{2}}{\sin^{2}\lambda'}$, $1+n''=\frac{b''^{2}}{\sin^{2}\lambda''}$, etc.

Remarquons qu'à partir de l'angle $\lambda$ qui est très-petit, les termes de la série $\lambda$, $\lambda'$, $\lambda''$, etc. diminuent continuellement, mais de manière que les rapports $\frac{\sin\lambda}{b}$, $\frac{\sin\lambda'}{b'}$, $\frac{\sin\lambda''}{b''}$, etc. convergent vers une quantité constante qui sera leur limite. En effet, des valeurs précédentes on déduit $\frac{\sin\lambda'}{b'}=\frac{\sin\lambda}{b}\cdot\frac{1+c}{1+\cos\lambda}$; on aurait de même $\frac{\sin\lambda''}{b''}=\frac{\sin\lambda'}{b'}\cdot\frac{1+c'}{1+\cos\lambda'}$, etc. Donc après un petit nombre de termes $\frac{\sin\lambda^{\mu}}{b^{\mu}}$ sera constant, et par conséquent aussi $1+n^{\mu}$; d'où l'on voit que la suite des paramètres $n$, $n'$, $n''$, etc., converge, comme dans le premier cas, vers une limite qu'elle atteint sensiblement au bout de quelques termes.

Le moyen très-simple que nous venons d'indiquer pour calculer la suite des paramètres $n$, $n'$, $n''$, etc., lorsque le premier $n$ est positif, peut aussi s'appliquer sans difficulté au cas où $n$ est de la forme $-c^{2}\sin^{2}\theta$; de sorte que nous pourrons supposer qu'il est employé généralement dans tous les cas; et d'après cette supposition, nous allons continuer les calculs nécessaires pour parvenir à l'expression de la valeur de $H$.

\subsection*{(90).}

Il faut d'abord avoir la loi des coefficiens $B^{\mu}$, $A^{\mu}$; pour cela, si l'on fait
\[ h'=\frac{(1+b')^{2}(1+n')^{2}}{(1+n')^{2}-b'^{2}},\quad h''=\frac{(1+b'')^{2}(1+n'')^{2}}{(1+n'')^{2}-b''^{2}},\quad\text{etc.}, \]
on aura successivement
\[ B'=2h'B,\quad B''=2h''B'=4h'h''B,\quad B'''=8h'h''h'''B,\quad\text{etc.} \]

\origpage{129}
Mais des équations $1+n'=b'\cot^{2}\frac{1}{2}\lambda$, $1+n=\frac{b^{2}}{\sin^{2}\lambda}$, $b^{2}=\frac{4b'}{(1+b')^{2}}$, on tire
\[ h'=\frac{1+n'}{1+n}\cdot\frac{1}{\cos\lambda}; \]
on aura semblablement $h''=\frac{1+n''}{1+n'}\cdot\frac{1}{\cos\lambda'}$, $h'''=\frac{1+n'''}{1+n''}\cdot\frac{1}{\cos\lambda''}$, etc.; donc
\[ B'=\frac{2B(1+n')}{(1+n)\cos\lambda},\quad B''=\frac{4B(1+n'')}{(1+n)\cos\lambda\cos\lambda'},\quad B'''=\frac{8B(1+n''')}{(1+n)\cos\lambda\cos\lambda'\cos\lambda''},\quad\text{etc.} \]
et en général $B^{\mu}=\frac{2^{\mu}B(1+n^{\mu})}{(1+n)\cos\lambda\cos\lambda'\ldots\cos\lambda^{\mu-1}}$.

Enfin des équations $A'=A-\frac{\frac{1}{2}B'}{1+n'}$, $A''=A'-\frac{\frac{1}{2}B''}{1+n''}$, etc., on déduit généralement
\[ A^{\mu}=A-\frac{\tfrac{1}{2}B'}{1+n'}-\frac{\tfrac{1}{2}B''}{1+n''}-\ldots-\frac{\tfrac{1}{2}B^{\mu}}{1+n^{\mu}}; \]
ou, en faisant les substitutions,
\[ A^{\mu}=A-\frac{B}{1+n}\left(\frac{1}{\cos\lambda}+\frac{2}{\cos\lambda\cos\lambda'}+\frac{4}{\cos\lambda\cos\lambda'\cos\lambda''}\ldots+\frac{2^{\mu-1}}{\cos\lambda\cos\lambda'\ldots\cos\lambda^{\mu-1}}\right). \]
Ces valeurs feront connaître celle du coefficient $A^{\mu}+\frac{B^{\mu}}{1+n^{\mu}}$ qui entre dans l'expression de $H^{\mu}$. Soit ce coefficient $=L^{\mu}$, on aura
\[ \begin{aligned} L^{\mu}=A&+\frac{2^{\mu}B}{(1+n)\cos\lambda\cos\lambda'\ldots\cos\lambda^{\mu-1}}\\ &-\frac{B}{1+n}\left(\frac{1}{\cos\lambda}+\frac{2}{\cos\lambda\cos\lambda'}+\frac{4}{\cos\lambda\cos\lambda'\cos\lambda''}\ldots+\frac{2^{\mu-1}}{\cos\lambda\cos\lambda'\ldots\cos\lambda^{\mu-1}}\right); \end{aligned} \]
Mais comme la suite comprise dans cette formule est divergente, il conviendra de lui donner une autre forme. Pour cet effet j'observe que la formule précédente donne
\[ L^{\mu+1}-L^{\mu}=\frac{2^{\mu}B(1-\cos\lambda^{\mu})}{(1+n)\cos\lambda\cos\lambda'\ldots\cos\lambda^{\mu}}. \]
Or on a vu qu'à mesure que $\mu$ augmente, la quantité $\frac{\sin\lambda^{\mu}}{b^{\mu}}$ s'approche

\origpage{130}
de plus en plus d'une certaine limite, et qu'ainsi $\sin\lambda^{\mu}$ est en général de l'ordre $b^{\mu}$; donc $1-\cos\lambda^{\mu}$ est de l'ordre $(b^{\mu})^{2}$, ou $b^{\mu+1}$; donc la différence $L^{\mu+1}-L^{\mu}$ est une quantité très-petite de l'ordre $2^{\mu}b^{\mu+1}$, intermédiaire entre $b^{\mu}$ et $b^{\mu+1}$, mais plus proche de ce dernier. On aura donc pour l'expression de $L^{\mu}$ cette suite fort convergente
\[ \begin{aligned} L^{\mu}=A+\frac{B}{1+n}&+\frac{B(1-\cos\lambda)}{(1+n)\cos\lambda}+\frac{2B(1-\cos\lambda')}{(1+n)\cos\lambda\cos\lambda'}\\ &+\frac{4B(1-\cos\lambda'')}{(1+n)\cos\lambda\cos\lambda'\cos\lambda''}+\text{etc.}, \end{aligned} \]
laquelle devra être prolongée jusqu'au terme \ldots\ $\frac{2^{\mu-1}B(1-\cos\lambda^{\mu-1})}{(1+n)\cos\lambda\cos\lambda'\cos\lambda''\ldots\cos\lambda^{\mu-1}}$ inclusivement, si on ne veut négliger que les quantités de l'ordre $2^{\mu}b^{\mu+1}$.

\subsection*{(91).}

Cela posé, les transformées successives étant
\[ \begin{aligned} H&=(1+b')H'+\frac{\tfrac{1}{2}B'}{1+n'}\cdot\frac{\operatorname{arc}\operatorname{tang}(\sqrt{n}\sin\varphi)}{\sqrt{n}}\\ H'&=(1+b'')H''+\frac{\tfrac{1}{2}B''}{1+n''}\cdot\frac{\operatorname{arc}\operatorname{tang}(\sqrt{n'}\sin\varphi')}{\sqrt{n'}}\\ &\text{etc.}, \end{aligned} \]
on pourra exprimer généralement la fonction $H$ par le moyen de $H^{\mu}$; et comme les séries sont prolongées jusqu'à un nombre de termes $\mu$ suffisant pour que la valeur de $H^{\mu}$ soit donnée avec toute l'exactitude nécessaire par la formule
\[ H^{\mu}=\left(A^{\mu}+\frac{B^{\mu}}{1+n^{\mu}}\right)\log\operatorname{tang}(45^{\circ}+\tfrac{1}{2}\varphi^{\mu})-\frac{B^{\mu}}{1+n^{\mu}}\cdot\frac{\operatorname{arc}\operatorname{tang}(\sqrt{n^{\mu}}\sin\varphi^{\mu})}{\sqrt{n^{\mu}}}, \]
il résultera des substitutions cette valeur générale de $H$, où l'on a fait pour abréger, $I^{\mu}=(1+b')(1+b'')(1+b''')\ldots(1+b^{\mu})$,

\origpage{131}
\[ \begin{aligned} H=I^{\mu}L^{\mu}\log\operatorname{tang}(45^{\circ}+\tfrac{1}{2}\varphi^{\mu})&-I^{\mu}\cdot\frac{B^{\mu}}{1+n^{\mu}}\cdot\frac{\operatorname{arc}\operatorname{tang}(\sqrt{n^{\mu}}\sin\varphi^{\mu})}{\sqrt{n^{\mu}}}\\ &+\frac{\tfrac{1}{2}B'}{1+n'}\cdot\frac{\operatorname{arc}\operatorname{tang}(\sqrt{n}\sin\varphi)}{\sqrt{n}}\\ &+(1+b')\cdot\frac{\tfrac{1}{2}B''}{1+n''}\cdot\frac{\operatorname{arc}\operatorname{tang}(\sqrt{n'}\sin\varphi')}{\sqrt{n'}}\\ &+(1+b')(1+b'')\cdot\frac{\tfrac{1}{2}B'''}{1+n'''}\cdot\frac{\operatorname{arc}\operatorname{tang}(\sqrt{n''}\sin\varphi'')}{\sqrt{n''}}\\ &\ \vdots\\ &+I^{\mu-1}\cdot\frac{\tfrac{1}{2}B^{\mu}}{1+n^{\mu}}\cdot\frac{\operatorname{arc}\operatorname{tang}(\sqrt{n^{\mu-1}}\sin\varphi^{\mu-1})}{\sqrt{n^{\mu-1}}}. \end{aligned} \]

\subsection*{(92).}

Cette formule a l'inconvénient d'offrir une suite divergente, et il importe de lui donner une autre forme. Séparons pour cet effet le premier terme $I^{\mu}L^{\mu}\log\operatorname{tang}(45^{\circ}+\tfrac{1}{2}\varphi^{\mu})$, et appelons le reste $Z^{\mu}$; désignons de plus par $T^{i}$ la quantité $\frac{\operatorname{arc}\operatorname{tang}(\sqrt{n^{i}}\sin\varphi^{i})}{\sqrt{n^{i}}}$, nous aurons
\[ Z^{\mu+1}-Z^{\mu}=\frac{B^{\mu}}{1+n^{\mu}}\cdot I^{\mu}T^{\mu}-\frac{B^{\mu+1}}{1+n^{\mu+1}}I^{\mu+1}T^{\mu+1}+\frac{\tfrac{1}{2}B^{\mu+1}}{1+n^{\mu+1}}I^{\mu}T^{\mu}. \]
Cette équation peut se mettre sous la forme
\[ Z^{\mu+1}-Z^{\mu}=I^{\mu}T^{\mu}\left(\frac{B^{\mu}}{1+n^{\mu}}+\frac{\tfrac{1}{2}B^{\mu+1}}{1+n^{\mu+1}}\right)-I^{\mu+1}T^{\mu+1}\frac{B^{\mu+1}}{1+n^{\mu+1}}; \]
mais si on fait pour un moment $\Lambda^{i}=\frac{1}{\cos\lambda}\cdot\frac{1}{\cos\lambda'}\ldots\frac{1}{\cos\lambda^{i}}$, on aura $\frac{B^{\mu}}{1+n^{\mu}}=\frac{2^{\mu}B}{1+n}\Lambda^{\mu-1}$ et $\frac{B^{\mu+1}}{1+n^{\mu+1}}=\frac{2^{\mu+1}B}{1+n}\Lambda^{\mu}$; donc
\[ Z^{\mu+1}-Z^{\mu}=\frac{2^{\mu+1}B}{1+n}I^{\mu}\Lambda^{\mu}\left[\frac{1+\cos\lambda^{\mu}}{2}T^{\mu}-(1+b^{\mu+1})T^{\mu+1}\right]. \]
Or d'une part la quantité $I^{\mu}\Lambda^{\mu}$ ne peut passer une certaine limite peu élevée au-dessus de l'unité; d'autre part les différences $1-\cos\lambda^{\mu}$, $T^{\mu}-T^{\mu+1}$ sont de l'ordre $b^{\mu+1}$; donc $Z^{\mu+1}-Z^{\mu}$ est une quantité très-petite de l'ordre $2^{\mu+1}b^{\mu+1}$ qu'on peut regarder comme intermédiaire entre $b^{\mu}$ et $b^{\mu+1}$, mais beaucoup plus rapprochée du dernier.

\origpage{132}
On pourra donc exprimer $Z^{\mu}$ par la suite fort convergente $Z+(Z'-Z)+(Z''-Z')+\text{etc.}$, et on aura pour l'expression générale de la fonction $H$, cette formule
\[ \begin{aligned} H=I^{\mu}L^{\mu}\log\operatorname{tang}(45^{\circ}+\tfrac{1}{2}\varphi^{\mu})&-\frac{B}{1+n}\cdot\frac{\operatorname{arc}\operatorname{tang}(\sqrt{n}\sin\varphi)}{\sqrt{n}}\\ &-\frac{2B}{1+n}\cdot\frac{1}{\cos\lambda}\left[(1+b')\frac{\operatorname{arc}\operatorname{tang}(\sqrt{n'}\sin\varphi')}{\sqrt{n'}}-\left(\frac{1+\cos\lambda}{2}\right)\cdot\frac{\operatorname{arc}\operatorname{tang}(\sqrt{n}\sin\varphi)}{\sqrt{n}}\right]\\ &-\frac{4B}{1+n}\cdot\frac{1+b'}{\cos\lambda\cos\lambda'}\left[(1+b'')\frac{\operatorname{arc}\operatorname{tang}(\sqrt{n''}\sin\varphi'')}{\sqrt{n''}}-\left(\frac{1+\cos\lambda'}{2}\right)\cdot\frac{\operatorname{arc}\operatorname{tang}(\sqrt{n'}\sin\varphi')}{\sqrt{n'}}\right]\\ &-\frac{8B}{1+n}\cdot\frac{1+b'\cdot 1+b''}{\cos\lambda\cos\lambda'\cos\lambda''}\left[(1+b''')\frac{\operatorname{arc}\operatorname{tang}(\sqrt{n'''}\sin\varphi''')}{\sqrt{n'''}}-\left(\frac{1+\cos\lambda''}{2}\right)\frac{\operatorname{arc}\operatorname{tang}(\sqrt{n''}\sin\varphi'')}{\sqrt{n''}}\right]\\ &-\text{etc.} \end{aligned} \]

Cette formule fondée sur celle que nous avons prise pour valeur de $H^{\mu}$, suppose que l'angle $\varphi^{\mu}$ est toujours moindre que la limite donnée par l'équation $\cot\varphi^{\mu}=\sqrt{b^{\mu}}$. Lorsque $\varphi^{\mu}$ est égal à cette limite, ou qu'il en est seulement approché, la valeur supposée de $H^{\mu}$ est exacte aux quantités près de l'ordre $b^{\mu}$, et alors il faut prolonger la série qui donne la valeur de $H$, jusqu'au $\mu^{\text{ième}}$ terme $-\frac{2^{\mu}B}{1+n}$, etc. inclusivement. Dans le cas où $\varphi^{\mu}$ sera éloigné de la limite dont il s'agit, ce qui arrivera lorsque $\varphi$ n'excède pas $90^{\circ}$, comme on peut toujours le supposer, on pourra ne calculer la série que jusqu'au terme $(\mu-1)^{\text{ième}}$ inclusivement, et le résultat sera toujours exact aux quantités près de l'ordre $b^{\mu}$.

Pour déduire de la formule générale la valeur de la fonction complète $H^{1}$, on pourrait faire $\varphi=90^{\circ}$, et prolonger la série, comme il vient d'être dit, jusqu'au $(\mu-1)^{\text{ième}}$ terme inclusivement. Mais il est plus simple de faire $\varphi=2^{\mu-1}\frac{\pi}{2}=2^{\mu-2}\pi$, ce qui donnera successivement $\varphi'=2^{\mu-3}\pi$, $\varphi''=2^{\mu-4}\pi\ldots\varphi^{\mu-1}=\tfrac{1}{2}\pi$, $\cot\varphi^{\mu}=\sqrt{b^{\mu}}$. On parvient ainsi à la limite des valeurs de $\varphi^{\mu}$, et en négligeant les quantités de l'ordre $b^{\mu}$, on aura $\sin\varphi^{\mu}=1$, $\log\operatorname{tang}(45^{\circ}+\tfrac{1}{2}\varphi^{\mu})=\tfrac{1}{2}\log\left(\frac{4}{b^{\mu}}\right)$, et par la substitution de toutes ces valeurs, on trouve
\[ \begin{aligned} 2^{\mu-1}H^{1}=I^{\mu}L^{\mu}\tfrac{1}{2}\log\frac{4}{b^{\mu}}&-\frac{2^{\mu-1}B}{1+n}\cdot\frac{1+b'\cdot 1+b''\ldots 1+b^{\mu-1}}{\cos\lambda\cos\lambda'\ldots\cos\lambda^{\mu-2}}\cdot\frac{\operatorname{arc}\operatorname{tang}\sqrt{n^{\mu-1}}}{\sqrt{n^{\mu-1}}}\\ &-\frac{2^{\mu}B}{1+n}\left(\frac{1+b'\cdot 1+b''\ldots 1+b^{\mu-1}}{\cos\lambda\cos\lambda'\ldots\ldots\cos\lambda^{\mu-1}}\left[(1+b^{\mu})\frac{\operatorname{arc}\operatorname{tang}\sqrt{n^{\mu}}}{\sqrt{n^{\mu}}}-\left(\frac{1+\cos\lambda^{\mu-1}}{2}\right)\frac{\operatorname{arc}\operatorname{tang}\sqrt{n^{\mu-1}}}{\sqrt{n^{\mu-1}}}\right]. \right. \end{aligned} \]
\ednote{The parenthesis opened before the second fraction of the last line is not closed in the print.}

\origpage{133}
Mais les quantités $1-\cos\lambda^{\mu-1}$, $n^{\mu}-n^{\mu-1}$ sont de l'ordre $b^{\mu}$, et par conséquent négligeables; on a donc plus simplement
\[ H^{1}=I^{\mu}\left[\frac{L^{\mu}}{2^{\mu}}\log\frac{4}{b^{\mu}}-\frac{B}{1+n}\cdot\frac{1}{\cos\lambda}\cdot\frac{1}{\cos\lambda'}\ldots\frac{1}{\cos\lambda^{\mu-2}}\cdot\frac{\operatorname{arc}\operatorname{tang}\sqrt{n^{\mu-1}}}{\sqrt{n^{\mu-1}}}\right], \]
formule où il faudra substituer la valeur connue de $L^{\mu}$ et celle de $I^{\mu}=(1+b')(1+b'')\ldots(1+b^{\mu-1})=\frac{2\sqrt{b'}}{b}\cdot\frac{2\sqrt{b''}}{b'}\cdot\frac{2\sqrt{b'''}}{b''}$, etc., ce produit étant continué jusqu'à ce que le nombre de ses facteurs soit $\mu-1$.

Telles sont les formules les plus simples par lesquelles on pourra calculer les valeurs des fonctions $H$ et $H^{1}$, lorsque le module $c$ est extrêmement près de l'unité; elles donneront tel degré d'approximation qu'on voudra obtenir, en prolongeant suffisamment les séries qui les composent.

\subsection*{(93).}

Si on compare maintenant cette seconde méthode avec la première, on ne manquera pas de remarquer que celle-ci n'est sujette à aucune exception, et qu'elle n'exige ni considérations de limites, ni transformations; il est donc préférable d'employer la première méthode, même dans les cas où $c$ serait fort près de l'unité, puisqu'alors cette méthode n'a d'autre inconvénient que celui d'employer quelques termes de plus, et le nombre de ces termes ne peut jamais être bien grand.

Nous remarquerons encore qu'on peut simplifier à quelques égards les formules de la première méthode, en employant des angles subsidiaires $\lambda^{0}$, $\lambda^{00}$, $\lambda^{000}$, etc. pour calculer tant la série des paramètres $n^{0}$, $n^{00}$, etc., que celle des quantités $k$, $k^{0}$, $k^{00}$, etc., à l'exemple de ce que nous avons pratiqué dans la seconde méthode. Pour cet effet on aura les formules successives
\[ \begin{aligned} 1+n&=b\cot^{2}\tfrac{1}{2}\lambda^{0} &&\ldots\ldots &\operatorname{tang}\tfrac{1}{2}\lambda^{0}&=\sqrt{\left(\frac{b}{1+n}\right)}\\ 1+n^{0}&=b^{0}\cot^{2}\tfrac{1}{2}\lambda^{00}=\frac{b^{02}}{\sin^{2}\lambda^{0}} &&\ldots\ldots &\operatorname{tang}\tfrac{1}{2}\lambda^{00}&=\frac{\sin\lambda^{0}}{\sqrt{b^{0}}}\\ 1+n^{00}&=b^{00}\cot^{2}\tfrac{1}{2}\lambda^{000}=\frac{b^{002}}{\sin^{2}\lambda^{00}} &&\ldots\ldots &\operatorname{tang}\tfrac{1}{2}\lambda^{000}&=\frac{\sin\lambda^{00}}{\sqrt{b^{00}}}\\ 1+n^{000}&=b^{000}\cot^{2}\tfrac{1}{2}\lambda^{0000}=\frac{b^{0002}}{\sin^{2}\lambda^{000}} &&\ldots\ldots &\operatorname{tang}\tfrac{1}{2}\lambda^{0000}&=\frac{\sin\lambda^{000}}{\sqrt{b^{000}}}\\ &\text{etc.} &&& &\text{etc.} \end{aligned} \]

\origpage{134}
d'où l'on déduit
\[ \begin{aligned} k&=\frac{1+n^{0}}{1+n}\cos\lambda^{0} &&\ldots\ldots & \frac{k}{1+n^{0}}&=\frac{1}{1+n}\cos\lambda^{0}\\ k^{0}&=\frac{1+n^{00}}{1+n^{0}}\cos\lambda^{00} &&\ldots\ldots & \frac{kk^{0}}{1+n^{00}}&=\frac{1}{1+n}\cos\lambda^{0}\cos\lambda^{00}\\ k^{00}&=\frac{1+n^{000}}{1+n^{00}}\cos\lambda^{000} &&\ldots\ldots & \frac{kk^{0}k^{00}}{1+n^{000}}&=\frac{1}{1+n}\cos\lambda^{0}\cos\lambda^{00}\cos\lambda^{000}\\ &\text{etc.} &&& &\text{etc.} \end{aligned} \]
Enfin si l'on fait
\[ \begin{aligned} A^{\mu}=A+\frac{\tfrac{1}{2}B}{1+n}(1&+\tfrac{1}{2}\cos\lambda^{0}+\tfrac{1}{4}\cos\lambda^{0}\cos\lambda^{00}\ldots\\ &\ldots+\frac{1}{2^{\mu}}\cos\lambda^{0}\cos\lambda^{00}\ldots\cos\lambda^{\mu}), \end{aligned} \]
et qu'on prenne toujours $K^{\mu}=(1+c^{0})(1+c^{00})\ldots(1+c^{\mu})$, la valeur de $H$ sera
\[ \begin{aligned} H=K^{\mu}A^{\mu}\Phi&-\frac{\sqrt{c^{0}}}{2c}\cdot\frac{B}{1+n}\cdot\frac{\operatorname{arc}\operatorname{tang}(\sqrt{n^{0}}\sin\varphi^{0})}{\sqrt{n^{0}}}\\ &-\frac{\sqrt{c^{0}}}{2c}\cdot\frac{\sqrt{c^{00}}}{2c^{0}}\cdot\frac{B}{1+n}\cdot\cos\lambda^{0}\cdot\frac{\operatorname{arc}\operatorname{tang}(\sqrt{n^{00}}\sin\varphi^{00})}{\sqrt{n^{00}}}\\ &-\frac{\sqrt{c^{0}}}{2c}\cdot\frac{\sqrt{c^{00}}}{2c^{0}}\cdot\frac{\sqrt{c^{000}}}{2c^{00}}\cdot\frac{B}{1+n}\cdot\cos\lambda^{0}\cos\lambda^{00}\cdot\frac{\operatorname{arc}\operatorname{tang}(\sqrt{n^{000}}\sin\varphi^{000})}{\sqrt{n^{000}}}\\ &-\text{etc.} \end{aligned} \]
La limite de $n^{\mu}$ étant zéro dans tous les cas, celle de $\sin\lambda^{\mu}$ sera $b^{\mu}$, et par conséquent celle de $\cos\lambda^{\mu}$ sera $c^{\mu}$. D'où l'on voit combien cette suite est convergente.

\section*{Des cas les plus généraux dans lesquels on peut opérer la réduction des fonctions elliptiques de la troisième espèce.}

\subsection*{(94).}

Si l'on différentie par rapport à $n$ chaque membre de l'équation $\Pi=\displaystyle\int\frac{d\varphi}{(1+n\sin^{2}\varphi)\Delta}$, on aura
\[ n\frac{d\Pi}{dn}=\int\frac{d\varphi}{(1+n\sin^{2}\varphi)^{2}\Delta}-\int\frac{d\varphi}{(1+n\sin^{2}\varphi)\Delta}: \]

\origpage{135}
or par la formule du n° 9 on trouve, en faisant $\alpha=(1+n)\left(1+\frac{c^{2}}{n}\right)$,
\[ \begin{aligned} \frac{2\alpha}{n}\int\frac{d\varphi}{(1+n\sin^{2}\varphi)^{2}\Delta}=\frac{\Delta\sin\varphi\cos\varphi}{1+n\sin^{2}\varphi}&-\frac{c^{2}}{n^{2}}\int(1+n\sin^{2}\varphi)\frac{d\varphi}{\Delta}\\ &+\left(1+\frac{2+2c^{2}}{n}+\frac{3c^{2}}{n^{2}}\right)\int\frac{d\varphi}{(1+n\sin^{2}\varphi)\Delta}. \end{aligned} \]
Combinant ces deux équations entre elles, et observant qu'on a
\[ c^{2}\int(1+n\sin^{2}\varphi)\frac{d\varphi}{\Delta}=(c^{2}+n)F-nE, \]
on parviendra à l'équation suivante, où l'on ne doit regarder que $n$ comme variable dans $\Pi$ et dans $\alpha$,
\[ \tfrac{1}{2}\Pi\,d\alpha+\alpha\,d\Pi=-\tfrac{1}{2}c^{2}F\cdot\frac{dn}{n^{2}}-\tfrac{1}{2}(F-E)\frac{dn}{n}+\tfrac{1}{2}\Delta\sin\varphi\cos\varphi\cdot\frac{dn}{1+n\sin^{2}\varphi}. \]
Multipliant de part et d'autre par $\alpha^{-\frac{1}{2}}$ et intégrant, on aura
\[ \sqrt{\alpha}\cdot\Pi=-\tfrac{1}{2}c^{2}F\int\frac{dn}{n^{2}\sqrt{\alpha}}-\tfrac{1}{2}(F-E)\int\frac{dn}{n\sqrt{\alpha}}+\tfrac{1}{2}\Delta\sin\varphi\cos\varphi\int\frac{dn}{(1+n\sin^{2}\varphi)\sqrt{\alpha}}. \]
Il faut maintenant, pour effectuer les intégrations, considérer successivement les trois cas de l'art.\ 50.

\section*{Premier cas $n=\cot^{2}\theta$.}

\subsection*{(95).}

Alors on a $dn=-\frac{2\cos\theta}{\sin^{3}\theta}\,d\theta$, $\sqrt{\alpha}=\frac{\Delta(b,\theta)}{\sin\theta\cos\theta}$, et l'équation générale devient
\[ \begin{aligned} \frac{\Delta(b,\theta)}{\sin\theta\cos\theta}\Pi=c^{2}F\int\frac{d\theta\sin^{2}\theta}{\cos^{2}\theta\Delta(b,\theta)}&+(F-E)\int\frac{d\theta}{\Delta(b,\theta)}\\ &-\Delta\sin\varphi\cos\varphi\int\frac{d\theta\cos^{2}\theta}{(\sin^{2}\theta+\sin^{2}\varphi\cos^{2}\theta)\Delta(b,\theta)}. \end{aligned} \]
L'intégrale $\displaystyle\int\frac{d\theta}{\Delta(b,\theta)}$ est représentée à l'ordinaire par $F(b,\theta)$; et suivant les réductions connues, on a
\[ \int\frac{c^{2}\,d\theta\sin^{2}\theta}{\cos^{2}\theta\Delta(b,\theta)}=\frac{\sin\theta}{\cos\theta}\Delta(b,\theta)-E(b,\theta); \]
enfin si l'on fait pour abréger $\cot^{2}\varphi=n'$, l'intégrale $\Delta\sin\varphi\cos\varphi\displaystyle\int\frac{d\theta\cos^{2}\theta}{(\sin^{2}\theta+\sin^{2}\varphi\cos^{2}\theta)\Delta(b,\theta)}$ pourra se décomposer ainsi:
\[ -\frac{\Delta\sin\varphi}{\cos\varphi}\int\frac{d\theta}{\Delta(b,\theta)}+\frac{\Delta}{\sin\varphi\cos\varphi}\int\frac{d\theta}{(1+n'\sin^{2}\theta)\Delta(b,\theta)}, \]

\origpage{136}
et sous cette forme elle peut être représentée par
\[ -\frac{\Delta\sin\varphi}{\cos\varphi}F(b,\theta)+\frac{\Delta}{\sin\varphi\cos\varphi}\Pi(n',b,\theta). \]
Substituant toutes ces valeurs, et désignant, pour plus de clarté, par $\Pi(n,c,\varphi)$, $\Delta(c,\varphi)$, $F(c,\varphi)$, $E(c,\varphi)$, les fonctions $\Pi$, $\Delta$, $F$, $E$; on aura, pour le premier cas, la formule générale qui suit:
\[ \begin{aligned} &\frac{\Delta(b,\theta)}{\sin\theta\cos\theta}\Pi(n,c,\varphi)+\frac{\Delta(c,\varphi)}{\sin\varphi\cos\varphi}\Pi(n',b,\theta)\\ &=\left\{\begin{aligned} &C+\frac{\sin\theta}{\cos\theta}\Delta(b,\theta)F(c,\varphi)+\frac{\sin\varphi}{\cos\varphi}\Delta(c,\varphi)F(b,\theta)\\ &+F(c,\varphi)F(b,\theta)-F(c,\varphi)E(b,\theta)-E(c,\varphi)F(b,\theta) \end{aligned}\right.\qquad(i'). \end{aligned} \]
La constante $C$, ajoutée à cette formule, pourrait être fonction de $\varphi$, puisqu'elle résulte d'une intégration où $\varphi$ a été regardée comme constante; mais comme la forme de l'intégrale est telle qu'on peut permuter entre elles les quantités $\varphi$ et $\theta$, pourvu qu'on fasse de même la permutation entre $c$ et $b$, il s'ensuit que $C$ ne contient ni $\theta$ ni $\varphi$, et qu'ainsi $C$ est une constante absolue.

Pour en déterminer la valeur, prenons un cas particulier, et supposons à la fois $\varphi$ et $\theta$ infiniment petits. Cette supposition donne, en rejetant les infiniment petits du second ordre $F(c,\varphi)=E(c,\varphi)=\varphi$, $F(b,\theta)=E(b,\theta)=\theta$, $\Delta(c,\varphi)=\Delta(b,\theta)=1$, $\Pi(n,c,\varphi)=\displaystyle\int\frac{d\varphi}{1+n\sin^{2}\varphi}=\frac{\operatorname{arc}\operatorname{tang}\varphi\sqrt{(1+n)}}{\sqrt{(1+n)}}=\theta\operatorname{arc}\operatorname{tang}\frac{\varphi}{\theta}$, et semblablement $\Pi(n',b,\theta)=\varphi\operatorname{arc}\operatorname{tang}\frac{\theta}{\varphi}$. Substituant toutes ces valeurs, il viendra
\[ C=\operatorname{arc}\operatorname{tang}\frac{\varphi}{\theta}+\operatorname{arc}\operatorname{tang}\frac{\theta}{\varphi}=\tfrac{1}{2}\pi. \]

\subsection*{(96).}

Si l'on veut appliquer la formule de l'art.\ précédent au cas où $\varphi=\tfrac{1}{2}\pi$, il faut, pour éviter les quantités infinies qui se rencontrent dans les deux membres, faire $\varphi=\tfrac{1}{2}\pi-\omega$, et supposer $\omega$ infiniment petit. On aura ainsi $n'=\cot^{2}\varphi=\operatorname{tang}^{2}\omega=\omega^{2}$; et puisque $n'$ est infiniment petit, la fonction $\Pi(n',b,\theta)$ s'exprimera de cette manière
\[ \Pi(n',b,\theta)=\int\frac{d\theta}{(1+n'\sin^{2}\theta)\Delta(b,\theta)}=\int\frac{d\theta}{\Delta(b,\theta)}-n'\int\frac{d\theta\sin^{2}\theta}{\Delta(b,\theta)}: \]

\origpage{137}
on aura donc
\[ \frac{\Delta}{\sin\varphi\cos\varphi}\Pi(n',b,\theta)-\frac{\Delta\sin\varphi}{\cos\varphi}F(b,\theta)=\omega\Delta F(b,\theta)-\omega\Delta\int\frac{d\theta\sin^{2}\theta}{\Delta(b,\theta)}, \]
d'où l'on voit que le premier membre se réduit à zéro, lorsqu'on fera $\omega=0$, et qu'ainsi la supposition de $\varphi=\tfrac{1}{2}\pi$ donne cette formule
\[ \frac{\Delta(b,\theta)}{\sin\theta\cos\theta}\Pi^{1}(n,c)=\tfrac{1}{2}\pi+\left.\begin{aligned} &\frac{\sin\theta}{\cos\theta}\Delta(b,\theta)F^{1}(c)+F^{1}(c)F(b,\theta)\\ &-F^{1}(c)E(b,\theta)-E^{1}(c)F(b,\theta) \end{aligned}\right\}\;(k'). \]
Toutes les fois donc que le paramètre $n$ est positif, la fonction complète de troisième espèce $\Pi^{1}(n,c)$, pourra toujours se déterminer par des fonctions de première et de seconde espèce, savoir, par les fonctions complètes $F^{1}(c)$, $E^{1}(c)$, qui se rapportent au module $c$, et par les fonctions non-complètes $F(b,\theta)$, $E(b,\theta)$, qui se rapportent au module complémentaire $b$, et dont l'amplitude $\theta$ se déduit du paramètre $n$ par l'équation $\cot\theta=\sqrt{n}$.

Ce théorème est très-important dans la théorie des fonctions elliptiques, et son usage sera d'autant plus étendu, que dans beaucoup de problèmes de géométrie ou de mécanique, qui dépendent des fonctions elliptiques, on n'a souvent besoin que des intégrales définies ou complètes. Nous en donnerons ci-après diverses applications.

\subsection*{(97).}

Si dans l'équation $(k')$ on met $\frac{c^{2}}{n}$ à la place de $n$, et qu'on fasse $\frac{c^{2}}{n}=\cot^{2}\lambda$, afin de mettre en même temps $\lambda$ au lieu de $\theta$, on aura
\[ \begin{aligned} \frac{\Delta(b,\lambda)}{\sin\lambda\cos\lambda}\Pi^{1}\left(\frac{c^{2}}{n},c\right)=\tfrac{1}{2}\pi&+\frac{\sin\lambda}{\cos\lambda}\Delta(b,\lambda)F^{1}(c)+F^{1}(c)F(b,\lambda)\\ &-F^{1}(c)E(b,\lambda)-E^{1}(c)F(b,\lambda). \end{aligned} \]
Mais, en vertu de l'équation $c^{2}=n\cot^{2}\lambda$ ou $c\operatorname{tang}\theta\operatorname{tang}\lambda=1$, qui se rapporte au module $b$, on a (n° 18) $F(b,\theta)+F(b,\lambda)=F^{1}(b)$, $E(b,\theta)+E(b,\lambda)=E^{1}(b)+b^{2}\sin\lambda\sin\theta$; on a en même temps $\sin\lambda=\frac{\cos\theta}{\Delta(b,\theta)}$, $\cos\lambda=\frac{c\sin\theta}{\Delta(b,\lambda)}$\ednote{Printed $\Delta(b,\lambda)$ in the denominator of $\cos\lambda$; the value $\cos\lambda=c\sin\theta:\Delta(b,\theta)$ is meant, an apparent misprint.}, $\Delta(b,\lambda)=\frac{c}{\Delta(b,\theta)}$, $\frac{\Delta(b,\lambda)}{\sin\lambda\cos\lambda}$, $=\frac{\Delta(b,\theta)}{\sin\theta\cos\theta}=\sqrt{\alpha}$,

\origpage{138}
Ajoutant donc l'équation précédente avec l'équation $(i')$, et observant que d'après la formule du n° 47, on a $\Pi^{1}(n,c)+\Pi^{1}\left(\frac{c^{2}}{n},c\right)=F^{1}(c)+\frac{\pi}{2\sqrt{\alpha}}$, on trouve que les angles $\theta$ et $\lambda$ disparaissent entièrement du calcul, et qu'il reste entre les fonctions complètes, pour les modules $c$ et $b$, cette relation:
\[ \frac{\pi}{2}=F^{1}(c)E^{1}(b)+F^{1}(b)E^{1}(c)-F^{1}(c)F^{1}(b), \]
ce qui est une nouvelle démonstration du théorème de l'art.\ 42.

\subsection*{(98).}

Puisqu'on peut exprimer la fonction complète $\Pi'(n,c)$\ednote{The print shows a prime here where the work writes $\Pi^{1}$ for the complete function throughout; probably a damaged or wrong superscript $1$.} par des fonctions de la première et de la seconde espèce, on pourra exprimer de la même manière toute fonction non-complète $\Pi(n,c,\varphi)$, dont l'amplitude $\varphi$ est telle qu'on ait $F(c,\varphi)=kF^{1}(c)$, $k$ étant un nombre rationnel. En effet, si cette relation a lieu, il suit des formules démontrées ci-dessus, qu'on a $\Pi(n,c,\varphi)=k\Pi^{1}(n,c)+W$, $W$ étant une quantité déterminable par arcs de cercle. Donc il y a une infinité de cas où la fonction $\Pi$, dont le paramètre est à volonté, pourvu qu'il soit positif, pourra se réduire aux fonctions de la première et de la seconde espèce, et ces cas sont déterminés par un symptôme général.

Etant donné l'amplitude $\varphi$, on peut trouver par un calcul assez simple, si la condition dont on vient de parler peut être remplie, c'est-à-dire, si la fonction $F(\varphi)$ est dans un rapport rationnel avec la fonction complète $F^{1}$. Il faut pour cela calculer la valeur approchée de la fonction $F(\varphi)$ par la méthode de l'article 65. Lorsqu'on aura déterminé l'angle $\Phi$, limite des angles $\varphi$, $\frac{\varphi^{0}}{2}$, $\frac{\varphi^{00}}{4}$, etc., il faudra examiner si cette limite $\Phi$ est avec l'angle droit dans un rapport à-la-fois simple et très-approché. Si le rapport était un peu composé, il n'y aurait lieu à aucune réduction, attendu que la réduction exige que le rapport soit exact, et que, pour peu qu'il fût composé, l'équation qui détermine $\sin\varphi$ serait d'un degré très-élevé, ce qui rendrait peu probable que la valeur donnée de $\varphi$ satisfît à cette équation. Mais si le rapport dont il s'agit est exprimé en nombres simples et qu'il paraisse au moins fort approché, il y aura lieu de croire qu'il est exact. Alors on s'assurera par les formules rigoureuses, si la valeur donnée de $\sin\varphi$ répond au

\origpage{139}
rapport $k$ trouvé entre $F(c,\varphi)$ et $F^{1}(c)$, ou entre $\Phi$ et $\tfrac{1}{2}\pi$. Lorsque cela a lieu, on trouvera également par les formules rigoureuses, la valeur de $W$ qui donne $\Pi(n,c,\varphi)=k\Pi^{1}(n,c)+W$; ainsi on pourra déterminer $\Pi(n,c,\varphi)$ par des fonctions de la première et de la seconde espèce.

La méthode qu'on vient d'indiquer n'est autre chose que la méthode d'approximation par laquelle on évalue la fonction de première espèce $F(c,\varphi)$. Mais on peut pour le même objet employer une autre méthode qui paraît conduire plus directement au but.

Etant donné l'amplitude $\varphi$, on calculera successivement les amplitudes $\varphi_{2}$, $\varphi_{3}$, $\varphi_{4}$, etc., qui donnent $F(\varphi_{2})=2F(\varphi)$, $F(\varphi_{3})=3F(\varphi)$, $F(\varphi_{4})=4F(\varphi)$, etc.; c'est ce qu'on fera par les formules
\[ \begin{aligned} \operatorname{tang}\tfrac{1}{2}\varphi_{2}&=\Delta\operatorname{tang}\varphi\\ \operatorname{tang}(\tfrac{1}{2}\varphi_{3}+\tfrac{1}{2}\varphi)&=\Delta\operatorname{tang}\varphi_{2}\\ \operatorname{tang}(\tfrac{1}{2}\varphi_{4}+\tfrac{1}{2}\varphi_{2})&=\Delta\operatorname{tang}\varphi_{3}\\ \operatorname{tang}(\tfrac{1}{2}\varphi_{5}+\tfrac{1}{2}\varphi_{3})&=\Delta\operatorname{tang}\varphi_{4}\\ &\text{etc.} \end{aligned} \]
Il est évident que si l'angle $\varphi$ est tel qu'on ait $F(c,\varphi)=kF^{1}(c)$, $k$ étant un nombre rationnel $\frac{\mu}{\nu}$, il y aura nécessairement dans la suite $\varphi_{2}$, $\varphi_{3}$, $\varphi_{4}$, etc. un terme $\varphi_{\nu}$ tel que $\varphi_{\nu}=\mu\cdot\frac{\pi}{2}$, c'est-à-dire qu'on devra rencontrer dans cette suite un terme $\varphi_{\nu}$ qui soit multiple de l'angle droit. Et comme d'ailleurs, par les raisons que nous avons exposées, le rapport $\frac{\mu}{\nu}$ devra être exprimé en nombres assez simples, on n'aura jamais qu'un petit nombre de termes à calculer dans la suite $\varphi_{2}$, $\varphi_{3}$, $\varphi_{4}$, etc., pour reconnaître si les fonctions $F(\varphi)$ et $F^{1}$ sont commensurables entre elles, et par suite si $\Pi(n,c,\varphi)$ peut se mesurer par $\Pi^{1}(n,c)$.

\subsection*{(99).}

Pour revenir à l'équation $(i')$, nous observerons que cette équation servira en général à déterminer l'une des fonctions $\Pi(n,c,\varphi)$, $\Pi(n',b,\theta)$, par le moyen de l'autre supposée connue. Ces deux fonctions ont des modules complémens l'un de l'autre, et leurs amplitudes se déduisent réciproquement de leurs paramètres par les équations $n=\cot^{2}\theta$, $n'=\cot^{2}\varphi$.

\origpage{140}
Supposons que l'angle $\theta$, déduit du paramètre $n$, soit tel qu'on ait $F(b,\theta)=kF^{1}(b)$, $k$ étant un nombre rationnel; on aura par les propriétés connues, $E(b,\theta)=kE^{1}(b)+V$, et $\Pi(n',b,\theta)=\Pi^{1}(n',b)+W$\ednote{Printed without the factor $k$; the corresponding statement in art.\ 102 reads $k\Pi^{1}(n',b)+W$, so the $k$ has apparently dropped out.}, les quantités $V$ et $W$ étant déterminables, la première algébriquement, la seconde par arcs de cercle.

Si maintenant on substitue ces valeurs dans l'équation $(i')$, et qu'ensuite on mette pour $\Pi^{1}(n',b)$ sa valeur déduite de la formule $(k')$, on aura, après les réductions, ce résultat très-simple
\[ \begin{aligned} \frac{\Delta(b,\theta)}{\sin\theta\cos\theta}\Pi(n,c,\varphi)=(1-k)\frac{\pi}{2}&+\frac{\sin\theta}{\cos\theta}\Delta(b,\theta)F(c,\varphi)\\ &-VF(c,\varphi)-\frac{W\Delta(c,\varphi)}{\sin\varphi\cos\varphi}. \end{aligned} \]
D'où l'on voit que la fonction $\Pi$ peut alors se réduire indéfiniment aux fonctions de la première espèce, et la seule condition nécessaire pour cette réduction, est que le paramètre $n$, représenté par $\cot^{2}\theta$, soit tel que le rapport de $F(b,\theta)$ à $F^{1}(b)$, soit rationnel. Ainsi nous avons encore pour ces réductions un symptôme général qui comprend une infinité de cas particuliers.

Le cas de $n=c$ qu'on a résolu n° 46, est compris dans le symptôme général, puisqu'en faisant $\cot^{2}\theta=c$, la valeur de $\theta$ qui en résulte, est celle qui, appliquée au module $b$, donne $F(b,\theta)=\tfrac{1}{2}F^{1}(b)$. Aussi en faisant $k=\tfrac{1}{2}$ et substituant les valeurs de $V$ et $W$ qui sont $V=\tfrac{1}{2}(1-c)$, $W=\frac{\sin\varphi\cos\varphi}{2\Delta}\operatorname{arc}\operatorname{tang}.\frac{\Delta\cos\varphi}{(1+c)\sin\varphi}=\frac{\sin\varphi\cos\varphi}{2\Delta}\times\left(\frac{\pi}{2}-\operatorname{arc}\operatorname{tang}.\frac{(1+c)\operatorname{tang}\varphi}{\Delta}\right)$, on retrouve la formule
\[ \Pi(c)=\tfrac{1}{2}F+\frac{\tfrac{1}{2}F}{1+c}\operatorname{arc}\operatorname{tang}.\frac{(1+c)\operatorname{tang}\varphi}{\Delta}. \]
\ednote{Printed $\frac{\frac{1}{2}F}{1+c}$; the arc-tangent term comes from $W$ alone and carries no factor $F$ (compare the analogous formula for $\Pi(-c)$ in art.\ 102, numerator $\tfrac{1}{2}$), so the $F$ in the numerator is apparently a misprint.}

\section*{Second cas, $n=-1+b^{2}\sin^{2}\theta$.}

\subsection*{(100).}

On aura dans ce cas, $\frac{dn}{\sqrt{\alpha}}=2\,d\theta\Delta(b,\theta)$, $\sqrt{\alpha}=\frac{b^{2}\sin\theta\cos\theta}{\Delta(b,\theta)}$, et l'équation générale du n° 94 deviendra
\[ \begin{aligned} \frac{b^{2}\sin\theta\cos\theta}{\Delta(b,\theta)}\Pi=(F-E)\int\frac{d\theta}{\Delta(b,\theta)}&-c^{2}F\int\frac{d\theta}{(1-b^{2}\sin^{2}\theta)^{\frac{3}{2}}}\\ &+\Delta\sin\varphi\cos\varphi\int\frac{d\theta\Delta(b,\theta)}{\cos^{2}\varphi+b^{2}\sin^{2}\theta\sin^{2}\varphi}. \end{aligned} \]

\origpage{141}
Or par les formules connues on a $\displaystyle\int\frac{d\theta}{\Delta(b,\theta)}=F(b,\theta)$ et $\displaystyle\int\frac{d\theta}{(1-b^{2}\sin^{2}\theta)^{\frac{3}{2}}}=\frac{1}{c^{2}}E(b,\theta)-\frac{b^{2}}{c^{2}}\cdot\frac{\sin\theta\cos\theta}{\Delta(b,\theta)}$; de plus en faisant $b^{2}\operatorname{tang}^{2}\varphi=n$\ednote{Printed $n$; the parameter so defined is $n'$, as in art.\ 102, so the prime has apparently dropped out.}, on a
\[ \Delta\sin\varphi\cos\varphi\int\frac{d\theta\Delta(b,\theta)}{\cos^{2}\varphi+b^{2}\sin^{2}\theta\sin^{2}\varphi}=\frac{\Delta}{\sin\varphi\cos\varphi}\Pi(n',b,\theta)-\frac{\Delta\cos\varphi}{\sin\varphi}F(b,\varphi); \]
\ednote{Printed $F(b,\varphi)$ at the end of this display; $F(b,\theta)$ is meant, as the formula $(l')$ below shows, an apparent misprint.}
donc en faisant toutes ces substitutions on aura pour le second cas, la formule générale
\[ \frac{b^{2}\sin\theta\cos\theta}{\Delta(b,\theta)}[\Pi(n,c,\varphi)-F(c,\varphi)]=\left\{\begin{aligned} &\frac{\Delta(c,\varphi)}{\sin\varphi\cos\varphi}[\Pi(n',b,\theta)-\cos^{2}\varphi F(b,\theta)]\\ &+F(c,\varphi)F(b,\theta)-E(c,\varphi)F(b,\theta)-F(c,\varphi)E(b,\theta) \end{aligned}\right\}\;(l'). \]
On n'a point ajouté de constante, parce qu'en faisant $\theta=0$, les deux membres s'évanouissent.

\subsection*{(101).}

Si on veut savoir ce que donne cette équation lorsque $\varphi=\tfrac{1}{2}\pi$, il faudra trouver la valeur de $\frac{\Delta(c,\varphi)}{\sin\varphi\cos\varphi}\Pi(n',b,\theta)$ à cette limite; et pour cela, il faut d'abord supposer $\varphi=\tfrac{1}{2}\pi-\omega$, $\omega$ étant infiniment petit. Or puisque $n'=b^{2}\operatorname{tang}^{2}\varphi=b^{2}\cot^{2}\omega$, il en résulte $\frac{b^{2}}{n'}=\operatorname{tang}^{2}\omega$, et la formule du n° 46 donne
\[ \Pi(n',b,\theta)+\Pi(\operatorname{tang}^{2}\omega,b,\theta)=F(b,\theta)+\frac{\sin\varphi\cos\varphi}{\Delta(c,\varphi)}\operatorname{arc}\operatorname{tang}.\frac{\Delta(c,\varphi)\operatorname{tang}\theta}{\sin\varphi\cos\varphi\Delta(b,\theta)}. \]
Mais $\omega$ étant infiniment petit, on a
\[ \Pi(\operatorname{tang}^{2}\omega,b,\theta)=F(b,\theta)-M\operatorname{tang}^{2}\omega, \]
$M$ étant une quantité finie; donc
\[ \frac{\Delta(c,\varphi)}{\sin\varphi\cos\varphi}\Pi(n',b,\theta)=\frac{\cos\varphi\Delta(c,\varphi)}{\sin^{3}\varphi}M+\operatorname{arc}\operatorname{tang}.\frac{\Delta(c,\varphi)\operatorname{tang}\theta}{\sin\varphi\cos\varphi\Delta(b,\theta)}. \]
Le second membre se réduit à $\tfrac{1}{2}\pi$ lorsqu'on fait $\varphi=\tfrac{1}{2}\pi$, donc on aura pour déterminer $\Pi^{1}(n,c)$, l'équation
\[ \frac{b^{2}\sin\theta\cos\theta}{\Delta(b,\theta)}[\Pi^{1}(n,c)-F^{1}(c)]=\left.\begin{aligned} &\frac{\pi}{2}+F^{1}(c)F(b,\theta)\\ &-E^{1}(c)F(b,\theta)-F^{1}(c)E(b,\theta) \end{aligned}\right\}\;(m'). \]

\origpage{142}
d'où l'on voit que la fonction complète de troisième espèce $\Pi^{1}(n,c)$, s'exprime généralement par des fonctions de la première et de la seconde espèce.

Nous observerons que les formules $(l')$ et $(m')$ auraient pu se déduire de celles qui ont été données art.\ 51, 95 et 96; mais il n'était pas inutile de faire voir comment on pouvait y parvenir directement.

La fonction non-complète $\Pi(n,c,\varphi)$ s'exprimera de même par des fonctions de la première et de la seconde espèce, si l'amplitude $\varphi$ est telle qu'on ait $F(c,\varphi)=kF^{1}(c)$, $k$ étant un nombre rationnel. Car alors on a $\Pi(n,c,\varphi)=k\Pi^{1}(n,c)+W$, $W$ étant une quantité déterminable par arcs de cercle.

\subsection*{(102).}

Supposons que l'angle $\theta$, déterminé par le paramètre $n=-1+b^{2}\sin^{2}\theta$, soit tel que pour le module $b$ on ait $F(b,\theta)=kF^{1}(b)$, $k$ étant rationnel; alors on aura par les propriétés connues $\Pi(n',b,\theta)=k\Pi^{1}(n',b)+W$, $W$ étant une quantité déterminable par des arcs de cercle. On aura en même temps $E(b,\theta)=kE^{1}(b)+V$, $V$ étant une quantité déterminable algébriquement. Il suffira donc d'avoir la valeur de $\Pi^{1}(n',b)$ où l'on a $n'=b^{2}\operatorname{tang}^{2}\varphi$.

Pour cet effet, j'observe que la formule du n° 47 donne
\[ \Pi^{1}(n',b)+\Pi^{1}(\cot^{2}\varphi,b)=F^{1}(b)+\frac{\pi\sin\varphi\cos\varphi}{2\Delta(c,\varphi)}; \]
ensuite par la formule de l'art.\ 96, on a
\[ \Pi^{1}(\cot^{2}\varphi,b)=\frac{\sin\varphi\cos\varphi}{\Delta(c,\varphi)}\left\{\begin{aligned} &\frac{\pi}{2}+\frac{\sin\varphi}{\cos\varphi}\Delta(c,\varphi)F^{1}(b)+F^{1}(b)F(c,\varphi)\\ &-E^{1}(b)F(c,\varphi)-F^{1}(b)E(c,\varphi) \end{aligned}\right\}. \]
Faisant toutes ces substitutions dans l'équation $(l')$ et réduisant, on aura
\[ \Pi(n,c,\varphi)=\left(1-\frac{\Delta(b,\theta)}{b^{2}\sin\theta\cos\theta}V\right)F(c,\varphi)+\frac{\Delta(c,\varphi)\Delta(b,\theta)}{b^{2}\sin\varphi\cos\varphi\sin\theta\cos\theta}W. \]
D'où l'on voit que la fonction de troisième espèce $\Pi(n,c,\varphi)$ se réduira indéfiniment à la fonction de première espèce $F(c,\varphi)$; et la seule condition nécessaire pour cette réduction, est que l'angle $\theta$

\origpage{143}
déduit du paramètre $n=-1+b^{2}\sin^{2}\theta$, soit tel qu'on ait $F(b,\theta)=kF^{1}(b)$, $k$ étant un nombre rationnel.

Ainsi dans le cas de $n=-c$, où l'on a $\sin^{2}\theta=\frac{1}{1+c}$, cette valeur de $\theta$ est celle qui donne $F(c,\theta)=\frac{1}{2}F^{1}(b)$\ednote{Printed $F(c,\theta)$; the angle $\theta$ belongs to the module $b$ (art.\ 102: $F(b,\theta)=kF^{1}(b)$), so $F(b,\theta)$ is meant.}; la réduction est donc possible. On trouve ensuite par les formules du n° 57, $V=\frac{1}{2}(1-c)$, $W=\frac{\sin\varphi\cos\varphi}{2\Delta}\operatorname{arc}\operatorname{tang}\frac{(1-c)\operatorname{tang}\varphi}{\Delta}$; et substituant ces valeurs dans la formule de l'article précédent, on en tire
\[ \Pi(-c)=\tfrac{1}{2}F+\frac{\frac{1}{2}}{1-c}\operatorname{arc}\operatorname{tang}.\frac{(1-c)\operatorname{tang}\varphi}{\Delta}, \]
ce qui s'accorde avec la formule du n° 46.

\subsection*{(103).}

Si dans l'équation $(l')$ on fait $n=-c^{2}$ ou $\theta=\frac{1}{2}\pi$, on trouvera encore la même équation qu'au n° 48. Mais si on fait $n=-1+\omega^{2}$, $\omega$ étant supposé infiniment petit, on aura $\sin\theta=\frac{\omega}{b}$, ou $\theta=\frac{\omega}{b}$, et la substitution faite en négligeant les puissances supérieures de $\omega$ donnera
\[ \Pi^{1}(n,c)=\frac{\frac{1}{2}\pi}{b\omega}+F^{1}(c)-\frac{1}{b^{2}}E^{1}(c). \]
Telle doit donc être la valeur de l'intégrale $\displaystyle\int\frac{d\varphi}{(\cos^{2}\varphi+\omega^{2}\sin^{2}\varphi)\Delta}$, lorsqu'on suppose $\omega$ infiniment petit et $\varphi=\frac{1}{2}\pi$.

Pour vérifier ce résultat par l'intégration directe, j'observe que $b$ étant la valeur de $\Delta$ lorsque $\varphi=\frac{1}{2}\pi$, on peut faire
\[ \Pi=\int\frac{d\varphi}{\cos^{2}\varphi+\omega^{2}\sin^{2}\varphi}\cdot\frac{1}{b}+\int\frac{d\varphi}{\cos^{2}\varphi+\omega^{2}\sin^{2}\varphi}\left(\frac{1}{\Delta}-\frac{1}{b}\right). \]
La première partie a pour intégrale $\frac{1}{b\omega}\operatorname{arc}\operatorname{tang}(\omega\operatorname{tang}\varphi)$, et en faisant $\varphi=\frac{1}{2}\pi$, elle devient $\frac{\frac{1}{2}\pi}{b\omega}$. La seconde partie se décompose en ces deux autres
\[ T=\int\frac{d\varphi}{\cos^{2}\varphi}\left(\frac{1}{\Delta}-\frac{1}{b}\right),\qquad V=\int\frac{\omega^{2}\sin^{2}\varphi\,d\varphi}{\cos^{2}\varphi(\cos^{2}\varphi+\omega^{2}\sin^{2}\varphi)}\left(\frac{1}{b}-\frac{1}{\Delta}\right). \]
Or en intégrant par parties on a $T=\operatorname{tang}\varphi\left(\frac{1}{\Delta}-\frac{1}{b}\right)-c^{2}\displaystyle\int\frac{d\varphi\sin^{2}\varphi}{\Delta^{3}}$,

\origpage{144}
expression dont le premier terme $\operatorname{tang}\varphi\left(\frac{1}{\Delta}-\frac{1}{b}\right)=\frac{(b-\Delta)\operatorname{tang}\varphi}{b\Delta}=-\frac{c^{2}\sin\varphi\cos\varphi}{b\Delta(b+\Delta)}$, s'évanouit lorsque $\varphi=\frac{1}{2}\pi$; donc on a simplement
\[ T=-\int\frac{c^{2}\sin^{2}\varphi\,d\varphi}{\Delta^{3}}=\int\frac{\Delta^{2}-1}{\Delta^{3}}\,d\varphi=F^{1}(c)-\frac{1}{b^{2}}E^{1}(c). \]

Nous avons déjà les deux premiers termes de la valeur de $\Pi^{1}(n,c)$; un troisième terme serait donné par l'intégrale $V$ qu'on peut mettre sous la forme
\[ V=\omega^{2}\int\frac{c^{2}\,d\varphi\sin^{2}\varphi}{(\cos^{2}\varphi+\omega^{2}\sin^{2}\varphi)(b+\Delta)b\Delta}; \]
et en procédant de la même manière, on trouve que le premier terme de $V$, développé suivant les puissances ascendantes de $\omega$, est en même temps le premier terme de l'intégrale plus simple $\omega^{2}\displaystyle\int\frac{c^{2}\,d\varphi\sin^{2}\varphi}{2b^{3}(\cos^{2}\varphi+\omega^{2}\sin^{2}\varphi)}$, lequel a pour expression $\frac{c^{2}\omega}{2b^{3}}\cdot\frac{\pi}{2}$. Réunissant toutes ces parties, on aura
\[ \Pi^{1}(n,c)=\frac{\pi}{2b\omega}+F^{1}(c)-\frac{1}{b^{2}}E^{1}(c)+\frac{\pi c^{2}}{4b^{3}}\omega+\text{etc.} \]

\section*{Troisième cas, $n=-c^{2}\sin^{2}\theta$.}

\subsection*{(104).}

La substitution de cette valeur de $n$ dans l'équation générale du n° 94, donne
\[ \begin{aligned} \frac{\cos\theta}{\sin\theta}\Delta(c,\theta)\Pi&=(F-E)\int\frac{d\theta}{\Delta(c,\theta)}-F\int\frac{d\theta}{\sin^{2}\theta\Delta(c,\theta)}\\ &\quad+\Delta\sin\varphi\cos\varphi\int\frac{c^{2}\,d\theta\sin^{2}\theta}{(1-c^{2}\sin^{2}\theta\sin^{2}\varphi)\Delta(c,\theta)}, \end{aligned} \]
Dans cette formule et dans toutes celles qui en seront déduites, les fonctions $\Delta$, $F$, $E$ étant toujours relatives au module $c$, nous n'y joindrons que la désignation de l'amplitude, c'est-à-dire, qu'au lieu de $\Delta(c,\theta)$, $F(c,\theta)$, $E(c,\theta)$, nous écrirons simplement $\Delta(\theta)$, $F(\theta)$, $E(\theta)$. Cela posé, on a par les formules connues,
\[ \int\frac{d\theta}{\sin^{2}\theta\Delta(\theta)}=F(\theta)-E(\theta)-\Delta(\theta)\cot\theta; \]
et si l'on fait $n'=-c^{2}\sin^{2}\varphi$, on aura
\[ \int\frac{c^{2}\sin^{2}\theta}{1-c^{2}\sin^{2}\varphi\sin^{2}\theta}\cdot\frac{d\theta}{\Delta(\theta)}=\frac{1}{\sin^{2}\varphi}[\Pi(n',\theta)-F(\theta)]. \]

\origpage{145}
Donc la formule générale pour le troisième cas, est
\[ \begin{aligned} \cot\theta\Delta(\theta)[\Pi(n,\varphi)-F(\varphi)]&=\cot\varphi\Delta(\varphi)[\Pi(n',\theta)-F(\theta)]\\ &\quad+E(\theta)F(\varphi)-F(\theta)E(\varphi)\ldots(n'). \end{aligned} \]
On n'a pas ajouté de constante, parce que la forme de l'équation fait voir que si on change $\varphi$ en $\theta$ et réciproquement, la constante devrait changer de signe, ce qui prouve qu'elle est nulle.

\subsection*{(105).}

Si dans la formule précédente on fait $\varphi=\frac{1}{2}\pi$, on aura immédiatement
\[ \Pi^{1}(n)=F^{1}+\frac{\operatorname{tang}\theta}{\Delta(\theta)}[F^{1}E(\theta)-E^{1}F(\theta)]\ldots(p'). \]
Ainsi dans le troisième cas, comme dans les deux premiers, la fonction complète de troisième espèce $\Pi^{1}(n)$, s'exprimera toujours par des fonctions de la première et de la seconde espèce.

Il en est de même de la fonction non-complète $\Pi(n,c,\varphi)$ ou $\Pi(n,\varphi)$, si l'amplitude $\varphi$ est telle qu'on ait $F(\varphi)=kF^{1}$, $k$ étant un nombre rationnel; car alors on aurait $\Pi(n,\varphi)=k\Pi^{1}(n)+W$, $W$ étant une quantité déterminable par logarithmes.

Il est remarquable que les formules $(n')$ et $(p')$ qu'on vient de trouver pour le troisième cas, sont plus simples que les formules analogues pour le premier et le second cas, puisqu'elles ne contiennent point les fonctions $F$ et $E$ relatives au module complémentaire $b$.

Remarquons encore que la valeur de $\Pi^{1}(n)$ serait indéterminée si on avait $\theta=\frac{1}{2}\pi$ ou $n=-c^{2}$. Mais alors on a généralement pour toute valeur de $\varphi$,
\[ \Pi(-c^{2},\varphi)=\frac{1}{b^{2}}E(\varphi)-\frac{c^{2}\sin\varphi\cos\varphi}{b^{2}\Delta(\varphi)}; \]
et par conséquent lorsque $\varphi=\frac{1}{2}\pi$, on a
\[ \Pi^{1}(-c^{2})=\frac{1}{b^{2}}E^{1}. \]

\subsection*{(106).}

Revenons à la formule générale $(n')$, et supposons que l'angle $\theta$, déduit du paramètre $n=-c^{2}\sin^{2}\theta$, soit tel qu'on ait $F(\theta)=kF^{1}$, $k$ étant rationnel; on aura alors $E(\theta)=kE^{1}+V$, et

\origpage{146}
$\Pi(n',\theta)=k\Pi^{1}(n')+W$, $V$ étant une quantité déterminable algébriquement, et $W$ étant déterminable par logarithmes. Ces valeurs et celle de la fonction $\Pi^{1}(n')$ tirée de la formule $(p')$ étant substituées dans l'équation $(n')$, il en résultera
\[ \Pi(n,\varphi)=\left(1+\frac{V\operatorname{tang}\theta}{\Delta(\theta)}\right)F(\varphi)+\frac{W\cot\varphi\Delta(\varphi)}{\cot\theta\Delta(\theta)}. \]
D'où il suit que la fonction $\Pi(n,\varphi)$ pourra être ramenée indéfiniment aux fonctions de la première espèce, si le paramètre satisfait à la condition mentionnée. Nous avons déja donné (n° 86) un symptôme semblable de réduction, mais celui qu'on vient d'indiquer est beaucoup plus général.

Soit, par exemple, $n=-1+b=-c^{2}\sin^{2}\theta$, on aura $\sin^{2}\theta=\frac{1}{1+b}$; ce qui donne $F(\theta)=\frac{1}{2}F^{1}$. On a alors par les formules du n° 57,
\[ E(\theta)=\tfrac{1}{2}E^{1}+\tfrac{1}{2}(1-b), \]
et
\[ \Pi(n',\theta)=\tfrac{1}{2}\Pi^{1}(n')+\frac{\sin\varphi}{4\Delta\cos\varphi}\log\left(\frac{\Delta-(1-b)\sin\varphi\cos\varphi}{\Delta+(1-b)\sin\varphi\cos\varphi}\right); \]
d'où résulte
\[ \Pi(n,\varphi)=\frac{1+b}{2b}F(\varphi)+\frac{1}{4b}\log\left(\frac{\Delta-(1-b)\sin\varphi\cos\varphi}{\Delta+(1-b)\sin\varphi\cos\varphi}\right), \]
ce qui s'accorde avec la formule du n° 52.

Nous terminerons ces recherches par une observation générale; c'est que toutes les fois que la fonction de troisième espèce $\Pi$ est réductible indéfiniment aux espèces inférieures, elle est toujours réductible à la première espèce, c'est-à-dire, qu'elle s'exprime par la seule fonction $F(c,\varphi)$. Il n'y a d'exception à cette règle générale que lorsque $n=-c^{2}$, et lorsque $n=-1$, seuls cas où la fonction de seconde espèce $E(\varphi)$ entre dans l'expression de $\Pi$, ainsi qu'on le voit par les formules du n° 48.

\origpage{147}
\section*{Réduction générale des fonctions elliptiques dont le paramètre est imaginaire.}

\subsection*{(107).}

Soit $p=\frac{\sin\varphi\cos\varphi}{(1+\zeta\sin^{2}\varphi)\Delta}$, $\zeta$ étant un coefficient constant, on aura par la différentiation et en introduisant un second coefficient $k$,
\[ \frac{dp}{1+kp^{2}}=\frac{d\varphi}{\Delta}\cdot\frac{1-(2+\zeta)\sin^{2}\varphi+(1+2\zeta)c^{2}\sin^{4}\varphi-\zeta c^{2}\sin^{6}\varphi}{(1+\zeta\sin^{2}\varphi)^{2}(1-c^{2}\sin^{2}\varphi)+k\sin^{2}\varphi(1-\sin^{2}\varphi)}. \]
Supposons que le dénominateur du second membre soit égal au produit de ces trois facteurs:
\[ (1+n\sin^{2}\varphi)(1+n'\sin^{2}\varphi)(1-m\sin^{2}\varphi); \]
il faudra satisfaire aux trois équations
\[ \begin{gathered} nn'm=c^{2}\zeta^{2}\\ (1+n)(1+n')(1-m)=b^{2}(1+\zeta)^{2}\\ (n+c^{2})(n'+c^{2})(m-c^{2})=b^{2}c^{2}k. \end{gathered} \]
Pour cela, supposons que $n$ et $n'$ soient connus, il faudra par leur moyen déterminer les trois autres quantités $\zeta$, $m$, et $k$. Et d'abord l'équation qui détermine $\zeta$ étant
\[ 1=\frac{c^{2}\zeta^{2}}{nn'}+\frac{b^{2}(1+\zeta)^{2}}{(1+n)(1+n')}; \]
si l'on fait pour abréger $M=b^{2}+c^{2}\left(\frac{1+n}{n}\right)\left(\frac{1+n'}{n'}\right)$, on aura
\[ M\zeta=-b^{2}+\sqrt{\left[\frac{1+n}{n'}\cdot\frac{1+n'}{n'}(c^{2}+n)(c^{2}+n')\right]}. \]
\ednote{Printed $\frac{1+n}{n'}$ as the first factor under the radical instead of $\frac{1+n}{n}$; an apparent misprint, since solving the preceding equation for $\zeta$ with $M$ as defined gives the radicand $\frac{1+n}{n}\cdot\frac{1+n'}{n'}(c^{2}+n)(c^{2}+n')$.}
$\zeta$ étant connu, on aura $m$ et $k$ par les équations
\[ \begin{aligned} m&=\frac{c^{2}\zeta^{2}}{nn'}\\ k&=(n+c^{2})(n'+c^{2})\cdot\frac{m-c^{2}}{b^{2}c^{2}}. \end{aligned} \]
Cela posé; on peut donner à l'équation différentielle la forme
\[ \frac{dp}{1+kp^{2}}=\frac{d\varphi}{\Delta}\left[\frac{A}{1+n\sin^{2}\varphi}+\frac{A'}{1+n'\sin^{2}\varphi}+\frac{B}{1-m\sin^{2}\varphi}+\frac{1}{\zeta}\right], \]

\origpage{148}
et les coefficiens $A$, $A'$, $B$ se détermineront de cette manière:
\[ \begin{aligned} A&=\frac{n^{3}+(2+\zeta)n^{2}+(1+2\zeta)c^{2}n+\zeta c^{2}}{n'(n-n')(n+m)}\\ A'&=\frac{n'^{3}+(2+\zeta)n'^{2}+(1+2\zeta)c^{2}n'+\zeta c^{2}}{n(n'-n)(n'+m)}\\ B&=\frac{m^{3}-(2+\zeta)m^{2}+(1+2\zeta)c^{2}m-\zeta c^{2}}{m(m+n)(m+n')}. \end{aligned} \]
\ednote{The denominators of $A$ and $A'$ are printed with $n'$ and $n$ as first factors; an apparent misprint for $n(n-n')(n+m)$ and $n'(n'-n)(n'+m)$: only these give $A+A'=1-\frac{1}{\zeta}-B$, as the value of $P$ below requires.}
On aura ensuite l'intégrale
\[ A\Pi(n)+A'\Pi(n')+B\Pi(-m)+\frac{1}{\zeta}F=\int\frac{dp}{1+kp^{2}}. \]

\subsection*{(108).}

Ce résultat a lieu quels que soient $n$ et $n'$. Supposons donc que ces paramètres sont imaginaires et qu'on a
\[ \begin{aligned} n&=\nu(\cos\theta+\sqrt{-1}\sin\theta)\\ n'&=\nu(\cos\theta-\sqrt{-1}\sin\theta). \end{aligned} \]
D'après ces valeurs, si on fait pour abréger $h=\frac{b^{2}\nu}{\nu^{2}+2c^{2}\nu\cos\theta+c^{2}}$\ednote{The exponent of the last $c$ is partly covered by a reader's ink stroke; it is read as $c^{2}$, which is also what the equation for $\zeta$ on p.~147 requires (with $b^{2}+c^{2}=1$).}, on trouvera
\[ \begin{aligned} m&=\frac{c^{2}\zeta^{2}}{\nu^{2}}\\ \zeta&=-\nu h+\nu\sqrt{(1+2h\cos\theta+h^{2})}\\ k&=(c^{4}+2c^{2}\nu\cos\theta+\nu^{2})\frac{m-c^{2}}{b^{2}c^{2}}; \end{aligned} \]
d'où l'on voit que les trois quantités $\zeta$, $m$, $k$ sont toujours réelles; il en sera de même du coefficient $B$, qu'on peut mettre sous la forme
\[ B=\frac{m^{3}-(2+\zeta)m^{2}+(1+2\zeta)c^{2}m-\zeta c^{2}}{m^{3}+2m^{2}\nu\cos\theta+m\nu^{2}}. \]
Quant aux coefficiens $A$ et $A'$, ils sont nécessairement imaginaires. Pour avoir plus facilement leur expression, soit $A=\frac{1}{2}(P+Q\sqrt{-1})$, $A'=\frac{1}{2}(P-Q\sqrt{-1})$, on trouvera que les quantités $P$ et $Q$ sont déterminées par les équations
\[ \begin{aligned} P&=1-\frac{1}{\zeta}-B\\ Q&=\frac{m(1-B)-2-\zeta}{\nu\sin\theta}-\frac{\cos\theta}{\sin\theta}\left(1+\frac{1}{\zeta}+\frac{1}{B}\right). \end{aligned} \]
\ednote{Printed $\frac{1}{B}$ as the last term in the bracket instead of $B$; an apparent misprint. The corresponding formula for $Q'$ on p. 149 has $+B'$, and only $+B$ reproduces the values $Q=-\frac{2}{c\sin\theta}(1+c\cos\theta)$ and $Q'=0$ given in art.\ 110.}

\origpage{149}
Tous les coefficiens étant ainsi connus, on aura enfin l'intégrale
\[ \begin{aligned} \frac{P+Q\sqrt{-1}}{2}\Pi(n)+\frac{P-Q\sqrt{-1}}{2}\Pi(n')&=-B\Pi(-m)-\frac{1}{\zeta}F\\ &\quad+\frac{1}{\sqrt{k}}\operatorname{arc}\operatorname{tang}.\frac{\sqrt{k}\sin\varphi\cos\varphi}{(1+\zeta\sin^{2}\varphi)\Delta}\ldots(q'). \end{aligned} \]

Cette intégrale ne suffit pas pour déterminer les deux fonctions $\Pi(n)$, $\Pi(n')$; mais si on fait attention que le radical contenu dans la valeur de $\zeta$ peut être pris avec le signe $-$, et qu'ainsi il y a une seconde valeur de $\zeta$, savoir,
\[ \zeta'=-\nu h-\nu\sqrt{(1+2h\cos\theta+h^{2})}, \]
on verra qu'il doit en résulter la seconde intégrale dont nous avons besoin.

Si on désigne les nouvelles valeurs de $m$, $k$, $B$, $P$, $Q$ par les mêmes lettres accentuées, on aura donc
\[ \begin{aligned} m'&=\frac{c^{2}\zeta'^{2}}{\nu^{2}}\\ k'&=(c^{4}+2c^{2}\nu\cos\theta+\nu^{2})\cdot\frac{m'-c^{2}}{b^{2}c^{2}}\\ B'&=\frac{m'^{3}-(2+\zeta')m'^{2}+(1+2\zeta')c^{2}m'-c^{2}\zeta'}{m'^{3}+2m'^{2}\nu\cos\theta+m'\nu^{2}}\\ P'&=1-\frac{1}{\zeta'}-B'\\ Q'&=\frac{m'(1-B')-2-\zeta'}{\nu\sin\theta}-\frac{\cos\theta}{\sin\theta}\left(1+\frac{1}{\zeta'}+B'\right), \end{aligned} \]
et la seconde intégrale sera
\[ \begin{aligned} \frac{P'+Q'\sqrt{-1}}{2}\Pi(n)+\frac{P'-Q'\sqrt{-1}}{2}\Pi(n')&=-B'\Pi(-m')-\frac{1}{\zeta'}F\ldots(r').\\ &\quad+\frac{1}{\sqrt{k'}}\operatorname{arc}\operatorname{tang}\frac{\sqrt{k'}\sin\varphi\cos\varphi}{(1+\zeta'\sin^{2}\varphi)\Delta} \end{aligned} \]
Il est visible maintenant que ces deux intégrales donnent les valeurs des fonctions $\Pi(n)$, $\Pi(n')$, exprimées chacune au moyen des fonctions $\Pi(-m)$ et $\Pi(-m')$. Donc en général \emph{toute fonction de troisième espèce dont le paramètre est imaginaire, peut s'exprimer par deux fonctions de la même espèce, dont les paramètres sont réels.}

Ce théorème résout pleinement la difficulté qu'on aurait pu élever

\origpage{150}
sur la division des fonctions elliptiques en trois espèces, s'il n'eût pas été constaté que les fonctions dont le paramètre est imaginaire, peuvent toujours se réduire à celles dont le paramètre est réel.

\subsection*{(109).}

Il faut maintenant entrer dans quelques détails sur les résultats de la solution précédente.

J'observe d'abord que les paramètres $-m$, $-m'$ sont tous deux plus petits que l'unité, puisque d'après l'équation qui détermine $\zeta$, on a
\[ 1-m=\frac{b^{2}(1+\zeta)^{2}}{1+2\nu\cos\theta+\nu^{2}}. \]
En second lieu, si on substitue la valeur de $\zeta$ dans l'équation $m=\frac{c^{2}\zeta^{2}}{\nu^{2}}$, on aura
\[ m-c^{2}=2c^{2}h[h+\cos\theta-\sqrt{(1+2h\cos\theta+h^{2})}]; \]
on aura semblablement
\[ m'-c^{2}=2c^{2}h[h+\cos\theta+\sqrt{(1+2h\cos\theta+h^{2})}]; \]
et le produit de ces équations donne
\[ (m-c^{2})(m'-c^{2})=-4c^{4}h^{2}\sin^{2}\theta: \]
donc on a toujours $m<c^{2}$ et $m'>c^{2}$. Ainsi des deux paramètres $-m$, $-m'$, l'un appartient à la forme $-c^{2}\sin^{2}\theta$, et l'autre à la forme $-1+b^{2}\sin^{2}\theta$, de sorte qu'on peut faire $m=c^{2}\sin^{2}\lambda$ et $m'=1-b^{2}\sin^{2}\mu$; on peut aussi déterminer directement les angles $\lambda$ et $\mu$ par les formules
\[ \begin{aligned} \sin\lambda&=-h+\sqrt{(1+2h\cos\theta+h^{2})}\\ \cos\mu&=\frac{2ch\sin\theta}{b\cos\lambda}. \end{aligned} \]
Les valeurs correspondantes de $k$ et $k'$ étant
\[ \begin{aligned} k&=(c^{4}+2c^{2}\nu\cos\theta+\nu^{2})\left(\frac{m-c^{2}}{b^{2}c^{2}}\right)\\ k'&=(c^{4}+2c^{2}\nu\cos\theta+\nu^{2})\left(\frac{m'-c^{2}}{b^{2}c^{2}}\right), \end{aligned} \]
on voit que la première est négative et la seconde positive; d'où il

\origpage{151}
suit que l'expression réduite de toute fonction elliptique dont le paramètre est imaginaire, contiendra toujours un arc de cercle et un logarithme concurremment avec les deux fonctions $\Pi(-m)$, $\Pi(-m')$ et la fonction de première espèce $F$.

\subsection*{(110).}

Les formules générales que nous venons de développer ne sont sujettes à aucune exception, mais il ne sera pas inutile d'en faire l'application à quelques cas qui présentent des réductions remarquables.

Soit d'abord $\nu=c$, on aura dans ce cas $\zeta'=-1$, $m'=1$, $B'=0$, $Q'=0$, $P'=2$, $k'=1+2c\cos\theta+c^{2}$, et alors la formule $(r')$ devient
\[ \Pi(n)+\Pi(n')=F+\frac{1}{\sqrt{k'}}\operatorname{arc}\operatorname{tang}.\frac{\sqrt{k'}\operatorname{tang}\varphi}{\Delta}. \]
Ce résultat se déduirait immédiatement de la formule du n° 46, puisqu'on a $nn'=c^{2}$.

Les données pour former l'autre intégrale se trouvent par notre analyse comme il suit:
\[ \begin{aligned} \zeta&=\frac{c^{2}+c\cos\theta}{1+c\cos\theta}\\ m&=\left(\frac{c^{2}+c\cos\theta}{1+c\cos\theta}\right)^{2}\\ k&=-\frac{c^{2}\sin^{2}\theta(1+2c\cos\theta+c^{2})}{(1+c\cos\theta)^{2}}\\ B&=1-\frac{1}{\zeta}=-\frac{b^{2}}{c(c+\cos\theta)}\\ P&=0\\ Q&=-\frac{2}{c\sin\theta}(1+c\cos\theta). \end{aligned} \]
Substituant toutes ces valeurs dans l'équation $(q')$ et faisant pour abréger $\varepsilon=\frac{c\sin\theta\sqrt{(1+2c\cos\theta+c^{2})}}{1+c\cos\theta}$, on aura
\[ \begin{aligned} \frac{1}{\sqrt{-1}}[\Pi(n)-\Pi(n')]&=\frac{b^{2}\sin\theta}{(1+c\cos\theta)(c+\cos\theta)}\Pi(-m)-\frac{\sin\theta}{c+\cos\theta}F\\ &\quad+\frac{1}{2\varepsilon}\log\left(\frac{(1+\zeta\sin^{2}\varphi)\Delta+\varepsilon\sin\varphi\cos\varphi}{(1+\zeta\sin^{2}\varphi)\Delta-\varepsilon\sin\varphi\cos\varphi}\right). \end{aligned} \]
Ces deux équations donnent la valeur de $\Pi(n)$ et celle de $\Pi(n')$

\origpage{152}
dans le cas assez général où l'on a $n=c(\cos\theta+\sqrt{-1}\sin\theta)$ et $n'=c(\cos\theta-\sqrt{-1}\sin\theta)$.

Si cependant on avait $\cos\theta=-c$, il y aurait un changement à faire à la seconde équation; alors on aurait $m=0$, et les deux termes $-B\Pi(-m)-\frac{1}{\zeta}F$ de l'équation $(q')$, se réduiraient au seul $-\left(B+\frac{1}{\zeta}\right)F$, ou simplement $-F$; de sorte que l'équation dont il s'agit deviendrait
\[ \Pi(n)-\Pi(n')=-\frac{c\sqrt{-1}}{b}F+\frac{\sqrt{-1}}{2b}\log\left(\frac{\Delta+c\sin\varphi\cos\varphi}{\Delta-c\sin\varphi\cos\varphi}\right); \]
on aurait en même temps
\[ \Pi(n)+\Pi(n')=F+\frac{1}{b}\operatorname{arc}\operatorname{tang}\frac{b\operatorname{tang}\varphi}{\Delta}. \]
Ainsi dans le cas de $n=-c^{2}+bc\sqrt{-1}$, la fonction $\Pi(n)$ se réduit indéfiniment à la première espèce. Mais quoique ce cas soit peut-être le plus simple de ceux où le paramètre est imagiuaire\ednote{Printed ``imagiuaire'' with a turned letter instead of ``imaginaire''; an apparent misprint.}, on voit néanmoins que l'expression de la fonction $\Pi(n)$ contient à-la-fois un arc de cercle et un logarithme.

\subsection*{(111).}

Un autre cas mérite d'être examiné avec soin; c'est celui où l'on a
\[ n=-1+b(\cos\lambda+\sqrt{-1}\sin\lambda). \]
Cette valeur donne $\nu\cos\theta=-1+b\cos\lambda$, $\nu\sin\theta=b\sin\lambda$, d'où résulte
\[ \begin{aligned} \nu^{2}&=1-2b\cos\lambda+b^{2}\\ h&=-\frac{1}{2\cos\theta}=\frac{\frac{1}{2}\nu}{1-b\cos\lambda}; \end{aligned} \]
On a donc $\zeta=-\nu h\pm\nu h$, c'est-à-dire que les deux valeurs de $\zeta$ sont $\zeta=0$, $\zeta'=-2\nu h$.

La valeur $\zeta=0$ donne $m=0$, ce qui fait la même difficulté que dans le dernier cas de l'article précédent. Elle se résoudra de la même manière; car ayant alors $\Pi(-m)=-F$\ednote{Printed clearly as $\Pi(-m)=-F$; with $m=0$ one has $\Pi(-m)=F$, which the reduction that follows requires. An apparent misprint.}, les deux termes $-B\Pi(-m)-\frac{1}{\zeta}F$ se réduisent à $-\left(B+\frac{1}{\zeta}\right)F$; or en supposant d'abord $\zeta$ infiniment petit et ensuite nul, on trouve

\origpage{153}
$B+\frac{1}{\zeta}=\frac{c^{2}}{\nu^{2}}$. D'ailleurs les substitutions donnent sans difficulté
\[ \begin{aligned} P&=\frac{2b}{\nu^{2}}(b-\cos\lambda)\\ Q&=-\frac{2b\sin\lambda}{\nu^{2}}\\ k&=-\nu^{2}; \end{aligned} \]
on aura donc pour la première équation
\[ \begin{aligned} &(b-\cos\lambda-\sqrt{-1}\sin\lambda)\Pi(n)+(b-\cos\lambda+\sqrt{-1}\sin\lambda)\Pi(n')\\ &\quad=-\frac{c^{2}}{b}F+\frac{\nu}{2b}\log\left(\frac{\Delta+\nu\sin\varphi\cos\varphi}{\Delta-\nu\sin\varphi\cos\varphi}\right); \end{aligned} \]
et on peut remarquer que cette équation se déduirait immédiatement de la formule du n° 51, en faisant $n=-1+b\cos\lambda+b\sqrt{-1}\sin\lambda$; car alors de l'équation de condition $(1+n)(1-m)=b^{2}$, on déduirait $1-m=b(\cos\lambda-\sqrt{-1}\sin\lambda)$, et par conséquent $-m=n'$.

Pour avoir la seconde équation, il faut continuer de faire les substitutions dans les valeurs des coefficiens: elles donnent après beaucoup de réductions, les résultats suivans:
\[ \begin{aligned} \zeta'&=-\frac{\nu^{2}}{1-b\cos\lambda}\\ m'&=\frac{c^{2}\nu^{2}}{(1-b\cos\lambda)^{2}}\\ k'&=\frac{b^{2}\nu^{2}\sin^{2}\lambda}{(1-b\cos\lambda)^{2}}\\ B'&=\frac{b}{\nu^{2}}(b-\cos\lambda)\\ \tfrac{1}{2}P'&=\frac{1-b\cos\lambda}{\nu^{2}}\\ \tfrac{1}{2}Q'&=\frac{(b-\cos\lambda)\cdot(1-b\cos\lambda)}{\nu^{2}\sin\lambda}, \end{aligned} \]
et la seconde équation devient
\[ \begin{aligned} &[\sin\lambda+\sqrt{-1}(b-\cos\lambda)]\Pi(n)+[\sin\lambda-\sqrt{-1}(b-\cos\lambda)]\Pi(n')\\ &\quad=-\frac{b\sin\lambda(b-\cos\lambda)}{1-b\cos\lambda}\Pi(-m')+\sin\lambda\cdot F+\frac{\nu}{b}\operatorname{arc}\operatorname{tang}\frac{b\nu\sin\lambda\sin\varphi\cos\varphi}{(1-b\cos\lambda-\nu^{2}\sin^{2}\varphi)\Delta}. \end{aligned} \]

\origpage{154}
Ainsi on connaîtra les fonctions $\Pi(n)$, $\Pi(n')$ par le moyen de la seule fonction de troisième espèce $\Pi(-m')$.

Si on avait $\cos\lambda=b$, les formules trouvées s'accorderaient avec celles qu'on a données dans l'article précédent pour le cas de $\cos\theta=-c$; et en effet dans les deux cas on a $n=-c^{2}+bc\sqrt{-1}$.

L'application de la formule précédente exige qu'on distingue soigneusement deux cas, selon que $\cos\lambda$ est $>b$ ou $<b$, afin d'évaluer convenablement l'arc de cercle $Z$ dont la tangente est
\[ \frac{\nu b\sin\lambda\sin\varphi\cos\varphi}{(1-b\cos\lambda-\nu^{2}\sin^{2}\varphi)\Delta}. \]

Soit 1°. $\cos\lambda>b$, le dénominateur $1-b\cos\lambda-\nu^{2}\sin^{2}\varphi$ sera positif pour toute valeur de $\varphi$, puisqu'en faisant $\sin\varphi=1$, il se réduit à $b(\cos\lambda-b)$, quantité positive. Alors l'arc $Z$ n'augmente que jusqu'à un certain terme qui est son \emph{maximum}, après quoi il diminue, devient zéro, puis négatif. Ainsi un mobile qui décrirait l'arc $Z$ pendant que $\varphi$ croît uniformément, ferait des oscillations autour du point de départ où $\varphi=0$, et ne s'en écarterait de part et d'autre que d'une quantité déterminée par la valeur $\operatorname{tang}\varphi=\sqrt[4]{\left(\frac{1-b\cos\lambda}{b^{3}\cos\lambda-b^{4}}\right)}$. L'arc $Z$ s'évanouira donc lorsqu'on aura $\varphi=\frac{1}{2}\pi$, $\varphi=\pi$, etc.

Soit 2°. $\cos\lambda<b$, alors le dénominateur $1-b\cos\lambda-\nu^{2}\sin^{2}\varphi$ s'évanouit lorsque $\sin\varphi=\frac{1-b\cos\lambda}{\nu^{2}}$, cas où l'on a $\cos\varphi=\frac{b(b-\cos\lambda)}{\nu^{2}}$\ednote{Printed $\sin\varphi$ and $\cos\varphi$ without the squares; the vanishing of $1-b\cos\lambda-\nu^{2}\sin^{2}\varphi$ gives $\sin^{2}\varphi=\frac{1-b\cos\lambda}{\nu^{2}}$ and $\cos^{2}\varphi=\frac{b(b-\cos\lambda)}{\nu^{2}}$. An apparent misprint.}, de sorte que cette valeur de $\varphi$ est réelle et $<\frac{1}{2}\pi$. Dans ce point l'arc $Z$ qui a une tangente infinie, est égal à $\frac{1}{2}\pi$: il augmente ensuite continuellement avec l'angle $\varphi$, et lorsque $\varphi=\frac{1}{2}\pi$, on a $Z=\pi$.

On commettrait donc une erreur dans l'application de la formule, si en faisant $\varphi=\frac{1}{2}\pi$, on prenait $Z=0$. La valeur $Z=0$ n'a lieu que lorsque $\cos\lambda>b$, mais si l'on a $\cos\lambda<b$, il faudra prendre $Z=\pi$.

\subsection*{(112).}

Il est très-remarquable que dans le cas d'un paramètre imaginaire de la forme $n=-1+b(\cos\lambda+\sqrt{-1}\sin\lambda)$, la transformation dont on a fait usage pour les approximations (art.\ 84), conduit immédiatement à la réduction de la fonction $\Pi(n)$.

\origpage{155}
En effet, d'après cette valeur de $n$, les formules de l'article cité donnent
\[ n^{0}=-\frac{\nu^{2}}{(1+b)^{2}}; \]
et puisque le nouveau paramètre $n^{0}$ est réel, la solution est donnée immédiatement par la première transformée
\[ \Pi(n)=\frac{1+c^{0}}{2}\left[p\Pi^{0}+qF^{0}+\frac{\tfrac{1}{2}n}{1+n}\int\frac{d\varphi^{0}\cos\varphi^{0}}{1+n^{0}\sin^{2}\varphi^{0}}\right], \]
où l'on a
\[ \begin{aligned} p&=\frac{\sin\lambda}{\nu^{2}}[\sin\lambda+\sqrt{-1}(\cos\lambda-b)]\\ q&=\frac{c^{2}}{2b\nu^{2}}[\cos\lambda-b-\sqrt{-1}\sin\lambda]\\ \frac{\tfrac{1}{2}n}{1+n}&=\frac{1}{2b}[b-\cos\lambda+\sqrt{-1}\sin\lambda]. \end{aligned} \]
La valeur de $\Pi(n)$ exprimée en fonction de $n^{0}$, $c^{0}$, $\varphi^{0}$, est donc
\[ \begin{aligned} \Pi(n)={}&\frac{1-b}{2b\nu^{2}}(\cos\lambda-b-\sqrt{-1}\sin\lambda)F(c^{0},\varphi^{0})\\ &+\frac{\sin\lambda}{(1+b)\nu^{2}}(\sin\lambda+(\cos\lambda-b)\sqrt{-1})\Pi(n^{0},c^{0},\varphi^{0})\\ &+\frac{b-\cos\lambda+\sqrt{-1}\sin\lambda}{4b\nu}\log\left(\frac{1+b+\nu\sin\varphi^{0}}{1+b-\nu\sin\varphi^{0}}\right). \end{aligned} \]
D'où l'on voit que la fonction $\Pi(n)$ dont le paramètre est imaginaire, se réduit immédiatement à la fonction $\Pi(n^{0},c^{0},\varphi^{0})$ dont le paramètre est réel, et qui se rapporte au second cas de l'article 50, puisqu'on a $n^{0}=-1+b^{02}\cos^{2}\frac{1}{2}\lambda$.

Connaissant la valeur de $\Pi(n)$, on aura en même temps celle de $\Pi(n')$ en changeant dans la formule précédente le signe de $\sqrt{-1}$. Si l'on veut comparer ces deux solutions, et qu'à cet effet on substitue les valeurs qui viennent d'être trouvées pour $\Pi(n)$ et $\Pi(n')$ dans la première formule du numéro précédent, on verra aisément, d'après les relations connues entre les amplitudes $\varphi$ et $\varphi^{0}$, que cette équation devient identique.

Il n'en sera pas de même si les substitutions sont faites dans la seconde équation. On tombe alors sur une équation entre les fonctions $\Pi(n^{0},c^{0},\varphi^{0})$, $\Pi(-m',c,\varphi)$, qui est vraie sans doute, mais qui

\origpage{156}
n'est pas identique. Cette équation est
\[ \begin{aligned} \Pi(n^{0},c^{0},\varphi^{0})={}&-\frac{\tfrac{1}{2}b(1+b)(b-\cos\lambda)}{1-b\cos\lambda}\Pi(-m',c,\varphi)+\frac{1+b}{2}F(c,\varphi)\\ &+\frac{(1+b)\nu}{2b\sin\lambda}\operatorname{arc}\operatorname{tang}\frac{b\nu\sin\lambda\sin\varphi\cos\varphi}{(1-b\cos\lambda-\nu^{2}\sin^{2}\varphi)\Delta}. \end{aligned} \]
Pour la vérifier, au moins dans un cas particulier, soit $\cos\lambda=b$, on aura $\sin\lambda=c$, $\nu=c$, et cette équation devient
\[ \Pi(n^{0})=\frac{1+b}{2}F(c,\varphi)+\frac{1+b}{2b}\operatorname{arc}\operatorname{tang}\frac{b\operatorname{tang}\varphi}{\Delta}: \]
mais on a dans ce même cas $n^{0}=-\frac{c^{2}}{(1+b)^{2}}=-\left(\frac{1-b}{1+b}\right)=-c^{0}$; donc la formule du n° 46 s'applique à la fonction $\Pi(n^{0})$ et donne
\[ \Pi(n^{0})=\tfrac{1}{2}F^{0}+\frac{1}{2(1-c^{0})}\operatorname{arc}\operatorname{tang}\frac{(1-c^{0})\operatorname{tang}\varphi^{0}}{\Delta^{0}}, \]
on a de plus $F(c,\varphi)=\frac{1+c^{0}}{2}F^{0}=\frac{1}{1+b}F^{0}$; donc pour que les deux valeurs de $\Pi(n^{0})$ s'accordent entre elles, il faut qu'on ait
\[ \operatorname{arc}\operatorname{tang}\left(\frac{b\operatorname{tang}\varphi}{\Delta}\right)=\tfrac{1}{2}\operatorname{arc}\operatorname{tang}\frac{(1-c^{0})\operatorname{tang}\varphi^{0}}{\Delta^{0}}, \]
équation qui a lieu en effet d'après les relations connues entre $\varphi$ et $\varphi^{0}$ (art.\ 58).

\subsection*{(113).}

Pour comparer encore mieux nos deux solutions, cherchons par l'une et l'autre méthode, la valeur de la fonction complète $\Pi^{1}(n)$. Dans ce cas on a $\varphi=\frac{1}{2}\pi$, $\varphi^{0}=\pi$, et comme en général $F(c,\varphi)=\frac{1+c^{0}}{2}F(c^{0},\varphi^{0})$, il faudra faire $F(c^{0},\varphi^{0})=\frac{2}{1+c^{0}}F^{1}(c)=(1+b)F^{1}(c)$. Substituant ces valeurs dans la formule du numéro précédent, on aura
\[ \begin{aligned} \Pi^{1}(n)={}&\frac{c^{2}}{2b\nu^{2}}(\cos\lambda-b-\sqrt{-1}\sin\lambda)F^{1}(c)\\ &+\frac{\sin\lambda}{(1+b)\nu^{2}}(\sin\lambda+(\cos\lambda-b)\sqrt{-1})\Pi(n^{0},c^{0},\pi \end{aligned} \]
\ednote{The closing parenthesis after $\pi$ is not printed at the end of the line; read $\Pi(n^{0},c^{0},\pi)$.}
Il s'agit maintenant d'avoir la valeur de $\Pi(n^{0},c^{0},\pi)$ qui est la même

\origpage{157}
que $2\Pi^{1}(n^{0},c^{0})$; celle-ci se trouvera par la formule de l'art.\ 101, parce que $n^{0}$ peut être mis sous la forme $n^{0}=-1+b^{02}\sin\theta^{2}$, en faisant $\theta=\frac{1}{2}(\pi-\lambda)$; on aura donc
\[ \begin{aligned} \Pi^{1}(n^{0},c^{0})=F^{1}(c^{0})+\frac{\Delta(b^{0},\theta)}{b^{02}\sin\theta\cos\theta}\Big[&\tfrac{1}{2}\pi+F^{1}(c^{0})F(b^{0},\theta)-E^{1}(c^{0})F(b^{0},\theta)\\ &-F^{1}(c^{0})E(b^{0},\theta)\Big]. \end{aligned} \]
Or par les formules connues on a $F^{1}(c^{0})=\frac{1+b}{2}F^{1}(c)$, $E^{1}(c^{0})=\frac{b}{1+b}F^{1}(c)+\frac{1}{1+b}E^{1}(c)$; d'un autre côté on a aussi $\Delta(b^{0},\theta)=\frac{\nu}{1+b}$, $b^{02}=\frac{4b}{(1+b)^{2}}$, $\frac{\Delta(b^{0},\theta)}{b^{02}\sin\theta\cos\theta}=\frac{\nu(1+b)}{2b\sin\lambda}$; donc
\[ \begin{aligned} \Pi^{1}(n^{0},c^{0})=\frac{1+b}{2}F^{1}(c)+\frac{(1+b)\nu}{2b\sin\lambda}\Big[&\tfrac{1}{2}\pi+\frac{1+b^{2}}{2+2b}F^{1}(c)F(b^{0},\theta)\\ &-\frac{1}{1+b}E^{1}(c)F(b^{0},\theta)-\left(\frac{1+b}{2}\right)F^{1}(c)E(b^{0},\theta)\Big]. \end{aligned} \]
Il faut maintenant ramener les fonctions $F(b^{0},\theta)$, $E(b^{0},\theta)$ à un module plus petit qui se déduira de $b^{0}$ suivant la même loi que $c^{0}$ se réduit de $c$. Mais, si l'on n'y prenait garde, la notation pourrait induire en erreur. En effet, la quantité $b^{0}$ n'entre dans la formule précédente que comme complément du module $c^{0}$, et non comme une quantité déduite de $b$ par la même loi que $c^{0}$ est déduit de $c$. Cette loi exige que $c^{0}$ soit plus petit que $c$, au lieu que $b^{0}$, complément de $c^{0}$, est plus grand que $b$.

Pour éviter donc toute méprise qui pourrait venir de cette source, nous ferons $b^{0}=C$, et nous désignerons par $B$ le complément de $C$, de sorte qu'on aura $B=c^{0}$. Maintenant il faut exprimer $F(C,\theta)$ et $E(C,\theta)$ par le moyen de $F(C^{0},\theta^{0})$ et $E(C^{0},\theta^{0})$, ce qui se fera par les formules des art.\ 60, 61, et en observant que la loi des modules décroissans donne $C^{0}=\frac{1-B}{1+B}=\frac{1-c^{0}}{1+c^{0}}=b$, on aura de cette manière les valeurs
\[ \begin{aligned} F(C,\theta)&=\frac{1+C^{0}}{2}F(C^{0},\theta^{0})=\frac{1+b}{2}F(b,\theta^{0})\\ E(C,\theta)&=-BF(C,\theta)+\frac{1+B}{2}E(C^{0},\theta^{0})+\frac{1-B}{2}\sin\theta^{0}\\ &=-\left(\frac{1-b}{2}\right)F(b,\theta^{0})+\frac{1}{1+b}E(b,\theta^{0})+\frac{b}{1+b}\sin\theta^{0}. \end{aligned} \]

\origpage{158}
A l'égard de $\theta^{0}$, cet angle se déduit de $\theta$ par la formule $\operatorname{tang}(\theta^{0}-\theta)=B\operatorname{tang}\theta=\frac{1-b}{1+b}\operatorname{tang}\theta$, qui donne
\[ \operatorname{tang}\theta^{0}=\frac{\sin 2\theta}{b+\cos 2\theta}=\frac{\sin\lambda}{b-\cos\lambda},\quad \sin\theta^{0}=\frac{\sin\lambda}{\nu}, \]
où il faut observer qu'on devra faire $\theta^{0}<\frac{1}{2}\pi$ si on a $\cos\lambda<b$, et $\theta^{0}>\frac{1}{2}\pi$ si on a $\cos\lambda>b$.

Cela posé, en faisant toutes les substitutions, il viendra
\[ \begin{aligned} \frac{\Pi^{1}(n^{0},c^{0})}{1+b}=\tfrac{1}{4}F^{1}(c)+\frac{\nu}{4b\sin\lambda}[&\pi+F^{1}(c)F(b,\theta^{0})-E^{1}(c)F(b,\theta^{0})\\ &-F^{1}(c)E(b,\theta^{0})]. \end{aligned} \]
Donc enfin la fonction complète
\[ \begin{aligned} \Pi^{1}(n)={}&\tfrac{1}{2}F^{1}(c)-\left(\frac{1-b\cos\lambda}{2b\nu^{2}}\right)(b-\cos\lambda+\sqrt{-1}\sin\lambda)F^{1}(c)\\ &+\frac{\sin\lambda+(\cos\lambda-b)\sqrt{-1}}{2b\nu}[\pi+F^{1}(c)F(b,\theta^{0})-E^{1}(c)F(b,\theta^{0})\\ &\qquad-F^{1}(c)E(b,\theta^{0})]. \end{aligned} \]
Pour avoir la valeur de cette même fonction par les formules de l'art.\ 111, il faudra d'abord mettre $m'$ sous la forme $1-b^{2}\sin^{2}\psi$, et chercher ensuite la valeur de $\Pi^{1}(-m')$ par la formule de l'art.\ 101; on observera d'ailleurs que les angles $\psi$ et $\theta^{0}$ ont entre eux la relation $1=c\operatorname{tang}\psi\operatorname{tang}\theta^{0}$, d'où résulte $F(b,\psi)+F(b,\theta^{0})=F^{1}(b)$, $E(b,\psi)+E(b,\theta^{0})=E^{1}(b)+b^{2}\sin\psi\sin\theta^{0}$. La substitution étant donc faite, on trouvera après les réductions, la même valeur de $\Pi^{1}(n)$ que la précédente, ce qui prouve l'accord des deux méthodes.

\subsection*{(114).}

Puisque toute fonction elliptique de troisième espèce dont le paramètre est imaginaire, peut se réduire à deux fonctions de la même espèce dont le paramètre est réel, il s'ensuit:

1°. Que toute fonction complète de troisième espèce dont le paramètre est imaginaire, peut toujours se réduire aux fonctions elliptiques de la première et de la seconde espèce.

2°. Que toute fonction de cette sorte non-complète, mais dont l'amplitude $\varphi$ est telle qu'on ait $F(c,\varphi)=kF^{1}(c)$, $k$ étant un nombre rationnel, pourra également s'exprimer par des fonctions de la première et de la seconde espèce.

\origpage{159}

3°. Que toute fonction $\Pi(n)$ dont le paramètre est imaginaire, pourra se réduire indéfiniment à la première espèce, si l'on peut satisfaire en même temps aux deux conditions nécessaires pour que les fonctions auxiliaires $\Pi(-m)$, $\Pi(-m')$ soient susceptibles d'une semblable réduction. On a déjà vu un exemple de ces réductions (art.\ 109) et il serait facile d'en produire beaucoup d'autres.

\section*{D'un symptôme général pour reconnaître si deux fonctions données de troisième espèce, qui ne diffèrent que par les paramètres, peuvent se réduire l'une à l'autre.}

\subsection*{(115).}

Nous avons admis trois formes dans les paramètres, savoir, $n=\cot^{2}\theta$, $n=-1+b^{2}\sin^{2}\theta$, $n=-c^{2}\sin^{2}\theta$; et nous avons fait voir que toute fonction qui se rapporte à la première forme, peut être convertie en une fonction qui se rapporte à la seconde forme, et réciproquement. Cette conversion se fait par la formule du n° 51, et sans rien changer aux deux autres élémens de chaque fonction, qui sont le module $c$ et l'amplitude $\varphi$. Mais une fonction dont le paramètre appartient à la troisième forme, ne peut jamais être réduite qu'à une fonction dont le paramètre est de la même forme. C'est pourquoi nous nous bornerons à comparer successivement, pour les trois formes du paramètre, deux fonctions qui se rapportent à une même forme.

Soit 1°. $n=\cot^{2}\theta$; la formule $(h')$ du n° 95 prouve que la fonction $\Pi(\cot^{2}\theta,c,\varphi)$ peut se réduire à la fonction $\Pi(b^{2}\operatorname{tang}^{2}\varphi,b,\theta)$ qui diffère de la première par ses trois élémens. Semblablement la fonction $\Pi(\cot^{2}\lambda,c,\varphi)$ pourra se réduire à la fonction $\Pi(b^{2}\operatorname{tang}^{2}\varphi,b,\lambda)$; mais les fonctions $\Pi(b^{2}\operatorname{tang}^{2}\varphi,b,\theta)$, $\Pi(b^{2}\operatorname{tang}^{2}\varphi,b,\lambda)$ qui ne diffèrent que par l'amplitude, peuvent se comparer entre elles et avec la fonction complète $\Pi^{1}(b^{2}\operatorname{tang}^{2}\varphi,b)$, d'une infinité de manières par les formules des n\textsuperscript{os} 53 et 57. Supposons en général qu'on ait l'équation
\[ iF(b,\lambda)\pm kF(b,\theta)=lF^{1}(b), \]
$i$, $k$, $l$ étant des nombres entiers; alors il s'ensuivra par les principes connus,
\[ i\Pi(b^{2}\operatorname{tang}^{2}\varphi,h,\lambda)\pm k\Pi(b^{2}\operatorname{tang}^{2}\varphi,b,\theta)=l\Pi^{1}(b^{2}\operatorname{tang}^{2}\varphi,b)+W, \]
\ednote{The second argument of the first $\Pi$ reads as $h$ in the scan, a misprint for $b$ (the other functions in this line and the text have $b$).}

\origpage{160}
$W$ étant une quantité déterminable par arcs de cercle. Mais la fonction complète de troisième espèce $\Pi^{1}(b^{2}\operatorname{tang}^{2}\varphi,b)$ est réductible aux fonctions de la première et de la seconde espèce; donc l'équation précédente donne le moyen de réduire l'une à l'autre les deux fonctions $\Pi(b^{2}\operatorname{tang}^{2}\varphi,b,\lambda)$, $\Pi(b^{2}\operatorname{tang}^{2}\varphi,b,\theta)$. Donc en général les deux fonctions $\Pi(\cot^{2}\theta,c,\varphi)$, $\Pi(\cot^{2}\lambda,c,\varphi)$ pourront se réduire l'une à l'autre, si les angles $\theta$ et $\lambda$ des paramètres sont tels qu'on ait $iF(b,\lambda)\pm kF(b,\theta)=lF^{1}(b)$, $i$, $k$, $l$ étant des nombres entiers à volonté.

Soit 2°. $n=-1+b^{2}\sin^{2}\theta$, la fonction $\Pi(n,c,\varphi)$ pourra se réduire à la fonction $\Pi(b^{2}\operatorname{tang}^{2}\varphi,b,\theta)$ par la formule du n° 101. De même si l'on a $n'=-1+b^{2}\sin^{2}\lambda$, la fonction $\Pi(n',c,\varphi)$ pourra se réduire à la fonction $\Pi(b^{2}\operatorname{tang}^{2}\varphi,b,\lambda)$. Mais on verra, comme dans le cas précédent, que les deux fonctions $\Pi(b^{2}\operatorname{tang}^{2}\varphi,b,\theta)$, $\Pi(b^{2}\operatorname{tang}^{2}\varphi,b,\lambda)$, peuvent se réduire l'une à l'autre, si les angles $\theta$ et $\lambda$, déduits des paramètres $n$ et $n'$, sont tels qu'on ait $iF(b,\lambda)\pm kF(b,\theta)=lF^{1}(b)$, $k$, $i$, $l$ étant des entiers. Donc cette condition ayant lieu, il sera toujours possible de réduire l'une à l'autre les deux fonctions $\Pi(n,c,\varphi)$, $\Pi(n',c,\varphi)$.

Soit 3°. $n=-c^{2}\sin^{2}\theta$, la fonction $\Pi(-c^{2}\sin^{2}\theta,c,\varphi)$ pourra se réduire à la fonction $\Pi(-c^{2}\sin^{2}\varphi,c,\theta)$ par la formule du n° 94; semblablement la fonction $\Pi(-c^{2}\sin^{2}\lambda,c,\varphi)$ pourra se réduire à la fonction $\Pi(-c^{2}\sin^{2}\varphi,c,\lambda)$. Supposons qu'entre les angles $\theta$ et $\lambda$, déduits des paramètres, on ait la relation $iF(c,\lambda)\pm kF(c,\theta)=lF^{1}(c)$, $i$, $k$, $l$ étant des entiers; alors on aura par les principes connus,
\[ i\Pi(-c^{2}\sin^{2}\varphi,c,\lambda)\pm k\Pi(-c^{2}\sin^{2}\varphi,c,\theta)=l\Pi^{1}(-c^{2}\sin^{2}\varphi,c)+W, \]
$W$ étant une quantité déterminable par logarithmes. Donc la fonction $\Pi(-c^{2}\sin^{2}\varphi,c,\lambda)$ pourra être exprimée par la fonction $\Pi(-c^{2}\sin^{2}\varphi,c,\theta)$, et réciproquement. Donc en général, si la condition mentionnée a lieu, les deux fonctions $\Pi(-c^{2}\sin^{2}\theta,c,\varphi)$, $\Pi(-c^{2}\sin^{2}\lambda,c,\varphi)$ seront réductibles l'une à l'autre.

On voit directement un exemple de cette réduction dans la formule du n° 52. Mais le symptôme général que nous venons de donner, donne les moyens de multiplier à l'infini ces sortes de comparaisons, et prouve que toute fonction donnée de troisième espèce peut être transformée en une infinité d'autres qui n'en différeront que par le paramètre.

\origpage{161}

La formule du n° 46, qui est la plus simple et la plus remarquable des formules de comparaison, est comprise dans le symptôme général donné pour les deux premières formes, puisqu'en faisant $n=\cot^{2}\theta$ et $\frac{c^{2}}{n}=\cot^{2}\lambda$, on a entre les angles $\theta$ et $\lambda$, la relation $c\operatorname{tang}\theta\operatorname{tang}\lambda=1$, d'où résulte $F(b,\theta)+F(b,\lambda)=F^{1}(b)$.

La même formule servirait à comparer les deux fonctions $\Pi(-1+b^{2}\sin^{2}\theta)$, $\Pi(-1+b^{2}\sin^{2}\lambda)$; car en faisant $n=-1+b^{2}\sin^{2}\theta$, $\frac{c^{2}}{n}=-1+b^{2}\sin^{2}\lambda$, il en résulte encore $c\operatorname{tang}\theta\operatorname{tang}\lambda=1$.

\section*{Exemple d'une transformation particulière de fonctions elliptiques de la première et de la seconde espèce.}

\subsection*{(116).}

Deux fonctions elliptiques de la première espèce peuvent être transformées l'une dans l'autre, si leurs modules appartiennent à une même suite ou échelle $\ldots c''$, $c'$, $c$, $c^{0}$, $c^{00}$, etc. formée suivant la loi connue. Cette propriété s'étend aux fonctions de la seconde et de la troisième espèce, pourvu qu'elles soient jointes à une fonction de la première espèce, dont le coefficient puisse changer à volonté. Mais si les modules des deux fonctions comparées ne sont pas compris dans la suite dont il s'agit, il y a très-peu de cas où la réduction d'une fonction à l'autre soit possible; nous ne connaissons qu'un seul exemple où la réduction indéfinie ait lieu entre deux fonctions dont les modules sont complémens l'un et l'autre, et c'est cet exemple que nous allons développer.

Considérons la formule $R=\displaystyle\int\frac{dz}{(1-z^{3})^{\frac{2}{3}}}$ qui doit être intégrée depuis $z=0$; je fais $1-z^{3}=\left(\frac{z}{u}\right)^{\frac{3}{2}}$, ce qui donne $z^{3}=-\frac{1}{2}u^{-\frac{3}{2}}+\sqrt{(1+\frac{1}{4}u^{-3})}$\ednote{Printed $z^{3}$; the quadratic in $z^{\frac{3}{2}}$ requires $z^{\frac{3}{2}}$ here, an apparent misprint.}, et j'ai la transformée
\[ R=\int\frac{du}{\sqrt{(4u^{3}+1)}}. \]
Soit $m=\sqrt[3]{4}$, et $mu=t^{2}-1$, on aura pour seconde transformée
\[ R=\frac{2}{m}\int\frac{dt}{\sqrt{(t^{4}-3t^{2}+3)}}. \]

\origpage{162}
Enfin soit $n=\sqrt[4]{3}$, $t=n\operatorname{tang}\frac{1}{2}\varphi$, $c^{2}=\frac{2+\sqrt{3}}{4}$, on aura pour troisième transformée
\[ R=\frac{1}{mn}\int\frac{d\varphi}{\sqrt{(1-c^{2}\sin^{2}\varphi)}}=\frac{1}{mn}(F(c,\varphi)+C). \]
Pour déterminer la constante, soit $z=0$, on aura $u=0$, $t=1$; et si l'on fait $\operatorname{tang}\frac{1}{2}\lambda=\frac{1}{n}$, on aura $\varphi=\lambda$, de sorte que l'intégrale cherchée sera
\[ R=\frac{1}{mn}[F(c,\varphi)-F(c,\lambda)]. \]
Si nous prenons un second angle $\mu$ tel que $b\operatorname{tang}\mu\operatorname{tang}\lambda=1$, afin qu'on ait $F(c,\lambda)+F(c,\mu)=F^{1}(c)$, il en résultera $\operatorname{tang}\mu=\frac{1}{b}\cot\lambda=\frac{1}{b}\cdot\frac{n^{2}-1}{2n}=\frac{\sqrt{2}}{n}=\sqrt{\left(\frac{2}{\sqrt{3}}\right)}$. Cette valeur est celle qui pour le module $c$ donne $F(c,\mu)=\frac{1}{3}F^{1}(c)$; donc $F(c,\lambda)=\frac{2}{3}F^{1}(c)$, et enfin
\[ R=\frac{1}{mn}[F(c,\varphi)-\tfrac{2}{3}F^{1}(c)]. \]
L'intégrale entière prise jusqu'à $z=1$ se trouvera en faisant $\varphi=\pi$, et sa valeur sera
\[ R'=\frac{1}{mn}\cdot\tfrac{4}{3}F^{1}(c). \]

\subsection*{(117).}

Faisons maintenant usage d'une autre transformation pour avoir la valeur de $R$. Soit $\sqrt[3]{(1-z^{3})}=1-\frac{z}{y}$, ce qui donne $(y^{3}-1)z^{2}+3yz=3y^{2}$, $z=\frac{y\sqrt{3}}{2}\cdot\frac{\sqrt{(4y^{3}-1)}-\sqrt{3}}{y^{3}-1}$; on aura d'abord $R=\displaystyle\int\frac{1}{zy}\left(\frac{dz}{z}-\frac{dy}{y}\right)$, et achevant la substitution, on trouve
\[ R=\sqrt{3}\cdot\int\frac{-dy}{\sqrt{(4y^{3}-1)}}. \]
Soit maintenant $my=1+x^{2}$, on aura
\[ R=\frac{2n^{2}}{m}\int\frac{-dx}{\sqrt{(x^{4}+3x^{2}+3)}}. \]
Soit enfin $x=n\cot\frac{1}{2}\omega$, et $b^{2}=\frac{2-\sqrt{3}}{4}$, on aura pour dernière

\origpage{163}
transformée
\[ R=\frac{n}{m}\int\frac{d\omega}{\sqrt{(1-b^{2}\sin^{2}\omega)}}=\frac{n}{m}F(b,\omega). \]
Il n'y a pas de constante à ajouter, parce qu'en faisant $z=0$, on a successivement $y=\infty$, $z=\infty$\ednote{Printed $z=\infty$; $x=\infty$ is evidently meant, an apparent misprint.}, $\omega=0$.

Si on fait $z=1$ pour avoir l'intégrale complète ou définie, on aura $y=1$, $x^{2}=m-1$, $\cot\frac{1}{2}\omega=\frac{\sqrt{(m-1)}}{n}$. Mais l'angle $\theta$ déterminé par l'équation $\cot\frac{1}{2}\theta=\frac{\sqrt{(m-1)}}{n}$, est tel qu'on a $F(b,\theta)=\frac{4}{3}F^{1}(b)$.

En effet nous avons trouvé par les formules de la trisection (art.\ 24) que pour le module $b=\frac{1}{2}\sqrt{(2-\sqrt{3})}$, on satisfait à l'équation $F(b,\gamma)=\frac{1}{3}F^{1}(b)$, en prenant $\cos\gamma=(m-1)\sqrt{\left(\frac{2+n^{2}}{n^{2}}\right)}$; soit $m-1=n^{2}r$, on aura $\cos\gamma=nr\sqrt{(2+n^{2})}$, et $\sin\gamma=\sqrt{(1-2n^{2}r^{2}-n^{4}r^{2})}$. Mais l'équation $m=1+n^{2}r$ étant élevée au cube, donne $1=n^{2}r+n^{4}r^{2}+n^{2}r^{3}$; donc $1-2n^{2}r^{2}-n^{4}r^{2}=n^{2}(r-2r^{2}+2r^{3})$\ednote{Printed $+2r^{3}$; the preceding equation gives $+r^{3}$, an apparent misprint.}, et par conséquent $\sin\gamma=n(1-r)\sqrt{r}$, $\operatorname{tang}\gamma=\frac{1-r}{\sqrt{r(2+n^{2})}}$. Soit $\delta$ l'angle qui donne $F(b,\delta)+F(b,\gamma)=F^{1}(b)$, ou $F(b,\delta)=\frac{2}{3}F^{1}(b)$, on aura $c\operatorname{tang}\delta\operatorname{tang}\gamma=1$; donc $\operatorname{tang}\delta=\frac{\cot\gamma}{c}=\frac{2\cot\gamma}{\sqrt{(2+n^{2})}}=\frac{2\sqrt{r}}{1-r}$, et $\operatorname{tang}\frac{1}{2}\delta=\sqrt{r}=\frac{\sqrt{(m-1)}}{n}$. Cette valeur étant celle de $\cot\frac{1}{2}\theta$, il s'ensuit qu'on a $\frac{1}{2}\delta+\frac{1}{2}\theta=\frac{1}{2}\pi$, ou $\theta=\pi-\delta$, donc $F(b,\theta)=2F^{1}(b)-\frac{2}{3}F^{1}(b)=\frac{4}{3}F^{1}(b)$. Donc enfin
\[ R^{1}=\frac{n}{m}\cdot\tfrac{4}{3}F^{1}(b). \]
\ednote{Printed $R^{1}$ here, while the next sentence and p.~162 write $R'$; possibly a damaged prime.}
Comparant les deux valeurs trouvées pour $R'$, on en tire
\[ F^{1}(c)=n^{2}F^{1}(b)=\sqrt{3}F^{1}(b), \]
ce que nous avons déjà trouvé (n° 41).

Et si l'on compare les deux valeurs de $R$, on aura, quel que soit $\varphi$, cette formule générale de réduction entre deux fonctions de différens modules
\[ F(c,\varphi)-\tfrac{2}{3}F^{1}(c)=\sqrt{3}F(b,\omega). \]

\origpage{164}

Quant à la relation entre les angles $\omega$ et $\varphi$, on peut la déduire de l'analyse précédente, et il en résulte
\[ \operatorname{tang}\tfrac{1}{2}\varphi=\frac{n\sin\tfrac{1}{2}\omega\Delta(b,\omega)+\cos^{3}\tfrac{1}{2}\omega}{n\cos\tfrac{1}{2}\omega\Delta(b,\omega)-\sin^{3}\tfrac{1}{2}\omega}. \]
Ainsi la simple substitution de cette valeur dans l'intégrale $\displaystyle\int\frac{d\varphi}{\sqrt{(1-c^{2}\sin^{2}\varphi)}}$, doit la transformer en $n^{2}\displaystyle\int\frac{d\omega}{\sqrt{(1-b^{2}\sin^{3}\omega)}}$\ednote{Printed $\sin^{3}\omega$ instead of $\sin^{2}\omega$; an apparent misprint.}, et c'est ce qu'on peut aisément vérifier.

La fonction $F(c,\varphi)$ s'exprime donc au moyen de la fonction $F(b,\omega)$ et de la fonction complète $F^{1}(b)$; car on déduit des équations précédentes,
\[ F(c,\varphi)=\sqrt{3}[F(b,\omega)+\tfrac{2}{3}F^{1}(b)]. \]
On peut ensuite réduire le second membre à une seule fonction; car ayant déjà $F(b,\delta)=\frac{2}{3}F^{1}(b)$, si l'on prend un angle $\psi$ tel que $F(b,\psi)=F(b,\omega)+F(b,\delta)$, ce qui se fera par les formules de l'art.\ 18, on aura
\[ F(c,\varphi)=\sqrt{3}F(b,\psi); \]
de sorte que les fonctions $F(c,\varphi)$, $F(b,\psi)$ dont les modules sont complémens l'un de l'autre, seront toujours entre elles dans un rapport constant.

On peut d'ailleurs avoir immédiatement la valeur de $\operatorname{tang}\frac{1}{2}\varphi$ exprimée en fonction de $\psi$; il suffit pour cela de substituer dans la valeur déjà trouvée de $\operatorname{tang}\frac{1}{2}\varphi$, celles de $\operatorname{tang}\frac{1}{2}\omega$ et $\Delta(b,\omega)$ qui sont
\[ \begin{aligned} \operatorname{tang}\tfrac{1}{2}\omega&=\frac{\sin\psi\Delta(b,\delta)-\sin\delta\Delta(b,\psi)}{\cos\psi+\cos\delta}\\ \Delta(b,\omega)&=\frac{\Delta(b,\psi)\Delta(b,\delta)+b^{2}\sin\psi\sin\delta\cos\psi\cos\delta}{1-c^{2}\sin^{2}\psi\sin^{2}\delta}, \end{aligned} \]
et cette substitution n'offre aucune difficulté.

\subsection*{(118).}

Cherchons maintenant la valeur de la fonction de seconde espèce $E(c,\varphi)$ exprimée par la fonction $E(b,\omega)$, ou, s'il est nécessaire, par les deux fonctions $E(b,\omega)$, $F(b,\omega)$. Le moyen qui se présente naturellement est de substituer la valeur connue de

\origpage{165}
$\operatorname{tang}\frac{1}{2}\varphi$ en fonction de $\omega$ dans l'intégrale $\displaystyle\int\Delta\,d\varphi$; et comme on a $\displaystyle\int\Delta\,d\varphi=\displaystyle\int\frac{d\varphi}{\Delta}(1-c^{2}\sin^{2}\varphi)=\displaystyle\int\frac{n^{2}\,d\omega}{\Delta(b,\omega)}(1-c^{2}\sin^{2}\varphi)$, la substitution ne resterait à faire que dans le facteur $1-c^{2}\sin^{2}\varphi$; or on trouve
\[ \sqrt{(1-c^{2}\sin^{2}\varphi)}=\frac{(1+\tfrac{1}{2}(n^{2}-1)\sin^{2}\omega)\Delta(b,\omega)-\tfrac{1}{2}n\sin\omega\cos\omega}{1+n^{2}-\tfrac{1}{2}n^{2}\sin^{2}\omega+n\sin\omega\cos\omega\Delta(b,\omega)}. \]
Mais on voit que le résultat de la substitution sera fort compliqué, et que si l'intégration n'offre pas des réductions inattendues et difficiles à apercevoir, il entrera nécessairement des fonctions de troisième espèce dans l'expression de $E(c,\varphi)$, ce qui rendrait cette transformation illusoire.

Il paraît donc nécessaire de suivre une autre route pour arriver au but qu'on se propose. Pour cet effet, considérons la formule intégrale
\[ T=\int\frac{x\,dx}{\sqrt[3]{(1-x^{3})}}: \]
si on fait $1-x^{3}=\left(\frac{x}{u}\right)^{\frac{3}{2}}$, on aura $x^{\frac{3}{2}}=-\frac{1}{2}u^{-\frac{3}{2}}+\sqrt{(1+\frac{1}{4}u^{-3})}$, $T=\displaystyle\int x^{\frac{1}{2}}u^{\frac{1}{2}}\,dx=-\frac{1}{2u}-\displaystyle\int\frac{du}{2u^{2}\sqrt{(4u^{3}+1)}}$. Mais l'intégration par parties donne $\displaystyle\int\frac{du}{u^{2}\sqrt{(4u^{3}+1)}}=-\frac{1}{u}\sqrt{(4u^{3}+1)}+\displaystyle\int\frac{2u\,du}{\sqrt{(4u^{3}+1)}}$. Donc
\[ T=\frac{\sqrt{(4u^{3}+1)}-1}{2u}-\int\frac{u\,du}{\sqrt{(4u^{3}+1)}}. \]
Tout se réduit donc à trouver l'intégrale $V=\displaystyle\int\frac{u\,du}{\sqrt{(4u^{3}+1)}}$. Pour cela soit $m^{3}=4$ et $mu=t^{2}-1$, on aura
\[ V=\frac{2}{m^{2}}\int\frac{(t^{2}-1)\,dt}{\sqrt{(t^{4}-3t^{2}+3)}}. \]
Soit encore $n^{4}=3$, $t=n\operatorname{tang}\frac{1}{2}\varphi$, $c=\frac{1}{2}\sqrt{(2+\sqrt{3})}$, et la transformée sera
\[ V=\frac{1}{m^{2}n}\int\frac{(n^{2}\operatorname{tang}^{2}\tfrac{1}{2}\varphi-1)\,d\varphi}{\sqrt{(1-c^{2}\sin^{2}\varphi)}}=\frac{n}{m^{2}}\int\frac{d\varphi}{\Delta\cos^{2}\tfrac{1}{2}\varphi}-\left(\frac{n^{2}+1}{m^{2}n}\right)F(c,\varphi). \]
Mais on a $\displaystyle\int\frac{d\varphi}{\Delta\cos^{2}\frac{1}{2}\varphi}=2\Delta\operatorname{tang}\frac{1}{2}\varphi+2F(c,\varphi)-2E(c,\varphi)$; donc

\origpage{166}
enfin
\[ V=\frac{2n}{m^{2}}\left[\Delta\operatorname{tang}\tfrac{1}{2}\varphi+\frac{n^{2}-1}{2n^{2}}F(c,\varphi)-E(c,\varphi)\right]; \]
et l'intégrale cherchée
\[ T=\frac{\sqrt{(4u^{3}+1)}-1}{2u}-\frac{2n}{m^{2}}\Delta\operatorname{tang}\tfrac{1}{2}\varphi+\frac{1-n^{2}}{m^{2}n}F+\frac{2n}{m^{2}}E+\text{const.} \]
Il faut déterminer la constante de manière que l'intégrale s'évanouisse lorsque $x=0$: alors $u=0$, $t=1$, $\operatorname{tang}\frac{1}{2}\varphi=\frac{1}{n}$, $\varphi=\lambda$, $F(c,\varphi)=F(c,\lambda)=\frac{2}{3}F^{1}(c)$, $E(c,\lambda)=\frac{2}{3}E^{1}(c)+\frac{c^{2}\sin\mu\sin\lambda}{3}(2-\sin\mu)=\frac{2}{3}E^{1}(c)+\frac{n^{2}+1}{2n}$, $\Delta\operatorname{tang}\frac{1}{2}\varphi=\frac{1}{n(n^{2}+1)}$. Donc on a
\[ \begin{aligned} T={}&\frac{\sqrt{(4u^{3}+1)}-1}{2u}-\frac{2n}{m^{2}}\Delta\operatorname{tang}\tfrac{1}{2}\varphi-\left(\frac{n^{2}-1}{m^{2}n}\right)(F(c,\varphi)-\tfrac{2}{3}F^{1}(c))\\ &+\frac{m}{2(n^{2}+1)}+\frac{2n}{m^{2}}(E(c,\varphi)-\tfrac{2}{3}E^{1}(c)), \end{aligned} \]
ou pour tout exprimer en fonction de $\varphi$:
\[ \begin{aligned} T={}&\tfrac{1}{2}m\cos\tfrac{1}{2}\varphi\cdot\frac{n(n^{2}+1)\Delta\sin\tfrac{1}{2}\varphi-\cos\tfrac{1}{2}\varphi}{n^{2}\sin^{2}\tfrac{1}{2}\varphi-\cos^{2}\tfrac{1}{2}\varphi}-\left(\frac{n^{2}-1}{m^{2}n}\right)[F(c,\varphi)-\tfrac{2}{3}F^{1}(c)]\\ &+\frac{2n}{m^{2}}[E(c,\varphi)-\tfrac{2}{3}E^{1}(c)]. \end{aligned} \]
Lorsque $x=1$, on a $\varphi=\pi$, et cette valeur devient
\[ T^{1}=\frac{2mn}{3}\left[E^{1}(c)-\left(\frac{n^{2}-1}{2n^{2}}\right)F^{1}(c)\right]. \]
Or par une propriété des fonctions $F^{1}(c)$, $E^{1}(c)$, qui a été démontrée n° 39, on a $E^{1}(c)-\frac{n^{2}-1}{2n^{2}}F^{1}(c)=\frac{\pi\sqrt{3}}{4F^{1}(c)}$; donc
\[ T^{1}=\frac{m}{2n}\cdot\frac{\pi}{F^{1}(c)}. \]

\subsection*{(119).}

Intégrons maintenant par un autre procédé la formule proposée $T=\displaystyle\int\frac{x\,dx}{\sqrt[3]{(1-x^{3})}}$, et soit comme ci-dessus $\sqrt[3]{(1-x^{3})}=1-\frac{x}{y}$.

\origpage{167}
la transformée sera
\[ T=\frac{9\sqrt{3}}{2}\int\frac{y^{4}\,dy}{(y^{3}-1)^{2}\sqrt{(4y^{3}-1)}}-\int\frac{3y\,dy\,(1+\tfrac{1}{2}y^{3})}{(y^{3}-1)^{2}}. \]
Cette transformée paraît au premier coup d'œil beaucoup plus composée que celle qui est résultée de la première méthode; et il semble qu'on ne peut éviter d'introduire dans l'intégrale des fonctions elliptiques de la troisième espèce, même de celles dont le paramètre est imaginaire. Mais ces craintes se dissipent en continuant le calcul qui offre des réductions très-heureuses.

J'observe d'abord que la partie rationnelle s'intègre algébriquement et qu'on a
\[ \int\frac{3y\,dy\,(1+\tfrac{1}{2}y^{3})}{(y^{3}-1)^{2}}=-\tfrac{3}{2}\cdot\frac{y^{2}}{y^{3}-1}. \]
Ensuite si on fait $Z=\dfrac{y^{2}\sqrt{(4y^{3}-1)}}{y^{3}-1}$, on aura par la différentiation,
\[ dZ=\frac{2y\,dy}{\sqrt{(4y^{3}-1)}}-\frac{9y^{4}\,dy}{(y^{3}-1)^{2}\sqrt{(4y^{3}-1)}}; \]
donc
\[ \int\frac{y^{4}\,dy}{(y^{3}-1)^{2}\sqrt{(4y^{3}-1)}}=\tfrac{2}{9}\int\frac{y\,dy}{\sqrt{(4y^{3}-1)}}-\tfrac{1}{9}Z. \]
Par ces diverses réductions on obtient
\[ T=\sqrt{3}\int\frac{y\,dy}{\sqrt{(4y^{3}-1)}}-\frac{\sqrt{3}}{2}\cdot\frac{y^{2}\sqrt{(4y^{3}-1)}}{y^{3}-1}+\tfrac{3}{2}\cdot\frac{y^{2}}{y^{3}-1}. \]
Tout se réduit donc à trouver l'intégrale $U=\displaystyle\int\frac{y\,dy}{\sqrt{(4y^{3}-1)}}$; or si l'on fait comme ci-dessus $my=1+\nu^{2}$, on aura
\[ U=\frac{m}{2}\int\frac{d\nu\,(1+\nu^{2})}{\sqrt{(\nu^{4}+3\nu^{2}+3)}}. \]
Si ensuite on fait $\nu=n\cot\frac{1}{2}\omega$, et $b=\frac{1}{2}\sqrt{(2-\sqrt{3})}$, on aura
\[ U=-\frac{m}{4n}\int\frac{d\omega\,(1+n^{2}\cot^{2}\tfrac{1}{2}\omega)}{\Delta(b,\omega)}, \]
ou
\[ U=-\frac{mn}{4}\int\frac{d\omega}{\sin^{2}\tfrac{1}{2}\omega\,\Delta(b,\omega)}+m\left(\frac{n^{2}-1}{4n}\right)\int\frac{d\omega}{\Delta(b,\omega)}, \]

\origpage{168}
et en appliquant les formules connues, on a enfin
\[ U=\frac{mn}{2}[E(b,\omega)+\cot\tfrac{1}{2}\omega\Delta(b,\omega)]-\frac{m(n^{2}+1)}{4n}F(b,\omega). \]
Maintenant si dans l'équation $T=U\sqrt{3}-yx$, on substitue la valeur de $U$ et celle de $yx$ en fonction de $\omega$, on aura l'intégrale cherchée
\[ \begin{aligned} T={}&\frac{mn^{3}}{2}E(b,\omega)-\frac{mn}{4}(n^{2}+1)F(b,\omega)\\ &+\frac{m\sin\tfrac{1}{2}\omega}{2}\cdot\frac{(4-2n^{2})\cos^{2}\tfrac{1}{2}\omega-3(2-n^{2})\cos^{4}\tfrac{1}{2}\omega-1+n^{3}\sin\tfrac{1}{2}\omega\cos\tfrac{1}{2}\omega\Delta(b,\omega)}{n\cos\tfrac{1}{2}\omega\Delta(b,\omega)+\sin^{3}\tfrac{1}{2}\omega}, \end{aligned} \]
intégrale où il n'y a pas de constante à ajouter, parce qu'elle s'évanouit lorsque $\omega=0$.

Pour avoir la valeur de l'intégrale complète $T^{1}$, soit $\omega=\theta$, on aura comme ci-dessus $F(b,\omega)=\frac{4}{3}F^{1}(b)$, $E(b,\omega)=\frac{4}{3}E^{1}(b)+W$, et on trouve $W=-\frac{b^{2}}{3}\sin\delta\sin\gamma\,(2-\sin\gamma)=-\frac{b^{2}}{3}\sin^{3}\delta$; donc
\[ T^{1}=\frac{2m}{n}\left[E^{1}(b)-\left(\frac{n^{2}+1}{2n^{2}}\right)F^{1}(b)\right]+\frac{3m}{2n}[W+\cot\tfrac{1}{2}\theta\Delta(b,\theta)]-1. \]
Mais par les valeurs déjà trouvées, on a $\cot\frac{1}{2}\theta=\sqrt{r}$, $\Delta(b,\theta)=\frac{1}{n(1+r)\sqrt{r}}$, $\sin\delta=\frac{2\sqrt{r}}{1+r}$, $\sin\gamma=1-\cos\delta=\frac{2r}{1+r}$; d'ailleurs les relations entre $m$, $n$, $r$ donnent $m=1+n^{2}r$, $n^{2}=\frac{1-3r^{2}}{r(1+r^{2})}$, $m=\frac{2(1-r^{2})}{1+r^{2}}$, $\sqrt{r}=\frac{n}{2}(1-r^{2})$. Substituant toutes ces valeurs en fonctions de $r$ dans la partie algébrique de $T^{1}$, et observant encore qu'on a $0=1-9r^{2}+3r^{4}-3r^{6}$, on trouve que cette partie se réduit à zéro, de sorte qu'on a simplement
\[ T^{1}=\frac{2m}{n}\left[E^{1}(b)-\left(\frac{n^{2}+1}{2n^{2}}\right)F^{1}(b)\right]. \]
Comparant cette valeur avec celle qui a été trouvée par l'autre méthode, il en résulte l'équation
\[ \frac{\pi}{4\sqrt{3}}=F^{1}(b)\left[E^{1}(b)-\left(\frac{1+\sqrt{3}}{2\sqrt{3}}\right)F^{1}(b)\right], \]
qui s'accorde avec la formule de l'art.\ 41.

\origpage{169}

Nous obtenons ainsi sur les fonctions définies $F^{1}(c)$, $F^{1}(b)$, $E^{1}(c)$, $E^{1}(b)$, les mêmes rapports que nous avons déjà obtenus par une voie plus simple. Quant aux fonctions indéfinies $F(c,\varphi)$, $F(b,\omega)$, $E(c,\varphi)$, $E(b,\omega)$, elles ont entre elles les relations comprises dans les deux équations
\[ \begin{gathered} F(c,\varphi)=n^{2}[F(b,\omega)+\tfrac{2}{3}F^{1}(b)],\\ \frac{mn}{2}[E(c,\varphi)-\tfrac{2}{3}E^{1}(c)]-\left(\frac{n^{2}-1}{m^{2}n}\right)[F(c,\varphi)-\tfrac{2}{3}F^{1}(c)]+\Phi\\ =\frac{mn^{3}}{2}E(b,\omega)-\frac{mn}{2}(n^{2}+1)F(b,\omega)+\Omega, \end{gathered} \]
$\Phi$ et $\Omega$ étant des fonctions algébriques connues, l'une de $\sin\frac{1}{2}\varphi$, l'autre de $\sin\frac{1}{2}\omega$; d'où l'on voit que les fonctions $F(c,\varphi)$, $E(c,\varphi)$ relatives au module $c$, peuvent s'exprimer indéfiniment par les fonctions $F(b,\omega)$, $E(b,\omega)$, relatives au module complémentaire $b$.

\section*{Des séries qui donnent, sans transformations, les valeurs approchées des fonctions elliptiques.}

\subsection*{(120).}

Il peut être nécessaire dans plusieurs cas, surtout dans les problèmes de mécanique, d'exprimer par la seule variable $\varphi$ les fonctions elliptiques dont ces problèmes dépendent. Alors le développement se fera de la manière connue; mais les méthodes précédentes serviront toujours à simplifier beaucoup la détermination des coefficiens.

Considérons d'abord la fonction $F=\displaystyle\int\frac{d\varphi}{\Delta}$, et supposons $\frac{1}{\Delta}=A-2B\cos2\varphi+4C\cos4\varphi-6D\cos6\varphi+\text{etc.}$, afin qu'il en résulte
\[ F=A\varphi-B\sin2\varphi+C\sin4\varphi-D\sin6\varphi+\text{etc.} \]
Il s'agit d'avoir le plus simplement qu'il est possible, les valeurs des coefficiens $A$, $B$, $C$, etc. Or par un premier développement on a
\[ \frac{1}{\Delta}=1+\tfrac{1}{2}c^{2}\sin^{2}\varphi+\frac{1\cdot3}{2\cdot4}c^{4}\sin^{4}\varphi+\frac{1\cdot3\cdot5}{2\cdot4\cdot6}c^{6}\sin^{6}\varphi+\text{etc.} \]
Mettant ensuite au lieu des puissances des sinus, leurs valeurs en

\origpage{170}
cosinus linéaires, on trouve une suite de la forme supposée, dont les coefficiens sont aisés à déduire les uns des autres. On a d'abord le premier
\[ A=1+\frac{1^{2}}{2^{2}}c^{2}+\frac{1^{2}\cdot3^{2}}{2^{2}\cdot4^{2}}c^{4}+\frac{1^{2}\cdot3^{2}\cdot5^{2}}{2^{2}\cdot4^{2}\cdot6^{2}}c^{6}+\frac{1^{2}\cdot3^{2}\cdot5^{2}\cdot7^{2}}{2^{2}\cdot4^{2}\cdot6^{2}\cdot8^{2}}c^{8}+\text{etc.} \]
Je le représente par
\[ A=1+m'c^{2}+m''c^{4}+m'''c^{6}+m^{\text{iv}}c^{8}+\text{etc.}; \]
les autres seront successivement
\[ \begin{aligned} B&=\tfrac{1}{2}m'c^{2}+\tfrac{2}{3}m''c^{4}+\tfrac{3}{4}m'''c^{6}+\tfrac{4}{5}m^{\text{iv}}c^{8}+\text{etc.}\\ C&=\tfrac{2}{3}\cdot\tfrac{1}{8}m''c^{4}+\tfrac{3}{4}\cdot\tfrac{2}{10}m'''c^{6}+\tfrac{4}{5}\cdot\tfrac{3}{12}m^{\text{iv}}c^{8}+\text{etc.}\\ D&=\tfrac{3}{4}\cdot\tfrac{2}{10}\cdot\tfrac{2}{18}m'''c^{6}+\tfrac{4}{5}\cdot\tfrac{3}{12}\cdot\tfrac{4}{21}m^{\text{iv}}c^{8}+\text{etc.}\\ E&=\tfrac{4}{5}\cdot\tfrac{3}{12}\cdot\tfrac{4}{21}\cdot\tfrac{3}{32}m^{\text{iv}}c^{8}+\text{etc.}\\ &\text{etc.} \end{aligned} \]
En général la $n^{\text{ieme}}$ des quantités $C$, $D$, $E$, etc., se déduira de la précédente, en omettant son premier terme, et multipliant les termes suivans par ceux de la suite $\frac{n}{(n+1)(2n+2)}$, $\frac{2n}{(n+1)(2n+3)}$, $\frac{3n}{(n+1)(2n+4)}$, etc. On pourra donc par ces suites, trouver les valeurs de $A$, $B$, $C$, etc., si le module $c$ est une quantité assez petite; mais lorsque $c$ diffère peu de l'unité, il faudrait calculer un grand nombre de termes pour n'avoir qu'une médiocre approximation.

\subsection*{(121).}

La valeur de $\frac{1}{\Delta}$ étant multipliée successivement par $d\varphi$, $d\varphi\cos2\varphi$, $d\varphi\cos4\varphi$, etc., puis intégrée depuis $\varphi=0$ jusqu'à $\varphi=\frac{1}{2}\pi$, on aura
\[ \begin{gathered} A\cdot\frac{\pi}{2}=\int\frac{d\varphi}{\Delta},\quad B\frac{\pi}{2}=\int-\frac{d\varphi\cos2\varphi}{\Delta},\\ 2C\cdot\frac{\pi}{2}=\int\frac{d\varphi\cos4\varphi}{\Delta},\quad 3D\cdot\frac{\pi}{2}=\int-\frac{d\varphi\cos6\varphi}{\Delta},\quad\text{etc.} \end{gathered} \]
Ainsi chaque coefficient peut se déterminer en particulier par une

\origpage{171}
intégrale définie, et toutes ces intégrales ne dépendent que des deux fonctions $F^{1}$, $E^{1}$, qu'on sait évaluer dans tous les cas, avec toute la précision nécessaire.

Les deux premiers $A$ et $B$ sont ceux qu'il importe de déterminer avec le plus de précision; ils se trouvent par les formules $\frac{\pi}{2}A=F^{1}$ et $\frac{\pi}{2}B=\frac{2}{c^{2}}(F^{1}-E^{1})-F^{1}$, et si on substitue pour $F^{1}$ et $E^{1}$ leurs valeurs données n\textsuperscript{os} 65 et 77, on aura
\[ \begin{gathered} A=\frac{2\sqrt{c^{0}}}{c}\cdot\frac{2\sqrt{c^{00}}}{c^{0}}\cdot\frac{2\sqrt{c^{000}}}{c^{00}}\cdot\text{etc.}\\ \frac{B}{A}=\frac{c^{0}}{2}+\frac{c^{0}c^{00}}{4}+\frac{c^{0}c^{00}c^{000}}{8}+\text{etc.} \end{gathered} \]
Ces deux premiers coefficiens étant trouvés, on peut en déduire tous les autres; car en différentiant l'équation $\frac{1}{\Delta}=A-2B\cos2\varphi+4C\cos4\varphi-\text{etc.}$, on a
\[ \frac{c^{2}\sin\varphi\cos\varphi}{\Delta^{3}}=4B\sin2\varphi-16C\sin4\varphi+36D\sin6\varphi-\text{etc.} \]
Multipliant le premier membre par $\frac{2\Delta^{2}}{c^{2}}$, et le second par la quantité équivalente $\frac{2}{c^{2}}-1+\cos2\varphi$; multipliant de même par $\sin2\varphi$, les deux membres de l'équation $\frac{1}{\Delta}=A-2B\cos2\varphi+4C\cos4\varphi-\text{etc.}$, comparant les deux produits, et réduisant tous les termes en sinus linéaires, on trouvera
\[ \begin{aligned} 2\cdot3C&=4B\left(\frac{2}{c^{2}}-1\right)-A\\ 3\cdot5D&=16C\left(\frac{2}{c^{2}}-1\right)-1\cdot3B\\ 4\cdot7E&=36D\left(\frac{2}{c^{2}}-1\right)-2\cdot5C\\ 5\cdot9F&=64E\left(\frac{2}{c^{2}}-1\right)-3\cdot7D\\ &\text{etc.} \end{aligned} \]
En général, si $A^{n}$ désigne le $n^{\text{ième}}$ terme de la suite $B$, $C$, $D$, $E$, etc., on aura

\origpage{172}
\[ n(2n-1)A^{n}=(2n-2)^{2}A^{n-1}\left(\frac{2}{c^{2}}-1\right)-(n-2)(2n-3)A^{n-2}: \]
c'est la loi suivant laquelle chaque coefficient, à compter de $D$, peut se déduire des deux précédens. Le coefficient $C$, excepté de la loi générale, se détermine par l'équation $6C=4B\left(\frac{2}{c^{2}}-1\right)-A$\ednote{The exponent of $c$ in this first occurrence is faintly printed and resembles $c^{4}$; it is read as $c^{2}$, as in the same equation later on this page and on p. 171.}, ou plus directement par la formule
\[ \begin{aligned} \frac{6C}{A}={}&c^{02}\left(1+\frac{c^{00}}{2}+\frac{c^{00}c^{000}}{4}+\frac{c^{00}c^{000}c^{0000}}{8}+\text{etc.}\right)\\ &+\frac{c^{00}}{2}\left(1+\frac{c^{000}}{2}+\frac{c^{000}c^{0000}}{4}+\text{etc.}\right). \end{aligned} \]
Lorsque $c$ est fort près de l'unité, on déterminera $E^{1}$ et $F^{1}$ par les formules qui conviennent à ce cas, on en déduira les coefficiens $A$ et $B$ par les formules $\frac{\pi}{2}A=F^{1}$, $\frac{\pi}{2}B=\frac{2}{c^{2}}(F^{1}-E^{1})-F^{1}$, on calculera ensuite $C$ par l'équation $6C=4B\left(\frac{2}{c^{2}}-1\right)-A$; les autres se déduiront chacun des deux précédens par la loi générale que nous avons exposée; et cette loi sera d'une application d'autant plus sûre, que le facteur $\frac{2}{c^{2}}-1$ se trouvera, dans le cas dont il s'agit, peu différent de l'unité.

\subsection*{(122).}

Les fonctions elliptiques de la seconde espèce pourront se développer de la même manière. En effet, si on considère généralement la fonction $G=\displaystyle\int(\alpha+\beta\sin^{2}\varphi)\frac{d\varphi}{\Delta}$, ou, ce qui revient au même, $G=\displaystyle\int(\alpha'+\beta'\cos2\varphi)\frac{d\varphi}{\Delta}$, les formules précédentes donneront
\[ \begin{aligned} G={}&\alpha'[A\varphi-B\sin2\varphi+C\sin4\varphi-D\sin6\varphi+\text{etc.}]\\ &-\beta'\left[B\varphi-\left(\frac{A+2C}{2}\right)\sin2\varphi+\frac{B+3D}{4}\sin4\varphi-\left(\frac{2C+4E}{6}\right)\sin6\varphi+\text{etc.}\right]. \end{aligned} \]
Ainsi ce développement ne présente aucune difficulté nouvelle, et s'exécute par les mêmes coefficiens que celui des fonctions de la première espèce.

On peut traiter de même les fonctions de la troisième espèce.

\origpage{173}
En effet, soit $\Pi=\displaystyle\int\frac{d\varphi}{(1+n\sin^{2}\varphi)\Delta}$; si on suppose
\[ \Pi=M\varphi-N\sin2\varphi+P\sin4\varphi-Q\sin6\varphi+\text{etc.}, \]
et qu'on différentie chaque membre par rapport à $\varphi$, il en résultera une valeur de $\frac{1}{\Delta}$ qui étant comparée à celle de l'article précédent, donne les équations
\[ \begin{aligned} N&=\frac{2}{n}(A-M)+M\\ 2P&=\frac{4}{n}(B-N)-2N-M\\ 3Q&=\frac{8}{n}(C-P)-4P-N\\ 4R&=\frac{12}{n}(D-Q)-6Q-2P\\ &\text{etc.} \end{aligned} \]
D'où il suit qu'en supposant toujours $A$, $B$, $C$, etc. connus, il suffit de déterminer le premier coefficient $M$ pour connaître tous les autres. Or en faisant $\varphi=\frac{1}{2}\pi$, on a $\Pi^{1}=M\cdot\frac{1}{2}\pi$. Ainsi $M$ se détermine par la fonction complète $\Pi^{1}$, laquelle ne dépend que des fonctions de la première et de la seconde espèce.

On pourrait aussi déduire la valeur de $M$ de celles des coefficiens $A$, $B$, $C$, etc., au moyen de la formule connue $\frac{\sqrt{(1+n)}}{1+n\sin^{2}\varphi}=1+2\alpha\cos2\varphi+2\alpha^{2}\cos4\varphi+\text{etc.}$, où l'on a $\alpha=\frac{\sqrt{(1+n)}-1}{\sqrt{(1+n)}+1}$; car en multipliant les deux membres par $\frac{d\varphi}{\Delta}$ et intégrant depuis $\varphi=0$ jusqu'à $\varphi=\frac{1}{2}\pi$, on aura
\[ M\sqrt{(1+n)}=A-2B\alpha+4C\alpha^{2}-6D\alpha^{3}+\text{etc.} \]

\section*{Surface du cône oblique.}

\subsection*{(123).}

\rmfigure{figures/fig-11.png}{Fig.~11.}{A right triangle SBO with the vertical side SO, the line SA to the point A on the base BO, the half-circle on the diameter BA below the base with centre C, and the radius CM to a point M on it.}

Considérons d'abord le cône oblique à base circulaire. Soit $S$ le sommet du cône, $C$ le centre de la base, $SO$ la perpendiculaire abaissée du sommet sur le plan de la base, $SAB$ la section faite dans le cône par le plan $SCO$. Si on fait le rayon $CA=1$, la hauteur $SO=h$, la distance $CO=f$, et qu'on appelle $\omega$ l'angle

\origpage{174}
$OCM$, l'aire $ASM$ correspondante à l'angle $\omega$ sera exprimée par l'intégrale
\[ Z=\int\tfrac{1}{2}d\omega\sqrt{[h^{2}+(1-f\cos\omega)^{2}]}. \]
Pour réduire cette intégrale à la forme ordinaire, soit $\operatorname{tang}\frac{1}{2}\omega=m\operatorname{tang}\frac{1}{2}\varphi$, et soit pris l'indéterminée $m$ de manière qu'on ait
\[ m^{4}=\frac{h^{2}+(1-f)^{2}}{h^{2}+(1+f)^{2}}; \]
si ensuite on fait $\alpha=\frac{1-m^{2}}{1+m^{2}}$ et $c^{2}=\frac{1}{2}-\frac{1}{2}m^{2}\cdot\frac{h^{2}+1-f^{2}}{h^{2}+(1-f)^{2}}$, on aura la transformée
\[ Z=\frac{2m\sqrt{[h^{2}+(1-f)^{2}]}}{(1+m^{2})^{2}}\int\frac{d\varphi\sqrt{(1-c^{2}\sin^{2}\varphi)}}{(1+\alpha\cos\varphi)^{2}}. \]
Pour avoir la surface totale du cône, il faudra prendre cette intégrale depuis $\varphi=0$, jusqu'à $\varphi=\pi$, et doubler le résultat; ou, ce qui revient au même, il faudra prendre l'intégrale suivante, depuis $\varphi=0$ jusqu'à $\varphi=\frac{1}{2}\pi$,
\[ Z=\frac{2m\sqrt{[h^{2}+(1-f)^{2}]}}{(1+m^{2})^{2}}\int\left(\frac{\Delta\,d\varphi}{(1+\alpha\cos\varphi)^{2}}+\frac{\Delta\,d\varphi}{(1-\alpha\cos\varphi)^{2}}\right). \]

Soit de nouveau $\alpha=\cos\theta$, ou $m=\operatorname{tang}\frac{1}{2}\theta$. Si on fait pour abréger $AS=k=\sqrt{[h^{2}+(1-f)^{2}]}$, on aura en réduisant,
\[ Z=2k\cot\tfrac{1}{2}\theta\int\frac{1+2\cot^{2}\theta-\cot^{2}\theta\sin^{2}\varphi}{(1+\cot^{2}\theta\sin^{2}\varphi)^{2}}\cdot\Delta\,d\varphi. \]
Mais par une première décomposition, on a
\[ \begin{aligned} \int\frac{1+2n-n\sin^{2}\varphi}{(1+n\sin^{2}\varphi)^{2}}\Delta\,d\varphi={}&\frac{c^{2}}{n}\int\frac{d\varphi}{\Delta}-\left(1+2c^{2}+\frac{3c^{2}}{n}\right)\int\frac{d\varphi}{(1+n\sin^{2}\varphi)\Delta}\\ &+(2+2n)\left(1+\frac{c^{2}}{n}\right)\int\frac{d\varphi}{(1+n\sin^{2}\varphi)^{2}\Delta}. \end{aligned} \]
Substituant ensuite pour la dernière intégrale la valeur que nous avons trouvée (art.\ 94), le second membre deviendra
\[ \frac{n\Delta\sin\varphi\cos\varphi}{1+n\sin^{2}\varphi}-\int\frac{c^{2}\,d\varphi\sin^{2}\varphi}{\Delta}+(n+1)\int\frac{d\varphi}{(1+n\sin^{2}\varphi)\Delta}. \]
Donc en prenant cette intégrale depuis $\varphi=0$ jusqu'à $\varphi=\frac{1}{2}\pi$, on

\origpage{175}
aura l'aire entière du cône
\[ Z^{1}=2k\cot\tfrac{1}{2}\theta[(n+1)\Pi^{1}(n,c)-F^{1}(c)+E^{1}(c)], \]
formule où l'on a $n=\cot^{2}\theta$.

Ayant déjà fait le plus petit apothème $AS=k$, si on fait le plus grand $BS=k'$, on aura immédiatement, en vertu des valeurs précédentes, $\operatorname{tang}\frac{1}{2}\theta=\sqrt{\left(\frac{k}{k'}\right)}$; de plus, si on appelle $\zeta$ l'angle $\frac{1}{2}(ASO+BSO)$ que fait la perpendiculaire $AO$\ednote{Printed $AO$; the perpendicular from the vertex is $SO$, an apparent misprint.} avec la ligne qui divise en deux également l'angle $ASB$, on aura $c=\sin\zeta$. Ainsi les quantités $c$, $k$, $\theta$, qui entrent comme élémens dans le résultat final, se déduisent immédiatement du triangle $SAB$.

Au reste, comme la fonction complète $\Pi^{1}(n,c)$ peut s'exprimer par des fonctions de la première et de la seconde espèce, on voit que la surface du cône oblique à base circulaire, ne dépend non plus que des fonctions elliptiques de la première et de la seconde espèce, de sorte qu'elle peut se mesurer par des arcs d'ellipse; théorème qui n'avait encore été démontré que dans quelques cas particuliers.

La surface du même cône étant développée sur un plan, il en résulte un secteur dont l'angle se détermine par la formule
\[ V=\int\frac{d\omega\sqrt{[h^{2}+(1-f\cos\omega)^{2}]}}{h^{2}+f^{2}+1-2f\cos\omega}, \]
et au moyen des mêmes substitutions, on a la transformée
\[ V=\frac{m}{k}\int\frac{\Delta\,d\varphi}{(1+n\sin^{2}\varphi)}=\frac{\operatorname{tang}\tfrac{1}{2}\theta}{k}\left[\left(1+\frac{c^{2}}{n}\right)\Pi(n,c,\varphi)-\frac{c^{2}}{n}F(c,\varphi)\right]. \]
Cette intégrale étant prise depuis $\varphi=0$ jusqu'à $\varphi=2\pi$, on aura l'angle total du secteur
\[ V^{1}=\frac{4\operatorname{tang}\tfrac{1}{2}\theta}{k}\left[\left(1+\frac{c^{2}}{n}\right)\Pi^{1}(n,c)-\frac{c^{2}}{n}F^{1}(c)\right]. \]

On peut remarquer que si dans l'expression de $V$ on fait $\varphi=\frac{1}{2}\pi$, la valeur de $V$ deviendra le quart de l'angle total $V^{1}$. On a en même temps $\omega=\theta$. Donc si on détermine le point $M$ de manière que l'angle $ACM=\theta$, l'angle correspondant $V$ sera précisément le quart de l'angle résultant du développement de la surface entière du cône.

\origpage{176}
\subsection*{(124).}

Supposons en second lieu que la base du cône soit une ellipse. Soient $h$ la hauteur du cône, $f$ et $g$ les coordonnées du point où la perpendiculaire $h$ rencontre le plan de la base, $f$ étant prise dans la direction du demi-grand axe de l'ellipse $1$, et $g$ dans celle du demi-axe conjugué $b$. Soit $\varphi$ l'amplitude d'un point quelconque de l'ellipse, les coordonnées de ce point seront $x=\sin\varphi$, $y=b\cos\varphi$, et l'élément de la courbe $ds=\Delta\,d\varphi$; si du sommet on mène une perpendiculaire sur la tangente en ce point, le quarré de cette perpendiculaire aura pour expression $h^{2}+\left(\frac{g\cos\varphi+bf\sin\varphi-b}{\Delta}\right)^{2}$, donc l'élément de la surface du cône sera
\[ dZ=\tfrac{1}{2}d\varphi\sqrt{[h^{2}(1-c^{2}\sin^{2}\varphi)+(g\cos\varphi+bf\sin\varphi-b)^{2}]}. \]
Cette différentielle paraît fort composée; cependant si on fait $\operatorname{tang}\frac{1}{2}\varphi=z$, ce qui donne $\sin\varphi=\frac{2z}{1+z^{2}}$, $\cos\varphi=\frac{1-z^{2}}{1+z^{2}}$, $d\varphi=\frac{2\,dz}{1+z^{2}}$, on aura la transformée
\[ dZ=\frac{dz}{(1+z^{2})^{2}}\sqrt{[h^{2}(1+z^{2})^{2}-4c^{2}h^{2}z^{2}+(g-b+2bfz-(g+b)z^{2})^{2}]}; \]
et puisque la variable sous le radical ne passe pas le quatrième degré, il est clair qu'on pourra trouver l'intégrale $Z$ au moyen des fonctions elliptiques. De plus, comme il suffit pour avoir la surface totale du cône, de connaître l'intégrale lorsque $\varphi=2\pi$, il semble que le résultat final ne doit dépendre que des fonctions complètes de la troisième espèce et des espèces inférieures, lesquelles, au moyen des réductions connues, ne dépendent elles-mêmes que des fonctions de la première et la seconde espèce; de sorte que l'aire entière du cône pourra encore être exprimée par des arcs d'ellipse. Mais il est nécessaire d'entrer dans quelques détails pour justifier cette conclusion.

\subsection*{(125).}

Pour faire disparaître les puissances impaires de la variable sous le radical dans la valeur de $dZ$, il faut d'abord supposer $z=\frac{p+qx}{1+x}$; et parce que la quantité sous le radical n'a, dans le cas dont il s'agit, que des facteurs simples imaginaires, le résultat de la substitution, après avoir déterminé $p$ et $q$, contiendra un nouveau radical de la

\origpage{177}
forme $\sqrt{[(1+\alpha^{2}x^{2})(1+\beta^{2}x^{2})]}$. Pour ramener ensuite la transformée aux fonctions elliptiques, il faudra, en supposant $\alpha>\beta$, faire $\alpha x=\operatorname{tang}\psi$, et $\psi$ deviendra l'amplitude de ces fonctions, tandis que leur module $c$ sera déterminé par l'équation $c^{2}=1-\frac{\beta^{2}}{\alpha^{2}}$. La relation entre $\varphi$ et $\psi$ sera donc telle qu'on aura $\operatorname{tang}\frac{1}{2}\varphi=\frac{p\alpha+q\operatorname{tang}\psi}{\alpha+\operatorname{tang}\psi}$. Or je dis que cette équation peut toujours se mettre sous la forme
\[ \operatorname{tang}(\psi-\nu)=A\operatorname{tang}\tfrac{1}{2}(\varphi-\mu), \]
$A$, $\mu$ et $\nu$ étant des quantités constantes. En effet si on fait $\operatorname{tang}\frac{1}{2}\mu=t$ et $\operatorname{tang}\nu=t'$, l'équation précédente donnera
\[ \frac{\operatorname{tang}\psi-t'}{1+t'\operatorname{tang}\psi}=\frac{A(\operatorname{tang}\tfrac{1}{2}\varphi-t)}{1+t\operatorname{tang}\tfrac{1}{2}\varphi}. \]
Substituant la valeur de $\operatorname{tang}\frac{1}{2}\varphi$ en fonction de $\operatorname{tang}\psi$, on aura
\[ \frac{\operatorname{tang}\psi-t'}{1+t'\operatorname{tang}\psi}=\frac{A(p\alpha+q\operatorname{tang}\psi-t\alpha-t\operatorname{tang}\psi)}{\alpha+\operatorname{tang}\psi+p\alpha t+qt\operatorname{tang}\psi}. \]
Cette équation devant avoir lieu quelle que soit $\operatorname{tang}\psi$, soit 1°. $\operatorname{tang}\psi=t'$, on aura
\[ t=\frac{p\alpha+qt'}{\alpha+t'}. \]
Soit 2°. $\operatorname{tang}\psi=-\frac{1}{t'}$, on aura
\[ t=\frac{1-\alpha t'}{p\alpha t'-q}. \]
Egalant ces deux valeurs de $t$, on aura pour déterminer $t'$, l'équation
\[ \frac{1-t'^{2}}{2t'}=\frac{\alpha^{2}(1+p^{2})-1-q^{2}}{2(1+pq)}=\cot2\nu. \]
\ednote{Printed denominator $2(1+pq)$; equating the two values of $t$ above gives $2\alpha(1+pq)$, an apparent misprint.}
On aurait semblablement pour déterminer $t$, l'équation
\[ \frac{1-t^{2}}{2t}=\frac{1-q^{2}+\alpha^{2}(1-p^{2})}{2(q-\alpha^{2}p)}=\cot\mu. \]
\ednote{Printed denominator $2(q-\alpha^{2}p)$; the two values of $t$ above give $2(q+\alpha^{2}p)$, an apparent misprint.}
3°. Enfin, pour que l'équation ci-dessus devienne entièrement identique,

\origpage{178}
soit $\operatorname{tang}\psi=\infty$, on aura
\[ A=\frac{1}{t'}\cdot\frac{1+qt}{q-t}. \]
Cela posé, l'équation $\operatorname{tang}(\psi-\nu)=A\operatorname{tang}\frac{1}{2}(\varphi-\mu)$ a lieu pour toutes les valeurs correspondantes de $\varphi$ et $\psi$; il en résulte qu'aux trois valeurs $\varphi=\mu$, $\varphi=\pi+\mu$, $\varphi=2\pi+\mu$, répondent ces trois valeurs de l'amplitude des fonctions elliptiques $\psi=\nu$, $\psi=\frac{1}{2}\pi+\nu$, $\psi=\pi+\nu$. Mais pour avoir l'aire entière du cône, on devra prendre l'intégrale $Z$ depuis $\varphi=0$ jusqu'à $\varphi=2\pi$; ou, ce qui revient au même, depuis $\varphi=\mu$ jusqu'à $\varphi=2\pi+\mu$; on devra donc prendre les fonctions elliptiques qui entrent dans l'expression de $Z$ depuis $\psi=\nu$ jusqu'à $\psi=\nu+\pi$; ou, ce qui revient au même, depuis $\psi=0$ jusqu'à $\psi=\pi$. Donc il n'entrera que des fonctions complètes dans l'expression de l'aire totale, et ces fonctions se réduisent toujours aux fonctions de la première et de la seconde espèce; d'où il suit que dans le cône oblique à base elliptique, l'aire entière peut encore s'exprimer par des arcs d'ellipse.

\section*{Construction de la ligne la plus courte sur la surface du sphéroïde.}

\subsection*{(126).}

Nous avons donné dans les Mémoires de l'Institut, an 1806, les équations nécessaires pour décrire la ligne la plus courte sur la surface d'un sphéroïde, et nous avons démontré que cette courbe, prolongée indéfiniment, est composée d'une infinité de spires égales et semblables, comprises entre deux parallèles également éloignés de l'équateur. Nous nous proposons maintenant de faire voir comment on détermine les différens points de cette courbe par le moyen des fonctions elliptiques, et pour cet objet, il suffira de construire une seule moitié de spire comprise depuis le parallèle auquel elle est tangente jusqu'à l'équateur.

Il faut d'abord se rappeler ce qui a été démontré dans le mémoire cité, que la ligne la plus courte menée entre deux points donnés sur la surface d'un sphéroïde, ligne qu'on appelle \emph{géodésique}, lorsqu'elle est tracée sur la surface du sphéroïde terrestre, fait toujours partie d'une ligne de cette nature perpendiculaire au méridien d'un

\origpage{179}
lieu déterminé, ou tangente au parallèle de ce lieu. Nous supposerons donc que par un point $A$ dont la latitude est donnée, on fasse passer une perpendiculaire au méridien du lieu $A$; il s'agit de construire cette courbe, c'est-à-dire de déterminer pour chaque point $M$ dont la latitude est donnée à volonté, les élémens inconnus du triangle $APM$, savoir, la distance $AM$ mesurée par l'arc de la ligne la plus courte, la longitude du point $M$ mesurée par l'angle $P$ du triangle $APM$, l'angle $M$ du même triangle qui représente en ce point la direction azimutale de la courbe; enfin l'arc $PM$ du méridien mobile qui sera connu par ses coordonnées. De ces quatre élémens, deux peuvent être déterminés algébriquement, savoir, l'angle $M$ et les coordonnées du point $M$ dans le plan $CMP$; les deux autres, savoir, l'arc $AM$ et l'angle $APM$ ne sont assignables que par les fonctions elliptiques.

\rmfigure{figures/fig-12.png}{Fig.~12.}{The vertical line PrTC and the horizontal line CDE, with an outer curve running from P through A and Q round E to N and on past I, a dashed curve from P to D, a curve from P through M to N, a curve from A through M to I, the straight line TM, the arc MQ, and a short dashed line from r ending at m.}

Soit le rayon de l'équateur $CA=A$, le demi-axe $CB=B$, la distance du centre au foyer de l'ellipse génératrice $\sqrt{(A^{2}-B^{2})}=A\delta$; soit $l$ la latitude du point $A$, $\lambda$ celle du point $M$, $P$ l'angle $APM$ qui mesure la longitude du point $M$, $M$ l'angle azimutal $AMP$; soient de plus les coordonnées $CT=t$, $TM=u$, et l'arc $AM=s$. Il s'agit de déterminer toutes les quantités $t$, $u$, $M$, $s$, $P$ en fonctions de la variable $\lambda$ et des quantités données.

Il est utile pour cet objet de substituer aux latitudes $l$ et $\lambda$ des latitudes réduites $l'$ et $\lambda'$, calculées par les formules
\[ \operatorname{tang} l'=\frac{B}{A}\operatorname{tang} l,\qquad \operatorname{tang}\lambda'=\frac{B}{A}\operatorname{tang}\lambda. \]
Cela posé, les élémens qui peuvent être déterminés algébriquement dans la question proposée, sont $t=B\sin\lambda'$, $u=A\cos\lambda'$, $\sin M=\frac{\cos l'}{\cos\lambda'}$. Il suit de la dernière équation que $l'$ est le \emph{maximum} de $\lambda'$, et $l$ celui de $\lambda$. On voit aussi par les valeurs des coordonnées que $\lambda'$ est l'amplitude de l'arc $PM$; de sorte que cet arc, si on en avait besoin, pourrait être exprimé par la formule $PM=A\cdot E(\delta,\lambda')$.

\subsection*{(127).}

Pour avoir les deux autres élémens $s$ et $P$, soit pris l'angle $\varphi$ tel que $\sin\lambda'=\sin l'\cos\varphi$, et on aura (Mémoire cité, n° 4) les équations différentielles

\origpage{180}
\[ \begin{aligned} ds&=d\varphi\sqrt{(B^{2}+A^{2}\delta^{2}\sin^{2}l'\cos^{2}\varphi)}\\ dP&=\frac{\cos l'}{A}\cdot\frac{d\varphi\sqrt{(B^{2}+A^{2}\delta^{2}\sin^{2}l'\cos^{2}\varphi)}}{1-\sin^{2}l'\cos^{2}\varphi}. \end{aligned} \]
Pour ramener ces expressions aux formes ordinaires des fonctions elliptiques, je fais $c^{2}=\frac{A^{2}\delta^{2}\sin^{2}l'}{B^{2}+A^{2}\delta^{2}\sin^{2}l'}$, ou plus simplement $c=\delta\sin l$, ce qui donne $\sqrt{(B^{2}+A^{2}\delta^{2}\sin^{2}l')}=\frac{B}{b}$, et j'ai d'abord $ds=\frac{B}{b}\,d\varphi\sqrt{(1-c^{2}\sin^{2}\varphi)}$; donc
\[ s=\frac{B}{b}E(c,\varphi). \]
Cette équation fait voir que tout arc $s$ de la ligne la plus courte peut être assimilé indéfiniment à un arc d'ellipse dont le demi-petit axe $CP=B$, et le demi-grand axe $CD=\frac{B}{b}$, quantité plus grande que $B$ et moindre que $A$, puisqu'on a $\frac{B}{b}=A\frac{\sin l'}{\sin l}$. Si sur ce quart d'ellipse $PmD$, on prend un point $m$ qui ait pour amplitude $\varphi$, ou qui soit déterminé par les coordonnées $Cr=B\cos\varphi$, $rm=\frac{B}{b}\sin\varphi$, l'arc $Pm$ de l'ellipse sera égal en longueur à l'arc $AM$ de la ligne la plus courte.

Ces déterminations ont lieu dans toute l'étendue de la courbe qu'on veut construire; mais il suffit d'examiner le cours de cette courbe depuis $\varphi=0$ jusqu'à $\varphi=90^{\circ}$, c'est-à-dire depuis le point $A$ où l'on a $\lambda=l$, jusqu'au point $I$ où la courbe rencontre l'équateur et où l'on a $\lambda=0$. Au point $A$ la courbe est perpendiculaire au méridien ou tangente au parallèle; au point $I$ elle coupe l'équateur sous un angle $I$, tel qu'on a $\cos I=\sin M=\cos l'$, de sorte que l'angle $I$ est égal à la latitude réduite $l'$.

Dans ce même point, la courbe $AMI$ est égale au quart d'ellipse $BmD$\ednote{Printed ``BmD''; the quarter-ellipse named above is $PmD$, an apparent misprint.}, et on a $s=\frac{B}{b}E^{1}(c)$.

\subsection*{(128).}

Venons maintenant à la seconde formule qui donne la longitude $P$, on aura d'abord, en faisant $n=\operatorname{tang}^{2}l'$, $M=\frac{B}{bA\cos l'}=\frac{B}{A\cos l}$
\[ dP=M\cdot\frac{\Delta\,d\varphi}{1+n\sin^{2}\varphi}, \]

\origpage{181}
donc en intégrant
\[ P=\frac{M}{n}[(n+c^{2})\Pi(n,c,\varphi)-c^{2}F(c,\varphi)]. \]
Par cette formule on déterminera pour chaque point $M$ dont la latitude est connue, la longitude $P$ de ce point; et les calculs pourront être faits avec tout le degré d'exactitude que comportent les tables de logarithmes, au moyen des méthodes que nous avons données pour évaluer les fonctions $F$ et $\Pi$.

Si on fait $\varphi=90^{\circ}$, on aura la position du point $I$ où la courbe coupe l'équateur par la formule
\[ P^{1}=\frac{M}{n}[(n+c^{2})\Pi^{1}(n,c)-c^{2}F^{1}(c)]; \]
et comme la fonction complète $\Pi^{1}(n,c)$ peut s'exprimer par des fonctions de la première et de la seconde espèce, il s'ensuit qu'on peut déterminer le point $I$ par ces seules fonctions; ou, ce qui revient au même, par les seuls arcs d'ellipse. On pourrait déterminer de même une infinité d'autres points de la ligne la plus courte; et lorsqu'on aura tracé la demi-spire $AMI$, on pourra de même tracer la continuation de cette courbe dans l'autre demi-sphéroïde, en observant que les points situés de part et d'autre à des latitudes égales, ont avec le point $I$ des différences égales en longitude.

\subsection*{(129).}

Il est remarquable que la formule qui donne la valeur de l'angle $P$, est absolument de la même forme que celle qui donne la valeur de l'angle résultant du développement sur un plan d'une portion de la surface d'un cône oblique à base circulaire (123).

Ce dernier angle est facile à trouver jusqu'à un certain degré d'approximation, par de simples constructions géométriques, déduites des pyramides inscrites et circonscrites. D'ailleurs comme le cône oblique n'a de courbure que dans un sens, cette surface est censée plus simple que celle du sphéroïde qui est partout à double courbure; ainsi on doit regarder le problème de déterminer la ligne la plus courte sur la surface du sphéroïde, comme étant réduit à une moindre difficulté par la construction suivante, qu'on tire aisément de nos formules.

\origpage{182}
\rmfigure{figures/fig-13.png}{Fig.~13.}{A triangle pfp' with the line pq from p to the point q on fp', the point c on pq, a dotted curve from p passing below the base pp' and up through n to q, and the dotted lines cn and fn.}

Faites un triangle isoscèle $pfp'$, tel que sa base $pp'$ représentant\ednote{Printed ``représentant''; ``représente'' is expected, an apparent misprint.} l'axe $2B$, le côté $pf$ ou $p'f$ représente le rayon de l'équateur $A$; tirez $pq$ de manière que l'angle $p'pq$ ait pour sinus $c$, ou qu'on ait $\sin p'pq=\sin l\sin fpp'$. Sur $pq$ comme diamètre, décrivez un cercle $qnp$ dont le plan soit perpendiculaire au plan du triangle $fpq$. Le cône qui a ce cercle pour base, et pour sommet le point $f$, est celui dont la surface étant développée sur un plan, servira à décrire la ligne la plus courte sur la surface du sphéroïde donné. Pour cela, soit $QM$ un parallèle dont la latitude est $\lambda$, on déterminera successivement les angles $\lambda'$, $\varphi$ et $\omega$, soit par les équations $\operatorname{tang}\lambda'=\frac{B}{A}\operatorname{tang}\lambda$, $\cos\varphi=\frac{\sin\lambda'}{\sin l'}$, $\operatorname{tang}\frac{1}{2}\omega=\operatorname{tang}(45^{\circ}-\frac{1}{2}l')\cdot\operatorname{tang}\frac{1}{2}\varphi$, soit par des constructions géométriques qui les représentent. L'angle $\omega$ étant trouvé, on mènera dans le plan de la base le rayon $cn$, de manière que l'angle $qcn=\omega$; tirant ensuite la droite $fn$, et développant sur un plan la surface $qfn$, l'angle $f$ du secteur produit par le développement, sera égal à l'angle $P$ du triangle sphéroïdique $APM$, et servira ainsi à déterminer la longitude du point $M$.

Lorsque $\varphi=90^{\circ}$, on aura $\omega=90^{\circ}-l'$; alors l'angle produit par le développement du secteur $qfn$ sera la moitié de l'angle produit par le développement de toute la surface $fpnq$; et comme ce dernier angle est toujours moindre que deux angles droits, il s'ensuit que l'angle $P$, mesuré alors par l'arc $EI$, est moindre qu'un angle droit, ce qui s'accorde avec les propriétés connues de la ligne la plus courte tracée sur un sphéroïde aplati.

\section*{Détermination de l'aire de l'ellipsoïde.}

\subsection*{(130).}

Soit l'équation de la surface proposée,
\[ \frac{x^{2}}{a^{2}}+\frac{y^{2}}{b^{2}}+\frac{z^{2}}{c^{2}}=1; \]
l'expression de l'aire, comprise entre des limites données, dépendra de la double intégrale
\[ S=\iint dx\,dy\sqrt{\left(\frac{1+\frac{c^{2}-a^{2}}{a^{2}}\cdot\frac{x^{2}}{a^{2}}+\frac{c^{2}-b^{2}}{b^{2}}\cdot\frac{y^{2}}{b^{2}}}{1-\frac{x^{2}}{a^{2}}-\frac{y^{2}}{b^{2}}}\right)}; \]

\origpage{183}
mais l'une ou l'autre des deux intégrations ne peut s'effectuer sans introduire des fonctions elliptiques qui ne permettraient pas d'obtenir la seconde intégrale. Pour rendre la première intégration possible, au moins par arcs de cercle ou par logarithmes, on pourrait appliquer la méthode que nous avons donnée dans les Mémoires de l'Académie des Sciences, année 1788; mais on parviendra plus directement au même but de la manière suivante.

Soit $z=c\cos\theta$, on aura $\frac{x^{2}}{a^{2}}+\frac{y^{2}}{b^{2}}=\sin^{2}\theta$; soit ensuite $y=b\sin\theta\cos\varphi$, on aura $x=a\sin\theta\sin\varphi$. D'après ces équations, qui équivalent à l'équation de la surface, il faudra exprimer l'élément de l'aire en fonction des deux nouvelles variables $\theta$ et $\varphi$. Et d'abord pour avoir la valeur de $dx\,dy$, je différentie l'équation $x=a\sin\theta\sin\varphi$, en supposant $\theta$ constant, j'ai $dx=a\,d\varphi\cos\varphi\sin\theta$; ensuite il faut avoir la valeur de $dy$, en supposant $x$ constante; c'est ce que donnera l'équation $\frac{x^{2}}{a^{2}}+\frac{y^{2}}{b^{2}}=\sin^{2}\theta$, d'où l'on tire $y\,dy=b^{2}d\theta\sin\theta\cos\theta$; donc $dx\,dy=ab\,d\varphi\,d\theta\sin\theta\cos\theta$.

Substituant cette valeur ainsi que celles de $x$ et $y$ dans l'expression de l'aire $S$; faisant de plus, pour abréger $\frac{a^{2}-c^{2}}{a^{2}}=\delta$, $\frac{b^{2}-c^{2}}{b^{2}}=\varepsilon$, on aura
\[ S=\iint ab\,d\varphi\,d\theta\sin\theta\sqrt{[1-(\delta\sin^{2}\varphi+\varepsilon\cos^{2}\varphi)\sin^{2}\theta]}. \]
L'intégrale, par rapport à $\theta$, est facile à trouver par logarithmes; mais comme la formule qui en résulterait pour la seconde intégration, n'est pas du nombre de celles qu'on sait ramener aux fonctions connues, nous nous contenterons d'intégrer par séries.

Soit donc $\delta\sin^{2}\varphi+\varepsilon\cos^{2}\varphi=p$, on aura, en développant le radical,
\[ S=\iint ab\,d\varphi\,d\theta\sin\theta\left(1-\tfrac{1}{2}p\sin^{2}\theta-\tfrac{1\cdot1}{2\cdot4}p^{2}\sin^{4}\theta-\tfrac{1\cdot1\cdot3}{2\cdot4\cdot6}p^{3}\sin^{6}\theta-\text{etc.}\right). \]
Or en intégrant depuis $\theta=0$ jusqu'à $\theta=\frac{1}{2}\pi$, on a
\[ \int d\theta\sin\theta=1,\qquad \int d\theta\sin^{3}\theta=\frac{2}{3},\qquad \int d\theta\sin^{5}\theta=\frac{2\cdot4}{3\cdot5},\qquad \text{etc.} \]
Donc il résulte de la première intégration,

\origpage{184}
\[ S=\int ab\,d\varphi\left(1-\tfrac{1}{3}p-\tfrac{1}{3\cdot5}p^{2}-\tfrac{1}{5\cdot7}p^{3}-\tfrac{1}{7\cdot9}p^{4}-\text{etc.}\right). \]
Cette formule doit être intégrée de nouveau depuis $\varphi=0$ jusqu'à $\varphi=\frac{1}{2}\pi$; soient donc dans ces limites $\displaystyle\int p\,d\varphi=\frac{\pi}{2}\cdot P'$, $\displaystyle\int p^{2}\,d\varphi=\frac{\pi}{2}\cdot P''$, $\displaystyle\int p^{3}\,d\varphi=\frac{\pi}{2}\cdot P'''$, etc.; et l'on aura, après avoir multiplié par 8, l'aire totale de l'ellipsoïde
\[ 8S=4\pi ab\left(1-\tfrac{1}{3}P'-\tfrac{1}{3\cdot5}P''-\tfrac{1}{5\cdot7}P'''-\tfrac{1}{7\cdot9}P^{\text{iv}}-\text{etc.}\right), \]
formule dans laquelle il reste à substituer les valeurs suivantes:
\[ \begin{aligned} P'&=\delta\cdot\tfrac{1}{2}+\varepsilon\cdot\tfrac{1}{2},\\ P''&=\delta^{2}\cdot\tfrac{1\cdot3}{2\cdot4}+\delta\varepsilon\cdot\tfrac{1}{2}\cdot\tfrac{1}{2}+\varepsilon^{2}\cdot\tfrac{1\cdot3}{2\cdot4},\\ P'''&=\delta^{3}\cdot\tfrac{1\cdot3\cdot5}{2\cdot4\cdot6}+\delta^{2}\varepsilon\cdot\tfrac{1\cdot3}{2\cdot4}\cdot\tfrac{1}{2}+\delta\varepsilon^{2}\cdot\tfrac{1}{2}\cdot\tfrac{1\cdot3}{2\cdot4}+\varepsilon^{3}\cdot\tfrac{1\cdot3\cdot5}{2\cdot4\cdot6},\\ &\text{etc.} \end{aligned} \]
La loi de ces expressions est manifeste, et on peut observer que le terme général $P^{(n)}$ est le coefficient de $z^{n}$ dans le développement de $(1-\delta z)^{-\frac{1}{2}}(1-\varepsilon z)^{-\frac{1}{2}}$, ou dans celui du produit
\[ \left(1+\tfrac{1}{2}\delta z+\tfrac{1\cdot3}{2\cdot4}\delta^{2}z^{2}+\tfrac{1\cdot3\cdot5}{2\cdot4\cdot6}\delta^{3}z^{3}+\text{etc.}\right)\left(1+\tfrac{1}{2}\varepsilon z+\tfrac{1\cdot3}{2\cdot4}\varepsilon^{2}z^{2}+\tfrac{1\cdot3\cdot5}{2\cdot4\cdot6}\varepsilon^{3}z^{3}+\text{etc.}\right). \]

Au reste, la suite $P'$, $P''$, $P'''$, etc. sera nécessairement convergente si $\delta$ et $\varepsilon$ sont tous deux plus petits que l'unité, ce qui aura toujours lieu en prenant pour $c$ la plus petite des trois quantités $a$, $b$, $c$.

\subsection*{(131).}

Le problème de déterminer l'aire totale de l'ellipsoïde est résolu par la suite qu'on vient de trouver, aussi simplement qu'il peut l'être quand on n'a pour but que d'obtenir une approximation. Mais il faut recourir à d'autres moyens, si on veut avoir un résultat dépendant seulement des fonctions elliptiques ou des quadratures algébriques.

\origpage{185}
Dans la méthode précédente, nous avons partagé l'ellipse qui sert de base à la surface proposée, en zônes elliptiques concentriques, déterminées par les projections des sections faites dans l'ellipsoïde, parallèlement à sa base. Toute autre manière de partager la base est également admissible et doit conduire au même résultat après les deux intégrations; mais il en est une surtout qui mérite d'être soumise au calcul, et qui semble promettre des résultats élégans. Je veux parler des projections qui naissent des lignes de plus grande et de plus petite courbure tracées sur la surface de l'ellipsoïde.

Monge, dans ses Feuilles d'Analyse appliquée (édit.\ de l'an 9, n° 19), a donné la construction de ces courbes, comme il suit.

Ayant supposé $a>b$ et $b>c$, soit pris pour plan de projection le plan des $x$ et $y$, qui est celui de la section principale dont $a$ et $b$ sont les demi-axes. Soient décrites dans ce plan une ellipse et une hyperbole auxiliaires, qui aient les demi-axes communs $A$ et $B$ ainsi déterminés,
\[ \begin{aligned} A&=a\sqrt{\left(\frac{a^{2}-b^{2}}{a^{2}-c^{2}}\right)},\\ B&=b\sqrt{\left(\frac{a^{2}-b^{2}}{b^{2}-c^{2}}\right)}, \end{aligned} \]
le demi-axe $A$ étant dirigé dans le sens du demi-axe $a$, et $B$ dans le sens de $b$.

L'équation de l'hyperbole sera $y^{2}=\frac{B^{2}}{A^{2}}(x^{2}-A^{2})$, et l'équation de l'ellipse, $y^{2}=\frac{B^{2}}{A^{2}}(A^{2}-x^{2})$; elles se toucheront par le sommet où $x=A$, et d'ailleurs on peut observer qu'en vertu des suppositions faites on a $A<a$.

Cela posé, soient $x=\alpha$, $y=\beta$, les coordonnées d'un point pris à volonté sur l'hyperbole, ensorte qu'on ait $\beta^{2}=\frac{B^{2}}{A^{2}}(\alpha^{2}-A^{2})$; si, avec les demi-axes $\alpha$ et $\beta$, pris dans les mêmes directions que $a$ et $b$, on décrit une ellipse, cette ellipse et toutes les ellipses décrites de la même manière suivant les diverses valeurs de $\alpha$ et $\beta$, avec la seule condition qu'on ait $\alpha<a$, seront les projections sur le plan des $x$ et $y$, de toutes les lignes de courbure d'une même espèce, tracées sur la surface de l'ellipsoïde.

\origpage{186}
Les projections des lignes de courbure de l'autre espèce, qui seront des hyperboles, se détermineront semblablement en prenant pour demi-axes de chaque hyperbole, les coordonnées $\alpha'$, $\beta'$ d'un même point de l'ellipse auxiliaire, lesquelles doivent satisfaire à l'équation $\beta'^{2}=\frac{B^{2}}{A^{2}}(A^{2}-\alpha'^{2})$.

Les projections des lignes de courbure de la première espèce seront donc les ellipses tracées d'après l'équation
\[ y^{2}=\frac{\beta^{2}}{\alpha^{2}}(\alpha^{2}-x^{2})=\frac{B^{2}}{A^{2}}(\alpha^{2}-x^{2})\left(1-\frac{A^{2}}{\alpha^{2}}\right), \]
en donnant à $\alpha$ toutes les valeurs depuis $\alpha=A$ jusqu'à $\alpha=a$.

Et les projections des lignes de courbure de la seconde espèce seront toutes les hyperboles décrites d'après l'équation
\[ y^{2}=\frac{\beta'^{2}}{\alpha'^{2}}(x^{2}-\alpha'^{2})=\frac{B^{2}}{A^{2}}(x^{2}-\alpha'^{2})\left(\frac{A^{2}}{\alpha'^{2}}-1\right), \]
en donnant à $\alpha'$ toutes les valeurs depuis $\alpha'=0$ jusqu'à $\alpha'=A$.

Considérons plus particulièrement les ellipses qui sont les projections des lignes de courbure de la première espèce.

Lorsqu'on fait $\alpha=A$, on a $\beta=0$, $y=0$, de sorte que l'ellipse est infiniment étroite, et se réduit à son grand axe. A mesure que le demi-axe $\alpha$ augmente, l'autre demi-axe $\beta$ augmente aussi; et enfin lorsqu'on fait $\alpha=a$, l'ellipse de projection se confond avec la base de l'ellipsoïde, ce qui est le dernier terme des projections.

\subsection*{(132).}

Cela posé, cherchons l'aire de l'ellipsoïde qui répond à l'élément de la base $dx\,dy$, compris entre deux ellipses de projection infiniment proches. L'une de ces ellipses ayant pour équation $y^{2}=\frac{\beta^{2}}{\alpha^{2}}(\alpha^{2}-x^{2})$, on pourra faire $x=\alpha\sin\psi$, ce qui donnera $y=\beta\cos\psi$. Différentiant $x$ en supposant $\alpha$ constante, on aura $dx=\alpha\,d\psi\cos\psi$; si on passe ensuite de cette ellipse à l'ellipse infiniment voisine, il faudra différentier l'équation
\[ y^{2}=\frac{B^{2}}{A^{2}}(\alpha^{2}-x^{2})\left(1-\frac{A^{2}}{\alpha^{2}}\right), \]
en faisant varier $\alpha$ seul, ce qui donnera $y\,dy=\frac{B^{2}}{A^{2}}\alpha\,d\alpha\left(1-\frac{A^{2}x^{2}}{\alpha^{4}}\right)$,

\origpage{187}
donc l'élément de la projection
\[ dx\,dy=\frac{B^{2}}{A^{2}}\cdot\frac{d\alpha\,d\psi}{\beta}(\alpha^{2}-A^{2}\sin^{2}\psi). \]
Lorsqu'on fait varier $\psi$ en supposant $\alpha$ constant, le point de projection ne varie que sur une même ellipse de projection; de même lorsqu'on fait varier $\alpha$ en supposant $\psi$ constant, le point de projection ne varie que sur une hyperbole de projection; car en éliminant $\alpha$ et $\beta$ des équations $\alpha=\frac{x}{\sin\psi}$, $\beta=\frac{y}{\cos\psi}$, $\beta^{2}=\frac{B^{2}}{A^{2}}(\alpha^{2}-A^{2})$, on a $y^{2}=\frac{B^{2}\cos^{2}\psi}{A^{2}\sin^{2}\psi}(x^{2}-A^{2}\sin^{2}\psi)$, ce qui est l'équation d'une hyperbole de projection, dont les demi-axes sont $\alpha'=A\sin\psi$, $\beta'=B\cos\psi$.

De là on voit que l'élément de projection dont nous venons de donner la valeur, représente le quadrilatère compris entre deux ellipses et deux hyperboles de projection infiniment voisines.

Si on substitue maintenant les valeurs de $x$, $y$ et $dx\,dy$ dans l'expression générale de $S$, on aura
\[ S=\iint\frac{B}{A}\cdot\frac{\alpha^{2}-A^{2}\sin^{2}\psi}{\sqrt{(\alpha^{2}-A^{2})}}\,d\alpha\,d\psi\sqrt{\left[\frac{1-\delta\frac{\alpha^{2}\sin^{2}\psi}{a^{2}}-\varepsilon\frac{\beta^{2}\cos^{2}\psi}{b^{2}}}{1-\frac{\alpha^{2}\sin^{2}\psi}{a^{2}}-\frac{\beta^{2}\cos^{2}\psi}{b^{2}}}\right]}. \]
Cette expression paraît au premier coup d'œil beaucoup plus compliquée que celle à laquelle nous avons été conduits par la première méthode; mais en l'examinant de plus près, on trouve qu'elle est susceptible d'une réduction fort remarquable.

En effet, si on substitue la valeur $\beta^{2}=\frac{B^{2}}{A^{2}}(\alpha^{2}-A^{2})$, on reconnaîtra que la quantité sous le radical se décompose en deux facteurs distincts, l'un qui ne dépend que de $\alpha$, l'autre qui ne dépend que de $\psi$, de sorte qu'on a
\[ S=\iint\frac{B}{A}\cdot\frac{\alpha^{2}-A^{2}\sin^{2}\psi}{\sqrt{(\alpha^{2}-A^{2})}}\,d\alpha\,d\psi\sqrt{\left[\frac{\sin^{2}\psi+\frac{a^{2}}{b^{2}}\cos^{2}\psi}{\sin^{2}\psi+\frac{B^{2}a^{2}}{A^{2}b^{2}}\cos^{2}\psi}\cdot\frac{1-\delta\frac{\alpha^{2}}{a^{2}}}{1-\frac{\alpha^{2}}{a^{2}}}\right]}. \]
Si on observe maintenant que pour l'objet proposé il faut prendre l'intégrale par rapport à $\psi$, depuis $\psi=0$ jusqu'à $\psi=\frac{1}{2}\pi$, et

\origpage{188}
l'intégrale par rapport à $\alpha$, depuis $\alpha=A$ jusqu'à $\alpha=a$; on verra que ces deux intégrales sont entièrement indépendantes l'une de l'autre, et qu'elles peuvent être cherchées isolément, ce qui permettra de réduire l'aire entière de l'ellipsoïde à des quadratures algébriques, et même, comme on le prouvera ensuite, aux seules fonctions elliptiques.

\subsection*{(133).}

Soient donc
\[ \begin{aligned} M&=\int d\psi\sqrt{\left(\frac{\sin^{2}\psi+\frac{a^{2}}{b^{2}}\cos^{2}\psi}{\sin^{2}\psi+\frac{B^{2}a^{2}}{A^{2}b^{2}}\cos^{2}\psi}\right)}\\ N&=\int d\psi\sin^{2}\psi\sqrt{\left(\frac{\sin^{2}\psi+\frac{a^{2}}{b^{2}}\cos^{2}\psi}{\sin^{2}\psi+\frac{B^{2}a^{2}}{A^{2}b^{2}}\cos^{2}\psi}\right)}; \end{aligned} \]
ces intégrales devant être prises depuis $\psi=0$ jusqu'à $\psi=\frac{1}{2}\pi$.

Soient encore
\[ \begin{aligned} P&=\int\frac{\alpha^{2}d\alpha}{\sqrt{(\alpha^{2}-A^{2})}}\sqrt{\left(\frac{a^{2}-\delta\alpha^{2}}{a^{2}-\alpha^{2}}\right)}\\ Q&=\int\frac{d\alpha}{\sqrt{(\alpha^{2}-A^{2})}}\sqrt{\left(\frac{a^{2}-\delta\alpha^{2}}{a^{2}-\alpha^{2}}\right)}, \end{aligned} \]
ces intégrales devant être prises depuis $\alpha=A$ jusqu'à $\alpha=a$.

La surface cherchée sera
\[ S=\frac{B}{A}PM-ABQN. \]
Ainsi tout se réduit à trouver les quatre intégrales $P$, $Q$, $M$, $N$, dans les limites désignées.

Pour simplifier les formules $M$ et $N$, soit $\operatorname{tang}\psi=\frac{a}{b}\operatorname{tang}\omega$, ensuite $n=\frac{a^{2}-b^{2}}{b^{2}}$ et $c'^{2}=1-\frac{A^{2}}{B^{2}}=\frac{c^{2}}{b^{2}}\cdot\frac{a^{2}-b^{2}}{a^{2}-c^{2}}$, il résultera de ces substitutions
\[ \begin{aligned} M&=\frac{aA}{bB}\int\frac{d\omega}{(1+n\sin^{2}\omega)\sqrt{(1-c'^{2}\sin^{2}\omega)}}=\frac{aA}{bB}\Pi(n,c',\omega)\\ N&=\frac{a^{3}A}{b^{3}B}\int\frac{d\omega\sin^{2}\omega}{(1+n\sin^{2}\omega)^{2}\sqrt{(1-c'^{2}\sin^{2}\omega)}}. \end{aligned} \]

\origpage{189}
Mais par les réductions connues, on a
\[ \begin{aligned} \int\frac{d\omega\sin^{2}\omega}{(1+n\sin^{2}\omega)^{2}\Delta(c',\omega)}&=-\frac{n\sin\omega\cos\omega\,\Delta(c',\omega)}{2(n+1)(n+c'^{2})(1+n\sin^{2}\theta)}+\frac{F(c',\theta)}{2n(n+1)}\\ &\quad-\frac{E(c',\theta)}{2(n+1)(n+c'^{2})}+\frac{n^{2}-c'^{2}}{2n(n+1)(n+c'^{2})}\Pi(n,c',\theta). \end{aligned} \]
Donc en étendant les intégrales jusqu'à $\theta=\frac{1}{2}\pi$, on aura pour les valeurs complètes de $M$ et $N$,
\[ \begin{aligned} M&=\frac{aA}{bB}\Pi^{1}(n,c')\\ N&=\frac{aA}{2bBn}\left[F^{1}(c')-\left(\frac{a^{2}-c^{2}}{a^{2}}\right)E^{1}(c')+\left(\frac{a^{2}-b^{2}-c^{2}}{b^{2}}\right)\Pi^{1}(n,c')\right]. \end{aligned} \]

\subsection*{(134).}

A l'égard des intégrales $P$, $Q$, comme $\alpha$ est toujours compris entre $A$ et $a$, si l'on fait $A=a\sin\mu$, on pourra prendre
\[ \alpha^{2}=\frac{a^{2}\sin^{2}\mu}{1-\cos^{2}\mu\sin^{2}\zeta}, \]
et la substitution donnera
\[ \begin{aligned} P&=ab\sin^{2}\mu\int\frac{d\zeta\sqrt{(1-b'^{2}\sin^{2}\zeta)}}{(1-\cos^{2}\mu\sin^{2}\zeta)^{2}}\\ Q&=\frac{b}{a}\int\frac{d\zeta\sqrt{(1-b'^{2}\sin^{2}\zeta)}}{1-\cos^{2}\mu\sin^{2}\zeta}. \end{aligned} \]
On a fait dans ces formules $b'^{2}=\frac{a^{2}}{b^{2}}\cdot\frac{b^{2}-c^{2}}{a^{2}-c^{2}}=\frac{A^{2}}{B^{2}}$, de sorte qu'on a $b'^{2}+c'^{2}=1$, et qu'ainsi les modules $b'$ et $c'$ sont complémens l'un de l'autre.

Soit le paramètre $n=\frac{a^{2}-b^{2}}{b^{2}}=\cot^{2}\lambda$, on aura $b=a\sin\lambda$, et $b'^{2}=\frac{a^{2}}{b^{2}}\cos^{2}\mu=\frac{\cos^{2}\mu}{\sin^{2}\lambda}$, ou $\cos^{2}\mu=b'^{2}\sin^{2}\lambda$; l'intégrale $Q$ se rapporte donc aux fonctions elliptiques de la troisième espèce, dont le paramètre $n'=-\cos^{2}\mu=-b'^{2}\sin^{2}\lambda$, et on trouvera
\[ Q=\frac{1}{\sin\lambda}F(b',\zeta)-\frac{\cos^{2}\lambda}{\sin\lambda}\Pi(n',b',\zeta). \]
Pour avoir la valeur de $P$, il faudra d'abord employer cette formule

\origpage{190}
de réduction
\[ \begin{aligned} \int\frac{d\zeta\sqrt{(1-b'^{2}\sin^{2}\zeta)}}{(1+n'\sin^{2}\zeta)^{2}}&=\frac{n'\sin\zeta\cos\zeta\,\Delta(b',\zeta)}{2(1+n')(1+n'\sin^{2}\zeta)}-\frac{(b'^{2}+n')}{2n'(1+n')}F(b',\zeta)\\ &\quad+\frac{E(b',\zeta)}{2(1+n')}+\frac{n'^{2}+2n'+b'^{2}}{2n'(1+n')}\Pi(n',b',\zeta). \end{aligned} \]
Faisait\ednote{Printed ``Faisait''; ``Faisant'' is expected, an apparent misprint.} ensuite $\zeta=\frac{1}{2}\pi$, on aura les intégrales complètes
\[ \begin{aligned} P&=\frac{a^{2}}{2\sin\lambda}[\cos^{2}\lambda F^{1}(b')+\sin^{2}\lambda E^{1}(b')-(1-2\sin^{2}\lambda+b'^{2}\sin^{4}\lambda)\Pi^{1}(n',b')]\\ Q&=\frac{1}{\sin\lambda}F^{1}(b')-\frac{\cos^{2}\lambda}{\sin\lambda}\Pi^{1}(n',b'). \end{aligned} \]

\subsection*{(135).}

Maintenant si dans l'expression de l'aire $S$ on substitue les valeurs trouvées de $M$, $N$, $P$, $Q$, il en résultera
\[ \begin{aligned} S={}&\frac{a^{2}}{2}\Pi^{1}(n,c')[\cot^{2}\lambda F^{1}(b')+E^{1}(b')+(\sin^{2}\mu-\cot^{2}\lambda)\Pi^{1}(n',b')]\\ &-\frac{a^{2}}{2}[F^{1}(b')-\cos^{2}\lambda\Pi^{1}(n',b')]\left[\frac{\sin^{2}\mu}{\cos^{2}\lambda}F^{1}(c')-E^{1}(c')+\frac{\cot^{2}\lambda-\sin^{2}\mu}{\cos^{2}\lambda}\Pi^{1}(n',c')\right]. \end{aligned} \]
Par un premier développement qui fait disparaître les termes où les deux fonctions $\Pi^{1}$ sont multipliées, on obtient
\[ \begin{aligned} S={}&\frac{a^{2}}{2}\Pi^{1}(n,c')\left[\left(\frac{\sin^{2}\mu}{\cos\lambda}-1\right)F^{1}(b')+E^{1}(b')\right]\\ &-\frac{a^{2}}{2}\left[\frac{\sin^{2}\mu}{\cos^{2}\lambda}F^{1}(c')-E^{1}(c')\right][F^{1}(b')-\cos^{2}\lambda\Pi(n',b')]. \end{aligned} \]
Mais par les valeurs connues des fonctions $\Pi^{1}$, on a
\[ \begin{aligned} \Pi^{1}(n,c')={}&\sin^{2}\lambda F^{1}(c')+\frac{\sin\lambda\cos\lambda}{\Delta(b',\lambda)}[\tfrac{1}{2}\pi+F^{1}(c)F(b',\lambda)-F^{1}(c)E(b',\lambda)]\\ &-E^{1}(c)F(b',\lambda)\\ F^{1}(b)-\cos^{2}\lambda\Pi^{1}(n',b')={}&\sin^{2}\lambda F^{1}(b')+\frac{\sin\lambda\cos\lambda}{\Delta(b',\lambda)}[F^{1}(b')E(b',\lambda)\\ &-E^{1}(b')F(b',\lambda)]. \end{aligned} \]
Substituant et faisant les réductions auxquelles donne lieu la formule connue $\frac{\pi}{2}=F^{1}(c')E^{1}(b')+F^{1}(b')E^{1}(c')-F^{1}(b')F^{1}(c')$, on

\origpage{191}
trouvera
\[ \begin{aligned} S={}&\frac{\pi a^{2}}{4}\cdot\frac{\sin\lambda\cos\lambda}{\Delta(b',\lambda)}\left[\left(\frac{\sin^{2}\mu}{\cos^{2}\lambda}-1\right)[F^{1}(b')-F(b',\lambda)]\right.\\ &\left.+E^{1}(b')-E(b',\lambda)+\frac{\sin\lambda}{\cos\lambda}\Delta(b',\lambda)\right]. \end{aligned} \]
J'observe qu'on peut encore simplifier ce résultat en prenant un nouvel angle $\nu$ tel que $F(b',\lambda)+F(b',\nu)=F^{1}(b')$; il faut pour cela qu'on ait $c'\operatorname{tang}\lambda\operatorname{tang}\nu=1$, ce qui donne $\cos\nu=\frac{c}{a}$; on aura en même temps $E(b',\lambda)+E(b',\nu)=E^{1}(b')+\frac{\sin\nu\cos\nu}{\Delta(b',\nu)}$. Donc enfin l'aire entière de l'ellipsoïde aura pour expression
\[ 8S=2\pi a^{2}\cdot\frac{\sin\nu\cos\nu}{\Delta(b',\nu)}\left[\frac{\cos^{2}\nu}{\sin^{2}\nu}F(b',\nu)+E(b',\nu)+\frac{\cos\nu}{\sin\nu}\Delta(b',\nu)\right]. \]
D'ailleurs comme on a $b'^{2}=\frac{a^{2}}{b^{2}}\cdot\frac{b^{2}-c^{2}}{a^{2}-c^{2}}=\frac{b^{2}-c^{2}}{b^{2}\sin^{2}\nu}$, il en résulte $\Delta(b',\nu)=\frac{c}{b}$, de sorte que la valeur de $8S$ se réduit à cette forme très-simple,
\[ 8S=2\pi c^{2}+\frac{2\pi ab}{\sin\nu}\left[\frac{c^{2}}{a^{2}}F(b',\nu)+\frac{a^{2}-c^{2}}{a^{2}}E(b',\nu)\right], \]
formule dont on pourra faire l'application immédiate, en se rappelant seulement qu'on a $\cos\nu=\frac{c}{a}$ et $b'^{2}=\frac{b^{2}-c^{2}}{b^{2}\sin^{2}\nu}$.

Il est facile de vérifier cette formule dans les deux cas où l'ellipsoïde devient un solide de révolution; car si l'on a $a=b$, la formule donne pour l'aire du sphéroïde aplati,
\[ 8S=2\pi a^{2}+\frac{\pi b^{2}}{\sin\nu}\log\left(\frac{1+\sin\nu}{1-\sin\nu}\right); \]
et si l'on a $b=c$, la formule donne pour l'aire du sphéroïde alongé,
\[ 8S=\frac{2\pi ab}{\sin\nu}(\nu+\sin\nu\cos\nu), \]
ce qui s'accorde avec les résultats connus.

Nous sommes donc parvenus à exprimer par des fonctions elliptiques très-simples, l'aire entière de l'ellipsoïde qui, d'après les

\origpage{192}
autres méthodes, semblait devoir être une transcendante beaucoup plus composée.

\subsection*{(136).}

Il résulte de cette solution que la série trouvée par la première méthode,
\[ 4\pi ab\left(1-\tfrac{1}{3}P'-\tfrac{1}{3\cdot5}P''-\tfrac{1}{5\cdot7}P'''-\text{etc.}\right), \]
est susceptible d'être sommée à l'aide des fonctions elliptiques; et c'est ce qu'on peut vérifier directement de la manière suivante.

Considérons la fonction $\Gamma(y)$ ou simplement $\Gamma$, dont la valeur développée serait
\[ \Gamma=y-\tfrac{1}{3}P'y^{3}-\tfrac{1}{3\cdot5}P''y^{5}-\tfrac{1}{5\cdot7}P'''y^{7}-\text{etc.}; \]
on aura par des différentiations successives,
\[ \begin{aligned} \frac{d\Gamma}{dy}&=1-P'y^{2}-\tfrac{1}{3}P''y^{4}-\tfrac{1}{5}P'''y^{6}-\text{etc.}\\ \frac{1}{y^{2}}\cdot\frac{d\Gamma}{dy}-\frac{1}{y}\cdot\frac{dd\Gamma}{dy^{2}}&=\frac{1}{y^{2}}+P'+P''y^{2}+P'''y^{4}+\text{etc.} \end{aligned} \]
Mais si on remonte à l'origine des quantités $P'$, $P''$, $P'''$, etc., on trouvera que la suite $\frac{\pi}{2}(1+P'y^{2}+P''y^{4}+P'''y^{6}+\text{etc.})$ n'est autre chose que l'intégrale prise, depuis $\varphi=0$ jusqu'à $\varphi=\frac{1}{2}\pi$, de la différentielle $d\varphi(1+py^{2}+p^{2}y^{4}+p^{3}y^{6}+\text{etc.})=\frac{d\varphi}{1-py^{2}}=\frac{d\varphi}{1-(\delta\sin^{2}\varphi+\varepsilon\cos^{2}\varphi)y^{2}}$. Cette intégrale a pour valeur
\[ \frac{\tfrac{1}{2}\pi}{\sqrt{[(1-\delta y^{2})(1-\varepsilon y^{2})]}}. \]
Donc si on fait pour abréger, $R=\sqrt{[1-(\delta+\varepsilon)y^{2}+\delta\varepsilon y^{4}]}$, on aura
\[ \frac{1}{y^{2}}\cdot\frac{d\Gamma}{dy}-\frac{1}{y}\cdot\frac{dd\Gamma}{dy^{2}}=\frac{1}{y^{2}R}; \]
d'où résulte en multipliant par $dy$ et intégrant,
\[ \frac{1}{y}\cdot\frac{d\Gamma}{dy}=\int-\frac{dy}{y^{2}R}=\frac{R}{y}-\delta\varepsilon\int\frac{y^{2}\,dy}{R}. \]

\origpage{193}
Multipliant par $y\,dy$ et intégrant de nouveau, on aura
\[ \Gamma=\int R\,dy-\delta\varepsilon\cdot\frac{y^{2}}{2}\int\frac{y^{2}\,dy}{R}+\frac{\delta\varepsilon}{2}\int\frac{y^{4}\,dy}{R}, \]
ou, ce qui revient au même,
\[ \Gamma=\frac{\delta\varepsilon}{2}(1-y^{2})\int\frac{y^{2}\,dy}{R}+\tfrac{1}{2}yR+\int\frac{dy}{2R}(1-\delta\varepsilon y^{2}). \]

Soit maintenant $y\sqrt{\delta}=\sin\psi$ et $k^{2}=\frac{\varepsilon}{\delta}=\frac{a^{2}}{b^{2}}\cdot\frac{b^{2}-c^{2}}{a^{2}-c^{2}}$, on aura $\frac{dy}{R}=\frac{1}{\sqrt{\delta}}\cdot\frac{d\psi}{\sqrt{(1-k^{2}\sin^{2}\psi)}}$, et
\[ \int\frac{dy}{R}(1-\delta\varepsilon y^{2})=\frac{1-\delta}{2\sqrt{\delta}}F(k,\psi)+\frac{\sqrt{\delta}}{2}E(k,\psi). \]
Il reste à faire dans ces expressions $y=1$, ce qui donnera $\sin\psi=\sqrt{\delta}$, $\cos\psi=\sqrt{(1-\delta)}=\frac{c}{a}$, $R=\sqrt{(1-\delta)}\cdot\sqrt{(1-\varepsilon)}=\frac{c^{2}}{ab}$; donc la valeur cherchée
\[ \Gamma=\tfrac{1}{2}\cdot\frac{c^{2}}{ab}+\frac{\cos^{2}\psi}{2\sin\psi}F(k,\psi)+\frac{\sin^{2}\psi}{2\sin\psi}E(k,\psi); \]
d'où résulte enfin l'aire de l'ellipsoïde
\[ 8S=2\pi c^{2}+\frac{2\pi ab}{\sin\psi}[\cos^{2}\psi F(k,\psi)+\sin^{2}\psi E(k,\psi)]; \]
et on voit que cette formule s'accorde entièrement avec celle que nous avons trouvée par la seconde méthode, puisque les quantités $k$ et $\psi$ ne diffèrent pas de celles qui ont été désignées ci-dessus par $b'$ et $\nu$.

Nous remarquerons enfin qu'on peut déterminer généralement l'aire d'un quadrilatère quelconque compris entre deux lignes de courbure d'une espèce, et deux lignes de courbure de l'autre espèce. Il suffit pour cela de substituer dans la formule
\[ S=\frac{B}{A}\cdot PM-ABQN, \]
les valeurs des intégrales $M$, $N$, $P$, $Q$, prises entre les limites données. Il faudra donc prendre les deux intégrales $M$, $N$ entre les deux valeurs de $\psi$ qui répondent aux côtés du quadrilatère dont les projections sont des hyperboles, et les deux intégrales $P$ et $Q$ entre

\origpage{194}
les deux valeurs de $x$ qui répondent aux côtés du quadrilatère dont les projections sont des ellipses. Ainsi tous les quadrilatères dont il s'agit peuvent être déterminés par des fonctions elliptiques, et les résultats deviendront plus simples si on les applique à la zône entière comprise entre deux lignes de courbure d'une même espèce.

La figure 14 pourra servir à faire mieux entendre les constructions que ce chapitre suppose.

\rmfigure{figures/fig-14.png}{Fig.~14.}{A quadrant with centre C and axes CA and CB, showing the ellipse BKA, the ellipse GNO, the hyperbola from O through M and K to I, the ellipses from b through n and m to a and from b' through n' and m' to a', the curves rising from e' through n, n' and N and from e through m and m', and the dotted lines BI, IA, fN, Ne, bM and Ma.}

$AKB$ est l'ellipse principale dont les demi-axes sont $CA=a$, $CB=b$.

L'ellipse et l'hyperbole auxiliaires sont représentées par $ONG$ et $OMI$; elles ont pour axes communs $CO=A$, $CG=B$.

Ayant pris sur l'hyperbole auxiliaire $OMI$ un point quelconque $M$ dont les coordonnées sont $Ca=\alpha$, $aM=\beta$, on décrit avec les demi-axes $Ca=\alpha$, $Cb=\beta$, l'ellipse $am\beta$\ednote{Printed $am\beta$ with the Greek letter; the ellipse through the points $a$, $m$, $b$ is presumably meant, $amb$, an apparent misprint, to be checked against the plate. Likewise $a'm'\beta'$ below.} qui sera la projection d'une des lignes de courbure de la première espèce

L'ellipse $a'm'\beta'$ représente la projection d'une autre ligne de courbure de la même espèce.

Ayant pris sur l'ellipse auxiliaire $ONG$ un point quelconque $N$ dont les coordonnées sont $Ce=\alpha'$, $eN=\beta'$, on décrit avec les demi-axes $Ce=\alpha'$, $Cf=\beta'$, une hyperbole $emm'$ qui sera la projection d'une des lignes de courbure de la seconde espèce.

L'hyperbole $e'nn'$ représente la projection d'une autre ligne de courbure de la seconde espèce.

Nous avons fait voir qu'on peut déterminer en général par les fonctions elliptiques, l'aire de tout quadrilatère formé sur la surface de l'ellipsoïde, qui a pour projection $mnm'n'$.

\section*{De quelques formules générales qui peuvent se ramener aux fonctions elliptiques.}

\subsection*{(137).}

La nature des fonctions elliptiques étant connue et approfondie, il importe de ramener à ces quantités le plus grand nombre de formules intégrales qu'il sera possible. Comme la multitude en est infinie et aussi variée que les substitutions qu'on peut mettre en usage, nous nous contenterons d'indiquer quelques cas où cette réduction a lieu.

I. On peut réduire aux fonctions elliptiques l'intégrale\ldots\ $\displaystyle\int\frac{P\,dx}{\sqrt{(\alpha+\beta x^{2}+\gamma x^{4}+\delta x^{6})}}$, dans laquelle $P$ est une fonction rationnelle de $x$.

\origpage{195}
Car si on fait $x^{2}=z$, on pourra supposer $P=M+N\sqrt{z}$, $M$ et $N$ étant des fonctions rationnelles de $z$, et l'intégrale proposée deviendra
\[ \int\frac{\tfrac{1}{2}M\,dz}{\sqrt{(\alpha z+\beta z^{2}+\gamma z^{3}+\delta z^{4})}}+\int\frac{\tfrac{1}{2}N\,dz}{\sqrt{(\alpha+\beta z+\gamma z^{2}+\delta z^{3})}}, \]
dont les deux parties sont comprises dans les fonctions elliptiques.

II. Toute formule $\displaystyle\int\frac{P\,dx}{\sqrt[4]{(\alpha+\beta x^{2}+\gamma x^{4})}}$, dans laquelle $P$ est une fonction rationnelle de $x$, peut se ramener aux fonctions elliptiques.

Car on peut toujours faire $P=M+Nx$, $M$ et $N$ étant des fonctions rationnelles paires de $x$. Considérons la partie
\[ \frac{Nx\,dx}{\sqrt[4]{(\alpha+\beta x^{2}+\gamma x^{4})}}. \]
Si on fait $\sqrt[4]{(\alpha+\beta x^{2}+\gamma x^{4})}=z$, ce qui donne
\[ x^{2}=-\frac{\beta+\sqrt{(\beta^{2}-4\alpha\gamma+4\gamma z^{4})}}{2\gamma}, \]
il est clair que par la substitution de la valeur de $x^{2}$, $Nx\,dx$ ne contiendra d'autre radical que $\sqrt{(\beta^{2}-4\alpha\gamma+4\gamma z^{4})}$; donc toute la difficulté se réduit à intégrer une quantité de la forme
\[ \frac{Qz^{3}\,dz}{\sqrt{(\beta^{2}-4\alpha\gamma+4\gamma z^{4})}}, \]
dans laquelle $Q$ est une fonction rationnelle de $z^{4}$.

Quant à la partie $\frac{M\,dx}{\sqrt[4]{(\alpha+\beta x^{2}+\gamma x^{4})}}$, si on fait $\sqrt[4]{(\alpha+\beta x^{2}+\gamma x^{4})}=xy$, on aura
\[ \begin{aligned} x^{2}&=\frac{-\beta+\sqrt{(\beta^{2}-4\alpha\gamma+4\alpha y^{4})}}{2(\gamma-y^{4})}\\ \frac{dx}{xy}&=\frac{2x^{2}y^{2}\,dy}{\sqrt{(\beta^{2}-4\alpha\gamma+4\alpha y^{4})}}. \end{aligned} \]
D'où l'on voit que la transformée en $y$ contiendra une partie rationnelle et une de la forme
\[ \frac{Ry^{2}\,dy}{\sqrt{(\beta^{2}-4\alpha\gamma+4\alpha y^{4})}}, \]

\origpage{196}
$R$ étant une fonction rationnelle de $y^{4}$; donc la formule proposée est toujours réductible aux fonctions elliptiques.

III. On résoudrait absolument de la même manière la formule $\displaystyle\int P\,dx\sqrt[4]{(\alpha+\beta x^{2}+\gamma x^{4})}$, $P$ étant une fonction rationnelle de $x$.

Ces deux cas comprennent la formule $\displaystyle\int P\,dx(\alpha+\beta x^{2}+\gamma x^{4})^{\pm\frac{3}{4}}$, et aussi la formule $\displaystyle\int Q\,dy(\alpha+\beta y+\gamma y^{2})^{\pm\frac{1}{4}}$, $Q$ étant une fonction rationnelle de $y$; car en faisant $y=x^{2}$, cette dernière devient un cas particulier des deux autres.

IV. On peut réduire aux fonctions elliptiques la formule $\displaystyle\int P\,dx(\alpha+\beta x+\gamma x^{2}+\delta x^{3})^{\pm\frac{1}{3}}$, $P$ étant une fonction rationnelle de $x$.

Cette réduction peut se faire de plusieurs manières; et d'abord si on fait l'une ou l'autre des suppositions
\[ \begin{aligned} \sqrt[3]{(\alpha+\beta x+\gamma x^{2}+\delta x^{3})}&=\sqrt[3]{\alpha}+z\\ \sqrt[3]{(\alpha+\beta x+\gamma x^{2}+\delta x^{3})}&=x\sqrt[3]{\delta}+z, \end{aligned} \]
\ednote{The first supposition is printed $\sqrt[3]{\alpha}+z$; the ``$\cdot x$'' after it in the scan is a reader's ink addition. The sense requires $\sqrt[3]{\alpha}+zx$.}
la transformée en $z$ sera comprise dans les fonctions elliptiques. On peut aussi faire disparaître à volonté le coefficient $\alpha$ ou le coefficient $\delta$ sous le radical; il faut faire pour cela $x=m+y$, ou $x=m+\frac{1}{y}$, et prendre pour $m$ une racine réelle de l'équation $\alpha+\beta m+\gamma m^{2}+\delta m^{3}=0$. Supposons qu'on ait fait disparaître de cette manière $\delta$, alors on fera $\alpha+\beta x+\gamma x^{2}=z^{3}$, ce qui donne
\[ x=-\frac{\beta+\sqrt{(\beta^{2}-4\gamma+4\gamma z^{3})}}{2\gamma}; \]
\ednote{The radicand is printed $\beta^{2}-4\gamma+4\gamma z^{3}$, with a blank space after the 4; $\beta^{2}-4\alpha\gamma+4\gamma z^{3}$ is expected, as in the formula that follows, an apparent dropped letter.}
et en substituant cette valeur dans la formule proposée, toute la difficulté se réduira à intégrer une différentielle de la forme $\frac{Q\,dz}{\sqrt{(\beta^{2}-4\alpha\gamma+4\gamma z^{3})}}$, $Q$ étant une fonction rationnelle de $z$.

V. On peut réduire aux fonctions elliptiques la formule\ldots\ $\displaystyle\int\frac{P\,dx}{\sqrt{(\alpha+\beta x+\gamma x^{2}+\delta x^{3}+\gamma x^{4}+\beta x^{5}+\alpha x^{6})}}$, $P$ étant une fonction rationnelle de $x$.

\origpage{197}

Car si on fait $x^{2}+1=xz$, cette formule deviendra d'abord
\[ \int\frac{Px^{-\frac{3}{2}}\,dx}{\sqrt{[\alpha(z^{3}-3z)+\beta(z^{2}-2)+\gamma z+\delta]}}; \]
ensuite comme on a $x=\frac{1}{2}z\pm\frac{1}{2}\sqrt{(z^{2}-4)}$, cette valeur étant substituée dans $P$, le résultat sera de la forme $M\pm N\sqrt{(z^{2}-4)}$, $M$ et $N$ étant des fonctions rationnelles de $z$. De plus on a
\[ \begin{aligned} x+1&=x^{\frac{1}{2}}\sqrt{(z+2)}\\ x-1&=\pm x^{\frac{1}{2}}\sqrt{(z-2)}\\ 2x^{-\frac{1}{2}}&=\sqrt{(z+2)}\mp\sqrt{(z-2)}\\ x^{-\frac{3}{2}}\,dx&=\frac{\tfrac{1}{2}dz}{\sqrt{(z+2)}}\mp\frac{\tfrac{1}{2}dz}{\sqrt{(z-2)}}. \end{aligned} \]
Donc toutes les substitutions étant faites, la formule proposée sera composée de deux parties, l'une de la forme $\displaystyle\int\frac{Z\,dz}{\sqrt{[Q(z+2)]}}$, $Z$ étant une fonction rationnelle de $z$, et $Q$ désignant le polynome $\alpha(z^{3}-3z)+\beta(z^{2}-2)+\gamma z+\delta$, l'autre étant de la forme $\displaystyle\int\frac{Z'\,dz}{\sqrt{[Q(z-2)]}}$, $Z'$ étant aussi une fonction rationnelle de $z$. Donc la transformée en $z$ sera intégrable par les fonctions elliptiques.

VI. On pourra donc intégrer de même la formule\ldots\ $\displaystyle\int\frac{P\,dy}{\sqrt{(\beta+\gamma y^{2}+\delta y^{4}+\gamma y^{6}+\beta y^{8})}}$, $P$ étant une fonction rationnelle de $y$.

Car si on fait $P=M+Ny$, $M$ et $N$ étant des fonctions paires de $y$, et qu'on appelle $R$ le radical, la partie $\frac{M\,dy}{R}$ rentrera dans le cas précédent, en faisant $\alpha=0$ et $y^{2}=x$. L'autre partie $\frac{Ny\,dy}{R}$ se réduit pareillement à une fonction elliptique en faisant $y^{2}=x$.

VII. On peut ramener aux fonctions elliptiques la formule $\displaystyle\int\frac{P\,dx}{\sqrt{(\alpha+\beta x^{4}+\gamma x^{8})}}$, $P$ étant une fonction rationnelle de $x$, pourvu que $\alpha$ et $\gamma$ soient de même signe.

Car en faisant $\gamma=\alpha m^{8}$ et $mx=y$, cette formule se trouvera comprise dans le cas précédent.

\origpage{198}

VIII. Soit proposé la formule $Z=\displaystyle\int\frac{Q\,d\varphi}{\sqrt{(1-m\sin^{2}\varphi)}}$, $Q$ étant une fonction rationnelle paire de $\sin\varphi$, et $m$ étant plus grand que l'unité.

Pour ramener cette formule aux fonctions elliptiques, soit d'abord $m=\frac{1}{\sin^{2}\alpha}$, on voit que $\alpha$ doit être la plus grande valeur de $\varphi$. Soit donc $\sin\varphi=\sin\alpha\sin\psi$, et $c=\sin\alpha$, la transformée sera $Z=\displaystyle\int\frac{cQ\,d\psi}{\sqrt{(1-c^{2}\sin^{2}\psi)}}$, et $Q$ deviendra par la substitution, une fonction rationnelle paire de $\sin\psi$, de sorte que la formule proposée sera ramenée aux fonctions elliptiques.

IX. Si on avait la formule $Z=\displaystyle\int\frac{Q\,d\varphi}{\sqrt{(1+m\sin^{2}\varphi)}}$, $m$ étant positif et d'une grandeur quelconque, il suffirait pour ramener cette quantité aux fonctions elliptiques, de faire $\varphi=\frac{1}{2}\pi-\psi$, et $\frac{m}{1+m}=c^{2}$.

X. Soit proposé la formule $Z=\displaystyle\int\frac{Q\,d\varphi}{\sqrt{(f+g\cos\varphi+h\sin\varphi)}}$, dans laquelle $f$, $g$, $h$ sont des constantes, et $Q$ une fonction rationnelle de $\sin\varphi$ et $\cos\varphi$.

On fera d'abord $\varphi=2\psi+m$, et on prendra l'indéterminée $m$ de manière que $\operatorname{tang}m=\frac{h}{g}$, alors la quantité $f+g\cos\varphi+h\sin\varphi$ deviendra de la forme $f'\pm g'\sin^{2}\psi$, et on aura $Q\,d\varphi=(M+N\sin\psi\cos\psi)\,d\psi$, $M$ et $N$ étant des fonctions paires de $\sin\psi$. La partie $\frac{N\,d\psi\sin\psi\cos\psi}{\sqrt{(f'\pm g'\sin^{2}\psi)}}$ peut être rendue rationnelle en faisant $f'\pm g'\sin^{2}\psi=x^{2}$, ainsi il ne restera à intégrer que la différentielle $\frac{M\,d\psi}{\sqrt{(f'\pm g'\sin^{2}\psi)}}$, et il est évident que cette intégrale ne dépend que des fonctions elliptiques.

XI. Si l'on a plus généralement
\[ Z=\int\frac{Q\,d\varphi}{\sqrt{(\alpha+\beta\cos\varphi+\gamma\sin\varphi+\delta\cos^{2}\varphi+\varepsilon\sin\varphi\cos\varphi+\zeta\sin^{2}\varphi)}}, \]
$Q$ étant une fonction rationnelle de $\sin\varphi$ et $\cos\varphi$; on fera $\operatorname{tang}\frac{1}{2}\varphi=z$, ce qui donne $\sin\varphi=\frac{2z}{1+z^{2}}$, $\cos\varphi=\frac{1-z^{2}}{1+z^{2}}$, $d\varphi=\frac{2\,dz}{1+z^{2}}$, et la transformée en $z$ ne contiendra qu'un radical sous lequel la variable ne

\origpage{199}
passe pas le quatrième degré; elle sera donc généralement réductible aux fonctions elliptiques.

XII. Soit enfin $Z=\displaystyle\int\frac{Q\,d\varphi}{\sqrt{[(\alpha^{2}+\beta^{2}\sin^{2}\varphi)(\gamma+\delta\sin^{2}\varphi)]}}$, $Q$ étant une fonction rationnelle paire de $\sin\varphi$; si on fait $\operatorname{tang}\varphi=m\operatorname{tang}\psi$, ce qui donne $\sin^{2}\varphi=\frac{m^{2}\sin^{2}\psi}{m^{2}\sin^{2}\psi+\cos^{2}\psi}$, $\cos^{2}\varphi=\frac{\cos^{2}\psi}{m^{2}\sin^{2}\psi+\cos^{2}\psi}$, $d\varphi=\frac{m\,d\psi}{m^{2}\sin^{2}\psi+\cos^{2}\psi}$, la transformée sera
\[ Z=\int\frac{Q\,d\psi}{\sqrt{(\alpha^{2}m^{2}\sin^{2}\psi+\alpha^{2}\cos^{2}\psi+\beta^{2}m^{2}\sin^{2}\psi)}\cdot\sqrt{(\gamma m^{2}\sin^{2}\psi+\gamma\cos^{2}\psi+\delta m^{2}\sin^{2}\psi)}}. \]
Prenons l'indéterminée $m$, de manière qu'on ait $\alpha^{2}m^{2}+\beta^{2}m^{2}=\alpha^{2}$, ou $m^{2}=\frac{\alpha^{2}}{\alpha^{2}+\beta^{2}}$, l'intégrale précédente se réduira à la forme $\displaystyle\int\frac{mQ\,d\psi}{\sqrt{(A+B\sin^{2}\psi)}}$; et comme par la même substitution $Q$ devient une fonction rationnelle paire de $\sin\psi$, la formule $Z$ se ramènera encore aux fonctions elliptiques.

\subsection*{(138).}

Nous ne multiplierons pas davantage ces exemples de réduction; ils suffisent pour indiquer les substitutions à faire dans les cas analogues à ceux que nous avons développés. Nous terminerons en réunissant dans un même tableau quelques-unes des formules les plus simples dont l'application se présente fréquemment dans les intégrales qui se rapportent aux fonctions elliptiques.
\[ \left.\begin{aligned} \int\frac{d\varphi}{\Delta}&=F(c,\varphi)\\ \int\Delta\,d\varphi&=E(c,\varphi) \end{aligned}\right\}\ \text{Formules principales.} \]
\[ \begin{aligned} \int\frac{d\varphi}{\Delta^{3}}&=\frac{1}{b^{2}}E(c,\varphi)-\frac{c^{2}\sin\varphi\cos\varphi}{b^{2}\Delta}\\ \int\frac{d\varphi\sin^{2}\varphi}{\Delta}&=\frac{1}{c^{2}}(F-E)\\ \int\frac{d\varphi\cos^{2}\varphi}{\Delta}&=\frac{1}{c^{2}}(E-b^{2}F)\\ \int\frac{d\varphi}{\Delta\cos^{2}\varphi}&=\frac{1}{b^{2}}(\Delta\operatorname{tang}\varphi+b^{2}F-E)\\ \int\frac{d\varphi\operatorname{tang}^{2}\varphi}{\Delta}&=\frac{1}{b^{2}}(\Delta\operatorname{tang}\varphi-E) \end{aligned} \]

\origpage{200}
\[ \begin{aligned} \int\frac{d\varphi\operatorname{tang}^{2}\frac{1}{2}\varphi}{\Delta}&=2\Delta\operatorname{tang}\tfrac{1}{2}\varphi+F-2E\\ \int\frac{d\varphi}{\Delta\cos^{2}\frac{1}{2}\varphi}&=2\Delta\operatorname{tang}\tfrac{1}{2}\varphi+2F-2E\\ \int\frac{\Delta\,d\varphi}{\cos^{2}\varphi}&=\Delta\operatorname{tang}\varphi+F-E\\ \int\Delta\,d\varphi\operatorname{tang}^{2}\varphi&=\Delta\operatorname{tang}\varphi+F-2E\\ \int\Delta^{3}\,d\varphi&=\frac{c^{2}}{3}\Delta\sin\varphi\cos\varphi+\frac{4-2c^{2}}{3}E-\frac{b^{2}}{3}F\\ \int\Delta\,d\varphi\sin^{2}\varphi&=-\tfrac{1}{3}\Delta\sin\varphi\cos\varphi+\frac{2c^{2}-1}{3c^{2}}E+\frac{b^{2}}{3c^{2}}F\\ \int\Delta\,d\varphi\cos^{2}\varphi&=\tfrac{1}{3}\Delta\sin\varphi\cos\varphi+\frac{1+c^{2}}{3c^{2}}E-\frac{b^{2}}{3c^{2}}F. \end{aligned} \]
Toutes ces intégrales sont prises à compter de $\varphi=0$.

\section*{De l'intégrale $V=\displaystyle\int\frac{(x^{2}+m)\,dx}{\sqrt{(\alpha+\beta x+\gamma x^{2}+\varepsilon x^{4})}}$.}

\subsection*{(139).}

Cette intégrale où le polynome sous le radical manque de second terme, est généralement réductible aux fonctions elliptiques de la première et de la seconde espèce.

En effet, si on décompose le polynome $\alpha+\beta x+\gamma x^{2}+\varepsilon x^{4}$ en facteurs réels du second degré, de cette manière:
\[ \alpha+\beta x+\gamma x^{2}+\varepsilon x^{4}=\varepsilon(x^{2}+\lambda x+\mu)(x^{2}-\lambda x+\nu), \]
et qu'ensuite pour faire disparaître les puissances impaires de la variable sous le radical, on suppose $x=\frac{p+qy}{1+y}$, on aura, pour déterminer $p$ et $q$, les équations
\[ \begin{aligned} 2pq+\lambda(p+q)+2\mu&=0\\ 2pq-\lambda(p+q)+2\nu&=0, \end{aligned} \]
lesquelles donneront toujours pour $p$ et $q$ des valeurs réelles, en prenant convenablement les valeurs de $\lambda$, $\mu$, $\nu$ (art.~5).

Cela posé, si l'on fait pour abréger,
\[ \begin{aligned} M&=p^{2}-\lambda p+\nu=(p-\tfrac{1}{2}\lambda)(p-q)\\ N&=q^{2}-\lambda q+\nu=-(q-\tfrac{1}{2}\lambda)(p-q)\\ P&=p^{2}+\lambda p+\mu=(p+\tfrac{1}{2}\lambda)(p-q)\\ Q&=q^{2}+\lambda q+\mu=-(q+\tfrac{1}{2}\lambda)(p-q), \end{aligned} \]

\origpage{201}
on aura
\[ \begin{aligned} x^{2}+\lambda x+\mu&=\frac{P+Qy^{2}}{(1+y)^{2}}\\ x^{2}-\lambda x+\nu&=\frac{M+Ny^{2}}{(1+y)^{2}}, \end{aligned} \]
et la transformée en $y$ sera
\[ V=(q-p)\int\frac{dy[(p+qy)^{2}+m(1+y)^{2}]}{(1+y)^{2}\sqrt{[\varepsilon(P+Qy^{2})(M+Ny^{2})]}}. \]
Soit pour abréger, $Y=\sqrt{[\varepsilon(P+Qy^{2})(M+Ny^{2})]}$, on aura par le développement de la quantité précédente,
\[ V=(m+q^{2})(q-p)\int\frac{dy}{Y}-2q(q-p)^{2}\int\frac{dy}{(1+y)Y}+(q-p)^{3}\int\frac{dy}{(1+y)^{2}Y}. \]
Mais par la différentielle de $\frac{Y}{1+y}$, on trouve
\[ d\left(\frac{Y}{1+y}\right)=\frac{\varepsilon\,dy}{Y}\left[QN(y^{2}-1)+\frac{MQ+NP+2QN}{1+y}-\frac{(M+N)(P+Q)}{(1+y)^{2}}\right]. \]
Substituant dans cette quantité les valeurs de $M$, $N$, $P$, $Q$, on aura
\[ d\left(\frac{Y}{1+y}\right)=(q^{2}-\tfrac{1}{4}\lambda^{2})(q-p)^{2}\cdot\frac{(y^{2}-1)\varepsilon\,dy}{Y}+2q(q-p)^{3}\cdot\frac{\varepsilon\,dy}{(1+y)Y}-(q-p)^{4}\cdot\frac{\varepsilon\,dy}{(1+y)^{2}Y}. \]
De là on tire
\[ V=-\frac{1}{\varepsilon(q-p)}\cdot\frac{Y}{1+y}+(q^{2}-\tfrac{1}{4}\lambda^{2})(q-p)\int\frac{y^{2}\,dy}{Y}+(m+\tfrac{1}{4}\lambda^{2})(q-p)\int\frac{dy}{Y}. \]
Or les intégrales $\displaystyle\int\frac{dy}{Y}$, $\displaystyle\int\frac{y^{2}\,dy}{Y}$ se ramènent par les transformations connues, aux fonctions elliptiques de la première et de la seconde espèce; donc l'intégrale $V$ ne dépend généralement que de ces deux fonctions, et par conséquent peut être déterminée par des arcs d'ellipse.

\section*{De l'intégrale $Z=\displaystyle\int\frac{dz}{(1+\mu z^{2})\sqrt[3]{(1+\nu z^{2})}}$.}

\subsection*{(140).}

Cette intégrale se ramène aisément aux fonctions elliptiques par les méthodes données dans l'article 137; mais le calcul mérite d'être développé dans quelques cas particuliers, à cause des

\origpage{202}
réductions qu'ils présentent et qui peuvent jeter un nouveau jour sur l'usage des fonctions elliptiques.

Supposons d'abord $\nu$ positif, $\mu$ étant à volonté positif ou négatif, si on fait $\sqrt[3]{(1+\nu z^{2})}=y$, et $\alpha^{3}=\frac{\nu}{\mu}-1$, on aura pour première transformée
\[ Z=\frac{3\sqrt{\nu}}{2\mu}\int\frac{y\,dy}{(\alpha^{3}+y^{3})\sqrt{(y^{3}-1)}}. \]
Cette intégrale peut se décomposer en deux autres; soit pour cet effet,
\[ \begin{aligned} P&=\int\frac{y\,dy}{(\alpha^{2}-\alpha y+y^{2})\sqrt{(y^{3}-1)}}\\ Q&=\int\frac{y^{2}\,dy}{(\alpha^{3}+y^{3})\sqrt{(y^{3}-1)}}, \end{aligned} \]
et on aura
\[ Z=\frac{3\sqrt{\nu}}{2\alpha\mu}(P-Q). \]
L'intégrale $Q$ se trouve immédiatement; car si on met au lieu de $y$ sa valeur en fonction de $z$, on aura
\[ Q=\frac{2\mu}{2\sqrt{\nu}}\int\frac{dz}{1+\mu z^{2}}=\frac{2\sqrt{\mu}}{3\sqrt{\nu}}\operatorname{arc}\operatorname{tang}(z\sqrt{\mu}). \]
\ednote{The first coefficient is printed $\frac{2\mu}{2\sqrt{\nu}}$; the substitution gives $\frac{2\mu}{3\sqrt{\nu}}$, which also agrees with the right-hand side, an apparent misprint.}
Ainsi la quantité $Q$ est donnée par un arc de cercle si $\mu$ est positif, et par un logarithme si $\mu$ est négatif. Tout se réduit donc à trouver la valeur de l'autre intégrale $P$.

Pour cela, soit $y=1+x^{2}$, on aura
\[ P=\int\frac{2(1+x^{2})\,dx}{(x^{4}+(2-\alpha)x^{2}+\alpha^{2}-\alpha+1)\sqrt{(x^{4}+3x^{2}+3)}}. \]
Cette intégrale dépend en général des fonctions elliptiques dont le paramètre est imaginaire, et nous avons donné pour cet objet les formules nécessaires. Mais dans le cas de $\alpha=2$, l'intégrale $P$ ne dépend que d'une fonction de troisième espèce dont le paramètre est réel, et c'est ce cas dont nous allons développer le calcul.

Ayant donc fait $\alpha=2$, ou $\nu=9\mu$, il s'agira de trouver l'intégrale
\[ P=\int\frac{2(1+x^{2})\,dx}{(x^{4}+3)\sqrt{(x^{4}+3x^{2}+3)}}. \]

\origpage{203}
Pour cela, soit $r=\sqrt[4]{3}$ et $x=r\operatorname{tang}\frac{1}{2}\varphi$; soit de plus $c=\frac{1}{2}\sqrt{(2-\sqrt{3})}=\sin 15^{\circ}$, on aura la transformée
\[ P=\int\frac{d\varphi}{\Delta}\left(\frac{1-r^{2}}{6r}+\frac{1}{6r}\cdot\frac{r^{2}+\cos\varphi}{1-\tfrac{1}{2}\sin^{2}\varphi}\right), \]
d'où l'on tire
\[ P=\frac{1-r^{2}}{6r}F+\frac{r}{6}\Pi(-\tfrac{1}{2})+\frac{1}{6r}\int\frac{d\varphi\cos\varphi}{(1-\tfrac{1}{2}\sin^{2}\varphi)\Delta}. \]
Pour avoir cette dernière intégrale, soit $\sqrt{(1-c^{2}\sin^{2}\varphi)}=u\sin\varphi$, la différentielle $\frac{d\varphi\cos\varphi}{(1-\frac{1}{2}\sin^{2}\varphi)\Delta}$ deviendra $-\frac{du}{u^{2}-\frac{1}{4}r^{2}}$, et son intégrale sera $\frac{1}{r}\log\left(\frac{u+\frac{1}{2}r}{u-\frac{1}{2}r}\right)$; on a donc
\[ P=\frac{1-r^{2}}{6r}F+\frac{r}{6}\Pi(-\tfrac{1}{2})+\frac{1}{6r^{2}}\log\left(\frac{\Delta+\frac{1}{2}r\sin\varphi}{\Delta-\frac{1}{2}r\sin\varphi}\right). \]
On n'ajoute point de constante, parce que cette intégrale s'évanouit lorsque $\varphi=0$, valeur qui répond à celle de $z=0$.

La valeur de $P$ dépend donc de la fonction elliptique de troisième espèce $\Pi(-\frac{1}{2})$, et il faut examiner si, à raison du paramètre $-\frac{1}{2}$, cette fonction ne pourrait pas se réduire indéfiniment à une espèce inférieure.

Ayant le module $c=\frac{1}{2}\sqrt{(2-\sqrt{3})}$, et son complément $b=\frac{1}{2}\sqrt{(2+\sqrt{3})}$, on voit immédiatement que le paramètre $n=-\frac{1}{2}$ se rapporte à la seconde forme du n° 50, et qu'en faisant $n=-1+b^{2}\sin^{2}\theta$, on aura $\sin^{2}\theta=\frac{n+1}{b^{2}}=4-2\sqrt{3}$, ou $\sin\theta=\sqrt{3}-1$. Or cette valeur de $\sin\theta$ est celle qui, suivant l'article 24, satisfait à l'équation $F(b,\theta)=\frac{1}{3}F^{1}(b)$. Donc la réduction de la fonction $\Pi(-\frac{1}{2})$ est possible, et elle s'effectuera par la formule du n° 102 qui donne
\[ \Pi(n)=\left(1-\frac{V\Delta(b,\theta)}{b^{2}\sin\theta\cos\theta}\right)F(c,\varphi)+\frac{W\Delta(b,\theta)\Delta(c,\varphi)}{b^{2}\sin\theta\cos\theta\sin\varphi\cos\varphi}. \]
Il s'agit donc de faire les substitutions dans cette formule; or la valeur connue de $\theta$ donne $\frac{\Delta(b,\theta)}{b^{2}\sin\theta\cos\theta}=\frac{2}{r}$, on a ensuite par les formules du n° 57, $V=\frac{1}{2}b^{2}\sin\theta\sin\theta_{2}(1+\sin\theta)$, $\cos\theta_{2}=1-\sin\theta$

\origpage{204}
$=2-\sqrt{3}$, $\sin\theta_{2}=r\sin\theta$; donc $V=\frac{1}{2r}$. On a en second lieu par les formules du même article,
\[ W=\frac{1}{3\sqrt{\alpha'}}\operatorname{arc}\operatorname{tang}\frac{n'\sqrt{\alpha'}\sin\theta\sin\theta_{2}}{1+n'}+\frac{1}{3\sqrt{\alpha'}}\operatorname{arc}\operatorname{tang}\frac{n'\sqrt{\alpha'}\sin^{2}\theta\sin\theta_{2}}{1+n'-n'\cos^{2}\theta\cos\theta_{2}}. \]
Substituant les valeurs de $\theta$ et $\theta_{2}$, faisant de plus $n'=b^{2}\operatorname{tang}^{2}\varphi$ et $\sqrt{\alpha'}=\frac{\Delta}{\sin\varphi\cos\varphi}$, on aura
\[ W=\frac{\sin\varphi\cos\varphi}{3\Delta}\operatorname{arc}\operatorname{tang}\frac{r\operatorname{tang}\varphi}{2\Delta}+\frac{\sin\varphi\cos\varphi}{3\Delta}\operatorname{arc}\operatorname{tang}\frac{r\Delta\operatorname{tang}\varphi}{1+r^{2}\Delta^{2}}. \]
Donc la fonction $\Pi(-\frac{1}{2})$ se réduit à la première espèce par la formule
\[ \Pi(-\tfrac{1}{2})=\left(1-\frac{1}{r}\right)F+\frac{2}{3r}\operatorname{arc}\operatorname{tang}\frac{r\operatorname{tang}\varphi}{2\Delta}+\frac{2}{3r}\operatorname{arc}\operatorname{tang}\frac{r\Delta\operatorname{tang}\varphi}{1+r^{2}\Delta^{2}}. \]
\ednote{Printed $\left(1-\frac{1}{r}\right)F$; the formula of n° 102 with $\frac{V\Delta(b,\theta)}{b^{2}\sin\theta\cos\theta}=\frac{1}{2r}\cdot\frac{2}{r}$ gives $1-\frac{1}{r^{2}}$, which is also needed for $F$ to cancel in $P$ below, an apparent misprint.}
Substituant cette valeur dans celle de l'intégrale $P$, on aura enfin
\[ P=\frac{1}{6r}\log\left(\frac{\Delta+\frac{1}{2}r\sin\varphi}{\Delta-\frac{1}{2}r\sin\varphi}\right)+\tfrac{1}{9}\operatorname{arc}\operatorname{tang}\frac{r\operatorname{tang}\varphi}{2\Delta}+\tfrac{1}{9}\operatorname{arc}\operatorname{tang}\frac{r\Delta\operatorname{tang}\varphi}{1+r^{2}\Delta^{2}}. \]
\ednote{The logarithmic term is printed with $\frac{1}{6r}$; the expression of $P$ on p.~203 has $\frac{1}{6r^{2}}$, an apparent misprint.}
D'où l'on voit que l'intégrale $P$ devient finalement indépendante des fonctions elliptiques, et que son expression ne renferme que des arcs de cercle et des logarithmes. Il en est par conséquent de même de l'intégrale proposée $Z$ dans le cas où l'on a $\alpha=2$, ou $\frac{\nu}{\mu}=9$.

Si on veut avoir la valeur de $Z$ lorsque $z=\infty$, il faudra faire $\varphi=\pi$, ce qui donne $P=\frac{2\pi}{9}$; d'ailleurs dans le même cas on a $Q=\frac{\pi}{9}$; donc $Z=\frac{\pi}{4\sqrt{\mu}}$.

\subsection*{(141).}

Les réductions nombreuses qui se sont offertes dans le calcul précédent, doivent faire présumer qu'il est possible de parvenir au même résultat sans le secours des fonctions elliptiques; c'est ce qu'on va mettre hors de doute de la manière suivante.

Remarquons d'abord que dans la formule $Z$, on peut changer à volonté le coefficient $\nu$ en un autre $\nu'$, pourvu qu'il soit de même

\origpage{205}
signe, et qu'il suffit pour cela de faire $z\sqrt{\nu}=z'\sqrt{\nu'}$. Soit donc $\nu=3$, ce qui donnera pour le cas que nous considérons $\mu=\frac{1}{3}$, et la formule proposée sera
\[ Z=\int\frac{3\,dz}{(3+z^{2})\sqrt[3]{(1+3z^{2})}}. \]
Je fais $z=\frac{1-x}{1+x}$, et j'ai la transformée
\[ \tfrac{1}{3}Z=-2^{\frac{5}{3}}\int\frac{(1-x^{2})\,dx}{(1-x^{3})\sqrt[3]{(1+x^{3})}}. \]
\ednote{The factor is printed $2^{\frac{5}{3}}$; the substitution gives $2^{-\frac{5}{3}}$, as in $\frac{1}{3}Z=2^{-\frac{5}{3}}(M-N)$ below, an apparent misprint.}
Je remarque ensuite qu'en faisant
\[ \begin{aligned} M&=\int\frac{-dx}{(1-x^{3})\sqrt[3]{(1+x^{3})}}\\ N&=\int\frac{-x^{2}\,dx}{(1-x^{3})\sqrt[3]{(1+x^{3})}}, \end{aligned} \]
on aura $\frac{1}{3}Z=2^{-\frac{5}{3}}(M-N)$. Il ne s'agit donc que d'avoir les intégrales $M$ et $N$. Pour cela,

Soit, 1°. $\sqrt[3]{(1+x^{3})}=\frac{x\sqrt[3]{2}}{t}$, on aura $M=2^{-\frac{1}{3}}\displaystyle\int\frac{-dt}{1-t^{3}}$.

Soit, 2°. $\sqrt[3]{(1+x^{3})}=\frac{\sqrt[3]{2}}{u}$, on aura $N=2^{-\frac{1}{3}}\displaystyle\int\frac{-du}{1-u^{3}}$.
Ces deux intégrales sont donc de la même forme, et peuvent s'exprimer par les arcs de cercle et les logarithmes.

S'il fallait les évaluer séparément, elles présenteraient une difficulté; car l'intégrale proposée devant être prise depuis $z=0$ qui donne $x=1$ et $t=1$, l'intégrale $M$ contiendrait une constante infinie. Il en serait de même de l'intégrale $N$; mais comme nous n'avons besoin que de la différence de ces intégrales, on peut éviter entièrement l'infini par un moyen fort simple.

On voit en effet par la nature de la question, qu'on aura la valeur de $Z$ par la seule formule
\[ Z=\frac{3}{4}\int\frac{-dv}{1-v^{3}}, \]
pourvu qu'on prenne l'intégrale depuis $v=t=\frac{x\sqrt[3]{2}}{\sqrt[3]{(1+x^{3})}}$ jusqu'à $v=u=\frac{\sqrt[3]{2}}{\sqrt[3]{(1+x^{3})}}$. L'intégration étant donc effectuée entre ces

\origpage{206}
limites toutes deux variables, on aura généralement
\[ \begin{aligned} Z={}&\tfrac{1}{8}\log\left(\frac{1+u+u^{2}}{1-2u+u^{2}}\right)+\frac{\sqrt{3}}{4}\operatorname{arc}\operatorname{tang}\frac{2u+1}{\sqrt{3}}\\ &-\tfrac{1}{8}\log\left(\frac{1+t+t^{2}}{1-2t+t^{2}}\right)-\frac{\sqrt{3}}{4}\operatorname{arc}\operatorname{tang}\frac{2t+1}{\sqrt{3}}. \end{aligned} \]
Si on veut avoir l'intégrale lorsque $z=\infty$, il faudra supposer $\sqrt[3]{(1+x^{3})}=\omega$, $\omega$ étant infiniment petit; on fera donc $t=-\frac{\sqrt[3]{2}}{\omega}$, et $u=\frac{\sqrt[3]{2}}{\omega}$, ce qui donnera $Z=\frac{\pi\sqrt{3}}{4}$; résultat conforme à celui que nous avons trouvé par l'autre méthode, puisqu'on a dans ce cas $\mu=\frac{1}{3}$.

\subsection*{(142).}

Revenons maintenant à la formule générale, et supposons que $\nu$ est négatif, ou, ce qui revient au même, considérons directement la formule
\[ Z=\int\frac{dz}{(1-\mu z^{2})\sqrt[3]{(1-\nu z^{2})}}, \]
dans laquelle nous supposerons $\nu$ positif, $\mu$ étant à volonté positif ou négatif. Si on fait $\sqrt[3]{(1-\nu z^{2})}=y$ et $\alpha^{3}=\frac{\nu}{\mu}-1$, on aura la transformée
\[ Z=\frac{3\sqrt{\nu}}{2\mu}\int\frac{-y\,dy}{(\alpha^{3}+y^{3})\sqrt{(1-y^{3})}}, \]
qui peut se mettre sous la forme $Z=\frac{3\sqrt{\nu}}{2\alpha\mu}(P-Q)$, en faisant
\[ \begin{aligned} P&=\int\frac{-y\,dy}{(\alpha^{2}-\alpha y+y^{2})\sqrt{(1-y^{3})}}\\ Q&=\int\frac{-y^{2}\,dy}{(\alpha^{3}+y^{3})\sqrt{(1-y^{3})}}. \end{aligned} \]
On trouve immédiatement la valeur de la seconde intégrale
\[ Q=\frac{2\mu}{3\sqrt{\nu}}\int\frac{dz}{1-\mu z^{2}}=\frac{\sqrt{\mu}}{3\sqrt{\nu}}\log\left(\frac{1+z\sqrt{\mu}}{1-z\sqrt{\mu}}\right); \]
pour avoir celle de la première, soit $y=1-x^{2}$, on aura la nouvelle

\origpage{207}
transformée
\[ P=\int\frac{2(1-x^{2})\,dx}{(x^{4}-(2-\alpha)x^{2}+\alpha^{2}-\alpha+1)\sqrt{(x^{4}-3x^{2}+3)}}. \]
Cette intégrale dépend en général des fonctions elliptiques dont le paramètre est imaginaire, et elle peut être ramenée à celles dont le paramètre est réel. Mais dans le cas de $\alpha=2$, la solution se simplifie beaucoup et devient tout à fait indépendante des fonctions elliptiques; c'est ce qu'on vérifiera aisément en traitant la formule
\[ P=\int\frac{2(1-x^{2})\,dx}{(x^{4}+3)\sqrt{(x^{4}-3x^{2}+3)}}, \]
comme celle de l'exemple précédent.

\subsection*{(143).}

On peut aussi, dans le même cas, parvenir directement au résultat de l'intégration sans le secours des fonctions elliptiques. Et d'abord comme la valeur de $\nu$ peut être prise à volonté, supposons $\nu=3$, ce qui donnera $\mu=\frac{1}{3}$, et la formule proposée deviendra
\[ Z=\int\frac{3\,dz}{(3-z^{2})\sqrt[3]{(1-3z^{2})}}. \]
Soit $\sqrt[3]{(1-3z^{2})}=1-\frac{x}{z}$, on aura la transformée
\[ Z=\frac{3}{2}\int\frac{(1-x^{2})\,dx}{(1+x^{2})(3x^{2}-1)}-\frac{\sqrt{3}}{2}\int\frac{3x^{4}+1}{3x^{4}+2x^{2}-1}\cdot\frac{dx}{\sqrt{(3x^{4}+6x^{2}-1)}}. \]
La première partie est rationnelle; la seconde semble devoir se décomposer en deux fonctions elliptiques de la troisième espèce, à cause des deux facteurs du dénominateur $3x^{4}+2x^{2}-1$. Mais si on examine les choses avec plus d'attention, on trouvera qu'en faisant $\sqrt{(3x^{4}+6x^{2}-1)}=px$, la seconde partie devient $\displaystyle-\frac{\sqrt{3}}{2}\int\frac{dp}{p^{2}-4}$, de sorte que la valeur totale de $Z$ est
\[ Z=\frac{3}{2}\int\frac{(1-x^{2})\,dx}{(1+x^{2})(3x^{2}-1)}-\frac{\sqrt{3}}{2}\int\frac{dp}{p^{2}-4}. \]
Ainsi cette intégrale se détermine entièrement par les arcs de cercle et les logarithmes.

\origpage{208}

Une pareille simplification aura lieu en général pour la formule
\[ \int\frac{(\gamma x^{4}-\alpha)\,dx}{(\gamma x^{4}+\varepsilon x^{2}+\alpha)\sqrt{(\alpha+\beta x^{2}+\gamma x^{4})}}; \]
car si on fait le radical $\sqrt{(\alpha+\beta x^{2}+\gamma x^{4})}=px$, cette formule deviendra $\displaystyle\int\frac{dp}{p^{2}-\beta+\varepsilon}$; réduction fort remarquable et analogue à celle qui sert de base à la formule du n° 46.

L'intégrale $\displaystyle\int\frac{dz}{(3\pm z^{2})\sqrt[3]{(1\pm 3z^{2})}}$ que nous venons de déterminer dans ses deux différens cas, se trouve mentionnée, page 116, dans la liste que Fuss a donnée des ouvrages d'Euler; mais le mémoire qui la concerne n'a pas encore été publié. Nous avons voulu faire voir par un exemple aussi remarquable, que la théorie des fonctions elliptiques conduit d'une manière sûre aux expressions les plus simples des intégrales qui s'y rapportent.

\subsection*{(144).}

Indépendamment des deux cas dont nous venons de parler, il y en a un troisième où l'intégrale $Z$ ne dépend que des arcs de cercle et des logarithmes. C'est celui où l'on aurait
\[ Z=\int\frac{3\,dz}{(3-z^{2})\sqrt[3]{(1+z^{2})}}, \]
intégrale qui se rapporte à la formule générale du n° 139, en faisant $\nu=1$, $\mu=-\frac{1}{3}$ et $\alpha^{3}=-4$. On ne saurait alors éviter de rencontrer dans l'intégrale $P$ des fonctions elliptiques dont le paramètre est imaginaire; mais la réduction de ces intégrales conduira à un résultat entièrement indépendant des fonctions elliptiques. Nous n'entrerons point dans le détail de ces réductions, et nous nous bornerons à prouver par une autre voie, que l'integrale $Z$ peut être déterminée par les arcs de cercle et les logarithmes.

Pour cet effet, soit $1+z^{2}=4y^{3}$, on aura la transformée
\[ Z=\frac{9}{2\sqrt[3]{4}}\int\frac{y\,dy}{(1-y^{3})\sqrt{(4y^{3}-1)}}: \]
Mais si on considère l'intégrale $T=\displaystyle\int\frac{t\,dt}{(t^{3}-1)^{\frac{2}{3}}}$, et qu'on fasse $\sqrt[3]{(t^{3}-1)}=\frac{t}{y}-1$, on aura $t=\frac{2y\sqrt{3}}{\sqrt{3}-\sqrt{(4y^{3}-1)}}$, $\frac{t^{2}\,dt}{(t^{3}-1)^{\frac{2}{3}}}=\frac{dt}{y}-\frac{t\,dy}{yy}$;

\origpage{209}
donc
\[ T=\frac{3}{2}\int\frac{y\,dy}{1-y^{3}}+\frac{3\sqrt{3}}{2}\int\frac{y\,dy}{(1-y^{3})\sqrt{(4y^{3}-1)}}. \]
Ainsi on connaîtra l'intégrale $Z$ si $T$ est connu; or il suffit de faire $t^{3}=\frac{1}{1-u^{3}}$, et on aura $T=\displaystyle\int\frac{du}{1-u^{3}}$; donc $Z$ se déterminera par l'équation
\[ \frac{\sqrt[3]{4}}{\sqrt{3}}Z=\int\frac{du}{1-u^{3}}-\frac{3}{2}\int\frac{y\,dy}{1-y^{3}}, \]
et par conséquent $Z$ peut s'exprimer par les arcs de cercle et les logarithmes.

\section*{Des cas principaux où l'on peut évaluer, par les fonctions elliptiques, l'intégrale $\left(\frac{p}{q}\right)=\displaystyle\int\frac{x^{p-1}\,dx}{\sqrt[n]{(1-x^{n})^{n-q}}}$, prise depuis $x=0$ jusqu'à $x=1$.}

\subsection*{(145).}

Ces intégrales dont nous traiterons fort au long dans la seconde partie, peuvent se déterminer dans plusieurs cas par des fonctions elliptiques complètes de la première espèce. Ces cas sont ceux où l'exposant $n$ est égal à l'un des nombres 3, 4, 6, 8 et 12. Nous allons les examiner successivement.

\section*{Premier cas, $n=3$.}

La seule transcendante à déterminer est
\[ \left(\frac{1}{1}\right)=2^{\frac{1}{3}}\int\frac{dx}{\sqrt{(1-x^{3})}}, \]
et sa valeur déjà trouvée (art.\ 39) est
\[ \left(\frac{1}{1}\right)=2^{\frac{4}{3}}3^{-\frac{3}{4}}F^{1}(\sin 15^{\circ}). \]

\section*{Second cas, $n=4$.}

\subsection*{(146).}

Il suffit encore dans ce cas, de déterminer la transcendante
\[ \left(\frac{1}{1}\right)=2^{\frac{1}{2}}\int\frac{dx}{\sqrt{(1-x^{4})}}=\int\frac{dz}{\sqrt{(1+z^{4})}}, \]

\origpage{210}
l'intégrale par rapport à $z$ devant être prise depuis $z=0$ jusqu'à $z=\infty$; or l'une ou l'autre des deux formes conduit au même résultat
\[ \left(\frac{1}{1}\right)=F^{1}(\sin 45^{\circ}). \]

\section*{Troisième cas, $n=6$.}

\subsection*{(147).}

Toutes les intégrales représentées par $\left(\frac{p}{q}\right)$ peuvent être déterminées dans ce cas par la seule transcendante
\[ \left(\frac{1}{1}\right)=2^{\frac{2}{3}}\int\frac{dx}{\sqrt{(1-x^{6})}}=2^{\frac{2}{3}}\cos\frac{\pi}{6}\int\frac{dz}{\sqrt{(1+z^{6})}}. \]
La première forme donnera, en faisant $\frac{1}{x^{2}}=1+z^{2}$, le résultat $\left(\frac{1}{1}\right)=2^{\frac{2}{3}}3^{-\frac{1}{4}}F^{1}(\sin 15^{\circ})$. La seconde donnerait, en faisant $\frac{1}{z^{2}}=y^{2}-1$, le résultat $\left(\frac{1}{1}\right)=2^{\frac{2}{3}}3^{-\frac{3}{4}}F^{1}(\sin 75^{\circ})$, et de ces deux valeurs on déduit $F^{1}(\sin 75^{\circ})=\sqrt{3}E^{1}(\sin 15^{\circ})$\ednote{Printed $E^{1}(\sin 15^{\circ})$; the two values just found give $F^{1}(\sin 75^{\circ})=\sqrt{3}F^{1}(\sin 15^{\circ})$, an apparent misprint for $F^{1}$.}, ce qui est une propriété connue de ces deux fonctions.

\section*{Quatrième cas, $n=8$.}

\subsection*{(148).}

Il faut dans ce cas évaluer les deux transcendantes
\[ \begin{aligned} \left(\frac{1}{1}\right)&=2^{\frac{3}{4}}\int\frac{dx}{\sqrt{(1-x^{8})}}=2^{\frac{3}{4}}\cos\frac{\pi}{8}\int\frac{dz}{\sqrt{(1+z^{8})}}\\ \left(\frac{2}{2}\right)&=2^{\frac{1}{2}}\int\frac{x\,dx}{\sqrt{(1-x^{8})}}=2^{\frac{1}{2}}\cos\frac{\pi}{4}\int\frac{z\,dz}{\sqrt{(1+z^{8})}}. \end{aligned} \]
La seconde, en mettant $x^{2}$ à la place de $x$, devient $2^{-\frac{1}{2}}\displaystyle\int\frac{dx}{\sqrt{(1-x^{4})}}$, et sa valeur est $\left(\frac{2}{2}\right)=\frac{1}{2}F^{1}(\sin 45^{\circ})$.

Pour évaluer l'intégrale $\left(\frac{1}{1}\right)$, nous choisirons la seconde forme où il s'agit d'intégrer la différentielle $\frac{dz}{\sqrt{(1+z^{8})}}$ entre les limites $z=0$, $z=\infty$. Supposons cette intégrale trouvée depuis $z=0$ jusqu'à $z=1$, et appelons $Q$ la partie comprise depuis $z=1$ jusqu'à $z=\infty$;

\origpage{211}
si l'on fait $z=\frac{1}{u}$, on aura la transformée $Q=\displaystyle\int\frac{-u^{2}\,du}{\sqrt{(1+u^{8})}}$, qui devra être prise depuis $u=1$ jusqu'à $u=0$; en changeant son signe, l'intégrale devra être prise depuis $u=0$ jusqu'à $u=1$; et comme alors on peut mettre $z$ à la place de $u$, il est clair que l'intégrale cherchée $\displaystyle\int\frac{dz}{\sqrt{(1+z^{8})}}$, prise depuis $z=0$ jusqu'à $z=\infty$, sera égale à l'intégrale $T=\displaystyle\int\frac{dz(1+z^{2})}{\sqrt{(1+z^{8})}}$, prise depuis $z=0$ jusqu'à $z=1$.

Soit $1+z^{4}=pz^{2}$, on aura $\sqrt{(1+z^{8})}=z^{2}\sqrt{(p^{2}-2)}$, et par conséquent $T=\displaystyle\int\frac{dz(1+z^{2})}{z^{2}\sqrt{(p^{2}-2)}}$; mais l'équation $1+z^{4}=pz^{2}$ donne $1-z^{2}=z\sqrt{(p-2)}$, $\frac{1}{z}-z=\sqrt{(p-2)}$, $\frac{dz}{zz}+dz=-\frac{\frac{1}{2}dp}{\sqrt{(p-2)}}$: donc on aura
\[ T=\int\frac{\tfrac{1}{2}dp}{\sqrt{(p-2)}\cdot\sqrt{(p^{2}-2)}}, \]
cette intégrale étant prise depuis $p=2$ jusqu'à $p=\infty$.

Soit $p=2+q^{2}$, on aura la nouvelle transformée
\[ T=\int\frac{dq}{\sqrt{[(q^{2}+2+\sqrt{2})(q^{2}+2-\sqrt{2})]}}, \]
qui devra être prise depuis $q=0$ jusqu'à $q=\infty$. Soit $m^{2}=2+\sqrt{2}$, $q=m\cot\varphi$ et $c^{2}=2\sqrt{2}-2$, on aura l'intégrale indéfinie $T=\displaystyle\frac{1}{m}\int\frac{d\varphi}{\sqrt{(1-c^{2}\sin^{2}\varphi)}}=\frac{1}{m}F(c,\varphi)$, et l'intégrale complète $T^{1}=\frac{1}{m}F^{1}(c)$; donc enfin la valeur de la transcendante cherchée $\left(\frac{1}{1}\right)$ sera
\[ \left(\frac{1}{1}\right)=2^{-\frac{1}{4}}F^{1}(c). \]
On a vu d'ailleurs (n° 64) que pour cette valeur du module $c$ à laquelle répond le module complémentaire $b=\sqrt{2}-1=\operatorname{tang}\frac{\pi}{8}$, on a $F^{1}(c)=\sqrt{2}F^{1}(b)$. Ainsi la transcendante $\left(\frac{1}{1}\right)$ s'exprime également par $2^{-\frac{1}{4}}F^{1}(c)$, ou par $2^{\frac{1}{4}}F^{1}(b)=2^{\frac{1}{4}}F^{1}\left(\operatorname{tang}\frac{\pi}{8}\right)$.

\origpage{212}

\section*{Cinquième cas, $n=12$.}

\subsection*{(149).}

Toutes les intégrales comprises dans ce cas seront déterminées, si on connaît les trois transcendantes
\[ \begin{aligned} \left(\frac{1}{1}\right)&=2^{\frac{5}{6}}\int\frac{dx}{\sqrt{(1-x^{12})}}=2^{\frac{5}{6}}\cos\frac{\pi}{12}\int\frac{dz}{\sqrt{(1+z^{12})}}\\ \left(\frac{2}{2}\right)&=2^{\frac{4}{6}}\int\frac{x\,dx}{\sqrt{(1-x^{12})}}=2^{\frac{4}{6}}\cos\frac{2\pi}{12}\int\frac{z\,dz}{\sqrt{(1+z^{12})}}\\ \left(\frac{3}{3}\right)&=2^{\frac{3}{6}}\int\frac{x^{2}\,dx}{\sqrt{(1-x^{12})}}=2^{\frac{3}{6}}\cos\frac{3\pi}{12}\int\frac{z^{2}\,dz}{\sqrt{(1+z^{12})}}, \end{aligned} \]
les intégrales en $x$ étant prises depuis $x=0$ jusqu'à $x=1$, et les intégrales en $z$ étant prises depuis $z=0$ jusqu'à $z=\infty$.

La seconde et la troisième se ramènent immédiatement aux formules déjà intégrées. En effet, si au lieu de $x^{2}$ on met $x$, la seconde formule donnera
\[ \left(\frac{2}{2}\right)=2^{-\frac{1}{3}}\int\frac{dx}{\sqrt{(1-x^{6})}}=2^{-\frac{1}{3}}3^{-\frac{1}{4}}F^{1}(\sin 15^{\circ}); \]
et si au lieu de $x^{3}$ on met $x$, la troisième donnera
\[ \left(\frac{3}{3}\right)=\frac{\sqrt{2}}{3}\int\frac{dx}{\sqrt{(1-x^{4})}}=\tfrac{1}{3}F^{1}(\sin 45^{\circ}). \]

Venons maintenant à la première intégrale qui paraît présenter d'assez grandes difficultés.

En choisissant la première forme, il s'agira de trouver l'intégrale $X=\displaystyle\int\frac{dx}{\sqrt{(1-x^{12})}}$ depuis $x=0$ jusqu'à $x=1$. Pour cela, je fais $x^{2}-\frac{1}{x^{2}}=q$\ednote{Printed thus; the lines that follow, and the equation $\frac{1}{x^{2}}-x^{2}=q$ further on, require $\frac{1}{x^{2}}-x^{2}=q$, an apparent misprint.}, d'où je tire successivement $\frac{1}{x^{4}}+x^{4}=q^{2}+2$. $\frac{1}{x^{6}}-x^{6}=q^{3}+3q$, $1-x^{12}=x^{6}(q^{3}+3q)$, $\frac{dx}{\sqrt{(1-x^{12})}}=\frac{dx}{x^{3}\sqrt{(q^{3}+3q)}}$; il reste à trouver la valeur de $\frac{dx}{x^{3}}$; or de l'équation $\frac{1}{x^{2}}-x^{2}=q$, on déduit $\frac{1}{x^{2}}=\frac{1}{2}q+\sqrt{(\frac{1}{4}q^{2}+1)}$; donc $-\frac{2\,dx}{x^{3}}=\frac{1}{2}dq+\frac{1}{2}\cdot\frac{q\,dq}{\sqrt{(q^{2}+4)}}$, donc enfin la transformée sera
\[ X=\tfrac{1}{4}\int\frac{dq}{\sqrt{(q^{3}+3q)}}+\tfrac{1}{4}\int\frac{dq\sqrt{q}}{\sqrt{(q^{2}+3)}\cdot\sqrt{(q^{2}+4)}}, \]

\origpage{213}
ces deux nouvelles intégrales devant être prises depuis $q=0$ jusqu'à $q=\infty$. Si dans la première on fait $q=p^{2}$, on aura $\displaystyle\int\frac{dq}{\sqrt{(q^{3}+3q)}}=\int\frac{2\,dp}{\sqrt{(p^{4}+3)}}$; cette intégrale qui doit être prise également depuis $p=0$ jusqu'à $p=\infty$, se trouvera par la formule du n° 38, et on aura pour résultat $\displaystyle\int\frac{dq}{\sqrt{(q^{3}+3q)}}=\frac{2}{\sqrt[4]{3}}F^{1}(\sin 45^{\circ})$.

\subsection*{(150).}

Il ne reste donc à trouver que l'intégrale
\[ Q=\int\frac{dq\sqrt{q}}{\sqrt{[(q^{2}+3)(q^{2}+4)]}}=\int\frac{dq\sqrt{q}}{\sqrt{(q^{4}+7q^{2}+12)}}; \]
cette formule rentre dans le cas VII de l'art.\ 137; on fera donc $m^{4}=12$ et $q^{2}+m^{2}=qz$, ce qui donnera d'abord $q^{4}+7q^{2}+12=q^{2}(z^{2}-2m^{2}+7)$, et $Q=\displaystyle\int\frac{q^{-\frac{1}{2}}\,dq}{\sqrt{(z^{2}-2m^{2}+7)}}$. Mais de l'équation $q^{2}+m^{2}=qz$, on déduit successivement
\[ \begin{aligned} q+m&=q^{\frac{1}{2}}\sqrt{(z+2m)}\\ q-m&=\pm q^{\frac{1}{2}}\sqrt{(z-2m)}\\ 2q^{\frac{1}{2}}&=\sqrt{(z+2m)}\pm\sqrt{(z-2m)}\\ q^{-\frac{1}{2}}\,dq&=\frac{\tfrac{1}{2}dz}{\sqrt{(z+2m)}}\pm\frac{\tfrac{1}{2}dz}{\sqrt{(z-2m)}}. \end{aligned} \]
Donc on aura la transformée
\[ Q=\int\frac{\tfrac{1}{2}dz}{\sqrt{(z+2m)}\cdot\sqrt{(z^{2}-2m^{2}+7)}}\pm\int\frac{\tfrac{1}{2}dz}{\sqrt{(z-2m)}\cdot\sqrt{(z^{2}-2m^{2}+7)}}; \]
le double signe $\pm$ est relatif aux deux valeurs de $q$ qui répondent à une même valeur de $z$. La quantité $z$ a pour \emph{maximum}, $z=2m$, lequel a lieu lorsque $q=m$; et la correspondance de ces deux variables est telle, que depuis $q=0$ jusqu'à $q=m$, la valeur de $z$ décroît depuis $z=\infty$ jusqu'à $z=2m$; et depuis $q=m$ jusqu'à $q=\infty$, la valeur de $z$ croît depuis $z=2m$ jusqu'à $z=\infty$.

De là on voit que depuis $q=0$ jusqu'à $q=m$, l'intégrale $Q$ sera représentée par la différence des intégrales
\[ \int\frac{\tfrac{1}{2}dz}{\sqrt{(z-2m)}\sqrt{(z^{2}-2m^{2}+7)}}-\int\frac{\tfrac{1}{2}dz}{\sqrt{(z+2m)}\sqrt{(z^{2}-2m^{2}+7)}}; \]

\origpage{214}
et la somme de ces intégrales représentera l'intégrale $Q$ prise depuis $q=m$ jusqu'à $q=\infty$. Donc en ajoutant ces deux parties, l'intégrale totale $Q$ sera exprimée par la seule intégrale
\[ \int\frac{dz}{\sqrt{(z-2m)}\sqrt{(z^{2}-2m^{2}+7)}}. \]
Soit $z=2m+p^{2}$, et on aura enfin la transformée
\[ Q=\int\frac{2\,dp}{\sqrt{(p^{4}+4mp^{2}+2m^{2}+7)}}, \]
laquelle devra être intégrée depuis $p=0$ jusqu'à $p=\infty$. Comparant cette intégrale à la formule du n° 38, et faisant $c=\frac{\sqrt{2}-\sqrt[4]{3}}{1+\sqrt{3}}$, on aura pour résultat $Q=4\sin 15^{\circ}\cdot F^{1}(c)$; donc
\[ X=\frac{1}{2\sqrt[4]{3}}F^{1}(\sin 45^{\circ})+\sin 15^{\circ}\cdot F^{1}(c); \]
et enfin la transcendante cherchée
\[ \left(\frac{1}{1}\right)=2^{\frac{5}{6}}X=2^{-\frac{1}{6}}3^{-\frac{1}{4}}F^{1}\sin(45^{\circ})+2^{\frac{5}{6}}\sin 15^{\circ}F^{1}(c). \]
Ainsi les transcendantes nécessaires pour faire connaître toutes les valeurs de l'intégrale $\left(\frac{p}{q}\right)$ dans le cas de $n=12$, se déterminent par les trois fonctions complètes de première espèce $F^{1}(\sin 45^{\circ})$, $F^{1}(\sin 15^{\circ})$, $F^{1}(c)$.

\subsection*{(151).}

Cherchons maintenant une autre valeur de la transcendante $\left(\frac{1}{1}\right)$, au moyen de l'intégrale $Z=\displaystyle\int\frac{dz}{\sqrt{(1+z^{12})}}$, prise depuis $z=0$ jusqu'à $z=\infty$. J'observe que si on prend d'abord cette intégrale depuis $z=0$ jusqu'à $z=1$, qu'ensuite pour avoir l'autre partie depuis $z=1$ jusqu'à $z=\infty$, on mette $\frac{1}{z}$ au lieu de $z$, on aura, par la réunion des deux parties,
\[ Z=\int\frac{dz(1+z^{4})}{\sqrt{(1+z^{12})}}, \]
intégrale qu'il faudra prendre depuis $z=0$ jusqu'à $z=1$.

\origpage{215}

Soit $1+z^{4}=pz^{2}$, on aura $1+z^{12}=z^{6}(p^{3}-3p)$, donc $\frac{dz(1+z^{4})}{\sqrt{(1+z^{12})}}=\frac{dz}{z}\cdot\frac{p}{\sqrt{(p^{3}-3p)}}$; mais de l'équation $1+z^{4}=pz^{2}$, on tire $\frac{dz}{z}=\frac{1}{2}\cdot\frac{dp}{\sqrt{(p^{2}-4)}}$; donc on a la transformée
\[ Z=\int\frac{\tfrac{1}{2}dp\sqrt{p}}{\sqrt{(p^{4}-7p^{2}+12)}}, \]
laquelle devra être intégrée depuis $p=2$ jusqu'à $p=\infty$.

Cette formule est semblable à celle que nous avons traitée dans l'article précédent; on fera donc de même $m^{4}=12$, $p^{2}+m^{2}=p\nu$, ce qui donnera
\[ \begin{aligned} p+m&=p^{\frac{1}{2}}\sqrt{(\nu+2m)}\\ p+m&=p^{\frac{1}{2}}\sqrt{(\nu-2m)}, \end{aligned} \]
\ednote{The second line is printed with $p+m$; from $p^{2}+m^{2}=p\nu$ it should read $p-m$, an apparent misprint.}
et il faut remarquer qu'il n'y a point ici d'ambiguité, parce qu'on a toujours $p>m$. On aura ensuite
\[ \begin{aligned} p^{4}-7p^{2}+12&=p^{2}(\nu^{2}-2m^{2}-7)\\ 2p^{\frac{1}{2}}&=\sqrt{(\nu+2m)}+\sqrt{(\nu-2m)}\\ p^{-\frac{1}{2}}dp&=\frac{\tfrac{1}{2}d\nu}{\sqrt{(\nu+2m)}}+\frac{\tfrac{1}{2}d\nu}{\sqrt{(\nu-2m)}}; \end{aligned} \]
donc on aura la transformée
\[ Z=\tfrac{1}{4}\int\frac{d\nu}{\sqrt{(\nu+2m)}\sqrt{(\nu^{2}-2m^{2}-7)}}+\tfrac{1}{4}\int\frac{d\nu}{\sqrt{(\nu-2m)}\sqrt{(\nu^{2}-2m^{2}+7)}}, \]
\ednote{The second radical is printed with $+7$; from $p^{4}-7p^{2}+12=p^{2}(\nu^{2}-2m^{2}-7)$ it should read $\nu^{2}-2m^{2}-7$, an apparent misprint.}
dont les deux parties doivent être intégrées depuis $\nu=\sqrt{(2m^{2}+7)}$ jusqu'à $\nu=\infty$. Soit $\sqrt{(2m^{2}+7)}=2+\sqrt{3}=\mu$, et $\nu+2m=x^{2}$, la première partie deviendra
\[ \int\frac{\tfrac{1}{2}dx}{\sqrt{[(x^{2}-2m+\mu)(x^{2}-2m-\mu)]}}. \]
Soit encore $x^{2}=\frac{2m+\mu}{\cos^{2}\varphi}$, et $c^{2}=\frac{\mu-2m}{2\mu}$, ou $c=\frac{\sqrt{2}-\sqrt[4]{3}}{1+\sqrt{3}}$, et la transformée sera $\frac{1}{2\sqrt{2\mu}}\displaystyle\int\frac{d\varphi}{\sqrt{(1-c^{2}\sin^{2}\varphi)}}$, dont la valeur dans les limites requises $=\frac{1}{2\sqrt{2\mu}}F^{1}(c)$.

\origpage{216}

On trouvera semblablement qu'en faisant $b=\frac{\sqrt{2}+\sqrt[4]{3}}{1+\sqrt{3}}$, la seconde intégrale a pour valeur $\frac{1}{2\sqrt{2\mu}}F^{1}(b)$; et le module de celle-ci a été désigné par $b$, parce qu'en effet d'après les valeurs trouvées, on a $b^{2}+c^{2}=1$, de sorte que ces modules sont complémens l'un de l'autre. Cela posé, en ajoutant les deux parties, on aura l'intégrale cherchée
\[ Z=\frac{1}{2\sqrt{2\mu}}[F^{1}(b)+F^{1}(c)]=\frac{\sqrt{3}-1}{4}[F^{1}(b)+F^{1}(c)], \]
ce qui donnera cette seconde valeur de la transcendante $\left(\frac{1}{1}\right)$,
\[ \left(\frac{1}{1}\right)=2^{-\frac{5}{3}}[F^{1}(b)+F^{1}(c)]. \]
Egalant les deux valeurs trouvées, on a une première relation entre les trois fonctions $F^{1}(\sin 45^{\circ})$, $F^{1}(b)$, $F^{1}(c)$. Mais on peut en obtenir une seconde qui servira à établir les rapports des fonctions $F^{1}(b)$, $F^{1}(c)$, tant entre elles qu'avec la fonction $F^{1}(\sin 45^{\circ})$.

\subsection*{(152).}

Pour parvenir à ce résultat, il faut chercher directement l'intégrale
\[ M=\int\frac{x^{3}dx}{\sqrt[4]{(1-x^{12})}}, \]
laquelle étant prise depuis $x=0$ jusqu'à $x=1$, sera la valeur de la transcendante $\left(\frac{4}{9}\right)$. Soit d'abord $x^{4}=1-y^{4}$, on aura la transformée
\[ M=\int\frac{y^{2}dy}{\sqrt[4]{(y^{8}-3y^{4}+3)}}, \]
qu'il faut intégrer depuis $y=0$ jusqu'à $y=1$. Si l'on fait ensuite $\nu=\sqrt[4]{3}$, $y^{4}+\nu^{2}=y^{2}z$, ce qui donnera $y^{2}=\frac{1}{2}z-\sqrt{(\frac{1}{4}z^{2}-\nu^{2})}$, $y^{8}-3y^{4}+3=y^{4}(z^{2}-2\nu^{2}-3)$, la transformée en $z$ sera
\[ M=-\tfrac{1}{4}\int\frac{dz}{\sqrt[4]{(z^{2}-2\nu^{2}-3)}}+\tfrac{1}{4}\int\frac{z\,dz}{\sqrt{(z^{2}-4\nu^{2})}\cdot\sqrt[4]{(z^{2}-2\nu^{2}-3)}}. \]
Soit $z^{2}-2\nu^{2}-3=u^{4}$, et on aura pour dernière transformée
\[ M=\int\frac{\tfrac{1}{2}u^{2}du}{\sqrt{(u^{4}-2\nu^{2}+3)}}-\int\frac{\tfrac{1}{2}u^{2}du}{\sqrt{(u^{4}+2\nu^{2}+3)}}, \]

\origpage{217}
d'où l'on voit que l'intégrale cherchée ne dépend que des fonctions elliptiques.

En vertu des limites fixées, les deux dernières intégrales devront être prises depuis $u=1$ jusqu'à $u=\infty$. Mais comme chacune de ces intégrales deviendrait infinie dans la dernière limite, il faut, pour éviter cet inconvénient, faire
\[ \begin{aligned} P&=\int\left(\frac{u^{2}du}{\sqrt{(u^{4}-2\nu^{2}+3)}}-du\right)\\ Q&=\int\left(du-\frac{u^{2}du}{\sqrt{(u^{4}+2\nu^{2}+3)}}\right). \end{aligned} \]
Par ce moyen, les deux intégrales $P$ et $Q$ seront toujours des quantités finies, et après avoir trouvé leurs valeurs complètes $P^{1}$, $Q^{1}$, on en déduira $M$ ou $\left(\frac{4}{9}\right)=\frac{1}{2}(P^{1}+Q^{1})$.

\subsection*{(153).}

Pour avoir l'intégrale $P$, soit $\beta^{2}=2\sqrt{3}-3$, et $u^{2}-\sqrt{(u^{4}-\beta^{2})}=p$, ou $u^{2}=\frac{1}{2}p+\frac{\beta^{2}}{2p}$, on aura la transformée $P=\displaystyle\int\frac{dp\sqrt{(\frac{1}{2}p)}}{\sqrt{(p^{2}+\beta^{2})}}$. Soit ensuite $p=\beta\operatorname{tang}^{2}\frac{1}{2}\varphi$ et $c^{2}=\frac{1}{2}$, on aura
\[ P=\sqrt{(\tfrac{1}{2}\beta)}\cdot\int\frac{d\varphi\operatorname{tang}^{2}\tfrac{1}{2}\varphi}{\sqrt{(1-c^{2}\sin^{2}\varphi)}}; \]
cette intégrale se trouve par les formules de l'art.\ 138, et on a
\[ P=\sqrt{(\tfrac{1}{2}\beta)}\cdot[2\Delta\operatorname{tang}\tfrac{1}{2}\varphi+F(\varphi)-2E(\varphi)]. \]
Il faut maintenant prendre cette intégrale depuis $\varphi=0$ qui répond à $u=\infty$, jusqu'à la valeur de $\varphi$ qui répond à $u=1$; on a dans cette dernière limite $\operatorname{tang}\frac{1}{2}\varphi=\frac{1-\sqrt{(1-\beta^{2})}}{\beta}$, ou $\cos\varphi=\sqrt{\left(\frac{1-\beta}{1+\beta}\right)}$. Or j'observe que l'angle $\varphi$ ainsi déterminé, est celui qui pour le module $c=\sqrt{\frac{1}{2}}$, donne $F(\varphi)=\frac{2}{3}F^{1}$; car nous avons trouvé (art.\ 24) qu'on satisfait à l'équation $F(\psi)=\frac{1}{3}F^{1}$, par les valeurs $\cos\psi=\sqrt{\beta}$, $\sin\psi=\sqrt{(1-\beta)}$. Si ensuite on suppose $F(\varphi)+F(\psi)=F^{1}$, ou $\operatorname{tang}\varphi\operatorname{tang}\psi=\frac{1}{b}=\sqrt{2}$, on en déduira $\operatorname{tang}\varphi=\sqrt{\left(\frac{2\beta}{1-\beta}\right)}$ ou $\cos\varphi=\sqrt{\left(\frac{1-\beta}{1+\beta}\right)}$, ce qui est la valeur de $\varphi$ dont il s'agit. Au reste en faisant $\alpha^{2}=2\sqrt{3}+3$, on aura plus simplement $\cos\varphi=\frac{\alpha-1}{\alpha+1}$, $\sin\varphi=\frac{2\sqrt{\alpha}}{\alpha+1}$, $\Delta(\varphi)=\frac{1}{2}(\alpha-1)$, $\operatorname{tang}\frac{1}{2}\varphi=\frac{1}{\sqrt{\alpha}}$, $\sin\psi=\frac{2}{\alpha+1}$. Mais

\origpage{218}
puisque la valeur de $\varphi$ est celle qui donne $F(\varphi)=\frac{2}{3}F^{1}$, on aura en même temps par les formules du n° 32, $3E(\varphi)=2E^{1}+c^{2}\sin\varphi\sin\psi\times(2-\sin\psi)$, ce qui donne
\[ E(\varphi)=\tfrac{2}{3}E^{1}+\frac{4}{3}\cdot\frac{\alpha\sqrt{\alpha}}{(\alpha+1)^{3}}. \]
Substituant toutes ces valeurs dans celle de $P$, il viendra
\[ P^{1}=\frac{\sqrt{2\beta}}{3}(F^{1}-2E^{1})+\sqrt{\left(\frac{\beta}{2\alpha}\right)}\cdot\left(\alpha-1-\frac{8}{3}\cdot\frac{\alpha^{2}}{(\alpha+1)^{3}}\right). \]

\subsection*{(154).}

Venons maintenant à l'intégrale $Q$. Si l'on fait $\sqrt{(u^{4}+\alpha^{2})}=u^{2}+q$, ou $u^{2}=\frac{\alpha^{2}-q^{2}}{2q}$, on aura d'abord $Q=\displaystyle\int\frac{dq\sqrt{(\frac{1}{2}q)}}{\sqrt{(\alpha^{2}-q^{2})}}$. Soit ensuite $q=\alpha\cos^{2}\omega$, et on aura $Q=\sqrt{\alpha}\displaystyle\int\frac{d\omega\cos^{2}\omega}{\sqrt{(1-\frac{1}{2}\sin^{2}\omega)}}$, d'où l'on tire
\[ Q=\sqrt{\alpha}\cdot[2E(\omega)-F(\omega)]+\text{const}, \]
le module de ces fonctions étant toujours $c=\sqrt{\frac{1}{2}}$.

Cette intégrale doit être prise entre les deux valeurs de $\omega$, qui donnent $u=1$ et $u=\infty$; dans la seconde limite on a $\omega=\frac{1}{2}\pi$, et dans la première $\cos^{2}\omega=\frac{\sqrt{(1+\alpha^{2})}-1}{\alpha}=\frac{\sqrt{3}}{\alpha}$, ou $\cos^{2}\omega=\beta$; car des deux équations $\alpha^{2}=2\sqrt{3}+3$, $\beta^{2}=2\sqrt{3}-3$, on déduit $\alpha^{2}\beta^{2}=3$, ou $\alpha\beta=\sqrt{3}$. Mais puisqu'on a $\cos\omega=\sqrt{\beta}$, l'angle $\omega$ est le même qui a été désigné dans l'article précédent par $\psi$, et on a par conséquent pour la valeur complète de $Q$,
\[ Q^{1}=\sqrt{\alpha}\cdot[2E^{1}-F^{1}-2E(\psi)+F(\psi)]. \]
Or on a $F(\psi)=\frac{1}{3}F^{1}$, $E(\psi)=\frac{1}{3}E^{1}+\frac{1}{6}\sin\psi\sin\varphi(1+\sin\psi)$; donc
\[ Q^{1}=\frac{2\sqrt{\alpha}}{3}(2E^{1}-F^{1})-\frac{4}{3}\cdot\frac{\alpha(\alpha+3)}{(\alpha+1)^{3}}. \]
Si on ajoute maintenant les valeurs de $P^{1}$ et $Q^{1}$, et qu'on ait égard aux réductions que fournissent les équations $\sqrt{(\frac{1}{2}\beta)}=\frac{2\sqrt{\alpha}}{\alpha^{2}-1}$, $\alpha^{4}=6\alpha^{2}+3$, on trouvera que les parties algébriques se détruisent mutuellement, et qu'on a
\[ P^{1}+Q^{1}=\tfrac{2}{3}(2E^{1}-F^{1})(\sqrt{\alpha}-\sqrt{\tfrac{1}{2}\beta}). \]

\origpage{219}
Mais le facteur $\sqrt{\alpha}-\sqrt{\frac{1}{2}\beta}$ se réduit à $\frac{3}{\nu\sqrt{2\alpha}}$; donc enfin la valeur de la transcendante cherchée $\left(\frac{4}{9}\right)=\frac{1}{2}(P^{1}+Q^{1})$ sera
\[ \left(\frac{4}{9}\right)=\frac{1}{\nu\sqrt{(2\alpha)}}(2E^{1}-F^{1}), \]
où l'on a $\nu=\sqrt[4]{3}$, et $\alpha=\nu\sqrt{(2+\sqrt{3})}=2\nu\cos 15^{\circ}$.

Ce résultat peut encore se simplifier; car d'après la propriété connue des fonctions $E^{1}$, $F^{1}$, dont le module est $\sin 45^{\circ}$, on a $\frac{\pi}{2}=F^{1}(2E^{1}-F^{1})$; donc la transcendante $\left(\frac{4}{9}\right)$ s'exprime par la seule fonction de première espèce $F^{1}$, au moyen de cette formule très-simple
\[ \left(\frac{4}{9}\right)=\frac{1}{2\nu\sqrt{2\alpha}}\cdot\frac{\pi}{F^{1}(\sin 45^{\circ})}. \]

\subsection*{(155).}

Soit pour abréger $F^{1}(\sin 45^{\circ})=B$, $F^{1}(\sin 15^{\circ})=C$; nous avons déjà trouvé $\left(\frac{2}{2}\right)=\frac{2^{-\frac{1}{3}}}{\nu}C$, $\left(\frac{3}{3}\right)=\frac{1}{3}B$, et le résultat précédent donne $\left(\frac{4}{9}\right)=\frac{1}{2\omega\sqrt{2\alpha}}\cdot\frac{\pi}{B}$\ednote{Printed $2\omega\sqrt{2\alpha}$; the formula just above has $2\nu\sqrt{2\alpha}$, so $\omega$ is an apparent misprint for $\nu$.}; nous allons faire voir qu'avec ces seules données, on pourra déterminer la transcendante $\left(\frac{1}{1}\right)$.

En effet il résulte des formules connues${}^{(1)}$ qu'on a $\left(\frac{4}{9}\right)\cdot\left(\frac{8}{3}\right)=\frac{\pi\sin\omega}{3\sqrt{6}}$, et $\left(\frac{8}{3}\right)=2^{\frac{1}{3}}\sin\omega\sqrt{\left(\frac{M_{1}B}{3\cos 2\omega}\right)}$, $\omega$ étant $\frac{\pi}{12}$ ou $15^{\circ}$, et $M_{1}$ représentant la transcendante $\left(\frac{1}{1}\right)$; de là on tirera
\[ M_{1}=\left(\frac{1}{1}\right)=\frac{2^{\frac{1}{3}}\alpha B}{3}=\frac{2^{\frac{4}{3}}\nu\cos 15^{\circ}}{3}\cdot B. \]
Ainsi la transcendante $\left(\frac{1}{1}\right)$ ne dépend que de la fonction $B$. Mais les autres valeurs trouvées de cette transcendante vont nous fournir de nouveaux résultats.

On a trouvé par la première méthode (art.\ 150),
\[ \left(\frac{1}{1}\right)=\frac{2^{-\frac{1}{6}}}{\nu}B+2^{\frac{5}{6}}\sin 15^{\circ}\cdot F^{1}(c). \]

\textbf{(1)} Voyez les articles 9 et 18 de la seconde Partie.

\origpage{220}

Donc en comparant ces deux valeurs, il en résultera
\[ F^{1}(c)=\frac{\alpha}{3}B=\frac{2\nu\cos 15^{\circ}}{3}F^{1}(\sin 45^{\circ}). \]
Si on compare la même valeur à celle qui est résultée de la seconde méthode (art.\ 151), on aura
\[ \frac{2^{\frac{4}{3}}\nu\cos 15^{\circ}}{3}B=2^{-\frac{5}{3}}[F^{1}(b)+F^{1}(c)]; \]
donc $F^{1}(b)+F^{1}(c)=\frac{8\nu\cos 15^{\circ}}{3}B=4F^{1}(c)$, ou
\[ F^{1}(b)=3F^{1}(c). \]
Ainsi l'égalité qui avait été trouvée par approximation (n° 84), se trouve confirmée par une démonstration rigoureuse, démonstration qui pouvait être regardée comme un problème d'analyse fort difficile.

Le module $c$ et son complément $b$ qui donnent lieu à cette réduction singulière, sont tels que si on fait $c=\sin\theta$, $b=\cos\theta$, on aura $\sin 2\theta=\operatorname{tang}^{2}15^{\circ}$, et $\operatorname{tang}(45^{\circ}+\theta)=\sqrt{\left(\frac{2}{\sqrt{3}}\right)}$, ou $\operatorname{tang}(45^{\circ}-\theta)=\sqrt{\cos 30^{\circ}}$.

On voit de plus que les fonctions $F^{1}(b)$, $F^{1}(c)$ s'expriment l'une et l'autre par la fonction $F^{1}(\sin 45^{\circ})$, de sorte qu'on a
\[ \left\{\begin{aligned} F^{1}(c)&=\frac{2\cos 15^{\circ}}{\sqrt[4]{27}}F^{1}(\sin 45^{\circ})\\ F^{1}(b)&=2\sqrt[4]{3}\cdot\cos 15^{\circ}\cdot F^{1}(\sin 45^{\circ}). \end{aligned}\right. \]
Ces rapports sont d'autant plus remarquables que les fonctions $F^{1}(b)$, $F^{1}(c)$ sont jusqu'à présent les seuls exemples de fonctions de la première espèce qui aient un rapport connu avec une autre fonction de même espèce $F^{1}(\sin 45^{\circ})$, sans que les deux fonctions comparées soient comprises dans une même série ou échelle formée d'après un module donné, et sans que leurs modules soient complémens l'un de l'autre.

FIN DE LA PREMIÈRE PARTIE.

\origpage{221}

\section*{SECONDE PARTIE.}

\section*{DES INTÉGRALES EULÉRIENNES.}

Quoique le nom d'Euler soit attaché à presque toutes les théories importantes du Calcul intégral, cependant j'ai cru qu'il me serait permis de donner plus spécialement le nom d'\emph{intégrales Eulériennes}, à deux sortes de transcendantes dont les propriétés ont fait le sujet de plusieurs beaux Mémoires d'Euler, et forment la théorie la plus complète que l'on connaisse jusqu'à présent sur les intégrales définies.

La première est l'intégrale $\displaystyle\int\frac{x^{p-1}dx}{\sqrt[n]{(1-x^{n})^{n-q}}}$ qu'on suppose prise entre les limites $x=0$, $x=1$. Nous la représenterons, comme Euler, par le caractère abrégé $\left(\frac{p}{q}\right)$.

La seconde est l'intégrale $\displaystyle\int dx\left(\log\frac{1}{x}\right)^{a-1}$, prise de même entre les limites $x=0$, $x=1$. Euler représente cette intégrale par le symbole $\left[\frac{p}{q}\right]$, en supposant que $a$ soit égal à la fraction rationnelle $\frac{p}{q}$; nous la représenterons plus généralement par $\Gamma(a)$, et nous regarderons $\Gamma(a)$ comme une fonction continue de $a$.

Ayant pour objet de réunir dans cet ouvrage tout ce que la théorie des transcendantes, et surtout celle des intégrales définies, offre de plus remarquable, j'ai dû y comprendre la théorie des intégrales Eulériennes. C'est pourquoi je la donne ici, à quelques additions près, telle que je l'ai déjà publiée dans les \emph{Mémoires de l'Institut}, pour l'année 1810.

\origpage{222}

\section*{Des intégrales Eulériennes de la première espèce.}

\subsection*{(1).}

L'expression $\left(\frac{p}{q}\right)$ par laquelle nous désignons l'intégrale $\displaystyle\int\frac{x^{p-1}dx}{\sqrt[n]{(1-x^{n})^{n-q}}}$, prise depuis $x=0$ jusqu'à $x=1$, est en général une fonction des deux exposans $p$ et $q$, et de l'exposant $n$, qu'on regarde comme des nombres entiers. Mais nous supposons $n$ constant, et notre but est de comparer entre elles les diverses valeurs de $\left(\frac{p}{q}\right)$ qui répondent à une même valeur de $n$, afin de réduire au moindre nombre possible les transcendantes que cette expression renferme.

Nous observerons d'abord que les deux exposans $p$ et $q$ peuvent être échangés entre eux. En effet, si on fait $x^{n}=1-y^{n}$, on aura
\[ \frac{x^{p-1}dx}{\sqrt[n]{(1-x^{n})^{n-q}}}=-\frac{y^{q-1}dy}{\sqrt[n]{(1-y^{n})^{n-p}}}, \]
et la transformée en $y$ devra être intégrée depuis $y=1$ jusqu'à $y=0$. En changeant son signe, elle devra être intégrée depuis $y=0$ jusqu'à $y=1$; et comme alors on peut mettre $x$ à la place de $y$, on aura
\[ \int\frac{x^{p-1}dx}{\sqrt[n]{(1-x^{n})^{n-q}}}=\int\frac{x^{q-1}dx}{\sqrt[n]{(1-x^{n})^{n-p}}}, \]
ou, suivant notre notation,
\[ \left(\frac{p}{q}\right)=\left(\frac{q}{p}\right),\qquad(a) \]
ce qui est la propriété énoncée.

\subsection*{(2).}

Il faut faire voir maintenant que si dans la formule $\left(\frac{p}{q}\right)$, l'un des deux nombres $p$ et $q$ est plus grand que $n$, la formule se ramène aisément au cas où $p$ et $q$ sont compris l'un et l'autre dans les limites 1 et $n$. Pour cela soit $Z=x^{k}(1-x^{n})^{r}$, on aura la différentielle
\[ dZ=kx^{k-1}dx(1-x^{n})^{r-1}-(k+rn)x^{k+n-1}dx(1-x^{n})^{r-1}, \]

\origpage{223}
d'où l'on tire en intégrant,
\[ Z=k\int x^{k-1}dx(1-x^{n})^{r-1}-(k+rn)\int x^{k+n-1}dx(1-x^{n})^{r-1}. \]
Donc si on prend les intégrales entre les limites $x=0$, $x=1$, et qu'en même temps on suppose $k>0$, et $r>0$, afin que $Z$ s'évanouisse dans les deux limites, on aura
\[ \int x^{k+n-1}dx(1-x^{n})^{r-1}=\frac{k}{k+rn}\int x^{k-1}dx(1-x^{n})^{r-1}; \]
d'où il suit qu'en faisant $k+n=p$ et $r=\frac{q}{n}$, il viendra
\[ \left(\frac{p}{q}\right)=\frac{p-n}{p+q-n}\cdot\left(\frac{p-n}{q}\right)\qquad(b). \]
Au moyen de cette formule, toute fonction $\left(\frac{p}{q}\right)$, dans laquelle on a $p>n$, se réduira successivement aux fonctions $\left(\frac{p-n}{q}\right)$, $\left(\frac{p-2n}{q}\right)$, etc., jusqu'à un terme $\left(\frac{p-in}{q}\right)$ ou $\left(\frac{p'}{q}\right)$, dans lequel $p'$ sera le reste de la division de $p$ par $n$.

Arrivé à ce terme, si de son côté $q$ est plus grand que $n$, la fonction $\left(\frac{p'}{q}\right)$ qui est la même que $\left(\frac{q}{p'}\right)$, se réduira successivement aux fonctions $\left(\frac{q-n}{p'}\right)$, $\left(\frac{q-2n}{p'}\right)$, etc.\ jusqu'à un terme $\left(\frac{q'}{p'}\right)$ ou $\left(\frac{p'}{q'}\right)$, dans lequel $q'$ sera le reste de la division de $q$ par $n$.

De là on voit qu'on peut toujours supposer la fonction $\left(\frac{p}{q}\right)$ réduite à une forme où $p$ et $q$ soient compris tous deux dans la suite $1, 2, 3\ldots n$.

\subsection*{(3).}

Cela posé, il y a deux cas principaux où on peut trouver immédiatement la valeur de $\left(\frac{p}{q}\right)$; c'est lorsque $n$ est égal à l'un des deux nombres $p$ et $q$, ou lorsqu'il est égal à leur somme.

Soit, 1°. $q=n$, on aura immédiatement $\left(\frac{p}{n}\right)=\displaystyle\int x^{p-1}dx=\frac{x^{p}}{p}+C$, intégrale qui, étant prise depuis $x=0$ jusqu'à $x=1$, se réduit à $\frac{1}{p}$; d'où résulte
\[ \left(\frac{p}{n}\right)=\frac{1}{p}\qquad(c). \]

\origpage{224}

Soit, 2°. $p+q=n$, ou $p=a$, $q=n-a$, on aura $\left(\frac{p}{q}\right)=\displaystyle\int x^{a-1}dx(1-x^{n})^{-\frac{a}{n}}$; faisant $1-x^{n}=x^{n}z^{n}$, on aura la transformée rationnelle
\[ \left(\frac{p}{q}\right)\text{ ou }\left(\frac{a}{n-a}\right)=\int\frac{z^{n-a-1}dz}{1+z^{n}}, \]
intégrale qui doit être prise depuis $z=0$ jusqu'à $z=\infty$; or par les formules connues (\emph{Eul.}, \emph{Calc.\ intég.}, tom.~I, pag.~252), cette intégrale $=\frac{\pi}{n\sin\frac{a\pi}{n}}$. Donc si on fait $\frac{\pi}{n}=\omega$, on aura
\[ \left(\frac{a}{n-a}\right)=\frac{\omega}{\sin a\omega}.\qquad(d) \]
Excepté ces deux cas généraux, toutes les quantités désignées par $\left(\frac{p}{q}\right)$ sont des transcendantes plus ou moins composées, selon la valeur de $n$, et ne sont point susceptibles d'une évaluation exacte. Mais pour chaque valeur de $n$, on peut exprimer toutes ces transcendantes au moyen d'un petit nombre d'entre elles, et c'est l'objet des recherches suivantes.

\subsection*{(4).}

Observons d'abord qu'en mettant $p+n$ à la place de $p$, l'équation (b) donne
\[ \left(\frac{p}{q}\right)=\frac{p+q}{p}\cdot\left(\frac{p+n}{q}\right). \]
On aurait semblablement
\[ \begin{aligned} \left(\frac{p+n}{q}\right)&=\frac{p+q+n}{p+n}\cdot\left(\frac{p+2n}{q}\right)\\ \left(\frac{p+2n}{q}\right)&=\frac{p+q+2n}{p+2n}\cdot\left(\frac{p+3n}{q}\right)\\ &\text{etc.} \end{aligned} \]
Donc en général, $i$ étant un nombre entier positif à volonté, on aura
\[ \left(\frac{p}{q}\right)=\frac{p+q}{p}\cdot\frac{p+q+n}{p+n}\cdot\frac{p+q+2n}{p+2n}\ldots\frac{p+q+in}{p+in}\cdot\left(\frac{p+\overline{i+1}\cdot n}{q}\right); \]

\origpage{225}
mettant $p+r$ au lieu de $p$, on aura semblablement
\[ \left(\frac{p+r}{q}\right)=\frac{p+q+r}{p+r}\cdot\frac{p+q+r+n}{p+r+n}\cdot\frac{p+q+r+2n}{p+r+2n}\ldots\frac{p+q+r+in}{p+r+in}\cdot\left(\frac{p+r+\overline{i+1}\cdot n}{q}\right). \]
Divisons ces deux formules l'une par l'autre, et faisons pour abréger,
\[ \begin{aligned} M^{0}&=\frac{(p+q)(p+r)}{p\cdot(p+q+r)}\\ M'&=\frac{(p+q+n)(p+r+n)}{(p+n)(p+q+r+n)}\\ M''&=\frac{(p+q+2n)(p+r+2n)}{(p+2n)(p+q+r+2n)}\\ &\text{etc.}, \end{aligned} \]
nous aurons
\[ \frac{\left(\frac{p}{q}\right)}{\left(\frac{p+r}{q}\right)}=M^{0}M'M''\ldots M^{(i)}\frac{\left(\frac{p+\overline{i+1}\cdot n}{q}\right)}{\left(\frac{p+r+\overline{i+1}\cdot n}{q}\right)}. \]
Supposons $r$ positif et $<n$; alors il est clair que la quantité $\left(\frac{p+r+\overline{i+1}\cdot n}{q}\right)$ sera plus petite que $\left(\frac{p+\overline{i+1}\cdot n}{q}\right)$ et plus grande que $\left(\frac{p+\overline{i+2}\cdot n}{q}\right)$. Car si on considère diverses formules $\left(\frac{p}{q}\right)$ qui ne diffèrent que par l'exposant $p$; comme ces formules désignent chacune l'intégrale, prise depuis $x=0$ jusqu'à $x=1$, d'une différentielle $\frac{x^{p-1}dx}{X}$ dont le dénominateur est le même pour toutes, il est évident que l'intégrale sera d'autant plus petite que $p$ sera plus grand.

Mais la formule (b) donne
\[ \left(\frac{p+\overline{i+1}\cdot n}{q}\right)=\frac{p+q+\overline{i+1}\cdot n}{p+\overline{i+1}\cdot n}\cdot\left(\frac{p+\overline{i+2}\cdot n}{q}\right). \]
Donc on a d'une part le rapport
\[ \frac{\left(\frac{p+\overline{i+1}\cdot n}{q}\right)}{\left(\frac{p+r+\overline{i+1}\cdot n}{q}\right)}>1, \]

\origpage{226}
et de l'autre ce même rapport $<\frac{p+q+\overline{i+1}\cdot n}{p+\overline{i+1}\cdot n}$. Si l'on veut donc que ce rapport soit compris entre les limites 1 et $1+\frac{1}{k}$, il faudra prendre $i+1>\frac{qk-p}{n}$, et alors on aura
\[ \frac{\left(\frac{p}{q}\right)}{\left(\frac{p+r}{q}\right)}=M^{0}M'M''\ldots M^{(i)}\left(1+\frac{1}{k'}\right), \]
$k'$ étant plus grand que $k$.

On sait par cette équation combien on approche du rapport de $\left(\frac{p}{q}\right)$ à $\left(\frac{p+r}{q}\right)$, en prolongeant le produit $M^{0}M'M''\ldots$ jusqu'à un terme $M^{(i)}$; et il est clair qu'en continuant ce produit à l'infini, on aura la vraie valeur de ce rapport, laquelle sera
\[ \frac{\left(\frac{p}{q}\right)}{\left(\frac{p+r}{q}\right)}=M^{0}M'M''M''',\ \text{etc.} \]
Maintenant si dans cette équation on échange entre elles les lettres $q$ et $r$, les quantités $M^{0}$, $M'$, $M''$, etc.\ resteront les mêmes, de sorte qu'on aura encore
\[ \frac{\left(\frac{p}{r}\right)}{\left(\frac{p+q}{r}\right)}=M^{0}M'M''M''',\ \text{etc.} \]
Donc par la comparaison de ces équations, on obtient la formule générale
\[ \left(\frac{p}{q}\right)\cdot\left(\frac{p+q}{r}\right)=\left(\frac{p}{r}\right)\cdot\left(\frac{p+r}{q}\right)\ldots\ldots(e). \]

Cette formule dont la découverte appartient à Euler, est une sorte d'équation aux différences finies, qui renferme presque toute la théorie des transcendantes $\left(\frac{p}{q}\right)$. Et d'abord nous allons en déduire l'expression générale des quantités $\left(\frac{p}{q}\right)$.

\origpage{227}

\subsection*{(5).}

Les formules (c) et (d) donnent les valeurs exactes de la fonction $\left(\frac{p}{q}\right)$, toutes les fois que l'un des deux nombres $p$ et $q$ ou leur somme est égale à $n$. Supposons maintenant qu'on connaisse de plus toutes les valeurs de $\left(\frac{p}{q}\right)$ lorsque $p+q=n-1$, et désignons en général par $A_{a}$ la fonction $\left(\frac{n-a-1}{a}\right)$, ensorte qu'on ait
\[ \left(\frac{n-a-1}{a}\right)=A_{a}\ldots\ldots(f). \]
On aura donc successivement $\left(\frac{n-2}{1}\right)=A_{1}$, $\left(\frac{n-3}{2}\right)=A_{2}$; $\left(\frac{n-4}{3}\right)=A_{3}$, etc.; et parce que $\left(\frac{q}{p}\right)=\left(\frac{p}{q}\right)$, on aura aussi $\left(\frac{1}{n-2}\right)=A_{1}$, $\left(\frac{2}{n-3}\right)=A_{2}$, etc.; donc en général
\[ A_{(n-1-a)}=A_{a}\ldots\ldots(g). \]
D'où l'on voit que le nombre des auxiliaires $A_{1}$, $A_{2}$, $A_{3}$, etc., se réduit toujours à $\frac{n-2}{2}$ ou $\frac{n-1}{2}$, selon que $n$ est pair ou impair.

Par exemple, si $n=7$, il y aura trois auxiliaires $A_{1}=\left(\frac{5}{1}\right)$, $A_{2}=\left(\frac{4}{2}\right)$, $A_{3}=\left(\frac{3}{3}\right)$, puisqu'on aurait $A_{4}=\left(\frac{2}{4}\right)=A_{2}$, et $A_{5}=\left(\frac{1}{5}\right)=A_{1}$.

Si $n=8$, il n'y aura encore que trois auxiliaires $A_{1}$, $A_{2}$, $A_{3}$; car on aurait, en vertu de l'équation (g), $A_{4}=A_{3}$, $A_{5}=A_{2}$, $A_{6}=A_{1}$.

Cela posé, au moyen des équations (c), (d), (e) et des auxiliaires données par l'équation (f), nous pourrons trouver l'expression générale de $\left(\frac{p}{q}\right)$ dans deux cas généraux, 1°. lorsque $p+q$ est $<n$; 2°. lorsque $p+q$ est $>n$. Voici comment on y parvient.

\subsection*{(6).}

L'équation (e) donne immédiatement $\left(\frac{n-a-1}{a}\right)\cdot\left(\frac{n-1}{1}\right)=\left(\frac{n-a-1}{1}\right)\cdot\left(\frac{n-a}{a}\right)$; substituant dans celle-ci les valeurs connues

\origpage{228}
par les équations (d) et (g), on aura
\[ \left(\frac{n-a-1}{1}\right)=\frac{A_{a}\sin a\omega}{\sin\omega}\ldots\ldots(h). \]
Ainsi on connaît toute fonction $\left(\frac{p}{q}\right)$ dans laquelle l'un des deux nombres $p$ et $q$ est égal à l'unité; si pour plus de simplicité, on met $a$ à la place de $n-a-1$, on aura pour le même objet la formule
\[ \left(\frac{a}{1}\right)=\frac{A_{a}\sin(a+1)\omega}{\sin\omega}\ldots\ldots(i). \]
La même équation (e) donne $\left(\frac{n-a-2}{a}\right)\cdot\left(\frac{n-2}{1}\right)=\left(\frac{n-a-2}{1}\right)\cdot\left(\frac{n-a-1}{a}\right)$, et substituant les valeurs connues, il en résulte:
\[ \left(\frac{n-a-2}{a}\right)=\frac{A_{a}A_{a+1}}{A_{1}}\cdot\frac{\sin(a+1)\omega}{\sin\omega}. \]
Ainsi on a la valeur de toute fonction $\left(\frac{p}{q}\right)$ dans laquelle $p+q=n-2$.

De l'équation (e) on déduit encore immédiatement
\[ \left(\frac{n-a-3}{a}\right)\cdot\left(\frac{n-3}{1}\right)=\left(\frac{n-a-3}{1}\right)\cdot\left(\frac{n-a-2}{a}\right); \]
ce qui donne
\[ \left(\frac{n-a-3}{a}\right)=\frac{A_{a}A_{a+1}A_{a+2}}{A_{1}A_{2}}\cdot\frac{\sin(a+1)\omega\sin(a+2)\omega}{\sin\omega\sin 2\omega}. \]
Ainsi on connaît la valeur de toute fonction $\left(\frac{p}{q}\right)$, dans laquelle $p+q=n-3$.

En général l'équation (e) donne
\[ \left(\frac{n-a-k}{a}\right)\cdot\left(\frac{n-k}{1}\right)=\left(\frac{n-a-k}{1}\right)\cdot\left(\frac{n-a-k+1}{a}\right), \]
et en mettant les valeurs tirées de l'équation (h), il en résulte
\[ \left(\frac{n-a-k}{a}\right)=\frac{A_{a+k-1}\sin(a+k-1)\omega}{A_{k-1}\sin(k-1)\omega}\cdot\left(\frac{n-a-k+1}{a}\right). \]
Donc on aura en général
\[ \left(\frac{n-a-k}{a}\right)=\frac{A_{a}A_{a+1}\ldots A_{a+k-1}}{A_{1}A_{2}\ldots A_{k-1}}\cdot\frac{\sin(a+1)\omega\sin(a+2)\omega\ldots\sin(a+k-1)\omega}{\sin\omega\sin 2\omega\ldots\sin(k-1)\omega}\ldots(k). \]

\origpage{229}
C'est la valeur de toute fonction $\left(\frac{p}{q}\right)$, dans laquelle $p+q$ est moindre que $n$.

\subsection*{(7).}

Pour avoir l'expression générale de $\left(\frac{p}{q}\right)$ lorsque $p+q$ est plus grand que $n$, observons que l'équation (e) donne aussi
\[ \left(\frac{n-a+k}{a}\right)\cdot\left(\frac{n+k}{1}\right)=\left(\frac{n-a+k}{1}\right)\cdot\left(\frac{n-a+k+1}{a}\right). \]
Or par l'équation (b) on a
\[ \left(\frac{n+k}{1}\right)=\frac{k}{k+1}\cdot\left(\frac{k}{1}\right); \]
par l'équation (i) on a
\[ \left(\frac{k}{1}\right)=\frac{A_{k}\sin(k+1)\omega}{\sin\omega}, \]
et par l'équation (h) on a
\[ \left(\frac{n-a+k}{1}\right)=\frac{A_{a-k-1}\sin(a-k-1)\omega}{\sin\omega}; \]
substituant toutes ces valeurs, il viendra
\[ \left(\frac{n-a+k+1}{a}\right)=\frac{k}{k+1}\cdot\frac{A_{k}\sin(k+1)}{A_{a-k-1}\sin(a-k-1)\omega}\cdot\left(\frac{n-a+k}{a}\right)\ldots(l). \]
\ednote{The numerator is printed $\sin(k+1)$ without $\omega$; $\sin(k+1)\omega$ is meant, as in equation (i).}
Tout se réduit donc à trouver la valeur de $\left(\frac{n-a+1}{a}\right)$.

Or l'équation (e) donne
\[ \left(\frac{n-a}{a}\right)\cdot\left(\frac{n}{1}\right)=\left(\frac{n-a}{1}\right)\cdot\left(\frac{n-a+1}{a}\right), \]
substituant les valeurs connues, il viendra
\[ \left(\frac{n-a+1}{a}\right)=\frac{1}{A_{a-1}}\cdot\frac{\omega\sin\omega}{\sin a\omega\sin(a-1)\omega}\ldots\ldots(m). \]
De là et de l'équation (l) on déduit successivement
\[ \begin{aligned} \left(\frac{n-a+2}{a}\right)&=\frac{1}{2}\cdot\frac{A_{1}}{A_{a-1}A_{a-2}}\cdot\frac{\omega\sin\omega\sin 2\omega}{\sin a\omega\sin(a-1)\omega\sin(a-2)\omega},\\ \left(\frac{n-a+3}{a}\right)&=\frac{1}{3}\cdot\frac{A_{1}A_{2}}{A_{a-1}A_{a-2}A_{a-3}}\cdot\frac{\omega\sin\omega\sin 2\omega\sin 3\omega}{\sin a\omega\sin(a-1)\omega\sin(a-2)\omega\sin(a-3)\omega}, \end{aligned} \]

\origpage{230}
et en général,
\[ \left(\frac{n-a+k}{a}\right)=\frac{1}{k}\cdot\frac{A_{1}A_{2}\ldots A_{k-1}}{A_{a-1}A_{a-2}\ldots A_{a-k}}\cdot\frac{\omega\sin\omega\sin 2\omega\ldots\sin k\omega}{\sin a\omega\sin(a+1)\omega\ldots\sin(a-k)\omega}\ldots(n). \]
\ednote{The denominator is printed $\sin(a+1)\omega$; $\sin(a-1)\omega$ is meant, as in the two preceding formulas.}
Cette formule donne la valeur de la fonction $\left(\frac{p}{q}\right)$ lorsque $p+q$ surpasse $n$; ainsi en la réunissant à la formule (k), on a généralement l'expression de toute fonction $\left(\frac{p}{q}\right)$, en supposant seulement connues les fonctions semblables pour lesquelles $p+q=n-1$, et on a déja remarqué que le nombre de ces auxiliaires est $\frac{n+2}{2}$\ednote{Printed $\frac{n+2}{2}$; $\frac{n-2}{2}$ is meant, as stated in article (5).} ou $\frac{n-1}{2}$, selon que $n$ est pair ou impair.

\subsection*{(8).}

La formule (n) pourrait être regardée comme une suite de la formule (k); et on n'aurait ainsi qu'une seule et même formule pour toutes les valeurs de $p$ et $q$; mais alors on aurait besoin de nouvelles auxiliaires $A_{0}$, $A_{-1}$, $A_{-2}$, etc., et il faudrait en fixer les valeurs ainsi: $A_{0}=\frac{\omega}{\sin 0\omega}$, ou plutôt $i$ étant infiniment petit, $A_{i}=\frac{\omega}{\sin i\omega}=\frac{1}{i}$, $A_{-1}=-A_{1}$, $A_{-2}=-A_{2}$, etc. C'est pour éviter l'embarras de ces substitutions, surtout dans le cas de $A_{0}$, que nous avons donné les deux formules séparément.

Il est assez étonnant que l'expression générale des fonctions $\left(\frac{p}{q}\right)$ ait échappé à Euler; on voit cependant qu'il s'était occupé spécialement de cette recherche, par le passage du tome V des \emph{Nova acta Petropol.}, pag.\ 125, où il dit: \emph{Neque tamen hinc adhuc elucet quânam lege omnes determinationes progrediantur, quandoquidem valores certarum formularum continuò magis evadunt complicati.}

Nous remarquerons au reste que les formules (k) et (n) qui contiennent l'expression générale dont il s'agit, peuvent être regardées comme l'intégrale complète de l'équation aux différences finies (e); de sorte qu'on ne peut tirer de cette équation aucune conséquence qui ne soit contenue dans les formules (k) et (n). C'est ce qu'il serait facile de démontrer par les méthodes que l'on suit dans ce genre d'analyse, et qui, pour la plupart, ont été indiquées

\origpage{231}
par Lagrange dans les \emph{Mémoires de l'Académie de Berlin}, an.\ 1775.

\subsection*{(9).}

Voici maintenant quelques formules particulières qui méritent d'être citées. De l'équation (e) on déduit les deux suivantes,
\[ \begin{gathered} \left(\frac{p}{q}\right)\cdot\left(\frac{p+q}{n-p-q}\right)=\left(\frac{p}{n-p-q}\right)\cdot\left(\frac{n-q}{q}\right),\\ \left(\frac{p}{n-p-q}\right)\cdot\left(\frac{n-q}{n-p}\right)=\left(\frac{p}{n-p}\right)\cdot\left(\frac{n}{n-p-q}\right); \end{gathered} \]
multipliant ces deux équations entre elles, et mettant dans le produit les valeurs connues par les formules (c) et (d), on aura
\[ \left(\frac{p}{q}\right)\cdot\left(\frac{n-p}{n-q}\right)=\frac{\omega\sin(p+q)\omega}{(n-p-q)\sin p\omega\sin q\omega}\ldots\ldots(p); \]
d'où il suit que la valeur de $\left(\frac{n-p}{n-q}\right)$ se déduit immédiatement de celle de $\left(\frac{p}{q}\right)$, qui est en quelque sorte son complément. On a en particulier,
\[ \left(\frac{a}{a}\right)\cdot\left(\frac{n-a}{n-a}\right)=\frac{2\omega\cot a\omega}{n-2a}\ldots\ldots(q). \]
Ainsi connaissant les valeurs de $\left(\frac{a}{a}\right)$ lorsque $a$ n'excède pas $\frac{1}{2}n$, on en déduit les valeurs de $\left(\frac{a}{a}\right)$ lorsque $a$ est plus grand que $\frac{1}{2}n$.

\subsection*{(10).}

Pour examiner plus particulièrement les fonctions de la forme $\left(\frac{a}{a}\right)$, reprenons la valeur primitive de ces fonctions, laquelle est
\[ \left(\frac{a}{a}\right)=\int\frac{x^{a-1}\,dx}{\sqrt[n]{(1-x^{n})^{n-a}}}; \]
si l'on fait
\[ 1-x^{n}=\frac{z^{n}}{4x^{n}},\quad\text{ou}\quad x^{n}=\tfrac{1}{2}\pm\tfrac{1}{2}\sqrt{(1-z^{n})}, \]
la transformée sera
\[ \left(\frac{a}{a}\right)=\mp 2^{-\frac{2a}{n}}\int\frac{z^{a-1}\,dz}{\sqrt{(1-z^{n})}}. \]

\origpage{232}
Quant aux limites de cette nouvelle intégrale, il faut observer que les valeurs $x^{n}=0$, $x^{n}=\frac{1}{2}$, $x^{n}=1$ donnent respectivement $z=0$, $z=1$, $z=0$. D'où l'on voit qu'il faut prendre deux fois l'intégrale en $z$, depuis $z=0$ jusqu'à $z=1$. Et comme alors rien n'empêche de mettre $x$ à la place de $z$, on aura
\[ \left(\frac{a}{a}\right)=2^{1-\frac{2a}{n}}\int\frac{x^{a-1}\,dx}{\sqrt{(1-x^{n})}}\ldots\ldots(r); \]
cette intégrale est ainsi réduite à la forme la plus simple dont elle soit susceptible, puisque le radical n'est plus que du second degré.

\subsection*{(11).}

Si dans cette formule on met $n-a$ à la place de $a$, on aura
\[ \left(\frac{n-a}{n-a}\right)=2^{-1+\frac{2a}{n}}\int\frac{x^{n-a-1}\,dx}{\sqrt{(1-x^{n})}}; \]
de là et de l'équation (q) résulte cette formule remarquable
\[ \int\frac{x^{a-1}\,dx}{\sqrt{(1-x^{n})}}\cdot\int\frac{x^{n-a-1}\,dx}{\sqrt{(1-x^{n})}}=\frac{2\omega\cot a\omega}{n-2a}\ldots\ldots(s). \]

\subsection*{(12).}

Puisque les fonctions $\left(\frac{a}{a}\right)$ sont les plus simples entre les fonctions $\left(\frac{p}{q}\right)$ non comprises dans les formules (c) et (l), il semblerait convenable de les substituer aux auxiliaires désignées par $A_{a}$, pour exprimer par leur moyen toutes les fonctions $\left(\frac{p}{q}\right)$.

Dans cette vue, désignons en général la fonction $\left(\frac{a}{a}\right)$ par $M_{a}$; comme on peut supposer $a<\frac{1}{2}n$, on aura par la formule (k),
\[ M_{a}=\frac{A_{a}A_{a+1}\ldots A_{n-a-1}}{A_{1}A_{2}\ldots A_{n-2a-1}}\cdot\frac{\sin(a+1)\omega\sin(a+2)\omega\ldots\sin(n-a-1)\omega}{\sin\omega\sin 2\omega\ldots\sin(n-2a-1)\omega}, \]
valeur qui, au moyen des équations $A_{n-1-k}=A_{k}$\ednote{The first letter of this subscript is badly printed; it is read as $n$, following equation (g).}, $\sin(n-k)\omega=\sin k\omega$, se réduit à cette forme
\[ M_{a}=\frac{A_{a}A_{a+1}\ldots A_{2a-1}}{A_{1}A_{2}\ldots A_{a-1}}\cdot\frac{\sin(a+1)\omega\sin(a+2)\omega\ldots\sin 2a\omega}{\sin\omega\sin 2\omega\ldots\sin a\omega}\ldots\ldots(t); \]

\origpage{233}
d'où l'on tire successivement
\[ \begin{aligned} M_{1}&=A_{1}\frac{\sin 2\omega}{\sin\omega}\\ M_{2}&=\frac{A_{2}A_{3}}{A_{1}}\cdot\frac{\sin 3\omega\sin 4\omega}{\sin\omega\sin 2\omega}\\ M_{3}&=\frac{A_{3}A_{4}A_{5}}{A_{1}A_{2}}\cdot\frac{\sin 4\omega\sin 5\omega\sin 6\omega}{\sin\omega\sin 2\omega\sin 3\omega}\\ &\text{etc.} \end{aligned} \]
Ces équations, qu'on peut mettre aussi sous la forme
\[ \left.\begin{aligned} A_{1}&=M_{1}\cdot\frac{\sin\omega}{\sin 2\omega}\\ \frac{A_{2}A_{3}}{A_{1}A_{1}}&=\frac{M_{2}}{M_{1}}\cdot\frac{\sin^{2}2\omega}{\sin 3\omega\sin 4\omega}\\ \frac{A_{4}A_{5}}{A_{2}A_{2}}&=\frac{M_{3}}{M_{2}}\cdot\frac{\sin^{2}3\omega}{\sin 5\omega\sin 6\omega}\\ &\text{etc.} \end{aligned}\right\}\qquad(u), \]
serviront à déterminer les auxiliaires $A_{1}$, $A_{2}$, $A_{3}$, etc.\ au moyen d'un égal nombre de quantités $M_{1}$, $M_{2}$, $M_{3}$, etc., prises dans l'ordre convenable. On pourra donc exprimer par ces dernières quantités toutes les fonctions $\left(\frac{p}{q}\right)$ qui répondent à une même valeur de $n$.

Mais il faut observer que ces substitutions ne peuvent s'effectuer que pour des valeurs particulières de $n$, et qu'ainsi par l'emploi des auxiliaires $M_{a}$, on ne peut parvenir à des formules aussi générales que le sont les formules (k) et (n).

\subsection*{(13).}

Considérons maintenant la formule
\[ \left(\frac{n-2a}{a}\right)=\int\frac{x^{n-2a-1}\,dx}{\sqrt[n]{(1-x^{n})^{n-a}}}: \]
si on fait $x^{-n}=\frac{1}{2}+\frac{1}{2}\sqrt{(1-z^{n})}$\ednote{Printed $1-z^{n}$ under the radical; the transformée (v) requires $x^{-n}=\frac{1}{2}+\frac{1}{2}\sqrt{(1+z^{n})}$.}, on aura la transformée
\[ \left(\frac{n-2a}{a}\right)=2^{-\frac{2a}{n}}\int\frac{z^{a-1}\,dz}{\sqrt{(1+z^{n})}}\qquad(v), \]
dans laquelle il faudra prendre l'intégrale depuis $z=0$ jusqu'à $z=\infty$, et qui d'ailleurs suppose $a<\frac{1}{2}n$. Mais en vertu des équations

\origpage{234}
(e) et (p), on a
\[ \begin{gathered} \left(\frac{n-2a}{a}\right)\cdot\left(\frac{2a}{n-a}\right)=\left(\frac{\omega}{a\sin 2a\omega}\right),\\ \left(\frac{a}{a}\right)\cdot\left(\frac{2a}{n-a}\right)=\left(\frac{a}{n-a}\right)\cdot\left(\frac{n}{a}\right)=\frac{\omega}{a\sin a\omega}; \end{gathered} \]
\ednote{In the first line the value $\frac{\omega}{a\sin 2a\omega}$ is printed inside the parentheses of the transcendent symbol; the plain quotient is meant, as in the second line.}
donc $\left(\frac{a}{a}\right)=2\cos a\omega\cdot\left(\frac{n-2a}{a}\right)$, ou, en substituant la valeur donnée par l'équation (v),
\[ \left(\frac{a}{a}\right)=2^{1-\frac{2a}{n}}\cos a\omega\cdot\int\frac{z^{a-1}\,dz}{\sqrt{(1+z^{n})}}\qquad(x). \]
Cette formule n'a lieu que lorsque $a$ est $<\frac{1}{2}n$; si $a$ est $>\frac{1}{2}n$, on commencera par prendre la valeur de $\left(\frac{n-a}{n-a}\right)$, laquelle sera
\[ \left(\frac{n-a}{n-a}\right)=2^{\frac{2a}{n}-1}\cos(n-a)\omega\cdot\int\frac{z^{n-a-1}\,dz}{\sqrt{(1+z^{n})}}, \]
et on en déduira celle de $\left(\frac{a}{a}\right)$ au moyen de l'équation (q). On aura ainsi, $a$ étant $>\frac{1}{2}n$:
\[ \left(\frac{a}{a}\right)=\frac{2^{2-\frac{2a}{n}}}{2a-n}\cdot\frac{\omega}{\sin a\omega}:\int\frac{z^{n-a-1}\,dz}{\sqrt{(1+z^{n})}}\qquad(y). \]

\subsection*{(14).}

Si l'on compare maintenant les équations (r), (x) et (y), on en tirera les formules
\[ \left.\begin{aligned} \int\frac{x^{a-1}\,dx}{\sqrt{(1-x^{n})}}&=\cos a\omega\cdot\int\frac{z^{a-1}\,dz}{\sqrt{(1+z^{n})}}\\ \int\frac{x^{a-1}\,dx}{\sqrt{(1-x^{n})}}&=\frac{2\omega}{(2a-n)\sin a\omega}:\int\frac{z^{n-a-1}\,dz}{\sqrt{(1+z^{n})}} \end{aligned}\right\}\qquad(z), \]
la première ayant lieu lorsque $a$ est $<\frac{1}{2}n$, et la seconde lorsque $a$ est $>\frac{1}{2}n$.

Lorsque $n$ est impair, si on fait dans la première équation, $x=1-y^{2}$ et $z=y^{2}-1$, les formules intégrales comprises dans les deux membres se réduiront l'une et l'autre à la forme
\[ \int\frac{(1-y^{2})^{a-1}\,dy}{\sqrt{\left(n-\frac{n\cdot n-1}{1\cdot 2}y^{2}+\frac{n\cdot n-1\cdot n-2}{1\cdot 2\cdot 3}y^{4}-\text{etc.}\right)}}: \]

\origpage{235}
on voit donc que la partie de cette dernière intégrale, prise depuis $y=0$ jusqu'à $y=1$, et la partie prise depuis $y=1$ jusqu'à $y=\infty$, sont entre elles $::\cos\frac{a\pi}{n}:1$.

\subsection*{(15).}

Considérons de nouveau la formule
\[ \left(\frac{p}{q}\right)=\int\frac{x^{p-1}\,dx}{\sqrt[n]{(1-x^{n})^{n-q}}}, \]
si l'on fait
\[ 1-x^{n}=\frac{x^{2n}}{4z^{n}}, \]
on obtient d'abord
\[ \left(\frac{p}{q}\right)=2^{2-\frac{2q}{n}}\int x^{p+2q-2n-1}\,dx\cdot z^{n-q}. \]
Soit $p=2a$ et $q=n-a$, on aura
\[ \left(\frac{2a}{n-a}\right)=2^{\frac{2a}{n}}\int z^{a}\cdot\frac{dx}{x}, \]
et en achevant les substitutions, il viendra
\[ \int\left(z^{a-1}\,dz-\frac{z^{a+\frac{1}{2}n-1}\,dz}{\sqrt{(1+z^{n})}}\right)=2^{1-\frac{2a}{n}}\cdot\left(\frac{2a}{n-a}\right); \]
cela posé, il faut distinguer deux cas, selon que $a$ est $<\frac{1}{2}n$ ou $>\frac{1}{2}n$.

Soit, 1°. $a<\frac{1}{2}n$, on aura par les formules du n° 13,
\[ \left(\frac{2a}{n-a}\right)=\frac{\omega}{a\sin 2a\omega}:\left(\frac{n-2a}{a}\right)=2^{\frac{2a}{n}}\cdot\frac{\omega}{a\sin 2a\omega}:\int\frac{z^{a-1}\,dz}{\sqrt{(1+z^{n})}}, \]
donc
\[ \int\frac{z^{a-1}\,dz}{\sqrt{(1+z^{n})}}\cdot\int\left(z^{a-1}\,dz-\frac{z^{a+\frac{1}{2}n-1}\,dz}{\sqrt{(1+z^{n})}}\right)=\frac{2\omega}{a\sin 2a\omega}\qquad(a'). \]
Soit, 2°. $a>\frac{1}{2}n$, on aura par l'équation (b),
\[ \left(\frac{2a}{n-a}\right)=\frac{2a-n}{a}\cdot\left(\frac{2a-n}{n-a}\right); \]
mais en faisant $a=n-c$, on a
\[ \left(\frac{2a-n}{n-a}\right)=\left(\frac{n-2c}{c}\right), \]

\origpage{236}
et par l'équation (v), on a
\[ \frac{n-2c}{c}=2^{-\frac{2c}{n}}\int\frac{z^{c-1}\,dz}{\sqrt{(1+z^{n})}}. \]
\ednote{The left side is printed as a plain fraction without the parentheses of the transcendent symbol; $\left(\frac{n-2c}{c}\right)$ is meant, as in the preceding line.}
Donc au lieu de l'équation (a'), on aura
\[ \int\left(z^{a-1}\,dz-\frac{z^{a+\frac{1}{2}n-1}\,dz}{\sqrt{(1+z^{n})}}\right)=\frac{2a-n}{2a}\int\frac{z^{n-a-1}\,dz}{\sqrt{(1+z^{n})}}\qquad(b'). \]
Au reste cette dernière équation se vérifie immédiatement au moyen de la fonction $P=z^{a-\frac{1}{2}n}\sqrt{(1+z^{n})}-z^{a}$, qui s'évanouit dans les deux limites, lorsque $z=0$ et lorsque $z=\infty$; car si on prend la différentielle de cette fonction, et qu'ensuite on l'intègre, on trouvera
\[ \int\left(z^{a-1}\,dz-\frac{z^{a+\frac{1}{2}n-1}\,dz}{\sqrt{(1+z^{n})}}\right)=\frac{2a-n}{2a}\cdot\int\frac{z^{a-\frac{n}{2}-1}\,dz}{\sqrt{(1+z^{n})}}; \]
formule qui ne diffère pas de la précédente, parce qu'en mettant $\frac{1}{z}$ au lieu de $z$, l'intégrale $\displaystyle\int\frac{z^{a-\frac{1}{2}n-1}\,dz}{\sqrt{(1+z^{n})}}$ se change en $\displaystyle\int\frac{z^{n-a-1}\,dz}{\sqrt{(1+z^{n})}}$, les limites étant toujours $z=0$, $z=\infty$.

\subsection*{(16).}

Considérons enfin, dans la supposition de $n$ pair, la formule
\[ \left(\frac{\frac{1}{2}n-a}{a}\right)=\int\frac{x^{\frac{1}{2}n-a-1}\,dx}{\sqrt[n]{(1-x^{n})^{n-a}}}; \]
si on fait $x^{-n}=1+z^{n}$, on aura par la substitution,
\[ \left(\frac{\frac{1}{2}n-a}{a}\right)=\int\frac{z^{a-1}\,dz}{\sqrt{(1+z^{n})}},\qquad(c) \]
\ednote{The label is printed (c) without a prime; the text below cites this equation as (c').}
et l'intégrale du second membre devra être prise entre les limites $z=0$, $z=\infty$.

Maintenant, par la combinaison des équations (x) et (c'), on obtient
\[ \left(\frac{a}{a}\right)=2^{1-\frac{2a}{n}}\cos a\omega\left(\frac{\frac{1}{2}n-a}{a}\right), \]
et par conséquent aussi,
\[ \left(\frac{\frac{1}{2}n-a}{\frac{1}{2}n-a}\right)=2^{\frac{2a}{n}}\sin a\omega\left(\frac{a}{\frac{1}{2}n-a}\right). \]

\origpage{237}
De ces deux-ci on conclura
\[ \left(\frac{a}{a}\right)=2^{1-\frac{4a}{n}}\operatorname{cotang}a\omega\cdot\left(\frac{\frac{1}{2}n-a}{\frac{1}{2}n-a}\right)\qquad(d'), \]
d'où l'on voit que $n$ étant pair, il suffit d'avoir les valeurs de $\left(\frac{a}{a}\right)$ pour tous les cas où $a$ ne surpasse pas $\frac{1}{4}n$, et qu'ainsi le nombre des auxiliaires nécessaires pour déterminer toutes les fonctions $\left(\frac{p}{q}\right)$ qui répondent à une même valeur de $n$, se réduit à $\frac{n}{4}$ ou $\frac{n-2}{4}$, selon que $n$ est de la forme $4m$ ou $4m+2$.

\subsection*{(17).}

A l'aide de l'équation (d'), on trouvera des relations entre les auxiliaires $A_{1}$, $A_{2}$, $A_{3}$, etc.\ qui réduiront leur nombre comme il vient d'être dit.

Pour cela reprenons la formule (t),
\[ \left(\frac{a}{a}\right)=\frac{A_{a}A_{a+1}\ldots A_{2a-1}}{A_{1}A_{2}\ldots\ldots A_{a-1}}\cdot\frac{\sin(a+1)\omega\sin(a+2)\omega\ldots\sin 2a\omega}{\sin\omega\sin 2\omega\ldots\sin a\omega}, \]
elle donne, en faisant $n=2m$,
\[ \left(\frac{m-a}{m-a}\right)=\frac{A_{m-a}\cdot A_{m-a+1}\ldots A_{2m-2a-1}}{A_{1}A_{2}\ldots A_{m-a-1}}\cdot\frac{\sin(m-a+1)\omega\sin(m-a+2)\omega\ldots\sin(2m-2a)\omega}{\sin\omega\sin 2\omega\ldots\sin(m-a)\omega}; \]
substituant ces valeurs dans l'équation (d'), et faisant les réductions dans l'hypothèse $a<\frac{1}{2}m$, on aura généralement
\[ \frac{A_{a}A_{a+1}\ldots A_{2a-1}}{A_{m-1}A_{m-2}\ldots A_{m-a}}=2^{-\frac{a}{m}}\cdot\frac{\sin(m-1)\omega\sin(m-2)\omega\ldots\sin(m-a+1)\omega}{\sin a\omega\sin(a+1)\omega\ldots\sin(2a-1)\omega}\qquad(e'). \]
De là résultent, en faisant successivement $a=1$, $2$, $3$, etc., des équations particulières qui peuvent être mises sous cette forme
\[ \left.\begin{aligned} A_{m-1}&=A_{1}\cdot 2^{\frac{1}{m}}\sin\omega\\ A_{m-2}&=\frac{A_{2}A_{3}}{A_{1}}\cdot 2^{1+\frac{1}{m}}\sin 3\omega\\ A_{m-3}&=\frac{A_{4}A_{5}}{A_{2}}\cdot 2^{1+\frac{1}{m}}\sin 5\omega\\ A_{m-4}&=\frac{A_{6}A_{7}}{A_{3}}\cdot 2^{1+\frac{1}{m}}\sin 7\omega\\ &\text{etc.} \end{aligned}\right\}\qquad(f'). \]

\origpage{238}
Elles devront être continuées jusqu'à ce que le nombre en soit $\frac{n-4}{4}$ ou $\frac{n-2}{4}$.

\subsection*{(18).}

Par exemple, lorsque $n=12$, il y a cinq auxiliaires $A_{1}=\left(\frac{10}{1}\right)$, $A_{2}=\left(\frac{9}{2}\right)$, $A_{3}=\left(\frac{8}{3}\right)$, $A_{4}=\left(\frac{7}{4}\right)$, $A_{5}=\left(\frac{6}{5}\right)$, entre lesquelles on a ces deux équations
\[ \begin{aligned} A_{5}&=A_{1}\cdot 2^{\frac{1}{6}}\sin\omega\\ A_{4}&=\frac{A_{2}A_{3}}{A_{1}}\cdot 2^{\frac{7}{6}}\sin 3\omega. \end{aligned} \]
De sorte que le nombre d'auxiliaires nécessaires se réduit à trois, pour lesquelles on peut prendre $A_{1}$, $A_{2}$, $A_{3}$.

Si on préférait de prendre pour auxiliaires trois des quantités $M_{a}$, il faudrait avoir recours aux équations (u), lesquelles donnent
\[ \begin{aligned} A_{1}&=M_{1}\cdot\frac{\sin\omega}{\sin 2\omega}\\ \frac{A_{2}A_{3}}{A_{1}A_{1}}&=\frac{M_{2}}{M_{1}}\cdot\frac{\sin^{2}.2\omega}{\sin 3\omega\sin 4\omega}\\ \frac{A_{4}A_{5}}{A_{2}A_{2}}&=\frac{M_{3}}{M_{2}}\cdot\frac{\sin^{2}3\omega}{\sin 5\omega\sin 6\omega}. \end{aligned} \]
Au moyen de ces équations et des précédentes, observant d'ailleurs qu'on a $\omega=\frac{\pi}{12}$, on trouve
\[ \begin{aligned} A_{1}&=2\sin\omega\cdot M_{1}\\ A_{2}&=2^{\frac{1}{6}}\sin\omega\cdot M_{2}\sqrt{\left(\frac{M_{1}}{M_{3}\cos 2\omega}\right)}\\ A_{3}&=2^{\frac{1}{3}}\sin\omega\cdot M_{3}\sqrt{\left(\frac{M_{1}}{M_{3}\cos 2\omega}\right)}\\ A_{4}&=2^{\frac{1}{6}}\sin\omega\cdot M_{2}\cdot\frac{1}{\cos 2\omega}\\ A_{5}&=2^{\frac{7}{6}}\sin^{2}\omega\cdot M_{1}. \end{aligned} \]

\subsection*{(19).}

Appliquons ici les résultats que nous avons trouvés dans la première partie, n° 155, et supposons connues les deux fonctions complètes de première espèce $F^{1}(\sin 45^{\circ})=B$, $F^{1}(\sin 15^{\circ})=C$. Si l'on fait comme dans l'article cité, $\nu=\sqrt[4]{3}$, on aura

\origpage{239}
\[ \begin{aligned} M_{1}&=\frac{2^{\frac{4}{3}}\nu\cos\omega}{3}B\\ M_{2}&=\frac{2^{-\frac{1}{3}}}{\nu}C\\ M_{3}&=\tfrac{1}{3}B. \end{aligned} \]
Au moyen de ces valeurs, on trouvera
\[ \begin{aligned} A_{1}&=\frac{2^{\frac{1}{3}}\nu}{3}\cdot B\\ A_{2}&=\sqrt{\left(\frac{\nu\sin\omega}{3}\right)}\cdot C\\ A_{3}&=\frac{\sqrt{(2\nu\sin\omega)}}{3\nu}\cdot B\\ A_{4}&=\frac{2^{\frac{5}{6}}\nu\sin\omega}{3}\cdot C\\ A_{5}&=\frac{2^{\frac{1}{2}}\nu\sin\omega}{3}\cdot B. \end{aligned} \]
Ainsi par les deux seules données $B$ et $C$, on pourra déterminer toutes les transcendantes désignées par $\left(\frac{p}{q}\right)$, dans le cas de $n=12$. Ces transcendantes, en excluant celles qui sont déterminables par arcs de cercle, sont au nombre de 60, toutes différentes les unes des autres. On voit de plus qu'elles se détermineront rationnellement en fonctions de $B$ et $C$ et du nombre $\pi$, puisqu'elles s'expriment rationnellement par les auxiliaires $A_{1}$, $A_{2}$, etc.

Si on veut, par exemple, avoir la valeur de la transcendante $\left(\frac{5}{9}\right)$, on cherchera d'abord par l'équation (n) sa valeur en fonction de $A$, laquelle est
\[ \left(\frac{5}{9}\right)=\tfrac{1}{2}\cdot\frac{A_{1}}{A_{3}A_{4}}\cdot\frac{\sin\omega\sin 2\omega}{\sin 5\omega\sin 4\omega\sin 3\omega}; \]
\ednote{The numerator is printed $\sin\omega\sin 2\omega$; formula (n) gives $\omega\sin\omega\sin 2\omega$, and only with the factor $\omega$ does the equation hold numerically (and lead to the final value below).}
faisant ensuite les substitutions et observant qu'on a $\omega=\frac{\pi}{12}$, il viendra
\[ \left(\frac{5}{9}\right)=\frac{\pi}{2C}\sqrt{\left(\frac{\nu\sin\omega}{6}\right)}. \]

\origpage{240}

\section*{Formules pour évaluer par approximation les intégrales $\left(\frac{p}{q}\right)$.}

\subsection*{(20).}

Ayant réduit au moindre nombre possible les transcendantes $\left(\frac{p}{q}\right)$, il ne reste plus qu'à faire voir comment on peut trouver par approximation, et d'une manière facile, la valeur de chacune de ces quantités.

Pour cet effet, considérons d'abord la formule
\[ \left(\frac{a}{a}\right)=2^{1-\frac{2a}{n}}\int\frac{x^{a-1}\,dx}{\sqrt{(1-x^{n})}}, \]
et soit $\sqrt{(1-x^{n})}=1-y$, ou $x^{n}=2y-y^{2}$, on aura pour transformée,
\[ \left(\frac{a}{a}\right)=2^{1-\frac{a}{n}}\int\frac{dy}{n}\left(y-\frac{y^{2}}{2}\right)^{\frac{a}{n}-1}; \]
cette différentielle étant développée et intégrée depuis $y=0$ jusqu'à $y=1$, on obtient
\[ \left(\frac{a}{a}\right)=2^{1-\frac{a}{n}}\left\{\begin{aligned} &1+\frac{n-a}{2n}\cdot\frac{a}{n+a}+\frac{n-a\cdot 2n-a}{2n\cdot 4n}\cdot\frac{a}{2n+a}\\ &\qquad+\frac{n-a\cdot 2n-a\cdot 3n-a}{2n\cdot 4n\cdot 6n}\cdot\frac{a}{3n+a}\\ &+\text{etc.} \end{aligned}\right\}:a\qquad(g'), \]
formule dont chaque terme est moindre que la moitié du précédent.

\subsection*{(21).}

En général si on veut avoir la valeur approchée de la quantité $\left(\frac{p}{q}\right)=\displaystyle\int\frac{x^{p-1}\,dx}{\sqrt[n]{(1-x^{n})^{n-q}}}$, il faut partager cette intégrale en deux parties, l'une depuis $x^{n}=0$ jusqu'à $x^{n}=\frac{1}{2}$, l'autre depuis $x^{n}=\frac{1}{2}$ jusqu'à $x^{n}=1$.

La première partie étant nommée $P$, on trouve par les développemens ordinaires,
\[ P=2^{-\frac{p}{n}}\left(\frac{1}{p}+\frac{n-q}{2n}\cdot\frac{1}{n+p}+\frac{n-q\cdot 2n-q}{2n\cdot 4n}\cdot\frac{1}{2n+p}+\text{etc.}\right). \]

\origpage{241}
Pour avoir la seconde partie que nous nommerons $Q$, il faut faire $x^{n}=1-y^{n}$, alors on aura
\[ \int\frac{x^{p-1}\,dx}{\sqrt[n]{(1-x^{n})^{n-q}}}=-\int\frac{y^{q-1}\,dy}{\sqrt[n]{(1-x^{n})^{n-q}}}, \]
\ednote{The second denominator is printed $\sqrt[n]{(1-x^{n})^{n-q}}$; the substitution gives $\sqrt[n]{(1-y^{n})^{n-p}}$, as the series for $Q$ below requires. An apparent misprint.}
et la transformée en $y$ devra être intégrée depuis $y^{n}=\frac{1}{2}$ jusqu'à $y^{n}=0$. Si on change son signe, elle devra être intégrée depuis $y^{n}=0$ jusqu'à $y^{n}=\frac{1}{2}$; on aura donc
\[ Q=2^{-\frac{q}{n}}\left(\frac{1}{q}+\frac{n-p}{2n}\cdot\frac{1}{n+q}+\frac{n-p\cdot 2n-p}{2n\cdot 4n}\cdot\frac{1}{2n+q}+\text{etc.}\right). \]
Il ne s'agit plus que de réunir ces deux parties, et on obtient
\[ \left(\frac{p}{q}\right)=\left\{\begin{aligned} &2^{-\frac{p}{n}}\left(\frac{1}{p}+\frac{n-q}{2n}\cdot\frac{1}{n+p}+\frac{n-q\cdot 2n-q}{2n\cdot 4n}\cdot\frac{1}{2n+p}+\frac{n-q\cdot 2n-q\cdot 3n-q}{2n\cdot 4n\cdot 6n}\cdot\frac{1}{3n+p}+\text{etc.}\right)\\ &+2^{-\frac{q}{n}}\left(\frac{1}{q}+\frac{n-p}{2n}\cdot\frac{1}{n+q}+\frac{n-p\cdot 2n-p}{2n\cdot 4n}\cdot\frac{1}{2n+q}+\frac{n-p\cdot 2n-p\cdot 3n-p}{2n\cdot 4n\cdot 6n}\cdot\frac{1}{3n+q}+\text{etc.}\right) \end{aligned}\right\}\qquad(h'). \]
Les deux séries comprises dans cette formule sont toujours convergentes, puisque chaque terme est moindre que la moitié du précédent: on obvie ainsi à l'inconvénient que présenterait la méthode ordinaire, si on voulait intégrer tout d'un coup la valeur de $\left(\frac{p}{q}\right)$; depuis $x=0$ jusqu'à $x=1$, ce qui donnerait la suite très-peu convergente
\[ \left(\frac{p}{q}\right)=\frac{1}{p}+\frac{n-q}{n}\cdot\frac{1}{n+p}+\frac{n-q\cdot 2n-q}{n\cdot 2n}\cdot\frac{1}{2n+p}+\text{etc.}\qquad(i'). \]
Au reste, lorsqu'on suppose $p=q=a$, la formule (h') se réduit précisément à la formule (g') trouvée par une autre voie.

\subsection*{(22).}

Il ne sera pas inutile de chercher la valeur de la fonction $\left(\frac{p}{q}\right)$, dans le cas où $n$ est très-grand par rapport aux nombres $p$ et $q$. Pour cela, soit $p=\alpha n$, $q=\beta n$, on pourra considérer $\alpha$ et $\beta$ comme des quantités très-petites du premier ordre, et il faudra développer jusqu'au degré convenable les suites $P$ et $Q$; on a d'abord,

\origpage{242}
en faisant $\frac{1}{2}=x$;
\[ P=\frac{2^{-\alpha}}{n\alpha}\left\{\begin{aligned} &1+\frac{1-\beta}{1}\cdot\frac{\alpha x}{1+\alpha}+1-\beta\cdot 1-\tfrac{1}{2}\beta\cdot\frac{\alpha x^{2}}{2+\alpha}\\ &\qquad+1-\beta\cdot 1-\tfrac{1}{2}\beta\cdot 1-\tfrac{1}{3}\beta\cdot\frac{\alpha x^{3}}{3+\alpha}\\ &+\text{etc.} \end{aligned}\right\}. \]

Développant la série renfermée dans la parenthèse, jusqu'aux quantités du second ordre inclusivement, et faisant pour abréger,
\[ \begin{aligned} A&=x+\frac{x^{2}}{2}+\frac{x^{3}}{3}+\frac{x^{4}}{4}+\text{etc.}\\ B&=x+\left(1+\tfrac{1}{2}\right)\frac{x^{2}}{2}+\left(1+\tfrac{1}{2}+\tfrac{1}{3}\right)\frac{x^{3}}{3}+\text{etc.}\\ C&=x+\frac{x^{2}}{2^{2}}+\frac{x^{3}}{3^{2}}+\frac{x^{4}}{4^{2}}+\text{etc.} \end{aligned} \]
on aura
\[ P=\frac{2^{-\alpha}}{n\alpha}\left[1+(A-B\beta)\alpha-C\alpha^{2}\right]. \]
Mais on a aussi $2^{-\alpha}=1-\alpha l2+\frac{1}{2}\alpha^{2}(l2)^{2}$, etc.; et d'ailleurs $A=-\log(1-x)=-\log\left(1-\frac{1}{2}\right)=\log 2$, de sorte qu'on peut supposer $2^{-\alpha}=1-A\alpha+\frac{1}{2}A^{2}\alpha^{2}$; ainsi en effectuant les développemens, on aura
\[ P=\frac{1}{n\alpha}\left[1-B\beta\alpha-\left(C+\tfrac{1}{2}A^{2}\right)\alpha^{2}\right]. \]
On aura semblablement
\[ Q=\frac{1}{n\beta}\left[1-B\alpha\beta-\left(C+\tfrac{1}{2}A^{2}\right)\beta^{2}\right]; \]
donc par la somme de ces quantités on trouve
\[ \left(\frac{p}{q}\right)=\frac{1}{n}\left[\frac{1}{\alpha}+\frac{1}{\beta}-\left(B+C+\tfrac{1}{2}A^{2}\right)(\alpha+\beta)\right], \]
ou
\[ \left(\frac{p}{q}\right)=\frac{p+q}{pq}\left[1-\left(B+C+\tfrac{1}{2}A^{2}\right)\frac{pq}{n^{2}}\right]. \]

\subsection*{(23).}

Il reste à trouver la valeur de $B+C$, et pour cela nous considérerons pour un moment $B$ et $C$ comme des fonctions d'une

\origpage{243}
variable $x$; nous aurons d'abord
\[ \begin{aligned} B-C&=\frac{x^{2}}{2}+\left(1+\tfrac{1}{2}\right)\frac{x^{3}}{3}+\left(1+\tfrac{1}{2}+\tfrac{1}{3}\right)\frac{x^{4}}{4}+\text{etc.}\\ \frac{d(B-C)}{dx}&=x+\left(1+\tfrac{1}{2}\right)x^{2}+\left(1+\tfrac{1}{2}+\tfrac{1}{3}\right)x^{3}+\text{etc.}\\ &=\frac{x}{1-x}+\frac{\frac{1}{2}x^{2}}{1-x}+\frac{\frac{1}{3}x^{3}}{1-x}+\text{etc.}=-\frac{\log(1-x)}{1-x}, \end{aligned} \]
donc $d(B-C)=-\frac{dx}{1-x}\log(1-x)$, et en intégrant on a
\[ B-C=\tfrac{1}{2}\log^{2}(1-x), \]
sans constante, parce que $B$ et $C$ s'évanouissent en même temps que $x$.

On a ensuite
\[ \frac{x\,dC}{dx}=x+\frac{x^{2}}{2}+\frac{x^{3}}{3}+\text{etc.}=-\log(1-x); \]
donc
\[ dC=-\frac{dx}{x}\log(1-x), \]
et en intégrant
\[ C=-\log x\log(1-x)-\int\frac{dx}{1-x}\log x. \]
Soit $x=1-y$, on aura
\[ \int\frac{dx}{1-x}\log x=\int\frac{-dy}{y}\log(1-y)=\text{const}+y+\frac{y^{2}}{2^{2}}+\frac{y^{3}}{3^{2}}+\text{etc.} \]
Si donc $C$ est une fonction de $x$ désignée par $\psi(x)$, on aura
\[ \int\frac{dx}{1-x}\log x=\text{const.}+\psi(y)=\text{const}+\psi(1-x); \]
donc
\[ \psi(x)+\psi(1-x)=\text{const}-\log x\log(1-x). \]
Si on fait $x=0$, on trouve la constante $=\psi(1)=1+\frac{1}{2^{2}}+\frac{1}{3^{2}}+\text{etc.}$, quantité dont on sait que la valeur est $\frac{\pi^{2}}{6}$, de sorte qu'on aura
\[ \psi(x)+\psi(1-x)=\frac{\pi^{2}}{6}-\log x\log(1-x)\qquad(k'). \]

\origpage{244}
Cette formule fait voir qu'étant connue la somme de la suite
\[ \psi(x)=x+\frac{x^{2}}{2^{2}}+\frac{x^{3}}{3^{2}}+\frac{x^{4}}{4^{2}}+\text{etc.}, \]
pour toute valeur de $x$, depuis $x=0$ jusqu'à $x=\frac{1}{2}$, on connaîtra la somme de la même suite pour toute valeur de $x$, depuis $x=\frac{1}{2}$ jusqu'à $x=1$.

\subsection*{(24).}

Dans le cas particulier où l'on fait $x=\frac{1}{2}$, l'équation (k') donne
\[ \psi\left(\frac{1}{2}\right)=\frac{\pi^{2}}{12}-\tfrac{1}{2}(\log 2)^{2}; \]
donc dans ce même cas, on aura
\[ C=\frac{\pi^{2}}{12}-\tfrac{1}{2}A^{2}; \]
mais on a trouvé $B-C=\frac{1}{2}\log^{2}(1-x)=\frac{1}{2}\log^{2}2=\frac{1}{2}A^{2}$; donc
\[ B=C+\tfrac{1}{2}A^{2}=\frac{\pi^{2}}{12}, \]
et enfin
\[ \left(\frac{p}{q}\right)=\frac{p+q}{pq}\left(1-\frac{\pi^{2}}{6}\cdot\frac{pq}{n^{2}}\right)\qquad(l'). \]
C'est la limite vers laquelle tend continuellement la fonction $\left(\frac{p}{q}\right)$, lorsque $n$ augmente de plus en plus, $p$ et $q$ restant les mêmes.

C'est en même temps une valeur approchée de $\left(\frac{p}{q}\right)$, lorsque $p$ et $q$ sont petits par rapport à $n$. Soit, par exemple, $p=1$, $q=1$, $n=12$, on aura à très-peu près $\left(\frac{1}{1}\right)=2\left(1-\frac{\pi^{2}}{864}\right)$.

\section*{Remarque sur quelques cas particuliers, où l'on peut sommer la suite désignée par $\psi(x)$, et deux autres de la même espèce.}

\subsection*{(25).}

On a vu dans les recherches précédentes, que la suite représentée par $\psi(x)$, est sommable tant lorsque $x=1$ que lorsque

\origpage{245}
$x=\frac{1}{2}$. Landen, dans ses \emph{Mathématical Memoirs}, pag.~112 et suiv., a indiqué deux autres cas où cette même suite est sommable: nous allons faire voir comment on parvient à ces résultats.

Ayant fait $\psi(x)=x+\frac{x^{2}}{2^{2}}+\frac{x^{3}}{3^{2}}+\frac{x^{4}}{4^{2}}+\text{etc.}$, on en déduit
\[ d\psi(x)=dx\left(1+\frac{x}{2}+\frac{x^{2}}{3}+\text{etc.}\right)=-\frac{dx}{x}\log(1-x) \]
et $\psi x)$\ednote{Printed without the opening parenthesis; $\psi(x)$ is meant.} $=\displaystyle\int-\frac{dx}{x}\log(1-x)$. Soit $\log x=p$ et $\log(1-x)=q$, on aurait à la fois
\[ \psi(x)=\int-q\,dp,\qquad\psi(1-x)=\int-p\,dq; \]
donc $\psi(x)+\psi(1-x)=C-pq=\psi(1)-\log x\log(1-x)$.
C'est l'équation (k') que nous avons déjà trouvée; mais on peut en trouver plusieurs autres.

En effet on aurait semblablement $\psi(-x)=\displaystyle\int-\frac{dx}{x}\log(1+x)$; mettant $\frac{x}{1-x}$ au lieu de $x$, on aura
\[ \psi\left(-\frac{x}{1-x}\right)=\int\left(\frac{dx}{x}+\frac{dx}{1-x}\right)\log(1-x); \]
donc
\[ \psi(x)+\psi\left(\frac{-x}{1-x}\right)=\int\frac{dx}{1-x}\log(1-x)=-\tfrac{1}{2}\log^{2}(1-x). \]
On n'ajoute point de constante, parce que les deux membres s'évanouissent lorsque $x=0$.

Mettant de nouveau $\frac{x}{1+x}$ au lieu de $x$, on aura
\[ \psi\left(\frac{x}{x+1}\right)+\psi(-x)=-\tfrac{1}{2}\log^{2}(1+x). \]
Cette équation donne une valeur toujours finie pour $\psi(-x)$, quelque grand que soit $x$, puisque $\psi(-x)$ ne dépend que de $\psi\left(\frac{x}{x+1}\right)$, où $\frac{x}{x+1}$ est plus petit que l'unité. Mais comme dans le cas de $x>1$, la suite désignée par $\psi(x)$ sera divergente, on ne voit pas quelle peut être son utilité, contentons-nous donc de considérer

\origpage{246}
les valeurs de $\psi(-x)$ lorsque $x$ est ${}<1$. Or on a directement
\[ \psi(x)+\psi(-x)=\tfrac{1}{2}\left(x^{2}+\frac{x^{4}}{2^{2}}+\frac{x^{6}}{3^{2}}+\text{etc.}\right)=\tfrac{1}{2}\psi(x^{2}), \]
l'équation précédente peut donc se mettre sous la forme
\[ \psi\left(\frac{x}{x+1}\right)+\tfrac{1}{2}\psi(x^{2})-\psi(x)=-\tfrac{1}{2}\log^{2}(1+x). \]
Donnons à $x$ une valeur particulière, telle qu'on ait $x^{2}=\frac{x}{x+1}$, ou $x^{2}+x=1$; si on fait $\beta=-\frac{1}{2}+\frac{1}{2}\sqrt{5}$, on aura $x=\beta$, $\frac{x}{x+1}=\beta^{2}=1-\beta$. Donc
\[ \tfrac{3}{2}\psi(1-\beta)-\psi(\beta)=-\tfrac{1}{2}\log^{2}(1+\beta)=-\tfrac{1}{2}\log^{2}\beta. \]
Mais on a d'ailleurs par l'équation (k'),
\[ \psi(1-\beta)+\psi(\beta)=\frac{\pi^{2}}{6}-2\log^{2}\beta; \]
donc
\[ \begin{aligned} \psi(\beta)&=\frac{\pi^{2}}{10}-\log^{2}\beta\\ \psi(1-\beta)&=\frac{\pi^{2}}{15}-\log^{2}\beta. \end{aligned} \]
Ainsi on connaît la transcendante $\psi(x)$, non-seulement lorsque $x=1$ et $x=\frac{1}{2}$, mais encore lorsque $x=\beta=-\frac{1}{2}+\frac{1}{2}\sqrt{5}$, et lorsque $x=1-\beta=\frac{3}{2}-\frac{1}{2}\sqrt{5}$.

\subsection*{(26).}

Considérons maintenant la fonction
\[ \varphi(x)=x+\frac{x^{3}}{3^{2}}+\frac{x^{5}}{5^{2}}+\text{etc.} \]
On aura d'abord $\varphi(x)=\psi(x)-\frac{1}{4}\psi(x^{2})$, ce qui se vérifie immédiatement par les expressions en série de $\psi(x)$ et $\psi(x^{2})$. Mais on connaît les valeurs de $\psi(x)$ et $\psi(x^{2})$ lorsque $x=\beta$, puisqu'alors $x^{2}=1-\beta$; donc on aura pour ce même cas,
\[ \varphi(\beta)=\frac{\pi^{2}}{12}-\tfrac{3}{4}\log^{2}\beta. \]
\ednote{Under $\pi^{2}$ only a digit $1$, crossed by a reader's pen stroke, is legible in the scan; read as $12$, the value that $\psi(\beta)-\frac{1}{4}\psi(1-\beta)$ gives.}
Pour avoir d'autres valeurs de la fonction $\varphi$, j'observe qu'on a

\origpage{247}
\[ d\varphi(x)=dx\left(1+\frac{x^{2}}{3}+\frac{x^{4}}{5}+\text{etc.}\right)=\frac{dx}{2x}\log\left(\frac{1+x}{1-x}\right), \]
d'où
\[ \varphi(x)=\int\frac{dx}{2x}\log\left(\frac{1+x}{1-x}\right). \]
Mettons $\frac{1-x}{1+x}$ à la place de $x$, nous aurons semblablement
\[ \varphi\left(\frac{1-x}{1+x}\right)=\int\frac{dx}{1-x^{2}}\log x. \]
Ajoutant ces deux équations et effectuant l'intégration indiquée, on aura
\[ \varphi(x)+\varphi\left(\frac{1-x}{1+x}\right)=\varphi(1)+\tfrac{1}{2}\log x\log\left(\frac{1+x}{1-x}\right). \]
On a ajouté la constante $\varphi(1)$, afin que les deux membres soient égaux lorsque $x=0$. Quant à la valeur de $\varphi(1)$, elle se déduit de l'équation $\varphi(x)=\psi(x)-\frac{1}{4}\psi(x^{2})$, qui donne $\varphi(1)=\frac{3}{4}\psi(1)=\frac{\pi^{2}}{8}$; de sorte qu'on aura
\[ \varphi(x)+\varphi\left(\frac{1-x}{1+x}\right)=\frac{\pi^{2}}{8}+\tfrac{1}{2}\log x\log\left(\frac{1+x}{1-x}\right). \]
Donnons maintenant à $x$ une valeur particulière telle que $\frac{1-x}{1+x}=x$, il en résultera $x=\sqrt{2}-1=\gamma$; et d'après cette valeur, on aura
\[ \varphi(\gamma)=\frac{\pi^{2}}{8}-\tfrac{1}{2}\log^{2}\gamma. \]
\ednote{As printed. The preceding equation with $x=\gamma$ gives $2\varphi(\gamma)=\frac{\pi^{2}}{8}-\frac{1}{2}\log^{2}\gamma$, that is $\varphi(\gamma)=\frac{\pi^{2}}{16}-\frac{1}{4}\log^{2}\gamma$; the halving appears to have been omitted.}
On connaît donc quatre valeurs diverses de $\varphi(x)$, comme on en connaît quatre de $\psi(x)$; car la valeur connue $\varphi(\beta)$ en fait connaître une autre $\varphi\left(\frac{1-\beta}{1+\beta}\right)$, ou $\varphi(2\beta-1)=\varphi(\sqrt{5}-2)$. Ainsi les valeurs de $x$ pour lesquelles $\varphi(x)$ est connu, sont $x=1$, $x=\beta=-\frac{1}{2}+\frac{1}{2}\sqrt{5}$, $x=2\beta-1=\sqrt{5}-2$, $x=\gamma=\sqrt{2}-1$.

\subsection*{(27).}

Pour pousser encore plus loin ces recherches, considérons la fonction
\[ \Lambda(x)=x+\frac{x^{2}}{2^{3}}+\frac{x^{3}}{3^{3}}+\frac{x^{4}}{4^{3}}+\text{etc.}, \]
qui en général est une transcendante fort composée, puisque c'est un

\origpage{248}
problème assez difficile de déterminer par une suite convergente, la valeur de $\Lambda(1)$. On aura d'abord par la différentiation
\[ d\Lambda(x)=dx\left(1+\frac{x}{2^{2}}+\frac{x^{2}}{3^{2}}+\text{etc.}\right)=\frac{dx}{x}\psi(x); \]
ce qui donne $\Lambda(x)=\displaystyle\int\frac{dx}{x}\psi(x)$, ou
\[ \Lambda(x)=\log x\cdot\psi(x)+\int\frac{dx}{x}\log x\log(1-x). \]
Faisant comme ci-dessus $\log x=p$, $\log(1-x)=q$, on aura
\[ \Lambda(x)=\log x\cdot\psi(x)+\int pq\,dp, \]
d'où l'on déduit, en mettant $1-x$ à la place de $x$,
\[ \Lambda(1-x)=\log(1-x)\psi(1-x)+\int pq\,dq. \]
On aurait semblablement $\Lambda(-x)=\displaystyle\int\frac{dx}{x}\psi(-x)$, ou
\[ \Lambda(-x)=\log x\cdot\psi(-x)+\int\frac{dx}{x}\log x\log(1+x). \]
Mettons dans cette équation $\frac{x}{1-x}$ à la place de $x$, il viendra
\[ \Lambda\left(-\frac{x}{1-x}\right)=\log\frac{x}{1-x}\cdot\psi\left(-\frac{x}{1-x}\right)+\int q(p-q)(dp-dq). \]
Ajoutant ces trois équations, et effectuant les intégrations indiquées, on aura la formule générale
\[ \left.\begin{aligned} &\Lambda(x)+\Lambda(1-x)\\ &+\Lambda\left(-\frac{x}{1-x}\right) \end{aligned}\right\}=\left\{\begin{aligned} &\Lambda(1)+\log x\cdot\psi(x)+\log(1-x)\psi(1-x)\\ &+\log\frac{x}{1-x}\psi\left(-\frac{x}{1-x}\right)+\log x\log^{2}(1-x)-\tfrac{1}{3}\log^{3}(1-x). \end{aligned}\right. \]
Soit $x=\frac{1}{2}$, le premier membre deviendra $2\Lambda\left(\frac{1}{2}\right)+\Lambda(-1)$; et comme on a en général $\Lambda(x)+\Lambda(-x)=\frac{1}{4}\Lambda(x^{2})$, ce qui donne $\Lambda(-1)=-\frac{3}{4}\Lambda(1)$, on trouvera
\[ 2\Lambda\left(\tfrac{1}{2}\right)-\tfrac{3}{4}\Lambda(1)=\Lambda(1)+2\log\tfrac{1}{2}\cdot\psi\left(\tfrac{1}{2}\right)+\tfrac{2}{3}\log^{3}\left(\tfrac{1}{2}\right). \]
Substituant la valeur connue de $\psi\left(\frac{1}{2}\right)$, il viendra
\[ \tfrac{7}{8}\Lambda(1)=\Lambda\left(\tfrac{1}{2}\right)+\frac{\pi^{2}}{12}\log(2)-\tfrac{1}{6}(\log 2)^{3}; \]

\origpage{249}
d'où l'on voit que $\Lambda\left(\frac{1}{2}\right)$ peut se déterminer par le moyen de $\Lambda(1)$; et réciproquement $\Lambda(1)$ peut se déterminer par le moyen de $\Lambda\left(\frac{1}{2}\right)$, ce qui est une première manière d'exprimer la transcendante $\Lambda(1)$ par une suite convergente.

En second lieu, si on fait $x=1-\beta=\beta^{2}$, $\beta$ ayant la même valeur que ci-dessus, la formule générale donnera
\[ \begin{aligned} \Lambda(\beta)+\Lambda(\beta^{2})+\Lambda(-\beta)&=\log\beta\left[2\psi(\beta^{2})+\psi(\beta)+\psi(-\beta)\right]\\ &\quad+\Lambda(1)+\tfrac{5}{3}\log^{3}\beta. \end{aligned} \]
Mais on a $\psi(\beta)+\psi(-\beta)=\frac{1}{2}\psi(\beta^{2})$ et $\Lambda(\beta)+\Lambda(-\beta)=\frac{1}{4}\Lambda(+\beta^{2})$; donc
\[ \tfrac{5}{4}\Lambda(\beta^{2})=\tfrac{5}{2}\log\beta\,\psi(\beta^{2})+\Lambda(1)+\tfrac{5}{3}\log^{3}\beta; \]
or on a trouvé $\psi(\beta^{2})=\frac{\pi^{2}}{15}-\log^{2}\beta$; donc enfin
\[ \Lambda(1)=\tfrac{5}{4}\Lambda(\beta^{2})-\frac{\pi^{2}}{6}\log\beta+\tfrac{5}{6}\log^{3}\beta. \]
Ainsi la quantité $\Lambda(1)$ qui représente la suite $1+\frac{1}{2^{3}}+\frac{1}{3^{3}}+\frac{1}{4^{3}}+\text{etc.}$, peut se déterminer par le moyen de $\Lambda(\beta^{2})$, c'est-à-dire par la suite convergente $\beta^{2}+\frac{\beta^{4}}{2^{3}}+\frac{\beta^{6}}{3^{3}}+\frac{\beta^{8}}{4^{3}}+\text{etc.}$, dans laquelle $\beta^{2}=\frac{1}{2}(3-\sqrt{5})$, quantité plus petite que 0,4. Ces formules ont été données sans démonstration dans l'ouvrage cité de Landen, pag.~118; et jusqu'à présent on n'est pas allé plus loin dans la théorie de ces sortes de transcendantes.

\section*{Considération des formules intégrales $\displaystyle\int\frac{x^{p-1}\,dx\log\frac{1}{x}}{\sqrt[n]{(1-x^{n})^{n-q}}}$, $\displaystyle\int\frac{x^{p-1}\,dx\log^{2}\frac{1}{x}}{\sqrt[n]{(1-x^{n})^{n-q}}}$, etc. Théorème très remarquable sur la première de ces formules.}

\subsection*{(28).}

Si on désigne par $Z$ la fonction $\left(\frac{p}{q}\right)$ ou l'intégrale $\displaystyle\int\frac{x^{p-1}\,dx}{\sqrt[n]{(1-x^{n})^{n-q}}}$, prise depuis $x=0$ jusqu'à $x=1$, les différentielles successives de $Z$, prises en faisant varier $p$ seule, donneront les

\origpage{250}
nouvelles formules
\[ \begin{aligned} \frac{dZ}{dp}&=-\int\frac{x^{p-1}\,dx\log\frac{1}{x}}{\sqrt[n]{(1-x^{n})^{n-q}}}\\ \frac{ddZ}{dp^{2}}&=\int\frac{x^{p-1}\,dx\log^{2}\frac{1}{x}}{\sqrt[n]{(1-x^{n})^{n-q}}}\\ \frac{d^{3}Z}{dp^{3}}&=-\int\frac{x^{p-1}\,dx\log^{3}\frac{1}{x}}{\sqrt[n]{(1-x^{n})^{n-q}}}\\ &\text{etc.}, \end{aligned} \]
ces intégrales étant encore prises depuis $x=0$ jusqu'à $x=1$.

Comme on a généralement par la formule (i'),
\[ Z=\frac{1}{p}+\frac{n-q}{n}\cdot\frac{1}{n+p}+\frac{n-q\cdot 2n-q}{n\cdot 2n}\cdot\frac{1}{2n+p}+\frac{n-q\cdot 2n-q\cdot 3n-q}{n\cdot 2n\cdot 3n}\cdot\frac{1}{3n+p}+\text{etc.} \]
on en déduit par des différentiations successives,
\[ \left.\begin{aligned} \int\frac{x^{p-1}\,dx\log\frac{1}{x}}{\sqrt[n]{(1-x^{n})^{n-q}}}&=\frac{1}{p^{2}}+\frac{n-q}{n}\cdot\frac{1}{(n+p)^{2}}+\frac{n-q\cdot 2n-q}{n\cdot 2n}\cdot\frac{1}{(2n+p)^{2}}+\text{etc.}\\ \frac{1}{2}\int\frac{x^{p-1}\,dx\log^{2}\frac{1}{x}}{\sqrt[n]{(1-x^{n})^{n-q}}}&=\frac{1}{p^{3}}+\frac{n-q}{n}\cdot\frac{1}{(n+p)^{3}}+\frac{n-q\cdot 2n-q}{n\cdot 2n}\cdot\frac{1}{(2n+p)^{3}}+\text{etc.}\\ \frac{1}{2\cdot 3}\int\frac{x^{p-1}\,dx\log^{3}\frac{1}{x}}{\sqrt[n]{(1-x^{n})^{n-q}}}&=\frac{1}{p^{4}}+\frac{n-q}{n}\cdot\frac{1}{(n+p)^{4}}+\frac{n-q\cdot 2n-q}{n\cdot 2n}\cdot\frac{1}{(2n+p)^{4}}+\text{etc.}\\ &\text{etc.} \end{aligned}\right\}\qquad(m'). \]
Mais ces suites ont l'inconvénient de n'être pas suffisamment convergentes.

\subsection*{(29).}

Pour avoir ces mêmes valeurs exprimées en suites plus convergentes, il faudrait partir de la formule (h'), et la différentier par rapport à $p$, autant de fois qu'il est nécessaire.

En la différentiant une fois, on aura

\origpage{251}
\[ \begin{aligned} \int\frac{x^{p-1}\,dx\log\frac{1}{x}}{\sqrt[n]{(1-x^{n})^{n-q}}}={}&\frac{2^{-\frac{p}{n}}\log 2}{n}\left[\frac{1}{p}+\frac{n-q}{2n}\cdot\frac{1}{n+p}+\frac{n-q\cdot 2n-q}{2n\cdot 4n}\cdot\frac{1}{2n+p}+\text{etc.}\right]\\ &+2^{-\frac{p}{n}}\left[\frac{1}{p^{2}}+\frac{n-q}{2n}\cdot\frac{1}{(n+p)^{2}}+\frac{n-q\cdot 2n-q}{2n\cdot 4n}\cdot\frac{1}{(2n+p)^{2}}+\text{etc}\right]\\ &+2^{-\frac{q}{n}}\left\{\begin{aligned} &\frac{1}{2n}\cdot\frac{1}{n+q}+\frac{n-p\cdot 2n-p}{2n\cdot 4n(2n+q)}\left(\frac{1}{n-p}+\frac{1}{2n-p}\right)\\ &+\frac{n-p\cdot 2n-p\cdot 3n-p}{2n\cdot 4n\cdot 6n\cdot(3n+q)}\left(\frac{1}{n-p}+\frac{1}{2n-p}+\frac{1}{3n-p}\right)\\ &+\text{etc.} \end{aligned}\right\} \end{aligned} \]
Ces suites sont convergentes, puisque chaque terme est moindre que la moitié du terme précédent; mais leur forme est compliquée, et elle le deviendrait davantage dans la différentielle du second ordre ou d'un ordre plus élevé. C'est pourquoi il convient d'avoir recours à d'autres moyens si l'on veut évaluer facilement les intégrales dont il s'agit.

\subsection*{(30).}

Désignons par $\varphi\left(\frac{p}{q}\right)$ l'intégrale $\displaystyle\int\frac{x^{p-1}\,dx\log\frac{1}{x}}{\sqrt[n]{(1-x^{n})^{n-q}}}$, et par $\psi\left(\frac{p}{q}\right)$ le rapport de cette intégrale à la fonction $\left(\frac{p}{q}\right)$, déjà représentée par $Z$, ensorte qu'on ait
\[ \psi\left(\frac{p}{q}\right)=\frac{\varphi\left(\frac{p}{q}\right)}{Z}; \]
suivant ce qui a été dit (art.~28), on aura
\[ \begin{aligned} \varphi\left(\frac{p}{q}\right)&=-\frac{dZ}{dp}\\ \text{et}\quad\psi\left(\frac{p}{q}\right)&=-\frac{dZ}{dp}\cdot\frac{1}{Z}=-\frac{d(\log Z)}{dp}. \end{aligned} \]
Mais puisqu'on a aussi $Z=\left(\frac{q}{p}\right)$, il en en\ednote{The word ``en'' is printed twice; an apparent misprint.} résulte
\[ \frac{dZ}{dq}=-\varphi\left(\frac{q}{p}\right), \]
où il faut observer que $\varphi\left(\frac{q}{p}\right)$ n'est pas la même chose que $\varphi\left(\frac{p}{q}\right)$.

\origpage{252}

La différentielle complète de $Z$ ou de $\left(\frac{p}{q}\right)$ pourra donc être exprimée ainsi:
\[ d\left(\frac{p}{q}\right)=-dp\,\varphi\left(\frac{p}{q}\right)-dq\,\varphi\left(\frac{q}{p}\right)\qquad(n'). \]

\subsection*{(31).}

Cela posé, si on prend l'équation générale
\[ \left(\frac{p}{q}\right)\cdot\left(\frac{p+q}{r}\right)=\left(\frac{p}{r}\right)\cdot\left(\frac{p+r}{q}\right), \]
qui peut être mise sous la forme
\[ \log\left(\frac{p}{q}\right)+\log\left(\frac{p+q}{r}\right)=\log\left(\frac{p}{r}\right)+\log\left(\frac{p+r}{q}\right), \]
et qu'on la différentie par rapport à $p$, on aura
\[ \frac{\varphi\left(\frac{p}{q}\right)}{\left(\frac{p}{q}\right)}+\frac{\varphi\left(\frac{p+q}{r}\right)}{\left(\frac{p+q}{r}\right)}=\frac{\varphi\left(\frac{p}{r}\right)}{\left(\frac{p}{r}\right)}+\frac{\varphi\left(\frac{p+r}{q}\right)}{\left(\frac{p+r}{q}\right)}, \]
ou suivant les dénominations établies,
\[ \psi\left(\frac{p}{q}\right)+\psi\left(\frac{p+q}{r}\right)=\psi\left(\frac{p}{r}\right)+\psi\left(\frac{p+r}{q}\right)\qquad(p'). \]
La même équation étant différentiée par rapport à $q$, donne
\[ \psi\left(\frac{q}{p}\right)+\psi\left(\frac{p+q}{r}\right)=\psi\left(\frac{q}{p+r}\right)\qquad(q'). \]
La différentielle par rapport à $r$ donnerait un résultat de la même forme que le précédent, et qui y serait par conséquent compris. On peut de plus faire voir que l'équation à quatre termes (p'), est comprise dans l'équation (q'); car de celle-ci on déduit, par la permutation des lettres $p$ et $q$,
\[ \psi\left(\frac{p}{q}\right)+\psi\left(\frac{p+q}{r}\right)=\psi\left(\frac{p}{q+r}\right), \]
et de cette dernière on conclut, par l'échange des lettres $q$ et $r$,
\[ \psi\left(\frac{p}{r}\right)+\psi\left(\frac{p+r}{q}\right)=\psi\left(\frac{p}{q+r}\right), \]

\origpage{253}
donc
\[ \psi\left(\frac{p}{q}\right)+\psi\left(\frac{p+q}{r}\right)=\psi\left(\frac{p}{r}\right)+\psi\left(\frac{p+r}{q}\right), \]
ce qui est l'équation (p').

\subsection*{(32).}

De là on voit qu'il suffit de considérer l'équation à trois termes (q'), et c'est de cette source que nous allons tirer toutes les relations qui existent entre les diverses quantités $\psi\left(\frac{p}{q}\right)$, qui répondent à une même valeur de $n$.

Parmi les quantités $\left(\frac{p}{q}\right)$, on distingue celles de la forme $\left(\frac{a}{n}\right)$, dont la valeur est $\left(\frac{a}{n}\right)=\frac{1}{a}$; il en résulte $\varphi\left(\frac{a}{n}\right)=\frac{1}{a^{2}}$, et par conséquent
\[ \psi\left(\frac{a}{n}\right)=\frac{1}{a}\qquad(r'), \]
première formule qui servira à la réduction des autres.

Parmi les mêmes quantités, on trouve en second lieu la formule remarquable
\[ \left(\frac{a}{n-a}\right)=\frac{\omega}{\sin a\omega}, \]
d'où l'on déduit, en prenant la différentielle de chaque membre par rapport à $a$,
\[ -\varphi\left(\frac{a}{n-a}\right)+\varphi\left(\frac{n-a}{a}\right)=-\frac{\omega^{2}\cos a\omega}{\sin^{2}a\omega}. \]
Cette équation étant divisée par $\left(\frac{a}{n-a}\right)$ donne
\[ \psi\left(\frac{a}{n-a}\right)-\psi\left(\frac{n-a}{a}\right)=\omega\cot a\omega. \]

La fonction $\psi\left(\frac{a}{n-a}\right)$ est remarquable dans l'ordre des fonctions $\psi$, comme la quantité $\left(\frac{a}{n-a}\right)$ l'est parmi celles de son espèce. Nous nous servirons donc de cette fonction pour exprimer toutes les autres, et nous ferons, pour abréger,
\[ \psi\left(\frac{a}{n-a}\right)=B_{a}.\qquad(s') \]

\origpage{254}
Cela posé, l'équation qu'on vient de trouver s'exprimera ainsi:
\[ B_{a}-B_{n-a}=\omega\cot a\omega.\qquad(t') \]
Elle fait voir que la valeur de $B_{a}$ se conclut de celle de $B_{n-a}$, et réciproquement; d'où il suit que dans les quantités $B_{1}$, $B_{2}$, $B_{3}$, etc., il suffit de connaître les premières jusqu'à celle dont le rang est $\frac{n}{2}$ ou $\frac{n-1}{2}$ inclusivement.

\subsection*{(33).}

De l'équation (q'), on déduit généralement
\[ \psi\left(\frac{a}{b}\right)-\psi\left(\frac{b}{a}\right)=\psi\left(\frac{a}{b-a}\right),\qquad(u') \]
donc en particulier,
\[ \psi\left(\frac{a}{n}\right)-\psi\left(\frac{n}{a}\right)=\psi\left(\frac{a}{n-a}\right)=B_{a}, \]
et par conséquent
\[ \psi\left(\frac{n}{a}\right)=\frac{a}{1}-B_{a}.\qquad(v'). \]
\ednote{Printed $\frac{a}{1}$ instead of $\frac{1}{a}$; an apparent misprint, since $\psi\left(\frac{a}{n}\right)=\frac{1}{a}$ by (r'), and (v') is later quoted as $\psi\left(\frac{n}{r}\right)=\frac{1}{r}-B_{r}$.}
Puisqu'on a $\left(\frac{p+n}{q}\right)=\frac{p}{p+q}\cdot\left(\frac{p}{q}\right)$, cette équation étant différentiée logarithmiquement par rapport à $p$, donne
\[ \psi\left(\frac{q}{p+n}\right)=\psi\left(\frac{p}{q}\right)-\frac{1}{p}+\frac{1}{p+q};\qquad(x') \]
\ednote{The left side is printed $\psi\left(\frac{q}{p+n}\right)$, as in (y'); differentiating $\left(\frac{p+n}{q}\right)$ logarithmically with respect to $p$ gives $\psi\left(\frac{p+n}{q}\right)$. An apparent misprint.}
et la même équation étant différentiée par rapport à $q$, donnera
\[ \psi\left(\frac{q}{p+n}\right)=\psi\left(\frac{q}{p}\right)+\frac{1}{p+q}.\qquad(y') \]
Ces deux équations serviront au besoin à transformer toute fonction $\psi\left(\frac{a}{b}\right)$ en une fonction semblable, dans laquelle $a$ et $b$ seraient plus petits que $n$: on en déduit, par exemple,
\[ \psi\left(\frac{n}{n+c}\right)=\psi\left(\frac{n}{c}\right)+\frac{1}{n+c}. \]
Mais on a $B_{c}=\frac{1}{c}-\psi\left(\frac{n}{c}\right)$ et $B_{n+c}=\frac{1}{n+c}-\psi\left(\frac{n}{n+c}\right)$; donc
\[ B_{n+c}=B_{c}-\frac{1}{c},\qquad(z') \]

\origpage{255}
formule qui servira à trouver la valeur de $B_{a}$, si $a$ est plus grand que $n$. D'ailleurs il est aisé de voir qu'on a $B_{n}=0$ et $B_{0}=\infty$.

Revenons à l'équation (q') et faisons $p+q+r=n$, afin qu'on ait à la fois $\psi\left(\frac{q}{p+r}\right)=B_{q}$ et $\psi\left(\frac{p+q}{r}\right)=B_{n+q}=B_{n-r}$\ednote{Printed $B_{n+q}$; $B_{p+q}$ is expected, since $\frac{p+q}{r}=\frac{p+q}{n-(p+q)}$. An apparent misprint.}, on aura
\[ \psi\left(\frac{q}{n-r-q}\right)=B_{q}-B_{n-r},\qquad(a'') \]
c'est la valeur de toute fonction $\psi\left(\frac{a}{b}\right)$, dans laquelle on a $a+b<n$.

Si dans la même équation (q') on fait $p=n-q$, on aura
\[ \psi\left(\frac{q}{n-q+r}\right)=\psi\left(\frac{q}{n-q}\right)+\psi\left(\frac{n}{r}\right). \]
Mais par l'équation (v') on a $\psi\left(\frac{n}{r}\right)=\frac{1}{r}-B_{r}$; donc
\[ \psi\left(\frac{q}{n+r-q}\right)=B_{q}-B_{r}+\frac{1}{r};\qquad(b'') \]
c'est la valeur de toute fonction $\psi\left(\frac{a}{b}\right)$, dans laquelle on a $a+b>n$.

\subsection*{(34).}

On pourra donc déterminer généralement la valeur de toute fonction $\psi\left(\frac{a}{b}\right)$, si on détermine celle de l'auxiliaire $B_{a}$, ou celle de la fonction $\psi\left(\frac{n}{a}\right)$, puisqu'on a $B_{a}=\frac{1}{a}-\psi\left(\frac{n}{a}\right)$.

Or on a $\left(\frac{n}{a}\right)=\frac{1}{a}$, et $\varphi\left(\frac{n}{a}\right)=\displaystyle\int\frac{x^{n-1}\,dx\log\frac{1}{x}}{\sqrt[n]{(1-x^{n})^{n-a}}}$; donc tout se réduit à trouver l'intégrale
\[ \psi\left(\frac{n}{a}\right)=a\int\frac{x^{n-1}\,dx\log\frac{1}{x}}{\sqrt[n]{(1-x^{n})^{n-a}}}, \]
prise à l'ordinaire depuis $x=0$ jusqu'à $x=1$.

Pour cela, soit $x^{n}=1-y^{n}$, l'intégrale précédente aura pour

\origpage{256}
transformée
\[ \begin{aligned} \frac{a}{n}\int-y^{a-1}\,dy\log(1-y^{n})&=-\frac{y^{a}}{n}\log(1-y^{n})-\int\frac{y^{a+n-1}\,dy}{1-y^{n}}\\ &=\frac{y^{a}}{a}-\frac{y^{a}}{n}\log(1-y^{n})-\int\frac{y^{a-1}\,dy}{1-y^{n}}. \end{aligned} \]
D'un autre côté, $y^{a}\log(1-y^{n})$ peut être mis sous la forme $(y^{a}-1)\log(1-y^{n})-\displaystyle\int\frac{ny^{n-1}\,dy}{1-y^{n}}$; on aura donc
\[ \psi\left(\frac{n}{a}\right)=\frac{y^{a}}{a}+\frac{1-y^{a}}{n}\log(1-y^{n})-\int\left(\frac{y^{a-1}-y^{n-1}}{1-y^{n}}\right)dy. \]
La partie hors du signe s'évanouit en faisant $y=0$, et elle se réduit à $\frac{1}{a}$ en faisant $y=1$; donc
\[ \psi\left(\frac{n}{a}\right)=\frac{1}{a}-\int\left(\frac{y^{a-1}-y^{n-1}}{1-y^{n}}\right)dy; \]
donc en général,
\[ B_{a}=\int\left(\frac{y^{a-1}-y^{n-1}}{1-y^{n}}\right)dy,\qquad(c'') \]
cette intégrale étant prise depuis $y=0$ jusqu'à $y=1$.

\subsection*{(35).}

Cela posé, en faisant toujours $\omega=\frac{\pi}{n}$, on trouve par les formules connues pour l'intégration des fractions rationnelles:
\[ \begin{aligned} B_{a}={}&\frac{1}{n}\log n+\frac{n-2}{n}\,\omega\sin 2a\omega-\frac{2}{n}\cos 2a\omega\log(2\sin\omega)\\ &+\frac{n-4}{n}\,\omega\sin 4a\omega-\frac{2}{n}\cos 4a\omega\log(2\sin 2\omega)\\ &+\frac{n-6}{n}\,\omega\sin 6a\omega-\frac{2}{n}\cos 6a\omega\log(2\sin 3\omega)\\ &+\text{etc.} \end{aligned} \]
Cette suite devra être prolongée jusqu'au terme
\[ \frac{1}{n}\,\omega\sin(n-1)a\omega-\frac{2}{n}\cos(n-1)a\omega\log\left(2\sin\frac{n-1}{2}\,\omega\right), \]
si $n$ est impair; mais si $n$ est pair, il faudra prendre seulement la moitié du dernier terme, laquelle sera $-\frac{1}{n}\cos a\pi\log 2$.

\origpage{257}

Si l'on se rappelle ensuite les formules
\[ \begin{aligned} \sin x+\sin 2x+\sin 3x\ldots+\sin mx&=\frac{\sin x+\sin mx-\sin(m+1)x}{2(1-\cos x)},\\ \sin x+2\sin 2x+3\sin 3x\ldots+m\sin mx&=\frac{(m+1)\sin mx-m\sin(m+1)x}{2(1-\cos x)}, \end{aligned} \]
on trouvera que la suite
\[ \frac{n-2}{n}\,\omega\sin 2a\omega+\frac{n-4}{n}\,\omega\sin 4a\omega+\frac{n-6}{n}\,\omega\sin 6a\omega+\text{etc.}, \]
prolongée jusqu'à un nombre de termes $\frac{n-1}{2}$ ou $\frac{n-2}{2}$, a pour somme $\frac{1}{2}\omega\cot a\omega$. Donc si $n$ est impair, on aura
\[ \left.\begin{aligned} B_{a}={}&\frac{1}{n}\log n+\tfrac{1}{2}\omega\cot a\omega-\frac{2}{n}\cos 2a\omega\log(2\sin\omega)\\ &-\frac{2}{n}\cos 4a\omega\log(2\sin 2\omega)\\ &-\frac{2}{n}\cos 6a\omega\log(2\sin 3\omega)\\ &\vdots\\ &-\frac{2}{n}\cos(n-1)a\omega\log\left(2\sin\frac{n-1}{2}\,\omega\right),\\ &\text{et si } n \text{ est pair, on aura}\\ B_{a}={}&\frac{1}{n}\log n+\tfrac{1}{2}\omega\cot a\omega-\frac{1}{n}\cos a\pi\log 2\\ &-\frac{2}{n}\cos 2a\omega\log(2\sin\omega)\\ &-\frac{2}{n}\cos 4a\omega\log(2\sin 2\omega)\\ &\vdots\\ &-\frac{2}{n}\cos(n-2)a\omega\log\left(2\sin\frac{n-2}{n}\,\omega\right). \end{aligned}\right\}\qquad(d'') \]
\ednote{The last row is printed $\log\left(2\sin\frac{n-2}{n}\,\omega\right)$; the general term $\log(2\sin k\omega)$ with $k=\frac{n-2}{2}$ requires $\frac{n-2}{2}$. An apparent misprint.}

\subsection*{(36).}

Dans le cas particulier où l'on a $a=\frac{1}{2}n$, $n$ étant pair, on trouve directement par la formule (c''),
\[ B_{a}=\int\frac{y^{a-1}\,dy}{1+y^{a}}=\frac{1}{a}\log(1+y^{a})=\frac{2}{n}\log 2. \]
Pour que cette valeur s'accorde avec celle que donne dans le même

\origpage{258}
cas la formule (d''), il faut qu'on ait
\[ \begin{aligned} &\log(2\sin\omega)-\log(2\sin 2\omega)+\log(2\sin 3\omega)\ldots\\ &\quad\pm\log\left[2\sin\left(\frac{n}{2}-1\right)\omega\right]=\left(1+\tfrac{1}{2}\cos a\pi\right)\log 2-\tfrac{1}{2}\log n; \end{aligned} \]
ou, ce qui revient au même, en passant aux nombres, et supposant soit $n=4i$, soit $n=4i+2$,
\[ \frac{\sin\omega\sin 3\omega\sin 5\omega\ldots\sin(2i-1)\omega}{\sin 2\omega\sin 4\omega\sin 6\omega\ldots\sin(2i\omega)}=\sqrt{\left(\frac{2}{n}\right)}. \]
Cette formule est facile à vérifier au moyen de la valeur de $\sin nz$ donnée par Euler dans son \emph{Introd. in anal.}, pag.~204; car en faisant successivement dans cette valeur, $z$ infiniment petit, et $nz=\frac{1}{2}\pi$, on en tire
\[ \begin{gathered} \sin\omega\sin 2\omega\sin 3\omega\ldots\sin\frac{n}{2}\,\omega=2^{\frac{1-n}{2}}\sqrt{n},\\ \sin\omega\sin 3\omega\sin 5\omega\ldots\sin(2i-1)\omega=2^{\frac{2-n}{4}}; \end{gathered} \]
$2i$ étant $\frac{n}{2}$ ou $\frac{n-2}{2}$, selon que $n$ est de la forme $4i$ ou $4i+2$, et de là résulte l'équation précédente.

\subsection*{(37).}

Ayant l'expression générale de $B_{a}$, on en déduit celle de $\psi\left(\frac{p}{q}\right)$, au moyen de l'une ou l'autre des formules
\[ \left.\begin{aligned} \psi\left(\frac{p}{q}\right)&=B_{p}-B_{p+q}\\ \psi\left(\frac{p}{q}\right)&=B_{p}-B_{p+q-n}+\frac{1}{p+q-n}, \end{aligned}\right\}\qquad(e'') \]
la première ayant lieu lorsque $p+q$ est $<n$, et la seconde lorsque $p+q$ est $>n$. En cas d'égalité, on a simplement $\psi\left(\frac{p}{q}\right)=B_{p}$.

Ces valeurs qui sont déduites des formules (a'') et (b'') peuvent aussi être mises sous une seule forme générale, qui est
\[ \psi\left(\frac{p}{q}\right)=\int\left(\frac{y^{p-1}-y^{p+q-1}}{1-y^{n}}\right)dy,\qquad(f'') \]
ce qui prouve immédiatement que la fonction $\psi\left(\frac{p}{q}\right)$ est toujours

\origpage{259}
déterminable par les arcs de cercle et les logarithmes. On parvient ainsi à ce théorème très-remarquable par son élégance et sa généralité.

\emph{Si l'on prend les trois intégrales}
\[ \int\frac{x^{p-1}\,dx\log\frac{1}{x}}{\sqrt[n]{(1-x^{n})^{n-q}}},\qquad\int\frac{x^{p-1}\,dx}{\sqrt[n]{(1-x^{n})^{n-q}}},\qquad\int\frac{x^{p-1}-x^{p+q-1}}{1-x^{n}}\,dx, \]
\emph{entre les limites} $x=0$, $x=1$, \emph{la première sera égale au produit des deux autres.}

\subsection*{(38).}

On trouve dans le tome IV du \emph{Calcul intégral} d'Euler, pag.~166, une démonstration du même théorème qui, à cause de sa grande simplicité, aurait dû être substituée à la précédente, s'il n'y avait pas quelque avantage à considérer un même objet sous différens points de vue. Voici cette démonstration${}^{(1)}$.

Par la formule de l'art.~4, on a
\[ Z=\frac{p+q}{p}\cdot\frac{p+q+n}{p+n}\cdot\frac{p+q+2n}{p+2n}\cdot\text{etc.} \]
Différentiant logarithmiquement les deux membres par rapport à $p$, il viendra
\[ \frac{dZ}{Z}=\left(\frac{dp}{p+q}-\frac{dp}{p}\right)+\left(\frac{dp}{p+q+n}-\frac{dp}{p+n}\right)+\left(\frac{dp}{p+q+2n}-\frac{dp}{p+2n}\right)+\text{etc.} \]
Soit $\Phi$ ou $\Phi(\nu)$ une fonction de $\nu$, telle qu'on ait
\[ \Phi=\frac{\nu^{p}}{p}-\frac{\nu^{p+q}}{p+q}+\frac{\nu^{p+n}}{p+n}-\frac{\nu^{p+q+n}}{p+q+n}+\frac{\nu^{p+2n}}{p+2n}-\frac{\nu^{p+q+2n}}{p+q+2n}+\text{etc.} \]
La fonction $\Phi(\nu)$ devenant $\Phi(1)$ lorsque $\nu=1$, on aura
\[ \frac{dZ}{Z}=-dp\,\Phi(1). \]

\textbf{(1)} Euler n'avait donné d'abord qu'un cas particulier de ce théorème, qu'on trouve dans le tome I de ses \emph{Opusc. anal.}, pag.~183; il est ensuite parvenu au théorème général, dont il a donné une démonstration que nous rapportons ici; mais cette circonstance m'avait échappé lorsque j'ai publié mon \emph{Mémoire sur les intégrales définies}, où je n'ai cité que le cas particulier des \emph{Opusc. anal.}

\origpage{260}
Mais $\frac{dZ}{dp}=-\varphi\left(\frac{p}{q}\right)$, donc $\frac{1}{Z}\varphi\left(\frac{p}{q}\right)$ ou $\psi\left(\frac{p}{q}\right)=\Phi(1)$. Ainsi pour avoir la valeur de la fonction $\psi\left(\frac{p}{q}\right)$, il ne s'agit que de trouver celle de $\Phi(1)$.

Or en différentiant par rapport à $\nu$ la valeur de $\Phi$, on a
\[ d\Phi=d\nu\left\{\begin{aligned} &\nu^{p-1}+\nu^{n+p-1}+\nu^{2n+p-1}+\text{etc.}\\ &-\nu^{p+q-1}-\nu^{n+p+q-1}-\nu^{2n+p+q-1}-\text{etc.} \end{aligned}\right\} \]
ou en sommant ces suites,
\[ d\Phi=d\nu\left(\frac{\nu^{p-1}-\nu^{p+q-1}}{1-\nu^{n}}\right); \]
donc
\[ \Phi=\int\left(\frac{\nu^{p-1}-\nu^{p+q-1}}{1-\nu^{n}}\right)d\nu. \]
Cette intégrale devra être prise de manière qu'elle s'évanouisse lorsque $\nu=0$, on fera ensuite $\nu=1$, et on aura la valeur cherchée de $\Phi(1)$. Donc comme on peut mettre $x$ à la place de $\nu$, on aura en général
\[ \psi\left(\frac{p}{q}\right)=\int\left(\frac{x^{p-1}-x^{p+q-1}}{1-x^{n}}\right)dx, \]
l'intégrale du second membre étant prise entre les limites $x=0$, $x=1$. De là résulte le théorème de l'article précédent.

\subsection*{(39).}

Pour étendre encore davantage cette théorie, considérons les deux suites d'intégrales en $Z$ et en $T$, prises depuis $x=0$ jusqu'à $x=1$, savoir:
\[ \begin{aligned} Z&=\int\frac{x^{p-1}\,dx}{\sqrt[n]{(1-x^{n})^{n-q}}}, &\qquad& T=\int\frac{x^{p-1}-x^{p+q-1}}{1-x^{n}}\,dx\\ Z'&=\int\frac{x^{p-1}\,dx\log\frac{1}{x}}{\sqrt[n]{(1-x^{n})^{n-q}}}, && T'=\int\frac{x^{p-1}-x^{p+q-1}}{1-x^{n}}\,dx\log\frac{1}{x}\\ Z''&=\int\frac{x^{p-1}\,dx\log^{2}\frac{1}{x}}{\sqrt[n]{(1-x^{n})^{n-q}}}, && T''=\int\frac{x^{p-1}-x^{p+q-1}}{1-x^{n}}\,dx\log^{2}\frac{1}{x}\\ &\text{etc.} && \text{etc.} \end{aligned} \]

\origpage{261}
Il résulte d'abord du théorème précédent qu'on a
\[ Z'=ZT. \]
Différentiant cette équation par rapport à $p$, et observant qu'on a $\frac{dZ}{dp}=-Z'$, $\frac{dZ'}{dp}=-Z''$, $\frac{dT}{dp}=-T'$, on aura $Z''=Z'T+ZT'$, ou
\[ Z''=Z(T^{2}+T'). \]
Celle-ci étant différentiée de nouveau par rapport à $p$, donne
\[ Z'''=Z(T^{3}+3TT'+T''), \]
et ainsi de suite, la loi de ces expressions étant analogue à celle des différentielles successives de la formule $ue^{\displaystyle\int u\,dp}$.

Si l'on veut donc avoir les valeurs des quantités $Z'$, $Z''$, $Z'''$, etc., ou simplement leur rapport à la fonction primitive $Z$, il faudra connaître les quantités $T$, $T'$, $T''$, etc. en pareil nombre. Mais comme celles-ci sont rationnelles, et contiennent des puissances moins élevées de $\log\frac{1}{x}$, on voit qu'au moins la difficulté est diminuée.

\subsection*{(40).}

Si l'on intègre la différentielle $x^{m+\alpha-1}\,dx$, depuis $x=0$ jusqu'à $x=1$, on aura
\[ \int x^{m+\alpha-1}\,dx=\frac{1}{m+\alpha}=\frac{1}{m}-\frac{\alpha}{m^{2}}+\frac{\alpha^{2}}{m^{3}}-\frac{\alpha^{3}}{m^{4}}+\text{etc.} \]
Si on met la même différentielle sous la forme $x^{m-1}\,dx\cdot x^{\alpha}$, ou
\[ x^{m-1}\,dx\left(1+\alpha\log x+\frac{\alpha^{2}}{1\cdot 2}\log^{2}x+\frac{\alpha^{3}}{1\cdot 2\cdot 3}\log^{3}x+\text{etc.}\right), \]
son intégrale prise entre les mêmes limites, sera
\[ \int x^{m-1}\,dx+\alpha\int x^{m-1}\,dx\log x+\frac{\alpha^{2}}{1\cdot 2}\int x^{m-1}\,dx\log^{2}x+\text{etc.} \]
L'identité de ces deux formules exige donc qu'on ait
\[ \begin{aligned} \int x^{m-1}\,dx&=\frac{1}{m}\\ \int x^{m-1}\,dx\log\left(\frac{1}{x}\right)&=\frac{1}{m^{2}}\\ \int x^{m-1}\,dx\log^{2}\left(\frac{1}{x}\right)&=\frac{1\cdot 2}{m^{3}}\\ \int x^{m-1}\,dx\log^{3}\left(\frac{1}{x}\right)&=\frac{1\cdot 2\cdot 3}{m^{4}}, \end{aligned} \]

\origpage{262}
et en général
\[ \int x^{m-1}\,dx\log^{n}\left(\frac{1}{x}\right)=\frac{1\cdot 2\cdot 3\ldots n}{m^{n+1}}.\qquad(g''). \]

\subsection*{(41).}

A l'aide de cette formule on peut exprimer les quantités $T'$, $T''$, etc., par des séries régulières, composées des puissances réciproques des nombres pris à des intervalles égaux dans la suite des nombres naturels.

Ainsi en développant d'abord la différentielle qu'il faut intégrer pour avoir $T'$, on a
\[ T'=\int dx\log\frac{1}{x}\left\{\begin{aligned} &x^{p-1}+x^{p+n-1}+x^{p+2n-1}+\text{etc.}\\ &-x^{p+q-1}-x^{p+q+n-1}-x^{p+q+2n-1}-\text{etc.} \end{aligned}\right\}, \]
et effectuant l'intégration entre les limites $x=0$, $x=1$, il vient
\[ \begin{aligned} T'={}&\frac{1}{p^{2}}+\frac{1}{(p+n)^{2}}+\frac{1}{(p+2n)^{2}}+\text{etc.}\\ &-\frac{1}{(p+q)^{2}}-\frac{1}{(p+q+n)^{2}}-\frac{1}{(p+q+2n)^{2}}+\text{etc.} \end{aligned} \]
\ednote{The second line closes with $+$ etc. as printed; $-$ etc. would be expected.}
Désignons en général par $(a,n)^{m}$, la somme de la suite
\[ \frac{1}{a^{m}}+\frac{1}{(a+n)^{m}}+\frac{1}{(a+2n)^{m}}+\frac{1}{(a+3n)^{m}}-\text{etc.},\qquad(h'') \]
\ednote{Printed $-$ before etc.; the series has all terms positive, so $+$ etc. is expected. An apparent misprint.}
laquelle deviendra $(1,n)^{m}$, $(2,n)^{m}$, $(3,n)^{m}$, etc., selon qu'on fera $a=1$, $2$, $3$, etc., $n$ étant constant: on aura suivant cette notation,
\[ \left.\begin{aligned} T'&=(p,n)^{2}-(p+q,n)^{2}\\ \tfrac{1}{2}T''&=(p,n)^{3}-(p+q,n)^{3}\\ \tfrac{1}{2\cdot 3}T'''&=(p,n)^{4}-(p+q,n)^{4}\\ &\text{etc.} \end{aligned}\right\}\qquad(i'') \]

\subsection*{(42).}

Il faut observer que $n$ et $m$ restant les mêmes, on n'aura besoin de considérer les valeurs de $(a,n)^{m}$, que depuis $a=1$ jusqu'à $a=n$; car il est visible, par exemple, que $(n+1,n)^{m}$ se réduit à $(1,n)^{m}-1$, et qu'en général on a
\[ (a+n,n)^{m}=(a,n)^{m}-\frac{1}{a^{m}}.\qquad(k'') \]
Par suite de cette formule, on aurait de même
\[ (a+2n,n)^{m}=(a,n)^{m}-\frac{1}{a^{m}}-\frac{1}{(a+n)^{m}}, \]

\origpage{263}
et en général on réduira toute quantité $(a,n)^{m}$, où $a$ est plus grand que $n$, à une quantité semblable où $a$ n'excédera pas $n$.

On voit encore que $(n,n)^{m}$ représentant la suite
\[ \frac{1}{n^{m}}+\frac{1}{(2n)^{m}}+\frac{1}{(3n)^{m}}+\text{etc.}, \]
cette quantité est la même chose que
\[ \frac{1}{n^{m}}\cdot\left(1+\frac{1}{2^{m}}+\frac{1}{3^{m}}+\text{etc.}\right). \]
Si donc on désigne par $S_{m}$ la somme des puissances de degré $-m$, des nombres naturels, ensorte qu'on ait
\[ S_{m}=1+\frac{1}{2^{m}}+\frac{1}{3^{m}}+\frac{1}{4^{m}}+\frac{1}{5^{m}}+\text{etc.},\qquad(l'') \]
on aura $(n,n)^{m}=\frac{1}{n^{m}}S_{m}$. Enfin il est visible qu'on aura aussi l'équation
\[ S_{m}=(1,n)^{m}+(2,n)^{m}\ldots+(n,n)^{m},\qquad(m'') \]
laquelle, en substituant la valeur de $(n,n)^{m}$, devient
\[ \left(1-\frac{1}{n^{m}}\right)S_{m}=(1,n)^{m}+(2,n)^{m}+(3,n)^{m}+\ldots+(n-1,n)^{m}. \]

\subsection*{(43).}

Considérons particulièrement le cas de $n=2$, $p=1$, $q=1$, alors on aura
\[ \begin{aligned} T'&=(1,2)^{2}-(2,2)^{2}=\left(1-\frac{2}{2^{2}}\right)S_{2}=\frac{1}{2}S_{2}\\ T''&=(1,2)^{3}-(2,2)^{3}=\left(1-\frac{2}{2^{3}}\right)S_{3}=\frac{3}{4}S_{3}\\ T'''&=(1,2)^{4}-(2,2)^{4}=\left(1-\frac{2}{2^{4}}\right)S_{4}=\frac{7}{8}S_{4}.\\ &\text{etc.} \end{aligned} \]
\ednote{Printed $T''$ and $T'''$ in the second and third rows; by $(i'')$ these values are $\tfrac{1}{2}T''$ and $\tfrac{1}{2\cdot 3}T'''$.}
On sait que les quantités $S_{2}$, $S_{4}$, $S_{6}$, etc. sont connues en fonctions de $\pi$, et qu'on a $S_{2}=\frac{1}{6}\pi^{2}$, $S_{4}=\frac{1}{90}\pi^{4}$, $S_{6}=\frac{1}{945}\pi^{6}$, etc. A l'égard des quantités $S_{3}$, $S_{5}$, etc., ce sont des transcendantes particulières qui ne se rattachent point aux autres transcendantes connues; il est facile néanmoins d'en trouver les valeurs avec une grande approximation,

\origpage{264}
par les belles méthodes qu'Euler a données pour cet objet (\emph{Calc.\ diff.}, pag.\ 451 et suiv.).

Au moyen de ces diverses quantités, on connaîtra donc les intégrales suivantes, déduites des équations de l'art.\ 39.
\[ \begin{aligned} \int\frac{dx}{\sqrt{(1-xx)}}&=\frac{\pi}{2}\\ \int\frac{dx\log\frac{1}{x}}{\sqrt{(1-xx)}}&=\frac{\pi}{2}\log 2\\ \int\frac{dx\log^{2}\frac{1}{x}}{\sqrt{(1-xx)}}&=\frac{\pi}{2}\left(\log^{2}2+\frac{\pi^{2}}{12}\right)\\ \int\frac{dx\log^{3}\frac{1}{x}}{\sqrt{(1-xx)}}&=\frac{\pi}{2}\left(\log^{3}2+3\log 2\cdot\frac{\pi^{2}}{12}+\frac{3}{2}S_{3}\right); \end{aligned} \]
et on pourra prolonger indéfiniment cette suite où tout est connu, excepté $S_{3}$, $S_{5}$, etc. dont on connaît au moins les valeurs très-approchées, jusqu'à $S_{15}$ (\emph{Calc.\ diff.}, pag.\ 456.).

\section*{De la réduction des transcendantes désignées par $(a,n)^{m}$.}

\subsection*{(44).}

Occupons-nous maintenant de réduire au plus petit nombre possible les quantités $(1,n)^{m}$, $(2,n)^{m}$, $(3,n)^{m}$, etc. qui répondent à une même valeur de $n$. Pour cet effet, reprenons l'équation (t'), et substituons-y la valeur de $B$ donnée par la formule (c''), nous aurons
\[ \int\left(\frac{y^{a-1}-y^{n-a-1}}{1-y^{n}}\right)dy=\omega\cot a\omega.\qquad(n'') \]
D'où l'on tire, en différentiant successivement par rapport à $a$,
\[ \left.\begin{aligned} \int\frac{y^{a-1}+y^{n-a-1}}{1-y^{n}}\,dy\log\frac{1}{y}&=\frac{\omega^{2}}{\sin^{2}a\omega}\\ \frac{1}{2}\int\frac{y^{a-1}-y^{n-a-1}}{1-y^{n}}\,dy\log^{2}\frac{1}{y}&=\frac{\omega^{3}}{\sin^{3}a\omega}\cos a\omega\\ \frac{1}{2\cdot 3}\int\frac{y^{a-1}+y^{n-a-1}}{1-y^{n}}\,dy\log^{3}\frac{1}{y}&=\frac{\omega^{4}}{\sin^{4}a\omega}\left(\frac{1}{3}+\frac{2}{3}\cos^{2}a\omega\right)\\ &\text{etc.} \end{aligned}\right\}\qquad(p'') \]
Mettant au lieu des premiers membres les valeurs qu'ils obtiennent

\origpage{265}
par le développement en série et l'application de la formule (g''), on aura
\[ \left.\begin{aligned} (a,n)^{2}+(n-a,n)^{2}&=\frac{\omega^{2}}{\sin^{2}a\omega}\\ (a,n)^{3}+(n-a,n)^{3}&=\frac{\omega^{3}}{\sin^{3}a\omega}\cos a\omega\\ (a,n)^{4}+(n-a,n)^{4}&=\frac{\omega^{4}}{\sin^{4}a\omega}\left(\frac{1}{3}+\frac{2}{3}\cos^{2}a\omega\right)\\ &\text{etc.} \end{aligned}\right\}\qquad(q'') \]
\ednote{Second row printed $(a,n)^{3}+(n-a,n)^{3}$; the second row of (p'') and art.\ 47 require $(a,n)^{3}-(n-a,n)^{3}$.}
Ces formules serviront à établir entre les diverses quantités $(a,n)^{m}$, qui répondent à une même valeur de $n$, toutes les relations qu'on peut obtenir par d'autres voies, et qui sont données sous diverses formes dans les ouvrages d'Euler. C'est ce que nous allons développer dans deux exemples.

\subsection*{(45).}

Soit $n=6$ et $m=2$, on aura l'équation générale
\[ (a,n)^{2}+(6-a,6)^{2}=\frac{\omega^{2}}{\sin^{2}a\omega}, \]
\ednote{Printed $(a,n)^{2}$ in the first term; with $n=6$ this is $(a,6)^{2}$.}
d'où l'on déduit successivement
\[ \left.\begin{aligned} (1,6)^{2}+(5,6)^{2}&=\frac{\omega^{2}}{\sin^{2}\omega}\\ (2,6)^{2}+(4,6)^{2}&=\frac{\omega^{2}}{\sin^{2}2\omega}\\ (3,6)^{2}+(3,6)^{2}&=\frac{\omega^{2}}{\sin^{2}3\omega} \end{aligned}\right\}\qquad(r'') \]
ce qui donne en substituant la valeur $\omega=\frac{\pi}{6}$,
\[ \begin{aligned} (1,6)^{2}+(5,6)^{2}&=\frac{\pi^{2}}{9}\\ (2,6)^{2}+(4,6)^{2}&=\frac{\pi^{2}}{27}\\ (3,6)^{2}&=\frac{\pi^{2}}{72}. \end{aligned} \]
La somme de ces équations est
\[ \left(1-\frac{1}{36}\right)S_{2}=\frac{35}{216}\pi^{2}, \]
d'où l'on déduit $S_{2}=\frac{1}{6}\pi^{2}$, ce qui est un résultat connu, mais la

\origpage{266}
démonstration qu'on en donne ici, n'en est pas moins remarquable.

Outre cette valeur de $S_{2}$, on connaît $(3,6)^{2}=\frac{1}{72}\pi^{2}$, et $(6,6)^{2}=\frac{1}{36}S=\frac{1}{216}\pi^{2}$\ednote{Printed $S$ without its index at the line end; $S_{2}$ is meant.}; mais les quatre autres quantités $(1,6)^{2}$, $(2,6)^{2}$, $(4,6)^{2}$, $(5,6)^{2}$ ne peuvent être déterminées par les équations précédentes.

Observons cependant que puisqu'on a
\[ (2,6)^{2}=\frac{1}{2^{2}}+\frac{1}{8^{2}}+\frac{1}{14^{2}}+\frac{1}{20^{2}}+\text{etc.}, \]
il en résulte
\[ (2,6)^{2}=\frac{1}{4}\left(1+\frac{1}{4^{2}}+\frac{1}{7^{2}}+\frac{1}{10^{2}}+\text{etc.}\right)=\frac{1}{4}(1,6)^{2}+\frac{1}{4}(4,6)^{2}; \]
donc les quatre quantités dont il s'agit, peuvent être déterminées au moyen de l'une d'entre elles; par exemple, au moyen de $(1,6)^{2}$, de la manière suivante
\[ \left.\begin{aligned} (1,6)^{2}&=\zeta\\ (2,6)^{2}&=\frac{1}{135}\pi^{2}+\frac{1}{5}\zeta\\ (4,6)^{2}&=\frac{4}{135}\pi^{2}-\frac{1}{5}\zeta\\ (5,6)^{2}&=\frac{1}{9}\pi^{2}-\zeta \end{aligned}\right\}\qquad(s'') \]

\subsection*{(46).}

Quant à la valeur absolue de $\zeta$, elle n'est déterminable exactement par aucune formule connue; mais on peut en trouver une valeur aussi approchée qu'on voudra par la méthode qu'Euler a donnée (\emph{Calc.\ diff.}, pag.\ 451).

Pour cela ayant fait
\[ s=\frac{1}{a^{2}}+\frac{1}{(a+n)^{2}}+\frac{1}{(a+2n)^{2}}+\ldots+\frac{1}{(a+nx)^{2}}, \]
on trouve en général
\[ \begin{aligned} s={}&C-\frac{1}{n}\cdot\frac{1}{a+nx}+\frac{1}{2}\cdot\frac{1}{(a+nx)^{2}}-\frac{A'n}{(a+nx)^{3}}+\frac{B'n^{3}}{(a+nx)^{5}}\\ &-\frac{C'n^{5}}{(a+nx)^{7}}+\frac{D'n^{7}}{(a+nx)^{9}}-\text{etc.} \end{aligned} \]
$A'$, $B'$, $C'$, $D'$, etc. étant la suite des nombres Bernoulliens.

Il résulte de cette formule que la somme de la suite prolongée à l'infini étant désignée par $(a,n)^{2}$, on a $(a,n)^{2}=C$; donc réciproquement on a

\origpage{267}
\[ \begin{aligned} (a,n)^{2}={}&s+\frac{1}{n}\cdot\frac{1}{a+nx}-\frac{1}{2}\cdot\frac{1}{(a+nx)^{2}}+\frac{A'n}{(a+nx)^{3}}\\ &-\frac{B'n^{3}}{(a+nx)^{5}}+\frac{C'n^{5}}{(a+nx)^{7}}-\text{etc.}\qquad(t'') \end{aligned} \]
On sait que la suite $A'$, $B'$, $C'$, $D'$ est divergente à compter du troisième terme, et le devient plus que toute progression géométrique donnée; d'où il suit que la suite contenue dans la formule (t''), deviendra nécessairement divergente après un certain nombre de termes. Mais ce qui est fort remarquable, c'est que cette formule n'en est pas moins propre à donner la valeur de $(a,n)^{2}$ avec tout le degré d'approximation qu'on peut desirer.

Pour cela il faut donner à $x$ une valeur arbitraire d'autant plus grande qu'on voudra obtenir une plus grande approximation (la valeur $x=10$ suffit pour donner 18 ou 20 décimales exactes). Au moyen de cette valeur, on commencera par prendre la somme effective de la suite
\[ s=\frac{1}{a^{2}}+\frac{1}{(a+n)^{2}}+\frac{1}{(a+2n)^{2}}\ldots+\frac{1}{(a+xn)^{2}}; \]
substituant ensuite cette valeur de $s$ ainsi que celle de $x$ dans l'équation (t''), on aura pour la valeur de $(a,n)^{2}$ une suite d'abord très-convergente, mais dont la convergence diminuera de plus en plus, jusqu'à un certain terme où elle deviendra divergente, et cette divergence augmenterait de plus en plus à l'infini.

Par le calcul des termes successifs, on obtiendra des résultats alternativement plus grands et plus petits que la valeur cherchée, et on devra s'arrêter aux termes où cesse la convergence.

Ces termes indiqueront deux limites fort rapprochées, entre lesquelles se trouve nécessairement la valeur de $a_{2}$\ednote{Printed $a_{2}$; the quantity sought is $(a,n)^{2}$.}.

Si ces deux limites ne donnaient pas encore une approximation suffisante, il ne resterait d'autre parti à prendre que de recommencer un nouveau calcul en donnant à $x$ une valeur plus grande. Mais pour l'ordinaire, une valeur médiocrement grande de $x$ donnera une très-grande approximation.

Nous donnerons ci-après un exemple du calcul de ces sortes de suites qu'on peut appeler \emph{suites demi-convergentes}.

\origpage{268}

\subsection*{(47).}

Soit maintenant $m=3$ et $n=6$, nous aurons l'équation générale
\[ (a,6)^{3}-(6-a,6)^{3}=\frac{\omega^{3}\cos a\omega}{\sin^{3}a\omega}, \]
d'où l'on déduit les deux suivantes:
\[ \left.\begin{aligned} (1,6)^{3}-(5,6)^{3}&=\frac{\omega^{3}\cos\omega}{\sin^{3}\omega}=4\omega^{3}\sqrt{3}\\ (2,6)^{3}-(4,6)^{3}&=\frac{\omega^{3}\cos 2\omega}{\sin^{3}2\omega}=\frac{4\omega^{3}}{3\sqrt{3}} \end{aligned}\right\}\qquad(u'') \]
D'ailleurs on a $(6,6)^{3}=\frac{1}{216}S_{3}$, et
\[ (1,6)^{3}+(2,6)^{3}+(3,6)^{3}+(4,6)^{3}+(5,6)^{3}=\frac{215}{216}S_{3}. \]
Ces équations sont insuffisantes pour déterminer toutes les inconnues; mais les diviseurs de 6 qui sont 2 et 3, en fournissent de nouvelles.

On a en effet
\[ \begin{aligned} (2,6)^{3}&=\frac{1}{8}\left(1+\frac{1}{4^{3}}+\frac{1}{7^{3}}+\frac{1}{10^{3}}+\text{etc.}\right)=\frac{1}{8}[(1,6)^{3}+(4,6)^{3}]\\ (3,6)^{3}&=\frac{1}{27}\left(1+\frac{1}{3^{3}}+\frac{1}{5^{3}}+\frac{1}{7^{3}}+\text{etc.}\right)=\frac{1}{27}[(1,6)^{3}+(3,6)^{3}+(5,6)^{3}]; \end{aligned} \]
de là on voit que les cinq qnantités\ednote{Printed qnantités, a turned letter; read quantités.} $(1,6)^{3}$, $(2,6)^{3}$, $(3,6)^{3}$, $(4,6)^{3}$, $(5,6)^{3}$ se détermineront en supposant connue l'une d'entre elles. Ainsi en faisant $(1,6)^{3}=\xi$, on aura
\[ \begin{aligned} (1,6)^{3}&=\xi\\ (2,6)^{3}&=\frac{1}{7}\left(\xi-\frac{4\omega^{3}}{3\sqrt{3}}\right)\\ (3,6)^{3}&=\frac{1}{13}\left(\xi-2\omega^{3}\sqrt{3}\right)\\ (4,6)^{3}&=\frac{1}{7}\left(\xi-\frac{32\omega^{3}}{3\sqrt{3}}\right)\\ (5,6)^{3}&=\xi-4\omega^{3}\sqrt{3}. \end{aligned} \]
La somme de ces quantités doit être égale à $\frac{215}{216}S_{3}$, ainsi on aura
\[ \frac{215}{216}S_{3}=\frac{215}{91}\xi-\frac{430}{91}\omega^{3}\sqrt{3}, \]
ou
\[ S_{3}=\frac{216}{91}\xi-\frac{432}{91}\omega^{3}\sqrt{3}. \]

\origpage{269}

Réciproquement on peut exprimer toutes les quantités $(1,6)^{3}$, $(2,6)^{3}$, etc., au moyen de $S_{3}$ et de $\omega=\frac{\pi}{6}$, et on trouve
\[ \left.\begin{aligned} (1,6)^{3}&=\frac{91}{216}S_{3}+2\omega^{3}\sqrt{3}\\ (2,6)^{3}&=\frac{13}{216}S_{3}+\frac{2}{9}\omega^{3}\sqrt{3}\\ (3,6)^{3}&=\frac{7}{216}S_{3}\\ (4,6)^{3}&=\frac{13}{216}S_{3}-\frac{2}{9}\omega^{3}\sqrt{3}\\ (5,6)^{3}&=\frac{91}{216}S_{3}-2\omega^{3}\sqrt{3}\\ (6,6)^{3}&=\frac{1}{216}S_{3} \end{aligned}\right\}\qquad(v'') \]

\subsection*{(48).}

On déduirait aisément des équations (q'') les autres propriétés connues des quantités $(a,n)^{m}$, c'est-à-dire des suites qui résultent de la décomposition de la suite générale
\[ S_{m}=1+\frac{1}{2^{m}}+\frac{1}{3^{m}}+\frac{1}{4^{m}}+\frac{1}{5^{m}}+\text{etc.}, \]
en prenant les termes de trois en trois, de quatre en quatre, ou en général de $n$ en $n$. On trouverait, par exemple, que la suite
\[ 1-\frac{1}{3^{m}}+\frac{1}{5^{m}}-\frac{1}{7^{m}}+\text{etc.} \]
est sommable lorsque $m$ est impair, et qu'elle ne l'est pas lorsque $m$ est pair.

On trouverait au contraire que la suite
\[ 1+\frac{1}{3^{m}}+\frac{1}{5^{m}}+\frac{1}{7^{m}}+\text{etc.} \]
est sommable lorsque $m$ est pair, et qu'elle ne l'est pas lorsque $m$ est impair. Le mot \emph{sommable} est ici entendu non dans un sens absolu, mais relativement aux méthodes connues jusqu'à présent.

Ces choses n'ayant point de difficulté, et ayant d'ailleurs été démontrées par d'autres voies, nous ne nous y arrêterons pas davantage. Nous ferons voir seulement comment on peut trouver, par des suites qui soient convergentes dans toute leur étendue,

\origpage{270}
les valeurs des quantités $\zeta$ et $S_{3}$ qui sont restées inconnues dans les art.\ 45 et 47; recherche plus difficile qu'elle ne paraît au premier coup d'œil, parce qu'en suivant les méthodes qui se présentent naturellement, on retombe sur la même difficulté qu'on voulait résoudre.

\subsection*{(49).}

Pour obtenir d'abord la valeur de $S_{3}$, j'observe qu'on peut supposer
\[ \left.\begin{aligned} (1,6)^{3}&=\frac{1}{2}\int\frac{dy}{1-y^{6}}\log^{2}y\\ (2,6)^{3}&=\frac{1}{2}\int\frac{y\,dy}{1-y^{6}}\log^{2}y\\ (3,6)^{3}&=\frac{1}{2}\int\frac{y^{2}dy}{1-y^{6}}\log^{2}y\\ (4,6)^{3}&=\frac{1}{2}\int\frac{y^{3}dy}{1-y^{6}}\log^{2}y\\ (5,6)^{3}&=\frac{1}{2}\int\frac{y^{4}dy}{1-y^{6}}\log^{2}y\\ (6,6)^{3}&=\frac{1}{2}\int\frac{y^{5}dy}{1-y^{6}}\log^{2}y \end{aligned}\right\};\qquad(x'') \]
car les seconds membres étant intégrés depuis $y=0$ jusqu'à $y=1$, au moyen de la formule (g''), ces équations deviennent identiques.

Or d'après les équations (u'') on a
\[ \begin{aligned} (1,6)^{3}-(5,6)^{3}&=\frac{1}{2}\int\frac{1-y^{4}}{1-y^{6}}\,dy\log^{2}y=4\omega^{3}\sqrt{3},\\ (2,6)^{3}-(4,6)^{3}&=\frac{1}{2}\int\frac{y-y^{3}}{1-y^{6}}\,dy\log^{2}y=\frac{4\omega^{3}}{3\sqrt{3}}. \end{aligned} \]
Ces deux équations ajoutées ensemble, puis soustraites l'une de l'autre, donnent les deux suivantes
\[ \begin{aligned} \frac{1}{2}\int\frac{dy\log^{2}y}{1-y+y^{2}}&=4\omega^{3}\sqrt{3}+\frac{4\omega^{3}}{3\sqrt{3}}=\frac{40\omega^{3}\sqrt{3}}{9},\\ \frac{1}{2}\int\frac{dy\log^{2}y}{1+y+y^{2}}&=4\omega^{3}\sqrt{3}-\frac{4\omega^{3}}{3\sqrt{3}}=\frac{32\omega^{3}\sqrt{3}}{9}, \end{aligned} \]
de sorte qu'on a entre ces intégrales le rapport très-simple
\[ \int\frac{dy\log^{2}y}{1-y+y^{2}}=\frac{5}{4}\int\frac{dy\log^{2}y}{1+y+y^{2}}; \]

\origpage{271}
mais cette formule est susceptible d'être généralisée ainsi:
\[ \int\frac{dy\log^{m}y}{1-y+y^{2}}=\frac{2^{m}+1}{2^{m}}\int\frac{dy\log^{m}y}{1+y+y^{2}}.\qquad(y'') \]
En effet, soit
\[ P=\int\frac{dy\log^{m}y}{1-y+y^{2}},\quad\text{et}\quad Q=\int\frac{dy\log^{m}y}{1+y+y^{2}}, \]
on aura
\[ \frac{1}{2}(P-Q)=\int\frac{y\,dy\log^{m}y}{1+y^{2}+y^{4}}; \]
mettant dans cette dernière $y^{\frac{1}{2}}$ au lieu de $y$, il viendra
\[ \frac{1}{2}(P-Q)=\left(\frac{1}{2}\right)^{m+1}\int\frac{dy\log^{m}y}{1+y+y^{2}}=\left(\frac{1}{2}\right)^{M+1}Q, \]
\ednote{The second exponent is printed $M+1$ with a capital; $m+1$ is meant.}
et par conséquent $P=\frac{2^{m}+1}{2^{m}}Q$.

Par la combinaison des équations (x''), et la substitution des valeurs données par les équations (v''), on obtient
\[ \begin{aligned} (1,6)^{3}+(4,6)^{3}&=\frac{1}{2}\int\frac{dy\log^{2}y}{1-y^{3}}=\frac{104}{216}S_{3}+\frac{16}{9}+\omega^{3}\sqrt{3},\\ (3,6)^{3}+(6,6)^{3}&=\frac{1}{2}\int\frac{y^{2}dy\log^{2}y}{1-y^{3}}=\frac{8}{216}S_{3}. \end{aligned} \]
\ednote{In the first line the print sets a stray $+$ between $\frac{16}{9}$ and $\omega^{3}\sqrt{3}$; from the values $(v'')$ the term is $\frac{16}{9}\omega^{3}\sqrt{3}$.}
Soustrayant la seconde de la première, on aura
\[ \frac{1}{2}\int\frac{(1+y)\,dy\log^{2}y}{1+y+y^{2}}=\frac{4}{9}S_{3}+\frac{19}{9}\omega^{3}\sqrt{3}; \]
\ednote{Printed $\frac{19}{9}$; the subtraction gives $\frac{16}{9}$, which the elimination below requires. An apparent misprint.}
d'ailleurs on a déjà trouvé
\[ \frac{1}{2}\int\frac{dy\log^{2}y}{1+y+y^{2}}=\frac{32}{9}\omega^{3}\sqrt{3}; \]
donc en éliminant $\omega^{3}$, on aura
\[ \frac{8}{9}S_{3}=\int\frac{(y+\frac{1}{2})\,dy\log^{2}y}{1+y+y^{2}}.\qquad(z'') \]
C'est de cette formule que nous allons tirer la valeur de $S_{3}$.

\subsection*{(50).}

Je remarque d'abord qu'on peut faire
\[ \int\frac{(y+\frac{1}{2})\,dy\log^{2}y}{1+y+y^{2}}=\int\frac{dy}{1+y}\log^{2}y-\frac{1}{2}\int\frac{(1-y)\,dy\log^{2}y}{(1+y)(1+y+y^{2})}. \]

\origpage{272}
La première partie
\[ \int\frac{dy}{1+y}\log^{2}y=2\left(1-\frac{1}{2^{3}}+\frac{1}{3^{3}}-\frac{1}{4^{3}}+\text{etc.}\right)=2S_{3}\left(1-\frac{2}{8}\right)=\frac{3}{2}S_{3}, \]
donc on aura
\[ \frac{11}{18}S_{3}=\int\frac{\frac{1}{2}(1-y)\,dy\log^{2}y}{(1+y)(1+y+y^{2})}: \]
soit $y=\frac{1-z}{1+z}$, et la transformée sera
\[ \frac{11}{18}S_{3}=\int\frac{z\,dz}{3+z^{2}}\log^{2}y, \]
formule qui doit toujours être intégrée depuis $z=0$ jusqu'à $z=1$.

On voit maintenant que comme $z^{2}$ est toujours plus petit que 1, on pourra réduire $\dfrac{z\,dz}{3+z^{2}}$ en une série convergente, de sorte que l'intégrale ne dépendra plus que de termes de la forme \ldots\ $\displaystyle\int z^{2m+1}dz\log^{2}y$. Mais pour rendre la série encore plus convergente, je fais $z^{2}=1-u^{2}$, et j'ai
\[ \frac{11}{18}S_{3}=\int\frac{u\,du}{4-u^{2}}\log^{2}y, \]
intégrale qui doit être encore prise depuis $u=0$ jusqu'à $u=1$.

Cette intégrale étant prise par parties, devient
\[ -\frac{1}{2}\log\left(1-\frac{u^{2}}{4}\right)\cdot\log^{2}y+\int\frac{dy}{y}\log y\cdot\log\left(1-\frac{u^{2}}{4}\right); \]
la première partie s'évanouit aux deux limites de l'intégrale; ainsi on a simplement
\[ \frac{11}{18}S_{3}=\int\frac{dy}{y}\log y\log\left(1-\frac{u^{2}}{4}\right). \]
Développant $\log\left(1-\dfrac{u^{2}}{4}\right)$, et substituant la valeur $\dfrac{dy}{y}=-\dfrac{2dz}{u^{2}}$, on aura
\[ \frac{11}{18}S_{3}=\frac{1}{2}\int dz\log\frac{1}{y}\cdot\left(1+\frac{1}{2}\cdot\frac{u^{2}}{4}+\frac{1}{3}\cdot\frac{u^{4}}{4^{2}}+\frac{1}{4}\cdot\frac{u^{6}}{4^{3}}+\text{etc.}\right). \]

\origpage{273}
Soit donc
\[ \begin{aligned} \frac{1}{2}\int dz\log\frac{1}{y}&=P^{0},\\ \frac{1}{2}\int u^{2}dz\log\frac{1}{y}&=P',\\ \frac{1}{2}\int u^{4}dz\log\frac{1}{y}&=P'',\\ &\text{etc.} \end{aligned} \]
ces intégrales étant prises depuis $z=0$, jusqu'à $z=1$, et on aura
\[ \frac{11}{18}S_{3}=P^{0}+\frac{1}{2}\cdot\frac{P'}{4}+\frac{1}{3}\cdot\frac{P''}{4^{2}}+\frac{1}{4}\cdot\frac{P'''}{4^{3}}+\text{etc.} \]

\subsection*{(51).}

Pour avoir maintenant les quantités $P^{0}$, $P'$, $P''$, etc. j'observe qu'on a en général
\[ \int u^{2m}dz\log\frac{1}{y}=\log\frac{1}{y}\int u^{2m}dz+\int\frac{dy}{y}\cdot\int u^{2m}dz. \]
Dans cette formule nous supposerons que l'intégrale $\displaystyle\int u^{2m}dz$ est prise de manière qu'elle s'évanouisse lorsque $z=1$ ou $u=0$, alors la partie $\displaystyle\log\frac{1}{y}\int u^{2m}dz$ s'évanouit aux deux limites de l'intégrale, et on a simplement
\[ \frac{1}{2}\int u^{2m}dz\log\frac{1}{y}=\frac{1}{2}\int\frac{dy}{y}\int u^{2m}dz=\int-\frac{dz}{u^{2}}\int u^{2m}dz, \]
où l'on voit que les logarithmes ont entièrement disparu, et qu'on ne doit plus tenir compte que de la relation $u^{2}=1-z^{2}$.

On a d'ailleurs
\[ \int u^{2m}dz=\frac{2m}{2m+1}\int u^{2m-2}dz+\frac{u^{2m}z}{2m+1}; \]
multipliant de part et d'autre par $\dfrac{dz}{u^{2}}$ et intégrant, on aura
\[ P^{(m)}=\frac{2m}{2m+1}P^{(m-1)}-\int\frac{u^{2m-2}z\,dz}{2m+1}. \]
Or $z\,dz=-u\,du$, et par conséquent
\[ -\int\frac{u^{2m-2}z\,dz}{2m+1}=\int\frac{u^{2m-1}du}{2m+1}=\frac{u^{2m}}{(2m+1)2m}+C; \]

\origpage{274}
cette intégrale devant être prise depuis $z=0$, jusqu'à $z=1$, elle se réduit à $-\dfrac{1}{2m(2m+1)}$; donc on aura
\[ P^{(m)}=\frac{2m}{2m+1}\cdot P^{(m-1)}-\frac{1}{2m(2m+1)}. \]
La première valeur
\[ P^{0}=-\int\frac{dz}{u^{2}}\int dz=-\int\frac{dz}{u^{2}}(z-1)=\int\frac{dz}{1+z}=L(1+z)=L2; \]
donc on aura successivement
\[ \begin{aligned} P^{0}&=\log 2,\\ P'&=\frac{2}{3}P^{0}-\frac{1}{2\cdot 3},\\ P''&=\frac{4}{5}P'-\frac{1}{4\cdot 5},\\ P'''&=\frac{6}{7}P''-\frac{1}{6\cdot 7},\\ &\text{etc.} \end{aligned} \]
Au moyen de cette loi très-simple on calculera aisément les différens termes de la suite décroissante $P^{0}$, $P'$, $P''$, $P'''$, etc. Ensuite on aura $S_{3}$ par la formule
\[ S_{3}=\frac{18}{11}\left(P^{0}+\frac{1}{2}\cdot\frac{P'}{4}+\frac{1}{3}\cdot\frac{P''}{4^{2}}+\frac{1}{4}\cdot\frac{P'''}{4^{3}}+\text{etc.}\right); \]
ce qui donne une approximation très-rapide, puisque chaque terme est moindre que le quart du précédent.

En calculant cette suite jusqu'au terme $P^{\text{vi}}$ inclusivement, on trouve $S_{3}=1.2020567$. Euler a trouvé par la méthode dont nous avons parlé, et en poussant l'approximation beaucoup plus loin,
\[ S_{3}=1.202056903159594281. \]
(Voyez \emph{Calc.\ diff.}, page 453).

\subsection*{(52).}

Pour trouver par des procédés semblables la valeur de la transcendante $\zeta$, demeurée inconnue dans les équations (s''), j'observe qu'on a

\origpage{275}
\[ \begin{aligned} (1,6)^{2}&=\int\frac{dy}{1-y^{6}}\log\frac{1}{y}=\zeta,\\ (2,6)^{2}&=\int\frac{y\,dy}{1-y^{6}}\log\frac{1}{y}=\tfrac{1}{135}\pi^{2}+\tfrac{1}{5}\zeta,\\ (4,6)^{2}&=\int\frac{y^{3}\,dy}{1-y^{6}}\log\frac{1}{y}=\tfrac{4}{135}\pi^{2}-\tfrac{1}{5}\zeta,\\ (5,6)^{2}&=\int\frac{y^{4}\,dy}{1-y^{6}}\log\frac{1}{y}=\tfrac{1}{9}\pi^{2}-\zeta: \end{aligned} \]
de là je tire
\[ (1,6)^{2}+(2,6)^{2}-(4,6)^{2}-(5,6)^{2}=\int\frac{dy\log\frac{1}{y}}{1-y+y^{2}}=\frac{12}{5}\zeta-\frac{2}{15}\pi^{2}; \]
donc
\[ \zeta=\tfrac{1}{18}\pi^{2}+\frac{5}{12}\int\frac{dy\log\frac{1}{y}}{1-y+y^{2}}; \]
tout se réduit donc à trouver la valeur de cette intégrale.

Si on fait successivement $y=\frac{1}{2}(1-u)$, $y=\frac{1}{2}(1+u)$, et qu'on ajoute les deux transformées prises positivement, on aura
\[ \int\frac{dy\log\frac{1}{y}}{1-y+y^{2}}=\int\frac{2\,du}{3+u^{2}}\log\left(\frac{4}{1-u^{2}}\right), \]
nouvelle intégrale qui doit encore être prise depuis $u=0$ jusqu'à $u=1$.

La première partie de cette intégrale
\[ \int\frac{2\,du\log 4}{3+u^{2}}=\frac{2}{\sqrt{3}}\log 4\cdot\operatorname{arc}\operatorname{tang}\frac{u}{\sqrt{3}}, \]
et en faisant $u=1$, elle se réduit à
\[ \frac{2\log 4}{\sqrt{3}}\cdot\frac{\pi}{6}\quad\text{ou}\quad\frac{2\pi\log 2}{3\sqrt{3}}. \]
L'autre partie
\[ \int-\frac{2\,du}{3+u^{2}}\log(1-u^{2}) \]
étant appelée $T$, on aura par le développement de la fraction,
\[ T=\int-\frac{2\,du}{3}\left(1-\frac{u^{2}}{3}+\frac{u^{4}}{3^{2}}-\frac{u^{6}}{3^{3}}+\text{etc.}\right)\log(1-u^{2}): \]
or en intégrant par parties on a
\[ \int-u^{2m}\,du\log(1-u^{2})=\frac{1-u^{2m+1}}{2m+1}\log(1-u^{2})+\int\frac{1-u^{2m+1}}{2m+1}\cdot\frac{2u\,du}{1-u^{2}}, \]

\origpage{276}
et puisque la quantité $(1-u^{2m+1})\log(1-u^{2})$ s'évanouit aux deux limites de l'intégrale, on a simplement
\[ \begin{aligned} \int-u^{2m}\,du\log(1-u^{2})&=\frac{1}{2m+1}\int\frac{1-u^{2m+1}}{1-u^{2}}\cdot 2u\,du\\ &=\frac{2}{2m+1}\int\left(\frac{1-u^{2m+2}}{1-u^{2}}\,du-\frac{du}{1+u}\right). \end{aligned} \]
Développant la quantité sous le signe et intégrant depuis $u=0$ jusqu'à $u=1$, on aura
\[ \int-u^{2m}\,du\log(1-u^{2})=\frac{2}{2m+1}\left(1+\frac{1}{3}+\frac{1}{5}\ldots+\frac{1}{2m+1}-\log 2\right); \]
donc enfin
\[ \begin{aligned} T={}&\frac{4}{3}\cdot(1-\log 2)\\ &-\frac{4}{3^{2}}\cdot\frac{1}{3}\left(1+\frac{1}{3}-\log 2\right)\\ &+\frac{4}{3^{3}}\cdot\frac{1}{5}\left(1+\frac{1}{3}+\frac{1}{5}-\log 2\right)\\ &-\frac{4}{3^{4}}\cdot\frac{1}{7}\left(1+\frac{1}{3}+\frac{1}{5}+\frac{1}{7}-\log 2\right)\\ &+\text{etc.}, \end{aligned} \]
série convergente, puisque chaque terme est moindre que le tiers du précédent. $T$ étant connu, on aura $\zeta$ par la valeur
\[ \zeta=\frac{\pi^{2}}{18}+\frac{5\pi\log 2}{18\sqrt{3}}+\frac{5}{12}T. \]

\section*{Des intégrales Eulériennes de la seconde espèce.}

\subsection*{(53).}

Considérons maintenant l'intégrale
\[ \int x^{m-1}\,dx\left(\log\frac{1}{x}\right)^{n}, \]
que nous supposerons toujours prise entre les limites $x=0$, $x=1$. On a d'abord en intégrant par parties,
\[ \int x^{m-1}\,dx\left(\log\frac{1}{x}\right)^{n}=\frac{x^{m}}{m}\left(\log\frac{1}{x}\right)^{n}+\frac{n}{m}\int x^{m-1}\,dx\left(\log\frac{1}{x}\right)^{n-1}, \]
et comme la partie hors du signe s'évanouit aux deux limites,

\origpage{277}
pourvu qu'on suppose $n>0$, on aura alors
\[ \int x^{m-1}\,dx\left(\log\frac{1}{x}\right)^{n}=\frac{n}{m}\int x^{m-1}\,dx\left(\log\frac{1}{x}\right)^{n-1}.\qquad(\alpha) \]
On aura donc en général, si $n$ est un nombre entier positif,
\[ \int x^{m-1}\,dx\left(\log\frac{1}{x}\right)^{n}=\frac{n\cdot n-1\cdot n-2\cdot n-3\ldots 1}{m^{n+1}}.\qquad(\beta) \]
Lorsque $n$ ne sera pas un nombre entier, l'intégrale $\displaystyle\int x^{m-1}\,dx\left(\log\frac{1}{x}\right)^{n}$ sera en général une transcendante dont il convient d'examiner les propriétés.

Et d'abord au moyen de la formule $(\alpha)$, on pourra toujours ramener cette transcendante au cas où l'exposant $n$ est compris entre 0 et $-1$.

De plus, j'observe que sans rien diminuer de la généralité du calcul, on peut faire $m=1$; car la formule $\displaystyle\int x^{m-1}\,dx\left(\log\frac{1}{x}\right)^{n}$ étant proposée, si l'on fait $x^{m}=z$, cette formule deviendra $\displaystyle\frac{1}{m^{n+1}}\int dz\left(\log\frac{1}{z}\right)^{n}$.

\subsection*{(54).}

Cela posé, il suffira de considérer l'intégrale $\displaystyle\int dx\left(\log\frac{1}{x}\right)^{a-1}$, dans laquelle nous supposerons que $a$ est positif et plus petit que l'unité. Cette quantité étant simplement fonction de $a$, nous la désignerons par $\Gamma(a)$, et nous ferons
\[ \Gamma(a)=\int dx\left(\log\frac{1}{x}\right)^{a-1}.\qquad(\gamma) \]
L'objet des recherches suivantes est d'évaluer la fonction $\Gamma(a)$, lorsque $a$ est une fraction rationnelle donnée, telle que $\frac{1}{2}$, $\frac{2}{3}$, $\frac{3}{4}$, etc., et nous nous proposons particulièrement de comparer entre elles les fonctions qui répondent à des valeurs de $a$ de même dénomination, telles que $\Gamma\left(\frac{1}{n}\right)$, $\Gamma\left(\frac{2}{n}\right)$, etc. Enfin nous chercherons aussi à déterminer par approximation la transcendante $\Gamma(a)$ pour toute valeur de $a$ rationnelle ou irrationnelle.

\subsection*{(55).}

En prenant les intégrales depuis $x=0$ jusqu'à $x=1$, et

\origpage{278}
supposant $m>1$, on a cette formule de réduction,
\[ \int x^{\alpha-1}\,dx(1-x^{\beta})^{m-1}=\frac{(m-1)\beta}{\alpha+(m-1)\beta}\int x^{\alpha-1}\,dx(1-x^{\beta})^{m-2}; \]
d'où il suit que si $m$ est un entier, on aura exactement
\[ \int x^{\alpha-1}\,dx(1-x^{\beta})^{m-1}=\frac{1\cdot 2\cdot 3\ldots(m-1)\beta^{m-1}}{\alpha\cdot\alpha+\beta\cdot\alpha+2\beta\ldots\alpha+(m-1)\beta}. \]

Désignons par $\Phi(\alpha,m)$ l'intégrale $\displaystyle\int x^{\alpha-1}\,dx(1-x^{\beta})^{m-1}$, prise depuis $x=0$ jusqu'à $x=1$, et dans l'hypothèse que $m$ est un entier positif, on aura donc
\[ \frac{1\cdot 2\cdot 3\ldots(m-1)}{\alpha+\beta\cdot\alpha+2\beta\ldots\alpha+(m-1)\beta}=\frac{\alpha}{\beta^{m-1}}\Phi(\alpha,m). \]
Dans cette équation, mettons successivement $\lambda m$ et $(\lambda+1)m$ à la place de $m$, nous aurons
\[ \begin{gathered} \frac{1\cdot 2\cdot 3\ldots\overline{\lambda m-1}}{\alpha+\beta\cdot\alpha+2\beta\ldots\alpha+\overline{\lambda m-1}\cdot\beta}=\frac{\alpha}{\beta^{\lambda m-1}}\cdot\Phi(\alpha,\lambda m)\\ \frac{1\cdot 2\cdot 3\ldots\overline{\lambda m+m-1}}{\alpha+\beta\cdot\alpha+2\beta\ldots\alpha+\overline{\lambda m+m-1}\cdot\beta}=\frac{\alpha}{\beta^{\lambda m+m-1}}\cdot\Phi(\alpha,\lambda m+m). \end{gathered} \]
Divisant la seconde par la première, et faisant pour abréger, $\alpha+\lambda m\beta=\alpha'$, il vient
\[ \frac{\alpha'\cdot\alpha'+\beta\cdot\alpha'+2\beta\ldots\alpha'+\overline{m-1}\cdot\beta}{\lambda m\cdot\lambda m+1\cdot\lambda m+2\ldots\lambda m+m-1}=\beta^{m}\cdot\frac{\Phi(\alpha,\lambda m)}{\Phi(\alpha,\lambda m+m)}; \]
mais en mettant $\alpha'$ à la place de $\alpha$ dans l'équation primitive, on a
\[ \frac{1\cdot 2\cdot 3\ldots m-1}{\alpha'+\beta\cdot\alpha'+2\beta\ldots\alpha'+\overline{m-1}\cdot\beta}=\frac{\alpha'}{\beta^{m-1}}\Phi(\alpha+\lambda m\beta,m), \]
et multipliant ces deux équations entre elles, on a pour produit
\[ \frac{1\cdot 2\cdot 3\ldots m-1}{\lambda m\cdot\lambda m+1\ldots\lambda m+m-1}=\beta\cdot\frac{\Phi(\alpha,\lambda m)\cdot\Phi(\alpha+\lambda m\beta,m)}{\Phi(\alpha,\lambda m+m)}. \]
Le premier membre, en vertu de l'équation primitive, peut être représenté par $\displaystyle\int x^{\lambda m-1}\,dx(1-x)^{m-1}$; et parce qu'on peut mettre $x^{n}$ à la place de $x$, sans changer les limites de l'intégrale, il peut être

\origpage{279}
aussi représenté par $\displaystyle n\int x^{\lambda nm-1}\,dx(1-x^{n})^{m-1}$; donc on a
\[ n\int x^{\lambda mn-1}\,dx(1-x^{n})^{m-1}=\beta\frac{\Phi(\alpha,\lambda m)\cdot\Phi(\alpha+\lambda m\beta,m)}{\Phi(\alpha,\lambda m+m)}.\qquad(\delta) \]
Cette équation ainsi exprimée en un nombre fini de termes, acquiert une plus grande généralité, et ne suppose plus que $m$ est un nombre entier.

En effet, les deux membres devant se réduire à une même fonction de $m$ et de $\lambda m$, laquelle est
\[ \frac{1}{\lambda m}-\frac{m-1}{1}\cdot\frac{1}{\lambda m+1}+\frac{m-1\cdot m-2}{1\cdot 2}\cdot\frac{1}{\lambda m+2}-\text{etc.}, \]
on est maître de donner à $m$ et $\lambda$ des valeurs positives quelconques, et à plus forte raison aux quantités $\alpha$, $\beta$, $n$, qui disparaissent dans les deux membres.

\subsection*{(56).}

Soit donc $\alpha=1$ et $\beta=$ à un infiniment petit, on aura
\[ 1-x^{\beta}=\beta\log\frac{1}{x},\quad\text{et}\quad\Phi(\alpha,m)=\beta^{m-1}\int dx\left(l\,\frac{1}{x}\right)^{m-1}, \]
de sorte qu'on aura en général
\[ \Phi(\alpha,k)=\beta^{k-1}\cdot\Gamma(k). \]
Au moyen de cette formule, l'équation $(\delta)$ devient
\[ n\int x^{\lambda mn-1}\,dx(1-x^{n})^{m-1}=\frac{\Gamma(\lambda m)\cdot\Gamma(m)}{\Gamma(\lambda m+m)}. \]
Soit maintenant $m=\frac{q}{n}$, $\lambda m=\frac{p}{n}$, on aura donc
\[ \int x^{p-1}\,dx(1-x^{n})^{\frac{q}{n}-1}=\frac{\Gamma\left(\frac{q}{n}\right)\cdot\Gamma\left(\frac{q}{n}\right)}{n\Gamma\left(\frac{p+q}{n}\right)}. \]
\ednote{The first factor of the numerator is printed $\Gamma\left(\frac{q}{n}\right)$ (confirmed by magnification); the sense requires $\Gamma\left(\frac{p}{n}\right)$, as in equation $(\varepsilon)$ below. A printer's misprint.}
Le premier membre n'est autre chose que la transcendante désignée ci-dessus par $\left(\frac{p}{q}\right)$; ainsi on aura cette équation remarquable
\[ \left(\frac{p}{q}\right)=\frac{\Gamma\left(\frac{p}{n}\right)\cdot\Gamma\left(\frac{q}{n}\right)}{n\Gamma\left(\frac{p+q}{n}\right)}.\qquad(\varepsilon) \]

\origpage{280}
D'où l'on voit que la transcendante $\left(\frac{p}{q}\right)$ serait connue, si on connaissait, pour la même valeur de $n$, les fonctions de la forme $\Gamma\left(\frac{a}{n}\right)$, $a$ étant entier.

\subsection*{(57).}

Il résulte d'abord de cette valeur de $\left(\frac{p}{q}\right)$, qu'on peut échanger entre eux les nombres $p$ et $q$, et qu'ainsi on a $\left(\frac{p}{q}\right)=\left(\frac{q}{p}\right)$, ce qui est une des principales propriétés de ces fonctions.

De plus on tire de cette formule
\[ \left(\frac{p}{q}\right)\cdot\left(\frac{p+q}{r}\right)=\frac{\Gamma\left(\frac{p}{n}\right)\cdot\Gamma\left(\frac{q}{n}\right)\Gamma\left(\frac{r}{n}\right)}{n^{2}\Gamma\left(\frac{p+q+r}{n}\right)} \]
Dans le second membre, il est visible qu'on peut faire la permutation entre deux des nombres $p$, $q$, $r$, à volonté, ce qui donne le théorème fondamental
\[ \left(\frac{p}{q}\right)\cdot\left(\frac{p+q}{r}\right)=\left(\frac{p}{r}\right)\cdot\left(\frac{p+r}{q}\right); \]
dont on a ainsi une nouvelle démonstration très-simple.

\subsection*{(58).}

Faisons voir maintenant comment les fonctions $\Gamma$ se déterminent au moyen des fonctions $\left(\frac{p}{q}\right)$.

Observons d'abord qu'au moyen de l'équation $(\alpha)$, on a en général
\[ \Gamma(n+1)=n\Gamma(n),\qquad(\zeta) \]
ce qui permettra de réduire les cas où $n$ est plus grand que l'unité à ceux où il est plus petit.

Si l'on a $n=1$, alors $\Gamma(n)$ se réduit à $\displaystyle\int dx=x=1$; ainsi on a
\[ \Gamma(1)=1.\qquad(\eta) \]
Cela posé, faisons $q=n-p$ dans l'équation $(\varepsilon)$, alors la valeur du premier membre est connue, et on a
\[ \frac{\omega}{\sin p\omega}=\frac{\Gamma\left(\frac{p}{n}\right)\Gamma\left(\frac{n-p}{n}\right)}{n\Gamma(1)}, \]

\origpage{281}
ou, en d'autres termes,
\[ \Gamma(a)\Gamma(1-a)=\frac{\pi}{\sin a\pi}.\qquad(\theta) \]
Lorsque $a=\frac{1}{2}$, on a $(\Gamma a)^{2}=\pi$; donc
\[ \Gamma\left(\tfrac{1}{2}\right)=\sqrt{\pi}.\qquad(\iota) \]

L'équation $(\theta)$, très-remarquable dans cette théorie, fait voir qu'il suffit de connaître la valeur de $\Gamma(x)$ depuis $x=\frac{1}{2}$ jusqu'à $x=1$, et on en déduira les autres valeurs de cette fonction depuis $x=\frac{1}{2}$ jusqu'à $x=0$. Au reste, la valeur de $\Gamma(x)$ depuis $x=\frac{1}{2}$ jusqu'à $x=1$, ne varie qu'entre les limites $\sqrt{\pi}$ et 1, ou $1{,}77245$ et 1, tandis que depuis $x=\frac{1}{2}$ jusqu'à $x=0$, elle varie depuis $1{,}77245$ jusqu'à l'infini; et en particulier lorsque $a$ est infiniment petit, la formule $(\theta)$ donne $\Gamma(a)=\frac{1}{a}$.

\subsection*{(59).}

Si dans l'équation $(\varepsilon)$ on fait $p=1$, et qu'on prenne successivement $q=1$, 2, 3\ldots\ jusqu'à $n=1$\ednote{Printed $n=1$; the sense requires $n-1$, the last equation of the series below being that for $\left(\frac{1}{n-1}\right)$. An apparent misprint.}, on aura une suite d'équations
\[ \begin{gathered} \left(\frac{1}{1}\right)=\frac{\Gamma\left(\frac{1}{n}\right)\Gamma\left(\frac{1}{n}\right)}{n\Gamma\left(\frac{2}{n}\right)}\\ \left(\frac{1}{2}\right)=\frac{\Gamma\left(\frac{1}{n}\right)\Gamma\left(\frac{2}{n}\right)}{n\Gamma\left(\frac{3}{n}\right)}\\ \left(\frac{1}{3}\right)=\frac{\Gamma\left(\frac{1}{n}\right)\Gamma\left(\frac{3}{n}\right)}{n\Gamma\left(\frac{4}{n}\right)}\\ \vdots\\ \left(\frac{1}{n-1}\right)=\frac{\Gamma\left(\frac{1}{n}\right)\Gamma\left(\frac{n-1}{n}\right)}{n\Gamma\left(\frac{n}{n}\right)}. \end{gathered} \]
Multipliant les $a-1$ premières équations entre elles, on aura le produit
\[ \left(\frac{1}{1}\right)\cdot\left(\frac{1}{2}\right)\cdot\left(\frac{1}{3}\right)\cdot\left(\frac{1}{4}\right)\ldots\left(\frac{1}{a-1}\right)=\frac{\left(\Gamma\frac{1}{n}\right)^{a}}{n^{a-1}\Gamma\left(\frac{a}{n}\right)}, \]

\origpage{282}
si on les multiplie toutes, ou qu'on fasse $a=n$, le produit donnera
\[ \Gamma\left(\frac{1}{n}\right)=n\sqrt[n]{\left[\frac{1}{n}\cdot\left(\frac{1}{1}\right)\cdot\left(\frac{1}{2}\right)\cdot\left(\frac{1}{3}\right)\ldots\left(\frac{1}{n-1}\right)\right]}. \]
Soit donc pour abréger
\[ T=\sqrt[n]{\left[\frac{1}{n}\cdot\left(\frac{1}{1}\right)\cdot\left(\frac{1}{2}\right)\cdot\left(\frac{1}{3}\right)\ldots\left(\frac{1}{n-1}\right)\right]}, \]
et on aura
\[ \Gamma\left(\frac{1}{n}\right)=nT. \]
Connaissant $\Gamma\left(\frac{1}{n}\right)$, on en déduira $\Gamma\left(\frac{a}{n}\right)$ par l'équation déjà trouvée, qui donne
\[ \Gamma\left(\frac{a}{n}\right)=\frac{nT^{a}}{\left(\frac{1}{1}\right)\left(\frac{1}{2}\right)\left(\frac{1}{3}\right)\ldots\left(\frac{1}{a-1}\right)}=\frac{\left(\frac{1}{a}\right)\left(\frac{1}{n+a}\right)\ldots\left(\frac{1}{n-1}\right)}{T^{n-a}}.\qquad(\kappa) \]
\ednote{Printed $\left(\frac{1}{n+a}\right)$ as the second factor; the sense requires $\left(\frac{1}{a+1}\right)$. An apparent misprint.}
Et comme les quantités $\left(\frac{1}{1}\right)$, $\left(\frac{1}{2}\right)$, $\left(\frac{1}{3}\right)$, etc. sont censées connues pour chaque valeur de $n$, il ne reste plus rien à desirer pour la détermination des fonctions $\Gamma\left(\frac{a}{n}\right)$.

\subsection*{(60).}

Réciproquement, si on connaissait la valeur de $\Gamma(x)$ pour toute valeur rationnelle de $x$, moindre que l'unité, il serait facile de déterminer l'intégrale $\left(\frac{p}{q}\right)$; car on a en général
\[ \left(\frac{p}{q}\right)=\frac{\Gamma\left(\frac{p}{n}\right)\Gamma\left(\frac{q}{n}\right)}{n\Gamma\left(\frac{p+q}{n}\right)};\qquad(\lambda) \]
et s'il arrive que $p+q$ soit $>n$, on changera, d'après l'équation $(\zeta)$, cette formule en celle-ci
\[ \left(\frac{p}{q}\right)=\frac{\Gamma\left(\frac{p}{n}\right)\Gamma\left(\frac{q}{n}\right)}{(p+q-n)\Gamma\left(\frac{p+q-n}{n}\right)}.\qquad(\lambda') \]
Nous remarquerons que ces formules sont propres à donner l'expression générale de $\left(\frac{p}{q}\right)$ au moyen des quantités de la même

\origpage{283}
espèce $\left(\frac{1}{1}\right)$, $\left(\frac{1}{2}\right)$, $\left(\frac{1}{3}\right)$, etc. En effet les valeurs de $\Gamma\left(\frac{p}{n}\right)$, $\Gamma\left(\frac{q}{n}\right)$, $\Gamma\left(\frac{p+q}{n}\right)$ étant tirées de la formule $(\kappa)$, on en déduira
\[ \left(\frac{p}{q}\right)=\frac{\left(\frac{1}{p}\right)\left(\frac{1}{p+1}\right)\left(\frac{1}{p+2}\right)\ldots\left(\frac{1}{p+q-1}\right)}{\left(\frac{1}{1}\right)\left(\frac{1}{2}\right)\left(\frac{1}{3}\right)\ldots\left(\frac{1}{q-1}\right)},\qquad(\mu) \]
formule qui servira tant que $p+q$ sera plus petit que $n$.

Si on a $p+q>n$, il faudra faire usage de la seconde formule, qui donnera par de semblables substitutions,
\[ \left(\frac{p}{q}\right)=\frac{1}{p+q-n}\cdot\frac{\left(\frac{1}{p}\right)\left(\frac{1}{p+1}\right)\left(\frac{1}{p+2}\right)\ldots\left(\frac{1}{n-1}\right)}{\left(\frac{1}{p+q-n}\right)\left(\frac{1}{p+q-n+1}\right)\ldots\left(\frac{1}{q-1}\right)},\qquad(\mu') \]
Ces formules répondent aux équations (k) et (n) trouvées ci-dessus, et on les ferait coïncider entièrement en substituant, au lieu de chaque quantité $\left(\frac{1}{a}\right)$, sa valeur
\[ \left(\frac{1}{a}\right)=\frac{A_{a}\sin(a+1)\omega}{\sin\omega}; \]
ainsi les fonctions $\Gamma$ offrent un nouveau moyen direct et très-simple de déterminer l'expression générale des quantités $\left(\frac{p}{q}\right)$.

\subsection*{(61).}

Pour revenir aux quantités $\Gamma(a)$, nous avons déjà trouvé l'équation
\[ \Gamma(a)\cdot\Gamma(1-a)=\frac{\pi}{\sin a\pi}, \]
au moyen de laquelle les valeurs de la fonction, depuis $a=0$ jusqu'à $a=\frac{1}{2}$, se déduisent des valeurs supposées connues, depuis $a=\frac{1}{2}$ jusqu'à $a=1$.

Nous allons prouver ultérieurement qu'il suffit de connaître les valeurs de la fonction dans la moitié de cet intervalle, c'est-à-dire seulement depuis $a=\frac{3}{4}$ jusqu'à $a=1$, et on en déduira toutes les autres valeurs.

En effet, si on suppose $a<\frac{1}{2}n$, l'équation $(\varepsilon)$ donne tout à la\ednote{The print breaks off here: p.~284 opens directly with the display, so a word (presumably ``fois'') is omitted. Kept as printed.}

\origpage{284}
\[ \begin{gathered} \left(\frac{a}{a}\right)=\frac{\Gamma\left(\frac{a}{n}\right)\Gamma\left(\frac{a}{n}\right)}{n\Gamma\left(\frac{2a}{n}\right)}\\ \left(\frac{\tfrac{1}{2}n-a}{\tfrac{1}{2}n-a}\right)=\frac{\Gamma\left(\frac{\tfrac{1}{2}n-a}{n}\right)\Gamma\left(\frac{\tfrac{1}{2}n-a}{n}\right)}{n\Gamma\left(\frac{n-2a}{n}\right)}. \end{gathered} \]
Substituant ces valeurs dans l'équation $(d')$, puis mettant simplement $a$ au lieu de $\frac{a}{n}$, et réduisant les fonctions d'après la formule $(\theta)$, on aura
\[ \Gamma(a)=\frac{2^{1-2a}}{\sqrt{\pi}}\cos a\pi\cdot\Gamma(2a)\cdot\Gamma\left(\tfrac{1}{2}-a\right), \]
équation qui suppose $a<\frac{1}{2}$.

Cette équation combinée avec l'équation $(\theta)$, donnera
\[ \Gamma(1-a)=\frac{2^{2a}}{\sqrt{\pi}}\cos a\pi\,\Gamma(1-2a)\Gamma\left(\tfrac{1}{2}+a\right); \]
enfin de celle-ci on déduit, en mettant $a-\frac{1}{2}$ au lieu de $a$,
\[ \Gamma(a)=\frac{2^{1-2a}\sqrt{\pi}}{\sin a\pi}\cdot\frac{\Gamma\left(\tfrac{3}{2}-a\right)}{\Gamma(2-2a)}\qquad(\nu) \]
Nous supposons connues les valeurs de $\Gamma(a)$ depuis $a=\frac{3}{4}$ jusqu'à $a=1$.

Cela posé, 1°. Le second membre de l'équation $(\nu)$ sera connu pour toute valeur de $a$, depuis $a=\frac{1}{2}$ jusqu'à $a=\frac{5}{8}$; donc on connaîtra $\Gamma(a)$ dans ce même intervalle depuis $a=\frac{1}{2}$ jusqu'à $a=\frac{5}{8}$.

2°. Au moyen de ce premier cas, le second membre sera connu si $2-2a$ est compris entre $\frac{1}{2}$ et $\frac{5}{8}$; on connaîtra donc $\Gamma(a)$ toutes les fois que $a$ est compris entre $\frac{11}{16}$ et $\frac{12}{16}$.

3°. Au moyen des deux premiers cas, le second membre de l'équation $(\nu)$ sera connu si $2-2a$ est compris entre $\frac{11}{16}$ et $\frac{12}{16}$; donc on connaîtra $\Gamma(a)$ pour toutes les valeurs de $a$ comprises depuis $a=\frac{5}{8}=\frac{20}{32}$ jusqu'à $a=\frac{21}{32}$.

4°. Le second membre sera encore connu si $2-2a$ est compris entre $\frac{20}{32}$ et $\frac{21}{32}$, donc $\Gamma(a)$ sera connu depuis $a=\frac{44}{64}=\frac{11}{16}$ jusqu'à $a=\frac{43}{64}$, et ainsi de suite.

Par ces diverses opérations les valeurs de $a$ pour lesquelles $\Gamma(a)$

\origpage{285}
devient connu, se rapprochent alternativement de la limite $\frac{2}{3}$, qu'elles n'atteignent cependant qu'à l'infini, puisque $\frac{2}{3}$ ne peut pas s'exprimer exactement en fractions dont le dénominateur soit une puissance de 2. Mais on voit que par quatre opérations seulement, l'intervalle où $\Gamma(a)$ reste à déterminer, ne s'étend plus que depuis $a=\frac{44}{64}$ jusqu'à $a=\frac{43}{64}$. Une cinquième opération resserrerait cet intervalle de $\frac{43}{64}$ ou $\frac{86}{128}$ à $\frac{85}{128}$, et ainsi de suite.

La limite commune de ces suites est $\frac{2}{3}$ et $\Gamma\left(\frac{2}{3}\right)$ se détermine directement en faisant $a=\frac{2}{3}$ dans la formnle\ednote{Printed ``formnle'' for ``formule''; a turned letter.} $(\nu)$, ce qui donne
\[ \Gamma\left(\frac{2}{3}\right)^{2}=\frac{2^{-\frac{1}{3}}\sqrt{\pi}}{\sin\frac{2}{3}\pi}\Gamma\left(\frac{5}{6}\right). \]

\subsection*{(62).}

Dans les cas particuliers où l'on chercherait à déterminer une valeur de $\Gamma(a)$, on ne doit s'embarrasser aucunement de la distinction des cas précédens, et l'application immédiate de la formule $(\nu)$, répétée autant de fois qu'il est nécessaire, ou jusqu'à ce qu'il n'y ait plus d'inconnue à déterminer, conduira toujours au résultat qu'on cherche.

Soit proposé, par exemple, de trouver la valeur de $\Gamma(0.675)$, on aura directement
\[ \Gamma(0.675)=\frac{2^{-0.35}\sqrt{\pi}}{\sin 38^{\circ}30'}\cdot\frac{\Gamma(0.825)}{\Gamma(0.65)}. \]
\ednote{Read $\sin 38^{\circ}30'$ at scan resolution; formula $(\nu)$ with $a=0.675$ requires $\sin 121^{\circ}30'=\sin 58^{\circ}30'$. Either a misprint or a 5 misread as 3; not confirmed by magnification.}
Dans le second membre $\Gamma(0.65)$ est inconnue; pour la trouver, il faudra faire une seconde application de la formule, et on aura
\[ \Gamma(0.65)=\frac{2^{-0.30}\sqrt{\pi}}{\sin 63^{\circ}}\cdot\frac{\Gamma(0.85)}{\Gamma(0.70)}; \]
une troisième application donnera
\[ \Gamma(0.70)=\frac{2^{-0.40}\sqrt{\pi}}{\sin 54^{\circ}}\cdot\frac{\Gamma(0.8)}{\Gamma(0.6)}; \]
enfin une quatrième
\[ \Gamma(0.60)=\frac{2^{-0.20}\sqrt{\pi}}{\sin 72^{\circ}}\cdot\frac{\Gamma(0.9)}{\Gamma(0.8)}, \]
d'où en remontant on conclura la valeur de $\Gamma(0.675)$ exprimée en quantités connues. Cette détermination est un peu longue dans ce cas, parce que 0.675 approche beaucoup de la limite $\frac{2}{3}$.

\origpage{286}

\subsection*{(63).}

Puisqu'on a $\Gamma(1)=1$ et $\Gamma(1+a)=a\Gamma(a)$, la fonction $\Gamma(a)$ se détermine toujours exactement lorsque $a$ sera un nombre entier, et on aura généralement
\[ \Gamma(a)=1\cdot 2\cdot 3\ldots a-1; \]
d'où l'on voit que la quantité $\Gamma(a)$, considérée comme une fonction continue de $a$, n'est autre chose que le terme général qui résulterait de l'interpolation de la suite 1, 1, $1\cdot 2$, $1\cdot 2\cdot 3$, $1\cdot 2\cdot 3\cdot 4$, etc.

Les deux premiers termes de cette suite répondent aux valeurs $a=1$, $a=2$, et puisqu'ils sont tous deux égaux à l'unité, il s'ensuit que dans l'intervalle depuis $a=1$ jusqu'à $a=2$, la fonction $\Gamma(a)$ doit devenir \emph{maximum} ou \emph{minimum}.

On reconnaît aisément après quelques essais, que c'est le \emph{minimum} qui a lieu, et dans ce cas on trouve
\[ \begin{aligned} a&=1.4616038,\\ \Gamma(a)&=0.8856033,\\ \log\Gamma(a)&=9.9472392. \end{aligned} \]

Passé la valeur $a=2$, la fonction $\Gamma(a)$ augmente de plus en plus, puisqu'on a $\Gamma(2+a)=(1+a)\Gamma(1+a)$. Donc la valeur que nous venons de trouver est la plus petite de toutes celles que peut prendre la fonction $\Gamma(a)$ depuis $a=0$ jusqu'à $a=\infty$.

Il suit de là que la meilleure manière de construire une table des valeurs de $\Gamma(a)$ est de la calculer pour l'intervalle entre $a=1$ et $a=2$. Car dans cet intervalle la fonction ne varie qu'entre les limites 1 et 0.8856033, de sorte que les différences seront très-petites, et la table très-facile à interpoler.

\subsection*{(64).}

D'après cette observation, et pour faciliter l'usage des fonctions $\Gamma$, nous joignons ici une table des logarithmes de la fonction $\Gamma$, calculés pour toutes les valeurs de $a$, de millième en millième, depuis $a=1.000$, jusqu'à $a=2.000$.

Cette table est aisée à interpoler pour toutes les valeurs intermédiaires, et au moyen de l'équation $\Gamma(1+a)=a\Gamma(a)$, on étendra son usage à tous les autres cas.

En effet, si $a$ est compris entre 0 et 1, on déterminera $\Gamma(a)$ par l'équation $\Gamma(a)=\frac{1}{a}\Gamma(1+a)$.

\origpage{287}
Si $a$ est compris entre 2 et 3, $\Gamma(a)$ se déterminera par l'équation $\Gamma(a)=(a-1)\Gamma(a-1)$. Si $a$ est compris entre 3 et 4, $\Gamma(a)$ se déterminera par l'équation $\Gamma(a)=(a-1)(a-2)\Gamma(a-2)$, et ainsi de suite.

\subsection*{(65).}

Au moyen de cette table, il sera facile d'évaluer dans tous les cas la transcendante $\left(\frac{p}{q}\right)$ qui répond à une valeur donnée de $n$. Pour cet effet, les équations $(\lambda)$ étant mises sous ces deux nouvelles formes,
\[ \begin{aligned} \left(\frac{p}{q}\right)&=\frac{n}{pq}\cdot\frac{\Gamma\left(1+\frac{p}{n}\right)\cdot\Gamma\left(1+\frac{q}{n}\right)}{\Gamma\left(\frac{p+q}{n}\right)}\\ \left(\frac{p}{q}\right)&=\frac{p+q}{pq}\cdot\frac{\Gamma\left(1+\frac{p}{n}\right)\cdot\Gamma\left(1+\frac{q}{n}\right)}{\Gamma\left(1+\frac{p+q}{n}\right)}, \end{aligned} \]
on fera usage de la première, si l'on a $p+q>n$, et de la seconde si l'on a $p+q<n$.

Soit proposé, par exemple, de trouver la valeur de la transcendante $Z=\left(\frac{3}{8}\right)$ lorsque $a=10$\ednote{Printed $a=10$ instead of $n=10$; an apparent misprint. A reader has struck the letter and written n in the margin.}; on se servira alors de la première formule, qui donne
\[ Z=\frac{\Gamma(1.300)\Gamma(1.800)}{2.4\Gamma(1.100)}; \]
et à l'aide de la table, on trouve immédiatement
\[ \log Z=9.5635972,\quad\text{et}\quad Z=0.3660979. \]

\subsection*{(66).}

Cherchons maintenant la valeur de $\log\Gamma(1+k)$, $k$ étant supposé très-petit. Cette valeur pourrait se déduire de l'interpolation des trois premiers termes de la table; mais il sera plus exact d'interpoler des termes moins rapprochés, afin que les différences secondes deviennent plus sensibles. Considérons pour cet effet les trois termes qui répondent aux valeurs $a=1.000$, $a=1.005$, $a=1.010$; on en déduira les différences premières et secondes, comme il suit:

\origpage{288}
\[ \begin{aligned} a && \log\Gamma && \text{Diff. }1^{\text{re}}. && \text{Diff. }2^{\text{e}}.\\ 1.000 && 0.0000000 && -12445 && +178\\ 1.005 && 9.9987555 && -12267 &&\\ 1.010 && 9.9975288 && && \end{aligned} \]
Soit $\Delta'=-0.0012445$, $\Delta''=0.0000178$; on aura, en négligeant les différences troisièmes, $\log\Gamma\left(1+\frac{5x}{1000}\right)=x\Delta'+\frac{x(x-1)}{2}\Delta''$, et en faisant $x=200k$,
\[ \log\Gamma(1+k)=-k(0.25068)+k^{2}(0.356). \]
Ce logarithme est un logarithme vulgaire; pour le rendre hyperbolique, il faut le multiplier par 2.3025, etc., ce qui donnera
\[ \log\Gamma(1+k)=-k(0.57721)+k^{2}(0.820). \]
Cette valeur n'est qu'approchée; supposons que la valeur exacte soit, en rejetant toujours les termes affectés de $k^{3}$ et des puissances supérieures de $k$,
\[ \log\Gamma(1+k)=-Pk+Qk^{2}, \]
on en déduira
\[ \Gamma(1+k)=1-Pk+(Q+\tfrac{1}{2}P^{2})k^{2}: \]
de là résulte, en changeant le signe de $k$,
\[ \Gamma(1-k)=1+Pk+(Q+\tfrac{1}{2}P^{2})k^{2}. \]
Mais l'équation $\Gamma(1+k)=k\Gamma(k)$ donne $\Gamma(k)=\frac{1}{k}\Gamma(1+k)$; donc
\[ \Gamma(k)=\frac{1}{k}-P+(Q+\tfrac{1}{2}P^{2})k. \]
Multipliant entre elles ces deux dernières équations, on aura
\[ \Gamma(k)\Gamma(1-k)=\frac{1}{k}(1+2Qk^{2}): \]
mais d'un autre côté on sait que le premier membre $=\frac{\pi}{\sin\pi k}=\frac{1}{k}\left(1+\frac{\pi^{2}k^{2}}{6}\right)$; on a donc exactement
\[ Q=\frac{\pi^{2}}{12}=0.822467. \]

\origpage{289}
Nous avions trouvé $Q=0.820$, mais cette valeur ne pouvait être qu'approchée.

Quant à la valeur de $P$, nous l'avons trouvée 0.57721; mais en poussant plus loin l'approximation, on trouverait
\[ P=0.5772156649, \]
ainsi qu'on le fera voir ci-après.

\subsection*{(67).}

Nous connaissons donc maintenant d'une manière fort approchée les valeurs de $\Gamma(1+k)$ et $\Gamma(k)$, lorsque $k$ est très-petit: ces valeurs sont
\[ \begin{aligned} \Gamma(1+k)&=1-Pk+\left(\tfrac{\pi^{2}}{12}+\tfrac{1}{2}P^{2}\right)k^{2},\\ \Gamma(k)&=\frac{1}{k}-P+\left(\tfrac{\pi^{2}}{12}+\tfrac{1}{2}P^{2}\right)k. \end{aligned} \]
Au moyen de la dernière formule, il est aisé de trouver la valeur de la transcendante $\left(\frac{p}{q}\right)$ lorsque $p$ et $q$ sont très-petits par rapport à $n$. En effet, si on substitue les valeurs de $\Gamma\left(\frac{p}{n}\right)$, $\Gamma\left(\frac{q}{n}\right)$, $\Gamma\left(\frac{p+q}{n}\right)$, données par la formule précédente, dans l'équation $(\varepsilon)$, on trouvera
\[ \left(\frac{p}{q}\right)=\frac{p+q}{pq}\left(1-\frac{\pi^{2}}{6}\cdot\frac{pq}{n^{2}}\right), \]
ce qui s'accorde avec l'équation $(l')$.

\subsection*{(68).}

Il reste à faire voir comment nous avons construit la table au moyen de laquelle on trouve si facilement, dans tous les cas, la valeur des fonctions $\Gamma$ et $\left(\frac{p}{q}\right)$. La méthode la plus simple qu'on puisse proposer pour cet objet est celle qui résulte d'une formule donnée par Euler dans son \emph{Calc.\ diff.}, page 465, et que nous allons rapporter.

Si on appelle $S$ la somme de la suite
\[ \log 1+\log 2+\log 3\ldots+\log k, \]
on aura
\[ S=k\log k+\tfrac{1}{2}\log(2\pi k)-k+\frac{A'}{1\cdot2k}-\frac{B'}{3\cdot4k^{3}}+\frac{C'}{5\cdot6k^{5}}-\text{etc.} \]
$A'$, $B'$, $C'$, etc. étant les nombres Bernoulliens. Soit donc $e$ le

\origpage{290}
nombre dont le logarithme est 1, et $R$ un nombre tel qu'on ait
\[ \log R=\frac{A'}{1\cdot2k}-\frac{B'}{3\cdot4k^{3}}+\frac{C'}{5\cdot6k^{5}}-\text{etc.} \]
on aura le produit
\[ 1\cdot2\cdot3\ldots k=\left(\frac{k}{e}\right)^{k}\cdot(2\pi k)^{\frac{1}{2}}\cdot R.\qquad(\rho) \]
Le premier membre est la valeur de $\Gamma(k+1)$, lorsque $k$ est un nombre entier; et comme le second membre est une fonction continue de $k$, on a généralement, quel que soit $k$,
\[ \Gamma(k+1)=\left(\frac{k}{e}\right)^{k}(2\pi k)^{\frac{1}{2}}R.\qquad(\sigma) \]
Telle est la formule par laquelle on pourra dans tous les cas déterminer la valeur approchée de $\Gamma(k+1)$; mais il est à propos de faire à ce sujet quelques observations.

\subsection*{(69).}

La quantité $R$ peut se développer suivant les puissances de $\frac{1}{k}$; car on a $R=1+\log R+\frac{1}{2}\log^{2}R+\frac{1}{2\cdot3}\log^{3}R+\text{etc.}$ Substituant donc la valeur de $\log R$, et mettant au lieu des coefficiens $A'$, $B'$, $C'$, etc. leurs valeurs connues $A'=\frac{1}{6}$, $B'=\frac{1}{30}$, $C'=\frac{1}{42}$, $D'=\frac{1}{30}$, etc., on aura
\[ R=1+\frac{1}{12k}+\frac{1}{2(12k)^{2}}-\frac{139}{30(12k)^{3}}-\frac{571}{120(12k)^{4}}-\text{etc.} \]
Dans cette suite, si on appelle $M$ la partie
\[ 1+\frac{1}{2(12k)^{2}}-\frac{571}{120(12k)^{4}}-\text{etc.} \]
où $k$ est élevé à des puissances paires, et $N$ l'autre partie, on aura $M^{2}-N^{2}=1$; de sorte qu'on peut prendre indifféremment $R=M+N$, ou $R=\frac{1}{M-N}$. En effet, comme toutes les puissances de $k$ sont impaires dans $\log R$, le changement du signe de $k$ donnera $\log(M+N)=-\log.(M-N)$, ou $\log(M^{2}-N^{2})=0$. Donc $M^{2}-N^{2}=1$.

\subsection*{(70).}

Il est à remarquer que la suite
\[ \frac{A'}{1\cdot2k}-\frac{B'}{3\cdot4k^{3}}+\frac{C'}{5\cdot6k^{3}}-\text{etc.} \]
\ednote{Printed $5\cdot6k^{3}$ in the last denominator instead of $5\cdot6k^{5}$; an apparent misprint, cf. the same series at the head of this page.}

\origpage{291}
même en supposant $k$ assez grand, n'est convergente que dans un certain nombre des premiers termes; car on sait que les nombres Bernoulliens, dont les expressions sont fort irrégulières, croissent continuellement, de manière que si $T$ et $V$ sont deux termes consécutifs fort éloignés, l'un du rang $n$, l'autre du rang $n+1$, le rapport $\frac{V}{T}$ a pour limite $\frac{n^{2}}{\pi^{2}}$. Cette suite, qui commence par être convergente pendant un assez grand nombre de termes, surtout si $k$ est un peu grand, finit donc par être divergente, et donnerait une valeur de $\log R$ d'autant plus fautive, qu'on prendrait plus de termes au-delà de ceux où elle cesse d'être convergente.

De là on voit que pour une valeur donnée de $k$, il y a un terme qu'on ne doit pas passer dans le calcul de la suite
\[ \frac{A'}{1\cdot2k}-\frac{B'}{3\cdot4k^{3}}+\text{etc.} \]
Le terme auquel il faudra s'arrêter est celui qui serait suivi d'un terme plus grand, alors l'approximation ne peut aller plus loin; mais elle sera tout aussi étendue qu'on voudra, en prenant $k$ suffisamment grand.

Il en serait de même de la série
\[ R=1+\frac{1}{12k}+\frac{1}{2(12k)^{2}}-\text{etc.}; \]
mais celle-ci n'est pas d'un usage aussi facile que la série
\[ \frac{A'}{1\cdot2k}-\frac{B'}{3\cdot4k^{3}}+\text{etc.} \]
dont la loi est manifeste, et ne dépend que des nombres Bernoulliens.

On peut fixer \emph{a priori} le nombre de termes après lequel la suite
\[ \frac{A'}{1\cdot2k}-\frac{B'}{3\cdot4k^{3}}+\text{etc.} \]
cesse d'être convergente; car en considérant les deux termes consécutifs
\[ \pm\frac{T^{(n)}}{2n\cdot2n-1\cdot k^{2n-1}}\mp\frac{T^{(n+1)}}{2n+2\cdot2n+1\cdot k^{2n+1}}, \]
et les supposant égaux, on aura
\[ \frac{T^{(n+1)}}{T^{(n)}}=\frac{2n+2\cdot2n+1}{2n\cdot2n-1}k^{2}; \]

\origpage{292}
mais plus $n$ est grand, plus le premier membre approche de la limite $\frac{2n+2\cdot2n+1}{4\pi^{2}}$ (\emph{Eul.\ Calc.\ diff.}, page 429). Donc on aura à très-peu près $n=\pi k$. Ainsi en faisant $k=5$, on a $n=15$ ou 16, c'est-à-dire que la série cesse d'être convergente vers le 15\textsuperscript{ème} terme; si on faisait $k=10$, la série ne cesserait d'être convergente que vers le 31\textsuperscript{ème} terme, et ainsi de suite.

\subsection*{(71).}

On peut en même temps avoir la mesure du degré d'approximation que l'on peut obtenir avec une valeur donnée de $k$. En effet, si on appelle $\Omega$ le $n^{\text{ème}}$ terme de la suite
\[ \frac{A'}{1\cdot2k^{2}}-\frac{B'}{3\cdot4k^{3}}+\text{etc.}, \]
\ednote{Printed $1\cdot2k^{2}$ in the first denominator instead of $1\cdot2k$; an apparent misprint.}
on aura
\[ \Omega=\frac{T^{(n)}}{2n(2n-1)k^{2n-1}}, \]
et comme on a à très-peu près
\[ T^{(n)}=\frac{1\cdot2\cdot3\ldots2n}{2^{2n-1}\pi^{2n}}, \]
on pourra faire
\[ \Omega=\frac{1\cdot2\cdot3\ldots2n-2}{\pi(2\pi k)^{2n-1}}. \]
Cette valeur, au moyen de la formule $(\rho)$, devient
\[ \Omega=\frac{\left(\frac{2n-2}{e}\right)^{2n-\frac{3}{2}}(2\pi e)^{\frac{1}{2}}}{\pi(2\pi k)^{2n-1}}, \]
et en mettant $n$ au lieu de $\pi k$, on en déduit aisément
\[ \log\Omega=-2n-\tfrac{1}{2}\log(\pi n). \]
Ainsi, faisant $k=5$ et $n=16$, on aura $\log\Omega=-33.96$. A ce logarithme hyperbolique répond le logarithme vulgaire $-14.75$, de sorte qu'on a $\Omega=10^{-14.75}$. Donc au moyen des 16 premiers termes de la suite $\frac{A'}{1\cdot2k}-\frac{B'}{3\cdot4k^{3}}+\text{etc.}$, on aura la valeur de $\log R$ approchée jusqu'à 15 décimales environ. Si on faisait $k=10$, on pourrait avoir 29 ou 30 décimales exactes, et ainsi de suite.

\subsection*{(72).}

Cette théorie est facile à vérifier, puisque toutes les fois que $k$ est un entier, la valeur de $\Gamma(k+1)$ est exactement $1\cdot2\cdot3\ldots k$.

\origpage{293}

Soit, par exemple, $k=3$, il résulte des formules précédentes que la série égale à $\log R$ cessera d'être convergente après un nombre de termes $n=k\pi=9$ ou 10, et que le nombre de décimales exactes obtenues par ces neuf ou dix termes, sera de 8 ou 9.

En effet, la vraie valeur de $\log R$ se déduit de l'équation $6=\left(\frac{3}{e}\right)^{3}(6\pi)^{\frac{1}{2}}R$, laquelle donne
\[ \log R=0.02767\ 79256\ 86. \]
Cette même valeur déduite de la suite
\[ \log R=\frac{A'}{1\cdot2k}-\frac{B'}{3\cdot4k^{3}}+\frac{C'}{5\cdot6k^{5}}-\text{etc.}, \]
se trouve en calculant successivement les différens termes comme il suit:
\[ \begin{aligned} 1^{\text{er}}\ \text{terme} && + && 0.02777\ 77777\ 78\\ 2^{\text{e}} && - && 0.00010\ 28806\ 58\\ && && 0.02767\ 48971\ 20\\ 3^{\text{e}} && + && 32660\ 53\\ && && 0.02767\ 81631\ 73\\ 4^{\text{e}} && - && 2721\ 71\\ && && 0.02767\ 78910\ 02\\ 5^{\text{e}} && + && 427\ 65\\ && && 0.02767\ 79337\ 67\\ 6^{\text{e}} && - && 108\ 24\\ && && 0.02767\ 79229\ 43\\ 7^{\text{e}} && + && 40\ 21\\ && && 0.02767\ 79269\ 64\\ 8^{\text{e}} && - && 20\ 59\\ && && 0.02767\ 79249\ 05\\ 9^{\text{e}} && + && 13\ 91\\ && && 0.02767\ 79262\ 96\\ 10^{\text{e}} && - && 11\ 98\\ && && 0.02767\ 79250\ 98\\ 11^{\text{e}} && + && 12\ 81\\ && && 0.02767\ 79263\ 79 \end{aligned} \]

\origpage{294}

On voit que conformément à la formule, la série cesse d'être convergente passé le 10\textsuperscript{ème} terme, et que la valeur de $\log R$ qui en est déduite, doit être comprise entre $0.02767\ 7926296$ et $0.02767\ 7925098$, ce qui donne par un milieu
\[ \log R=0.02767\ 7925697, \]
valeur exacte presque jusqu'à la onzième décimale. Mais en continuant la suite plus loin, on s'éloignerait de plus en plus du vrai résultat.

Cet exemple met dans tout son jour la manière de tirer tout le parti possible pour les approximations, des suites \emph{demi-convergentes}, c'est-à-dire des suites qui sont convergentes dans les premiers termes, et qui deviennent ensuite divergentes.

\subsection*{(73).}

Au moyen de la formule $(\sigma)$, on peut développer en série la fonction $\Gamma(k)$ lorsque $k$ est très-petit. Pour cela, observons d'abord que $\Gamma(k+1)=k\Gamma(k)$, et qu'ainsi on aura
\[ \Gamma(k)=c^{-k}k^{k-\frac{1}{2}}(2\pi)^{\frac{1}{2}}R; \]
\ednote{The base of the first factor reads $c$ instead of $e$ cf. formula $(\sigma)$; a misprint, or an $e$ whose bar failed to print.}
d'où
\[ \log\Gamma(k)=(k-\tfrac{1}{2})\log k-k+\tfrac{1}{2}l(2\pi)+\frac{A'}{1\cdot2k}-\frac{B'}{3\cdot4k^{3}}+\text{etc.}\ldots(\tau) \]
Cette formule ne peut servir que pour des valeurs de $k$ plus grandes que l'unité; mais si l'on met $1+k$ au lieu de $k$, et qu'au lieu du premier membre qui deviendra $\log\Gamma(1+k)$, on mette sa valeur $\log k+\log\Gamma(k)$, on en tirera de nouveau
\[ \begin{aligned} \log\Gamma(k)={}&-\log k+(\tfrac{1}{2}+k)\log(1+k)-1-k+\tfrac{1}{2}\log(2\pi)\\ &+\frac{A'}{1\cdot2(1+k)}-\frac{B'}{3\cdot4(1+k)^{3}}+\frac{C'}{5\cdot6(1+k)^{5}}-\text{etc.} \end{aligned} \]
et le développement étant fait jusqu'aux quantités de l'ordre $k^{2}$ exclusivement, on aura
\[ \begin{aligned} \log\Gamma(k)={}&-\log k+\tfrac{1}{2}\log(2\pi)-1+\frac{A'}{1\cdot2}-\frac{B'}{3\cdot4}+\frac{C'}{5\cdot6}-\text{etc.}\\ &-k\left(\tfrac{1}{2}+\frac{A'}{2}-\frac{B'}{4}+\frac{C'}{6}-\text{etc.}\right) \end{aligned} \]
Lorsqu'on fait $k=0$, on doit avoir $\log\Gamma(k)=-\log k$, parce que $k\Gamma(k)=\Gamma(1+k)$, et qu'en faisant $k=0$ le second membre

\origpage{295}
$=\Gamma(1)=1$; donc on doit avoir
\[ \tfrac{1}{2}\log(2\pi)-1+\frac{A'}{1\cdot2}-\frac{B'}{3\cdot4}+\frac{C'}{5\cdot6}-\text{etc.}=0. \]
Et en effet, Euler a trouvé cette égalité, page 466 de son \emph{Calc.\ diff.} Cela posé, la valeur précédente se réduit à celle-ci,
\[ \log\Gamma(k)=-\log k-k\left(\tfrac{1}{2}+\frac{A'}{2}-\frac{B'}{4}+\frac{C'}{6}-\text{etc.}\right); \]
or dans l'ouvrage cité, page 444, on trouve encore l'égalité
\[ C=\tfrac{1}{2}+\frac{A'}{2}-\frac{B'}{4}+\frac{C'}{6}-\text{etc.}, \]
$C$ étant une constante dont la valeur calculée avec précision par une autre voie, est
\[ C=0.5772156649015325; \]
donc enfin on aura, $k$ étant très-petit,
\[ \log\Gamma(k)=-\log k-Ck. \]
Nous avons trouvé ci-dessus, en poussant l'approximation plus loin,
\[ \Gamma(k)=\frac{1}{k}\left[1-Pk+\left(\tfrac{1}{2}P^{2}+\frac{\pi^{2}}{12}\right)k^{2}\right]; \]
de là on voit que $P=C$, et qu'ainsi $P$ n'est autre chose que la constante $C$ dont on vient de donner la valeur approchée jusqu'à seize décimales.

Nous connaissons donc maintenant la valeur de $\Gamma(k)$ très-approchée, lorsque $k$ est très-petit, et on pourrait en approcher encore davantage en poussant le développement plus loin.

\subsection*{(74).}

Mais voici des considérations qui mènent plus généralement et plus directement au même but. Soit
\[ \begin{aligned} M&=1+\tfrac{1}{2}+\tfrac{1}{3}\ldots+\frac{1}{x},\\ N&=\frac{1}{x+1}+\frac{1}{x+2}+\frac{1}{x+3}+\text{etc.}, \end{aligned} \]
cette dernière suite étant prolongée à l'infini.

Nous regarderons la quantité $M$ comme une fonction continue

\origpage{296}
de $x$, puisqu'en effet on a (\emph{Calc.\ diff.}, pag.\ 443),
\[ M=lx+C+\frac{1}{2x}-\frac{A'}{2x^{2}}+\frac{B'}{4x^{4}}-\frac{C'}{6x^{6}}+\text{etc.}; \]
formule qui est propre à donner la valeur approchée de $M$, quel que soit $x$, pourvu qu'on suppose $x$ plus grand que l'unité. Cette même formule donnerait la valeur de $M$, lorsque $x$ est plus petit que l'unité; car en représentant par $M(x)$ et par $M(x+1)$ des fonctions semblables de $x$ et de $x+1$, on a évidemment
\[ M(x)=M(x+1)-\frac{1}{x+1}, \]
ou encore
\[ \begin{aligned} M(x)&=M(x+2)-\frac{1}{x+1}-\frac{1}{x+2},\\ &\text{etc.}; \end{aligned} \]
de sorte que $x$ étant plus petit que l'unité, on aura la fonction $M(x)$ ou $M$, au moyen d'une semblable fonction où $x$ sera augmenté de plusieurs unités, et qui rendra la suite précédente assez convergente dans les premiers termes, pour qu'elle puisse donner toute l'approximation qu'on peut desirer.

\subsection*{(75).}

Cela posé, si $x$ devient $x+\omega$ ($\omega$ étant plus petit que l'unité), les fonctions $M$ et $N$ qui peuvent être représentées par $M(x)$ et $N(x)$, deviendront $M(x+\omega)$ et $N(x+\omega)$. Or, je dis que la somme $M(x)+N(x)=M(x+\omega)+N(x+\omega)$, et que les deux sommes sont égales à une même constante.

En effet, si la suite qui a pour somme $N(x)$, au lieu d'être prolongée à l'infini, était continuée seulement jusqu'au terme $\frac{1}{x+m}$, $m$ étant un nombre très-grand, la somme $M(x)+N(x)$ et la somme $M(x+\omega)+N(x+\omega)$ ne pourraient différer entre elles que d'une quantité moindre que $\frac{1}{x+m+1}$, puisque cette différence est celle qui a lieu lorsque $\omega=1$. Or $m$ étant très-grand, la différence $\frac{1}{x+m+1}$ est censée nulle; à plus forte raison le sera-t-elle lorsque la suite $N(x)$ sera prolongée à l'infini. On aura donc
\[ M+N=\text{const}=C'; \]

\origpage{297}
mais puisque la valeur de $N$ est
\[ \frac{1}{1+x}+\frac{1}{2+x}+\frac{1}{3+x}+\text{etc.}, \]
si on développe chacune de ces fractions dans l'hypothèse que $x$ est plus petit que l'unité, on aura
\[ \begin{aligned} N={}&1+\tfrac{1}{2}+\tfrac{1}{3}+\tfrac{1}{4}+\text{etc.}\\ &-x\left(1+\frac{1}{2^{2}}+\frac{1}{3^{2}}+\frac{1}{4^{2}}+\text{etc.}\right)\\ &+x^{2}\left(1+\frac{1}{2^{3}}+\frac{1}{3^{3}}+\frac{1}{4^{3}}+\text{etc.}\right)\\ &-\text{etc.} \end{aligned} \]
La première partie $1+\frac{1}{2}+\frac{1}{3}+\frac{1}{4}+\text{etc.}$ n'est autre chose que la constante $C'$; donc on aura
\[ M=S_{2}x-S_{3}x^{2}+S_{4}x^{3}-S_{5}x^{4}+\text{etc.},\qquad(\upsilon) \]
$S_{n}$ représentant en général la somme des puissances réciproques de degré $n$ des nombres naturels.

\subsection*{(76).}

Mais dans l'hypothèse de $x>1$, on a, d'après l'équation $(\tau)$,
\[ \frac{d\log\Gamma(x)}{dx}=\log x-\frac{1}{2x}-\frac{A'}{2x^{2}}+\frac{B'}{4x^{4}}-\text{etc.}, \]
et dans la même hypothèse on a
\[ M=\log x+C+\frac{1}{2x}-\frac{A'}{2x^{2}}+\frac{B'}{4x^{4}}-\text{etc.} \]
Donc
\[ \frac{d\log\Gamma(x)}{dx}=-\frac{1}{x}+M-C,\qquad(\varphi) \]
équation qui doit avoir lieu quel que soit $x$, puisque $M$ peut être regardée comme une fonction continue de $x$.

Si maintenant on suppose $x<1$, ce qui permettra d'employer la suite $(\upsilon)$, on aura
\[ \frac{d\log\Gamma(x)}{dx}=-\frac{1}{x}-C+S_{2}x-S_{3}x^{2}+S_{4}x^{3}-\text{etc.}; \]
d'où résulte en intégrant,
\[ \log\Gamma(x)=-\log x-Cx+\tfrac{1}{2}S_{2}x^{2}-\tfrac{1}{3}S_{3}x^{3}+\tfrac{1}{4}S_{4}x^{4}-\text{etc.}\qquad(\psi) \]

\origpage{298}
On n'ajoute pas de constante, parce que $x\Gamma(x)$ ou $\Gamma(1+x)$ se réduit à l'unité lorsque $x=0$.

\subsection*{(77).}

La formule $(\psi)$ qui donne un développement complet de $\log\Gamma(x)$, servira à exprimer la valeur de toute fonction proposée $\Gamma(z)$, puisque cette fonction peut toujours être ramenée au cas où $z$ est plus petit que $\frac{1}{2}$, et même à celui où $z$ est plus petit que $\frac{1}{4}$.

De la formule $(\psi)$ on déduit immédiatement la suivante,
\[ \log\Gamma(1+x)=-Cx+\tfrac{1}{2}S_{2}x^{2}-\tfrac{1}{3}S_{3}x^{3}+\tfrac{1}{4}S_{4}x^{4}-\text{etc.},\qquad(\omega) \]
et en changeant dans celle-ci le signe de $x$, on aurait
\[ \log\Gamma(1-x)=Cx+\tfrac{1}{2}S_{2}x^{2}+\tfrac{1}{3}S_{3}x^{3}+\tfrac{1}{4}S_{4}x^{4}+\text{etc.},\qquad(\alpha') \]
formules également propres à faire trouver dans chaque cas proposé la valeur de $\Gamma(z)$, puisqu'on peut toujours ramener cette fonction soit à la forme $\Gamma(1+x)$, soit à la forme $\Gamma(1-x)$, $x$ n'excédant pas $\frac{1}{4}$.

Ainsi on a des séries régulières et toujours convergentes pour calculer la valeur d'une fonction donnée $\Gamma(z)$; elles supposent seulement l'emploi des quantités $S_{2}$, $S_{3}$, $S_{4}$, etc., dont Euler a donné les valeurs numériques fort approchées jusqu'à $S_{15}$ (\emph{Calc.\ diff.}, page 456).

Nous devons observer que la formule $(\omega)$ revient à celle qu'on voit dans l'ouvrage cité, page 800; il paraît seulement qu'Euler n'a pas aperçu l'usage que cette suite pouvait avoir dans la détermination des fonctions $\Gamma(x)$ dont il s'est occupé dans d'autres endroits de ses ouvrages.

Remarquons encore qu'on peut déduire de la formule $(\omega)$, cette valeur de $C$,
\[ C=-\frac{1}{x}\log\Gamma(1+x)+\frac{1}{2}S_{2}x-\frac{1}{3}S_{3}x^{2}+\frac{1}{4}S_{4}x^{3}-\text{etc.}, \]
d'où l'on voit qu'il suffit de connaître une valeur particulière de $\Gamma(1+x)$, et qu'on aura la valeur de $C$ exprimée par une suite d'autant plus convergente que $x$ sera plus petit. Si l'on fait $x=\frac{1}{2}$, on aura
\[ \Gamma\left(1+\tfrac{1}{2}\right)=\tfrac{1}{2}\Gamma\tfrac{1}{2}=\tfrac{1}{2}\sqrt{\pi}, \]

\origpage{299}
d'où résulte
\[ C=\log\frac{4}{\pi}+\tfrac{1}{2}S_{2}\cdot\tfrac{1}{2}-\tfrac{1}{3}S_{3}\cdot\tfrac{1}{4}+\tfrac{1}{4}S_{4}\cdot\tfrac{1}{8}-\text{etc.} \]

\subsection*{(78).}

Les deux équations $\Gamma(1+x)=x\Gamma(x)$, et $\Gamma(x)\Gamma(1-x)=\frac{\pi}{\sin\pi x}$, donnent
\[ \Gamma(1+x)\Gamma(1-x)=\frac{\pi x}{\sin\pi x};\qquad(\beta') \]
prenant les logarithmes des deux membres, et substituant les valeurs données par les formules $(\omega)$ et $(\alpha')$, on aura
\[ \log\left(\frac{\pi x}{\sin\pi x}\right)=S_{2}x^{2}+\tfrac{1}{2}S_{4}x^{4}+\tfrac{1}{3}S_{6}x^{6}+\text{etc.}, \]
formule connue, et qui par sa différentielle, sert à déterminer les valeurs des quantités $S_{2}$, $S_{4}$, $S_{6}$, etc.

On peut faire usage de cette formule pour rendre encore plus convergentes les suites contenues dans les équations $(\omega)$ et $(\alpha')$: on aura ainsi
\[ \left.\begin{aligned} \log\Gamma(1+x)&=\tfrac{1}{2}\log\left(\frac{\pi x}{\sin\pi x}\right)-Cx-\tfrac{1}{3}S_{3}x^{3}-\tfrac{1}{5}S_{5}x^{5}-\text{etc.}\\ \log\Gamma(1-x)&=\tfrac{1}{2}\log\left(\frac{\pi x}{\sin\pi x}\right)+Cx+\tfrac{1}{3}S_{3}x^{3}+\tfrac{1}{5}S_{5}x^{5}+\text{etc.} \end{aligned}\right\}\qquad(\gamma') \]
La dernière, en faisant $x=\frac{1}{2}$, donne
\[ C=L2-\tfrac{1}{3}S_{3}\cdot\tfrac{1}{4}-\tfrac{1}{5}S_{5}\cdot\tfrac{1}{16}-\text{etc.} \]
valeur plus convergente que celle de l'art.\ 77.

\subsection*{(79).}

Ayant l'expression développée de $\log\Gamma(x)$, par la formule $(\psi)$, on peut pareillement avoir celle du logarithme de la fonction $\left(\frac{p}{q}\right)$; car puisqu'on a trouvé
\[ \left(\frac{p}{q}\right)=\frac{\Gamma\left(\frac{p}{n}\right)\Gamma\left(\frac{p}{n}\right)}{n\Gamma\left(\frac{p+q}{n}\right)}, \]
\ednote{The second factor of the numerator is printed $\Gamma\left(\frac{p}{n}\right)$; read $\Gamma\left(\frac{q}{n}\right)$.}
si on prend les logarithmes de chaque membre, et qu'on désigne par $Z$ la fonction $\left(\frac{p}{q}\right)$ qui répond à une valeur donnée de $n$, on

\origpage{300}
aura
\[ \left.\begin{aligned} \log Z=\log.\left(\frac{p+q}{pq}\right)&-\tfrac{1}{2}S_{2}\cdot\frac{(p+q)^{2}-p^{2}-q^{2}}{n^{2}}\\ &+\tfrac{1}{3}S_{3}\cdot\frac{(p+q)^{3}-p^{3}-q^{3}}{n^{3}}\\ &-\tfrac{1}{4}S_{4}\cdot\frac{(p+q)^{4}-p^{4}-q^{4}}{n^{4}}\\ &+\text{etc.} \end{aligned}\right\}\qquad(\delta') \]
Les deux premiers termes de cette formule s'accordent avec l'équation $(l')$; mais on voit ici la loi générale du développement qu'on peut continuer à volonté, et qui donne une suite d'autant plus convergente, que $p$ et $q$ seront plus petits par rapport à $n$.

Ainsi pour la fonction $\left(\frac{a}{a}\right)$ désignée ci-dessus par $M_{a}$, on aura
\[ \log.M_{a}=\log\left(\frac{2}{a}\right)-\tfrac{1}{2}S_{2}(2^{2}-2)\frac{a^{2}}{n^{2}}+\tfrac{1}{3}S_{3}(2^{3}-2)\frac{a^{3}}{n^{3}}-\text{etc.} \]

\subsection*{(80).}

D'après ces formules, il a été facile de calculer la table ci-jointe des valeurs de $\log\Gamma(a)$. Pour cela, nous avons fait $k=4+a$ dans la formule $(\tau)$, (on aurait pu prendre également $k=3+a$, $k=5+a$, etc.) Alors le premier membre donnant la valeur de $\log\Gamma(4+a)$, nous en avons déduit $\log\Gamma(a)$ par la relation connue entre ces quantités; savoir:
\[ \Gamma(4+a)=(3+a)(2+a)(1+a)a\Gamma(a). \]
Nous avons donc eu à calculer la formule
\[ \begin{aligned} \log\Gamma(a)={}&\left(k-\tfrac{1}{2}\right)\log k+\tfrac{1}{2}l(2\pi)-mk+\frac{mA'}{1\cdot2k}-\frac{mB'}{3\cdot4k^{3}}+\frac{mC'}{5\cdot6k^{5}}-\text{etc.}\\ &-\log[a(1+a)(2+a)(3+a)], \end{aligned} \]
dans laquelle on a introduit le facteur $m=0.43429$, etc. afin de réduire tout aux logarithmes des tables.

De cette manière on n'a jamais eu besoin de calculer plus de deux ou trois termes de la série $\frac{mA'}{1\cdot2k}-\frac{mB'}{3\cdot4k^{3}}+\frac{mC'}{5\cdot6k^{5}}-\text{etc.}$, pour avoir $\log\Gamma(a)$ approché jusqu'à sept décimales, dans tout l'intervalle depuis $a=1$ jusqu'à $a=2$.

\subsection*{(81).}

Nous remarquerons en finissant que les intégrales de la forme

\origpage{301}
$\displaystyle\int e^{-t^{m}}t^{n}\,dt$, prises depuis $t=0$ jusqu'à $t=\infty$, peuvent être ramenées aux fonctions $\Gamma(a)$. En effet, soit d'abord $t^{n+1}=z$, et $\frac{m}{n+1}=\alpha$, l'intégrale précédente deviendra $\displaystyle\int\frac{m}{n+1}\cdot e^{-z^{\alpha}}\,dz$,\ednote{Printed $\frac{m}{n+1}$; the substitution gives the factor $\frac{1}{n+1}$; an apparent misprint.} celle-ci devant encore être prise depuis $z=0$ jusqu'à $z=\infty$.

De là on voit que l'intégrale proposée ne perd pas de sa généralité en faisant $n=0$, et qu'ainsi on peut se proposer simplement l'intégrale $\displaystyle\int e^{-t^{m}}\,dt$. Si dans celle-ci on fait $e^{-t^{m}}=x$, ou $t=\left(\log\frac{1}{x}\right)^{\frac{1}{m}}$, la transformée sera $\displaystyle\frac{1}{m}\int dx\left(\log\frac{1}{x}\right)^{\frac{1}{m}-1}$, et cette nouvelle intégrale devra être prise depuis $x=0$ jusqu'à $x=1$; on aura donc généralement
\[ \int e^{-t^{m}}\,dt=\frac{1}{m}\Gamma\left(\frac{1}{m}\right). \]

Ainsi les intégrales dont il s'agit n'offrent point une nouvelle espèce de transcendantes, et se rapportent aux fonctions $\Gamma$.

Il n'est pas inutile pour l'histoire de la science, d'observer que l'intégrale $\displaystyle\int dx\left(\log\frac{1}{x}\right)^{-\frac{1}{2}}=\sqrt{\pi}$ que l'on trouve § 28 du mémoire d'Euler, imprimé dans le tom.\ XVI des \emph{Novi Comm.\ Petrop.}, avait été donnée long-temps auparavant par le même auteur, dans le tom.\ V des anciens Mémoires de Pétersbourg, pag.\ 44. C'est donc à cette époque que remonte la découverte de l'intégrale $\displaystyle\int e^{-x^{2}}\,dx=\frac{1}{2}\sqrt{\pi}$, prise entre les limites $x=0$, $x=\infty$, puisque la simple substitution $e^{-x^{2}}=z$ suffit pour ramener l'intégrale $\displaystyle\int e^{-x^{2}}\,dx$ à la forme $\displaystyle\frac{1}{2}\int dz\left(\log\frac{1}{z}\right)^{-\frac{1}{2}}$.

FIN DE LA SECONDE PARTIE.

\origpage{302}

\section*{TABLE DES LOGARITHMES DE LA FONCTION $\Gamma(a)$.}

\[ \begin{aligned} a && \text{Log }\Gamma &\qquad& a && \text{Log }\Gamma. &\qquad& a. && \text{Log }\Gamma. &\qquad& a. && \text{Log }\Gamma.\\ 1.000 && 0.0000000 && 1.050 && 9.9883379 && 1.100 && 9.9783407 && 1.150 && 9.9699006\\ 1.001 && 9.9997497 && 1.051 && 9.9881220 && 1.101 && 9.9781570 && 1.151 && 9.9697470\\ 1.002 && 9.9995001 && 1.052 && 9.9879068 && 1.102 && 9.9779739 && 1.152 && 9.9695941\\ 1.003 && 9.9992512 && 1.053 && 9.9876922 && 1.103 && 9.9777914 && 1.153 && 9.9694417\\ 1.004 && 9.9990030 && 1.054 && 9.9874783 && 1.104 && 9.9776095 && 1.154 && 9.9692898\\ 1.005 && 9.9987555 && 1.055 && 9.9872651 && 1.105 && 9.9774283 && 1.155 && 9.9691386\\ 1.006 && 9.9985088 && 1.056 && 9.9870525 && 1.106 && 9.9772477 && 1.156 && 9.9689879\\ 1.007 && 9.9982627 && 1.057 && 9.9868406 && 1.107 && 9.9770677 && 1.157 && 9.9688378\\ 1.008 && 9.9980174 && 1.058 && 9.9866293 && 1.108 && 9.9768883 && 1.158 && 9.9686883\\ 1.009 && 9.9977727 && 1.059 && 9.9864187 && 1.109 && 9.9767096 && 1.159 && 9.9685394\\ 1.010 && 9.9975288 && 1.060 && 9.9862088 && 1.110 && 9.9765314 && 1.160 && 9.9683910\\ 1.011 && 9.9972855 && 1.061 && 9.9859996 && 1.111 && 9.9763539 && 1.161 && 9.9682432\\ 1.012 && 9.9970430 && 1.062 && 9.9857910 && 1.112 && 9.9761769 && 1.162 && 9.9680960\\ 1.013 && 9.9968011 && 1.063 && 9.9855830 && 1.113 && 9.9760006 && 1.163 && 9.9679494\\ 1.014 && 9.9965599 && 1.064 && 9.9853758 && 1.114 && 9.9758249 && 1.164 && 9.9678033\\ 1.015 && 9.9963195 && 1.065 && 9.9851691 && 1.115 && 9.9756498 && 1.165 && 9.9676578\\ 1.016 && 9.9960798 && 1.066 && 9.9849631 && 1.116 && 9.9754753 && 1.166 && 9.9675129\\ 1.017 && 9.9958408 && 1.067 && 9.9847577 && 1.117 && 9.9753014 && 1.167 && 9.9673686\\ 1.018 && 9.9956025 && 1.068 && 9.9845530 && 1.118 && 9.9751281 && 1.168 && 9.9672248\\ 1.019 && 9.9953648 && 1.069 && 9.9843489 && 1.119 && 9.9749555 && 1.169 && 9.9670816\\ 1.020 && 9.9951279 && 1.070 && 9.9841455 && 1.120 && 9.9747834 && 1.170 && 9.9669390\\ 1.021 && 9.9948917 && 1.071 && 9.9839427 && 1.121 && 9.9746120 && 1.171 && 9.9667969\\ 1.022 && 9.9946561 && 1.072 && 9.9837406 && 1.122 && 9.9744411 && 1.172 && 9.9666554\\ 1.023 && 9.9944213 && 1.073 && 9.9835391 && 1.123 && 9.9742709 && 1.173 && 9.9665145\\ 1.024 && 9.9941871 && 1.074 && 9.9833383 && 1.124 && 9.9741012 && 1.174 && 9.9663742\\ 1.025 && 9.9939536 && 1.075 && 9.9831381 && 1.125 && 9.9739322 && 1.175 && 9.9662344\\ 1.026 && 9.9937208 && 1.076 && 9.9829386 && 1.126 && 9.9737638 && 1.176 && 9.9660952\\ 1.027 && 9.9934887 && 1.077 && 9.9827398 && 1.127 && 9.9735960 && 1.177 && 9.9659566\\ 1.028 && 9.9932572 && 1.078 && 9.9825415 && 1.128 && 9.9734287 && 1.178 && 9.9658186\\ 1.029 && 9.9930265 && 1.079 && 9.9823440 && 1.129 && 9.9732621 && 1.179 && 9.9656810\\ 1.030 && 9.9927964 && 1.080 && 9.9821470 && 1.130 && 9.9730961 && 1.180 && 9.9655441\\ 1.031 && 9.9925670 && 1.081 && 9.9819507 && 1.131 && 9.9729307 && 1.181 && 9.9654077\\ 1.032 && 9.9923383 && 1.082 && 9.9817550 && 1.132 && 9.9727659 && 1.182 && 9.9652718\\ 1.033 && 9.9921103 && 1.083 && 9.9815599 && 1.133 && 9.9726016 && 1.183 && 9.9651366\\ 1.034 && 9.9918830 && 1.084 && 9.9813655 && 1.134 && 9.9724380 && 1.184 && 9.9650019\\ 1.035 && 9.9916564 && 1.085 && 9.9811717 && 1.135 && 9.9722750 && 1.185 && 9.9648677\\ 1.036 && 9.9914305 && 1.086 && 9.9809786 && 1.136 && 9.9721126 && 1.186 && 9.9647341\\ 1.037 && 9.9912052 && 1.087 && 9.9807860 && 1.137 && 9.9719508 && 1.187 && 9.9646011\\ 1.038 && 9.9909807 && 1.088 && 9.9805942 && 1.138 && 9.9717896 && 1.188 && 9.9644687\\ 1.039 && 9.9907567 && 1.089 && 9.9804029 && 1.139 && 9.9716289 && 1.189 && 9.9643367\\ 1.040 && 9.9905335 && 1.090 && 9.9802123 && 1.140 && 9.9714689 && 1.190 && 9.9642054\\ 1.041 && 9.9903109 && 1.091 && 9.9800223 && 1.141 && 9.9713095 && 1.191 && 9.9640746\\ 1.042 && 9.9900890 && 1.092 && 9.9798330 && 1.142 && 9.9711506 && 1.192 && 9.9639444\\ 1.043 && 9.9898678 && 1.093 && 9.9796443 && 1.143 && 9.9709923 && 1.193 && 9.9638147\\ 1.044 && 9.9896472 && 1.094 && 9.9794562 && 1.144 && 9.9708346 && 1.194 && 9.9636856\\ 1.045 && 9.9894273 && 1.095 && 9.9792687 && 1.145 && 9.9706775 && 1.195 && 9.9635570\\ 1.046 && 9.9892081 && 1.096 && 9.9790819 && 1.146 && 9.9705209 && 1.196 && 9.9634290\\ 1.047 && 9.9889895 && 1.097 && 9.9788956 && 1.147 && 9.9703650 && 1.197 && 9.9633016\\ 1.048 && 9.9887717 && 1.098 && 9.9787100 && 1.148 && 9.9702096 && 1.198 && 9.9631747\\ 1.049 && 9.9885545 && 1.099 && 9.9785251 && 1.149 && 9.9700548 && 1.199 && 9.9630483\\ 1.050 && 9.9883379 && 1.100 && 9.9783407 && 1.150 && 9.9699006 && 1.200 && 9.9629225 \end{aligned} \]

\origpage{303}

\[ \begin{aligned} a && \text{Log }\Gamma &\qquad& a && \text{Log }\Gamma. &\qquad& a. && \text{Log }\Gamma. &\qquad& a. && \text{Log }\Gamma.\\ 1.200 && 9.9629225 && 1.250 && 9.9573211 && 1.300 && 9.9530204 && 1.350 && 9.9499515\\ 1.201 && 9.9627972 && 1.251 && 9.9572226 && 1.301 && 9.9529472 && 1.351 && 9.9499022\\ 1.202 && 9.9626725 && 1.252 && 9.9571246 && 1.302 && 9.9528744 && 1.352 && 9.9498535\\ 1.203 && 9.9625484 && 1.253 && 9.9570271 && 1.303 && 9.9528021 && 1.353 && 9.9498051\\ 1.204 && 9.9624248 && 1.254 && 9.9569302 && 1.304 && 9.9527304 && 1.354 && 9.9497573\\ 1.205 && 9.9623017 && 1.255 && 9.9568337 && 1.305 && 9.9526591 && 1.355 && 9.9497099\\ 1.206 && 9.9621792 && 1.256 && 9.9567377 && 1.306 && 9.9525883 && 1.356 && 9.9496630\\ 1.207 && 9.9620573 && 1.257 && 9.9566423 && 1.307 && 9.9525180 && 1.357 && 9.9496165\\ 1.208 && 9.9619359 && 1.258 && 9.9565474 && 1.308 && 9.9524482 && 1.358 && 9.9495706\\ 1.209 && 9.9618150 && 1.259 && 9.9564530 && 1.309 && 9.9523789 && 1.359 && 9.9495251\\ 1.210 && 9.9616947 && 1.260 && 9.9563591 && 1.310 && 9.9523101 && 1.360 && 9.9494800\\ 1.211 && 9.9615749 && 1.261 && 9.9562657 && 1.311 && 9.9522417 && 1.361 && 9.9494354\\ 1.212 && 9.9614556 && 1.262 && 9.9561729 && 1.312 && 9.9521739 && 1.362 && 9.9493913\\ 1.213 && 9.9613369 && 1.263 && 9.9560805 && 1.313 && 9.9521065 && 1.363 && 9.9493476\\ 1.214 && 9.9612187 && 1.264 && 9.9559887 && 1.314 && 9.9520396 && 1.364 && 9.9493044\\ 1.215 && 9.9611011 && 1.265 && 9.9558974 && 1.315 && 9.9519732 && 1.365 && 9.9492617\\ 1.216 && 9.9609840 && 1.266 && 9.9558066 && 1.316 && 9.9519073 && 1.366 && 9.9492194\\ 1.217 && 9.9608675 && 1.267 && 9.9557163 && 1.317 && 9.9518419 && 1.367 && 9.9491776\\ 1.218 && 9.9607516 && 1.268 && 9.9556266 && 1.318 && 9.9517769 && 1.368 && 9.9491363\\ 1.219 && 9.9606361 && 1.269 && 9.9555373 && 1.319 && 9.9517125 && 1.369 && 9.9490953\\ 1.220 && 9.9605212 && 1.270 && 9.9554486 && 1.320 && 9.9516485 && 1.370 && 9.9490549\\ 1.221 && 9.9604068 && 1.271 && 9.9553604 && 1.321 && 9.9515850 && 1.371 && 9.9490149\\ 1.222 && 9.9602930 && 1.272 && 9.9552727 && 1.322 && 9.9515220 && 1.372 && 9.9489754\\ 1.223 && 9.9601797 && 1.273 && 9.9551855 && 1.323 && 9.9514594 && 1.373 && 9.9489363\\ 1.224 && 9.9599669 && 1.274 && 9.9550988 && 1.324 && 9.9513974 && 1.374 && 9.9488977\\ 1.225 && 9.9599547 && 1.275 && 9.9550126 && 1.325 && 9.9513358 && 1.375 && 9.9488595\\ 1.226 && 9.9598430 && 1.276 && 9.9549269 && 1.326 && 9.9512747 && 1.376 && 9.9488218\\ 1.227 && 9.9597318 && 1.277 && 9.9548417 && 1.327 && 9.9512141 && 1.377 && 9.9487845\\ 1.228 && 9.9596212 && 1.278 && 9.9547570 && 1.328 && 9.9511540 && 1.378 && 9.9487477\\ 1.229 && 9.9595111 && 1.279 && 9.9546728 && 1.329 && 9.9510944 && 1.379 && 9.9487114\\ 1.230 && 9.9594015 && 1.280 && 9.9545891 && 1.330 && 9.9510352 && 1.380 && 9.9486755\\ 1.231 && 9.9592925 && 1.281 && 9.9545059 && 1.331 && 9.9509765 && 1.381 && 9.9486401\\ 1.232 && 9.9591839 && 1.282 && 9.9544232 && 1.332 && 9.9509183 && 1.382 && 9.9486051\\ 1.233 && 9.9590760 && 1.283 && 9.9543410 && 1.333 && 9.9508605 && 1.383 && 9.9485706\\ 1.234 && 9.9589685 && 1.284 && 9.9542594 && 1.334 && 9.9508033 && 1.384 && 9.9485365\\ 1.235 && 9.9588616 && 1.285 && 9.9541782 && 1.335 && 9.9507465 && 1.385 && 9.9485029\\ 1.236 && 9.9587552 && 1.286 && 9.9540976 && 1.336 && 9.9506902 && 1.386 && 9.9484698\\ 1.237 && 9.9586494 && 1.287 && 9.9540174 && 1.337 && 9.9506344 && 1.387 && 9.9484371\\ 1.238 && 9.9585441 && 1.288 && 9.9539377 && 1.338 && 9.9505791 && 1.388 && 9.9484048\\ 1.239 && 9.9584393 && 1.289 && 9.9538586 && 1.339 && 9.9505242 && 1.389 && 9.9483730\\ 1.240 && 9.9583350 && 1.290 && 9.9537799 && 1.340 && 9.9504698 && 1.390 && 9.9483417\\ 1.241 && 9.9582313 && 1.291 && 9.9537017 && 1.341 && 9.9504159 && 1.391 && 9.9483108\\ 1.242 && 9.9581280 && 1.292 && 9.9536240 && 1.342 && 9.9503624 && 1.392 && 9.9482804\\ 1.243 && 9.9580253 && 1.293 && 9.9535468 && 1.343 && 9.9503094 && 1.393 && 9.9482504\\ 1.244 && 9.9579232 && 1.294 && 9.9534701 && 1.344 && 9.9502569 && 1.394 && 9.9482208\\ 1.245 && 9.9578215 && 1.295 && 9.9533939 && 1.345 && 9.9502048 && 1.395 && 9.9481917\\ 1.246 && 9.9577204 && 1.296 && 9.9533182 && 1.346 && 9.9501532 && 1.396 && 9.9481631\\ 1.247 && 9.9576198 && 1.297 && 9.9532430 && 1.347 && 9.9501021 && 1.397 && 9.9481349\\ 1.278 && 9.9575197 && 1.298 && 9.9531683 && 1.348 && 9.9500514 && 1.398 && 9.9481071\\ 1.249 && 9.9574201 && 1.299 && 9.9530941 && 1.349 && 9.9500012 && 1.399 && 9.9480798\\ 1.250 && 9.9573211 && 1.300 && 9.9530204 && 1.350 && 9.9499515 && 1.400 && 9.9480529 \end{aligned} \]

\ednote{In the row $a=1.224$ the table prints $9.9599669$ instead of $9.9600669$, and the argument $1.248$ is printed $1.278$; apparent misprints.}

\origpage{304}

\[ \begin{aligned} a. && \text{Log }\Gamma. &\qquad& a. && \text{Log }\Gamma. &\qquad& a. && \text{Log }\Gamma. &\qquad& a. && \text{Log }\Gamma.\\ 1.400 && 9.9480529 && 1.450 && 9.9472677 && 1.500 && 9.9475449 && 1.550 && 9.9488375\\ 1.401 && 9.9480264 && 1.451 && 9.9472630 && 1.501 && 9.9475610 && 1.551 && 9.9488734\\ 1.402 && 9.9480004 && 1.452 && 9.9472587 && 1.502 && 9.9475775 && 1.552 && 9.9489097\\ 1.403 && 9.9479748 && 1.453 && 9.9472548 && 1.503 && 9.9475943 && 1.553 && 9.9489464\\ 1.404 && 9.9479497 && 1.454 && 9.9472514 && 1.504 && 9.9476116 && 1.554 && 9.9489835\\ 1.405 && 9.9479250 && 1.455 && 9.9472484 && 1.505 && 9.9476293 && 1.555 && 9.9490209\\ 1.406 && 9.9479008 && 1.456 && 9.9472458 && 1.506 && 9.9476474 && 1.556 && 9.9490587\\ 1.407 && 9.9478770 && 1.457 && 9.9472436 && 1.507 && 9.9476659 && 1.557 && 9.9490969\\ 1.408 && 9.9478536 && 1.458 && 9.9472419 && 1.508 && 9.9476847 && 1.558 && 9.9491355\\ 1.409 && 9.9478307 && 1.459 && 9.9472406 && 1.509 && 9.9477040 && 1.559 && 9.9491745\\ 1.410 && 9.9478083 && 1.460 && 9.9472397 && 1.510 && 9.9477237 && 1.560 && 9.9492139\\ 1.411 && 9.9477863 && 1.461 && 9.9472392 && 1.511 && 9.9477438 && 1.561 && 9.9492536\\ 1.412 && 9.9477648 && 1.462 && 9.9472392 && 1.512 && 9.9477643 && 1.562 && 9.9492938\\ 1.413 && 9.9477437 && 1.463 && 9.9472396 && 1.513 && 9.9477851 && 1.563 && 9.9493343\\ 1.414 && 9.9477230 && 1.464 && 9.9472404 && 1.514 && 9.9478064 && 1.564 && 9.9493752\\ 1.415 && 9.9477028 && 1.465 && 9.9472416 && 1.515 && 9.9478281 && 1.565 && 9.9494165\\ 1.416 && 9.9476830 && 1.466 && 9.9472432 && 1.516 && 9.9478502 && 1.566 && 9.9494582\\ 1.417 && 9.9476637 && 1.467 && 9.9472453 && 1.517 && 9.9478726 && 1.567 && 9.9495003\\ 1.418 && 9.9476447 && 1.468 && 9.9472477 && 1.518 && 9.9478955 && 1.568 && 9.9495428\\ 1.419 && 9.9476263 && 1.469 && 9.9472506 && 1.519 && 9.9479189 && 1.569 && 9.9495857\\ 1.420 && 9.9476082 && 1.470 && 9.9472539 && 1.520 && 9.9479425 && 1.570 && 9.9496289\\ 1.421 && 9.9475906 && 1.471 && 9.9472576 && 1.521 && 9.9479666 && 1.571 && 9.9496725\\ 1.422 && 9.9475734 && 1.472 && 9.9472618 && 1.522 && 9.9479911 && 1.572 && 9.9497165\\ 1.423 && 9.9475566 && 1.473 && 9.9472664 && 1.523 && 9.9480160 && 1.573 && 9.9497608\\ 1.424 && 9.9475403 && 1.474 && 9.9472713 && 1.524 && 9.9480413 && 1.574 && 9.9498056\\ 1.425 && 9.9475244 && 1.475 && 9.9472767 && 1.525 && 9.9480670 && 1.575 && 9.9498507\\ 1.426 && 9.9475089 && 1.476 && 9.9472825 && 1.526 && 9.9480931 && 1.576 && 9.9498962\\ 1.427 && 9.9474939 && 1.477 && 9.9472887 && 1.527 && 9.9481196 && 1.577 && 9.9499421\\ 1.428 && 9.9474794 && 1.478 && 9.9472953 && 1.528 && 9.9481465 && 1.578 && 9.9499884\\ 1.429 && 9.9474652 && 1.479 && 9.9473023 && 1.529 && 9.9481737 && 1.579 && 9.9500350\\ 1.430 && 9.9474515 && 1.480 && 9.9473097 && 1.530 && 9.9482014 && 1.580 && 9.9500821\\ 1.431 && 9.9474382 && 1.481 && 9.9473175 && 1.531 && 9.9482295 && 1.581 && 9.9501295\\ 1.432 && 9.9474254 && 1.482 && 9.9473258 && 1.532 && 9.9482579 && 1.582 && 9.9501773\\ 1.433 && 9.9474130 && 1.483 && 9.9473345 && 1.533 && 9.9482868 && 1.583 && 9.9502255\\ 1.434 && 9.9474010 && 1.484 && 9.9473436 && 1.534 && 9.9483160 && 1.584 && 9.9502741\\ 1.435 && 9.9473895 && 1.485 && 9.9473531 && 1.535 && 9.9483457 && 1.585 && 9.9503230\\ 1.436 && 9.9473784 && 1.486 && 9.9473631 && 1.536 && 9.9483757 && 1.586 && 9.9503723\\ 1.437 && 9.9473677 && 1.487 && 9.9473734 && 1.537 && 9.9484062 && 1.587 && 9.9504220\\ 1.438 && 9.9473574 && 1.488 && 9.9473842 && 1.538 && 9.9484370 && 1.588 && 9.9504721\\ 1.439 && 9.9473476 && 1.489 && 9.9473953 && 1.539 && 9.9484683 && 1.589 && 9.9505225\\ 1.440 && 9.9473382 && 1.490 && 9.9474069 && 1.540 && 9.9484999 && 1.590 && 9.9505733\\ 1.441 && 9.9473292 && 1.491 && 9.9474189 && 1.541 && 9.9485319 && 1.591 && 9.9506245\\ 1.442 && 9.9473207 && 1.492 && 9.9474312 && 1.542 && 9.9485643 && 1.592 && 9.9506761\\ 1.443 && 9.9473126 && 1.493 && 9.9474440 && 1.543 && 9.9485971 && 1.593 && 9.9507280\\ 1.444 && 9.9473049 && 1.494 && 9.9474572 && 1.544 && 9.9486302 && 1.594 && 9.9507803\\ 1.445 && 9.9472976 && 1.495 && 9.9474708 && 1.545 && 9.9486638 && 1.595 && 9.9508330\\ 1.446 && 9.9472908 && 1.496 && 9.9474848 && 1.546 && 9.9486978 && 1.596 && 9.9508860\\ 1.447 && 9.9472844 && 1.497 && 9.9474992 && 1.547 && 9.9487321 && 1.597 && 9.9509395\\ 1.448 && 9.9472784 && 1.498 && 9.9475140 && 1.548 && 9.9487669 && 1.598 && 9.9509933\\ 1.449 && 9.9472728 && 1.499 && 9.9475293 && 1.549 && 9.9488020 && 1.599 && 9.9510474\\ 1.450 && 9.9472677 && 1.500 && 9.9475449 && 1.550 && 9.9488375 && 1.600 && 9.9511020 \end{aligned} \]

\origpage{305}

\[ \begin{aligned} a. && \text{Log }\Gamma. &\qquad& a. && \text{Log }\Gamma. &\qquad& a. && \text{Log }\Gamma. &\qquad& a. && \text{Log }\Gamma.\\ 1.600 && 9.9511020 && 1.650 && 9.9542989 && 1.700 && 9.9583911 && 1.750 && 9.9633451\\ 1.601 && 9.9511569 && 1.651 && 9.9543721 && 1.701 && 9.9584819 && 1.751 && 9.9634527\\ 1.602 && 9.9512123 && 1.652 && 9.9544456 && 1.702 && 9.9585730 && 1.752 && 9.9635607\\ 1.603 && 9.9512679 && 1.653 && 9.9545195 && 1.703 && 9.9586644 && 1.753 && 9.9636689\\ 1.604 && 9.9513240 && 1.654 && 9.9545938 && 1.704 && 9.9587562 && 1.754 && 9.9637775\\ 1.605 && 9.9513804 && 1.655 && 9.9546684 && 1.705 && 9.9588483 && 1.755 && 9.9638865\\ 1.606 && 9.9514372 && 1.656 && 9.9547434 && 1.706 && 9.9589408 && 1.756 && 9.9639958\\ 1.607 && 9.9514943 && 1.657 && 9.9548187 && 1.707 && 9.9590335 && 1.757 && 9.9641054\\ 1.608 && 9.9515518 && 1.658 && 9.9548944 && 1.708 && 9.9591267 && 1.758 && 9.9642154\\ 1.609 && 9.9516097 && 1.659 && 9.9549704 && 1.709 && 9.9592202 && 1.759 && 9.9643257\\ 1.610 && 9.9516680 && 1.660 && 9.9550468 && 1.710 && 9.9593140 && 1.760 && 9.9644363\\ 1.611 && 9.9517267 && 1.661 && 9.9551236 && 1.711 && 9.9594082 && 1.761 && 9.9645473\\ 1.612 && 9.9517857 && 1.662 && 9.9552007 && 1.712 && 9.9595027 && 1.762 && 9.9646586\\ 1.613 && 9.9518451 && 1.663 && 9.9552782 && 1.713 && 9.9595975 && 1.763 && 9.9647702\\ 1.614 && 9.9519049 && 1.664 && 9.9553560 && 1.714 && 9.9596928 && 1.764 && 9.9648821\\ 1.615 && 9.9519650 && 1.665 && 9.9554342 && 1.715 && 9.9597883 && 1.765 && 9.9649944\\ 1.616 && 9.9520255 && 1.666 && 9.9555127 && 1.716 && 9.9598842 && 1.766 && 9.9651070\\ 1.617 && 9.9520863 && 1.667 && 9.9565916 && 1.717 && 9.9599805 && 1.767 && 9.9652199\\ 1.618 && 9.9521476 && 1.668 && 9.9556709 && 1.718 && 9.9600770 && 1.768 && 9.9653331\\ 1.619 && 9.9522092 && 1.669 && 9.9557505 && 1.719 && 9.9601739 && 1.769 && 9.9654467\\ 1.620 && 9.9522711 && 1.670 && 9.9558304 && 1.720 && 9.9602712 && 1.770 && 9.9655606\\ 1.621 && 9.9523334 && 1.671 && 9.9559107 && 1.721 && 9.9603688 && 1.771 && 9.9656748\\ 1.622 && 9.9523961 && 1.672 && 9.9559914 && 1.722 && 9.9604667 && 1.772 && 9.9657894\\ 1.623 && 9.9524591 && 1.673 && 9.9560724 && 1.723 && 9.9605650 && 1.773 && 9.9659043\\ 1.624 && 9.9525225 && 1.674 && 9.9561537 && 1.724 && 9.9606636 && 1.774 && 9.9660195\\ 1.625 && 9.9525863 && 1.675 && 9.9562354 && 1.725 && 9.9607625 && 1.775 && 9.9661350\\ 1.626 && 9.9526504 && 1.676 && 9.9563174 && 1.726 && 9.9608618 && 1.776 && 9.9662509\\ 1.627 && 9.9527149 && 1.677 && 9.9563998 && 1.727 && 9.9609614 && 1.777 && 9.9663671\\ 1.628 && 9.9527798 && 1.678 && 9.9564826 && 1.728 && 9.9610614 && 1.778 && 9.9664836\\ 1.629 && 9.9528450 && 1.679 && 9.9565657 && 1.729 && 9.9611617 && 1.779 && 9.9666004\\ 1.630 && 9.9529106 && 1.680 && 9.9566491 && 1.730 && 9.9612623 && 1.780 && 9.9667176\\ 1.631 && 9.9529766 && 1.681 && 9.9567329 && 1.731 && 9.9613632 && 1.781 && 9.9668351\\ 1.632 && 9.9530429 && 1.682 && 9.9568171 && 1.732 && 9.9614645 && 1.782 && 9.9669529\\ 1.633 && 9.9531096 && 1.683 && 9.9569016 && 1.733 && 9.9615661 && 1.783 && 9.9670711\\ 1.634 && 9.9531767 && 1.684 && 9.9569864 && 1.734 && 9.9616681 && 1.784 && 9.9671895\\ 1.635 && 9.9532441 && 1.685 && 9.9570716 && 1.735 && 9.9617704 && 1.785 && 9.9673083\\ 1.636 && 9.9533119 && 1.686 && 9.9571572 && 1.736 && 9.9618731 && 1.786 && 9.9674274\\ 1.637 && 9.9533801 && 1.687 && 9.9572430 && 1.737 && 9.9619760 && 1.787 && 9.9675468\\ 1.638 && 9.9534487 && 1.688 && 9.9573293 && 1.738 && 9.9620794 && 1.788 && 9.9676665\\ 1.639 && 9.9535176 && 1.689 && 9.9574159 && 1.739 && 9.9621830 && 1.789 && 9.9677866\\ 1.640 && 9.9535868 && 1.690 && 9.9575028 && 1.740 && 9.9622870 && 1.790 && 9.9679070\\ 1.641 && 9.9536564 && 1.691 && 9.9575901 && 1.741 && 9.9623913 && 1.791 && 9.9680277\\ 1.642 && 9.9537263 && 1.692 && 9.9576777 && 1.742 && 9.9624960 && 1.792 && 9.9681487\\ 1.643 && 9.9537966 && 1.693 && 9.9577656 && 1.743 && 9.9626009 && 1.793 && 9.9682701\\ 1.644 && 9.9538673 && 1.694 && 9.9578539 && 1.744 && 9.9627063 && 1.794 && 9.9683918\\ 1.645 && 9.9539383 && 1.695 && 9.9579426 && 1.745 && 9.9628119 && 1.795 && 9.9685138\\ 1.646 && 9.9540097 && 1.696 && 9.9580316 && 1.746 && 9.9629179 && 1.796 && 9.9686361\\ 1.647 && 9.9540815 && 1.697 && 9.9581210 && 1.747 && 9.9630242 && 1.797 && 9.9687588\\ 1.648 && 9.9541536 && 1.698 && 9.9582107 && 1.748 && 9.9631308 && 1.798 && 9.9688818\\ 1.649 && 9.9542261 && 1.699 && 9.9583007 && 1.749 && 9.9632378 && 1.799 && 9.9690051\\ 1.650 && 9.9542989 && 1.700 && 9.9583911 && 1.750 && 9.9633451 && 1.800 && 9.9691287 \end{aligned} \]

\ednote{In the row $a=1.667$ the table prints $9.9565916$ instead of $9.9555916$; an apparent misprint.}

\origpage{306}

\[ \begin{aligned} a. && \text{Log }\Gamma. &\qquad& a. && \text{Log }\Gamma. &\qquad& a. && \text{Log }\Gamma. &\qquad& a. && \text{Log }\Gamma.\\ 1.800 && 9.9691287 && 1.850 && 9.9757126 && 1.900 && 9.9830693 && 1.950 && 9.9911732\\ 1.801 && 9.9692526 && 1.851 && 9.9758522 && 1.901 && 9.9832241 && 1.951 && 9.9913427\\ 1.802 && 9.9693769 && 1.852 && 9.9759922 && 1.902 && 9.9833793 && 1.952 && 9.9915125\\ 1.803 && 9.9695015 && 1.853 && 9.9761325 && 1.903 && 9.9835347 && 1.953 && 9.9916826\\ 1.804 && 9.9696264 && 1.854 && 9.9762730 && 1.904 && 9.9836905 && 1.954 && 9.9918530\\ 1.805 && 9.9697516 && 1.855 && 9.9764139 && 1.905 && 9.9838465 && 1.955 && 9.9920237\\ 1.806 && 9.9698771 && 1.856 && 9.9765551 && 1.906 && 9.9840028 && 1.956 && 9.9921947\\ 1.807 && 9.9700030 && 1.857 && 9.9766966 && 1.907 && 9.9841594 && 1.957 && 9.9923659\\ 1.808 && 9.9701291 && 1.858 && 9.9768384 && 1.908 && 9.9843164 && 1.958 && 9.9925374\\ 1.809 && 9.9702556 && 1.859 && 9.9769805 && 1.909 && 9.9844736 && 1.959 && 9.9927093\\ 1.810 && 9.9703824 && 1.860 && 9.9771229 && 1.910 && 9.9846311 && 1.960 && 9.9928814\\ 1.811 && 9.9705095 && 1.861 && 9.9772657 && 1.911 && 9.9847889 && 1.961 && 9.9930538\\ 1.812 && 9.9706369 && 1.862 && 9.9774087 && 1.912 && 9.9849471 && 1.962 && 9.9932266\\ 1.813 && 9.9707647 && 1.863 && 9.9775521 && 1.913 && 9.9851055 && 1.963 && 9.9933996\\ 1.814 && 9.9708927 && 1.864 && 9.9776957 && 1.914 && 9.9852642 && 1.964 && 9.9935728\\ 1.815 && 9.9710211 && 1.865 && 9.9778397 && 1.915 && 9.9854232 && 1.965 && 9.9937464\\ 1.816 && 9.9711498 && 1.866 && 9.9779840 && 1.916 && 9.9855825 && 1.966 && 9.9939202\\ 1.817 && 9.9712788 && 1.867 && 9.9781286 && 1.917 && 9.9857421 && 1.967 && 9.9940943\\ 1.818 && 9.9714081 && 1.868 && 9.9782734 && 1.918 && 9.9859020 && 1.968 && 9.9942687\\ 1.819 && 9.9715377 && 1.869 && 9.9784186 && 1.919 && 9.9860621 && 1.969 && 9.9944434\\ 1.820 && 9.9716677 && 1.870 && 9.9785641 && 1.920 && 9.9862226 && 1.970 && 9.9946184\\ 1.821 && 9.9717980 && 1.871 && 9.9787099 && 1.921 && 9.9863834 && 1.971 && 9.9947937\\ 1.822 && 9.9719286 && 1.872 && 9.9788559 && 1.922 && 9.9865444 && 1.972 && 9.9949692\\ 1.823 && 9.9720595 && 1.873 && 9.9790023 && 1.923 && 9.9867058 && 1.973 && 9.9951451\\ 1.824 && 9.9721907 && 1.874 && 9.9791490 && 1.924 && 9.9868675 && 1.974 && 9.9953212\\ 1.825 && 9.9723223 && 1.875 && 9.9792960 && 1.925 && 9.9870294 && 1.975 && 9.9954976\\ 1.826 && 9.9724542 && 1.876 && 9.9794433 && 1.926 && 9.9871916 && 1.976 && 9.9956743\\ 1.827 && 9.9725864 && 1.877 && 9.9795909 && 1.927 && 9.9873542 && 1.977 && 9.9958513\\ 1.828 && 9.9727189 && 1.878 && 9.9797389 && 1.928 && 9.9875170 && 1.978 && 9.9960286\\ 1.829 && 9.9728517 && 1.879 && 9.9798871 && 1.929 && 9.9876801 && 1.979 && 9.9962062\\ 1.830 && 9.9729848 && 1.880 && 9.9800356 && 1.930 && 9.9878435 && 1.980 && 9.9963840\\ 1.831 && 9.9731182 && 1.881 && 9.9801844 && 1.931 && 9.9880072 && 1.981 && 9.9965621\\ 1.832 && 9.9732520 && 1.882 && 9.9803335 && 1.932 && 9.9881712 && 1.982 && 9.9967405\\ 1.833 && 9.9733860 && 1.883 && 9.9804829 && 1.933 && 9.9883355 && 1.983 && 9.9969192\\ 1.834 && 9.9735204 && 1.884 && 9.9806327 && 1.934 && 9.9885001 && 1.984 && 9.9970982\\ 1.835 && 9.9736551 && 1.885 && 9.9807827 && 1.935 && 9.9886650 && 1.985 && 9.9972774\\ 1.836 && 9.9737901 && 1.886 && 9.9809330 && 1.936 && 9.9888302 && 1.986 && 9.9974569\\ 1.837 && 9.9739254 && 1.887 && 9.9810837 && 1.937 && 9.9889956 && 1.987 && 9.9976367\\ 1.838 && 9.9740610 && 1.888 && 9.9812346 && 1.938 && 9.9891614 && 1.988 && 9.9978168\\ 1.839 && 9.9741970 && 1.889 && 9.9813859 && 1.939 && 9.9893275 && 1.989 && 9.9979972\\ 1.840 && 9.9743332 && 1.890 && 9.9815374 && 1.940 && 9.9894938 && 1.990 && 9.9981779\\ 1.841 && 9.9744697 && 1.891 && 9.9816892 && 1.941 && 9.9896604 && 1.991 && 9.9983589\\ 1.842 && 9.9746066 && 1.892 && 9.9818414 && 1.942 && 9.9898274 && 1.992 && 9.9985401\\ 1.843 && 9.9747437 && 1.893 && 9.9819938 && 1.943 && 9.9899946 && 1.993 && 9.9987216\\ 1.844 && 9.9748812 && 1.894 && 9.9821466 && 1.944 && 9.9901621 && 1.994 && 9.9989034\\ 1.845 && 9.9750190 && 1.895 && 9.9822996 && 1.945 && 9.9903299 && 1.995 && 9.9990855\\ 1.846 && 9.9751571 && 1.896 && 9.9824529 && 1.946 && 9.9904980 && 1.996 && 9.9992678\\ 1.847 && 9.9752955 && 1.897 && 9.9826066 && 1.947 && 9.9906664 && 1.997 && 9.9994504\\ 1.848 && 9.9754342 && 1.898 && 9.9827605 && 1.948 && 9.9908350 && 1.998 && 9.9996333\\ 1.849 && 9.9755733 && 1.899 && 9.9829148 && 1.949 && 9.9910040 && 1.999 && 9.9998165\\ 1.850 && 9.9757126 && 1.900 && 9.9830693 && 1.950 && 9.9911732 && 2.000 && 0.0000000 \end{aligned} \]

\origpage{307}

\section*{TROISIÈME PARTIE.}

\section*{DES QUADRATURES.}

Diverses méthodes relatives aux quadratures, sont l'objet de cette troisième partie: ces méthodes complètent naturellement la théorie des transcendantes, puisque si les transcendantes sont trop composées, ou si elles ne peuvent se réduire à celles qui sont données par les Tables, il faut nécessairement, pour les évaluer dans les cas particuliers, avoir recours aux quadratures.

La méthode que je propose comme la plus générale et la plus sûre, consiste à exprimer l'intégrale cherchée aux différences infiniment petites, par une intégrale aux différences finies à laquelle on ajoute les corrections que l'analyse indique et qui servent à diriger l'approximation. On peut parvenir de cette manière à des résultats dont l'exactitude s'étende jusqu'à tel ordre de décimales qu'on voudra.

La même méthode considérée sous un autre point de vue, peut servir à construire une courbe dont les coordonnées dépendent chacune d'une quadrature particulière. J'ai donné pour exemple en ce genre, le calcul de la trajectoire d'un projectile dans un milieu résistant, et j'ai par ce moyen poussé l'approximation plus loin que je ne l'avais fait dans la pièce couronnée par l'Académie de Berlin en 1782.

On trouve dans les Mémoires de l'Académie des Sciences, années 1778 et 1782, une méthode fort ingénieuse pour avoir la valeur de l'intégrale $\displaystyle\int y\,dx$, dans le cas où la fonction $y$ est nulle aux deux limites de l'intégrale. J'ai exposé cette méthode avec les développemens qu'elle exige dans quelques exemples, et principalement

\origpage{308}
dans son application au cas où les limites de l'intégrale sont supposées imaginaires.

Enfin j'ai profité de la liberté que me donnait le titre de cet ouvrage, pour traiter de diverses sortes d'intégrales définies qui appartiennent à la théorie des transcendantes. De ce nombre sont plusieurs formules que Laplace a données dans différens recueils, et qui tiennent un rang distingué dans cette théorie; les autres sont, pour la plupart, extraites des ouvrages d'Euler, et particulièrement des Supplémens qui composent le tome IV de son \emph{Calcul intégral}.

Je me flatte que ces divers objets réunis sous un même point de vue, pourront intéresser les Géomètres; ils ont été d'ailleurs choisis de manière que chacun d'eux puisse offrir de nouvelles formules ou de nouvelles démonstrations.

\section*{Formule générale pour les quadratures.}

\subsection*{(1).}

On est censé avoir résolu un problème lorsqu'il est réduit aux quadratures, c'est-à-dire lorsqu'il ne dépend plus que d'une ou de plusieurs intégrales de la forme $\displaystyle\int y\,dx$, où l'on connaît $y$ en fonction de $x$.

Si les intégrales dont il s'agit sont susceptibles d'être exprimées par des suites convergentes, la solution est aussi complète qu'on peut le desirer, eu égard à la nature de la question; mais le plus souvent on a des expressions trop compliquées pour les intégrer par des suites convergentes, et il ne reste d'autre parti à prendre que de chercher leurs valeurs dans les cas particuliers seulement, c'est-à-dire pour une partie donnée de la ligne des abscisses.

Soit donc proposé de trouver la valeur de l'intégrale $\displaystyle Z=\int y\,dx$, depuis $x=a$ jusqu'à $x=b$, ou seulement depuis $x=0$ jusqu'à $x=a$; car on peut supposer que l'origine des abscisses est prise au point où commence l'aire, de manière que $Z$ et $x$ s'évanouissent en même temps.

Pour obtenir plus facilement le degré d'approximation qu'on peut desirer, nous supposerons,

1°. Que dans l'intervalle donné depuis $x=0$ jusqu'à $x=a$, la

\origpage{309}
fonction $y$ reste constamment de même signe, et nous regarderons ce signe comme positif. S'il en était autrement, on chercherait successivement les différentes parties de l'aire $\displaystyle\int y\,dx$, situées de différens côtés de la ligne des abscisses, et on retrancherait la somme des aires négatives de la somme des aires positives.

2°. Que la courbure de la courbe, depuis $x=0$ jusqu'à $x=a$, n'éprouve aucune variation assez brusque pour que l'un des coefficiens $\frac{dy}{dx}$, $\frac{ddy}{dx^{2}}$, $\frac{d^{3}y}{dx^{3}}$ devienne infini. Ainsi dans toute l'étendue de la portion de courbe que nous voulons quarrer, l'ordonnée qui répond à une abscisse $x+\alpha$, très-peu différente de $x$, pourra s'exprimer avec une exactitude suffisante par la suite $y+\alpha\frac{dy}{dx}+\frac{\alpha^{2}}{2}\frac{ddy}{dx^{2}}+\frac{\alpha^{3}}{2\cdot3}\cdot\frac{d^{3}y}{dx^{3}}+\text{etc.}$, sans qu'il y ait lieu à exception pour aucun point. Cette condition est surtout nécessaire aux deux limites de l'intégrale; elle est moins importante dans les points intermédiaires, parce que l'erreur due à cette cause est limitée, et qu'elle peut être corrigée assez facilement, comme on le verra ci-après.

Cela posé, soit $y=F(x)$, et $x=n\omega$, $n$ étant un nombre entier d'autant plus grand, qu'on voudra obtenir une plus grande approximation; si l'on désigne par $\Sigma F(x+\frac{1}{2}\omega)$ la somme de la suite
\[ F(\tfrac{1}{2}\omega)+F(\tfrac{3}{2}\omega)+F(\tfrac{5}{2}\omega)\ldots+F(x-\tfrac{1}{2}\omega), \]
il est visible qu'on aura d'abord à très-peu près $Z=\omega\Sigma F(x+\frac{1}{2}\omega)$. Pour avoir une valeur plus exacte, soit
\[ Z=\omega\Sigma F(x+\tfrac{1}{2}\omega)+\zeta, \]
on aura en faisant varier $x$ de $\omega$, et prenant les différences finies de chaque membre,
\[ \Delta\zeta=\Delta Z-\omega F(x+\tfrac{1}{2}\omega). \]
Développant le second membre suivant le théorème de Taylor, et observant que $\frac{dZ}{dx}=y=F(x)$, on aura
\[ \begin{aligned} \Delta\zeta={}&\omega F+\frac{\omega^{2}}{2}\cdot\frac{dF}{dx}+\frac{\omega^{3}}{2\cdot3}\cdot\frac{ddF}{dx^{2}}+\frac{\omega^{4}}{2\cdot3\cdot4}\cdot\frac{d^{3}F}{dx^{3}}+\text{etc.}\\ &-\omega F-\omega^{2}\cdot\frac{1}{2}\frac{dF}{dx}-\frac{\omega^{3}}{2}\cdot\frac{1}{4}\cdot\frac{ddF}{dx^{2}}-\frac{\omega^{4}}{2\cdot3}\cdot\frac{1}{8}\cdot\frac{d^{3}F}{dx^{3}}-\text{etc.}, \end{aligned} \]

\origpage{310}
ou en réduisant,
\[ \Delta\zeta=\frac{\omega^{3}}{2}\left(\frac{1}{3}-\frac{1}{4}\right)\frac{ddF}{dx^{2}}+\frac{\omega^{4}}{2\cdot3}\left(\frac{1}{4}-\frac{1}{8}\right)\frac{d^{3}F}{dx^{3}}+\frac{\omega^{5}}{2\cdot3\cdot4}\left(\frac{1}{5}-\frac{1}{16}\right)\frac{d^{4}F}{dx^{4}}+\text{etc.} \]
On voit que le premier terme de la valeur de $\zeta$ est $\frac{\omega^{2}}{2}\cdot\frac{1}{12}\cdot\frac{dF}{dx}$; ensorte qu'on peut faire $\zeta=\frac{\omega^{2}}{24}\cdot\frac{dF}{dx}+\zeta'$; on trouverait ensuite que le premier terme de la valeur de $\zeta'$ est $-\frac{7\omega^{4}}{5760}\cdot\frac{d^{3}F}{dx^{3}}$. Mais pour connaître la loi de la suite qui exprime $\zeta$, nous laisserons indéterminés les coefficiens de ses différens termes, et nous supposerons
\[ \begin{aligned} Z={}&\omega\Sigma F(x+\tfrac{1}{2}\omega)+\text{const.}\\ &+a'\omega^{2}\frac{dF}{dx}+b'\omega^{3}\frac{ddF}{dx^{2}}+c'\omega^{4}\frac{d^{3}F}{dx^{3}}+f'\omega^{5}\frac{d^{4}F}{dx^{4}}+\text{etc.} \end{aligned} \]

\subsection*{(2).}

Il suffira de déterminer les coefficiens $a'$, $b'$, $c'$, etc. dans un cas particulier; soit donc $F(x)=e^{x}$, on aura $\frac{dF}{dx}=e^{x}$, $\frac{ddF}{dx^{2}}=e^{x}$, etc., $\displaystyle Z=\int e^{x}\,dx=e^{x}-1$, $\Sigma F(x+\frac{1}{2}\omega)=e^{\frac{1}{2}\omega}\Sigma e^{x}=\frac{\omega e^{\frac{1}{2}\omega}}{e^{\omega}-1}(e^{x}-1)$,\ednote{The factor $\omega$ in this numerator belongs to $\omega\Sigma F(x+\frac{1}{2}\omega)$, the quantity substituted next, not to $\Sigma F(x+\frac{1}{2}\omega)$; as printed.} et la substitution de ces valeurs donnera l'équation suivante qui doit être identique:
\[ e^{x}-1=\frac{\omega(e^{x}-1)}{e^{\frac{1}{2}\omega}-e^{-\frac{1}{2}\omega}}+\text{const.}+(a'\omega^{2}+b'\omega^{3}+c'\omega^{4}+\text{etc.})e^{x}. \]
Faisant $x=0$, on trouve la constante $=-a'\omega^{2}-b'\omega^{3}-c'\omega^{4}-\text{etc.}$; de sorte qu'en divisant toute l'équation par $e^{x}-1$, il viendra
\[ \frac{\omega}{e^{\frac{1}{2}\omega}-e^{-\frac{1}{2}\omega}}=1-a'\omega^{2}-b'\omega^{3}-c'\omega^{4}-f'\omega^{5}-\text{etc.} \]

Le premier membre est une fonction paire de $\omega$, puisqu'il reste le même en changeant le signe de $\omega$; donc dans le second membre, tous les coefficiens des puissances impaires de $\omega$ sont nuls. Pour avoir égard à cette circonstance, nous ferons de nouveau
\[ \frac{\omega}{e^{\frac{1}{2}\omega}-e^{-\frac{1}{2}\omega}}=1-A\omega^{2}+B\omega^{4}-C\omega^{6}+D\omega^{8}-\text{etc.}, \tag{1} \]

\origpage{311}
et la valeur de $Z$ sera en général
\[ \begin{aligned} Z={}&\omega\Sigma F(x+\tfrac{1}{2}\omega)+A\omega^{2}\frac{dF}{dx}-B\omega^{4}\frac{d^{3}F}{dx^{3}}+C\omega^{6}\frac{d^{5}F}{dx^{5}}-\text{etc.}\\ &+\text{const.} \end{aligned} \]
Désignons par $\frac{dF^{0}}{dx}$, $\frac{ddF^{0}}{dx^{2}}$, etc., ce que deviennent les coefficiens $\frac{dF}{dx}$, $\frac{ddF}{dx^{2}}$, etc. lorsque $x=0$, et la constante étant déterminée de manière que $Z$ et $x$ s'évanouissent en même temps, on aura
\[ \begin{aligned} Z={}&\omega\Sigma F(x+\tfrac{1}{2}\omega)+A\omega^{2}\left(\frac{dF}{dx}-\frac{dF^{0}}{dx^{0}}\right)\\ &-B\omega^{4}\left(\frac{d^{3}F}{dx^{3}}-\frac{d^{3}F^{0}}{dx^{3}}\right)\\ &+C\omega^{6}\left(\frac{d^{5}F}{dx^{5}}-\frac{d^{5}F^{0}}{dx^{5}}\right)\\ &-\text{etc.}: \end{aligned} \tag{2} \]
\ednote{Printed $dx^{0}$ in the denominator of the second fraction of the first line; $dx$ is meant, as in the other terms. An apparent misprint.}
c'est l'intégrale demandée prise depuis $x=0$ jusqu'à $x=n\omega$.

\subsection*{(3).}

Si dans l'équation (1) on met $-\omega^{2}$ à la place de $\omega^{2}$, il en résultera
\[ \frac{\tfrac{1}{2}\omega}{\sin\tfrac{1}{2}\omega}=1+A\omega^{2}+B\omega^{4}+C\omega^{6}+D\omega^{8}+\text{etc.} \]
Le premier membre peut se mettre sous la forme
\[ 1+\frac{2\omega^{2}}{4\pi^{2}-\omega^{2}}-\frac{2\omega^{2}}{16\pi^{2}-\omega^{2}}+\frac{2\omega^{2}}{36\pi^{2}-\omega^{2}}-\text{etc.}, \]
(\emph{Introd. in An. inf.}, page 142), et par son développement on a
\[ \begin{aligned} A&=\frac{2}{4\pi^{2}}\left(1-\frac{1}{2^{2}}+\frac{1}{3^{2}}-\frac{1}{4^{2}}+\text{etc.}\right)=\frac{1}{24}\\ B&=\frac{2}{16\pi^{4}}\left(1-\frac{1}{2^{4}}+\frac{1}{3^{4}}-\frac{1}{4^{4}}+\text{etc.}\right)=\frac{7}{5760}\\ C&=\frac{2}{64\pi^{6}}\left(1-\frac{1}{2^{6}}+\frac{1}{3^{6}}-\frac{1}{4^{6}}+\text{etc.}\right)=\frac{31}{967680}\\ D&=\frac{2}{256\pi^{8}}\left(1-\frac{1}{2^{8}}+\frac{1}{3^{8}}-\frac{1}{4^{8}}+\text{etc.}\right)=\frac{127}{(128)^{2}\cdot9450}\\ &\text{etc.} \end{aligned} \]

\origpage{312}
D'où l'on voit que tous les coefficiens $A$, $B$, $C$, etc. sont positifs, et qu'ils décroissent suivant une raison qui approche de plus en plus de $\frac{1}{4\pi^{2}}$, de sorte que chaque terme sera environ le quarantième du précédent.

La manière la plus simple de calculer les valeurs numériques de ces coefficiens, est de les déduire du développement de la fraction
\[ \frac{1}{1-\frac{1}{2\cdot3}\cdot\frac{\omega^{2}}{4}+\frac{1}{2\cdot3\cdot4\cdot5}\cdot\frac{\omega^{4}}{16}-\text{etc.}}=1+A\omega^{2}+B\omega^{4}+C\omega^{6}+\text{etc.}, \]
d'où l'on tire
\[ \begin{aligned} A&=\frac{1}{2\cdot3}\cdot\frac{1}{4}\\ B&=\frac{1}{2\cdot3}\cdot\frac{1}{4}A-\frac{1}{2\cdot3\cdot4\cdot5}\cdot\frac{1}{16}\\ C&=\frac{1}{2\cdot3}\cdot\frac{1}{4}B-\frac{1}{2\cdot3\cdot4\cdot5}\cdot\frac{1}{16}A+\frac{1}{2\cdot3\cdot4\cdot5\cdot6\cdot7}\cdot\frac{1}{64}\\ &\text{etc.} \end{aligned} \]
Si on appelle $S_{n}$ la somme de la suite $1+\frac{1}{2^{n}}+\frac{1}{3^{n}}+\frac{1}{4^{n}}+\text{etc.}$, on aura encore
\[ \begin{aligned} A&=\frac{1}{2\pi^{2}}\left(1-\frac{1}{2}\right)S_{2}\\ B&=\frac{1}{2^{3}\pi^{4}}\left(1-\frac{1}{2^{3}}\right)S_{4}\\ C&=\frac{1}{2^{5}\pi^{6}}\left(1-\frac{1}{2^{5}}\right)S_{6}\\ &\text{etc.} \end{aligned} \]
Et en général, $N$ étant le $n^{\text{ieme}}$ terme de la suite $A$, $B$, $C$, etc., on aura
\[ N=\frac{1}{2^{2n-1}\pi^{2n}}\left(1-\frac{1}{2^{2n-1}}\right)S_{2n}. \]

Nous connaissons maintenant la loi des termes de l'équation (2); mais pour que cette formule soit employée avec succès à la détermination de l'aire $Z$, il faudra prendre $\omega$ assez petit pour que le terme $B\omega^{4}\left(\frac{d^{3}F}{dx^{3}}-\frac{d^{3}F^{0}}{dx^{3}}\right)$ et les suivans puissent être négligés;

\origpage{313}
et c'est ce qu'il est toujours facile de faire, à moins que $\frac{d^{3}F}{dx^{3}}$ ne soit infini à l'une ou à l'autre limite de l'intégrale.

Ainsi ce cas excepté, et à plus forte raison celui où $\frac{dF}{dx}$ deviendrait infini à l'une de ces limites, on déterminera l'aire $Z$ par la formule
\[ Z=\omega\Sigma F(x+\tfrac{1}{2}\omega)+\frac{\omega^{2}}{24}\left(\frac{dF}{dx}-\frac{dF^{0}}{dx}\right), \]
la seconde partie étant la correction de la première.

\subsection*{(4).}

Si le coefficient $\frac{dF}{dx}$ était infini à l'une des limites de l'intégrale, c'est-à-dire, si dans l'un de ces points, l'ordonnée était tangente à la courbe, il faudrait chercher par un autre procédé l'aire comprise entre cette ordonnée et une ordonnée peu éloignée; le reste de l'aire ne serait sujet à aucune difficulté.

Supposons, par exemple, que $\frac{dy}{dx}$ soit infini à la seconde limite de l'intégrale, et soit $y=b$ l'ordonnée qui est tangente à la courbe en ce point. Si en faisant $x=a-\alpha$, on a $y=c$, alors l'aire comprise entre les deux ordonnées $c$, $b$, sera à très-peu près $\frac{\alpha}{3}(2c+b)$. C'est ce qui résulte de la supposition que l'arc de courbe dont il s'agit peut être assimilé à un arc de parabole dont l'axe est parallèle à la ligne des abscisses.

Il faudra donc joindre à la quantité $\frac{\alpha}{3}(2c+b)$, l'aire comprise depuis $x=0$ jusqu'à $x=a-\alpha$.

Une connaissance plus intime de la nature de la courbe, pourra conduire à une valeur plus approchée de l'aire comprise depuis $x=a-\alpha$ jusqu'à $x=a$; mais la détermination précédente suffira dans presque tous les cas, et le procédé serait le même si $\frac{dF}{dx}$ était infini à la première limite.

De même si $\frac{dF}{dx}$ n'est pas infini à l'une des limites de l'intégrale, mais que $\frac{d^{3}F}{dx^{3}}$ le soit, ou qu'il ait une valeur très-grande; alors il conviendra de calculer d'une manière particulière l'aire de la portion

\origpage{314}
de courbe, dans laquelle ce coefficient a une valeur trop grande pour que la formule (2) s'y applique avec sureté.

\subsection*{(5).}

Cherchons maintenant quelle erreur peut résulter de notre formule, lorsqu'entre les limites de l'intégrale il y a un point où la valeur de $x$ rend infini soit le coefficient $\frac{dy}{dx}$, soit l'un des deux suivans $\frac{ddy}{dx^{2}}$, $\frac{d^{3}y}{dx^{3}}$; j'observe que si $x=\alpha$ et $y=\beta$ sont les coordonnées de ce point singulier, on devra avoir en général
\[ y=\beta+A(x-\alpha)^{\mu}+B(x-\alpha)^{\mu+\nu}+C(x-\alpha)^{\mu+2\nu}+\text{etc.}; \]
$\mu$ étant un nombre fractionnaire positif, et $\nu$ un nombre également positif, mais qui peut être supposé égal à l'unité, excepté dans des cas extraordinaires. J'observe de plus que comme la courbe est supposée former une branche continue depuis $x=0$ jusqu'à $x=a$, $a$ étant $>\alpha$, il faut que $\mu$ soit représenté par une fraction $\frac{p}{q}$ dont le dénominateur sera toujours impair, afin que la supposition de $x-\alpha$ négatif ne rende pas imaginaire $y-\beta$.

Supposons maintenant que l'intervalle $\omega$ qui est une partie aliquote de $a$, en soit une aussi de $\alpha$, ce qui est toujours possible en retranchant une petite partie de l'aire proposée, et ajoutant ensuite la partie retranchée; l'intégrale finie $\Sigma\omega F(x+\frac{1}{2}\omega)$ contiendra les deux termes $\omega F(\alpha-\frac{1}{2}\omega)+\omega F(\alpha+\frac{1}{2}\omega)$ qui répondront à la partie d'aire comprise depuis $x=\alpha-\omega$ jusqu'à $x=\alpha+\omega$, et la formule (2) supposera que cette partie d'aire $\zeta$ est représentée par la formule
\[ \zeta=\omega F(\alpha-\tfrac{1}{2}\omega)+\omega F(\alpha+\tfrac{1}{2}\omega)+\frac{\omega^{2}}{24}\left(\frac{dF''}{dx}-\frac{dF'}{dx}\right), \]
$\frac{dF'}{dx}$ et $\frac{dF''}{dx}$ étant les valeurs de $\frac{dy}{dx}$ qui répondent aux abscisses $\alpha-\omega$ et $\alpha+\omega$. Or la valeur supposée de $y$ donne
\[ \begin{aligned} F(\alpha+\tfrac{1}{2}\omega)&=\beta+A(\tfrac{1}{2}\omega)^{\mu}+B(\tfrac{1}{2}\omega)^{\mu+\nu}+\text{etc.}\\ F(\alpha-\tfrac{1}{2}\omega)&=\beta+A(-\tfrac{1}{2}\omega)^{\mu}+B(-\tfrac{1}{2}\omega)^{\mu+\nu}+\text{etc.}\\ \frac{dF''}{dx}&=A\mu\omega^{\mu-1}+B(\mu+\nu)\omega^{\mu+\nu-1}+\text{etc.}\\ \frac{dF'}{dx}&=A\mu(-\omega)^{\mu-1}+B(\mu+\nu)(-\omega)^{\mu+\nu-1}+\text{etc.} \end{aligned} \]

\origpage{315}
Donc
\[ \begin{aligned} \zeta={}&2\beta\omega+A\omega^{\mu+1}\left(\frac{1}{2^{\mu}}+\frac{\mu}{24}\right)[1+(-1)^{\mu}]\\ &+B\omega^{\mu+\nu+1}\left(\frac{1}{2^{\mu+\nu}}+\frac{\mu+\nu}{24}\right)[1+(-1)^{\mu+\nu}]+\text{etc.} \end{aligned} \]
Mais l'intégrale $\displaystyle\int y\,dx$, prise depuis $x=\alpha-\omega$ jusqu'à $x=\alpha+\omega$, a pour valeur exacte
\[ 2\beta\omega+A\omega^{\mu+1}\left[\frac{1+(-1)^{\mu}}{\mu+1}\right]+B\omega^{\mu+\nu+1}\left[\frac{1+(-1)^{\mu+\nu}}{\mu+\nu+1}\right]+\text{etc.} \]
Ainsi la correction qu'il faut appliquer à $\zeta$ ou à $Z$ pour avoir la vraie valeur de l'aire que l'on cherche, sera
\[ \begin{aligned} \delta Z={}&A\omega^{\mu+1}\left(\frac{1}{\mu+1}-\frac{1}{2^{\mu}}-\frac{\mu}{24}\right)[1+(-1)^{\mu}]\\ &+B\omega^{\mu+\nu+1}\left(\frac{1}{\mu+\nu+1}-\frac{1}{2^{\mu+\nu}}-\frac{(\mu+\nu)}{24}\right)[1+(-1)^{\mu+\nu}]. \end{aligned} \]
Il faut dans cette formule distinguer deux cas, selon que $\mu$ est de la forme $\frac{2k}{2i+1}$ ou $\frac{2k+1}{2i+1}$; supposant d'ailleurs $\nu=1$, ce qui ne peut souffrir que très-peu d'exceptions, on aura dans le premier cas,
\[ \delta Z=2A\omega^{\mu+1}\left(\frac{1}{\mu+1}-\frac{1}{2^{\mu}}-\frac{\mu}{24}\right), \]
et dans le second cas,
\[ \delta Z=2B\omega^{\mu+2}\left[\frac{1}{\mu+2}-\frac{1}{2^{\mu+1}}-\left(\frac{\mu+1}{24}\right)\right]. \]
Soit $F(\alpha+\frac{1}{2}\omega)-2F(\alpha)+F(\alpha-\frac{1}{2}\omega)=M$, la quantité $M$ sera connue par les termes qui composent $\Sigma\omega F(x+\frac{1}{2}\omega)$, et au moyen de cette quantité, on aura dans le premier cas,
\[ \delta Z=M\omega\left(\frac{2^{\mu}}{\mu+1}-1-\frac{\mu\cdot2^{\mu}}{24}\right), \]
et dans le second
\[ \delta Z=M\omega\left(\frac{2^{\mu+1}}{\mu+2}-1-\frac{(\mu+1)2^{\mu+1}}{24}\right). \]
Ainsi on connaîtra la correction à appliquer à la somme trouvée,

\origpage{316}
% The print numbers this page "163" (digits transposed); 316 by sequence.
pourvu qu'on connaisse l'exposant $\mu$ qui sert à caractériser le point singulier qui rend infini un ou plusieurs des coefficiens $\frac{dy}{dx}$, $\frac{d^{2}y}{dx^{2}}$, $\frac{d^{3}y}{dx^{3}}$. On voit en même temps que la correction sera beaucoup plus petite lorsque $\mu$ sera de la forme $\frac{2k+1}{2i+1}$ que lorsque $\mu$ sera de la forme $\frac{2k}{2i+1}$; ainsi le cas où $y$ devient un \emph{maximum} ou un \emph{minimum} lorsque $x=\alpha$, est celui qui exige la plus forte correction. L'autre cas n'apporte à la formule qu'une correction très-légère et souvent négligeable.

\subsection*{(6).}

Lorsqu'on aura reconnu que la portion de courbe qu'on veut quarrer, a un point singulier déterminé par l'équation
\[ y-\beta=A(x-\alpha)^{\mu}+B(x-a)^{\mu+\nu}+\text{etc.}; \]
\ednote{Printed $(x-a)$ in the second term; $(x-\alpha)$ is meant, as in the first term and in art. (5). An apparent misprint.}
on pourra, par une transformation fort simple, prévenir l'inconvénient qui naît de la valeur infinie des coefficiens $\frac{dy}{dx}$, $\frac{ddy}{dx^{2}}$, $\frac{d^{3}y}{dx^{3}}$. Il suffit pour cela de faire $x=\alpha+(z-f)^{m}$, $m$ étant le dénominateur impair de la fraction égale à $\mu$; alors $y$ deviendra une fonction connue de $z$, et si l'on prend $f=\sqrt[m]{\alpha}$, on aura $\displaystyle\int y\,dx=\displaystyle\int m(z-f)^{m-1}y\,dz$. Cette nouvelle intégrale devra être prise depuis $z=0$ jusqu'à $z=\sqrt[m]{\alpha}+\sqrt[m]{(a-\alpha)}$; et dans tout cet intervalle, la nouvelle courbe qu'il faut quarrer et qui a pour ordonnée la fonction $m(z-f)^{m-1}y$, n'aura aucun point singulier; de sorte que l'application de la formule (2) ne sera sujette à aucune difficulté.

\subsection*{(7).}

La plupart des méthodes qu'on donne ordinairement pour trouver par approximation l'aire d'une courbe, sont moins exactes ou moins commodes dans la pratique que celle que nous venons d'exposer. Je regarde surtout comme l'une des plus défectueuses, celle qui suppose que l'ordonnée de la courbe est représentée dans toute son étendue par la formule $y=a+bx+cx^{2}+ex^{3}+\text{etc.}$, ou par une formule équivalente; car de ce qu'une courbe passe par un grand nombre de points d'une courbe donnée, il ne s'ensuit pas que les deux courbes soient fort approchées l'une de l'autre; il peut arriver au contraire que les deux aires, malgré tous les points communs, soient aussi différentes entre elles qu'on voudra.

\origpage{317}
La seule méthode pratique qui soit à la fois simple et exacte pour les quadratures, consiste à diviser la courbe par des ordonnées équidistantes, de manière que les arcs compris de deux en deux ordonnées, soient assimilés à des arcs de parabole. Désignons par $c$, $c'$, $c''\ldots c^{(2n)}$ la suite des ordonnées qui répondent aux abscisses $0$, $\frac{1}{2}\omega$, $\omega$, $\frac{3}{2}\omega\ldots n\omega$; l'aire parabolique depuis $x=0$ jusqu'à $x=\omega$ sera exprimée par $\frac{\omega}{6}(c+4c'+c'')$; de même l'aire parabolique depuis $x=\omega$ jusqu'à $x=2\omega$, sera $\frac{\omega}{6}(c''+4c'''+c^{\text{iv}})$; ainsi de suite. Ajoutant toutes ces parties, on a pour l'aire entière $S$, comprise depuis $x=0$ jusqu'à $x=n\omega$, la formule
\[ S=\frac{\omega}{6}(c+4c'+2c''+4c'''+2c^{\text{iv}}\ldots+4c^{(2n-1)}+c^{(2n)}), \tag{3} \]
ou, ce qui revient au même, suivant les dénominations précédentes,
\[ S=\frac{\omega}{6}[F(0)+4F(\tfrac{1}{2}\omega)+2F(\omega)+4F(\tfrac{3}{2}\omega)\ldots+4F(x-\tfrac{1}{2}\omega)+F(x)]. \]
Cette formule peut s'écrire ainsi:
\[ \begin{aligned} S={}&\frac{\omega}{6}[F(0)-F(x)]\\ &+\frac{2\omega}{3}[F(\tfrac{1}{2}\omega)+F(\tfrac{3}{2}\omega)+F(\tfrac{5}{2}\omega)\ldots+F(x-\tfrac{1}{2}\omega)]\\ &+\frac{\omega}{3}[F(\omega)+F(2\omega)+F(3\omega)\ldots+F(x-\omega)+F(x)]. \end{aligned} \]
Mais d'après la formule (2), l'aire comprise depuis $x=0$ jusqu'à $x=n\omega$ étant appelée $Z$, on aura
\[ \begin{aligned} Z={}&\omega[F(\tfrac{1}{2}\omega)+F(\tfrac{3}{2}\omega)+F(\tfrac{5}{2}\omega)\ldots F(x-\tfrac{1}{2}\omega)]\\ &+A\omega^{2}P-B\omega^{4}Q+C\omega^{6}R-\text{etc.}, \end{aligned} \]
$P$, $Q$, $R$, etc. dépendant de la valeur des coefficiens différentiels $\frac{dF}{dx}$, $\frac{d^{3}F}{dx^{3}}$, $\frac{d^{5}F}{dx^{5}}$, etc. aux limites de l'intégrale.

Si on appelle $Z'$ l'aire comprise depuis $x=\frac{1}{2}\omega$ jusqu'à $x=(n+\frac{1}{2})\omega$, l'expression de cette aire sera semblablement
\[ \begin{aligned} Z'={}&\omega[F(\omega)+F(2\omega)+F(3\omega)\ldots+F(n\omega)]\\ &+A\omega^{2}P'-B\omega^{4}Q'+C\omega^{6}R'-\text{etc.}, \end{aligned} \]

\origpage{318}
en désignant par $P'$, $Q'$, $R'$, etc. ce que deviennent $P$, $Q$, $R$, etc. lorsqu'au lieu de $x$ on met $x+\frac{1}{2}\omega$. Mais si on désigne par $\psi$ l'accroissement de l'aire $Z$ lorsque $x$ devient $x+\frac{1}{2}\omega$, et par $\psi^{0}$ ce que devient $\psi$ lorsque $x=0$, on aura $Z'=Z+\psi-\psi^{0}$.

Or ayant $Z=\displaystyle\int dx\,F(x)$, on en déduit
\[ \psi=\frac{\omega}{2}F-\frac{\omega^{2}}{2\cdot4}\frac{dF}{dx}+\frac{\omega^{3}}{2\cdot3\cdot8}\cdot\frac{ddF}{dx^{2}}-\text{etc.}, \]
et par conséquent,
\[ \psi^{0}=\frac{\omega}{2}F^{0}-\frac{\omega^{2}}{2\cdot4}\cdot\frac{dF^{0}}{dx}+\frac{\omega^{3}}{2\cdot3\cdot8}\cdot\frac{ddF^{0}}{dx^{2}}-\text{etc.} \]
On tire de ces équations,
\[ \begin{aligned} \omega[F(\tfrac{1}{2}\omega)+F(\tfrac{3}{2}\omega)\ldots+F(x-\tfrac{1}{2}\omega)]&=Z-A\omega^{2}P+B\omega^{4}Q-\text{etc.}\\ \omega[F(\omega)+F(2\omega)\ldots+F(x)]&=Z+\psi-\psi^{0}-A\omega^{2}P'+B\omega^{4}Q'-\text{etc.} \end{aligned} \]
Substituant ces valeurs dans celle de $S$, il viendra
\[ S=\frac{\omega}{6}(F^{0}-F)+Z+\tfrac{1}{3}(\psi-\psi^{0})-A\omega^{2}\left(\frac{2P+P'}{3}\right)+B\omega^{4}\left(\frac{2Q+Q'}{3}\right)-\text{etc.}, \]
ou, en substituant la valeur de $\psi-\psi^{0}$,
\[ \begin{aligned} S=Z&-\frac{\omega^{2}}{2\cdot4}\cdot\frac{1}{3}\left(\frac{dF}{dx}-\frac{dF^{0}}{dx}\right)+\frac{\omega^{2}}{2\cdot3\cdot8}\cdot\frac{1}{3}\left(\frac{ddF}{dx^{2}}-\frac{ddF^{0}}{dx^{2}}\right)\\ &-\frac{\omega^{4}}{2\cdot3\cdot4\cdot16}\cdot\frac{1}{3}\left(\frac{d^{3}F}{dx^{3}}-\frac{d^{3}F^{0}}{dx^{3}}\right)+\text{etc.}\\ &-A\omega^{2}\left(\frac{2P+P'}{3}\right)+B\omega^{4}\left(\frac{2Q+Q'}{3}\right)-\text{etc.} \end{aligned} \]
\ednote{Printed $\omega^{2}$ over $2\cdot3\cdot8$ in the second term of the first line; $\omega^{3}$ is meant, as in the value of $\psi$ above. An apparent misprint.}
Mais on a $P=\frac{dF}{dx}-\frac{dF^{0}}{dx}$; si dans cette quantité on met $x+\frac{1}{2}\omega$ à la place de $x$, $\frac{dF}{dx}$ deviendra $\frac{dF}{dx}+\frac{\omega}{2}\cdot\frac{ddF}{dx^{2}}+\frac{1}{2}\cdot\frac{\omega^{2}}{4}\cdot\frac{d^{3}F}{dx^{3}}+\text{etc.}$; donc
\[ P'=\frac{dF}{dx}-\frac{dF^{0}}{dx}+\frac{\omega}{2}\left(\frac{ddF}{dx^{2}}-\frac{ddF^{0}}{dx^{2}}\right)+\frac{1}{2}\cdot\frac{\omega^{2}}{4}\left(\frac{d^{3}F}{dx^{3}}-\frac{d^{3}F^{0}}{dx^{3}}\right)+\text{etc.}; \]
de même ayant $Q=\frac{d^{3}F}{dx^{3}}-\frac{d^{3}F^{0}}{dx^{3}}$, on en déduira
\[ Q'=\frac{d^{3}F}{dx^{3}}-\frac{d^{3}F^{0}}{dx^{3}}+\frac{1}{2}\omega\left(\frac{d^{4}F}{dx^{4}}-\frac{d^{4}F^{0}}{dx^{4}}\right)+\text{etc.} \]

\origpage{319}
Cela posé, en faisant les substitutions et rejetant les puissances de $\omega$ au-delà de la quatrième, on aura
\[ S=Z+\frac{\omega^{4}}{2880}\left(\frac{d^{3}F}{dx^{3}}-\frac{d^{3}F^{0}}{dx^{3}}\right); \]
d'où l'on voit combien est petite la différence entre la valeur exacte de l'aire $Z$, et sa valeur $S$ approchée par une somme d'aires paraboliques.

Ainsi dans la pratique on pourra, si l'on veut, s'en tenir à la formule (3) qui donne la valeur de $S$, et y ajouter la petite correction que fournit la formule précédente pour avoir une valeur plus exacte de l'aire $Z$, ce qui donnera
\[ Z=S-\frac{\omega^{4}}{2880}\left(\frac{d^{3}F}{dx^{3}}-\frac{d^{3}F^{0}}{dx^{3}}\right). \]
Au reste la formule (3) suppose le calcul de deux fois plus de termes que la formule (2); ainsi à ce titre seul, l'usage de la formule (2) est préférable.

\subsection*{(8).}

Nous avons suffisamment expliqué dans ce qui précède, comment on peut trouver une intégrale proposée $\displaystyle\int y\,dx$ entre deux limites données $x=0$, $x=a$; et nous avons donné les moyens de pousser l'approximation jusqu'à tel degré qu'on voudra fixer. Mais si l'intégrale devait s'étendre depuis $x=0$ jusqu'à $x=\infty$, on pourrait craindre que la méthode ne fût impraticable, puisque l'intégrale finie $\Sigma\omega F(x+\frac{1}{2}\omega)$ contiendrait alors une infinité de termes. La réponse à cette objection est que si l'intégrale cherchée doit être une quantité finie, seul cas où il y a lieu d'en faire le calcul, l'ordonnée $F(x)$ doit décroître avec beaucoup de rapidité lorsque $x$ devient un peu grand, de sorte que les termes qui composent $\Sigma\omega F(x+\frac{1}{2}\omega)$ ne tarderont pas à devenir très-petits et entièrement négligeables.

Il est facile au reste de remédier à cet inconvénient par une transformation. En effet si on fait, par exemple, $x=\frac{z}{1-z}$, l'intégrale $\displaystyle\int y\,dx$ sera transformée en une autre $\displaystyle\int P\,dz$, dans laquelle $P$ sera une fonction connue de $z$ et qui devra être intégrée depuis $z=0$ jusqu'à $z=1$.

On pourrait encore faire $x=k\operatorname{tang}\varphi$, $k$ étant à volonté; et la

\origpage{320}
transformée en $\varphi$, qui serait de la forme $\displaystyle\int Q\,d\varphi$, devrait être intégrée depuis $\varphi=0$ jusqu'à $\varphi=\frac{1}{2}\pi$.

En général c'est par la nature de la question qu'on jugera de la substitution qu'il convient d'employer, préférablement à toute autre, pour transformer l'intégrale proposée en une autre qui soit comprise entre des limites finies.

\section*{Construction de la courbe dans laquelle l'arc $s$ est donné en fonction de la quantité $\frac{dy}{dx}$:}

\subsection*{(9).}

Soit $s$ l'arc d'une courbe, $\theta$ l'angle que la tangente à l'extrémité de l'arc fait avec la ligne des abscisses, ensorte qu'on ait $\operatorname{tang}\theta=\frac{dy}{dx}$; si l'équation de la courbe n'est pas donnée, mais qu'on ait simplement l'expression de l'arc $s$ en fonction de l'angle $\theta$, savoir $s=F(\theta)$, et que de cette expression on ne puisse déduire les valeurs des coordonnées $x$ et $y$, parce qu'elles dépendent d'intégrales trop compliquées, il s'agit de trouver au moins par approximation les valeurs de $x$ et $y$ qui correspondent à une valeur donnée de l'angle $\theta$.

Puisque $s$ est une fonction connue de $\theta$, on pourra supposer $ds=Q\,d\theta$, et alors les valeurs de $x$ et $y$ dépendent immédiatement des quadratures, puisqu'on a $x=\displaystyle\int Q\,d\theta\cos\theta$, $y=\displaystyle\int Q\,d\theta\sin\theta$. On pourrait donc faire usage de la méthode donnée dans le chapitre précédent pour calculer les valeurs de $x$ et de $y$ qui correspondent à une valeur donnée de $\theta$. Mais la courbe se construira plus facilement par une méthode particulière que nous allons exposer.

Supposons que depuis un point donné où $\theta=\alpha$, jusqu'à un autre point quelconque où $\theta$ a une valeur donnée, on veuille connaître les valeurs de $x$ et de $y$; on divisera l'intervalle $\theta-\alpha$ en un certain nombre de parties égales, d'autant plus petites qu'on voudra pousser plus loin l'approximation. Soit $\omega$ une de ces parties, et $n$ leur nombre, ensorte qu'on ait $\theta=\alpha+n\omega$.

Au moyen de l'équation donnée entre $s$ et $\theta$, on calculera successivement les valeurs de $s$ qui répondent aux angles $\alpha$, $\alpha+\omega$, $\alpha+2\omega\ldots\alpha+n\omega$, et on prendra leurs différences consécutives.

\origpage{321}
Soit $\Delta s$ l'une de ces différences supposée correspondante à l'angle $\theta$, ensorte qu'on ait $s=F(\theta)$, $s+\Delta s=F(\theta+\omega)$, on aura pour premières valeurs approchées de $x$ et $y$, les formules
\[ \begin{aligned} x&=\Sigma\Delta s\cos(\theta+\tfrac{1}{2}\omega)\\ y&=\Sigma\Delta s\sin(\theta+\tfrac{1}{2}\omega), \end{aligned} \]
où l'on voit que chaque $\Delta s$ est multiplié par le cosinus ou le sinus de l'angle $\theta+\frac{1}{2}\omega$, moyen entre l'angle $\theta$ qui répond à l'arc $s$, et l'angle $\theta+\omega$ qui répond à l'arc $s+\Delta s$.

\subsection*{(10).}

Il faut voir maintenant quelles sont les corrections qu'il faut appliquer à ces premières valeurs approchées. Pour cela cherchons en général la valeur de $\xi$ d'après l'équation
\[ x=\Sigma\Delta s\cos(\theta+\tfrac{1}{2}\omega)+\xi, \]
on aura
\[ \Delta\xi=\Delta x-\Delta s\cos(\theta+\tfrac{1}{2}\omega); \]
Or en regardant $x$ et $s$ comme fonctions de $\theta$, on a
\[ \begin{aligned} \Delta s&=\omega\frac{ds}{d\theta}+\frac{\omega^{2}}{2}\cdot\frac{dds}{d\theta^{2}}+\frac{\omega^{3}}{2\cdot3}\cdot\frac{d^{3}s}{d\theta^{3}}+\text{etc.}\\ \Delta x&=\omega\frac{dx}{d\theta}+\frac{\omega^{2}}{2}\cdot\frac{ddx}{d\theta^{2}}+\frac{\omega^{3}}{2\cdot3}\cdot\frac{d^{3}x}{d\theta^{3}}+\text{etc.} \end{aligned} \]
On a en même temps
\[ \cos(\theta+\tfrac{1}{2}\omega)=\cos\theta-\frac{\omega}{2}\sin\theta-\frac{1}{1\cdot2}\cdot\frac{\omega^{2}}{4}\cos\theta+\frac{1}{1\cdot2\cdot3}\cdot\frac{\omega^{3}}{8}\sin\theta+\text{etc.} \]
Mais de ce que $dx=ds\cos\theta$, on déduit
\[ \begin{aligned} \frac{dx}{d\theta}&=\frac{ds}{d\theta}\cos\theta\\ \frac{ddx}{d\theta^{2}}&=\frac{dds}{d\theta^{2}}\cos\theta-\frac{ds}{d\theta}\sin\theta\\ \frac{d^{3}x}{d\theta^{3}}&=\frac{d^{3}s}{d\theta^{3}}\cos\theta-2\frac{dds}{d\theta^{2}}\sin\theta-\frac{ds}{d\theta}\cos\theta\\ \frac{d^{4}x}{d\theta^{4}}&=\frac{d^{4}s}{d\theta^{4}}\cos\theta-3\frac{d^{3}s}{d\theta^{3}}\sin\theta-3\frac{d^{2}s}{d\theta^{2}}\cos\theta+\frac{ds}{d\theta}\sin\theta.\\ &\text{etc.} \end{aligned} \]

\origpage{322}
Ainsi on a d'une part,
\[ \begin{aligned} \Delta x={}&\frac{ds}{d\theta}\left(\omega\cos\theta-\frac{\omega^{2}}{2}\sin\theta-\frac{\omega^{3}}{2\cdot3}\cos\theta+\frac{\omega^{4}}{2\cdot3\cdot4}\sin\theta+\text{etc.}\right)\\ &+\frac{dds}{d\theta^{2}}\left(\frac{\omega^{2}}{2}\cos\theta-2\frac{\omega^{3}}{2\cdot3}\sin\theta-3\frac{\omega^{4}}{2\cdot3\cdot4}\cos\theta+\text{etc.}\right)\\ &+\frac{d^{3}s}{d\theta^{3}}\left(\frac{\omega^{3}}{2\cdot3}\cos\theta-3\frac{\omega^{4}}{2\cdot3\cdot4}\sin\theta-\text{etc.}\right)\\ &+\frac{d^{4}s}{d\theta^{4}}\left(\frac{\omega^{4}}{2\cdot3\cdot4}\cos\theta-\text{etc.}\right)+\text{etc.}; \end{aligned} \]
et d'autre part,
\[ \begin{aligned} \Delta s\cos(\theta+\tfrac{1}{2}\omega)={}&\frac{ds}{d\theta}\left(\omega\cos\theta-\frac{\omega^{2}}{2}\sin\theta-\frac{\omega^{3}}{2\cdot4}\cos\theta+\frac{\omega^{4}}{2\cdot4\cdot6}\sin\theta+\text{etc.}\right)\\ &+\frac{dds}{d\theta^{2}}\left(\frac{\omega^{2}}{2}\cos\theta-\frac{\omega^{3}}{4}\sin\theta-\frac{\omega^{4}}{2\cdot4}\cdot\frac{1}{2}\cos\theta+\text{etc.}\right)\\ &+\frac{d^{3}s}{d\theta^{3}}\left(\frac{\omega^{3}}{2\cdot3}\cos\theta-\frac{\omega^{4}}{2\cdot3}\cdot\frac{1}{2}\sin\theta-\text{etc.}\right)\\ &+\frac{d^{4}s}{d\theta^{4}}\left(\frac{\omega^{4}}{2\cdot3\cdot4}\cos\theta-\text{etc.}\right)+\text{etc.} \end{aligned} \]
La différence de ces deux quantités donnera, en s'arrêtant aux $\omega^{4}$,
\[ \begin{aligned} \Delta\xi={}&-\frac{ds}{d\theta}\left(\frac{\omega^{3}}{2\cdot3\cdot4}\cos\theta-\frac{\omega^{4}}{2\cdot4\cdot6}\sin\theta\right)\\ &-\frac{dds}{d\theta^{2}}\left(\frac{\omega^{3}}{3\cdot4}\sin\theta+\frac{\omega^{4}}{16}\cos\theta\right)\\ &-\frac{d^{3}s}{d\theta^{3}}\left(\frac{\omega^{4}}{2\cdot3\cdot4}\sin\theta\right). \end{aligned} \]
Soit $\xi=P\omega^{2}+Q\omega^{3}+R\omega^{4}$, on aura
\[ \Delta\xi=\omega^{3}\frac{dP}{d\theta}+\frac{\omega^{4}}{2}\cdot\frac{ddP}{d\theta^{2}}+\omega^{4}\frac{dQ}{d\theta}. \]
Ainsi les coefficiens $P$ et $Q$ devront être déterminés par les équations
\[ \begin{aligned} \frac{dP}{d\theta}&=-\frac{ds}{d\theta}\cdot\frac{\cos\theta}{2\cdot3\cdot4}-\frac{dds}{d\theta^{2}}\cdot\frac{\sin\theta}{3\cdot4}\\ \frac{1}{2}\frac{ddP}{d\theta^{2}}+\frac{dQ}{d\theta}&=\frac{ds}{d\theta}\cdot\frac{\sin\theta}{48}-\frac{dds}{d\theta^{2}}\cdot\frac{\cos\theta}{16}-\frac{d^{3}s}{d\theta^{3}}\cdot\frac{\sin\theta}{24}. \end{aligned} \]
La première donne en intégrant,
\[ P=-\frac{ds}{d\theta}\cdot\frac{\sin\theta}{2\cdot3\cdot4}+\frac{1}{2\cdot3\cdot4}\int ds\cos\theta, \]
\ednote{The first term is printed with the denominator $2\cdot3\cdot4$ here and in the next formula; integrating the preceding equation gives $3\cdot4$, which the corrected value of $x$ below (coefficient $\frac{\omega^{2}}{12}$) presupposes. An apparent misprint.}

\origpage{323}
ou
\[ P=-\frac{ds}{d\theta}\cdot\frac{\sin\theta}{2\cdot3\cdot4}+\frac{x}{2\cdot3\cdot4}+\text{constante}. \]
L'autre donne $\frac{dQ}{d\theta}=0$; on a donc pour première valeur corrigée
\[ x=\Sigma\Delta s\cos(\theta+\tfrac{1}{2}\omega)-\frac{\omega^{2}}{12}\left(\frac{ds}{d\theta}\sin\theta-\tfrac{1}{2}x+\text{const.}\right); \]
et puisqu'on a trouvé $Q=0$, il est clair que ce résultat est exact, aux quantités près de l'ordre $\omega^{4}$. En poussant plus loin l'approximation, on trouverait
\[ \begin{aligned} x={}&\Sigma\Delta s\cos(\theta+\tfrac{1}{2}\omega)-\frac{\omega^{2}}{2}\left(\frac{ds}{d\theta}\sin\theta-\tfrac{1}{2}x+\text{const.}\right)\\ &+\frac{\omega^{4}}{720}\left(\frac{d^{3}s}{d\theta^{3}}\sin\theta+3\frac{d^{2}s}{d\theta^{2}}\cos\theta-\frac{1}{2}\cdot\frac{ds}{d\theta}\sin\theta-\frac{3}{8}x+\text{const.}\right); \end{aligned} \]
\ednote{The first parenthesis is printed with the coefficient $\frac{\omega^{2}}{2}$, where the preceding formula has $\frac{\omega^{2}}{12}$; an apparent misprint.}
et ce nouveau résultat est exact aux quantités près de l'ordre $\omega^{6}$; car la série du second membre ne contient aucune puissance impaire de $\omega$, ainsi qu'il sera démontré ci-après.

En examinant de plus près cette valeur de $x$, on voit qu'elle peut être mise sous la forme
\[ \begin{aligned} \frac{x\sin\frac{1}{2}\omega}{\frac{1}{2}\omega}={}&\Sigma\Delta s\cos(\theta+\tfrac{1}{2}\omega)-\frac{\omega^{2}}{12}\left(\frac{ds}{d\theta}\sin\theta+\text{const.}\right)\\ &+\frac{\omega^{4}}{720}\left(\frac{d^{3}s}{d\theta^{3}}\sin\theta+3\frac{dds}{d\theta^{2}}\cos\theta-\frac{1}{2}\frac{ds}{d\theta}\sin\theta+\text{const.}\right); \end{aligned} \]
d'où l'on tire
\[ \begin{aligned} x={}&\frac{\frac{1}{2}\omega}{\sin\frac{1}{2}\omega}\Sigma\Delta s\cos(\theta+\tfrac{1}{2}\omega)-\frac{\omega^{2}}{12}\left(\frac{ds}{d\theta}\sin\theta+\text{const.}\right)\\ &+\frac{\omega^{4}}{720}\left(\frac{d^{3}s}{d\theta^{3}}\sin\theta+3\frac{dds}{d\theta^{2}}\cos\theta-3\frac{ds}{d\theta}\sin\theta+\text{const.}\right)\\ &-\text{etc.} \end{aligned} \]
Les constantes renfermées dans les parenthèses sont telles qu'elles doivent faire disparaître les quantités jointes, lorsqu'on suppose dans celles-ci $\theta=\alpha$.

La forme de cette expression est assez évidente, mais il importe de déterminer la loi de ses coefficiens, afin de pouvoir la prolonger indéfiniment. Pour cet effet, nous prendrons une forme particulière

\origpage{324}
de $s$, telle qu'elle puisse indiquer, aussi simplement qu'il est possible, les différens termes de ce développement.

\subsection*{(11).}

Soit donc $s=e^{m\theta}-e^{m\alpha}$, afin qu'on ait $s=0$ lorsque $\theta=\alpha$; les coordonnées $x$ et $y$ se détermineront par les équations différentielles $dx=ds\cos\theta=md\theta e^{m\theta}\cos\theta$, $dy=ds\sin\theta=md\theta e^{m\theta}\sin\theta$, dont les intégrales sont
\[ \begin{aligned} (m^{2}+1)x&=m^{2}e^{m\theta}\cos\theta+me^{m\theta}\sin\theta+\text{const.}\\ (m^{2}+1)y&=m^{2}e^{m\theta}\sin\theta-me^{m\theta}\cos\theta+\text{const.} \end{aligned} \]
Dans ce même cas, on aura $\Delta s=e^{m\theta}(e^{m\omega}-1)$, et
\[ \begin{aligned} \Sigma\Delta s\cos(\theta+\tfrac{1}{2}\omega)&=(e^{m\omega}-1)\Sigma e^{m\theta}\cos(\theta+\tfrac{1}{2}\omega)\\ &=\frac{(e^{m\omega}-1)^{2}\cos\frac{1}{2}\omega}{Re^{m\omega}}e^{m\theta}\cos\theta+\frac{(e^{2m\omega}-1)\sin\frac{1}{2}\omega}{Re^{m\omega}}e^{m\theta}\sin\theta+\text{const.}, \end{aligned} \]
formule où l'on a fait pour abréger, $R=e^{m\omega}+e^{-m\omega}-2\cos\omega$. Cette valeur de $\Sigma$ peut se mettre sous la forme
\[ \cos\tfrac{1}{2}\omega\cdot e^{m\theta}\cos\theta-\frac{2\sin\omega\sin\frac{1}{2}\omega}{R}e^{m\theta}\cos\theta+\frac{e^{m\omega}-e^{-m\omega}}{R}\sin\tfrac{1}{2}\omega e^{m\theta}\sin\theta+\text{const.} \]
Soit donc
\[ x=\frac{\frac{1}{2}\omega}{\sin\frac{1}{2}\omega}\Sigma\Delta s\cos(\theta+\tfrac{1}{2}\omega)-\xi', \]
et on aura en faisant les substitutions,
\[ \begin{aligned} \xi'={}&\left(\tfrac{1}{2}\omega\cot\tfrac{1}{2}\omega-\frac{m^{2}}{m^{2}+1}-\frac{\omega\sin\omega}{R}\right)e^{m\theta}\cos\theta\\ &+\left(\frac{\frac{1}{2}\omega(e^{m\omega}-e^{-m\omega})}{R}-\frac{m}{m^{2}+1}\right)e^{m\theta}\sin\theta+\text{const.} \end{aligned} \]
On remarquera d'abord que cette valeur de $\xi'$ ne change pas de signe lorsque $\omega$ en change, et qu'ainsi le développement de $\xi'$ ne contient que des puissances paires de $\omega$, comme nous l'avons déjà annoncé.

\subsection*{(12).}

Il faut maintenant développer la valeur de $\xi'$ suivant les puissances de $\omega$; et pour cela j'observe qu'on a
\[ \frac{1}{m^{2}+1}-\frac{\omega\sin\omega}{R}=\frac{\frac{1}{3\cdot4}\omega^{2}+\frac{m^{2}-2}{3\cdot4\cdot5\cdot6}\omega^{4}+\frac{m^{4}-2m^{2}+3}{3\cdot4\ldots8}\omega^{6}+\frac{m^{6}-2m^{4}+3m^{2}-4}{3\cdot4\ldots10}\omega^{8}+\text{etc.}}{1+\frac{m^{2}-1}{3\cdot4}\omega^{2}+\frac{m^{4}-m^{2}+1}{3\cdot4\cdot5\cdot6}\omega^{4}+\frac{m^{6}-m^{4}+m^{2}-1}{3\cdot4\ldots8}\omega^{6}+\text{etc.}} \]

\origpage{325}
Cette quantité étant fonction de $m$, je la désigne par $\Psi(m)$, et je remarque qu'en faisant $m=0$, on aura $1-\frac{1}{2}\omega\cot\frac{1}{2}\omega=\Psi(0)$. Supposons qu'en effectuant le développement on ait
\[ \begin{aligned} \Psi(m)&=A^{m}\omega^{2}+B^{m}\omega^{4}+C^{m}\omega^{6}+\text{etc.}\\ \Psi(0)&=A^{0}\omega^{2}+B^{0}\omega^{4}+C^{0}\omega^{6}+\text{etc.} \end{aligned} \]
Comme le premier terme $A^{m}=\frac{1}{3\cdot4}=A^{0}$, on aura
\[ \tfrac{1}{2}\omega\cot\tfrac{1}{2}\omega-1+\frac{1}{m^{2}+1}-\frac{\omega\sin\omega}{R}=(B^{m}-B^{0})\omega^{4}+(C^{m}-C^{0})\omega^{6}+\text{etc.}; \]
c'est le coefficient de $e^{m\theta}\cos\theta$ dans la valeur de $\xi'$: on aura semblablement
\[ \frac{\frac{1}{2}\omega(e^{m\omega}-e^{-m\omega})}{mR}-\frac{1}{m^{2}+1}=\frac{\frac{\omega^{2}}{3\cdot4}+\frac{2m^{2}-1}{3\cdot4\cdot5\cdot6}\omega^{4}+\frac{3m^{4}-2m^{2}+1}{3\cdot4\cdot5\ldots8}\omega^{6}+\text{etc.}}{1+\frac{m^{2}-1}{3\cdot4}\omega^{2}+\frac{m^{4}-m^{2}+1}{3\cdot4\cdot5\cdot6}\omega^{4}+\text{etc.}}. \]
Appelons cette quantité $\Phi(m)$, et supposons que son développement donne
\[ \Phi(m)=\alpha^{m}\omega^{2}+\beta^{m}\omega^{4}+\gamma^{m}\omega^{6}+\delta^{m}\omega^{8}+\text{etc.}, \]
on aura enfin
\[ \begin{aligned} \xi'={}&[(B^{m}-B^{0})\omega^{4}+(C^{m}-C^{0})\omega^{6}+(D^{m}-D^{0})\omega^{8}+\text{etc.}]e^{m\theta}\cos\theta\\ &+(\alpha^{m}\omega^{2}+\beta^{m}\omega^{4}+\gamma^{m}\omega^{6}+\text{etc.})me^{m\theta}\sin\theta+\text{const.} \end{aligned} \]
Ces suites étant développées jusqu'à telle puissance de $\omega$ qu'on voudra, et les coefficiens $B^{m}$, $C^{m}$, $D^{m}$, $\alpha^{m}$, $\beta^{m}$, $\gamma^{m}$, etc. étant exprimés en fonctions de $m$, on fera les substitutions $me^{m\theta}=\frac{ds}{d\theta}$, $m^{2}e^{m\theta}=\frac{dds}{d\theta^{2}}$, $m^{3}e^{m\theta}=\frac{d^{3}s}{d\theta^{3}}$, etc., et on aura la valeur de $\xi'$, quelle que soit la fonction de $\theta$ égale à $s$.

\subsection*{(13).}

Si par exemple on veut développer la valeur de $\xi'$ jusqu'aux $\omega^{6}$ inclusivement, on fera d'abord le développement des fonctions $\Psi(m)$ et $\Phi(m)$ en s'arrêtant aux $\omega^{6}$; ce qui donnera
\[ \begin{aligned} \Psi(m)&=\frac{1}{12}\omega^{2}-\frac{(3m^{2}-1)}{720}\omega^{4}+\frac{5m^{4}-10m^{2}+1}{1440\cdot21}\omega^{6}\\ \Phi(m)&=\frac{1}{12}\omega^{2}-\frac{(m^{2}-3)}{720}\omega^{4}+\frac{m^{4}-10m^{2}+5}{1440\cdot21}\omega^{6}; \end{aligned} \]

\origpage{326}
Ces deux valeurs se déduiraient l'une de l'autre en mettant $\frac{1}{m}$ au lieu de $m$, supprimant les dénominateurs et changeant les signes des termes pris alternativement; c'est aussi ce qu'on déduirait de la forme des fractions dont le développement donne les fonctions $\Psi(m)$ et $\Phi(m)$.

Substituant ces valeurs dans la formule
\[ \xi'=[\Psi(m)-\Psi(0)]e^{m\theta}\cos\theta+m\Phi(m)e^{m\theta}\sin\theta+\text{const.}, \]
on aura dans le cas supposé,
\[ \begin{aligned} \xi'={}&-\frac{3\omega^{4}}{720}m^{2}e^{m\theta}\cos\theta+\frac{5m^{4}-10m^{2}}{30240}\omega^{6}e^{m\theta}\cos\theta\\ &+\frac{\omega^{2}}{12}me^{m\theta}\sin\theta-\frac{\omega^{4}}{720}(m^{3}-3m)e^{m\theta}\sin\theta\\ &+\frac{\omega^{6}}{30240}(m^{5}-10m^{3}+5m)e^{m\theta}\sin\theta+\text{const.}; \end{aligned} \]
et par conséquent la valeur générale de $\xi'$ est
\[ \begin{aligned} \xi'={}&\frac{\omega^{2}}{12}\cdot\frac{ds}{d\theta}\sin\theta-\frac{\omega^{4}}{720}\left(\frac{d^{3}s}{d\theta^{3}}\sin\theta+3\frac{dds}{d\theta^{2}}\cos\theta-3\frac{ds}{d\theta}\sin\theta\right)\\ &+\frac{\omega^{6}}{30240}\left(\frac{d^{5}s}{d\theta^{5}}\sin\theta+5\frac{d^{4}s}{d\theta^{4}}\cos\theta-10\frac{d^{3}s}{d\theta^{3}}\sin\theta-10\frac{d^{2}s}{d\theta^{2}}\cos\theta+5\frac{ds}{d\theta}\sin\theta\right)\\ &+\text{const.} \end{aligned} \]
Cette formule est telle, qu'on voit au premier coup d'œil la loi que suivent les facteurs différentiels; quant aux coefficiens constans $\frac{1}{12}$, $\frac{1}{720}$, $\frac{1}{30240}$, etc., ils ne sont autre chose que ceux qu'on déduit du développement de la quantité $1-\frac{1}{2}\omega\cot\frac{1}{2}\omega$; de sorte que si on fait
\[ 1-\tfrac{1}{2}\omega\cot\tfrac{1}{2}\omega=A^{0}\omega^{2}+B^{0}\omega^{4}+C^{0}\omega^{6}+D^{0}\omega^{8}+\text{etc.}, \]
on aura généralement
\[ \begin{aligned} \xi'={}&\text{const}+A^{0}\omega^{2}\left(\frac{d(s\sin\theta)}{d\theta}-s\cos\theta\right)\\ &-B^{0}\omega^{4}\left(\frac{d^{3}(s\sin\theta)}{d\theta^{3}}+s\cos\theta\right)\\ &+C^{0}\omega^{6}\left(\frac{d^{5}(s\sin\theta)}{d\theta^{5}}-s\cos\theta\right)\\ &-\text{etc.} \end{aligned} \]

\origpage{327}
La constante est telle, que la valeur de $\xi'$ devra s'évanouir au commencement de l'intégrale où $s=0$ et $\theta=\alpha$; il faudra donc de chaque coefficient différentiel réduit en quantités finies, retrancher ce que devient ce terme lorsque $s=0$ et $\theta=\alpha$.

\subsection*{(14).}

L'état de simplicité où nous avons réduit la valeur de $\xi'$, fait présumer qu'il est possible de parvenir à cette formule par une voie plus directe et moins laborieuse. Mais sans nous arrêter à cette recherche, nous nous contenterons de vérifier la formule trouvée, dans toute son étendue, au moyen d'une valeur de $s$ qui permettra de trouver généralement, et d'une manière fort simple, la différence d'un ordre quelconque de $s\sin\theta$.

Nous choisirons pour cet objet la valeur $s=\sin a\theta$ qui donne $s\sin\theta=\frac{1}{2}\cos(a-1)\theta-\frac{1}{2}\cos(a+1)\theta$. On aura donc par des différentiations successives,
\[ \begin{aligned} \frac{d(s\sin\theta)}{d\theta}&=-\tfrac{1}{2}(a-1)\sin(a-1)\theta+\tfrac{1}{2}(a+1)\sin(a+1)\theta\\ \frac{d^{3}(s\sin\theta)}{d\theta^{3}}&=\tfrac{1}{2}(a-1)^{3}\sin(a-1)\theta-\tfrac{1}{2}(a+1)^{3}\sin(a+1)\theta\\ \frac{d^{5}(s\sin\theta)}{d\theta^{5}}&=-\tfrac{1}{2}(a-1)^{5}\sin(a-1)\theta+\tfrac{1}{2}(a+1)^{5}\sin(a+1)\theta\\ &\text{etc.} \end{aligned} \]
Substituant ces valeurs dans celle de $\xi'$, on aura
\[ \begin{aligned} \xi'={}&-(A^{0}\omega^{2}+B^{0}\omega^{4}(a-1)^{2}+C^{0}\omega^{6}(a-1)^{4}+\text{etc.})\frac{a-1}{2}\sin(a-1)\theta\\ &+(A^{0}\omega^{2}+B^{0}\omega^{4}(a+1)^{2}+C^{0}\omega^{6}(a+1)^{4}+\text{etc.})\frac{a+1}{2}\sin(a+1)\theta\\ &-(A^{0}\omega^{2}+B^{0}\omega^{4}+C^{0}\omega^{6}+\text{etc.})s\cos\theta. \end{aligned} \]
Cela posé, puisqu'on a en général $A^{0}z^{2}+B^{0}z^{4}+C^{0}z^{6}+\text{etc.}=1-\frac{1}{2}z\cot\frac{1}{2}z$, les suites comprises dans la valeur de $\xi'$ sont faciles à sommer, et il en résulte
\[ \begin{aligned} \xi'={}&\left[1-\tfrac{1}{2}(a+1)\omega\cot(a+1)\frac{\omega}{2}\right]\frac{\sin(a+1)\theta}{2(a+1)}\\ &-\left[1-\tfrac{1}{2}(a-1)\omega\cot(a-1)\frac{\omega}{2}\right]\frac{\sin(a-1)\theta}{2(a-1)}\\ &-(1-\tfrac{1}{2}\omega\cot\tfrac{1}{2}\omega)[\tfrac{1}{2}\sin(a+1)\theta+\tfrac{1}{2}\sin(a-1)\theta]. \end{aligned} \]

\origpage{328}
Je n'ajoute pas de constante, parce que je suppose $\theta=0$ au commencement de l'intégrale.

Mais de la valeur $s=\sin a\theta$ on déduit $ds=ad\theta\cos a\theta$, $dx=ad\theta\cos a\theta\cos\theta=\frac{1}{2}ad\theta\cos(a+1)\theta+\frac{1}{2}ad\theta\cos(a-1)\theta$, et par conséquent
\[ x=\frac{\frac{1}{2}a}{a+1}\sin(a+1)\theta+\frac{\frac{1}{2}a}{a-1}\sin(a-1)\theta. \]
Or nous avons fait
\[ x=\frac{\frac{1}{2}\omega}{\sin\frac{1}{2}\omega}\Sigma\Delta s\cos(\theta+\tfrac{1}{2}\omega)-\xi'; \]
ainsi il ne reste plus à vérifier que l'équation
\[ \Sigma\Delta s\cos(\theta+\tfrac{1}{2}\omega)=\frac{\frac{1}{2}\sin\frac{a\omega}{2}}{\sin(a+1)\frac{\omega}{2}}\sin(a+1)\theta+\frac{\frac{1}{2}\sin\frac{a\omega}{2}}{\sin(a-1)\frac{\omega}{2}}\sin(a-1)\theta. \]
Prenant les différences de part et d'autre en supposant que $\theta$ devienne $\theta+\omega$, on trouve que l'équation est entièrement identique.

\subsection*{(15).}

La formule générale trouvée pour la valeur de $\xi'$ est établie par là d'une manière certaine; car s'il y avait dans la suite générale un seul terme qui ne fût pas conforme à la loi observée, ce même terme se retrouverait dans l'application au cas de $s=\sin a\theta$, puisqu'il n'y a pas deux termes qui soient affectés à la fois d'une même puissance de $\omega$ et d'une même puissance de $a$, et que d'ailleurs il n'existe aucun terme dans la formule générale qui n'ait son correspondant dans la formule propre au cas particulier.

On aura donc généralement
\[ x=\frac{\frac{1}{2}\omega}{\sin\frac{1}{2}\omega}\Sigma\Delta s\cos(\theta+\tfrac{1}{2}\omega)-X(\theta)+X(\alpha), \]
$X$ ou $X(\theta)$ étant une fonction de $\theta$ représentée par la formule
\[ \begin{aligned} X(\theta)={}&A^{0}\omega^{2}\left(\frac{d(s\sin\theta)}{d\theta}-s\cos\theta\right)\\ &-B^{0}\omega^{4}\left(\frac{d^{3}(s\sin\theta)}{d\theta^{3}}+s\cos\theta\right)\\ &+C^{0}\omega^{6}\left(\frac{d^{5}(s\sin\theta)}{d\theta^{5}}-s\cos\theta\right)\\ &-\text{etc.}, \end{aligned} \]
et $X(\alpha)$ représentant une fonction semblable de $\alpha$.

\origpage{329}

Si l'on change $\theta$ en $90^{\circ}-\theta$ et $\omega$ en $-\omega$, la valeur de $x$ deviendra celle de $y$; on aura donc
\[ y=\frac{\frac{1}{2}\omega}{\sin\frac{1}{2}\omega}\Sigma\Delta s\sin(\theta+\tfrac{1}{2}\omega)+Y(\theta)-Y(\alpha), \]
$Y(\theta)$ étant une fonction de $\theta$ donnée par l'équation
\[ \begin{aligned} Y(\theta)={}&A^{0}\omega^{2}\left(\frac{d(s\cos\theta)}{d\theta}+s\sin\theta\right)\\ &-B^{0}\omega^{4}\left(\frac{d^{3}(s\cos\theta)}{d\theta^{3}}-s\sin\theta\right)\\ &+C^{0}\omega^{6}\left(\frac{d^{5}(s\cos\theta)}{d\theta^{5}}+s\sin\theta\right)\\ &-\text{etc.}, \end{aligned} \]
et $Y(\alpha)$ une fonction semblable de $\alpha$.

\subsection*{(16).}

Nous remarquerons que les coefficiens $A^{0}$, $B^{0}$, $C^{0}$, etc. ont des rapports fort simples avec les coefficiens $A$, $B$, $C$, etc. dont nous avons fait usage dans le chapitre précédent (art.~3). En effet, si dans la formule
\[ 1-\tfrac{1}{2}\omega\cot\tfrac{1}{2}\omega=A^{0}\omega^{2}+B^{0}\omega^{4}+C^{0}\omega^{6}+\text{etc.}, \]
on met $2\omega$ à la place de $\omega$, on aura
\[ 1-\omega\cot\omega=A^{0}\cdot2^{2}\omega^{2}+B^{0}\cdot2^{4}\omega^{4}+C^{0}\cdot2^{6}\omega^{6}+\text{etc.} \]
Mais on a $\frac{1}{\sin\omega}=\cot\frac{1}{2}\omega-\cot\omega$; donc
\[ \frac{\omega}{\sin\omega}=1+(2^{2}-2)A^{0}\omega^{2}+(2^{4}-2)B^{0}\omega^{4}+(2^{6}-2)C^{0}\omega^{6}+\text{etc.}, \]
et il en résulte
\[ \frac{\frac{1}{2}\omega}{\sin\frac{1}{2}\omega}=1+\left(1-\frac{1}{2}\right)A^{0}\omega^{2}+\left(1-\frac{1}{2^{3}}\right)B^{0}\omega^{4}+\left(1-\frac{1}{2^{5}}\right)C^{0}\omega^{6}+\text{etc.} \]
Mais dans l'article cité, on a fait $\frac{\frac{1}{2}\omega}{\sin\frac{1}{2}\omega}=1+A\omega^{2}+B\omega^{4}+C\omega^{6}+\text{etc.}$; donc

\origpage{330}
\[ \begin{aligned} A&=\left(1-\frac{1}{2}\right)A^{0} &&\text{ou}\qquad & A^{0}&=\frac{2A}{2-1}=\frac{1}{12}\\ B&=\left(1-\frac{1}{2^{3}}\right)B^{0} &&& B^{0}&=\frac{2^{3}B}{2^{3}-1}=\frac{1}{720}\\ C&=\left(1-\frac{1}{2^{5}}\right)C^{0} &&& C^{0}&=\frac{2^{5}C}{2^{5}-1}=\frac{1}{30240}\\ D&=\left(1-\frac{1}{2^{7}}\right)D^{0} &&& D^{0}&=\frac{2^{7}D}{2^{7}-1}=\frac{1}{1209600}\\ &\text{etc.} &&& &\text{etc.} \end{aligned} \]
La suite $A^{0}$, $B^{0}$, $C^{0}$, $D^{0}$, etc. a cela de commun avec la suite $A$, $B$, $C$, $D$, etc., que le rapport de deux termes consécutifs converge vers la limite $\frac{1}{4\pi^{2}}$; mais ce rapport est moindre encore dans les premiers termes, puisqu'on a $\frac{B^{0}}{A^{0}}=\frac{1}{60}$, $\frac{C^{0}}{B^{0}}=\frac{1}{42}$, $\frac{D^{0}}{C^{0}}=\frac{1}{40}$, et on voit que dès le quatrième terme, le rapport est presque égal à sa limite.

Dans l'application de ces formules, il conviendra de prendre $\omega$ assez petit pour que le premier terme ou les deux premiers au plus des corrections $X(\theta)$, $Y(\theta)$ suffisent pour le degré d'approximation qu'on a en vue; ainsi tout dépend de la grandeur des coefficiens $\frac{ds}{d\theta}$, $\frac{dds}{d\theta^{2}}$, $\frac{d^{3}s}{d\theta^{3}}$, qu'il faudra calculer d'avance pour le premier et le dernier point de l'arc de courbe dont on veut connaître les coordonnées. Lorsque la courbe a très-peu de courbure, ces coefficiens sont très-grands; alors il faudra un plus grand arc de courbe $\Delta s$ pour répondre à une même différence $\Delta\theta$, et il n'est pas étonnant que les corrections à faire à la somme des différences finies, pour avoir la somme des différences infiniment petites, soient plus considérables. Il faudra donc prendre dans ce cas $\omega$ plus petit que si la courbure était plus sensible. Voici au reste un exemple dans lequel se trouvent réunies toutes les difficultés qu'on peut rencontrer dans ces sortes de calculs, avec les moyens de les surmonter.

\section*{Application de la méthode précédente au calcul de la trajectoire d'un projectile.}

\subsection*{(17).}

Nous prendrons pour exemple la courbe décrite par un projectile dans un milieu de densité constante, et dont la résistance

\origpage{331}
est proportionnelle au quarré de la vitesse. L'équation de cette courbe ne peut s'obtenir en termes finis que dans le cas d'une résistance très-petite ou d'un angle de projection très-petit; mais on peut y suppléer par une équation entre l'arc $s$ parcouru depuis le commencement du mouvement et l'angle $\theta$ que la tangente à la courbe fait avec l'horizon. Cette équation est
\[ e^{\frac{s}{k}}-1=\frac{h\cos^{2}\alpha}{k}[f(\alpha)-f(\theta)], \]
$f(\theta)$ désignant la fonction $\frac{\sin\theta}{\cos^{2}\theta}+\log\operatorname{tang}(45^{\circ}+\frac{1}{2}\theta)$, et $f(\alpha)$ une fonction semblable de l'angle $\alpha$, valeur initiale de $\theta$.

Supposons l'angle de projection $\alpha=45^{\circ}$, et la vitesse de projection telle qu'on ait $\frac{h}{k}=10$, nous aurons, en prenant $k$ pour l'unité,
\[ e^{s}=1+5f(\alpha)-5f(\theta). \]
D'après cette équation, il s'agit de trouver les valeurs de $x$ et $y$ au sommet de la courbe; on se proposera ensuite de calculer l'amplitude de la branche descendante.

Il semble d'abord qu'il suffit de calculer les valeurs successives de $\Delta x$ et $\Delta y$ en faisant varier $\theta$ de $5^{\circ}$, depuis $\theta=45^{\circ}$ jusqu'à $\theta=0^{\circ}$. Mais dans la partie de la courbe comprise depuis $\theta=45^{\circ}$ jusqu'à $\theta=40^{\circ}$, la valeur de $\Delta s$ se trouverait très-grande, parce que la courbure est très-petite dans cette partie, et c'est ce qui résulte des valeurs de $\frac{ds}{d\theta}$, $\frac{dds}{d\theta^{2}}$, etc., qui sont très-grandes lorsque $s=0$ ou $\theta=45^{\circ}$.

En effet de l'équation donnée on tire
\[ \begin{aligned} \frac{ds}{d\theta}&=-\frac{10e^{-s}}{\cos^{3}\theta}\\ \frac{dds}{d\theta^{2}}&=3\operatorname{tang}\theta\cdot\frac{ds}{d\theta}-\frac{ds^{2}}{d\theta^{2}}\\ \frac{d^{3}s}{d\theta^{3}}&=\left(3\operatorname{tang}\theta-2\frac{ds}{d\theta}\right)\frac{dds}{d\theta^{2}}+\frac{3}{\cos^{2}\theta}\cdot\frac{ds}{d\theta}. \end{aligned} \]
Si l'on fait dans ces valeurs $\theta=45^{\circ}$ et $s=0$, on aura

\origpage{332}
\[ \begin{aligned} \log\left(-\frac{ds}{d\theta}\right)&=1.451545\\ \log\left(-\frac{dds}{d\theta^{2}}\right)&=2.946871\\ \log\left(-\frac{d^{3}s}{d\theta^{3}}\right)&=4.723285. \end{aligned} \]
Ces quantités vont en augmentant assez rapidement, et en faisant varier $\theta$ de $5^{\circ}$, les corrections se trouveraient trop considérables. C'est pourquoi il convient de ne le faire varier que de $1^{\circ}$, et alors le second terme de la correction $X(\theta)-X(\alpha)$ n'influera que sur le septième ordre de décimales.

\subsection*{(18).}

Soit donc $\omega=-1^{\circ}=-\frac{\pi}{180}$, et proposons-nous d'abord de déterminer la portion de courbe comprise depuis $\theta=45^{\circ}$ jusqu'à $\theta=40^{\circ}$. Voici le calcul des valeurs successives de $\Delta s$, $\Delta s\cos(\theta+\frac{1}{2}\omega)$, $\Delta s\sin(\theta+\frac{1}{2}\omega)$.
\[ \begin{aligned} &\theta. && s. && \Delta s. && \Delta s\cos(\theta+\tfrac{1}{2}\omega). && \Delta s\sin(\theta+\tfrac{1}{2}\omega).\\ &45^{\circ} && 0.000000 && 0.392782 && 0.280152 && 0.275305\\ &44 && 0.392782 && 0.269115 && 0.195209 && 0.185247\\ &43 && 0.661897 && 0.202695 && 0.149442 && 0.136938\\ &42 && 0.864592 && 0.161284 && 0.120795 && 0.106870\\ &41 && 1.025876 && 0.133064 && 0.101183 && 0.086418\\ &40 && 1.158940 && \text{Sommes}\ldots && 0.846781 && 0.790778 \end{aligned} \]
Ces sommes doivent être multipliées par le facteur $\frac{\frac{1}{2}\omega}{\sin\frac{1}{2}\omega}$; or on a $\frac{\frac{1}{2}\omega}{\sin\frac{1}{2}\omega}=1+\frac{\omega^{2}}{24}+\frac{7\omega^{4}}{5760}$, et faisant $\Omega=\frac{\omega^{2}}{24}+\frac{7\omega^{4}}{5760}$, puis $\omega=-\frac{\pi}{180}$, on trouve $\log\Omega=5.103547$; ce qui donne
\[ \begin{aligned} &\text{Pour les sommes}\ldots\ldots && 0.846781\ldots\ldots && 0.790778\\ &\text{les corrections}\ldots\ldots && +\ 0.000011 && +\ 0.000010\\ & && 0.846792 && 0.790788. \end{aligned} \]
Pour avoir les corrections désignées par $X(\alpha)-X(\theta)$ et $Y(\alpha)-Y(\theta)$, il faut d'abord calculer les coefficiens $\frac{ds}{d\theta}$, $\frac{dds}{d\theta^{2}}$, $\frac{d^{3}s}{d\theta^{3}}$ pour la valeur

\origpage{333}
$\theta=40^{\circ}$. On trouvera par les formules ci-dessus,
\[ \begin{aligned} L\left(-\frac{ds}{d\theta}\right)&=0.843916\\ L\left(-\frac{dds}{d\theta^{2}}\right)&=1.821562\\ L\left(-\frac{d^{3}s}{d\theta^{3}}\right)&=3.052452. \end{aligned} \]
D'ailleurs on a $L\left(\frac{\omega^{2}}{12}\right)=5.404573$ et $L\left(\frac{\omega^{4}}{720}\right)=0.110176$; de là résulte
\[ \begin{aligned} X(\theta)&=-0.0001139 &\qquad Y(\theta)&=-0.0001357\\ X(\alpha)&=-0.0005026 & Y(\alpha)&=-0.0005031\\ X(\theta)-X(\alpha)&=\phantom{-}0.0003887 & Y(\theta)-Y(\alpha)&=\phantom{-}0.0003674. \end{aligned} \]
Appliquant ces corrections aux sommes trouvées,
\[ \begin{aligned} &\text{Somme des }\Delta x\ldots && 0.846792 &\qquad& \text{des }\Delta y\ldots && 0.790788\\ &\text{Correction}\ldots\ldots && -\ 0.000389 && \ldots\ldots && +\ 0.000367\\ &x= && 0.846403 && y= && 0.791155 \end{aligned} \]
on a les valeurs de $x$ et de $y$ qui répondent au point où $\theta=40^{\circ}$.

\subsection*{(19).}

On peut maintenant partir de ce point pour faire varier par de plus grands intervalles, les angles $\theta$: en prenant ces intervalles de $5^{\circ}$, on formera la table suivante:
\[ \begin{aligned} &\theta. && s. && \Delta s. && \Delta s\cos(\theta+\tfrac{1}{2}\omega). && \Delta s\sin(\theta+\tfrac{1}{2}\omega).\\ &40^{\circ} && 1.158940 && 0.438379 && 0.347789 && 0.266868\\ &35 && 1.597319 && 0.258678 && 0.218167 && 0.138988\\ &30 && 1.855997 && 0.178839 && 0.158632 && 0.082579\\ &25 && 2.034836 && 0.135291 && 0.124993 && 0.051774\\ &20 && 2.170127 && 0.108854 && 0.103816 && 0.032733\\ &15 && 2.278981 && 0.091783 && 0.089607 && 0.019865\\ &10 && 2.370764 && 0.080408 && 0.079720 && 0.010495\\ &5 && 2.451172 && 0.072792 && 0.072723 && 0.003175\\ &0 && 2.523964 && \text{Sommes}\ldots\ldots && 1.195447 && 0.606477 \end{aligned} \]
Lorsque $\omega=-5^{\circ}=-\frac{\pi}{36}$, on a $\log\Omega=6.501579$, ainsi à raison

\origpage{334}
du facteur $\frac{\frac{1}{2}\omega}{\sin\frac{1}{2}\omega}$, on aura les corrections suivantes:
\[ \begin{aligned} & 1.195447 &\qquad& 0.606477\\ +\ & 0.0003794 && 0.0001925\\ & 1.1958264 && 0.6066695. \end{aligned} \]
On a ensuite, en faisant $\theta=0$,
\[ \begin{aligned} L\left(-\frac{ds}{d\theta}\right)&=9.903857 &\qquad L\left(\frac{\omega^{2}}{12}\right)&=6.802513\\ L\left(-\frac{dds}{d\theta^{2}}\right)&=9.807714 & L\left(\frac{\omega^{4}}{720}\right)&=2.906056\\ L\left(-\frac{d^{3}s}{d\theta^{3}}\right)&=0.053577; \end{aligned} \]
d'où résulte
\[ \begin{aligned} X(0^{\circ})&=+0.0000002 &\qquad Y(0^{\circ})&=-0.0005085\\ X(40^{\circ})&=-0.0027781 & Y(40^{\circ})&=-0.0033358\\ X(0^{\circ})-X(40^{\circ})&=+0.0027783 & Y(0^{\circ})-Y(40^{\circ})&=+0.0028273\\ &\phantom{=+}1.1958264 & &\phantom{=+}0.6066695\\ x&=\phantom{+}1.193048 & y&=\phantom{+}0.609497 \end{aligned} \]
Ce sont les valeurs corrigées de $x$ et $y$, depuis $\theta=40^{\circ}$ jusqu'à $\theta=0$; si on leur ajoute les valeurs des mêmes coordonnées, depuis $\theta=45^{\circ}$ jusqu'à $\theta=40^{\circ}$, on aura les coordonnées qui répondent au sommet de la courbe
\[ x=2.039451,\quad y=1.400652. \]

\subsection*{(20).}

Il faut maintenant calculer la branche descendante, et pour cet effet, il faut changer le signe de $\theta$ dans l'équation des arcs, ce qui donnera
\[ e^{s}=1+5f(\alpha)+sf(\theta). \]
\ednote{Printed $sf(\theta)$ with an italic $s$; read $5f(\theta)$, as the constant $2.495587$ below requires.}

Lorsque $\theta$ deviendra $\theta+\omega$, l'arc $s$ deviendra $s+\Delta s$, et on aura $\Delta s$ par l'équation
\[ \Delta s=\log\left(\frac{2.495587+f(\theta+\omega)}{2.495587+f(\theta)}\right). \]
Faisant donc successivement $\theta=0^{\circ}$, $\theta=5^{\circ}$, $\theta=10^{\circ}$, etc., et $\omega=5^{\circ}$, on formera la table suivante:

\origpage{335}
\[ \begin{aligned} &\theta. && \Delta s. && \Delta s\cos(\theta+\tfrac{1}{2}\omega). && \Delta s\sin(\theta+\tfrac{1}{2}\omega).\\ &0^{\circ} && 0.067849 && 0.067785 && 0.002960\\ &5 && 0.064967 && 0.064411 && 0.008480\\ &10 && 0.063804 && 0.062291 && 0.013810\\ &15 && 0.064212 && 0.061240 && 0.019309\\ &20 && 0.066195 && 0.061156 && 0.025332\\ &25 && 0.069902 && 0.062003 && 0.032277\\ &30 && 0.075647 && 0.063800 && 0.040645\\ &35 && 0.083965 && 0.066614 && 0.051114\\ &40 && 0.095712 && 0.070566 && 0.064662\\ &45 && 0.112505 && 0.075833 && 0.082757\\ &50 && 0.135758 && 0.082644 && 0.107704\\ &55 && 0.169891 && 0.091282 && 0.143284\\ &60 && 0.221005 && 0.102049 && 0.196034\\ &65 && 0.300999 && 0.115187 && 0.278087\\ &70 && 0.434744 && 0.130730 && 0.414623\\ & && \text{Sommes}\ldots && 1.177591 && 1.481078\\ & && \text{Produit par }\Omega\ldots && 0.000374 && 0.000470\\ & && && 1.177965 && 1.481548 \end{aligned} \]
\ednote{Printed $0.112505$ in the $\Delta s$ column for $\theta=45^{\circ}$; the two products in the same row, and the formula for $\Delta s$, give $0.112246$, an apparent misprint. The column sums are unaffected.}

Pour corriger ces sommes, il faut avoir les valeurs des coefficiens $\frac{ds}{d\theta}$, $\frac{dds}{d\theta^{2}}$, $\frac{d^{3}s}{d\theta^{3}}$ lorsque $\theta=75^{\circ}$; alors ces coefficiens deviennent positifs, et on trouve leurs logarithmes comme il suit:
\[ \begin{aligned} \log\frac{ds}{d\theta}&=0.784599\\ \log\frac{dds}{d\theta^{2}}&=1.492717\\ \log\frac{d^{3}s}{d\theta^{3}}&=2.521532 \end{aligned} \]
Au moyen de ces valeurs, le calcul des corrections donne

\origpage{336}
\[ \begin{aligned} X(75^{\circ})&=0.0037066 &\qquad Y(75^{\circ})&=0.0010010\\ -X(0^{\circ})&=0.0000002 & Y(0^{\circ})&=0.0005085\\ X(75^{\circ})-X(0^{\circ})&=0.0037068 & Y(75^{\circ})-Y(0^{\circ})&=0.0004925\\ &\phantom{=}1.177965 & &\phantom{=}1.481548\\ x&=1.174258 & y&=1.482040. \end{aligned} \]
Ce sont les valeurs de $x$ et $y$ comprises depuis le sommet jusqu'au point où $\theta=75^{\circ}$.

Mais puisque la hauteur trouvée $y$ est plus grande que la hauteur de la branche ascendante $1.400652$, leur différence $0.081388$ exigera qu'il soit fait une diminution proportionnelle sur la valeur de $x$, pour avoir la vraie amplitude de la branche descendante. Et puisque la différence $0.414623$ répond à $5^{\circ}$ de différence dans l'angle $\theta$, on trouvera que $0.081388$ répond à une différence de $58'54''$; de sorte que l'angle de chute doit être environ $74^{\circ}1'6''$. Prenant le milieu entre cet angle et $75^{\circ}$, on aura $74^{\circ}30'33''$, et la quantité dont il faudra diminuer $x$ sera $0.081388\cot 74^{\circ}30'33''=0.022557$, d'où l'on conclura
\[ \begin{aligned} & && 1.174258\\ & && 0.022557\\ &\text{Ampl. de la branc. desc}\ldots\ldots && 1.151701\\ &\text{Ampl. de la branc. asc}\ldots\ldots && 2.039451\\ &\text{Ampl. totale}\ldots\ldots\ldots\ldots && 3.191152; \end{aligned} \]
ces résultats doivent être exacts, au moins jusqu'à la cinquième décimale.

\subsection*{(21).}

Ayant suivi le cours de la trajectoire jusqu'au point où $\theta=75^{\circ}$, il ne sera pas inutile de déterminer la position de l'asymptote verticale, c'est-à-dire, de chercher la valeur de $x$ lorsque $\theta=90^{\circ}$.

Pour cela nous compterons les $x$ du point où $\theta=75^{\circ}$, et en faisant $\operatorname{tang}\theta=p$, nous aurons à intégrer l'équation
\[ \tfrac{1}{2}\,dx=\frac{dp}{A+p\sqrt{(1+pp)}+\log[p+\sqrt{(1+pp)}]}, \]
dans laquelle $A=\frac{1}{5}+f(\alpha)=2.495587$, et il faudra que l'intégrale s'étende depuis $p=\operatorname{tang}75^{\circ}$ jusqu'à $p=\infty$.

\origpage{337}

Le moyen que nous emploierons pour cet objet, consiste à prendre successivement pour le dénominateur du second membre, une quantité toujours plus grande, ou toujours plus petite que ce dénominateur, et on conclura que la vraie valeur de $x$ est comprise entre les deux qui résulteront de chaque hypothèse.

Soit $D=A+p\sqrt{(1+pp)}+\log[p+\sqrt{(1+pp)}]$ et $m=\operatorname{tang}75^{\circ}$; puisque $p$ est toujours compris entre $\operatorname{tang}75^{\circ}$ et $\operatorname{tang}90^{\circ}$, on aura toujours $D>A+\log[m+\sqrt{(1+mm)}]+m\sqrt{(1+mm)}-m^{2}+p^{2}$; donc si on fait $c^{2}=A+f(75^{\circ})-\operatorname{tang}^{2}75^{\circ}=5.014525$, on aura
\[ \tfrac{1}{2}\,dx<\frac{dp}{c^{2}+p^{2}}. \]
Celle-ci donne en intégrant, $\frac{1}{2}cx<\operatorname{arc}\operatorname{tang}\frac{p}{c}-\operatorname{arc}\operatorname{tang}\frac{m}{c}$, et en faisant $p=\infty$, on aura $\frac{1}{2}cx<\frac{\pi}{2}-\operatorname{arc}\operatorname{tang}\frac{m}{c}$, ou $\operatorname{tang}\frac{1}{2}cx<\frac{c}{m}$; d'où résulte $x<0.48268$.

Pour avoir l'autre limite de $x$, je fais $n=\frac{1}{2m}\log[m+\sqrt{(1+mm)}]$, ce qui donne $\log n=9.434002$; j'observe ensuite qu'on aura dans toute l'étendue de l'intégrale, $p\sqrt{(1+pp)}<p^{2}+\frac{1}{2}$, $\log[p+\sqrt{(1+pp)}]<2np$; donc on a constamment
\[ \tfrac{1}{2}\,dx>\frac{dp}{A+\frac{1}{2}+2np+p^{2}}. \]
Soit $A+\frac{1}{2}-n^{2}=g^{2}$, ou $\log g=0.232825$, on aura
\[ \tfrac{1}{2}\,dx>\frac{dp}{g^{2}+(p+n)^{2}}; \]
d'où résulte en intégrant, $\frac{1}{2}gx>\operatorname{arc}\operatorname{tang}\frac{p+n}{g}-\operatorname{arc}\operatorname{tang}\frac{m+n}{g}$; Faisant $p=\infty$, on aura $\frac{1}{2}gx>\frac{\pi}{2}-\operatorname{arc}\operatorname{tang}\frac{m+n}{g}>\operatorname{arc}\operatorname{tang}\frac{g}{m+n}$, ou $x>0.47212$.

On voit donc que la valeur de $x$ qui à partir du point où $\theta=75^{\circ}$, répond à l'asymptote verticale, est certainement comprise entre deux limites assez rapprochées, savoir, entre $0.47212$ et $0.48268$. Par un milieu, on trouve $x=0.4774$, et cette valeur ne peut pas être en erreur de plus de $0.0053$.

\subsection*{(22).}

Il reste à connaître la position de l'autre asymptote du côté

\origpage{338}
des $x$ négatives. Or il suit de l'équation $e^{s}=1+5f(\alpha)-5f(\theta)$; qu'en faisant $s$ infini négatif, on aura $f(\theta)=\frac{1}{5}+f(\alpha)=2.495587$; d'après cette valeur, on trouve $\theta=46^{\circ}55'23''4$; c'est l'angle que fait l'asymptote dont il s'agit avec la ligne des abscisses.

En appelant cet angle $\beta$, si on considère à la fois $s$, $x$ et $y$ comme positifs pour tous les points de la branche qui descend vers l'asymptote; si ensuite on fait comme ci-dessus, $\operatorname{tang}\theta=p$, $\operatorname{tang}\beta=b$, $p\sqrt{(1+pp)}+\log[p+\sqrt{(1+pp)}]=\psi(p)$, on aura pour déterminer les coordonnées $x$ et $y$, les équations
\[ \tfrac{1}{2}\,dx=\frac{dp}{\psi(b)-\psi(p)},\quad \tfrac{1}{2}\,dy=\frac{p\,dp}{\psi(b)-\psi(p)}. \]
Faisons $x-\frac{y}{b}=z$, et la valeur de $z$ calculée pour le point où $p=b$, sera la distance de l'origine des abscisses au point où l'asymptote rencontre la ligne des abscisses. On aura donc à intégrer la formule suivante, depuis $p=1$ jusqu'à $p=b$,
\[ \tfrac{1}{2}b\,dz=\tfrac{1}{2}(b\,dx-dy)=\frac{(b-p)\,dp}{\psi(b)-\psi(p)}. \]
Comme $p$ diffère peu de $b$ dans l'intervalle où nous devons étendre l'intégrale, on peut faire $p=b-u$, et on aura la série fort convergente
\[ \psi(p)=\psi(b)-\psi'(b)u+\psi''(b)\cdot\frac{u^{2}}{2}-\psi'''(b)\cdot\frac{u^{3}}{2\cdot3}+\text{etc.}, \]
où l'on a $\psi'(b)=2\sqrt{(1+bb)}=\frac{2}{\cos\beta}$, $\psi''(b)=\frac{2b}{\sqrt{(1+bb)}}=2\sin\beta$, etc., et il faudra intégrer l'équation
\[ \tfrac{1}{2}b\,dz=\frac{-du}{\psi'(b)-\frac{u}{2}\psi''(b)+\frac{u^{2}}{2\cdot3}\psi'''(b)-\text{etc.}}. \]
Pour avoir par approximation cette intégrale, je fais
\[ \tfrac{1}{2}b\psi'(b)\,dz=-\frac{du}{1-Au}, \]
et j'observe que $A$ sera toujours compris entre deux valeurs $A'$ et $A''$, la première qui répond au commencement de l'intégrale, lorsque $u=b-1$, et la seconde qui répond à la fin de l'intégrale, lorsque $u=0$. Cette dernière se trouve exactement par le développement

\origpage{339}
indiqué qui donne $A''=\frac{1}{2}\cdot\frac{\psi''(b)}{\psi'(b)}=\frac{1}{2}\sin\beta\cos\beta$; l'autre valeur $A'$ se trouvera par l'équation
\[ 1-A'(b-1)=\frac{\psi(b)-\psi(1)}{(b-1)\psi'(b)}=\frac{0.2}{(b-1)\psi'(b)}. \]
Ainsi on aura $\log A'=9.3921695$ et $\log A''=9.3969607$.

Mais en regardant $A$ comme constante, l'intégrale de l'équation précédente est $\frac{1}{2}b\psi'(b)\cdot z=\frac{1}{A}\log(1-Au)+\text{const.}$ Prenant donc l'intégrale entre les limites $u=b-1$, $u=0$, on aura
\[ z=-\frac{\log[1-A(b-1)]}{\frac{1}{2}bA\psi'(b)}. \]
Substituant successivement au lieu de $A$ les deux valeurs $A'$ et $A''$, on trouvera $z=0.0447606$ et $z=0.0447650$. Le peu de différence qu'il y a entre ces deux valeurs, prouve combien cette détermination est exacte, et par un milieu pris entre elles, on aura encore plus exactement $z=0.0447628$. C'est l'abscisse cherchée du point où l'asymptote dont il s'agit rencontre la ligne des abscisses.

\section*{De l'intégrale indéfinie $\displaystyle\int dx\left(\log\frac{1}{x}\right)^{a-1}$, prise depuis $x=0$.}

\subsection*{(23).}

Nous avons fait voir dans la seconde partie, comment on trouve l'intégrale $\displaystyle\int dx\left(l\frac{1}{x}\right)^{a-1}$ lorsque $x=1$, et nous l'avons désignée dans ce cas par la fonction $\Gamma(a)$. Mais il peut être utile de déterminer cette intégrale pour une valeur quelconque de $x$: nous la représenterons généralement par $\Gamma(a,x)$, de sorte que $\Gamma(a,1)$ sera la même chose que $\Gamma(a)$.

Observons d'abord que l'intégration par parties donne la formule
\[ \Gamma(a,x)=x\left(\log\frac{1}{x}\right)^{a-1}+(a-1)\Gamma(a-1,x), \tag{1} \]
d'où il suit que l'intégrale $\Gamma(a,x)$ peut toujours se ramener à une intégrale semblable, dans laquelle $a$ est compris entre $1$ et $2$, ainsi que nous l'avons fait pour l'intégrale définie $\Gamma(a)$.

Soit $\log\frac{1}{x}=z$, on aura la transformée
\[ \Gamma(a,x)=\int -z^{a-1}\,dz\,e^{-z}=\int -z^{a-1}\,dz\left(1-z+\frac{z^{2}}{2}-\frac{z^{3}}{2\cdot3}+\text{etc.}\right); \]

\origpage{340}
d'où l'on tire en intégrant,
\[ \Gamma(a,x)=\Gamma(a)-z^{a}\left(\frac{1}{a}-\frac{z}{a+1}+\frac{1}{2}\cdot\frac{z^{2}}{a+2}-\frac{1}{2\cdot3}\cdot\frac{z^{3}}{a+2}+\text{etc.}\right). \tag{2} \]
\ednote{Printed $\frac{z^{3}}{a+2}$ in the fourth term; the integration requires $\frac{z^{3}}{a+3}$, an apparent misprint.}
Cette formule donne la valeur de $\Gamma(a,x)$ par une suite qui peut être divergente dans les premiers termes, mais qui finit toujours par être convergente.

Cette suite est convergente dès les premiers termes, si $x$ n'est pas plus petit que $\frac{1}{e}$; elle peut néanmoins être employée avec succès pour des valeurs beaucoup plus petites de $x$, telles que $x=\frac{1}{10}$, ou même $x=\frac{1}{20}$; mais alors il faudra calculer un assez grand nombre de termes de la série; par exemple, si $x=\frac{1}{10}$, il faudra calculer douze à treize termes de la série pour avoir la valeur de l'intégrale approchée jusqu'à la cinquième ou la sixième décimale. On trouve de cette manière $\Gamma\left(\frac{1}{2},\frac{1}{10}\right)=0.056497$.

A mesure que $x$ devient plus petit, il faudra prolonger plus loin la suite pour obtenir un égal degré d'approximation; de sorte qu'il convient de recourir à un autre moyen pour évaluer l'intégrale avec précision, lorsque $x$ est très-petit, tel que $\frac{1}{100}$, $\frac{1}{1000}$, etc.

\subsection*{(24).}

La meilleure méthode pour calculer l'intégrale $\Gamma(a,x)$ lorsque $x$ est très-petit, est de la déduire de la formule (1) qui donne, par des transformations successives,
\[ \Gamma(a,x)=x[z^{a-1}+(a-1)z^{a-2}+(a-1)(a-2)z^{a-3}+\text{etc.}] \tag{3} \]
Cette série étant continuée suffisamment, ses termes deviendront alternativement positifs et négatifs, de sorte qu'on aura des valeurs alternativement plus grandes et plus petites que l'intégrale cherchée. Mais comme la suite qui est d'abord convergente, devient nécessairement divergente après un certain nombre de termes, il faudra s'arrêter au point où la divergence commence, et on n'aura ainsi qu'une approximation bornée.

Si on appelle $P^{n}$ le terme de rang $n$ dans la formule (3), et $P^{n+1}$ le terme suivant, on aura $P^{n+1}=\frac{a-n}{z}P^{n}$; ainsi la divergence commencera

\origpage{341}
lorsqu'on aura $n=$ ou $>a+z$, et alors la suite ne devra pas être prolongée plus loin.

\subsection*{(25).}

Soit, par exemple, $a=\frac{1}{2}$, $x=\frac{1}{100}$, on aura $z=\log 100=4.60517$, et $a+z=5.1$. Ainsi la suite ne devra pas être prolongée au-delà des 5 ou 6 premiers termes. En effet, on a en faisant $z=\log 10$\ednote{Printed $\log 10$; the value used, $z=4.60517$, is $\log 100$ as in the line above, an apparent misprint.},
\[ \Gamma\left(\frac{1}{2},\frac{1}{100}\right)=\frac{1}{100}\left(z^{-\frac{1}{2}}-\frac{1}{2}z^{-\frac{3}{2}}+\frac{1\cdot3}{2\cdot2}z^{-\frac{5}{2}}-\frac{1\cdot3\cdot5}{2\cdot2\cdot2}z^{-\frac{7}{2}}+\text{etc.}\right), \]
et par le calcul de termes successifs, on trouve
\[ \begin{aligned} z^{-\frac{1}{2}}&=\phantom{+\ }0.465991\\ -\frac{1}{2}z^{-\frac{3}{2}}&=-\ 0.050594\\ &\phantom{=+\ }0.415397\\ +\frac{1\cdot3}{2\cdot2}z^{-\frac{5}{2}}&=+\ 0.016479\\ &\phantom{=+\ }0.431876\\ -\frac{1\cdot3\cdot5}{2\cdot2\cdot2}z^{-\frac{7}{2}}&=-\ 0.008946\\ &\phantom{=+\ }0.422930\\ +\frac{1\cdot3\cdot5\cdot7}{2\cdot2\cdot2\cdot2}z^{-\frac{9}{2}}&=+\ 0.006799\\ &\phantom{=+\ }0.429729\\ -\frac{1\cdot3\cdot5\cdot7\cdot9}{2\cdot2\cdot2\cdot2\cdot2}z^{-\frac{11}{2}}&=-\ 0.006644\\ &\phantom{=+\ }0.423085\\ +\frac{1\cdot3\ldots11}{2\cdot2\ldots2}z^{-\frac{13}{2}}&=+\ 0.007935 \end{aligned} \]
On voit qu'il est inutile d'aller plus loin que le sixième terme, puisque la suite devient divergente à compter de ce terme.

Il résulte de ce calcul que la somme cherchée est plus grande que $0.423085$, et plus petite que $0.429729$. Par un milieu, on trouve la somme de la suite $=0.4264$, et on voit qu'on ne peut compter sur plus de quatre décimales exactes dans cette évaluation; il en résulte pour la valeur de l'intégrale cherchée, $\Gamma(a,x)=0.004264$.

\origpage{342}
\subsection*{(26).}

Si en prenant toujours $a=\frac{1}{2}$, on eût fait $x=\frac{1}{1000}$, on aurait vu d'abord que la suite ne doit être prolongée que jusqu'au neuvième terme, et par le calcul effectif des termes, on aurait trouvé que l'intégrale cherchée est comprise entre $x\times0.357708$ et $x\times0.357173$. Le milieu est $x\times0.35744$; d'où résulte l'intégrale $\Gamma(a,x)=0.00035744$.

En général la formule (3) donnera d'autant plus de précision, que $x$ sera plus petit; mais cette précision ne s'obtiendra que par le calcul d'un plus grand nombre de termes, puisque pour tirer de la suite toute l'approximation qu'elle peut offrir, il faut la prolonger jusqu'à ce que le nombre de ses termes soit $a+z$ ou $a+\log\frac{1}{x}$.

Si on appliquait la formule (3) au cas de $x=\frac{1}{10}$ déjà résolu par la formule (2), on trouverait que la suite cesse d'être convergente au quatrième terme. Les deux premiers termes donnent $\Gamma>x\times0.5159$, et les trois premiers donnent $\Gamma<x\times0.6091$; il en résulte par un milieu, $\Gamma=x\times0.5625=0.05625$, tandis que la vraie valeur est $0.056497$. Il convient donc de préférer la formule (2) lorsqu'on voudra avoir au moins quatre chiffres significatifs exacts, et que la valeur de $x$ ne sera pas plus petite que $0.01$.

\subsection*{(27).}

Jusqu'ici nous avons supposé tacitement que $x$ ne surpasse pas l'unité; mais on peut aussi demander l'intégrale $\displaystyle\int dx\left(l\frac{1}{x}\right)^{a-1}$ pour une valeur de $x$ plus grande que l'unité. La partie comprise depuis $x=0$ jusqu'à $x=1$, est connue et représentée par $\Gamma(a)$; ainsi tout se réduit à trouver l'intégrale depuis $x=1$ jusqu'à une valeur quelconque de $x>1$.

Remarquons d'abord que dans tout cet intervalle, $l\left(\frac{1}{x}\right)$ étant négatif, il faudra mettre l'intégrale sous la forme $(-1)^{a-1}\displaystyle\int dx(lx)^{a-1}$. Je fais abstraction du facteur $(-1)^{a-1}$ qui peut être réel ou imaginaire, suivant les diverses valeurs de $a$, et je considère simplement l'intégrale $\displaystyle\int dx(lx)^{a-1}$ que je désigne par $\psi(a,x)$, et qui est supposée nulle lorsque $x=1$.

Si on fait $lx=u$, on aura $x=e^{u}$, et l'intégrale dont il s'agit

\origpage{343}
deviendra $\displaystyle\int u^{a-1}\,du\cdot e^{u}=\displaystyle\int u^{a-1}\,du\left(1+u+\frac{u^{2}}{2}+\text{etc.}\right)$; d'où l'on tire
\[ \psi(a,x)=\frac{u^{a}}{a}+\frac{u^{a+1}}{a+1}+\frac{1}{2}\cdot\frac{u^{a+2}}{a+2}+\frac{1}{2\cdot3}\cdot\frac{u^{a+3}}{a+3}+\text{etc.}, \tag{4} \]
formule où il n'y a pas de constante à ajouter, parce qu'elle s'évanouit lorsque $x=1$.

La suite comprise dans cette intégrale pourra être divergente dans les premiers termes; mais elle finira toujours par être convergente. Ainsi on en tirera dans tous les cas une valeur aussi approchée qu'on voudra de la vraie intégrale, et il est visible que cette valeur deviendra infinie si on fait $x=\infty$.

\subsection*{(28).}

On peut encore mettre $\psi(a,x)$ sous une forme plus commode. Soit $\displaystyle\int u^{a-1}\,du\cdot e^{u}=e^{u}P$, on aura $u^{a-1}=\frac{dP}{du}+P$; soit ensuite $P=\frac{u^{a}}{a}+Au^{a+1}+Bu^{a+2}+\text{etc.}$, on trouvera $A=-\frac{1}{a\cdot a+1}$, $B=\frac{1}{a\cdot a+1\cdot a+2}$, etc.; donc
\[ \psi(a,x)=x\left(\frac{u^{a}}{a}-\frac{u^{a+1}}{a\cdot a+1}+\frac{u^{a+2}}{a\cdot a+1\cdot a+2}-\text{etc.}\right), \tag{5} \]
série qui aura, comme la précédente, la propriété de devenir toujours convergente, mais qui aura de plus l'avantage de donner des valeurs alternativement plus grandes et plus petites que l'intégrale cherchée; de sorte que le degré d'approximation sera connu dans chaque cas par la différence de deux termes consécutifs.

\section*{De l'intégrale $\displaystyle\int y\,dx$ prise entre deux limites qui rendent nulle la fonction $y$.}

\subsection*{(29).}

Si la fonction $y$ est nulle aux deux limites de l'intégrale; si on suppose en même temps qu'elle ne change point de signe d'une limite à l'autre, et que dans cet intervalle elle ne soit susceptible que d'un seul \emph{maximum}, on pourra faire usage de la méthode suivante pour déterminer l'intégrale $Z=\displaystyle\int y\,dx$~${}^{(1)}$.

Soit $m$ la valeur de $x$ qui rend $y$ un \emph{maximum}, et soit ce \emph{maximum}

\textbf{(1)} Mém.\ de l'Acad.\ des Sc., ann.\ 1778 et 1782.

\origpage{344}
$=M$, on fera en général
\[ y=Me^{-t^{2}}, \]
$t$ étant une nouvelle variable qui s'étendra depuis $t=-\infty$, jusqu'à $t=+\infty$, ces valeurs répondant aux deux limites de l'intégrale où l'on a $y=0$.

D'après la valeur supposée de $y$, on trouvera une expression de $x$ en fonction de $t$, qui sera de la forme
\[ x=m+At+Bt^{2}+Ct^{3}+Dt^{4}+\text{etc.}, \]
et on aura
\[ \int y\,dx=M\int e^{-t^{2}}\,dt(A+2Bt+3Ct^{2}+4Dt^{3}+\text{etc.}). \]
Il faut distinguer dans cette intégrale deux parties; l'une depuis $t=-\infty$ jusqu'à $t=0$; l'autre depuis $t=0$ jusqu'à $t=\infty$. Si on change le signe de la première, elle devra être prise depuis $t=0$ jusqu'à $t=\infty$, et sera représentée par la formule
\[ M\int e^{-t^{2}}\,dt(A-2Bt+3Ct^{2}-4Dt^{3}+\text{etc.}); \]
la seconde, qui devra être prise entre les mêmes limites, sera
\[ M\int e^{-t^{2}}\,dt(A+2Bt+3Ct^{2}+4Dt^{3}+\text{etc.}). \]
Ajoutant ces deux parties, on a l'intégrale totale
\[ Z=2M\int e^{-t^{2}}\,dt(A+3Ct^{2}+5Et^{4}+\text{etc.}). \]
Or en faisant $e^{-t^{2}}=z$, l'intégrale $\displaystyle\int t^{2n}e^{-t^{2}}\,dt$ aura pour transformée $\frac{1}{2}\displaystyle\int dz\left(l\frac{1}{z}\right)^{\frac{2n-1}{2}}$, et celle-ci devant être prise depuis $z=0$ jusqu'à $z=1$, sera représentée par $\frac{1}{2}\Gamma\left(\frac{2n+1}{2}\right)$. Ainsi on aura généralement $\displaystyle\int e^{-t^{2}}t^{2n}\,dt=\frac{1}{2}\Gamma\left(\frac{2n+1}{2}\right)$, d'où résulte en particulier
\[ \int e^{-t^{2}}\,dt=\tfrac{1}{2}\sqrt{\pi},\quad \int t^{2}e^{-t^{2}}\,dt=\tfrac{1}{2}\sqrt{\pi}\cdot\tfrac{1}{2},\quad \int t^{4}e^{-t^{2}}\,dt=\tfrac{1}{2}\sqrt{\pi}\cdot\frac{1\cdot3}{2\cdot2},\ \text{etc.} \]
Donc enfin l'intégrale cherchée,
\[ Z=M\sqrt{\pi}\cdot\left(A+\frac{3}{2}C+\frac{3\cdot5}{2,2}E+\frac{3\cdot5\cdot7}{2\cdot2\cdot2}G+\text{etc.}\right); \tag{1} \]
\ednote{The first denominator is printed $2,2$ with a comma, where the numerators and the next denominator have product dots; kept as printed. The separators in $2\cdot2\cdot2$ are faint and could also be commas.}

\origpage{345}
et en considérant $x$ comme une fonction connue de $t$, faisant ensuite dans les coefficiens différentiels $t=0$, on aura
\[ Z=M\sqrt{\pi}\left(\frac{dx}{dt}+\frac{1}{2^{2}}\cdot\frac{d^{3}x}{dt^{3}}+\frac{1}{1\cdot2}\cdot\frac{1}{2^{4}}\cdot\frac{d^{5}x}{dt^{5}}+\frac{1}{1\cdot2\cdot3}\cdot\frac{1}{2^{6}}\cdot\frac{d^{7}x}{dt^{7}}+\text{etc.}\right). \]

\subsection*{(30).}

Si dans l'équation $y=Me^{-t^{2}}$, ou $\log M-\log y=t^{2}$, on substitue la valeur de $y$ en $x$, on aura une équation entre $x$ et $t$, d'où on devra déduire les valeurs des coefficiens $\frac{dx}{dt}$, $\frac{d^{3}x}{dt^{3}}$, etc. qui entrent dans la formule précédente, et dans lesquels on fera ensuite $t=0$ ou $x=m$. C'est dans les exemples particuliers qu'il convient de déterminer ces coefficiens; cependant on peut aussi les déduire généralement des coefficiens $\frac{ddy}{dx^{2}}$, $\frac{d^{3}y}{dx^{3}}$, etc. supposés connus au point du \emph{maximum}.

En effet l'équation $\log M-\log y=t^{2}$ étant différentiée, donne $\frac{dy}{dx}\cdot\frac{dx}{dt}+2ty=0$. Si on y fait $x=m$ ou $t=0$, on n'en tire aucun résultat, parce qu'alors $\frac{dy}{dx}=0$; mais si on différentie une seconde fois, on aura
\[ \frac{ddy}{dx^{2}}\cdot\frac{dx^{2}}{dt^{2}}+\frac{dy}{dx}\cdot\frac{ddx}{dt^{2}}+2y+2t\frac{dy}{dx}\cdot\frac{dx}{dt}=0; \]
faisant dans celle-ci $t=0$, on aura $\frac{ddy}{dx^{2}}\cdot\frac{dx^{2}}{dt^{2}}+2y=0$; d'où résulte $\frac{dx}{dt}$, ou
\[ A=\sqrt{\left(\frac{2y}{-\frac{ddy}{dx^{2}}}\right)}. \]
Si donc la série (1) est assez convergente pour qu'on puisse la réduire à son premier terme, on aura l'intégrale cherchée
\[ Z=y\sqrt{\left(\frac{2\pi y}{-\frac{ddy}{dx^{2}}}\right)}, \]
$y$ et $\frac{ddy}{dx^{2}}$ étant les valeurs de ces quantités au point du \emph{maximum} où $x=m$.

En différentiant ultérieurement l'équation dont nous avons tiré

\origpage{346}
la valeur de $\frac{dx}{dt}$, on trouverait les valeurs des autres coefficiens $\frac{ddx}{dt^{2}}$, $\frac{d^{3}x}{dt^{3}}$, etc., exprimées en fonctions des quantités $y$, $\frac{ddy}{dx^{2}}$, $\frac{d^{3}y}{dx^{3}}$, etc. au point du \emph{maximum}. Mais ces valeurs se compliquent à mesure qu'on pousse le calcul plus loin, et il vaut mieux, comme nous l'avons dit, déterminer ces coefficiens dans les cas particuliers.

Pour peu qu'on réfléchisse sur l'esprit de la méthode précédente, on verra qu'elle doit donner un résultat d'autant plus approché, que la fonction $y$ décroîtra plus promptement en s'éloignant du \emph{maximum}; et c'est ce qui arrivera si les facteurs dont $y$ est composé sont des puissances d'un ordre fort élevé. Cette circonstance permet de n'avoir égard qu'à une petite partie de l'intégrale, depuis $x=m-\alpha$ jusqu'à $x=m+\alpha$; elle contribuerait aussi à simplifier l'usage des autres méthodes; mais celles-ci ne donneraient pas un résultat aussi élégant, parce qu'il resterait toujours quelque chose de vague dans la détermination de $\alpha$. Au reste, on prendra une idée plus juste de cette méthode par les exemples suivans.

\section*{EXEMPLE I.}

\subsection*{(31).}

Soit proposé de trouver l'intégrale $Z=\displaystyle\int x^{\alpha}e^{-x}\,dx$, depuis $x=0$ jusqu'à $x=\infty$, $\alpha$ étant un nombre positif.

Dans ce cas, la fonction $x^{\alpha}e^{-x}$ est nulle aux deux limites de l'intégrale; sa valeur dans cet intervalle est toujours positive, et elle n'est susceptible que d'un seul \emph{maximum}, qui a lieu lorsque $x=\alpha$; ainsi la méthode précédente peut être appliquée à cette intégrale.

Soit donc $y=x^{\alpha}e^{-x}=\left(\frac{\alpha}{e}\right)^{\alpha}e^{-t^{2}}$, il en résultera
\[ \alpha\log\frac{x}{\alpha}-(x-\alpha)=-t^{2}, \]
et en supposant $x=\alpha+u$, on aura par le développement de cette équation,
\[ \frac{u^{2}}{2\alpha^{2}}-\frac{u^{3}}{3\alpha^{3}}+\frac{u^{4}}{4\alpha^{4}}-\text{etc.}=\frac{t^{2}}{\alpha}. \]
Soit, comme ci-dessus, $u=At+Bt^{2}+Ct^{3}+Dt^{4}+\text{etc.}$, on trouvera
\[ A=\sqrt{2\alpha},\quad B=\frac{2}{3},\quad C=\frac{1}{9A},\quad D=-\frac{4}{135A^{2}},\quad E=\frac{1}{270A^{3}},\ \text{etc.} \]

\origpage{347}
Substituant ces valeurs et celle de $M=\left(\frac{\alpha}{e}\right)^{\alpha}$ dans la formule (1), on aura l'intégrale cherchée
\[ Z=\left(\frac{\alpha}{e}\right)^{\alpha}\sqrt{(2\alpha\pi)}\left[1+\frac{1}{12\alpha}+\frac{1}{288\alpha^{2}}+\text{etc.}\right]\ldots\ldots \tag{2} \]

Si dans l'intégrale proposée $Z=\displaystyle\int x^{\alpha}e^{-x}\,dx$, on fait $e^{-x}=z$, on aura la transformée $Z=\displaystyle\int dz\left(l\frac{1}{z}\right)^{\alpha}$ qui devra être intégrée depuis $z=0$ jusqu'à $z=1$. On a donc $Z=\Gamma(\alpha+1)$; et en effet, la valeur que nous venons de trouver pour $Z$ s'accorde avec la formule $(\sigma)$ du n° 67, seconde partie.

\subsection*{(32).}

Cette formule résout très-bien le cas où $\alpha$ est un grand nombre, puisqu'alors elle donne une suite fort convergente (au moins dans les premiers termes); mais elle ne résout qu'imparfaitement le cas où $\alpha$ est compris entre 0 et 1, et elle donne un résultat imaginaire lorsque $\alpha$ est négatif et plus petit que l'unité, quoiqu'alors la valeur de $\Gamma(\alpha+1)$ soit réelle.

Il est facile de remédier à l'inconvénient que présentent les deux derniers cas; car l'intégration par parties donne
\[ \int x^{\alpha}e^{-x}\,dx=\frac{1}{\alpha+1}x^{\alpha+1}e^{-x}+\frac{1}{\alpha+1}\int x^{\alpha+1}e^{-x}\,dx; \]
et comme le terme hors du signe s'évanouit dans les deux limites de l'intégrale, on a simplement
\[ \int x^{\alpha}e^{-x}\,dx=\frac{1}{\alpha+1}\int x^{\alpha+1}e^{-x}\,dx; \]
d'où il suit qu'étant proposée l'intégrale $\displaystyle\int x^{\alpha}e^{-x}\,dx$, on peut la transformer en une autre où l'exposant de $x$ sera aussi grand qu'on voudra; et alors on déterminera celle-ci d'une manière fort approchée par la formule (2).

Il est facile de voir pourquoi dans cet exemple le résultat de la formule est d'autant plus exact, que $\alpha$ est plus grand; c'est que le facteur $x^{\alpha}$ décroît d'autant plus rapidement dans le voisinage du \emph{maximum}, que $\alpha$ est plus grand. Il n'y a plus de \emph{maximum} lorsque $\alpha$ est négatif; c'est pourquoi le résultat de la formule est entièrement fautif dans ce cas, quoique la vraie valeur de l'intégrale soit réelle, tant que $1+\alpha$ est positif.

\origpage{348}

\section*{EXEMPLE II.}

\subsection*{(33).}

Soit proposé de trouver l'intégrale $Z=\displaystyle\int x^{\alpha}(1-x)^{\beta}\,dx$, entre les limites $x=0$, $x=1$, les nombres $\alpha$ et $\beta$ étant positifs.

Ayant fait $y=x^{\alpha}(1-x)^{\beta}$, on trouve que $y$ est un \emph{maximum} lorsque $x=\frac{\alpha}{\alpha+\beta}=m$. Soit donc en général
\[ x^{\alpha}(1-x)^{\beta}=m^{\alpha}(1-m)^{\alpha}e^{-t^{2}}; \]
\ednote{Printed $(1-m)^{\alpha}$ instead of $(1-m)^{\beta}$; an apparent misprint.}
Si on prend les logarithmes de part et d'autre, et qu'on fasse $x=m+u$, on aura
\[ \alpha\log\left(1+\frac{u}{m}\right)+\beta\log\left(1-\frac{u}{1-m}\right)=-t^{2}, \]
ou
\[ \frac{u^{2}}{2}\left(\frac{\alpha}{m^{2}}+\frac{\beta}{(1-m)^{2}}\right)-\frac{u^{3}}{3}\left(\frac{\alpha}{m^{3}}-\frac{\beta}{(1-m)^{3}}\right)+\frac{u^{4}}{4}\left(\frac{\alpha}{m^{4}}+\frac{\beta}{(1-m)^{4}}\right)-\text{etc.}=t^{2}; \]
Soit, comme ci-dessus, $u=At+Bt^{2}+Ct^{3}+\text{etc.}$, on trouvera
\[ A=\frac{1}{\alpha+\beta}\sqrt{\left(\frac{2\alpha\beta}{\alpha+\beta}\right)},\quad B=\frac{2}{3}\cdot\frac{\beta-\alpha}{(\alpha+\beta)^{2}},\quad C=\frac{\alpha^{2}-11\alpha\beta+\beta^{2}}{18\alpha\beta(\alpha+\beta)}A,\ \text{etc.}; \]
d'où résulte l'intégrale cherchée
\[ Z=\frac{\alpha^{\alpha}\beta^{\beta}\sqrt{(2\pi\alpha\beta)}}{(\alpha+\beta)^{\alpha+\beta+\frac{3}{2}}}\left(1+\frac{\alpha^{2}-11\alpha\beta+\beta^{2}}{12\alpha\beta(\alpha+\beta)}+\text{etc.}\right). \]

Si les exposans $\alpha$ et $\beta$ sont tous les deux de grands nombres, cette suite sera fort convergente; mais si l'un des deux seulement est un grand nombre, la suite ne sera que peu convergente, de sorte qu'il faudrait en calculer beaucoup de termes pour n'avoir qu'une médiocre approximation: or ces termes sont difficiles à calculer par la méthode que nous venons de suivre.

\subsection*{(34).}

Examinons particulièrement le cas où $\alpha$ et $\beta$ sont tous deux de grands nombres; alors la série se réduira sensiblement à son premier terme, et on aura
\[ Z=\frac{\alpha^{\alpha+\frac{1}{2}}\beta^{\beta+\frac{1}{2}}\sqrt{2\pi}}{(\alpha+\beta)^{\alpha+\beta+\frac{3}{2}}}. \]

\origpage{349}
Mais l'intégrale exacte, lorsque $\alpha$ et $\beta$ sont des entiers, est
\[ Z^{1}=\frac{1\cdot2\cdot3\ldots\alpha}{\beta+1\cdot\beta+2\ldots\beta+\alpha+1}=\frac{1\cdot2\cdot3\ldots\alpha\cdot1\cdot2\cdot3\ldots\beta}{1\cdot2\cdot3\ldots\alpha+\beta+1}; \]
Cette intégrale peut s'exprimer dans le même cas par la formule
\[ Z^{1}=\frac{\Gamma(\alpha+1)\cdot\Gamma(\beta+1)}{(\alpha+\beta+1)\,\Gamma(\alpha+\beta+1)}; \]
et alors elle a lieu pour toutes valeurs de $\alpha$ et $\beta$. Si ensuite on substitue les valeurs des fonctions $\Gamma$ d'après la formule
\[ \Gamma(z+1)=\left(\frac{z}{e}\right)^{z+\frac{1}{2}}\sqrt{(2e\pi)}\cdot\Phi(z), \]
$\Phi(z)$ étant mis pour la fonction $1+\frac{1}{12z}+\frac{1}{288z^{2}}+\text{etc.}$, on trouve l'intégrale cherchée
\[ Z^{1}=\frac{\alpha^{\alpha+\frac{1}{2}}\beta^{\beta+\frac{1}{2}}\sqrt{2\pi}}{(\alpha+\beta+1)(\alpha+\beta)^{\alpha+\beta+\frac{1}{2}}}\cdot\frac{\Phi(\alpha)\cdot\Phi(\beta)}{\Phi(\alpha+\beta)}. \]
Mais en développant la quantité $\frac{\Phi(\alpha)\,\Phi(\beta)}{\Phi(\alpha+\beta)}$, ou, ce qui revient au même, $\Phi(\alpha)\cdot\Phi(\beta)\cdot\Phi(-\alpha-\beta)$, on a pour les trois premiers termes du développement
\[ 1+\frac{\alpha^{2}+\alpha\beta+\beta^{2}}{12\alpha\beta(\alpha+\beta)}+\frac{(\alpha^{2}+\alpha\beta+\beta^{2})^{2}}{288\alpha^{2}\beta^{2}(\alpha+\beta)^{2}}. \]
Donc la valeur de $Z^{1}$, développée jusqu'au troisième terme de la série, sera
\[ Z^{1}=\frac{\alpha^{\alpha+\frac{1}{2}}\beta^{\beta+\frac{1}{2}}\sqrt{2\pi}}{(\alpha+\beta)^{\alpha+\beta+\frac{3}{2}}}\left(1+\frac{\alpha^{2}-11\alpha\beta+\beta^{2}}{12\alpha\beta(\alpha+\beta)}+\frac{(\alpha^{2}-11\alpha\beta+\beta^{2})^{2}+144\alpha^{2}\beta^{2}}{288\alpha^{2}\beta^{2}(\alpha+\beta)^{2}}\right), \]
formule qui s'accorde entièrement avec la valeur trouvée pour $Z$, laquelle n'a été calculée que jusqu'au second terme.

On voit maintenant \emph{a priori} pourquoi la valeur trouvée pour $Z$ ne donne pas \uncertain{une} suite convergente lorsque l'un seulement des exposans $\alpha$ et $\beta$ \uncertain{est} un grand nombre; c'est que la suite dont il s'agit étant égale au développement de la fonction $\frac{\Phi(\alpha)\,\Phi(\beta)}{\Phi(\alpha+\beta)}\left(\frac{\alpha+\beta}{\alpha+\beta+1}\right)$, est bien convergente par rapport à la fonction $\Phi(\alpha+\beta)$, et à l'une des fonctions $\Phi(\alpha)$, $\Phi(\beta)$, mais ne l'est pas par rapport à

\origpage{350}
l'autre; et il en résulte que la suite totale n'est pas plus convergente que celle des deux suites $\Phi(\alpha)$, $\Phi(\beta)$, qui répond au moindre des deux nombres $\alpha$ et $\beta$.

\subsection*{(35).}

L'hypothèse de l'art.\ précédent ayant toujours lieu, si on fait de plus $\alpha=\beta$, on aura la valeur très-approchée
\[ Z=\left(\tfrac{1}{2}\right)^{2\alpha+1}\sqrt{\left(\frac{\pi}{\alpha}\right)}\cdot\left(1-\frac{3}{8\alpha}+\frac{25}{144\alpha^{2}}-\text{etc.}\right) \]
\ednote{Printed $\frac{25}{144\alpha^{2}}$; the expansion (art.~34 with $\alpha=\beta$) gives $\frac{25}{128\alpha^{2}}$, an apparent misprint.}
La valeur exacte est
\[ Z^{1}=\frac{1\cdot2\ldots\alpha}{\alpha+1\cdot\alpha+2\ldots2\alpha+1}. \]
Ainsi il faut qu'on ait, lorsque $\alpha$ est très-grand,
\[ \frac{1\cdot2\cdot3\ldots\alpha}{\alpha+1\cdot\alpha+2\ldots2\alpha+1}=\left(\tfrac{1}{2}\right)^{2\alpha+1}\sqrt{\left(\frac{\pi}{\alpha}\right)}. \]
Voici comment on pourra vérifier cette formule.

Soit $\psi(m)=\frac{1\cdot2\cdot3\ldots m}{m+1\cdot m+2\ldots2m+1}$, on aura
\[ \psi(m+1)=\psi(m)\cdot\frac{(m+1)^{2}}{(2m+2)(2m+3)}=\tfrac{1}{4}\psi(m)\cdot\frac{2m+2}{2m+3}; \]
de là résulte
\[ \begin{aligned} &\psi(0)=1,\\ &\psi(1)=\tfrac{1}{4}\cdot\tfrac{2}{3},\\ &\psi(2)=\left(\tfrac{1}{4}\right)^{2}\cdot\frac{2\cdot4}{3\cdot5},\\ &\psi(3)=\left(\tfrac{1}{4}\right)^{3}\cdot\frac{2\cdot4\cdot6}{3\cdot5\cdot7},\\ &\text{etc.} \end{aligned} \]
Mais par l'expression de \emph{Wallis}, on a $\frac{\pi}{2}=\frac{2\cdot2}{1\cdot3}\cdot\frac{4\cdot4}{3\cdot5}\cdot\frac{6\cdot6}{5\cdot7}\cdot\frac{8\cdot8}{7\cdot9}$, etc. c'est-à-dire que $m$ étant un très-grand nombre, on peut supposer
\[ \frac{\pi}{2}=\frac{2^{2}\cdot4^{2}\cdot6^{2}\ldots(2m)^{2}}{3^{2}\cdot5^{2}\cdot7^{2}\ldots(2m+1)^{2}}(2m+1); \]
il en résulte $\sqrt{\left(\frac{\pi}{4m+2}\right)}=\frac{2\cdot4\cdot6\ldots2m}{3\cdot5\cdot7\ldots2m+1}$; donc $\psi(m)=\left(\frac{1}{4}\right)^{m}\sqrt{\left(\frac{\pi}{4m+2}\right)}$, ou, puisque $m$ est très-grand, $\psi(m)=\left(\frac{1}{2}\right)^{2m+1}\sqrt{\left(\frac{\pi}{m}\right)}$, ce qui s'accorde avec la formule précédente.

\origpage{351}

Notre formule est donc vérifiée en général, et on peut être assuré qu'elle donnera une grande approximation, lorsque les exposans $\alpha$ et $\beta$ seront tous deux de grands nombres; mais si l'un des deux seulement est un grand nombre, la série n'aura que le degré de convergence que comporte le développement de $\Phi(\alpha)$, $\alpha$ étant le plus petit des deux exposans. Au reste, comme la valeur de $Z$ est donnée généralement par les fonctions $\Gamma$, les observations que nous venons de faire sont plutôt relatives à la méthode générale, qu'à l'exemple particulier dont on aura toujours la solution aussi approchée qu'on voudra par les propriétés des fonctions $\Gamma$.

\subsection*{(36).}

Nous remarquerons encore que dans l'exemple II se trouve comprise la détermination générale des intégrales définies que nous avons désignées ci-dessus par la formule $\left(\frac{p}{q}\right)=\displaystyle\int x^{p-1}\,dx\,(1-x^{n})^{\frac{q-n}{n}}$. En effet, si au lieu de $x^{n}$ on met $x$, on aura $\left(\frac{p}{q}\right)=\frac{1}{n}\displaystyle\int x^{\frac{p}{n}-1}(1-x)^{\frac{q}{n}-1}\,dx$; soit encore $\frac{p}{n}=\alpha$, et $\frac{q}{n}=\beta$; on aura $\left(\frac{p}{q}\right)=\frac{1}{n}\displaystyle\int x^{\alpha-1}(1-x)^{\beta-1}\,dx$. Mais pourvu que $\alpha$ et $\beta$ soient positifs, on a entre les limites $x=0$, $x=1$,
\[ \int x^{\alpha-1}(1-x)^{\beta-1}\,dx=\frac{(\alpha+\beta)(\alpha+\beta+1)}{\alpha\beta}\int x^{\alpha}(1-x)^{\beta}\,dx; \]
et l'intégrale du second membre $=\frac{\Gamma(\alpha+1)\Gamma(\beta+1)}{(\alpha+\beta+1)\Gamma(\alpha+\beta+1)}$; donc
\[ \int x^{\alpha-1}(1-x)^{\beta-1}\,dx=\frac{(\alpha+\beta)\Gamma(\alpha+1)\Gamma(\beta+1)}{\alpha\beta\Gamma(\alpha+\beta+1)}=\frac{\Gamma(\alpha)\Gamma(\beta)}{\Gamma(\alpha+\beta)}, \]
et par conséquent
\[ \left(\frac{p}{q}\right)=\frac{\Gamma(\alpha)\Gamma(\beta)}{n\Gamma(\alpha+\beta)}=\frac{\Gamma\left(\frac{p}{n}\right)\Gamma\left(\frac{q}{n}\right)}{n\Gamma\left(\frac{p+q}{n}\right)}, \]
ce qui est la formule du n° 56, seconde partie; on a ainsi une démonstration très-simple de la formule qui sert à exprimer les intégrales $\left(\frac{p}{q}\right)$ par le moyen des fonctions $\Gamma$.

\origpage{352}

\section*{De l'intégrale $Z=\displaystyle\int x^{-\alpha}e^{-x}\,dx$, prise entre les limites imaginaires qui rendent nulle $x^{-\alpha}e^{-x}$.}

\subsection*{(37).}

La quantité $x^{-\alpha}e^{-x}$, où l'on suppose $\alpha$ positif, devient un \emph{minimum} lorsque $x=-\alpha$; mais ce \emph{minimum} se change en \emph{maximum}, si au lieu de donner à $x$ des valeurs réelles, on suppose
\[ x=-\alpha+z\sqrt{-1}, \]
et qu'on donne à $z$ des valeurs quelconques, depuis $-\infty$ jusqu'à $+\infty$. Cette supposition est admissible analytiquement, et les conséquences qu'on en déduit méritent d'être remarquées.

D'après cette valeur de $x$, l'ordonnée $x^{-\alpha}e^{-x}$ sera nulle aux deux limites de l'intégrale, savoir, lorsque $z=-\infty$ et lorsque $z=+\infty$. On pourra donc appliquer à l'intégrale $Z$ la méthode donnée dans le chapitre précédent.

Avant tout, j'observe que $\alpha$ peut être supposé plus grand que l'unité, et même aussi grand qu'on voudra; car on a $d(x^{-m}e^{-x})=-mx^{-m-1}e^{-x}\,dx-x^{-m}e^{-x}\,dx$; et par conséquent $x^{-m}e^{-x}=-m\displaystyle\int x^{-m-1}e^{-x}\,dx-\displaystyle\int x^{-m}e^{-x}\,dx$; or $m$ étant positif, la quantité $x^{-m}e^{-x}$ s'évanouit aux deux limites de l'intégrale; on a donc
\[ \int x^{-m}e^{-x}\,dx=-m\int x^{-m-1}e^{-x}\,dx; \]
d'où l'on voit que si $\alpha$ était plus petit que l'unité, on pourrait transformer la formule proposée en une autre, où $\alpha$ serait plus grand d'une unité, et ainsi de suite. La formule proposée peut donc être préparée de manière que $\alpha$ soit assez grand pour que les suites qui résultent de l'intégration soient convergentes.

\subsection*{(38).}

Cela posé, si on remarque que le calcul nécessaire pour avoir la valeur de l'intégrale $\displaystyle\int x^{-\alpha}e^{-x}\,dx$ dans le cas des limites imaginaires, est absolument le même que celui que nous avons fait pour avoir l'intégrale $\displaystyle\int x^{\alpha}e^{-x}\,dx$ dans le cas des limites réelles, on verra qu'il suffit de changer le signe de $\alpha$ dans l'intégrale déjà trouvée, afin d'avoir celle que nous cherchons. Ainsi puisque nous avons trouvé
\[ \int x^{\alpha}e^{-x}\,dx=\left(\frac{\alpha}{e}\right)^{\alpha+\frac{1}{2}}\sqrt{(2e\pi)}\cdot\Phi(\alpha), \]

\origpage{353}
$\Phi(\alpha)$ désignant la suite $1+\frac{1}{12\alpha}+\frac{1}{288\alpha^{2}}+\text{etc.}$, nous en déduirons, en changeant simplement le signe de $\alpha$,
\[ \int x^{-\alpha}e^{-x}\,dx=\left(\frac{-\alpha}{e}\right)^{-\alpha+\frac{1}{2}}\sqrt{(2e\pi)}\cdot\Phi(-\alpha). \]
Multipliant ces deux équations l'une par l'autre, et observant que par la propriété de la fonction $\Phi$, on a $\Phi(\alpha)\times\Phi(-\alpha)=1$, il viendra
\[ \int x^{-\alpha}e^{-x}\,dx\times\int x^{\alpha}e^{-x}\,dx=2\alpha\pi(-1)^{-\alpha+\frac{1}{2}}, \]
équation par laquelle on déduira généralement l'intégrale $\displaystyle\int x^{-\alpha}e^{-x}\,dx$, dont les limites sont imaginaires, de l'intégrale $\displaystyle\int x^{\alpha}e^{-x}\,dx$, dont les limites sont réelles.

Cette dernière intégrale est représentée par $\Gamma(\alpha+1)$, et on a $\Gamma(\alpha+1)=\alpha\Gamma(\alpha)$; donc si on désigne par $\Gamma'(\alpha)$ l'intégrale $\displaystyle\int x^{-\alpha}e^{-x}\,dx$ dont les limites sont imaginaires, on aura
\[ \Gamma'(\alpha)=\frac{2\pi(-1)^{\frac{1}{2}-\alpha}}{\Gamma(\alpha)}. \tag{1} \]

\subsection*{(39).}

Dans le cas particulier où $\alpha$ est un nombre entier, on a $\Gamma(\alpha)=1\cdot2\cdot3\ldots(\alpha-1)$; on aura donc dans ce même cas,
\[ \Gamma'(\alpha)=\frac{2\pi(-1)^{\alpha}\sqrt{-1}}{1\cdot2\cdot3\ldots\alpha-1}; \]
ce qui donne successivement $\Gamma'(1)=-2\pi\sqrt{-1}$, $\Gamma'(2)=2\pi\sqrt{-1}$,
\[ \Gamma'(3)=-\frac{2\pi\sqrt{-1}}{1\cdot2},\quad\Gamma'(4)=\frac{2\pi\sqrt{-1}}{1\cdot2\cdot3},\quad\Gamma'(5)=-\frac{2\pi\sqrt{-1}}{1\cdot2\cdot3\cdot4},\ \text{etc.} \]

Si on fait $\alpha=\frac{1}{2}$ (car la formule trouvée étant indépendante des suites, n'est plus assujétie aux conditions qui concernent la convergence de ces suites, et elle suppose seulement $\alpha$ positif), on aura $\Gamma'\left(\frac{1}{2}\right)=\frac{2\pi}{\Gamma\left(\frac{1}{2}\right)}=2\sqrt{\pi}$; d'où l'on voit que $\Gamma'\left(\frac{1}{2}\right)$ est réel et double de $\Gamma\left(\frac{1}{2}\right)$.

Il est remarquable que $\Gamma'\left(\frac{1}{2}\right)$ et $\Gamma\left(\frac{1}{2}\right)$ représentent toutes deux l'intégrale $\displaystyle\int x^{-\frac{1}{2}}e^{-x}\,dx$; mais la première est prise entre les limites imaginaires qui rendent nulle $x^{-\frac{1}{2}}e^{-x}$, et la seconde est prise entre

\origpage{354}
les limites réelles $x=0$, $x=\infty$; il n'est donc pas surprenant qu'elles soient inégales.

On distinguerait de même $\Gamma'\left(\frac{1}{3}\right)$ qui représente $\displaystyle\int x^{-\frac{1}{3}}e^{-x}\,dx$, prise entre des limites imaginaires, de $\Gamma\left(\frac{2}{3}\right)$ qui représente l'intégrale $\displaystyle\int x^{-\frac{1}{3}}e^{-x}\,dx$ prise entre des limites réelles. Pour les comparer entre elles, on appliquera la formule générale qui donne
\[ \Gamma'\left(\tfrac{1}{3}\right)=\frac{2\pi\sqrt{-1}}{\Gamma\left(\frac{1}{3}\right)}(-1)^{-\frac{1}{3}}; \]
ensuite, comme on a $\Gamma\left(\frac{1}{3}\right)\Gamma\left(\frac{2}{3}\right)=\frac{\pi}{\sin\frac{1}{3}\pi}$, il en résulte
\[ \frac{\Gamma'\left(\frac{1}{3}\right)}{\Gamma\left(\frac{2}{3}\right)}=2\sin\frac{\pi}{3}\cdot(-1)^{\frac{1}{3}}\sqrt{-1}=2\sqrt{-1}\sin\tfrac{1}{3}\pi\left(\cos\frac{2k+1}{3}\pi+\sqrt{-1}\sin\frac{2k+1}{3}\pi\right), \]
ce qui donne trois valeurs pour le rapport cherché de $\Gamma'\left(\frac{1}{3}\right)$ à $\Gamma\left(\frac{2}{3}\right)$. Ces trois valeurs répondent à celles que peut prendre $x^{\frac{1}{3}}$ dans la formule $\displaystyle\int x^{-\frac{1}{3}}e^{-x}\,dx$, lesquelles sont $x^{\frac{1}{3}}$, $x^{\frac{1}{3}}\left(-\frac{1}{2}+\frac{1}{2}\sqrt{-3}\right)$, $x^{\frac{1}{3}}\left(-\frac{1}{2}-\frac{1}{2}\sqrt{-3}\right)$; si on se borne à la première, on aura simplement
\[ F'\left(\tfrac{1}{3}\right)=-\sqrt{(-3)}\cdot\Gamma\left(\tfrac{2}{3}\right). \]
\ednote{Printed $F'$ (a clear capital F) where $\Gamma'$ is meant; a misprint.}

\subsection*{(40).}

Reprenons maintenant la formule $Z=\displaystyle\int x^{-\alpha}e^{-x}\,dx$; puisque nous avons fait $x=-\alpha+z\sqrt{-1}$, et que l'intégrale doit être prise entre les limites $z=-\infty$, $z=+\infty$, il s'ensuit que si on fait pour abréger $\left(-\frac{\alpha}{e}\right)^{-\alpha}=M$, on aura
\[ Z=M\int\left[\left(1-\frac{z}{\alpha}\sqrt{-1}\right)^{-\alpha}e^{-z\sqrt{-1}}+\left(1+\frac{z}{\alpha}\sqrt{-1}\right)^{-\alpha}e^{z\sqrt{-1}}\right]dz\sqrt{-1}, \]
cette intégrale étant prise depuis $z=0$ jusqu'à $z=\infty$.

Soit $z=\alpha\operatorname{tang}\varphi$, on aura la transformée
\[ Z=2M\alpha\sqrt{-1}\cdot\int d\varphi\cos^{\alpha-2}\varphi\cos(\alpha\operatorname{tang}\varphi-\alpha\varphi), \]
nouvelle intégrale qu'il faudra prendre depuis $\varphi=0$ jusqu'à $\varphi=\frac{1}{2}\pi$.

Dans les applications, on pourra supposer $\alpha>2$; ainsi l'intégrale $\displaystyle\int dx\cos^{\alpha-2}\varphi\cos(\alpha\operatorname{tang}\varphi-\alpha\varphi)$\ednote{Printed $dx$ instead of $d\varphi$; an apparent misprint.} ne présentera que des difficultés ordinaires, lorsqu'on voudra l'évaluer par la méthode des quadratures.

\origpage{355}
On connaîtra donc avec le degré d'approximation qu'on voudra, l'intégrale $Z$ ou la quantité que nous avons désignée par $\Gamma'(\alpha)$.

Réciproquement, comme la valeur de l'intégrale $\Gamma'(\alpha)$ est connue et représentée par l'expression $\frac{2\pi(-1)^{\frac{1}{2}-\alpha}}{\Gamma(\alpha)}$, si on fait
\[ \int d\varphi\cos^{\alpha-2}\varphi\cos(\alpha\operatorname{tang}\varphi-\alpha\varphi)=Q(\alpha), \]
on aura
\[ \frac{2\pi(-1)^{\frac{1}{2}-\alpha}}{\Gamma(\alpha)}=2\alpha\sqrt{-1}\cdot\left(-\frac{\alpha}{e}\right)^{-\alpha}Q(\alpha), \]
ce qui donne l'intégrale
\[ Q(\alpha)=\frac{\pi}{\Gamma(\alpha+1)}\cdot\left(\frac{\alpha}{e}\right)^{\alpha}. \tag{2} \]

\subsection*{(41).}

Cette formule étant indépendante des suites, doit avoir lieu quel que soit $\alpha$, pourvu qu'il soit positif; on connaîtra donc $Q(\alpha)$ dans tous les cas où $\Gamma(\alpha+1)$ est connu; ces cas sont ceux où $2\alpha$ est un nombre entier.

Ainsi en faisant successivement $\alpha=1$, 2, 3, 4, etc., on aura
\[ \begin{aligned} Q(1)&=\frac{\pi}{e}, &\qquad Q(2)&=\frac{\pi}{2}\cdot\left(\frac{2}{e}\right)^{2}=\frac{2\pi}{e^{2}},\\ Q(3)&=\frac{\pi}{2\cdot3}\cdot\left(\frac{3}{e}\right)^{3}=\frac{9\pi}{2e^{3}}, & Q(4)&=\frac{\pi}{2\cdot3\cdot4}\left(\frac{4}{e}\right)^{4}=\frac{32\pi}{3e^{4}},\ \text{etc.} \end{aligned} \]
De même, en faisant $\alpha=\frac{1}{2}$, $\frac{3}{2}$, $\frac{5}{2}$, etc., on aura
\[ \begin{aligned} Q\left(\frac{1}{2}\right)&=\sqrt{\left(\frac{2\pi}{e}\right)}, &\qquad Q\left(\frac{3}{2}\right)&=\frac{\sqrt{2\pi}}{1\cdot3}\left(\frac{3}{e}\right)^{\frac{3}{2}},\\ Q\left(\frac{5}{2}\right)&=\frac{\sqrt{2\pi}}{1\cdot3\cdot5}\left(\frac{5}{e}\right)^{\frac{5}{2}}, & Q\left(\frac{7}{2}\right)&=\frac{\sqrt{2\pi}}{1\cdot3\cdot5\cdot7}\left(\frac{7}{e}\right)^{\frac{7}{2}},\ \text{etc.} \end{aligned} \]
En général, si dans la valeur de $Q(\alpha)$ on substitue la valeur connue de $\Gamma(\alpha+1)$ développée en série, on aura
\[ Q(\alpha)=\left(\frac{\pi}{2\alpha}\right)^{\frac{1}{2}}\left(1-\frac{1}{12\alpha}+\frac{1}{288\alpha^{2}}-\text{etc.}\right). \tag{3} \]
Mais cette formule suppose $\alpha>1$, et elle donnera un résultat d'autant plus exact que $\alpha$ sera plus grand. Lors donc qu'on pourra

\origpage{356}
négliger $\frac{1}{12\alpha}$ par rapport à l'unité, on aura $Q(\alpha)=\sqrt{\left(\frac{\pi}{2\alpha}\right)}$, ou
\[ \int d\varphi\cos^{\alpha-2}\varphi\cos(\alpha\operatorname{tang}\varphi-\alpha\varphi)=\sqrt{\left(\frac{\pi}{2\alpha}\right)}; \]
théorème qu'il serait peut-être fort difficile de démontrer par une autre voie.

Les théorèmes particuliers concernant $Q(1)$, $Q(2)$, etc. peuvent être présentés sous une autre forme. Soit $\operatorname{tang}\varphi=z$, et supposons que les intégrales suivantes soient prises depuis $z=0$ jusqu'à $z=\infty$, on aura
\[ \begin{aligned} Q(1)&=\int\frac{dz\,(\cos z+z\sin z)}{1+zz}=\frac{\pi}{e}\\ Q(2)&=\int\frac{dz\,[(1-z^{2})\cos 2z+2z\sin 2z]}{(1+zz)^{2}}=\frac{2\pi}{e^{2}}\\ &\text{etc.} \end{aligned} \tag{4} \]
Comme ces résultats sont déduits d'une analyse fort épineuse, on doit être curieux de les vérifier, au moins dans un cas particulier. Soit, par exemple, $\alpha=2$, la formule donne $Q(2)=\frac{2\pi}{e^{2}}=0.850337$; il faut donc voir si cette valeur est celle de l'intégrale $\displaystyle\int d\varphi\cos(2\operatorname{tang}\varphi-2\varphi)$, prise depuis $\varphi=0$ jusqu'à $\varphi=\frac{1}{2}\pi$.

Si l'on fait $y=\cos(2\operatorname{tang}\varphi-2\varphi)$, l'intégrale $\displaystyle\int y\,d\varphi$ sera facile à trouver par approximation, depuis $\varphi=0$ jusqu'au premier point où l'on a $y=0$. Ce premier point a lieu lorsque $\operatorname{tang}\varphi-\varphi=\frac{1}{4}\pi$, et alors on trouve $\varphi=61^{\circ}46'40''$. Les points suivans où $y=0$, sont ceux où $\operatorname{tang}\varphi-\varphi$ a les valeurs successives $\frac{3}{4}\pi$, $\frac{5}{4}\pi$, $\frac{7}{4}\pi$, etc., et le nombre en est visiblement infini. On voit donc que depuis $\varphi=61^{\circ}46'40''$ jusqu'à $\varphi=90^{\circ}$, l'aire de la courbe est composée d'une infinité de parties alternativement positives et négatives. Ces parties sont difficiles à calculer avec un certain degré de précision; mais elles échappent bientôt par leur petitesse, et on trouve qu'en effet le résultat total s'approche beaucoup du nombre donné par la formule précédente.

Au reste on verra dans le chapitre précédent, qu'on peut par une intégration directe, vérifier les valeurs de $Q(1)$, $Q(2)$, etc. ce qui achevera de dissiper tous les doutes sur l'exactitude des formules précédentes. Nous aurons en même temps occasion de considérer

\origpage{357}
une nouvelle suite d'intégrales qui peuvent être déterminées généralement par les nombres $e$ et $\pi$.

\section*{De l'intégrale $Z=\displaystyle\int\frac{dx\cos ax}{1+xx}$, et autres semblables, prises depuis $x=0$ jusqu'à $x=\infty$.}

\subsection*{(42).}

Supposons qu'on prenne l'intégrale depuis $x=0$ jusqu'à $x=\frac{2k\pi}{a}$, $k$ étant un nombre entier, on aura en différentiant par rapport à $a$, et ayant égard à la variabilité de la seconde limite,
\[ \frac{dZ}{da}=-\int\frac{x\,dx\sin ax}{1+xx}-\frac{2k\pi}{a^{2}+4k^{2}\pi^{2}}. \]
Différentiant une seconde fois par rapport à $a$, on aura
\[ \frac{ddZ}{da^{2}}=-\int\frac{x^{2}\,dx\cos ax}{1+xx}+\frac{4ka\pi}{(a^{2}+4k^{2}\pi^{2})^{2}}; \]
d'où résulte
\[ Z-\frac{ddZ}{da^{2}}=\int dx\cos ax-\frac{4ka\pi}{(a^{2}+4k^{2}\pi^{2})^{2}}; \]
Mais on a $\displaystyle\int dx\cos ax=\frac{1}{a}\sin ax$, et cette intégrale s'évanouit à la limite supposée; donc on a simplement
\[ Z-\frac{ddZ}{da^{2}}=-\frac{4ka\pi}{(a^{2}+4k^{2}\pi^{2})^{2}}. \]
Supposons maintenant que $k$ soit un nombre très-grand par rapport à $a$, ensorte qu'on puisse négliger les quantités de l'ordre $\frac{a}{k}$; à plus forte raison pourra-t-on négliger celles de l'ordre $\frac{a}{k^{3}}$; ainsi l'équation précédente se réduit à celle-ci
\[ Z-\frac{ddZ}{da^{2}}=0, \]
et il en résulte $Z=Ae^{a}+Be^{-a}$, $A$ et $B$ étant deux constantes arbitraires.

Dans le cas de $a=0$, l'intégrale $Z=\displaystyle\int\frac{dx}{1+xx}=\operatorname{arc}\operatorname{tang}x$, et

\origpage{358}
en faisant $x=\frac{2k\pi}{a}$, on a $Z=\frac{1}{2}\pi-\frac{a}{2k\pi}$, ou simplement $Z=\frac{1}{2}\pi$; donc $A+B=\frac{1}{2}\pi$.

Je suppose maintenant qu'on donne à $a$ des valeurs de plus en plus grandes; puisque $\cos ax$ est toujours plus petit que l'unité, on aura toujours $Z<\displaystyle\int\frac{dx}{1+xx}$ ou $Z<\frac{1}{2}\pi$. Mais si $A$ n'était pas zéro, la valeur $Ae^{a}+Be^{-a}$, lorsque $a$ est devenu un nombre très-grand, se réduirait à $Ae^{a}$, quantité infiniment plus grande que $\frac{1}{2}\pi$; donc on a généralement $A=0$; donc $B=\frac{1}{2}\pi$, et enfin l'intégrale cherchée
\[ \int\frac{dx\cos ax}{1+xx}=\tfrac{1}{2}\pi e^{-a}. \]
Si dans cette formule on met $\frac{x}{m}$ au lieu de $x$, et $am$ au lieu de $a$, on aura plus généralement
\[ \int\frac{dx\cos ax}{m^{2}+x^{2}}=\frac{\pi}{2m}e^{-am}; \tag{1} \]
et cette formule étant différentiée par rapport à $a$, en donne une seconde non moins remarquable, savoir,
\[ \int\frac{x\,dx\sin ax}{m^{2}+x^{2}}=\frac{\pi}{2}e^{-am}. \tag{2} \]

\subsection*{(43).}

Si on différentie par rapport à $m$ la formule (1), et qu'on répète les différentiations, on aura successivement
\[ \begin{aligned} \int\frac{dx\cos ax}{(m^{2}+x^{2})^{2}}&=\frac{\pi e^{-am}}{2^{2}}\left(\frac{a}{m^{2}}+\frac{1}{m^{3}}\right)\\ \int\frac{dx\cos ax}{(m^{2}+x^{2})^{3}}&=\frac{1}{2}\cdot\frac{\pi e^{-am}}{2^{3}}\left(\frac{a^{2}}{m^{3}}+\frac{3a}{m^{4}}+\frac{3}{m^{5}}\right)\\ \int\frac{dx\cos ax}{(m^{2}+x^{2})^{4}}&=\frac{1}{2\cdot3}\cdot\frac{\pi e^{-am}}{2^{4}}\left(\frac{a^{3}}{m^{4}}+\frac{6a^{2}}{m^{5}}+\frac{15a}{m^{6}}+\frac{15}{m^{7}}\right).\\ &\text{etc.} \end{aligned} \]
La loi de ces expressions est facile à trouver, et si l'on fait en général,
\[ \int\frac{dx\cos ax}{(m^{2}+x^{2})^{k}}=\frac{A^{k}}{1\cdot2\cdot3\ldots k-1}\cdot\frac{\pi e^{-am}}{2^{k}}, \tag{3} \]
le coefficient $A^{k}$ aura pour valeur

\origpage{359}
\[ \begin{aligned} A^{k}&=\frac{a^{k-1}}{m^{k}}+\frac{k\cdot k-1}{2}\cdot\frac{a^{k-2}}{m^{k+1}}+\frac{k+1\cdot k\cdot k-1\cdot k-2}{2\cdot4}\cdot\frac{a^{k-3}}{m^{k+2}}\\ &\quad+\frac{k+2\cdot k+1\cdot k\cdot k-1\cdot k-2\cdot k-3}{2\cdot4\cdot6}\cdot\frac{a^{k-4}}{m^{k+3}}+\text{etc.} \end{aligned} \]
De même si on différentie successivement par rapport à $m$ l'équation (2), on en déduira cette suite de formules
\[ \begin{aligned} \int\frac{x\,dx\sin ax}{(m^{2}+x^{2})^{2}}&=\frac{\pi e^{-am}}{2^{2}}\cdot\frac{a}{m}\\ \int\frac{x\,dx\sin ax}{(m^{2}+x^{2})^{3}}&=\frac{1}{1\cdot2}\cdot\frac{\pi e^{-am}}{2^{3}}\left(\frac{a^{2}}{m^{2}}+\frac{a}{m^{3}}\right)\\ \int\frac{x\,dx\sin ax}{(m^{2}+x^{2})^{4}}&=\frac{1}{1\cdot2\cdot3}\cdot\frac{\pi e^{-am}}{2^{4}}\left(\frac{a^{3}}{m^{3}}+\frac{3a^{2}}{m^{4}}+\frac{3a}{m^{5}}\right)\\ &\text{etc.}, \end{aligned} \]
et en général,
\[ \int\frac{x\,dx\sin ax}{(m^{2}+x^{2})^{k}}=\frac{A^{k-1}}{1\cdot2\cdot3\ldots k-1}\cdot\frac{\pi a}{2^{k}}e^{-am}, \tag{4} \]
$A^{k-1}$ étant déterminé suivant la même loi que $A^{k}$.

\subsection*{(44).}

Si dans l'équation (1) on fait $m=n(\cos\theta+\sqrt{-1}\sin\theta)$, on en déduira les deux formules suivantes:
\[ \begin{aligned} \int\frac{x^{2}dx\cos ax}{x^{4}+2n^{2}x^{2}\cos2\theta+n^{4}}&=\frac{\pi}{2n\sin2\theta}e^{-an\cos\theta}\sin(\theta-an\sin\theta)\\ \int\frac{dx\cos ax}{x^{4}+2n^{2}x^{2}\cos2\theta+n^{4}}&=\frac{\pi}{2n^{3}\sin2\theta}e^{-an\cos\theta}\sin(\theta+an\sin\theta). \end{aligned} \tag{5} \]
La même substitution étant faite dans la formule (2), il en résultera ces deux autres formules:
\[ \begin{aligned} \int\frac{x^{3}dx\sin ax}{x^{4}+2n^{2}x^{2}\cos2\theta+n^{4}}&=\frac{\pi}{2\sin2\theta}e^{-an\cos\theta}\sin(2\theta-an\sin\theta)\\ \int\frac{x\,dx\sin ax}{x^{4}+2n^{2}x^{2}\cos2\theta+n^{4}}&=\frac{\pi}{2n^{2}\sin2\theta}e^{-an\cos\theta}\sin(an\sin\theta). \end{aligned} \tag{6} \]
On trouverait semblablement, au moyen des formules (3) et (4), les valeurs générales des intégrales
\[ \begin{aligned} &\int\frac{dx\cos ax}{(x^{2}+n^{2}\cos2\theta+n^{2}\sin2\theta\cdot\sqrt{-1})^{k}}\\ &\int\frac{x\,dx\sin ax}{(x^{2}+n^{2}\cos2\theta+n^{2}\sin2\theta\cdot\sqrt{-1})^{k}}. \end{aligned} \tag{7} \]

\origpage{360}

\subsection*{(45).}

Soient $M$, $N$, $P$ des fonctions rationnelles et entières de $x^{2}$, si l'on suppose que le plus haut exposant de $x$ dans $P$ est plus grand que dans $M$ et $N$, et qu'en outre $P$ n'a aucun facteur de la forme $x^{2}-m^{2}$, c'est-à-dire qu'il n'y a aucune valeur réelle de $x$ qui rende $P$ égal à zéro; alors il est visible qu'on pourra généralement trouver, au moyen des formules précédentes, l'intégrale
\[ T=\int\left(\frac{M\cos ax+Nx\sin ax}{P}\right)dx, \tag{8} \]
prise depuis $x=0$ jusqu'à $x=\infty$. L'opération à faire pour cela, est entièrement semblable à celle qu'on pratique pour l'intégration des fractions rationnelles.

Il est facile maintenant d'appliquer ces formules aux intégrales désignées (art 41) par $Q(1)$, $Q(2)$, etc., et on trouvera que les résultats s'accordent parfaitement avec ceux que nous avons déduits d'une méthode fort différente et beaucoup moins directe. En général on pourra vérifier la valeur de $Q(\alpha)$, $\alpha$ étant un nombre entier quelconque; car en faisant $\operatorname{tang}\varphi=x$, cette intégrale pourra toujours se réduire à une somme de termes de la forme $\displaystyle\int\frac{A\,dx\cos ax+Bx\,dx\sin ax}{(1+xx)^{k}}$, et sera ainsi comprise dans la formule (8).

\subsection*{(46).}

Si on multiplie par $da$ les deux membres de l'équation $\displaystyle\int\frac{dx\cos ax}{1+xx}=\tfrac{1}{2}\pi e^{-a}$, et qu'on prenne ensuite l'intégrale par rapport à $a$, depuis $a=0$, on aura
\[ \int\frac{dx\sin ax}{x(1+x^{2})}=\tfrac{1}{2}\pi(1-e^{-a}). \]
Ajoutant cette équation à l'équation $\displaystyle\int\frac{x\,dx\sin ax}{1+xx}=\tfrac{1}{2}\pi e^{-a}$, on aura ce résultat remarquable
\[ \int\frac{dx}{x}\sin ax=\tfrac{1}{2}\pi. \tag{9} \]
Il paraît d'abord étonnant que cette intégrale soit indépendante de $a$; mais si on met $\frac{x}{a}$ à la place de $x$, ce qui ne change pas les limites $0$ et $\infty$ entre lesquelles l'intégrale doit être comprise; on trouve $\displaystyle\int\frac{dx}{x}\sin ax=\int\frac{dx}{x}\sin x$; d'où il suit qu'en effet l'intégrale doit être indépendante de $a$.

\origpage{361}

On a donc aussi $\displaystyle\int\frac{dx}{x}\sin ax\cos bx=\frac{\pi}{2}$ ou zéro, selon que $a$ est plus grand ou plus petit que $b$; car dans le premier cas on a $\displaystyle\int\frac{dx}{x}\sin ax\cos bx=\tfrac{1}{2}\int\frac{dx}{x}\sin(a+b)x+\tfrac{1}{2}\int\frac{dx}{x}\sin(a-b)x=\tfrac{1}{4}\pi+\tfrac{1}{4}\pi=\tfrac{1}{2}\pi$, et dans le second on a $\displaystyle\int\frac{dx}{x}\sin ax\cos bx=\tfrac{1}{2}\int\frac{dx}{x}\sin(a+b)x-\tfrac{1}{2}\int\frac{dx}{x}\sin(b-a)x=\tfrac{1}{4}\pi-\tfrac{1}{4}\pi=0$. Quant à l'intégrale $\displaystyle\int\frac{dx}{x}\cos ax$, il est visible qu'elle est infinie.

\subsection*{(47).}

La formule (1) et les principales conséquences qui en résultent, sont dues à Laplace qui les a publiées dans le Bulletin de la Société philomatique (avril 1811). La démonstration qu'en donne cet illustre auteur est différente de celle que nous avons rapportée; et comme elle est fondée sur un principe qui peut avoir d'autres applications, nous croyons qu'on sera bien aise de la trouver ici.

Considérons la double intégrale
\[ Z=\int2y\,dy\,e^{-y^{2}(1+x^{2})}dx\cos ax, \]
et supposons qu'elle doive être prise depuis $x=0$ jusqu'à $x=\infty$, et depuis $y=0$ jusqu'à $y=\infty$. Si on intègre d'abord par rapport à $y$, on aura
\[ Z=\int\frac{dx\cos ax}{1+xx}, \]
c'est l'intégrale qu'il s'agit de trouver.

Si on revient ensuite à la double intégrale, et qu'on veuille exécuter l'intégration par rapport à $x$, il faudra chercher la valeur de $\displaystyle\int e^{-x^{2}y^{2}}dx\cos ax$. Or par une formule qui sera démontrée dans l'article suivant, on a $\displaystyle\int e^{-x^{2}}dx\cos ax=\tfrac{1}{2}\sqrt{\pi}\cdot e^{-\frac{a^{2}}{4}}$; donc en mettant $\frac{a}{y}$ au lieu de $a$, et $xy$ au lieu de $x$, on aura
\[ \int e^{-x^{2}y^{2}}dx\cos ax=\frac{\sqrt{\pi}}{2y}e^{-\frac{a^{2}}{4y^{2}}}. \]
Il reste donc à intégrer la formule
\[ Z=\sqrt{\pi}\int dy\,e^{-y^{2}-\frac{a^{2}}{4y^{2}}}. \]
Mais on démontrera ci-après (art.\ 50) qu'entre les limites $z=0$ et

\origpage{362}
$z=\infty$, on a $\displaystyle\int z^{-\frac{1}{2}}dz\,e^{-\left(\frac{1+z^{2}}{2nz}\right)}=e^{-\frac{1}{n}}\sqrt{(2n\pi)}$; donc en faisant $a=\frac{1}{n}$, et $y^{2}=\frac{z}{2n}$, on trouve $Z=\frac{1}{2}\pi e^{-a}$; donc enfin
\[ \int\frac{dx\cos ax}{1+xx}=\frac{\pi}{2}e^{-a}; \]
c'est la formule principale d'où les autres sont faciles à déduire.

\section*{De l'intégrale $\displaystyle Z=\int e^{-x^{2}}dx\cos ax$, prise depuis $x=0$ jusqu'à $x=\infty$.}

\subsection*{(48).}

Si on différentie cette formule par rapport à $a$, on aura
\[ \frac{dZ}{da}=-\int e^{-x^{2}}x\,dx\sin ax=\tfrac{1}{2}e^{-x^{2}}\sin ax-\frac{a}{2}\int e^{-x^{2}}dx\cos ax. \]
La partie hors du signe s'évanouit aux deux limites de l'intégrale; ainsi on a $\frac{dZ}{da}=-\frac{a}{2}Z$, et en intégrant $Z=Ae^{-\frac{a^{2}}{4}}$.

Pour déterminer la constante $A$, soit $a=0$, alors $\displaystyle Z=\int e^{-x^{2}}dx=\tfrac{1}{2}\sqrt{\pi}$; donc $Z$ ou
\[ \int e^{-x^{2}}dx\cos ax=\frac{\sqrt{\pi}}{2}e^{-\frac{a^{2}}{4}}. \tag{1} \]
Laplace, qui a donné cette formule dans les Mémoires de l'Institut, ann.\ 1809, pag.\ 367, la démontre de la manière suivante.

Si l'on substitue au lieu de $\cos ax$ sa valeur développée en série, on aura
\[ Z=\int e^{-x^{2}}dx\left(1-\frac{a^{2}x^{2}}{2}+\frac{a^{4}x^{4}}{2\cdot3\cdot4}-\text{etc.}\right) \]
Or en général $\displaystyle\int x^{2n}e^{-x^{2}}dx=\tfrac{1}{2}\Gamma\left(\frac{2n+1}{2}\right)$; donc on a
\[ Z=\frac{\sqrt{\pi}}{2}\left(1-\frac{a^{2}}{4}+\frac{1}{2}\cdot\frac{a^{4}}{16}-\frac{1}{2\cdot3}\cdot\frac{a^{6}}{64}+\text{etc.}\right), \]
ou
\[ Z=\frac{\sqrt{\pi}}{2}e^{-\frac{a^{2}}{4}}. \]

\subsection*{(49).}

Si on prend les différentielles successives de l'équation (1) par rapport à $a$, on en déduira cette suite d'intégrales,

\origpage{363}
\[ \begin{aligned} \int e^{-x^{2}}x\,dx\sin ax&=\frac{\sqrt{\pi}}{4}e^{-\frac{a^{2}}{4}}a\\ \int e^{-x^{2}}x^{2}dx\cos ax&=\frac{\sqrt{\pi}}{4}e^{-\frac{a^{2}}{4}}\left(1-\frac{a^{2}}{2}\right)\\ \int e^{-x^{2}}x^{3}dx\sin ax&=\frac{\sqrt{\pi}}{8}e^{-\frac{a^{2}}{4}}\left(3a-\frac{a^{3}}{2}\right)\\ \int e^{-x^{2}}x^{4}dx\cos ax&=\frac{\sqrt{\pi}}{8}e^{-\frac{a^{2}}{4}}\left(3-3a^{2}+\frac{a^{4}}{4}\right)\\ \int e^{-x^{2}}x^{5}dx\sin ax&=\frac{\sqrt{\pi}}{16}e^{-\frac{a^{2}}{4}}\left(15a-5a^{3}+\frac{a^{5}}{4}\right)\\ &\text{etc.} \end{aligned} \tag{2} \]
d'où l'on voit qu'on peut trouver en général l'intégrale
\[ \int(M\cos ax+Nx\sin ax)e^{-x^{2}}dx, \]
$M$ et $N$ étant des fonctions rationnelles et entières de $x^{2}$.

Si on proposait de trouver entre les mêmes limites l'intégrale $\displaystyle T=\int e^{-x^{2}}dx\sin ax$, on aurait d'abord $\displaystyle\frac{dT}{da}=\int e^{-x^{2}}x\,dx\cos ax=\tfrac{1}{2}-\tfrac{1}{2}e^{-x^{2}}\cos ax-\tfrac{1}{2}a\int e^{-x^{2}}dx\sin ax$. Faisant $x=\infty$ dans la la\ednote{The word ``la'' is printed twice across the line break; an apparent misprint.} partie hors du signe, il viendrait $\frac{dT}{da}=\frac{1}{2}-\frac{1}{2}aT$, d'où $\displaystyle T=\tfrac{1}{2}e^{-\frac{a^{2}}{4}}\int e^{\frac{a^{2}}{4}}da$. Cette intégrale prise depuis $a=0$, est plus simple que la proposée; mais on ne peut en trouver l'expression que par cette suite convergente
\[ T=\frac{1}{2}\left(a-\frac{a^{3}}{2\cdot3}+\frac{a^{5}}{3\cdot4\cdot5}-\frac{a^{7}}{4\cdot5\cdot6\cdot7}+\text{etc.}\right), \]
laquelle ne paraît pas susceptible d'être réduite aux transcendantes connues.

$T$ étant supposé connu, on en déduira par des différentiations successives,
\[ \begin{aligned} \int e^{-x^{2}}dx\sin ax&=T\\ \int e^{-x^{2}}x\,dx\cos ax&=\frac{dT}{da}=\frac{1}{2}-\frac{1}{2}aT\\ \int e^{-x^{2}}x^{2}dx\sin ax&=-\frac{ddT}{da^{2}}=\frac{1}{4}a+\frac{1}{2}T\left(1-\frac{a^{2}}{2}\right)\\ \int e^{-x^{2}}x^{3}dx\cos ax&=-\frac{d^{3}T}{da^{3}}=\frac{1}{2}-\frac{a^{2}}{8}-\frac{1}{4}T\left(3a-\frac{a^{3}}{2}\right)\\ \int e^{-x^{2}}x^{4}dx\sin ax&=\frac{d^{4}T}{da^{4}}=\frac{5a}{8}-\frac{a^{3}}{16}+\frac{1}{4}T\left(3-3a^{2}+\frac{a^{4}}{4}\right)\\ &\text{etc.}; \end{aligned} \]

\origpage{364}
et comme $T$ est exprimé par une suite fort convergente, on voit qu'il sera facile de trouver dans tous les cas les valeurs fort approchées de ces diverses intégrales.

Au reste, par la combinaison de ces formules, on pourrait trouver quelques résultats assez remarquables, tels que les deux suivans:
\[ \begin{aligned} \int e^{-x^{2}}dx\left(\tfrac{1}{2}a\sin ax+x\cos ax\right)&=\tfrac{1}{2}\\ \int e^{-x^{2}}dx\left(1-\tfrac{1}{2}a^{2}-2x^{2}\right)\sin ax&=-\tfrac{1}{2}a. \end{aligned} \]

\section*{De l'intégrale $\displaystyle Z(k)=\int x^{\frac{2k-1}{2}}dx\,e^{-\left(\frac{1+xx}{2nx}\right)}$, prise depuis $x=0$ jusqu'à $x=\infty$.}

\subsection*{(50).}

Considérons d'abord l'intégrale
\[ Z(0)=\int x^{-\frac{1}{2}}dx\cdot e^{-\left(\frac{1+zz}{2nx}\right)}, \]
\ednote{Printed $1+zz$ in the numerator of the exponent instead of $1+xx$; an apparent misprint.} et divisons-la en deux parties, l'une depuis $x=0$ jusqu'à $x=1$; l'autre depuis $x=1$ jusqu'à $x=\infty$. Pour avoir la première partie, soit $x=\frac{1}{z}$, on aura la transformée $\displaystyle\int-z^{\frac{3}{2}}dz\,e^{-\left(\frac{1+2z}{2nz}\right)}$\ednote{Printed $z^{\frac{3}{2}}$ instead of $z^{-\frac{3}{2}}$, and $1+2z$ in the numerator of the exponent where $1+zz$ is meant (the printed Errata, p.\ 387, corrects it to $1+zz$); apparent misprints.}, qu'il faudra intégrer depuis $z=\infty$ jusqu'à $z=1$; si on change son signe, il faudra l'intégrer depuis $z=1$ jusqu'à $z=\infty$; ainsi en réunissant les deux parties, on voit que tout se réduit à trouver l'intégrale
\[ \int\left(x^{-\frac{1}{2}}+x^{-\frac{3}{2}}\right)dx\cdot e^{-\left(\frac{1+xx}{2nx}\right)} \]
entre les limites $x=1$ et $x=\infty$.

Soit $1+x^{2}=2xz$, on aura $x-1=x^{\frac{1}{2}}\sqrt{(2z-2)}$, et par conséquent
\[ \left(x^{-\frac{1}{2}}+x^{-\frac{3}{2}}\right)dx=\frac{dz\sqrt{2}}{\sqrt{(z-1)}}. \]
Donc l'intégrale dont il s'agit aura pour transformée
\[ Z(0)=\sqrt{2}\int\frac{dz}{\sqrt{(z-1)}}e^{-\frac{z}{n}}, \]
celle-ci devant être prise depuis $z=1$ jusqu'à $z=\infty$.

Soit $z=1+y^{2}$, et on aura une nouvelle transformée
\[ Z(0)=2^{\frac{3}{2}}e^{-\frac{1}{n}}\int dy\,e^{-\frac{y^{2}}{n}}, \]

\origpage{365}
qui devra être intégrée depuis $y=0$ jusqu'à $y=\infty$. Mais en faisant, suivant notre usage, $e^{-\frac{y^{2}}{n}}=u$, on aura
\[ Z(0)=(2n)^{\frac{1}{2}}e^{-\frac{1}{n}}\int du\left(l\frac{1}{u}\right)^{-\frac{1}{2}}; \]
et comme cette intégrale doit être prise depuis $u=0$ jusqu'à $u=1$, on aura enfin
\[ Z(0)=e^{-\frac{1}{n}}\sqrt{(2n\pi)}. \tag{1} \]

\subsection*{(51).}

Considérons en second lieu la formule
\[ Z(1)=\int x^{\frac{1}{2}}dx\,e^{-\left(\frac{1+xx}{2nx}\right)}; \]
si on la divise en deux parties, comme dans l'article précédent, on aura
\[ Z(1)=\int\left(x^{\frac{1}{2}}+x^{-\frac{5}{2}}\right)dx\,e^{-\left(\frac{1+xx}{2nx}\right)}, \]
qu'il faudra intégrer depuis $x=1$ jusqu'à $x=\infty$. Or en faisant $1+x^{2}=2xz$, puis $z=1+y^{2}$, on aura successivement
\[ \begin{aligned} x^{\frac{3}{2}}&=(2z-1)\sqrt{\left(\frac{z+1}{2}\right)}+(2z+1)\sqrt{\left(\frac{z-1}{2}\right)}\\ x^{-\frac{3}{2}}&=(2z-1)\sqrt{\left(\frac{z+1}{2}\right)}-(2z+1)\sqrt{\left(\frac{z-1}{2}\right)}\\ x^{\frac{3}{2}}-x^{-\frac{3}{2}}&=(2z+1)\sqrt{(2z-2)}=(2y^{3}+3y)\sqrt{2}\\ \left(x^{\frac{1}{2}}+x^{-\frac{5}{2}}\right)dx&=2^{\frac{3}{2}}(2y^{2}+1)\,dy. \end{aligned} \]
Donc la transformée en $y$ sera
\[ Z(1)=2^{\frac{3}{2}}e^{-\frac{1}{n}}\int(1+2y^{2})e^{-\frac{y^{2}}{n}}dy. \]
Soit encore $e^{-\frac{y^{2}}{n}}=u$, et on aura la nouvelle transformée
\[ Z(1)=(2n)^{\frac{1}{2}}e^{-\frac{1}{n}}\int du\left[\left(l\frac{1}{u}\right)^{-\frac{1}{2}}+2n\left(l\frac{1}{u}\right)^{\frac{1}{2}}\right]; \]
d'où résulte en intégrant,
\[ Z(1)=(2n)^{\frac{1}{2}}e^{-\frac{1}{n}}\left[\Gamma(\tfrac{1}{2})+2n\Gamma(\tfrac{3}{2})\right]. \]
Substituant les valeurs connues $\Gamma(\frac{1}{2})=\sqrt{\pi}$, $\Gamma(\frac{3}{2})=\frac{1}{2}\sqrt{\pi}$, on

\origpage{366}
aura enfin
\[ Z(1)=(1+n)e^{-\frac{1}{n}}\sqrt{(2n\pi)}=(1+n)Z(0). \tag{2} \]

\subsection*{(52).}

On pourrait calculer de la même manière les valeurs de $Z(2)$, $Z(3)$, etc. Mais la loi de progression de ces quantités est facile à trouver; en effet, si on différentie la quantité\ldots\ $T=x^{\frac{2k+1}{2}}e^{-\left(\frac{1+xx}{2nx}\right)}$, on aura
\[ dT=\frac{2k+1}{2}x^{\frac{2k-1}{2}}dx\cdot e^{-\left(\frac{1+xx}{2nx}\right)}+x^{\frac{2k+1}{2}}\left(\frac{1}{2nx^{2}}-\frac{1}{2n}\right)dx\cdot e^{-\left(\frac{1+xx}{2nx}\right)}. \]
Intégrant et observant que $T$ s'évanouit aux deux limites de l'intégrale, on trouve
\[ 0=\frac{2k+1}{2}Z(k)+\frac{1}{2n}Z(k-1)-\frac{1}{2n}Z(k+1); \]
d'où résulte
\[ Z(k+1)=(2k+1)nZ(k)+Z(k-1). \tag{3} \]
Ainsi on aura successivement
\[ \begin{aligned} Z(2)&=3nZ(1)+Z(0)=Z(0)(1+3n+3n^{2})\\ Z(3)&=5nZ(2)+Z(1)=Z(0)(1+6n+15n^{2}+15n^{3})\\ &\text{etc.} \end{aligned} \]
L'expression générale de ces quantités est
\[ \begin{aligned} Z(k)=Z(0)\cdot\Big[1&+\frac{k+1\cdot k}{2}n+\frac{k+2\cdot k+1\cdot k\cdot k-1}{2\cdot4}n^{2}\\ &+\frac{k+3\cdot k+2\cdot k+1\cdot k\cdot k-1\cdot k-2}{2\cdot4\cdot6}n^{3}+\text{etc.}\Big]; \end{aligned} \tag{4} \]
ce qu'on peut aisément vérifier par la substitution dans l'équation (3). La série que nous rencontrons ici suit donc la même loi que celle qui a été représentée ci-dessus par $A^{k}$ (art.\ 43).

Cette formule a également lieu lorsque $k$ est négatif, et elle donne immédiatement $Z(-k-1)=Z(k)$; c'est aussi ce qui résulte de la considération directe des intégrales. En effet on a généralement
\[ Z(k)=\int\left(x^{\frac{2k-1}{2}}+x^{-\left(\frac{2k+3}{2}\right)}\right)dx\,e^{-\left(\frac{1+xx}{2nx}\right)}, \]
cette intégrale étant prise depuis $x=1$ jusqu'à $x=\infty$.

\origpage{367}

On aurait semblablement
\[ Z(-k-1)=\int\left(x^{-\left(\frac{2k+3}{2}\right)}+x^{\frac{2k-1}{2}}\right)dx\,e^{-\left(\frac{1+xx}{2nx}\right)}; \]
donc $Z(-k-1)=Z(k)$.

\subsection*{(53).}

Euler, dans le tome IV des \emph{Supplémens au Calcul intégral}, pag.\ 415, fait mention des intégrales
\[ A=\frac{1}{2n}\int x^{-\frac{5}{2}}dx\,e^{-\left(\frac{1+xx}{2nx}\right)},\qquad B=\frac{1}{2n}\int x^{-\frac{3}{2}}dx\,e^{-\left(\frac{1+xx}{2nx}\right)}, \]
qui lui semblent ne pouvoir être ramenées aux méthodes connues. Cette difficulté est résolue par les formules précédentes qui donnent
\[ \begin{aligned} A&=\frac{1}{2n}Z(-2)=\frac{1}{2n}Z(1)=(1+n)e^{-\frac{1}{n}}\sqrt{\left(\frac{\pi}{2n}\right)}\\ B&=\frac{1}{2n}Z(-1)=\frac{1}{2n}Z(0)=e^{-\frac{1}{n}}\sqrt{\left(\frac{\pi}{2n}\right)}, \end{aligned} \]
d'où $\frac{A}{B}=1+n$. Le même auteur ajoute que si on ne peut pas trouver séparément les valeurs de ces deux intégrales, on connaît au moins leur rapport, savoir,
\[ \frac{A}{B}=\frac{1+e^{\frac{2}{n}}}{1-e^{\frac{2}{n}}}. \]
Mais il y a évidemment erreur dans les calculs qui ont conduit à ce résultat, puisqu'il donnerait une valeur négative de $\frac{A}{B}$, tandis que $A$ et $B$ sont positifs.

\section*{Des intégrales $\displaystyle\int x^{a-1}e^{-mx}dx\cos nx$, $\displaystyle\int x^{a-1}e^{-mx}dx\sin nx$, prises entre les limites $x=0$, $x=\infty$.}

\subsection*{(54).}

Nous supposerons que $a$ et $m$ sont positifs, condition nécessaire pour que les intégrales dont il s'agit soient des quantités finies. Pour en trouver les valeurs, considérons d'abord la formule
\[ Z=\int x^{a-1}dx\,e^{-(m-n\sqrt{-1})x}; \]

\origpage{368}
si on fait $e^{-(m-n\sqrt{-1})x}=z$, on aura la transformée
\[ Z=\int\frac{dz\left(l\frac{1}{z}\right)^{a-1}}{(m-n\sqrt{-1})^{a}}, \]
laquelle devra être intégrée depuis $z=0$ jusqu'à $z=1$. Soit donc $\frac{n}{m}=\operatorname{tang}\theta$ et $\sqrt{(m^{2}+n^{2})}=r$, on aura
\[ Z=\frac{\cos a\theta+\sqrt{-1}\sin a\theta}{r^{a}}\Gamma(a). \]
Il suffit maintenant de substituer dans les formules proposées les valeurs de $\cos nx$ et $\sin nx$ en exponentielles imaginaires, et on en conclura immédiatement
\[ \begin{aligned} \int x^{a-1}e^{-mx}dx\cos nx&=\frac{\cos a\theta}{r^{a}}\Gamma(a)\\ \int x^{a-1}e^{-mx}dx\sin nx&=\frac{\sin a\theta}{r^{a}}\Gamma(a). \end{aligned} \tag{1} \]
Au moyen de ces formules, on trouvera les valeurs des intégrales $\displaystyle\int x^{a-1}e^{-mx}dx\sin^{k}nx$, $\displaystyle\int x^{a-1}e^{-mx}dx\cos^{k}nx$, et en général celle de l'intégrale $\displaystyle\int Tx^{a-1}e^{-mx}dx$, $T$ étant une fonction de sinus et cosinus qui peut se développer en une suite finie de sinus et cosinus linéaires de la forme $A\sin(\alpha x+\beta)+$ etc.

Si le nombre $a$ est entier, les intégrales (1) pourront se trouver par les procédés ordinaires de l'intégration, et on aura d'ailleurs $\Gamma(a)=1\cdot2\cdot3\ldots a-1$, ce qui permettra de vérifier ces formules.

\subsection*{(55).}

Si l'on fait $m=0$ ou $\theta=\frac{1}{2}\pi$, les formules (1) se réduiront aux deux suivantes:
\[ \begin{aligned} \int x^{a-1}dx\cos nx&=\frac{\cos\frac{1}{2}a\pi}{n^{a}}\Gamma(a)\\ \int x^{a-1}dx\sin nx&=\frac{\sin\frac{1}{2}a\pi}{n^{a}}\Gamma(a); \end{aligned} \tag{2} \]
et en particulier lorsque $a=\frac{1}{2}$, on aura
\[ \begin{aligned} \int\frac{dx\cos nx}{\sqrt{x}}&=\sqrt{\left(\frac{\pi}{2n}\right)}\\ \int\frac{dx\sin nx}{\sqrt{x}}&=\sqrt{\left(\frac{\pi}{2n}\right)}. \end{aligned} \tag{3} \]
Si dans les formules (1) on suppose $a$ infiniment petit, ce qui

\origpage{369}
donne $\Gamma(a)=\frac{1}{a}$, on aura
\[ \begin{aligned} \int\frac{dx}{x}e^{-mx}\cos nx&=\infty,\\ \int\frac{dx}{x}e^{-mx}\sin nx&=\theta=\operatorname{arc}\operatorname{tang}\frac{n}{m}. \end{aligned} \tag{4} \]

Si dans cette dernière on fait $m=0$, on retrouve la formule du n° 46.

Les intégrales dont nous venons de nous occuper se trouvent dans le quatrième volume du Calcul intégral d'Euler, pag.\ 337 et suiv.\ Elles ont été traitées aussi par Laplace, dans le tome VIII du Journal de l'École Polytechnique, pag.\ 244 et suiv.

\subsection*{(56).}

Si on différentie par rapport à $a$ les formules (1), on en déduira
\[ \begin{aligned} \int x^{a-1}e^{-mx}dx\cos nx\log x&=\frac{\cos a\theta}{r^{a}}\cdot\frac{d\Gamma(a)}{da}-\left(\frac{\theta\sin a\theta+\cos a\theta\log r}{r^{a}}\right)\Gamma(a)\\ \int x^{a-1}e^{-mx}dx\sin nx\log x&=\frac{\sin a\theta}{r^{a}}\cdot\frac{d\Gamma(a)}{da}+\left(\frac{\theta\cos a\theta-\sin a\theta\log r}{r^{a}}\right)\Gamma(a). \end{aligned} \tag{5} \]
La nouvelle transcendante $\frac{d\Gamma(a)}{da}$ qui entre dans ces formules, peut se trouver d'une manière approchée par les tables; on a aussi sa valeur en séries convergentes par les formules du n° 76, II\textsuperscript{e} partie.

Si de ces deux dernières équations on élimine $\frac{d\Gamma(a)}{da}$, on en tire ce résultat remarquable:
\[ \int x^{a-1}e^{-mx}dx\sin(a\theta-nx)\log\frac{1}{x}=\frac{\theta}{r^{a}}\Gamma(a), \tag{6} \]
et en particulier lorsque $a=1$,
\[ \int e^{-mx}dx\sin(\theta-nx)\log\frac{1}{x}=\frac{\theta}{r}; \tag{7} \]
formule où l'on a $\operatorname{tang}\theta=\frac{n}{m}$ et $r=\sqrt{(m^{2}+n^{2})}$.

\origpage{370}

\section*{De l'intégrale $\displaystyle\int\frac{(x^{n}-1)\,dx}{\log x}$ et autres semblables, prises depuis $x=0$ jusqu'à $x=1$.}

\subsection*{(57).}

L'intégrale $\displaystyle\int\frac{dx}{\log x}$ prise entre ces limites, est infinie; l'intégrale $\displaystyle\int\frac{x^{n}dx}{\log x}$ se réduit à la même forme $\displaystyle\int\frac{du}{\log u}$, en faisant $x^{n+1}=u$; elle est donc aussi infinie; mais la différence de ces infinis est une quantité finie qu'il s'agit de déterminer.

Pour cela, soit $\displaystyle P=\int\frac{(x^{n}-1)\,dx}{lx}$; si on différentie cette équation par rapport à $n$, on aura
\[ \frac{dP}{dn}=\int x^{n}dx=\frac{1}{n+1}; \]
donc $dP=\frac{dn}{n+1}$, et par conséquent $P=\log(n+1)$, sans constante, parce que $P$ doit s'évanouir lorsque $n=0$. On a donc entre les limites données,
\[ \int(x^{n}-1)\frac{dx}{lx}=\log(n+1). \tag{1} \]
De là il est facile de trouver, entre les mêmes limites, l'intégrale $\displaystyle\int\frac{(x^{n}-1)x^{m}dx}{lx}$. Soit pour cet effet $x^{m+1}=u$, et $\frac{n}{m+1}=\alpha$, on aura
\[ \int\frac{(x^{n}-1)x^{m}dx}{lx}=\int\frac{(u^{\alpha}-1)\,du}{lu}=\log(\alpha+1)=\log\left(\frac{m+n+1}{m+1}\right); \]
donc
\[ \int\frac{(x^{n}-1)x^{m}dx}{lx}=\log\left(\frac{m+n+1}{m+1}\right). \tag{2} \]
On pourrait parvenir immédiatement à ce résultat en observant que $(x^{n}-1)x^{m}=x^{m+n}-1-(x^{m}-1)$; ce qui donne
\[ \int\frac{(x^{n}-1)x^{m}dx}{lx}=\int(x^{m+n}-1)\frac{dx}{lx}-\int(x^{m}-1)\frac{dx}{lx}=\log\left(\frac{m+n+1}{m+1}\right). \]
En général, si on a un polynome
\[ X=Ax^{n}+Bx^{n-1}+Cx^{n-2}+\text{etc.}, \]
qui se réduise à zéro lorsque $x=1$, on pourra le mettre sous la forme
\[ X=A(x^{n}-1)+B(x^{n-1}-1)+C(x^{n-2}-1)+\text{etc.}; \]

\origpage{371}
et alors il est clair qu'on aura
\[ \int\frac{X\,dx}{lx}=A\log(n+1)+B\log(n)+C\log(n-1)+\text{etc.} \]

\subsection*{(58).}

Supposons maintenant qu'on veuille avoir l'intégrale
\[ Q=\int\frac{(x^{n}-1)^{2}dx}{(lx)^{2}}; \]
en différentiant par rapport à $n$, on aura
\[ \frac{dQ}{dn}=\int\frac{2(x^{n}-1)x^{n}dx}{lx}=2\log(2n+1)-2\log(n+1); \]
d'où résulte en intégrant,
\[ Q=(2n+1)\log(2n+1)-2(n+1)\log(n+1). \tag{3} \]
Soit proposée plus généralement l'intégrale $\displaystyle\int\frac{(x^{n}-1)^{2}x^{m}dx}{(lx)^{2}}$, on fera $x^{m+1}=u$, $\frac{n}{m+1}=\alpha$, et la transformée sera $\displaystyle(m+1)\int\frac{(u^{\alpha}-1)^{2}du}{lu}$.\ednote{Printed $lu$ in the denominator; the substitution gives $(lu)^{2}$, as the value that follows requires. An apparent misprint.} Celle-ci, d'après la formule précédente, a pour valeur,
\[ (m+1)(2\alpha+1)\log(2\alpha+1)-2(m+1)(\alpha+1)\log(\alpha+1). \]
Donc en remettant la valeur de $\alpha$, on aura
\[ \begin{aligned} \int\frac{(x^{n}-1)^{2}x^{m}dx}{(lx)^{2}}={}&(2n+m+1)\log(2n+m+1)-2(n+m+1)\log(n+m+1)\\ &+(m+1)\log(m+1), \end{aligned} \]
intégrale qui pourrait être représentée plus simplement par la formule
\[ \int\left(\frac{x^{n}-1}{lx}\right)^{2}x^{m}dx=\frac{1}{2}\Delta^{3}[(m+1)^{2}\log(m+1)], \tag{4} \]
\ednote{The right side of (4) is printed as read here, the same form as in formula (5); by the general formulae of art.\ 60 it should be $\Delta^{2}[(m+1)\log(m+1)]$. An apparent misprint. A reader has struck the $\frac{1}{2}$ and the exponents in pencil, so the printed exponents under the pencil are read with some doubt.}
en supposant que la différence $\Delta m$ soit égale à $n$.

\subsection*{(59).}

On trouvera semblablement
\[ \int\left(\frac{x^{n}-1}{lx}\right)^{3}x^{m}dx=\frac{1}{2}\Delta^{3}[(m+1)^{2}\log(m+1)]. \]
En effet, si on fait le premier membre $=R$, on aura
\[ \frac{dR}{dn}=3\int\left(\frac{x^{n}-1}{lx}\right)^{2}x^{m+n}\,dx, \]
ou
\[ \begin{aligned} \frac{dR}{dn}={}&3(3n+m+1)\log(3n+m+1)-6(2n+m+1)\log(2n+m+1)\\ &+3(n+m+1)\log(n+m+1). \end{aligned} \]

\origpage{372}
Multipliant par $dn$, et intégrant le second membre d'après la formule $\displaystyle\int x\,dx\log x=\frac{x^{2}}{2}(\log x-\tfrac{1}{2})$; déterminant ensuite la constante de manière que l'intégrale soit nulle lorsque $n=0$, on aura
\[ \begin{aligned} R={}&\tfrac{1}{2}(3n+m+1)^{2}\log(3n+m+1)-\tfrac{3}{2}(2n+m+1)^{2}\log(2n+m+1)\\ &+\tfrac{3}{2}(n+m+1)^{2}\log(n+m+1)-\tfrac{1}{2}(m+1)^{2}\log(m+1), \end{aligned} \]
quantité qu'on peut mettre sous la forme $\frac{1}{2}\Delta^{3}[(m+1)^{2}\log(m+1)]$; on aura donc
\[ \int\left(\frac{x^{n}-1}{lx}\right)^{3}x^{m}dx=\frac{1}{2}\Delta^{3}[(m+1)^{2}\log(m+1)]. \tag{5} \]

\subsection*{(60).}

Ces formules peuvent être présentées de la manière la plus simple et la plus générale, comme il suit:
\[ \begin{aligned} \int\left(\frac{x^{n}-1}{lx}\right)x^{m-1}dx&=\Delta(\log m)\\ \int\left(\frac{x^{n}-1}{lx}\right)^{2}x^{m-1}dx&=\Delta^{2}(m\log m)\\ \int\left(\frac{x^{n}-1}{lx}\right)^{3}x^{m-1}dx&=\frac{1}{2}\Delta^{3}(m^{2}\log m)\\ \int\left(\frac{x^{n}-1}{lx}\right)^{4}x^{m-1}dx&=\frac{1}{2\cdot3}\Delta^{4}(m^{3}\log m)\\ \text{etc.}\qquad& \end{aligned} \]
les différences indiquées par $\Delta$ étant prises en supposant $\Delta m=n$.

Euler a considéré ces intégrales dans quelques cas seulement; voyez le tome IV de son \emph{Calcul intégral}, pag.~271 et suiv.

\section*{Des intégrales $\displaystyle\int\frac{d\varphi\cos\lambda\varphi}{\Delta^{n}}$ et $\displaystyle\int\Delta^{n}d\varphi\cos\lambda\varphi$, prises depuis $\varphi=0$ jusqu'à $\varphi=\pi$, en supposant les nombres $\lambda$ et $n$ entiers, et $\Delta=1+a^{2}-2a\cos\varphi$.}

\subsection*{(61).}

D'après la théorie des suites récurrentes, il est facile de voir qu'on a l'équation
\[ \frac{1-a\cos\varphi}{1-2a\cos\varphi+a^{2}}=1+a\cos\varphi+a^{2}\cos2\varphi+a^{3}\cos3\varphi+\text{etc.} \]
Multipliant chaque membre par 2 et retranchant de part et d'autre l'unité, il viendra
\[ \frac{1-a^{2}}{1-2a\cos\varphi+a^{2}}=1+2a\cos\varphi+2a^{2}\cos2\varphi+2a^{3}\cos3\varphi+\text{etc.} \]

\origpage{373}
donc
\[ \frac{d\varphi\cos\lambda\varphi}{\Delta}=\frac{d\varphi\cos\lambda\varphi}{1-a^{2}}(1+2a\cos\varphi+2a^{2}\cos2\varphi+2a^{3}\cos3\varphi+\text{etc.}). \]
Appelons en général $P^{n}$ l'intégrale $\displaystyle\int\frac{d\varphi\cos\lambda\varphi}{\Delta^{n}}$, prise entre les limites $\varphi=0$, $\varphi=\pi$; il faudra pour avoir la valeur de $P^{1}$, intégrer le second membre de l'équation précédente. Or il est visible que pour cet objet, on pourra réduire la suite $1+2a\cos\varphi+2a^{2}\cos2\varphi+\text{etc.}$, au seul terme $2a^{\lambda}\cos\lambda\varphi$; car $k$ étant différent de $\lambda$, on a\ldots\ $\displaystyle\int d\varphi\cos\lambda\varphi\cos k\varphi=\frac{1}{2}\int d\varphi\cos(\lambda-k)\varphi+\frac{1}{2}\int d\varphi\cos(\lambda+k)\varphi=\frac{\sin(\lambda-k)\varphi}{2(\lambda-k)}+\frac{\sin(\lambda+k)\varphi}{2(\lambda+k)}$, intégrale qui s'évanouit dans les limites données, puisque $\lambda$ et $k$ sont des nombres entiers différens l'un de l'autre. Nous aurons donc simplement
\[ P^{1}=\int\frac{d\varphi\cos\lambda\varphi}{1-a^{2}}\cdot2a^{\lambda}\cos\lambda\varphi=\frac{a^{\lambda}}{1-a^{2}}\int d\varphi(1+\cos2\lambda\varphi). \]
Mais l'intégrale $\displaystyle\int d\varphi(1+\cos2\lambda\varphi)$, prise entre les limites $\varphi=0$, $\varphi=\pi$, se réduit à $\pi$; donc enfin
\[ P^{1}=\frac{\pi a^{\lambda}}{1-a^{2}}. \tag{1} \]
Nous avons supposé tacitement que $a$ était plus petit que l'unité; s'il était plus grand, on ferait $a=\frac{1}{\alpha}$, et alors $\Delta$ deviendrait $\frac{1}{\alpha^{2}}(1+\alpha^{2}-2\alpha\cos\varphi)$, de sorte que l'intégrale $\displaystyle\int\frac{d\varphi\cos\lambda\varphi}{\Delta^{n}}$ se ramènerait toujours à une intégrale semblable où l'on aurait $a<1$.

\subsection*{(62).}

Connaissant ainsi la valeur de $P^{1}$, on peut en déduire par des différentiations successives, les valeurs de $P^{2}$, $P^{3}$, etc. En effet si on différentie par rapport à $a$ l'équation $\displaystyle P^{n}=\int\frac{d\varphi\cos\lambda\varphi}{\Delta^{n}}$, on en tire
\[ \frac{dP^{n}}{da}=\int\frac{n\,d\varphi\cos\lambda\varphi}{\Delta^{n+1}}\left(\frac{1-a^{2}-\Delta}{a}\right), \]
ou
\[ \frac{dP^{n}}{da}=\frac{n}{a}(1-a^{2})P^{n+1}-\frac{n}{a}P^{n}; \]
donc
\[ P^{n+1}=\frac{1}{1-a^{2}}\cdot\left(P^{n}+\frac{n}{a}\cdot\frac{dP^{n}}{da}\right)=\frac{d(a^{n}P^{n})}{(1-a^{2})\,d(a^{n})}. \tag{2} \]
\ednote{Printed $\frac{n}{a}\cdot\frac{dP^{n}}{da}$; the preceding equation and the last member of (2) require $\frac{a}{n}\cdot\frac{dP^{n}}{da}$. An apparent misprint, corrected in pencil by a reader.}
Au moyen de cette équation et de la valeur connue de $P^{1}$, on trouvera successivement

\origpage{374}
\[ \begin{aligned} P^{2}&=\frac{\pi a^{\lambda}}{(1-a^{2})^{3}}\{\lambda+1-a^{2}(\lambda-1)\}\\ P^{3}&=\frac{1}{1\cdot2}\cdot\frac{\pi a^{\lambda}}{(1-a^{2})^{5}}\{(\lambda+1)(\lambda+2)-2a^{2}(\lambda+2)(\lambda-2)+a^{4}(\lambda-2)(\lambda-1)\}\\ P^{4}&=\frac{1}{1\cdot2\cdot3}\cdot\frac{\pi a^{\lambda}}{(1-a^{2})^{7}}\{(\lambda+1)(\lambda+2)(\lambda+3)-3a^{2}(\lambda+2)(\lambda+3)(\lambda-3)\\ &\qquad+3a^{4}(\lambda+3)(\lambda-3)(\lambda-2)-a^{6}(\lambda-3)(\lambda-2)(\lambda-1)\}\\ P^{5}&=\frac{1}{1\cdot2\cdot3\cdot4}\cdot\frac{\pi a^{\lambda}}{(1-a^{2})^{9}}\{(\lambda+1)(\lambda+2)(\lambda+3)(\lambda+4)-4a^{2}(\lambda+2)(\lambda+3)(\lambda+4)(\lambda-4)\\ &\qquad+6a^{4}(\lambda+3)(\lambda+4)(\lambda-4)(\lambda-3)-4a^{6}(\lambda+4)(\lambda-4)(\lambda-3)(\lambda-2)\\ &\qquad+a^{8}(\lambda-4)(\lambda-3)(\lambda-2)(\lambda-1)\}\\ &\text{etc.} \end{aligned} \]
La loi de ces expressions est facile à saisir, et en général si on fait
\[ \left.\begin{aligned} P^{n}&=\frac{\pi a^{\lambda}}{(1-a^{2})^{2n-1}}\cdot\frac{\lambda+1\cdot\lambda+2\cdot\lambda+3\ldots\lambda+n-1}{1\cdot2\cdot3\ldots\ldots n-1}A^{n},\\ &\text{le coefficient }A^{n}\text{ aura pour valeur}\\ A^{n}&=1+\frac{n-1}{1}a^{2}\cdot\frac{n-\lambda-1}{\lambda+1}+\frac{n-1\cdot n-2}{1\cdot2}a^{4}\cdot\frac{n-\lambda-1\cdot n-\lambda-2}{\lambda+1\cdot\lambda+2}\\ &\qquad+\frac{n-1\cdot n-2\cdot n-3}{1\cdot2\cdot3}a^{6}\cdot\frac{n-\lambda-1\cdot n-\lambda-2\cdot n-\lambda-3}{\lambda+1\cdot\lambda+2\cdot\lambda+3}+\text{etc.} \end{aligned}\right\} \tag{3} \]

Pour s'assurer de l'exactitude de cette formule, il suffit de substituer la valeur générale de $P^{n}$ dans le second membre de l'équation (2), on trouvera, après les réductions, une valeur de $P^{n+1}$ qui ne sera autre chose que ce que devient $P^{n}$ en mettant $n+1$ à la place de $n$. Tout autre mode de vérification serait moins facile que celui que nous indiquons.

Le cas le plus simple est celui où $n=\lambda+1$; alors $A^{n}$ se réduit à son premier terme 1, et on a
\[ P^{\lambda+1}=\int\frac{d\varphi\cos\lambda\varphi}{\Delta^{\lambda+1}}=\frac{\pi a^{\lambda}}{(1-a^{2})^{2\lambda+1}}\cdot\frac{\lambda+1\cdot\lambda+2\cdot\lambda+3\ldots2\lambda}{1\cdot2\cdot3\ldots\lambda}. \tag{4} \]
Un autre cas qui mérite d'être remarqué, est celui de $\lambda=0$; alors on a
\[ \begin{aligned} \int\frac{d\varphi}{\Delta^{n}}=\frac{\pi}{(1-a^{2})^{2n-1}}\Big[&1+\left(\frac{n-1}{1}\right)^{2}a^{2}+\left(\frac{n-1\cdot n-2}{1\cdot2}\right)^{2}a^{4}\\ &+\left(\frac{n-1\cdot n-2\cdot n-3}{1\cdot2\cdot3}\right)^{2}a^{6}+\text{etc.}\Big], \end{aligned} \tag{5} \]
formule dont les coefficiens sont les quarrés des coefficiens du binome élevé à la puissance $n-1$. On aurait, par exemple,

\origpage{375}
\[ \int\frac{d\varphi}{\Delta^{5}}=\frac{\pi}{(1-a^{2})^{9}}[1+4^{2}a^{2}+6^{2}a^{4}+4^{2}a^{6}+a^{8}]. \]

\subsection*{(63).}

Venons maintenant à l'intégrale $\displaystyle\int\Delta^{n}d\varphi\cos\lambda\varphi$ que nous désignerons par $Q^{n}$. Sans passer par les cas particuliers, on peut déterminer tout d'un coup la valeur générale de $Q^{n}$. En effet, si on décompose le polynome $1-2a\cos\varphi+a^{2}$ en ses deux facteurs imaginaires, on aura
\[ \Delta=(1-ae^{\varphi\sqrt{-1}})(1-ae^{-\varphi\sqrt{-1}}). \]
Si ensuite on suppose
\[ \begin{aligned} (1-ae^{\varphi\sqrt{-1}})^{n}={}&1+K_{1}e^{\varphi\sqrt{-1}}+K_{2}e^{2\varphi\sqrt{-1}}\ldots+K_{\lambda}e^{\lambda\varphi\sqrt{-1}}\\ &+K_{\lambda+1}e^{(\lambda+1)\varphi\sqrt{-1}}+\text{etc.} \end{aligned} \]
on aura semblablement
\[ \begin{aligned} (1-ae^{-\varphi\sqrt{-1}})^{n}={}&1+K_{1}e^{-\varphi\sqrt{-1}}+K_{2}e^{-2\varphi\sqrt{-1}}\ldots+K_{\lambda}e^{-\lambda\varphi\sqrt{-1}}\\ &+K_{\lambda+1}e^{-(\lambda+1)\varphi\sqrt{-1}}+\text{etc.} \end{aligned} \]
Multipliant ces deux équations l'une par l'autre, afin d'avoir la valeur de $\Delta^{n}$, et observant qu'il suffit de conserver dans le second membre les termes affectés de $e^{\lambda\varphi\sqrt{-1}}$ et de $e^{-\lambda\varphi\sqrt{-1}}$, on aura le produit cherché
\[ \Delta^{n}=(K_{\lambda}+K_{\lambda+1}K_{1}+K_{\lambda+2}K_{2}+K_{\lambda+3}K_{3}+\text{etc.})(e^{\lambda\varphi\sqrt{-1}}+e^{-\lambda\varphi\sqrt{-1}}). \]
Mais le facteur $e^{\lambda\varphi\sqrt{-1}}+e^{-\lambda\varphi\sqrt{-1}}$ se réduit à $2\cos\lambda\varphi$, et on a entre les limites données $\displaystyle\int d\varphi\cdot2\cos^{2}\lambda\varphi=\pi$; donc la valeur générale de $Q^{n}$ est
\[ Q^{n}=\pi[K_{\lambda}+K_{\lambda+1}K_{1}+K_{\lambda+2}K_{2}+K_{\lambda+3}K_{3}+\text{etc.}]. \]
Or on a par la formule du binome,
\[ \begin{aligned} K_{\lambda}&=\frac{n\cdot n-1\cdot n-2\ldots n+1-\lambda}{1\cdot2\cdot3\ldots\lambda}(-a)^{\lambda}\\ K_{\lambda+1}&=K_{\lambda}\cdot\frac{n-\lambda}{\lambda+1}(-a) &\qquad K_{1}&=\frac{n}{1}(-a)\\ K_{\lambda+2}&=K_{\lambda}\cdot\frac{n-\lambda\cdot n-\lambda-1}{\lambda+1\cdot\lambda+2}(-a)^{2} & K_{2}&=\frac{n\cdot n-1}{1\cdot2}(-a)^{2}\\ K_{\lambda+3}&=K_{\lambda}\cdot\frac{n-\lambda\cdot n-\lambda-1\cdot n-\lambda-2}{\lambda+1\cdot\lambda+2\cdot\lambda+3}(-a)^{3} & K_{3}&=\frac{n\cdot n-1\cdot n-2}{1\cdot2\cdot3}(-a)^{3}\\ &\text{etc.} && \text{etc.} \end{aligned} \]
Donc en faisant $Q^{n}$ ou
\[ \int\Delta^{n}d\varphi\cos\lambda\varphi=\pi(-a)^{\lambda}\cdot\frac{n\cdot n-1\cdot n-2\ldots n+1-\lambda}{1\cdot2\cdot3\ldots\lambda}B^{n}, \]

\origpage{376}
on aura pour la valeur du coefficient $B^{n}$,
\[ \begin{aligned} B^{n}=1&+\frac{n}{1}a^{2}\cdot\frac{n-\lambda}{\lambda+1}+\frac{n\cdot n-1}{1\cdot2}a^{4}\cdot\frac{n-\lambda\cdot n-\lambda-1}{\lambda+1\cdot\lambda+2}\\ &+\frac{n\cdot n-1\cdot n-2}{1\cdot2\cdot3}a^{6}\cdot\frac{n-\lambda\cdot n-\lambda-1\cdot n-\lambda-2}{\lambda+1\cdot\lambda+2\cdot\lambda+3}+\text{etc.} \end{aligned} \tag{6} \]

\subsection*{(64).}

Remarquons que dans l'application de cette formule il faut supposer $n=$ ou $>\lambda$. En effet, si $n$ était $<\lambda$, le développement de $\Delta^{n}$ ne pourrait donner aucun terme semblable à $\cos\lambda\varphi$, ainsi l'intégrale $\displaystyle\int\Delta^{n}d\varphi\cos\lambda\varphi$, prise entre les limites données, serait nulle.

La plus petite valeur qu'on puisse donner à $n$ est donc $n=\lambda$, alors on a $\displaystyle\int\Delta^{\lambda}d\varphi\cos\lambda\varphi=\pi(-a)^{\lambda}$, ce qui se vérifie immédiatement.

Le cas de $\lambda=0$ mérite encore d'être remarqué; alors on a
\[ \int\Delta^{n}d\varphi=\pi\left[1+n^{2}a^{2}+\left(\frac{n\cdot n-1}{1\cdot2}\right)^{2}a^{4}+\left(\frac{n\cdot n-1\cdot n-2}{1\cdot2\cdot3}\right)^{2}a^{6}+\text{etc.}\right]. \tag{7} \]
Comparant cette formule à la formule (5), on en tire
\[ \int\Delta^{n}d\varphi=(1-a^{2})^{2n+1}\int\frac{d\varphi}{\Delta^{n+1}}, \tag{8} \]
ce qui établit un rapport remarquable entre ces deux intégrales, rapport qui aurait lieu quand même $n$ ne serait pas un nombre entier, et c'est ce qu'on vérifiera aisément pour le cas de $n=\frac{1}{2}$, par les formules de la première partie.

En général si l'on observe que la valeur du coefficient $B^{n}$ n'est autre chose que celle du coefficient $A^{n}$, dans laquelle on aurait mis $n+1$ au lieu de $n$, on en conclura que les deux intégrales désignées par $P^{n+1}$ et $Q^{n}$, ont entre elles un rapport très-simple; on a en effet d'après les équations (3) et (6), cette formule très-remarquable
\[ \int\Delta^{n}d\varphi\cos\lambda\varphi=\frac{n\cdot n-1\cdot n-2\ldots n-\lambda+1}{n+1\cdot n+2\cdot n+3\ldots n+\lambda}(-1)^{\lambda}(1-a^{2})^{2n+1}\cdot\int\frac{d\varphi\cos\lambda\varphi}{\Delta^{n+1}} \]
et l'équation (8) n'est qu'un cas particulier de celle-ci.

Toutes ces formules sont dues à Euler; mais les démonstrations de cet illustre Auteur sont, si je ne me trompe, beaucoup moins simples que celles que je viens d'exposer. Voyez le tome IV de son \emph{Calcul intégral}, pag.~194 et suiv.

FIN DE LA TROISIÈME PARTIE.

\origpage{377}
% No page number is printed on this first page of the Table (only the signature "48"); 377 by sequence (leaf 384 prints 378).

\section*{TABLE DES MATIÈRES.}

\section*{PREMIÈRE PARTIE.}

\section*{DES FONCTIONS ELLIPTIQUES.}

§ I. \emph{Idée générale des différentes sortes de transcendantes contenues dans la formule intégrale $\displaystyle\int\frac{P\,dx}{R}$}, pag.~4

On démontre que $P$ étant une fonction rationnelle de $x$, et $R$ un radical de la forme $\sqrt{(\alpha+\beta x+\gamma x^{2}+\delta x^{3}+\varepsilon x^{4})}$, l'intégrale $\displaystyle\int\frac{P\,dx}{R}$ ne contient que deux sortes de transcendantes, l'une de la forme $\displaystyle\int(A+Bx+Cx^{2})\frac{dx}{R}$, l'autre de la forme $\displaystyle\int\frac{D}{1+nx}\cdot\frac{dx}{R}$, $A$, $B$, $C$, $D$, $n$ étant des coefficiens constans.

§ II. \emph{Manière de faire disparaître les puissances impaires de la variable sous le radical}, 7

§ III. \emph{Réduction de la différentielle $\frac{P\,dx}{R}$ à la forme $\frac{Q\,d\varphi}{\sqrt{(1-c^{2}\sin^{2}\varphi)}}$}, 9

Les principaux avantages de cette transformation sont, 1°. que le radical $R$, qui dans sa généralité contient quatre coefficiens arbitraires, est remplacé par le radical $\Delta=\sqrt{(1-c^{2}\sin^{2}\varphi)}$, où il n'y en a qu'un; 2°. que la variable $x$, qui pourrait être assujétie à certaines limites, est remplacée par l'amplitude $\varphi$ qui croît indéfiniment; 3°. que l'intégrale, dans sa nouvelle forme, peut être déterminée pour toute valeur de $\varphi$, si elle est connue seulement depuis $\varphi=0$ jusqu'à $\varphi=\frac{1}{2}\pi$.

§ IV. \emph{Développement de la formule $\displaystyle\int\frac{Q\,d\varphi}{\Delta}$}, 12

On démontre que la formule $\displaystyle\int\frac{Q\,d\varphi}{\Delta}$, où $Q$ est une fonction rationnelle paire de $\sin\varphi$, ne contient que deux sortes de transcendantes; l'une de la forme

\origpage{378}
$\displaystyle\int(A+B\sin^{2}\varphi)\frac{d\varphi}{\Delta}$, l'autre de la forme $\displaystyle\int\frac{N}{1+n\sin^{2}\varphi}\cdot\frac{d\varphi}{\Delta}$, $N$ et $n$ ayant des valeurs quelconques réelles ou imaginaires.

§ V. \emph{Définition des fonctions elliptiques, et leur division en trois espèces}, page 14

La plus simple des fonctions elliptiques est représentée par $\displaystyle\int\frac{d\varphi}{\Delta}$, elle constitue la première espèce; l'arc d'ellipse désigné par $\displaystyle\int\Delta\,d\varphi$ est la seconde espèce, et l'intégrale $\displaystyle\int\frac{d\varphi}{(1+n\sin^{2}\varphi)\Delta}$ est la troisième.

§ VI. \emph{Comparaison des fonctions elliptiques de la première espèce}, 20

On donne, d'après Euler, l'intégrale algébrique complète de l'équation $\frac{d\varphi}{\Delta(\varphi)}+\frac{d\psi}{\Delta(\psi)}=0$, cette intégrale répond à l'équation transcendante $F(\varphi)+F(\psi)=F(\mu)$; d'où il suit qu'on peut résoudre algébriquement tous les problèmes relatifs à l'addition, la soustraction, la multiplication et la division des fonctions $F$, comme on le fait à l'égard des arcs de cercle.

§ VII. \emph{Application à la lemniscate}, 37

Les arcs de cette courbe représentent la fonction $F(c,\varphi)$ dans le cas où le module $c=\sin45^{\circ}$.

§ VIII. \emph{Application au mouvement du pendule simple}, 40

Le tems du mouvement par un arc quelconque, soit que le pendule ne fasse que des oscillations, soit qu'il tourne toujours dans le même sens autour du point de suspension, est en général une fonction elliptique de la première espèce, et jouit de toutes les propriétés de cette fonction.

§ IX. \emph{Comparaison des fonctions elliptiques de la seconde espèce}, 41

La même équation algébrique, qui donne $F(\varphi)+F(\psi)-F(\mu)=0$, s'applique aux fonctions de la seconde espèce, et donne $E(\varphi)+E(\psi)-E(\mu)=$ à une quantité algébrique $c^{2}\sin\varphi\sin\psi\sin\mu$. Ainsi toutes les comparaisons des arcs d'ellipse s'établissent sur la même base que celles des fonctions $F$.

§ X. \emph{Comparaison des arcs d'hyperbole}, 52

§ XI. \emph{Développement particulier de la formule $\displaystyle Z=\int\frac{(f+gx^{2})\,dx}{\sqrt{(\alpha^{2}+2\alpha\beta x^{2}\cos\theta+\beta^{2}x^{4})}}$}, 54

L'application de cette formule à diverses intégrales définies, offre différens résultats

\origpage{379}
remarquables sur les relations qui existent entre les fonctions complètes $F^{1}(c)$, $F^{1}(b)$, $E^{1}(c)$, $E^{1}(b)$, dans le cas où $c=\sin15^{\circ}$.

§ XII. \emph{Théorème sur les fonctions complètes de première et de seconde espèce, dont les modules sont complémens l'un de l'autre}, page 61

§ XIII. \emph{Équations différentielles qui expriment la liaison mutuelle des fonctions $E$ et $F$}, 62

On remarque qu'une table des fonctions $E$ servirait à déterminer les fonctions $F$, et réciproquement.

§ XIV. \emph{Développement des fonctions $F^{1}$ et $E^{1}$ en séries}, 65

Ces fonctions dépendent chacune d'une équation différentielle linéaire du second ordre: elles se développent en séries, de deux manières, l'une suivant les puissances du module $c$, l'autre suivant les puissances du module complémentaire $b$.

§ XV. \emph{Des changemens qu'on peut faire subir au paramètre dans les fonctions elliptiques de troisième espèce}, 68

Formule pour transformer une fonction de troisième espèce dont le paramètre est $n$, en une autre dont le paramètre est $\frac{c^{2}}{n}$, 69

On distingue dans les fonctions elliptiques de troisième espèce, trois formes du paramètre; savoir, $n=\cot^{2}\theta$, $n=-1+b^{2}\sin^{2}\theta$, $n=-c^{2}\sin^{2}\theta$, 72

Formule pour transformer l'une dans l'autre les fonctions elliptiques dont les paramètres $n$ et $n'$ satisfont à l'équation $(1+n)(1+n')=b^{2}$, 73

Au moyen de cette formule on peut réduire toute fonction dont le paramètre est de la forme $n=\cot^{2}\theta$, à une fonction dont le paramètre est de la forme $n=-1+b^{2}\sin^{2}\theta$, et réciproquement; mais celles qui sont relatives à la troisième forme $n=-c^{2}\sin^{2}\theta$, ne sont réductibles qu'à des fonctions dont le paramètre est de la même forme, 74

§ XVI. \emph{Comparaison des fonctions elliptiques de la troisième espèce}, 75

La même équation algébrique, qui donne $F(\varphi)+F(\psi)-F(\mu)=0$ et $E(\varphi)+E(\psi)-E(\mu)=c^{2}\sin\varphi\sin\psi\sin\mu$, s'applique aux fonctions de troisième espèce, et donne $\Pi(\varphi)+\Pi(\psi)-\Pi(\mu)=$ à une quantité déterminable par arcs de cercle ou par logarithmes. Ainsi les fonctions de la troisième espèce se comparent comme celles de la première espèce, aux différences près que comporte la nature de ces fonctions.

\origpage{380}

§ XVII. \emph{Formation d'une suite infinie de fonctions elliptiques de la première espèce, liées entre elles par des rapports constans}, pag.~81

Le module d'une de ces fonctions sert à déterminer tous les autres; il en résulte une suite ou échelle de modules qui se prolonge à l'infini dans les deux sens, depuis zéro jusqu'à l'unité. De même, l'amplitude d'une de ces fonctions sert à déterminer toutes les autres qui sont croissantes à l'infini dans un sens, et décroissantes dans l'autre, jusqu'à une limite déterminée.

§ XVIII. \emph{Application de la même loi aux fonctions elliptiques de la seconde espèce}, 85

Le premier résultat de cette application est que la fonction de première espèce $F(c,\varphi)$ peut s'exprimer généralement par deux fonctions de seconde espèce $E(c,\varphi)$, $E(c',\varphi')$; d'où il suit que tout arc d'hyperbole peut être déterminé par deux arcs d'ellipse, ce qui est le théorème de Landen.

On peut former, d'après la même loi, une suite infinie d'ellipses telles qu'en supposant connus les arcs de deux de ces ellipses, on pourra trouver les arcs de toutes les autres.

§ XIX. \emph{Méthode d'approximation appliquée aux fonctions elliptiques de la première espèce}, 89

Deux formules remplissent cet objet, l'une pour le cas où le module $c$ n'est pas trop près de l'unité, l'autre pour le cas où la différence $1-c$ est extrêmement petite.

Les mêmes formules servent à résoudre le problème inverse, c'est-à-dire, à déterminer l'amplitude $\varphi$, quand on connaît la fonction $F(c,\varphi)$.

§. XX. \emph{Propriétés particulières des fonctions $F(c)$, $F(b)$, dont les modules sont complémens l'un de l'autre}, 97

Ces propriétés conduisent à des expressions singulièrement approchées du nombre $\pi$ en quantités logarithmiques.

§ XXI. \emph{Méthode d'approximation appliquée aux fonctions elliptiques de la seconde espèce}, 102

On donne encore sur cet objet deux formules, l'une pour le cas où $c$ n'est pas trop près de l'unité, l'autre pour le cas où la différence $1-c$ est extrêmement petite.

Ces formules sont ce que l'analyse peut offrir de plus simple pour la rectification de l'ellipse et celle de l'hyperbole.

§ XXII. \emph{Formules remarquables pour déterminer les fonctions complètes}

\origpage{381}
\emph{$F^{1}(c)$, $E^{1}(c)$, lorsque $c$ est peu éloigné de l'une de ses limites}, page 113

On a ajouté une table des logarithmes des fonctions $F^{1}(c)$, $E^{1}(c)$, calculés pour tous les angles du module, de degré en degré, depuis zéro jusqu'à $90^{\circ}$, 118

§ XXIII. \emph{Méthode d'approximation appliquée aux fonctions elliptiques de la troisième espèce}, 119

Après avoir déterminé la loi que suivent les paramètres et les coefficiens des transformées successives, on donne l'expression générale de la fonction, pour tous les cas où le module n'est pas trop près de l'unité.

On discute ensuite, d'une manière fort étendue, le cas où le module diffère très-peu de l'unité, et on en donne la solution générale par une suite convergente.

§ XXIV. \emph{Des cas les plus généraux dans lesquels on peut opérer la réduction des fonctions elliptiques de la troisième espèce}, 134

On démontre généralement, 1°. que toute fonction complète de troisième espèce peut s'exprimer par des fonctions de la première et de la seconde espèce.

2°. Qu'il en est de même de toutes les fonctions non complètes en nombre infini, qui peuvent se mesurer par la fonction complète.

3°. Que la fonction de troisième espèce $\Pi(n,c,\varphi)$ peut se réduire indéfiniment à la première espèce, dans une infinité de cas pour lesquels on a un symptôme général.

§ XXV. \emph{Réduction générale des fonctions elliptiques dont le paramètre est imaginaire}, 147

On démontre généralement que toute fonction elliptique de troisième espèce dont le paramètre est imaginaire, peut se réduire à deux fonctions de la même espèce dont les paramètres sont réels, l'un étant de la forme $-1+b^{2}\sin^{2}\theta$, l'autre de la forme $-c^{2}\sin^{2}\theta$; ce qui confirme pleinement la division qui a été faite des fonctions elliptiques en trois espèces.

On discute particulièrement le cas où $n=c(\cos\theta+\sqrt{-1}\sin\theta)$, et celui où $1+n=b(\cos\lambda+\sqrt{-1}\sin\lambda)$.

§ XXVI. \emph{D'un symptôme général pour reconnaître si deux fonctions données de troisième espèce, qui ne diffèrent que par les paramètres, peuvent se réduire l'une à l'autre}, 159

§ XXVII. \emph{Exemple d'une transformation particulière de fonctions elliptiques de la première et de la seconde espèce}, 161

Cet exemple se rapporte aux fonctions dont les modules sont $\sin75^{\circ}$ et $\sin15^{\circ}$, et sont par conséquent complémens l'un de l'autre.

\origpage{382}

§ XXVIII. \emph{Des séries qui donnent, sans transformations, les valeurs approchées des fonctions elliptiques}, page 169

Ces séries procèdent suivant les sinus des arcs multiples de l'amplitude, et leurs coefficiens peuvent être déterminés généralement par les formules relatives aux fonctions complètes de la première et de la seconde espèce.

§ XXIX. \emph{Surface du cône oblique}, 173

On démontre que la surface totale est toujours déterminable par des arcs d'ellipse, soit que la base soit un cercle, soit qu'elle soit une ellipse.

§ XXX. \emph{Construction de la ligne la plus courte sur la surface du sphéroïde}, 178

La latitude d'un point quelconque de cette courbe étant donnée, on trouve la longueur d'un arc de la courbe par celle d'un arc égal d'ellipse; quant à la longitude, elle dépend d'une fonction elliptique de troisième espèce, et d'une de première espèce: on peut aussi la construire par le développement d'un cône oblique à base circulaire. Les points où la courbe rencontre l'équateur et ceux où elle touche les parallèles entre lesquels elle est contenue, se déterminent par les seuls arcs d'ellipse.

§ XXXI. \emph{Détermination de l'aire de l'ellipsoïde}, 182

On donne d'abord la valeur approchée de cette aire par une série régulière et convergente.

On détermine ensuite la valeur exacte de l'aire en fonctions elliptiques, par la considération des lignes de plus grande et de plus petite courbure tracées sur la surface de l'ellipsoïde, et on trouve que l'aire totale peut s'exprimer par des arcs d'ellipse.

§ XXXII. \emph{De quelques formules intégrales qui peuvent se ramener aux fonctions elliptiques}, 194

§ XXXIII. \emph{De l'intégrale $\displaystyle V=\int\frac{(x^{2}+m)\,dx}{\sqrt{(\alpha+\beta x+\gamma x^{2}+\varepsilon x^{4})}}$}, 200

On prouve que cette intégrale ne dépend que des fonctions elliptiques de la première et de la seconde espèce.

§ XXXIV. \emph{De l'intégrale $\displaystyle Z=\int\frac{dz}{(1+\mu z^{2})\sqrt[3]{(1+\nu z^{2})}}$}, 201

On examine successivement le cas de $\nu=3$ et $\mu=\frac{1}{3}$, celui de $\nu=-3$ et $\mu=-\frac{1}{3}$, et enfin celui de $\nu=1$ et $\mu=-\frac{1}{8}$,\ednote{The denominator of this fraction is heavily inked; read as $8$.} tous trois susceptibles d'être résolus par les arcs de cercle et les logarithmes.

§ XXXV. \emph{Des cas principaux où l'on peut évaluer, par les fonctions}

\origpage{383}
\emph{elliptiques, l'intégrale $\left(\frac{p}{q}\right)=\displaystyle\int x^{p-1}\,dx(1-x^{n})^{\frac{q}{n}-1}$, prise depuis $x=0$ jusqu'à $x=1$}, page 209

Ces cas sont ceux où $n$ est égal à l'un des nombres 3, 4, 6, 8 et 12. Dans le cas de $n=12$, on parvient à des résultats très-remarquables sur la comparaison des fonctions complètes $F^{1}(c)$, $F^{1}(b)$, tant entre elles qu'avec la fonction $F^{1}(\sin 45^{\circ})$, l'angle $\theta$ du module $c$ étant tel qu'on a $\sin 2\theta=\operatorname{tang}^{2}15^{\circ}$.

\section*{SECONDE PARTIE.}

\section*{DES INTÉGRALES EULÉRIENNES.}

§ I. \emph{Des intégrales Eulériennes de la première espèce}, 222

Formule générale qui comprend presque toute la théorie de ces fonctions, et qui peut être regardée comme une sorte d'équation aux différences finies partielles, 227

Au moyen des auxiliaires désignées par $A_{a}$, on donne l'expression générale des fonctions $\left(\frac{p}{q}\right)$ sous deux formes différentes, selon que $p+q$ est $>$ ou $<n$, 230

On fait voir que si $n$ est pair, le nombre des auxiliaires peut être réduit à moitié par la formule $(d')$, 237

§ II. \emph{Formule pour évaluer par approximation les intégrales $\left(\frac{p}{q}\right)$}, 240

On cherche particulièrement la valeur approchée de cette intégrale lorsque $p$ et $q$ sont très-petits par rapport à $n$.

§ III. \emph{Remarque sur quelques cas particuliers où l'on peut sommer la suite désignée par $\psi(x)$, et deux autres de la même espèce}, 244

On donne sur ces suites plusieurs théorèmes remarquables qui sont dus à \emph{Landen}, et on ajoute les démonstrations de quelques-uns d'eux, que l'auteur a supprimées.

§ IV. \emph{Considération des formules intégrales $\displaystyle\int\frac{x^{p-1}\,dx}{\sqrt[n]{(1-x^{n})^{n-q}}}\log\frac{1}{x}$, $\displaystyle\int\frac{x^{p-1}\,dx}{\sqrt[n]{(1-x^{n})^{n-q}}}\log^{2}\frac{1}{x}$, etc. Théorème très-remarquable sur la première de ces formules}, 249

Ce théorème sert à déterminer exactement la première de ces intégrales, en

\origpage{384}
supposant connue l'intégrale $\left(\frac{p}{q}\right)$; quant aux formules suivantes, leur détermination dépend en général de la transcendante $(a,n)^{m}$, c'est-à-dire, de la somme des termes pris de $n$ en $n$, dans la suite $1+\frac{1}{2^{m}}+\frac{1}{3^{m}}+\frac{1}{4^{m}}+\text{etc.}$

§. V. \emph{De la réduction des transcendantes désignées par $(a,n)^{m}$}, p.~265

On traite particulièrement le cas de $m=2$, $n=6$, et celui de $m=3$, $n=6$; ils dépendent de deux transcendantes dont on cherche l'expression en séries convergentes.

§ VI. \emph{Des intégrales Eulériennes de la seconde espèce}, 276

Toute intégrale de la première espèce $\left(\frac{p}{q}\right)$ peut s'exprimer très-simplement par les fonctions $\Gamma$; réciproquement toute fonction $\Gamma(a)$ dans laquelle $a$ est un nombre rationnel, peut s'exprimer par les intégrales $\left(\frac{p}{q}\right)$, 283

Expression de la fonction $\Gamma(k)$ lorsque $k$ est très-petit, 289

On prouve que la fonction $\Gamma(a)$ sera connue pour toute valeur de $a$, si on connaît seulement cette fonction depuis $a=\frac{3}{4}$ jusqu'à $a=1$, \emph{ibid.}

Formule pour calculer directement la valeur de chaque fonction $\Gamma$, 290

Cette formule est du nombre des suites \emph{demi-convergentes}; on cherche, \emph{a priori}, le point où il faut s'arrêter dans le calcul de cette suite, et on fixe le degré d'approximation qu'elle peut donner dans chaque cas, 292

Considérations générales d'où l'on déduit des suites convergentes et régulières, pour exprimer les valeurs de $\log\Gamma(1+x)$, $\log\Gamma(1-x)$, etc., 298

Application de ces suites à l'expression générale du logarithme de la fonction $\left(\frac{p}{q}\right)$, 300

L'intégrale $\displaystyle\int t^{n}e^{-t^{m}}\,dt$, prise depuis $t=0$ jusqu'à $t=\infty$, se ramène immédiatement aux fonctions $\Gamma$, 301

Table des logarithmes de la fonction $\Gamma(a)$ pour toute valeur de $a$, de millième en millième, depuis $a=1.000$ jusqu'à $a=2.000$, 302

\section*{TROISIÈME PARTIE.}

\section*{DES QUADRATURES.}

§ I. \emph{Formule générale pour les quadratures}, 308

Cette formule se compose d'une intégrale aux différences finies et d'une suite de corrections dont la loi est connue, 311

\origpage{385}
Moyens d'obvier à l'inconvénient que présente la formule, lorsque quelqu'un des coefficiens $\frac{dy}{dx}$, $\frac{ddy}{dx^{2}}$, $\frac{d^{3}y}{dx^{3}}$ devient infini dans la partie de courbe qu'on veut quarrer, page 313

Le résultat de la formule générale se rapproche beaucoup de celui qu'on obtient en considérant l'aire de la courbe proposée comme une somme d'aires paraboliques, déterminées chacune par trois ordonnées consécutives et équidistantes, 319

§ II. \emph{Construction de la courbe dans laquelle l'arc $s$ est donné en fonction de la quantité $\frac{dy}{dx}$}, 320

On trouve généralement l'expression de chacune des coordonnées par une intégrale aux différences finies, jointe à une suite de corrections dont la loi est connue.

§ III. \emph{Application de la méthode précédente au calcul de la trajectoire d'un projectile}, 330

On donne les détails du calcul pour trouver, dans un exemple particulier, la hauteur du jet, l'amplitude de la branche ascendante et celle de la branche descendante. On détermine ensuite la position de l'asymptote verticale et celle de l'asymptote inclinée.

§ IV. \emph{De l'intégrale indéfinie $\displaystyle\int dx\left(\log\frac{1}{x}\right)^{a-1}$, prise depuis $x=0$}, 339

On donne deux formules pour évaluer cette intégrale par approximation, l'une pour tous les cas, l'autre pour celui seulement où $x$ est très-petit. On donne aussi une formule pour trouver l'intégrale lorsque $x$ est plus grand que l'unité.

§ V. \emph{De l'intégrale $\displaystyle\int y\,dx$ prise entre deux limites qui rendent nulle la fonction $y$}, 343

Application à l'intégrale $\displaystyle\int x^{\alpha}e^{-x}\,dx$, prise depuis $x=0$ jusqu'à $x=\infty$, $\alpha$ étant un nombre positif. Il en résulte la même formule qui a été trouvée dans la seconde partie pour la valeur de $\Gamma(\alpha+1)$, 346

Application à l'intégrale $\displaystyle\int x^{\alpha}(1-x)^{\beta}\,dx$, prise entre les limites $x=0$, $x=1$, $\alpha$ et $\beta$ étant positifs, 348

Cette intégrale comprend la détermination générale des fonctions $\left(\frac{p}{q}\right)$, 351

§ VI. \emph{De l'intégrale $\displaystyle\int x^{-\alpha}e^{-x}\,dx$, prise entre les limites imaginaires qui rendent nulle $x^{-\alpha}e^{-x}$}, 352\ednote{In this copy the last digit is overwritten by a reader's ink; read as $352$, since the section begins on p.~352 (reviewer's reading from context).}

La valeur de cette intégrale désignée par $\Gamma'(a)$, se trouve généralement en supposant connue la fonction $\Gamma(a)$. 353

\origpage{386}
On en déduit l'expression générale de l'intégrale $Q(\alpha)=\displaystyle\int d\varphi\cos^{\alpha-2}\varphi\ldots\ \cos(\alpha\operatorname{tang}\varphi-\alpha\varphi)$, prise depuis $\varphi=0$ jusqu'à $\varphi=\frac{1}{2}\pi$. Il s'ensuit que cette intégrale peut être déterminée exactement en fonction de $\pi$ et $e$, toutes les fois que $2\alpha$ est un nombre entier, page 355

§ VII. \emph{De l'intégrale $Z=\displaystyle\int\frac{dx\cos ax}{1+xx}$ et autres semblables, prises depuis $x=0$ jusqu'à $x=\infty$.} 357

On déduit de la formule principale deux séries d'intégrales qui peuvent toutes être déterminées en fonction des nombres $e$ et $\pi$. En général on peut trouver la valeur exacte de l'intégrale $\displaystyle\int\left(\frac{M\cos ax+Nx\sin ax}{P}\right)dx$, depuis $x=0$ jusqu'à $x=\infty$, si $M$, $N$, $P$ sont des fonctions rationnelles et entières de $x^{2}$, $P$ étant la plus élevée, et si en outre $P$ n'a aucun facteur de la forme $x^{2}-m^{2}$, 360

§ VIII. \emph{De l'intégrale $Z=\displaystyle\int e^{-x^{2}}\,dx\cos ax$, prise depuis $x=0$ jusqu'à $x=\infty$.} 362

Autres intégrales qu'on déduit de la même formule, 363

§ IX. \emph{De l'intégrale $Z(k)=\displaystyle\int x^{\frac{2k-1}{2}}\,dx\,e^{-\left(\frac{1+xx}{2nx}\right)}$, prise depuis $x=0$ jusqu'a $x=\infty$}, 364

§ X. \emph{Des intégrales $\displaystyle\int x^{a-1}e^{-mx}\,dx\cos nx$, $\displaystyle\int x^{a-1}e^{-mx}\,dx\sin nx$, prises entre les limites $x=0$, $x=\infty$}, 367

§ XI. \emph{De l'intégrale $\displaystyle\int\frac{(x^{n}-1)\,dx}{\log x}$ et autres semblables, prises depuis $x=0$ jusqu'à $x=1$}, 370

Loi très-simple que suivent les intégrales $\displaystyle\int\left(\frac{x^{n}-1}{\log x}\right)^{k}x^{m-1}\,dx$, en donnant à $k$ les valeurs successives 1; 2, 3, etc.

§ XII. \emph{Des intégrales $\displaystyle\int\frac{d\varphi\cos\lambda\varphi}{\Delta^{n}}$ et $\displaystyle\int\Delta^{n}\,d\varphi\cos\lambda\varphi$, prises depuis $\varphi=0$ jusqu'à $\varphi=\pi$, en supposant les nombres $\lambda$ et $n$ entiers, et $\Delta=1+a^{2}-2a\cos\varphi$}, 372

FIN DE LA TABLE.

\origpage{387}
% No page number is printed on the Errata page; 387 by sequence (the Table ends on printed p. 386, leaf 392).

\section*{ERRATA.}

\[
\begin{aligned}
&\text{Pages.} && \text{Lignes.} && \text{Fautes.} && \text{Corrections.}\\
&26 && 2 && \cot\tfrac{1}{2}\beta\ldots\ldots && \cos\tfrac{1}{2}\beta.\\
&38 && 18 && \text{retranchés}\ldots\ldots && \text{retranchés, multipliés.}\\
&61 && \text{titre} && \textit{fonctions}\ldots\ldots && \textit{fonctions complètes.}\\
&62 && \text{fin de la pag.} && F-c\frac{dE}{dc}\ldots\ldots && E-c\frac{dE}{dc}.\\
& && && b^{2}\left(E+c\frac{dF}{dc}\right)\ldots && b^{2}\left(F+c\frac{dF}{dc}\right).\\
&66 && 2 && -\frac{1^{2}}{2^{2}}c^{2}\ldots\ldots && +\frac{1^{2}}{2^{2}}c^{2}.\\
&\textit{Ibid.} && 3 && -\frac{1^{2}}{2^{2}\cdot4^{2}}3c^{2}\ldots\ldots && -\frac{1^{2}}{2^{2}\cdot4^{2}}3c^{4}.\\
&85 && 15 && (2+c)E(c',\varphi')\ldots && (2+2c)E(c',\varphi').\\
&104 && \text{avant-dern.} && \tfrac{1}{2}B''\sin\varphi''\ldots\ldots && \tfrac{1}{2}B''\sin\varphi'.\\
&156 && 15 && \text{art.\ 58}\ldots\ldots && \text{art.\ 60.}\\
&179 && 15 && CA=A,\ CB=B. && CE=A,\ CP=B.\\
&194 && \text{13 et 15} && am\beta\ldots am'\beta'\ldots && amb\ldots am'b'.\\
&287 && 14 && a=10\ldots\ldots && n=10.\\
&364 && 14 && \frac{1+2z}{2nz}\ldots\ldots && \frac{1+zz}{2nz}.
\end{aligned}
\]

\end{document}
